% WFDB-x: User's manual.
% Built with pdflatex (latexmk; see README.md). The source is ASCII: special characters are
% written as LaTeX commands. Every change to WFDB-x is reflected here (see
% docs/porting-guide.md), and check_manual.py checks this manual against the code.
\documentclass[11pt,a4paper,oneside,openany]{book}
% Preamble of the WFDB-x user's manual (wfdb-x-user-manual.tex).
% Built with pdflatex (latexmk; see README.md). The source is ASCII: special characters are
% written as LaTeX commands.

% ---- Engine, encoding and fonts ----------------------------------------------------------
\usepackage{iftex}
\ifPDFTeX
  \usepackage[T1]{fontenc}
  \usepackage[utf8]{inputenc}
\else
  \usepackage{fontspec}
\fi
\usepackage{XCharter}                               % text: Charter
\IfFileExists{sourcesans.sty}                       % headings, labels: Source Sans (\sffamily)
  {\usepackage{sourcesans}}{\usepackage{sourcesanspro}}  % (older TeX distributions)
\usepackage[varqu,varl,scaled=0.9]{zi4}             % code: Inconsolata, upright quotes, bold
\usepackage{textcomp}
\usepackage{amsmath,amssymb}
\usepackage{microtype}
\DisableLigatures[-,<,>]{family=tt*}           % "--", "<<", ">>" stay as typed in code

% ---- Colours (colour-blind-safe; red is never used: it is for alarms) --------------------
\usepackage[table]{xcolor}
\definecolor{titlecolor}{HTML}{12324A}               % headings, title page
\definecolor{accent}{HTML}{0072B2}                   % Okabe-Ito blue: numbers, links, rules
\definecolor{okskyblue}{HTML}{56B4E9}
\definecolor{okgreen}{HTML}{009E73}
\definecolor{okorange}{HTML}{E69F00}
\definecolor{okviolet}{HTML}{7B3294}                 % identifying data (PHI)
\definecolor{rulecolor}{HTML}{B8C4CE}
\definecolor{mutedtext}{HTML}{5A6B7B}
\definecolor{codebg}{HTML}{F5F7F9}
\definecolor{codecomment}{HTML}{5A6B7B}
\definecolor{codekeyword}{HTML}{005A8C}
\definecolor{codestring}{HTML}{00785A}
\definecolor{noteback}{HTML}{EBF4FA}
\definecolor{phiback}{HTML}{F4EDF8}
\definecolor{tipback}{HTML}{E9F6F1}
\definecolor{tableheadback}{HTML}{EEF2F5}

% ---- Page --------------------------------------------------------------------------------
\usepackage[a4paper,top=28mm,bottom=27mm,inner=27mm,outer=25mm,headheight=15pt,
  headsep=9mm,footskip=11mm]{geometry}
\usepackage[skip=0.55em plus 0.2em minus 0.1em]{parskip}
\raggedbottom
\clubpenalty=10000
\widowpenalty=10000
\setcounter{tocdepth}{2}
\setcounter{secnumdepth}{3}

% ---- Headings and the table of contents --------------------------------------------------
\usepackage{titlesec}
\usepackage{titletoc}
\titleformat{\part}[display]
  {\normalfont\sffamily\centering\color{titlecolor}}
  {\large\bfseries\color{accent}\MakeUppercase{\partname}\enspace\thepart}
  {14pt}
  {\Huge\bfseries}
\titleformat{\chapter}[display]
  {\normalfont\sffamily\color{titlecolor}}
  {\large\bfseries\color{accent}\MakeUppercase{\chaptertitlename}\enspace\thechapter}
  {4pt}
  {{\color{accent}\titlerule[1.2pt]}\vspace{10pt}\Huge\bfseries}
\titleformat{name=\chapter,numberless}[display]
  {\normalfont\sffamily\color{titlecolor}}
  {}
  {0pt}
  {{\color{accent}\titlerule[1.2pt]}\vspace{10pt}\Huge\bfseries}
\titlespacing*{\chapter}{0pt}{0pt}{26pt}
\titleformat{\section}
  {\normalfont\sffamily\Large\bfseries\color{titlecolor}}{\color{accent}\thesection}{0.75em}{}
\titleformat{\subsection}
  {\normalfont\sffamily\large\bfseries\color{titlecolor}}{\color{accent}\thesubsection}{0.75em}{}
\titleformat{\subsubsection}
  {\normalfont\sffamily\normalsize\bfseries\color{titlecolor}}{\thesubsubsection}{0.75em}{}
\titleformat{\paragraph}[runin]
  {\normalfont\sffamily\normalsize\bfseries\color{titlecolor}}{}{0pt}{}
\titlespacing*{\section}{0pt}{1.6em plus 0.4em minus 0.2em}{0.5em}
\titlespacing*{\subsection}{0pt}{1.3em plus 0.3em minus 0.2em}{0.4em}
\titlespacing*{\subsubsection}{0pt}{1.1em plus 0.3em minus 0.2em}{0.3em}
\titlespacing*{\paragraph}{0pt}{0.8em plus 0.2em}{0.6em}

\titlecontents{part}[0pt]
  {\addvspace{1.3em}\sffamily\bfseries\large\color{titlecolor}}
  {\partname\ \thecontentslabel\quad}{}{\hfill\contentspage}[\addvspace{0.3em}]
\titlecontents{chapter}[2.2em]
  {\addvspace{0.55em}\sffamily\bfseries}
  {\contentslabel{2.2em}}{\hspace*{-2.2em}}{\hfill\contentspage}
\titlecontents{section}[4.9em]
  {\small}{\contentslabel{2.7em}}{\hspace*{-2.7em}}
  {\titlerule*[0.7pc]{.}\contentspage}
\titlecontents{subsection}[8.4em]
  {\small\color{mutedtext}}{\contentslabel{3.5em}}{\hspace*{-3.5em}}
  {\titlerule*[0.7pc]{.}\contentspage}
\titlecontents{figure}[2.6em]
  {\small}{\contentslabel{2.6em}}{\hspace*{-2.6em}}{\titlerule*[0.7pc]{.}\contentspage}

% ---- Running heads -----------------------------------------------------------------------
\usepackage{fancyhdr}
\pagestyle{fancy}
\fancyhf{}
\fancyhead[L]{\sffamily\small\color{mutedtext}\nouppercase{\leftmark}}
\fancyhead[R]{\sffamily\small\bfseries\color{titlecolor}\thepage}
\fancyfoot[L]{\sffamily\footnotesize\color{mutedtext}WFDB-x\enspace\textbar\enspace
  User's manual}
\fancyfoot[R]{\sffamily\footnotesize\color{mutedtext}Version \wfdbxversion}
\renewcommand{\headrule}{{\color{rulecolor}\hrule height 0.4pt width\headwidth}}
\renewcommand{\footrulewidth}{0pt}
\renewcommand{\chaptermark}[1]{\markboth{%
  \ifnum\value{secnumdepth}>-1 \ifnum\value{chapter}>0 \thechapter\quad\fi\fi #1}{}}
\fancypagestyle{plain}{\fancyhf{}%
  \fancyfoot[L]{\sffamily\footnotesize\color{mutedtext}WFDB-x\enspace\textbar\enspace
    User's manual}%
  \fancyfoot[R]{\sffamily\footnotesize\bfseries\color{titlecolor}\thepage}%
  \renewcommand{\headrulewidth}{0pt}}

% ---- Lists, figures, tables --------------------------------------------------------------
\usepackage{enumitem}
\setlist{itemsep=0.15em,topsep=0.3em,parsep=0.1em}
\setlist[itemize,1]{label={\color{accent}\textbullet}}
\usepackage{graphicx}
\usepackage{caption}
\captionsetup{font=small,labelfont={bf,sf,color=accent},textfont=sf,labelsep=quad,
  skip=6pt}
\captionsetup[table]{position=top}
\usepackage{float}
\usepackage{booktabs}
\usepackage{array}
\usepackage{tabularx}
\usepackage{longtable}
\usepackage{xltabular}
\usepackage{multirow}
\usepackage{etoolbox}
\newcolumntype{L}{>{\raggedright\arraybackslash}X}       % ragged-right X column
\newcolumntype{P}[1]{>{\raggedright\arraybackslash}p{#1}} % ragged-right fixed-width column
\renewcommand{\arraystretch}{1.2}
\setlength{\aboverulesep}{0.5ex}
\setlength{\belowrulesep}{0.6ex}
\setlength{\heavyrulewidth}{0.9pt}
\setlength{\lightrulewidth}{0.45pt}
\AtBeginEnvironment{tabularx}{\small}
\AtBeginEnvironment{xltabular}{\small}
\newcommand{\tabhead}{\leavevmode\sffamily\bfseries\color{titlecolor}}   % the header cells of a table

% ---- Code --------------------------------------------------------------------------------
\usepackage{listings}
\lstdefinelanguage{Rust}{
  morekeywords={as,break,const,continue,crate,else,enum,extern,false,fn,for,if,impl,in,let,
    loop,match,mod,move,mut,pub,ref,return,self,Self,static,struct,super,trait,true,type,
    unsafe,use,where,while,dyn,Some,None,Ok,Err,Option,Result,Vec,String,u8,u16,u32,u64,
    i16,i32,i64,f32,f64,usize,isize,bool,str},
  sensitive=true,
  morecomment=[l]{//},
  morecomment=[s]{/*}{*/},
  morestring=[b]",
}
\lstdefinelanguage{TOML}{
  morecomment=[l]{\#},
  morestring=[b]",
  morestring=[b]',
  morekeywords={true,false},
}
\lstdefinelanguage{CMake}{
  morekeywords={add_subdirectory,target_link_libraries,set,cmake_minimum_required,project,
    add_executable,find_package,PRIVATE,PUBLIC,CACHE,BOOL,FORCE,OFF,ON,EXCLUDE_FROM_ALL},
  sensitive=false,
  morecomment=[l]{\#},
  morestring=[b]",
}
\lstdefinelanguage{Julia}{
  morekeywords={abstract,baremodule,begin,break,catch,const,continue,do,else,elseif,end,
    export,false,finally,for,function,global,if,import,let,local,macro,module,mutable,
    nothing,quote,return,struct,true,try,using,while,where,missing},
  sensitive=true,
  morecomment=[l]\#,
  morecomment=[n]{\#=}{=\#},
  morestring=[b]",
}
\lstdefinelanguage{Kotlin}{
  morekeywords={as,break,class,continue,do,else,false,for,fun,if,in,interface,is,null,
    object,package,return,super,this,throw,true,try,typealias,val,var,when,while,by,
    catch,constructor,data,enum,finally,import,internal,private,public,vararg,
    Int,Long,Double,String,Boolean,List,Map,Pair},
  sensitive=true,
  morecomment=[l]{//},
  morecomment=[s]{/*}{*/},
  morestring=[b]",
}
\lstdefinelanguage{JSON}{
  morestring=[b]",
  morekeywords={true,false,null},
}
\lstset{
  basicstyle=\ttfamily\small,
  columns=fullflexible,
  keepspaces=true,
  upquote=true,
  showstringspaces=false,
  breaklines=true,
  breakatwhitespace=false,
  postbreak=\mbox{\textcolor{codecomment}{$\hookrightarrow$}\space},
  backgroundcolor=\color{codebg},
  frame=l,
  framerule=2pt,
  rulecolor=\color{okskyblue},
  framesep=6pt,
  xleftmargin=8pt,
  xrightmargin=0pt,
  aboveskip=0.8em,
  belowskip=0.6em,
  tabsize=4,
  keywordstyle=\color{codekeyword}\bfseries,
  commentstyle=\color{codecomment}\itshape,
  stringstyle=\color{codestring},
}
\lstdefinestyle{text}{language={},keywordstyle={},commentstyle={},stringstyle={}}
% The styles of the listings of the reference chapters. docs/user-manual/check_manual.py
% finds them by these names: "signature" is the call syntax of a function (one listing
% right after each \apientry), "example" is an example that the check runs (Python), and
% "shell" is a command line.
\lstdefinestyle{signature}{backgroundcolor=\color{white},rulecolor=\color{accent},
  basicstyle=\ttfamily\small\bfseries,keywordstyle=\color{codekeyword}}
\lstdefinestyle{example}{}
\lstdefinestyle{shell}{language=bash,keywordstyle={},commentstyle=\color{codecomment},
  stringstyle={}}

% ---- Boxes -------------------------------------------------------------------------------
\usepackage{tikz}
\usepackage[most]{tcolorbox}
\tcbset{wfdbxbox/.style={enhanced,breakable,boxrule=0pt,leftrule=3pt,arc=0pt,outer arc=0pt,
  left=8pt,right=8pt,top=5pt,bottom=5pt,before skip=0.8em,after skip=0.8em,
  fonttitle=\sffamily\bfseries\small,attach title to upper={\quad}}}
\newtcolorbox{note}[1][Note]{wfdbxbox,colback=noteback,colframe=accent,coltitle=accent,
  title={#1}}
\newtcolorbox{tip}[1][Tip]{wfdbxbox,colback=tipback,colframe=okgreen,coltitle=okgreen,
  title={#1}}
\newtcolorbox{phibox}[1][Patient data]{wfdbxbox,colback=phiback,colframe=okviolet,
  coltitle=okviolet,title={#1}}

% ---- Links and references ----------------------------------------------------------------
\usepackage{xurl}
\usepackage[bookmarksnumbered,bookmarksopen,bookmarksopenlevel=0]{hyperref}
\hypersetup{colorlinks=true,linkcolor=accent,urlcolor=accent,citecolor=accent,
  pdftitle={WFDB-x: User's manual},
  pdfauthor={Reza Sameni},
  pdfsubject={The WFDB Software Package in Python, C++, Rust, Julia, R, MATLAB/Octave
    and Kotlin},
  pdfkeywords={WFDB, PhysioNet, physiological signals, ECG, annotations, Python, C++,
    Rust, Julia, R, MATLAB, Octave, Kotlin}}
\usepackage[nameinlink,capitalise,noabbrev]{cleveref}

% ---- Manual macros -----------------------------------------------------------------------
\newcommand{\wfdbxversion}{1.2.1}      % the version of WFDB-x described
\newcommand{\wfdbxdate}{November 2026}
\newcommand{\wfdbversion}{10.7.0}     % the version of the WFDB Software Package implemented
\newcommand{\wfdbx}{WFDB-x}

% An underscore that allows a line break after it, so long identifiers wrap in tables.
\DeclareRobustCommand{\_}{\ifmmode\text{\textunderscore}\else\textunderscore\allowbreak\fi}

\newcommand{\code}[1]{\texttt{#1}}          % inline code, identifiers, values (escape _ # % & { })
\newcommand{\cmd}[1]{\texttt{#1}}           % a command line or a subcommand
\newcommand{\opt}[1]{\texttt{#1}}           % a command-line option
% A file or folder of the repository or an output: may break after each /.
\makeatletter
\DeclareRobustCommand{\file}[1]{\texttt{\wfdbx@file#1/\@nil}}
\def\wfdbx@file#1/#2\@nil{#1\if\relax\detokenize{#2}\relax\else/\allowbreak\wfdbx@file#2\@nil\fi}
\makeatother
\newcommand{\phimark}{\textsuperscript{\textcolor{okviolet}{\sffamily\bfseries PHI}}}
\pdfstringdefDisableCommands{\def\_{\textunderscore}\def\allowbreak{}\def\code#1{#1}%
  \def\file#1{#1}\def\cmd#1{#1}\def\opt#1{#1}\def\phimark{}}

% A function or command entry of a reference chapter: \apientry{name}{one-line summary}
\newcommand{\apientry}[2]{\par\addvspace{0.9em}\noindent
  {\color{rulecolor}\rule{\linewidth}{0.4pt}}\par\nopagebreak\vspace{0.15em}\noindent
  {\large\ttfamily\bfseries\color{titlecolor}#1}\hfill{\small\sffamily\color{mutedtext}#2}\par
  \nopagebreak\smallskip}
% Parameter lists of reference entries: \begin{params} \item[\code{name}] meaning \end{params}
\newlist{params}{description}{1}
\setlist[params]{font=\normalfont\ttfamily\bfseries,leftmargin=2.2em,style=nextline,
  itemsep=0.1em}

% Return value and errors of a reference entry: \returns{...} \raises{...}
\newcommand{\returns}[1]{\par\noindent{\sffamily\bfseries\small\color{titlecolor}Returns}
  \quad #1\par}
\newcommand{\raises}[1]{\par\noindent{\sffamily\bfseries\small\color{titlecolor}Errors}
  \quad #1\par}
% The language of a chapter, for the running heads and the check (\chapterlanguage{python})
\newcommand{\chapterlanguage}[1]{}



\begin{document}

% ---- Title page --------------------------------------------------------------------------
\begin{titlepage}
  \thispagestyle{empty}
  \vspace*{0.18\textheight}
  {\sffamily\color{accent}\large\bfseries USER'S MANUAL\par}
  \vspace{0.6em}
  {\color{accent}\rule{\linewidth}{1.2pt}\par}
  \vspace{1.0em}
  {\sffamily\bfseries\color{titlecolor}\fontsize{40}{44}\selectfont WFDB-x\par}
  \vspace{0.8em}
  {\sffamily\Large\color{titlecolor}The WFDB Software Package in Python, C++, Rust, Julia,
    R, MATLAB/Octave and Kotlin\par}
  \vspace{2.2em}
  {\sffamily\large\color{mutedtext}Version \wfdbxversion\quad\textbar\quad \wfdbxdate\par}
  \vfill
  {\sffamily\large\color{titlecolor}Reza Sameni\par}
  \vspace{0.3em}
  {\sffamily\small\color{mutedtext}\url{https://github.com/alphanumericslab/wfdb-x}\par}
  \vspace{1.2em}
  {\small\color{mutedtext}
  WFDB-x implements version \wfdbversion{} of the WFDB Software Package. Its library is
  distributed under the GNU Lesser General Public License, version 2.1 or later,
  and its command-line tools under the GNU General Public License, version 3 or
  later (\cref{sec:license}). To cite WFDB-x, see
  \cref{sec:cite}.\par}
\end{titlepage}

\frontmatter
\tableofcontents

\mainmatter

% ---- Part I: Concepts --------------------------------------------------------------------
\part{Concepts}
\label{part:concepts}
% Part I, chapter 1: Introduction.
\chapter{Introduction}
\label{chap:introduction}

\wfdbx{} is a reimplementation of the WFDB Software Package, version
\wfdbversion{}: the C library and the command-line tools with which
physiologic signal and annotation files in the WFDB format, such as those of
the PhysioNet databases, are read, written and converted. \wfdbx{}
reimplements the package in seven programming languages: Python, C++, Rust,
Julia, R, MATLAB/Octave and Kotlin. In each of them it reads and writes the
same files as the original C library, provides the same functions under
their C names, and provides the same command-line tools, with the same
options, output, messages and exit status. The seven implementations are
tested to give identical results, byte for byte, to the original C programs
and to each other.

\wfdbx{} is an independent project. It is not part of the WFDB project
(\url{https://wfdb.io/}) or of PhysioNet (\url{https://physionet.org/}), and
is not endorsed by them. It follows the WFDB format and the behavior of the
original C package, which serves as its specification: the expected results
of all its tests were produced by the C programs.

\section{What WFDB-x is}
\label{sec:what}

Each of the seven implementations
\begin{itemize}
\item reads and writes exactly the same files as the original C library:
  headers (\file{.hea}), signal files in every WFDB storage format (8, 16,
  61, 80, 160, 212, 310, 311, 24, 32, and the FLAC formats 508, 516 and
  524), annotation files (in MIT and AHA format), EDF and EDF+ files,
  calibration files, info files, and multi-segment records (see
  \cref{chap:records});
\item provides the same 33 command-line tools as the C package (listed in
  \cref{sec:three-ways}), with the same options;
\item provides the WFDB library functions under their C names
  (\code{isigopen}, \code{getvec}, \code{getann}, \code{putann},
  \code{strtim}, \ldots), with the behavior of the C functions, and a
  high-level API in the idiom of the language;
\item is written natively in its language: none of the implementations calls
  the C library or another language. It needs nothing beyond the standard
  library of its language, except NumPy for Python. FLAC files are read and
  written without libFLAC. Records are read from URLs with the system
  \cmd{curl} program in C++ and Rust (and in R and MATLAB/Octave when
  \cmd{curl} is installed), and with the standard facilities of the
  language otherwise (see \cref{sec:network}).
\end{itemize}

\wfdbx{} behaves exactly as the original package, with a few deliberate
exceptions, which are the same in all seven languages and are tested like
everything else: it never assumes a sampling frequency, gain, ADC resolution
or units that a record does not give; it reads and writes files only where
they are named, and opens URLs only when network access is turned on; it
sorts annotation files itself, without losing annotations; a tool never
writes a signal file over a signal file it is reading; and a tool that cannot
write its output exits with status 1. \Cref{chap:behavior} describes each of
them, and \emph{legacy mode}, which restores the original behavior exactly.

The interactive viewer of the C package (WAVE), its plotting programs, and
the analysis programs of its \file{psd/} directory are outside the scope of
\wfdbx{}.

The speed and memory use of the seven implementations, compared with those
of the C library and tools, are measured in \file{docs/BENCHMARKS.md}. In
short, Rust, Julia and C++ read an hour of a two-signal ECG in 4 to 6\,ms,
4 to 7 times faster than the C library, which reads one frame per
\code{getvec} call; the high-level functions of every language read whole
blocks of samples, and are the fastest way to read a record in that
language.

\section{The seven implementations}
\label{sec:implementations}

All seven implement the complete library and all 33 tools. Each has a
directory of the repository, with a \file{README.md} that gives its
requirements, installation instructions and examples, and a chapter of
\cref{part:languages}. \Cref{tab:implementations} shows how each is
installed, from a copy of the repository (or from PyPI for Python), and how
a tool is run in it. In code, the packages are named \code{wfdbx} (Python,
MATLAB/Octave, Rust, R, Kotlin, and the C++ namespace) or \code{WFDBX}
(Julia).

\begin{table}[htbp]
\caption{The seven implementations: installation, and how a tool is run.}
\label{tab:implementations}
\begin{tabularx}{\linewidth}{@{}P{2.3cm}P{1.5cm}LP{3.0cm}@{}}
\toprule
{\tabhead Language} & {\tabhead Directory} & {\tabhead Installation} &
  {\tabhead Running a tool} \\
\midrule
Python & \file{python/} & \cmd{pip install wfdb-x} (PyPI), or
  \cmd{pip install ./python} & \cmd{wfdbx rdsamp -r 100s} \\
C++ & \file{cpp/} & \cmd{cmake -S cpp -B cpp/build \&\& cmake --build
  cpp/build} & \cmd{rdsamp -r 100s} \\
Rust & \file{rust/} & \cmd{cargo install --path rust} &
  \cmd{rdsamp -r 100s} \\
Julia & \file{julia/} & \cmd{using Pkg;
  Pkg.develop(path="/path/to/wfdb-x/julia")} & \cmd{wfdbx rdsamp -r 100s} \\
R & \file{r/} & \cmd{R CMD INSTALL r} & \cmd{wfdbx rdsamp -r 100s} \\
MATLAB / GNU Octave & \file{matlab/} &
  \cmd{addpath('/path/to/wfdb-x/matlab')} (tested with GNU Octave; MATLAB
  untested) & \cmd{wfdbx rdsamp -r 100s} \\
Kotlin (Android, iOS, JVM) & \file{kotlin/} & \cmd{cd kotlin \&\&
  ./gradlew publishToMavenLocal}, then
  \cmd{implementation("io.github.\allowbreak alphanumericslab:\allowbreak wfdbx:1.2.1")} (Gradle,
  with \code{mavenLocal()}) & \cmd{kotlin/bin/wfdbx rdsamp -r 100s} \\
\bottomrule
\end{tabularx}
\end{table}

In Julia, R and MATLAB/Octave, \cmd{wfdbx} is a short alias for the
language's own launcher of the tools; the chapter of each language in
\cref{part:languages} shows how to set it up. The Python package is
distributed on PyPI as \code{wfdb-x} and installs the \cmd{wfdbx} command;
\cmd{wfdbx --list} lists the tools, \cmd{wfdbx --version} prints the version,
and \cmd{wfdbx --examples} copies the example notebooks of the package.

The Kotlin implementation runs on the JVM, on Android, and on iOS through
Kotlin Multiplatform; it is tested on the JVM. The MATLAB/Octave
implementation is written for both MATLAB and GNU Octave, and is tested with
GNU Octave.

\paragraph{The sample record.} Record \code{100s}, the first minute of record
100 of the MIT-BIH Arrhythmia Database (two ECG signals sampled at 360\,Hz,
and its beat annotations), is in the directory \file{tests/data} of the
repository, and the Python package includes a copy of it (in the directory
that \code{wfdbx.sample\_data\_dir()} returns). The examples of this manual
use it. To run them, run the tools in that directory, or set the input
directory to it (see \cref{chap:files}), for example:
\begin{lstlisting}[style=shell]
cd /path/to/wfdb-x/tests/data
wfdbx rdsamp -r 100s -t s3
\end{lstlisting}
which prints the first three samples of each signal, in ADC units:
\begin{lstlisting}[style=text]
              0	    995	   1011
              1	    995	   1011
              2	    995	   1011
\end{lstlisting}

\section{Three ways to use WFDB-x}
\label{sec:three-ways}

Every language offers the same three ways in, summarized in
\cref{tab:three-ways}. All three are written natively in the language.

\begin{table}[htbp]
\caption{The three ways to use \wfdbx{}.}
\label{tab:three-ways}
\begin{tabularx}{\linewidth}{@{}P{2.8cm}LL@{}}
\toprule
& {\tabhead What it is} & {\tabhead Use it for} \\
\midrule
1.~The high-level API (recommended) &
  A few functions in the idiom of the language: \code{read\_record},
  \code{read\_header}, \code{read\_annotations}, \code{write\_record},
  \code{write\_annotations}, \code{check\_record}, and the header metadata
  functions \code{read\_metadata}, \code{get\_metadata},
  \code{set\_metadata}, \code{remove\_metadata} and \code{update\_metadata}
  (\code{readRecord}, \ldots{} in Kotlin). One call reads a whole record or
  annotation file into arrays, with named options, errors reported as
  exceptions or error values, and fast block reads. It has no counterpart in
  the C package. &
  New code, analysis, notebooks, machine learning: anything that needs a
  record as arrays. \\
\addlinespace
2.~The WFDB library functions &
  The C library API under its C names (\code{isigopen}, \code{getvec},
  \code{getann}, \code{putann}, \code{strtim}, \ldots), with the C behavior,
  frame by frame and annotation by annotation. &
  Porting a program written for the C library, or processing a record as a
  stream with the exact C semantics. \\
\addlinespace
3.~The command-line tools &
  The 33 tools (\cmd{rdsamp}, \cmd{rdann}, \cmd{wrsamp}, \cmd{xform},
  \ldots) with the options, output, messages and exit status of the C
  tools, implemented in each language. &
  Shell scripts and pipelines, and reproducing work done with the original
  tools, without the C package. \\
\bottomrule
\end{tabularx}
\end{table}

The high-level API is built on the library functions, and returns exactly
the samples and annotations that they, and the C library, give; the tools are
built on the library functions too. The three can be mixed in one program,
with one caution: the high-level readers and writers call \code{wfdbquit},
which closes the files opened with the library functions.

In Python, for example, the first three samples of record \code{100s} are
read in physical units with the high-level API as follows:
\begin{lstlisting}[language=Python,style=example]
import wfdbx
rec = wfdbx.read_record("100s", sampto=3, physical=True)
print(rec.fs, rec.names, rec.signals.shape)   # 360.0 ['MLII', 'V5'] (3, 2)
print(rec.signals[0])                         # [-0.145 -0.065]
\end{lstlisting}
and the first sample of each signal is read in ADC units, and converted, with
the library functions:
\begin{lstlisting}[language=Python,style=example]
import wfdbx
si = []
n = wfdbx.isigopen("100s", si, 2)     # 2 signals, described in si
v = [0] * n
wfdbx.getvec(v)                       # v is [995, 1011]
print(wfdbx.aduphys(0, v[0]))         # -0.145
wfdbx.wfdbquit()
\end{lstlisting}

The 33 tools are those of the \file{app/} and \file{convert/} directories of
the C package:
\begin{itemize}
\item signals and annotations: \cmd{rdsamp}, \cmd{wrsamp}, \cmd{rdann},
  \cmd{wrann}, \cmd{sortann}, \cmd{ann2rr}, \cmd{rr2ann}, \cmd{wfdbdesc},
  \cmd{wfdbtime}, \cmd{time2sec}, \cmd{sampfreq}, \cmd{signame},
  \cmd{signum}, \cmd{sigamp}, \cmd{snip}, \cmd{xform}, \cmd{wfdbcat},
  \cmd{wfdbwhich};
\item format converters: \cmd{edf2mit}, \cmd{mit2edf}, \cmd{rdedfann},
  \cmd{wav2mit}, \cmd{mit2wav}, \cmd{wfdb2mat}, \cmd{a2m}, \cmd{m2a},
  \cmd{ad2m}, \cmd{md2a}, \cmd{ahaecg2mit}, \cmd{makeid}, \cmd{readid},
  \cmd{revise}, \cmd{parsescp}.
\end{itemize}
\Cref{part:reference} describes each of them, and each library function.

\section{How the results are checked}
\label{sec:checked}

For the same input, a \wfdbx{} tool produces byte for byte the same standard
output, standard error (including its error messages), exit status and
output files as the original C tool, and the library functions return the
same values as their C counterparts. The exceptions are the deliberate
differences of \cref{chap:behavior}; the reference results are produced by
the C code with these changes applied, so that they are tested in the same
way. This is verified by tests whose expected results were produced by the
original C code, not written by hand:
\begin{itemize}
\item \textbf{The conformance suite} (\file{tests/conformance/}, described in
  \file{docs/conformance.md}): 1135 command-line test cases, run by a
  harness that does not depend on the language of the implementation. Their
  expected results were produced by the C programs, built from the original
  sources (\file{third\_party/wfdb-main.zip}) with FLAC support and with
  the patch \file{tests/reference/divergences.patch}, which defines the
  deliberate differences, applied. The exceptions are 20 cases whose input
  makes the C programs crash; their expected results are the behavior that
  \wfdbx{} defines for that input (\cref{sec:crashes}). A second suite, for
  legacy mode, has 883 cases: the cases of the original package, run with
  \opt{--legacy} and compared with the results of the unmodified C package.
  Every implementation passes all the cases of both suites on the reference
  platform, x86-64 Linux with glibc.
\item \textbf{Test records} (\file{tests/data/}), written by the C library
  itself (\file{tests/reference/gen\_records.c}). They cover every storage
  format, multi-frequency records, multi-segment records of fixed and of
  variable layout, skewed signals, byte offsets, block-structured files, EDF,
  and annotation files that use every feature of the annotation format.
\item \textbf{Unit tests} of the library and of the high-level functions, in
  each language directory, with expected values printed by the C library
  (\file{tests/reference/api\_expect.c}).
\item \textbf{Metadata tests} (\file{tests/metadata/}): 49 cases of the
  header metadata functions, which every implementation must pass with
  byte-identical headers, and a check that the original C programs read the
  changed headers as before (\cref{chap:metadata}).
\end{itemize}
Because all seven implementations must pass the same cases, each is also
checked against the other six. The MATLAB/Octave code is tested with GNU
Octave, also through the MATLAB driver \file{matlab/bin/wfdbx\_matlab.sh};
MATLAB itself has not been available to run the tests.

A few behaviors of the C code are undefined (crashes, reads of uninitialized
memory), and a few depend on the platform. They are listed, with the
behavior that \wfdbx{} chose, in \file{docs/KNOWN\_C\_BEHAVIORS.md} and
summarized in \cref{sec:known-c}.

\section{Versions}
\label{sec:versions}

\wfdbx{} has its own version numbers, shared by the seven implementations,
which follow semantic versioning (\url{https://semver.org}): a new patch
number (1.1.x) for corrections only, a new minor number (1.x.0) for new
features that do not change existing behavior, and a new major number
(x.0.0) for changes that can break existing code. This manual describes
version \wfdbxversion{}. The file \file{CHANGELOG.md} lists the changes in
each version:
\begin{itemize}
\item 1.0.0 (16 October 2026): the WFDB library, its 33 command-line tools
  and a high-level API in Python, C++, MATLAB/Octave, Rust, Julia and R;
\item 1.0.1 (27 October 2026): the Kotlin Multiplatform implementation, the
  seventh language;
\item 1.1.0 (28 October 2026): faster reading and writing of records, signal
  names without surrounding spaces in the high-level readers, and write
  errors and frame checks in the high-level writers;
\item 1.2.0: the header metadata functions (\cref{chap:metadata}), this
  manual, and the reference to cite (\cref{sec:cite}). No existing function
  or tool changes behavior, except for three corrections: the files that a
  tool reads or creates itself follow the input and output directories,
  \code{osigopen(record, NULL, 0)} returns 0 in every language, and the
  high-level functions leave the annotation time resolution and the getvec
  mode as they were (\cref{sec:corrections-1-2-0});
\item 1.2.1: the library under the GNU Lesser General Public License,
  version 2.1 or later, with the command-line tools still under the GNU
  General Public License, version 3 or later (\cref{sec:license}); and the
  C++ library builds on Windows again.
\end{itemize}

The version numbers of \wfdbx{} are independent of the version of the WFDB
Software Package that it implements, \wfdbversion{}, in every version. The
library reports that version as the C library does: \code{wfdbversion()}
returns \code{"10.7.0"}, and \code{wfdberror()}, before any error, returns
\code{"WFDB library version 10.7.0 (Jun 10 2022).\textbackslash{}n"}.

\section{How this manual is organized}
\label{sec:organization}

\begin{itemize}
\item \textbf{\Cref{part:concepts}, Concepts}, describes what the seven
  implementations share: this introduction; the WFDB records and their files
  (\cref{chap:records}); where files are read and written, network access
  and the environment variables (\cref{chap:files}); the differences from
  the original package, legacy mode, and the reporting of errors
  (\cref{chap:behavior}); and the header metadata (\cref{chap:metadata}).
\item \textbf{\Cref{part:languages}, The seven languages}, has one chapter
  for the users of each language: installation, the high-level API with
  every function, its arguments, results, errors and examples, and how the
  library functions and the tools are used in that language
  (\cref{chap:python,chap:cpp,chap:rust,chap:julia,chap:r,chap:matlab,chap:kotlin}).
\item \textbf{\Cref{part:reference}, Reference}, describes the WFDB library
  functions and the command-line tools, which are the same in every
  language, and the storage formats and annotation codes
  (\cref{chap:formats}).
\end{itemize}

\section{Typographic conventions}
\label{sec:conventions}

\begin{itemize}
\item \code{read\_record}, \code{WFDB\_INVALID\_SAMPLE}, \code{"0:10"}:
  identifiers, values and code, in a typewriter font.
\item \cmd{rdsamp -r 100s}: a command line; \opt{-r}: an option of a tool;
  \file{tests/data/100s.hea}: a file or a directory.
\item The WFDB library functions and the tools are named by their C names,
  which are their names in every language.
\item Listings of code and of command lines are set on a shaded background.
  A listing on a white background, with a bold blue rule, is the call syntax
  of a function, copied from the code. The Python examples marked as
  examples are run, as part of the checks of this manual, in a directory
  that holds a copy of \file{tests/data}, which is the current directory and
  the input directory; their results are given in comments.
\end{itemize}
Three kinds of boxes set off remarks:
\begin{note}
A note gives details that matter in some cases, such as a behavior of the C
package that \wfdbx{} reproduces.
\end{note}
\begin{tip}
A tip gives a way to do something faster or more simply.
\end{tip}
\begin{phibox}
A box of this kind concerns data that can identify patients, such as names,
dates or ages stored in a header.
\end{phibox}

\section{How to cite}
\label{sec:cite}

To refer to \wfdbx{} in a publication, cite the preprint that describes it:
\begin{quote}
Reza Sameni. \emph{WFDB-x: Conformance-Tested Implementations of the WFDB
Software Package in Seven Programming Languages.} arXiv:2611.01358, 2026.\\
DOI: \href{https://doi.org/10.48550/arXiv.2611.01358}{10.48550/ARXIV.2611.01358}.\\
\url{https://arxiv.org/abs/2611.01358}
\end{quote}
As a BibTeX entry:
\begin{lstlisting}[style=text]
@misc{wfdb-x,
  doi = {10.48550/ARXIV.2611.01358},
  url = {https://arxiv.org/abs/2611.01358},
  author = {Sameni,  Reza},
  title = {{WFDB-x: Conformance-Tested Implementations of the WFDB Software Package in Seven Programming Languages}},
  publisher = {arXiv},
  year = {2026},
  copyright = {Creative Commons Attribution Share Alike 4.0 International}
}
\end{lstlisting}
The same reference is in the file \file{CITATION.cff} of the repository.
Work that uses the WFDB format or PhysioNet data should also cite the WFDB
Software Package and PhysioNet, as PhysioNet asks.

\section{Authors and license}
\label{sec:license}

The WFDB Software Package was written by George B.~Moody, Benjamin Moody and
the WFDB contributors; it is Copyright \textcopyright{} 1983--2022 George
B.~Moody and the WFDB contributors. WFDB was originally developed at MIT by
George B.~Moody, and builds on decades of work by him, Roger G.~Mark and the
broader research community. The authors of the C code that each part of
\wfdbx{} translates are credited, file by file, in each \wfdbx{} source file
and in \file{docs/AUTHORS.md}.

\wfdbx{} was developed by Reza Sameni (2026). It is an independent
reimplementation of the WFDB Software Package: its code in each language is a
translation of the package's C code, written for \wfdbx{}, not code by the
original WFDB developers.

The \wfdbx{} library is distributed under the GNU Lesser General Public
License, version 2.1 or (at the user's option) any later version
(LGPL-2.1-or-later; see the file \file{LICENSE}): programs under any license,
proprietary ones included, may use it on the LGPL's terms. The command-line
tools (each language's \code{apps} and the \cmd{wfdbx} command), translations
of the tools of the original package, are distributed under the GNU General
Public License, version 3 or (at the user's option) any later version
(GPL-3.0-or-later; see the file \file{LICENSE-GPL}). Each source file states
its license. The original C sources, kept unmodified in
\file{third\_party/}, keep their own licenses: the GNU Library General Public
License, version 2 or later, for the library, and the GNU General Public
License, version 2 or later, for the tools; their ``or later'' terms allow
those of \wfdbx{}.

% Part I, chapter 2: WFDB records.
\chapter{WFDB records}
\label{chap:records}

A WFDB \emph{record} is a recording of one or more signals, such as the
leads of an ECG, and the files that describe and annotate it: a
\emph{header file}, which describes the record and its signals; one or more
\emph{signal files}, which hold the samples; and any number of
\emph{annotation files}, which label events of the signals, such as beats.
A record may also have an \emph{info file} of comments. This chapter
describes these files as \wfdbx{} reads and writes them, which is as the
WFDB library 10.7.0 does, except where \cref{chap:behavior} says otherwise.
The exact layout of each storage format, and the annotation codes, are in
\cref{chap:formats}.

\section{Records and record names}
\label{sec:record-names}

A record is designated by its \emph{record name}, such as \code{100s}. The
files of a record are found from it:
\begin{itemize}
\item the header is the file \emph{record}\file{.hea} (\file{100s.hea});
\item the signal files are named in the header (\file{100s.dat});
\item an annotation file is the file \emph{record}\code{.}\emph{annotator},
  where the \emph{annotator name} says which set of annotations it holds
  (\file{100s.atr}, the reference annotations of record \code{100s},
  annotator \code{atr});
\item the info file is \emph{record}\file{.info}.
\end{itemize}

A record name may include a directory (\code{mitdb/100},
\code{/data/mitdb/100}), and is then found as \cref{chap:files} describes. A
trailing \file{.hea} is removed from a record name by the library functions
that open a record, so \code{100s.hea} designates record \code{100s}. A name
that ends with \code{/} stands for a record of the same name in that
directory: \code{abc/123/} is \code{abc/123/123}. A record name that ends
with another extension, such as \code{100s.edf}, names an EDF file
(\cref{sec:edf}). The record name \code{-} stands for the standard input or
output (\cref{sec:stdio}).

When a header or an annotation file is created, the record name, and the
annotator name, may contain only letters, digits, \code{\_}, \code{-},
\code{\textasciitilde} and \code{/}; any other character is an error (for
example, \code{init: illegal character 46 in record name}). In a record name
written under a path, such as \code{out.d/rec}, only the last component, the
record's own name, is checked (the C library checked the whole name; see
\cref{sec:resolve}).

\section{The header file}
\label{sec:header}

A header is a text file of lines separated by line feeds; a carriage return
may precede each line feed. Its fields are separated by spaces or tabs. The
lines are of four kinds:
\begin{itemize}
\item the \emph{record line}, the first line that is neither empty nor a
  comment;
\item one \emph{signal specification line} per signal, for an ordinary
  record, or one \emph{segment specification line} per segment, for a
  multi-segment record (\cref{sec:multisegment});
\item \emph{comment lines}, whose first character is \code{\#}; those that
  follow the record line are the \emph{info strings} of the record;
\item empty lines, which are ignored.
\end{itemize}
Comment lines and empty lines may appear anywhere. If the first line of a
header begins with \code{\#wfdb} followed by a version number, such as
\code{\#wfdb 10.7.1}, a library older than that version refuses the header
(\code{init: reading record REC requires WFDB library version 10.7.1 or
later}); \wfdbx{} has the version of the C library that it implements,
10.7.0.

\subsection{The record line}

The record line gives the fields of \cref{tab:record-line}, in this order,
separated by spaces or tabs. Only the record name and the number of signals
are required; each optional field may be given only if the fields before it
are.

\begin{table}[htbp]
\caption{The fields of the record line.}
\label{tab:record-line}
\begin{tabularx}{\linewidth}{@{}P{3.5cm}L@{}}
\toprule
{\tabhead Field} & {\tabhead Meaning} \\
\midrule
record name & The name of the record. It must match the last component of
  the name under which the header is opened, or the header is refused
  (\code{init: record name in record REC header is incorrect}). \\
\code{/}\emph{segments} & (Optional, attached to the record name.) The
  number of segments of a multi-segment record, as in \code{multi/3}. \\
number of signals & An integer; 0 is allowed. Several signals can share one
  signal file. \\
sampling frequency & (Optional.) The number of frames per second, a
  floating-point number greater than 0 (\code{360}, \code{360.0} and
  \code{3.6e2} are equivalent). \wfdbx{} assumes none: if it is missing,
  the frequency is unspecified (0), unless one was set with
  \code{setsampfreq} before the record was opened (the C library assumed
  250\,Hz; see \cref{sec:no-defaults}). \\
\code{/}\emph{counter frequency} & (Optional, attached to the sampling
  frequency.) The frequency of a counter, such as the page numbers of a
  chart recording, used by the time strings that begin with \code{c}
  (\cref{sec:times}). If it is absent or not positive, it is the sampling
  frequency. \\
\code{(}\emph{base counter}\code{)} & (Optional, attached to the counter
  frequency.) The counter value of sample 0; 0 if absent. \\
number of samples & (Optional.) The number of samples per signal. 0 or a
  missing value means that the length is unspecified, and then the
  checksums of the signals are not verified. \\
base time & (Optional.) The time of day of sample 0, as
  \code{hh:mm:ss}, with optional milliseconds (\code{12:34:56.789}). \\
base date & (Optional, after the base time.) The date of sample 0, as
  \code{dd/mm/yyyy}. \\
\bottomrule
\end{tabularx}
\end{table}

\subsection{Signal specification lines}

Each signal specification line gives the fields of
\cref{tab:signal-line}, in this order; again each optional field may be
given only if the fields before it are. The lines of the signals that share
a signal file must be consecutive; such signals form a \emph{signal group},
and must have the same format and block size (otherwise
\code{init: error in specification of signal N or M}).

\begin{table}[htbp]
\caption{The fields of a signal specification line, and the values that
\wfdbx{} uses when a field is missing.}
\label{tab:signal-line}
\begin{tabularx}{\linewidth}{@{}P{3.3cm}L@{}}
\toprule
{\tabhead Field} & {\tabhead Meaning, and value when missing} \\
\midrule
file name & The signal file, relative to the directory of the header
  (\cref{sec:header-files}); \code{-} is the standard input. \\
format & The storage format (\cref{sec:signal-files}), an integer such as
  \code{212}. Required. \\
\code{x}\emph{spf} & (Attached to the format.) The number of samples per
  frame (\cref{sec:frames}); 1 if missing (or less than 1). \\
\code{:}\emph{skew} & (Attached.) The number of samples of the signal that
  precede sample 0, to correct a skew between signals; 0 if missing. \\
\code{+}\emph{offset} & (Attached.) The number of bytes at the beginning of
  the signal file that precede the samples (a prolog); 0 if missing. \\
gain & The number of ADC units per physical unit, a floating-point number.
  0, or a missing field, means that the signal is \emph{uncalibrated}
  (\cref{sec:units}). \\
\code{(}\emph{baseline}\code{)} & (Attached to the gain.) The ADC value
  that corresponds to 0 physical units; the ADC zero if missing. \\
\code{/}\emph{units} & (Attached to the gain or baseline.) The physical
  units, such as \code{mV}, without spaces (at most 50 characters). If
  missing, the units are unspecified; \wfdbx{} never takes them to be
  millivolts, and its tools show them as \code{nd}. \\
ADC resolution & The resolution of the analog-to-digital converter, in
  bits. If missing, it is unspecified (0); the C library inferred it from the
  format or assumed 12 bits. \\
ADC zero & The ADC value at the middle of the ADC's input range; 0 if
  missing. \\
initial value & The value of sample 0, used by the difference format 8;
  the ADC zero if missing. \\
checksum & The sum of all the samples of the signal, as a 16-bit signed
  integer, verified when the end of the signal is read (if the record line
  gives the number of samples); 0 if missing. \\
block size & For signal files that must be read in blocks of a given size
  (character special files); 0 (none) if missing. \\
description & The rest of the line, which may contain spaces (at most 100
  characters), such as \code{MLII}. If missing:
  \code{record REC, signal N}. \\
\bottomrule
\end{tabularx}
\end{table}

\subsection{Info strings and info files}
\label{sec:info-strings}

The \emph{info strings} of a record are the comment lines of its header
that follow the record line, without their initial \code{\#}, as the
library function \code{getinfo} reads them: every line after the record line
whose first character is \code{\#} is an info string, even one between two
signal lines. A line that begins with a space before its \code{\#} is not
an info string. The comment lines before the record line are not info
strings. \code{getinfo} then returns the comment lines of the record's info
file, \emph{record}\file{.info}, if there is one. \code{putinfo} writes an
info string, as \code{\#} followed by the string and a CR LF line ending,
into the header being written, or into the info file opened by
\code{setinfo}. \Cref{chap:metadata} describes the info strings of the form
\code{Key: value}, which hold the metadata of a record.

\subsection{An example: the header of record 100s}

The header of the sample record \code{100s}, \file{tests/data/100s.hea}, is:
\begin{lstlisting}[style=text]
100s 2 360 21600
100s.dat 212 200 11 1024 995 21537 0 MLII
100s.dat 212 200 11 1024 1011 -3962 0 V5
# 69 M 1085 1629 x1
# Aldomet, Inderal
#Produced by xform from record 100, beginning at 0:0
\end{lstlisting}
\begin{itemize}
\item The record line says that record \code{100s} has 2 signals, sampled at
  360 frames per second, of 21600 samples each (60 seconds). It gives no
  base time or date.
\item Each signal line says that the signal is in the file
  \file{100s.dat} (both signals share it: they form signal group 0), in
  format 212, with a gain of 200 ADC units per physical unit, no units, an
  ADC resolution of 11 bits, an ADC zero of 1024, an initial value of 995
  (signal 0) or 1011 (signal 1), a checksum of 21537 or $-3962$, a block
  size of 0, and the description \code{MLII} or \code{V5}. The baseline is
  not given, so it is the ADC zero, 1024. A sample of 995 ADC units in
  signal 0 is thus $(995 - 1024)/200 = -0.145$ physical units.
\item The last three lines are the info strings
  \code{" 69 M 1085 1629 x1"}, \code{" Aldomet, Inderal"} and
  \code{"Produced by xform from record 100, beginning at 0:0"}. The first
  two are free text; the third contains a colon, and is therefore read as a
  metadata entry (\cref{chap:metadata}).
\end{itemize}

\section{Signal files and storage formats}
\label{sec:signal-files}

A signal file holds the samples of the signals of one signal group, as a
sequence of frames (\cref{sec:frames}): the samples of each frame are stored
one signal after the other. The \emph{storage format} of the group says how
each sample is encoded. \Cref{tab:formats-summary} summarizes the formats;
\cref{chap:formats} gives their exact layout.

\begin{table}[htbp]
\caption{The storage formats (summary).}
\label{tab:formats-summary}
\begin{tabularx}{\linewidth}{@{}P{1.4cm}LP{2.6cm}@{}}
\toprule
{\tabhead Format} & {\tabhead Encoding} & {\tabhead Invalid sample} \\
\midrule
0 & no signal file: a null signal, whose samples are all invalid & (all) \\
8 & 8-bit first differences, from the initial value of the header & --- \\
16 & 16-bit two's complement, least significant byte first & $-32768$ \\
24 & 24-bit two's complement, least significant byte first & $-2^{23}$ \\
32 & 32-bit two's complement, least significant byte first & $-2^{31}$ \\
61 & 16-bit two's complement, most significant byte first & $-32768$ \\
80 & 8-bit offset binary (128 is subtracted from each byte) & $-128$ \\
160 & 16-bit offset binary, least significant byte first & $-32768$ \\
212 & 12-bit two's complement, two samples packed in 3 bytes & $-2048$ \\
310 & 10-bit two's complement, three samples packed in 4 bytes (two 16-bit
  words) & $-512$ \\
311 & 10-bit two's complement, three samples packed in a 32-bit word & $-512$ \\
508, 516, 524 & FLAC-compressed, with 8, 16 or 24 bits per sample (at most
  8 signals per file, all with the same number of samples per frame) &
  $-128$, $-32768$, $-2^{23}$ \\
\bottomrule
\end{tabularx}
\end{table}

The last column is the value that marks an invalid (missing) sample in the
file: the most negative value of the format. The library returns an invalid
sample, whatever its format, as \code{WFDB\_INVALID\_SAMPLE} ($-32768$), and
the high-level readers give it as NaN in physical units. A file in format
212, 310 or 311 may end with a partial group of samples. FLAC files are read
and written by \wfdbx{} without libFLAC: the files it writes decode to the
same samples as those of the C library, but their compressed bytes differ.

Three fields of the signal line change how a signal file is read: the
\emph{byte offset} skips a prolog at the beginning of the file; the
\emph{skew} discards the first samples of a signal, to align it with the
others (the skipped samples are counted in the checksum); and the
\emph{block size} makes the file be read in blocks of that size.

\section{ADC units and physical units}
\label{sec:units}

A signal file stores samples as integers, in \emph{ADC units} (adu), the
output of the analog-to-digital converter. The header gives, for each
signal, the values that relate them to \emph{physical units}, such as
millivolts:
\begin{description}
\item[gain] $g$, the number of ADC units per physical unit (0: the signal
  is uncalibrated);
\item[baseline] $b$, the ADC value that corresponds to 0 physical units;
\item[units] the physical units, such as \code{mV} (unspecified if the
  header does not give them);
\item[ADC resolution] $r$, in bits, and \textbf{ADC zero} $z$, the value at
  the middle of the ADC's range. Together they give the range of the ADC:
  from $z - 2^{r-1}$ to $z + 2^{r-1} - 1$.
\end{description}

The conversions are those of the C library, and of the high-level
functions:
\begin{itemize}
\item From ADC units to physical units (\code{aduphys(s, a)}, and
  \code{read\_record} with physical units): $v = (a - b)/g$, as a
  floating-point number. An invalid sample becomes NaN in the high-level
  readers.
\item From physical units to ADC units (\code{physadu(s, v)}, and
  \code{write\_record} with physical units): $a = \mathrm{round}(v g) + b$,
  where $\mathrm{round}(x)$ is $\lfloor x + 0.5\rfloor$ for $x \ge 0$ and
  $\lceil x - 0.5\rceil$ for $x < 0$ (halves are rounded away from zero).
  In \code{write\_record}, NaN becomes an invalid sample.
\item \code{adumuv(s, a)} converts $a$ ADC units to microvolts, as
  $\mathrm{round}(1000\,a/g)$, and \code{muvadu(s, v)} converts $v$
  microvolts to ADC units, as $\mathrm{round}(g\,v/1000)$. Neither uses the
  baseline; both take the physical units to be millivolts.
\end{itemize}
In these functions, $g$ and $b$ are those of input signal \code{s} of the
open record. \wfdbx{} assumes no gain: for an uncalibrated signal (gain 0)
or a signal number that is not that of an input signal, the four functions
report an error (\code{aduphys: signal 0 is uncalibrated},
\code{aduphys: invalid signal number 5}) and return NaN (\code{aduphys}) or
\code{WFDB\_INVALID\_SAMPLE} (the other three); the C library assumed a gain
of 200 (\cref{sec:no-defaults}).

\section{Frames, samples per frame and multi-frequency records}
\label{sec:frames}

A \emph{frame} is the set of samples of all the signals that belong to one
sampling interval of the record. The sampling frequency of the record line
is the frequency of the frames. Usually each signal has one sample per
frame. A signal sampled at an integer multiple $n$ of that frequency has
$n$ samples per frame (the \code{x}\emph{n} suffix of its format), and its
own sampling frequency is $n$ times that of the frames. A record with such
signals is a \emph{multi-frequency} record. For example, the test record
\file{tests/data/mfreq.hea}:
\begin{lstlisting}[style=text]
mfreq 3 100 600
mfreq.dat 16 200/mV 16 0 138 -16869 0 slow
mfreq.dat 16x2 200/mV 16 0 218 5410 0 medium
mfreq.dat 16x4 200/mV 16 0 219 -15996 0 fast
\end{lstlisting}
has 600 frames at 100 frames per second; the signals are sampled at 100,
200 and 400\,Hz, and each frame of its signal file holds 1 + 2 + 4 = 7
samples.

The library reads such a record in one of two modes, set by
\code{setgvmode}:
\begin{itemize}
\item \textbf{Low resolution} (\code{WFDB\_LOWRES}, the default):
  \code{getvec} returns one sample per signal per frame. The samples of a
  signal that has more than one per frame are averaged: their mean is
  rounded to the nearest integer, halves upward, and is invalid if one of
  them is invalid. \code{sampfreq} returns the frame frequency.
\item \textbf{High resolution} (\code{WFDB\_HIGHRES}): \code{getvec}
  returns as many sample vectors per frame as the largest number of samples
  per frame (4 for \code{mfreq}). Each holds the next sample of each signal
  that has that many samples per frame; the samples of the other signals
  are repeated. \code{sampfreq} returns the frame frequency times that
  number (400\,Hz for \code{mfreq}), and times (\cref{sec:times}) count
  these sample vectors.
\end{itemize}
\code{WFDB\_GVPAD}, which can be combined with either mode, replaces each
invalid sample by the previous valid sample of the same signal.
\code{getframe} returns all the samples of a frame. In the high-level API,
\code{read\_record} reads in low resolution by default, and in high
resolution with its \code{highres} option; its \code{pad} option is
\code{WFDB\_GVPAD}. For \code{mfreq}, the first frame is read as
\code{138 1459 499} in low resolution (1459 is the mean of 218 and 2699,
rounded up), and as these four rows in high resolution:
\begin{lstlisting}[style=text]
138   218    219
138   218  -1316
138  2699   1120
138  2699   1973
\end{lstlisting}
The library can also resample all the signals of a record to another
frequency, by linear interpolation (\code{setifreq}).

\section{Multi-segment records}
\label{sec:multisegment}

A multi-segment record is a sequence of ordinary records, its
\emph{segments}, read as if their signals were continuous. Its header (the
\emph{master header}) has a record line whose record name is followed by
\code{/} and the number of segments, and then one segment specification
line per segment: the record name of the segment and its length in samples.
The segments' own headers and signal files are found in the directory of the
master header (\cref{sec:header-files}). A segment cannot itself be a
multi-segment record (\code{init: record REC cannot be nested in another
multi-segment record}), and the sampling frequency must be the same in all
segments. If the record line gives a number of samples that differs from the
sum of the segment lengths, the library warns
(\code{warning (init): in record REC, stated record length (N)} \ldots).

In a \emph{fixed-layout} record, every segment has the same signals, in
the same order, with the same gains and baselines (their storage formats
can differ). The test record \file{tests/data/multi.hea} is one: two copies
of record \code{100s} with a null record between them.
\begin{lstlisting}[style=text]
multi/3 2 360 45000
100s 21600
null 1800
100s 21600
\end{lstlisting}

In a \emph{variable-layout} record, the signals may differ from segment to
segment. Segment 0, of length 0, is then the \emph{layout segment}: its
header, whose signal file names are \code{\textasciitilde}, gives the
signals of the record, with their gains, baselines and units. When the
record is read, the signals of each segment are rearranged, shifted and
rescaled to the gains and baselines of the layout (the results rounded to
the nearest integer, halves upward); a signal absent from a segment reads as
invalid samples. If a segment's gain is
larger than that of the layout, the library warns
(\code{sigmap\_init: loss of precision in signal N in segment SEG}). A
segment whose record name is \code{\textasciitilde} is a \emph{null
segment}: it has no files, and all its samples are invalid. The test record
\file{tests/data/vlay.hea} and its layout segment \file{vlay\_0.hea}:
\begin{lstlisting}[style=text]
vlay/4 3 250 3201
vlay_0 0
fmt16 1500
~ 200
fmt212 1501
\end{lstlisting}
\begin{lstlisting}[style=text]
vlay_0 3 250 0
~ 0 400(10)/mV 12 0 0 0 0 sig0 fmt16
~ 0 400/uV 12 0 0 0 0 sig1 fmt212
~ 0 100/mmHg 16 0 0 0 0 sig2 fmt16
\end{lstlisting}

\section{Annotation files}
\label{sec:annotations}

An annotation file holds a sequence of \emph{annotations}, each of which
labels a point of the record, such as a beat. It is the file
\emph{record}\code{.}\emph{annotator}. Each annotation has the fields of
\cref{tab:annotation-fields}, those of the C structure
\code{WFDB\_Annotation}.

\begin{table}[htbp]
\caption{The fields of an annotation.}
\label{tab:annotation-fields}
\begin{tabularx}{\linewidth}{@{}P{2.2cm}L@{}}
\toprule
{\tabhead Field} & {\tabhead Meaning} \\
\midrule
\code{time} & The time of the annotation, in sample intervals from the
  beginning of the record (a 64-bit integer). \\
\code{anntyp} & The \emph{annotation code}, from 1 to \code{ACMAX} (49),
  such as \code{NORMAL} (1) or \code{PVC} (5). Each code has a
  \emph{mnemonic}, such as \code{N} or \code{V}, and a description
  (\code{annstr}, \code{anndesc}); \cref{chap:formats} lists them. \\
\code{subtyp} & The subtype, a signed 8-bit integer ($-128$ to 127). \\
\code{chan} & The channel (signal) number, an unsigned 8-bit integer (0 to
  255). \\
\code{num} & The annotator number, a signed 8-bit integer. \\
\code{aux} & Optional auxiliary data, up to 255 bytes, such as the rhythm
  \code{(N} of a \code{RHYTHM} annotation, or the text of a \code{NOTE}. In
  the library functions, its first byte is its length, as in C. \\
\bottomrule
\end{tabularx}
\end{table}

The high-level \code{read\_annotations} returns these fields as arrays:
the times (\code{samples}), the codes, their mnemonics (\code{symbols}),
the subtypes, the channels, the annotator numbers and the aux strings.

Annotation files are in one of two binary formats, which the library tells
apart when it opens a file (with a warning, \code{file appears to be in AHA
format}, if the caller asked for the other):
\begin{itemize}
\item \textbf{MIT format}, the usual one: each annotation is a 16-bit word,
  least significant byte first, whose 6 high bits are the code and 10 low
  bits the interval from the previous annotation, followed by optional
  words for the fields that change (pseudo-codes \code{SKIP} for a long
  interval, \code{NUM}, \code{SUB}, \code{CHN} and \code{AUX}); a word of 0
  ends the file.
\item \textbf{AHA format}: 16 bytes per annotation, as in the distribution
  files of the AHA database.
\end{itemize}

An annotation file is in \emph{canonical order} when its annotations are
sorted by time, then annotator number, then channel. An output file that is
not is sorted when it is closed (\cref{sec:sorting}). When a program has set
a time resolution with \code{setafreq} (as \cmd{wrann} does), the library
writes it at the beginning of the file, as a \code{NOTE} annotation whose
aux string is \code{\#\# time resolution: F}. Annotation files can also
begin with \emph{modification labels}, \code{NOTE} annotations at time 0
that redefine the mnemonics and descriptions of codes; the library reads
them when it opens the file, so \code{getann} does not return them.

\section{Times and dates}
\label{sec:times}

Times are counted in \emph{sample intervals} from the beginning of the
record (in high resolution, in intervals of the fastest signal;
\cref{sec:frames}). The library function \code{strtim} converts a time
string to a number of sample intervals, and \code{timstr} and
\code{mstimstr} convert back; the tools take their time arguments
(\opt{-f}, \opt{-t}, \ldots) as \code{strtim} strings, and the high-level
readers accept them for \code{sampfrom} and \code{sampto}.

\subsection{Time strings}

\code{strtim} first skips spaces, tabs, line feeds and carriage returns, and
then reads the string according to its first character, as in
\cref{tab:strtim}. In the table, $f$ is the sampling frequency (the one set
by \code{setifreq}, if any), and a number is rounded to an integer by adding
0.5 and truncating toward zero.

\begin{table}[htbp]
\caption{The time strings that \code{strtim} accepts.}
\label{tab:strtim}
\begin{tabularx}{\linewidth}{@{}P{3.6cm}L@{}}
\toprule
{\tabhead String} & {\tabhead Value (sample intervals)} \\
\midrule
\code{s}\emph{n}, as \code{s123} & Sample number $n$. \\
\code{e} & The end of the record: its number of samples (times the number
  of samples per frame of the fastest signal, in high resolution); 0 if it
  is not known. \\
\code{i} & The time of the next sample to be read from the input
  signals. \\
\code{o} & The time of the next sample to be written to the output
  signals. \\
\code{c}\emph{x}, as \code{c350.5} & Counter value $x$: $(x - x_0)\,f/f_c$,
  with the counter frequency $f_c$ and base counter $x_0$ of the header.
  With no counter frequency, $x$ is read as a sample number. \\
\code{f}\emph{n}, as \code{f100} & Frame number $n$: $n f/f_0$, where $f_0$
  is the frame frequency (in high resolution, the number of the first sample
  of that frame). \\
\code{[}\emph{time}\code{]}, as \code{[13:06:00]} & A time of day: the
  value is the \emph{negated} number of sample intervals from the base time
  of the record to that time, so that it is 0 or negative. \\
\code{[}\emph{time} \emph{d}\code{]}, as \code{[8:00:00 1]} & A time of day
  $d$ days after the base date. \\
\code{[}\emph{time} \emph{date}\code{]}, as \code{[12:00:00 1/3/1992]} & A
  time of day on a date (\code{dd/mm/yyyy}). \\
\emph{s}, as \code{143} or \code{14.875} & An elapsed time of $s$ seconds:
  $s f$. \\
\emph{m}\code{:}\emph{s}, as \code{2:14.875} & $m$ minutes and $s$
  seconds. \\
\emph{h}\code{:}\emph{m}\code{:}\emph{s}, as \code{4:02:01} & $h$ hours, $m$
  minutes and $s$ seconds. \\
\bottomrule
\end{tabularx}
\end{table}

A string that cannot be read gives 0, which is also a valid result
(\code{strtim("0")}). The high-level readers make the value of a time of day
positive, so that it is the number of the sample at that time. For record
\code{100s} (360\,Hz), \code{"0:10"} and \code{"10"} are 3600, \code{"s123"}
is 123, \code{"1:02:03.5"} is 1340460, \code{"e"} is 21600, and
\code{"[00:00:10]"} is $-3600$.

\wfdbx{} assumes no sampling frequency: when none is known (no record is
open, or its header gives none, and none was set with \code{setsampfreq}),
\code{strtim} reports \code{strtim: sampling frequency is not specified} and
returns 0 for every string that needs one, that is, for all the strings of
\cref{tab:strtim} except those that begin with \code{s}, \code{e},
\code{i} or \code{o}, and those that begin with \code{c} when there is no
counter frequency. The function \code{wfdbx\_time\_needs\_freq(string)}
returns 1 if \code{strtim} needs the frequency for a string, and 0
otherwise. The C library used 1 sample per second instead
(\cref{sec:no-defaults}).

\subsection{Formatted times}

\code{timstr(t)} converts a time $t$ to text:
\begin{itemize}
\item if $t > 0$, or if the record has no base time or date, as an elapsed
  time $|t|/f$, truncated to whole seconds: \code{h:mm:ss}, or, under one
  hour, \code{mm:ss}, right-aligned in 8 characters (\code{"    0:10"} for
  10 seconds);
\item if $t \le 0$ and the record has a base time or date, as the time of
  day at $-t$ sample intervals after the base time, in brackets:
  \code{[hh:mm:ss dd/mm/yyyy]}, with the date if the record has a base date
  (or the number of days after the base time if it has none).
\end{itemize}
\code{mstimstr(t)} does the same with milliseconds (\code{"    0:10.003"},
\code{[12:34:56.789 31/12/1999]}). Without a sampling frequency, both report
\code{timstr: sampling frequency is not specified} (or \code{mstimstr: \ldots})
and return the time as a sample number, \code{s} followed by $|t|$, which
\code{strtim} reads back.

\subsection{Base time and dates}

The base time and date of a record come from its header, or from
\code{setbasetime}, which takes a string \code{hh:mm:ss}, with optional
milliseconds and an optional date \code{dd/mm/yyyy} after a space (an empty
string sets the current local date and time). \code{strdat} and
\code{datstr} convert between dates \code{dd/mm/yyyy} and day numbers. The
high-level readers give the base time as text (\code{base\_time}, as
\code{timstr(0)} gives it, to the second), as a number of milliseconds after
midnight (\code{base\_time\_ms}), and the base date as year, month and day
(\code{base\_date}).

\section{Calibration files}
\label{sec:calibration}

A calibration file lists, for types of signals, the calibration pulses with
which they were recorded and their customary plotting scale. Each line that
is not a comment (\code{\#}) has six fields:
\begin{lstlisting}[style=text]
DESC<tab>LOW HIGH TYPE SCALE UNITS
\end{lstlisting}
\code{DESC} is a signal type, matched against the beginning of a signal's
description (\code{*} matches any signal); \code{LOW} and \code{HIGH} are
the levels of the calibration pulse in physical units (\code{-} for
\code{LOW}: the signal is AC-coupled; \code{-} for \code{HIGH}: not
defined); \code{TYPE} is \code{sine}, \code{square}, \code{sawtooth} or
\code{undefined}; \code{SCALE} is in physical units per centimeter; and
\code{UNITS} are the physical units. For example:
\begin{lstlisting}[style=text]
ECG	- 1 sine 1 mV
BP	0 100 square 100 mmHg
\end{lstlisting}
\code{calopen(name)} reads a calibration file into the library's
calibration list (emptied first, unless the name begins with \code{+});
without a name, it reads the file named by the environment variable
\code{WFDBCAL}, or \file{wfdbcal}. The file is found as any other file
(\cref{chap:files}). \code{getcal} finds the first entry for a signal's
description and units, \code{putcal} adds one, and \code{newcal} writes the
list to a file. The file \file{tests/data/wfdbcal} is the default
calibration file of the C package.

\section{EDF and EDF+ records}
\label{sec:edf}

The library reads EDF and EDF+ files (and BDF files, with 24-bit samples)
directly, as records. A record name with an extension other than
\file{.hea}, such as \code{100s.edf}, is opened as an EDF file: its header
is read from the EDF header, and its samples from the data records that
follow. Each EDF data record is one frame, so a signal has as many samples
per frame as it has samples per data record, and the frame frequency is the
inverse of the duration of a data record. The signals are in format 16 (24
for BDF). The gain of a signal is
$(d_{\max} - d_{\min})/(p_{\max} - p_{\min})$, from its digital and physical
ranges (0 if the physical range is empty), its baseline is
$d_{\max} - p_{\max} \times \mathrm{gain}$ rounded to the nearest integer,
and its ADC zero is the middle of the digital range. The record's base time
and date are the start time and date of the EDF file.

Because frames are data records, the library opens an EDF file in high
resolution (\cref{sec:frames}): the sample intervals are those of the
fastest signal, and \code{sampfreq} gives its frequency. A program that sets
another mode reads one sample per data record per signal; the high-level
\code{read\_record}, which sets the mode, returns every sample of an EDF file
only with its \code{highres} option.

The command-line tools \cmd{edf2mit} and \cmd{mit2edf} convert EDF files to
WFDB records and back, and \cmd{rdedfann} prints the annotations of an EDF+
file in the format of \cmd{rdann}.

% Part I, chapter 3: Where files are read and written.
\chapter{Where files are read and written}
\label{chap:files}

\wfdbx{} reads and writes a file exactly where its name says, and nowhere
else. A full path, a \code{\textasciitilde} path or a URL is used as given;
any other name is read from the \emph{input directory} and written to the
\emph{output directory}, which are both the current directory unless they
are set. Nothing is searched, and URLs are opened only when network access
has been turned on. This differs from the C library, which searched a list
of places, the WFDB path, for every file it read (see
\cref{sec:named-files} for the reasons); legacy mode restores the search
(\cref{sec:legacy}).

\section{The input and output directories}
\label{sec:directories}

\Cref{tab:directories} shows the three settings that decide where files
are read and written. Each can be set by a library function, which has the
same name in every language, and by an environment variable; a call of the
function overrides the environment variable.

\begin{table}[htbp]
\caption{The settings that locate files.}
\label{tab:directories}
\begin{tabularx}{\linewidth}{@{}P{3.1cm}LP{2.6cm}P{2.2cm}@{}}
\toprule
{\tabhead Setting} & {\tabhead Functions} & {\tabhead Environment variable}
  & {\tabhead Default} \\
\midrule
input directory (a directory, a \code{\textasciitilde} path, or a URL) &
  \code{wfdbx\_set\_input\_dir}, \code{wfdbx\_get\_input\_dir} &
  \code{WFDBX\_INPUT\_DIR} & the current directory \\
output directory (a local directory or a \code{\textasciitilde} path) &
  \code{wfdbx\_set\_output\_dir}, \code{wfdbx\_get\_output\_dir} &
  \code{WFDBX\_OUTPUT\_DIR} & the current directory \\
network access & \code{wfdbx\_set\_network}, \code{wfdbx\_get\_network} &
  \code{WFDBX\_NETWORK} & off \\
\bottomrule
\end{tabularx}
\end{table}

\begin{itemize}
\item \code{wfdbx\_set\_input\_dir(d)} sets the input directory, and returns
  the previous one (or none, if none was set). \code{d} is a local
  directory, a \code{\textasciitilde} path, or the URL of a directory, such
  as \code{https://physionet.org/files/mitdb/1.0.0/}; reading from a URL also
  needs network access. None (null, \code{NULL}, \ldots, depending on the
  language) or an empty string restores the default: the value of
  \code{WFDBX\_INPUT\_DIR} if it is set, and otherwise the current directory.
  \code{wfdbx\_get\_input\_dir()} returns the input directory, or none.
\item \code{wfdbx\_set\_output\_dir(d)} sets the output directory in the
  same way, and returns the previous one; the default is the value of
  \code{WFDBX\_OUTPUT\_DIR}, or the current directory. The output directory
  must be local: a name that begins with a URL scheme, such as
  \code{https://} or \code{ftp://} (given to the function, or in the
  environment variable), is refused with the message
  \code{wfdbx\_set\_output\_dir: the output directory must be a local
  directory}, and the setting does not change.
\item \code{wfdbx\_set\_network(on)} turns network access on (1, or true) or
  off (0, or false), and returns the previous setting (1 or 0); none or a
  negative number restores the default: on if \code{WFDBX\_NETWORK} is set
  to a nonzero number (such as \code{1}), off otherwise.
  \code{wfdbx\_get\_network()} returns 1 or 0.
\end{itemize}
In Kotlin, these functions are also available under camelCase names
(\code{wfdbxSetInputDir}, \ldots). The settings apply to the whole library:
to the high-level functions, the library functions and the tools, which use
the environment variables, also for the files that they open themselves
(\cref{sec:tool-files}). \code{wfdbquit} does not reset them.

The high-level functions also take these settings as options, which apply
for the duration of one call, after which the previous settings are
restored: the readers take an input directory (\code{input\_dir}) and a
network permission (\code{network}), and the writers an output directory
(\code{dir}). Each of these options is optional; \code{input\_dir} applies
to relative record names only, and \code{dir} to plain record names
(\cref{sec:record-paths}). In Kotlin the options are \code{inputDir},
\code{network} and \code{dir}.

\paragraph{The C functions setwfdb and getwfdb.} In \wfdbx{},
\code{setwfdb} and \code{getwfdb} set and return the input directory.
\code{setwfdb} splits its argument as the C library splits a WFDB path (at
spaces, semicolons, and colons that are not part of a URL), and accepts a
single directory (\code{setwfdb(".")}); a path of several directories is an
error (\code{setwfdb: WFDB-x reads from a single input directory; use
wfdbx\_set\_input\_dir (or legacy mode for the WFDB search path)}), and
nothing changes. A directory whose name contains a space must therefore be
set with \code{wfdbx\_set\_input\_dir}, which takes the name as it is.
\code{getwfdb} returns \code{.} when no input directory is set.
\code{resetwfdb} restores the directory set by the last \code{setwfdb}.

\section{Record names}
\label{sec:record-paths}

A record name is given without the extension \file{.hea} of its header
(\code{100s}, not \code{100s.hea}; a trailing \file{.hea} is removed by the
functions that read a record and by \code{write\_record}). It is one of:
\begin{itemize}
\item a full path, absolute (\code{/data/mitdb/100}) or in the home
  directory (\code{\textasciitilde/Documents/100s}), used as given;
\item a relative path (\code{100s}, \code{data/100s}), read from the input
  directory if one is set, otherwise from the current directory, and
  written to the output directory if one is set, otherwise to the current
  directory (missing directories are created);
\item a URL (\code{https://physionet.org/files/mitdb/1.0.0/100}), read only
  if network access is on (\cref{sec:network}); a URL cannot be written.
  Only \code{http://} and \code{https://} URLs can be read: a name with
  another URL scheme, such as \code{ftp://} or \code{sftp://}, is an error,
  and is never read as a local file name.
\end{itemize}
A writer given a record name with a path (\code{out/rec}) creates the files
in that directory, and the header names the record and its signal file
without the path (\code{rec}, \code{rec.dat}), so that the record opens
under any path that leads to it (\cref{sec:resolve}). With the
\code{dir} option of a writer, the record name must be a plain name
(\code{rec}), without a directory: the files are created in \code{dir},
which must exist, and the record opens from that directory under its plain
name. The characters allowed in record names are given in
\cref{sec:record-names}.

\section{How a name is resolved}
\label{sec:resolve}

The name of a file is the record name followed by a dot and the file type
(\file{100s.hea}, \file{100s.atr}), or a file name given by itself (a signal
file named in a header, a calibration file). It designates exactly one
location, as \cref{tab:resolve} shows.

\begin{table}[htbp]
\caption{Where a name is read from and written to.}
\label{tab:resolve}
\begin{tabularx}{\linewidth}{@{}P{4.4cm}LL@{}}
\toprule
{\tabhead The name is} & {\tabhead Read from} & {\tabhead Written to} \\
\midrule
an absolute path: \code{/data/100}; on MS-Windows also
  \code{C:\textbackslash{}data\textbackslash{}100}, \code{C:/data/100} or
  \code{\textbackslash\textbackslash{}server\textbackslash{}share\textbackslash{}100} &
  that path & that path \\
a home path: \code{\textasciitilde/data/100}, \code{\textasciitilde
  user/data/100} & the home directory (expanded by \wfdbx{}) & the same \\
a URL: \code{http://\ldots} or \code{https://\ldots} & that URL, if network
  access is on & (not possible) \\
a name with another URL scheme: \code{ftp://\ldots}, \code{sftp://\ldots},
  \code{file://\ldots} & (not possible: an error) & (not possible) \\
\code{-} (as a record name) & the standard input & the standard output \\
anything else: \code{100}, \code{mitdb/100}, \code{../other/100} & the input
  directory followed by the name & the output directory followed by the
  name \\
\bottomrule
\end{tabularx}
\end{table}

\begin{itemize}
\item A leading \code{\textasciitilde} or \code{\textasciitilde/} is
  replaced by the home directory (the environment variable \code{HOME}, or
  the home directory of the current user if it is not set; on MS-Windows,
  \code{USERPROFILE}), and, on POSIX systems, \code{\textasciitilde
  user} by the home directory of that user. The input and output
  directories may themselves be \code{\textasciitilde} paths.
\item The input or output directory and the name are joined with one
  directory separator (\code{/} in a URL). \code{..} means one level up;
  the C library refused every name that contains \code{..}.
\item Nothing else is tried: a file that is not at that location is not
  found, and the operation fails (for example, \code{init: can't open header
  for record 100}). The only other name tried at the same location is the
  one whose file type is cut to its first three characters, as in C.
\item When an output file cannot be created because a directory of its path
  does not exist, the missing directories are created, as in C.
\end{itemize}
A record written under a path, such as \code{out/rec}, gets its files in
\file{out/} (under the output directory, if it is relative). Only the
record's own name, the last component, must be made of the characters of
record names (\cref{sec:record-names}); the directories may contain any
character. The header names the record, its signal files and its segments
relative to its own directory (\code{rec}, \code{rec.dat}), which is where
they are read from; the C library wrote the names as given, with their path
(\code{out/rec} and \code{out/rec.dat} in \file{out/rec.hea}), as legacy
mode does, so that such a record could not be read back by finding its
files relative to its header.

\subsection{Files named in a header}
\label{sec:header-files}

The files that a header names, the signal files of a record and the headers
of the segments of a multi-segment record, are found relative to the
directory of the header, as they are written in it (\code{100s.dat}, or
\code{data/100s.dat}): on disk, or at the URL from which the header was read.
A file missing there is an error, not a search elsewhere. Such a name must
stay inside the header's directory: a header that names a file with a
\code{..} component, an absolute path, a \code{\textasciitilde} path or a
URL (or, on MS-Windows, a drive letter) is rejected, with the message
\begin{lstlisting}[style=text]
init: file NAME named in the header of record REC is not in the header's directory
\end{lstlisting}
The signal file name \code{-} (the standard input) is accepted, and the null
segment \code{\textasciitilde} of a multi-segment record keeps its meaning.
For a header read from the standard input, which has no directory, the
files it names are resolved as any relative name, under the input
directory. This rule keeps a downloaded or crafted header from making the
library read other local files or contact another server.

\subsection{Files that a tool opens itself}
\label{sec:tool-files}

A few tools open a file themselves, by the name given on the command line,
instead of through the library. Such a file is found, or created, exactly
where the library finds, or creates, a file of that name: a relative name
under the input directory (a file that the tool reads) or the output
directory (a file that it creates, whose missing directories under the
output directory are created), a \code{\textasciitilde} name expanded, and
an absolute name as given; \code{-}, where the tool takes it for the
standard input or output, keeps that meaning. A URL is not accepted for
these files, also when the input directory is a URL: the tool reports that
it cannot read, or create, the file, as for a file that does not exist.
Messages name the file as it was given. \Cref{tab:tool-files} lists these
files.

\begin{table}[htbp]
\caption{The files that the tools open themselves.}
\label{tab:tool-files}
\begin{tabularx}{\linewidth}{@{}P{3.4cm}LL@{}}
\toprule
{\tabhead Tool} & {\tabhead Reads} & {\tabhead Creates} \\
\midrule
\cmd{wrsamp}, \cmd{edf2mit}, \cmd{wav2mit}, \cmd{a2m}, \cmd{ad2m} & the
  \opt{-i} file & \\
\cmd{ahaecg2mit} & the input files (\file{.ecg}, \file{.txt}) & \\
\cmd{xform} & the script of \opt{-S} & \\
\cmd{wfdbdesc} & the first signal file, to compute the length of a record
  whose header gives none: in the header's directory, where the library
  reads it & \\
\cmd{mit2wav} & & the WAV file (\opt{-o}), completed where the library
  wrote its samples \\
\cmd{mit2edf} & & the EDF file (\opt{-o}, or \textit{RECORD}\file{.edf}) \\
\cmd{snip} & & the master header of a multi-segment excerpt (\opt{-m}) \\
\cmd{parsescp} & & the \file{.des}, \file{.ecg}, \file{.key} and
  \file{.txt} files \\
\bottomrule
\end{tabularx}
\end{table}

The library functions \code{wfdbx\_input\_path} and
\code{wfdbx\_output\_path} (\cref{sec:library-additions}) return these
locations, for any program that opens a file itself. The C tools opened
these files relative to the current directory, so that, with an input or
output directory set, a tool read its records from one place and its own
files from another (\cmd{mit2wav} left the WAV header in the current
directory and the samples, without a header, in the output directory).
Without an input or output directory, nothing changes, except that a
\code{\textasciitilde} name is expanded and that \cmd{wfdbdesc} looks for
the signal file of a record in another directory in that directory.

\subsection{Legacy mode}

In legacy mode (\cref{sec:legacy}), files are found as by the original
reference build of the C library, which has no network support: an input
file is searched for in each entry of the WFDB path in turn, given by the
\code{WFDB} environment variable or \code{setwfdb} (by default
\code{. /usr/local/database http://physionet.org/physiobank/database}), with
the \code{\%r} substitution of the record name and \code{@file} indirect
paths; absolute names are joined to the entries of the path (an empty entry,
as in \code{WFDB=";"}, makes names be used as given); names that contain
\code{..} are refused, and URLs are not opened. The input directory and
\code{WFDBX\_NETWORK} are not used. The output directory is used, as a plain
prefix of every output name, absolute names included. The tools open the
files that they read themselves (\cref{sec:tool-files}) as named, relative
to the current directory, as the original tools do, and the output
directory, if one is set, is a plain prefix of the names of the files that
they create themselves, as of every other output file.

\section{Network access}
\label{sec:network}

Network access is off by default. A URL is opened only when it has been
turned on, by \code{wfdbx\_set\_network(1)}, by \code{WFDBX\_NETWORK=1}, or
by the \code{network} option of a high-level reader. Otherwise the library
says once, the first time it is asked to open a URL:
\begin{lstlisting}[style=text]
WFDB-x: network access is disabled; set WFDBX_NETWORK=1 or call wfdbx_set_network(1) to open URL
\end{lstlisting}
and the file is not found. The files of PhysioNet are reached in the same
way as any other remote files: by setting the input directory to the URL of
the database, \code{https://\allowbreak physionet.org/\allowbreak files/}\emph{database}\code{/}%
\emph{version}\code{/}, and turning network access on. No URL is built in:
\wfdbx{} never contacts a server it was not told to contact. For example, to
print the first 10 seconds of record 100 of the MIT-BIH Arrhythmia Database:
\begin{lstlisting}[style=shell]
export WFDBX_INPUT_DIR=https://physionet.org/files/mitdb/1.0.0/
export WFDBX_NETWORK=1
rdsamp -r 100 -t 10
\end{lstlisting}
or, in Python:
\begin{lstlisting}[language=Python]
import wfdbx
wfdbx.wfdbx_set_network(1)
wfdbx.wfdbx_set_input_dir("https://physionet.org/files/mitdb/1.0.0/")
rec = wfdbx.read_record("100", sampto="0:10")
\end{lstlisting}
or, for one call only:
\begin{lstlisting}[language=Python]
rec = wfdbx.read_record("100", input_dir="https://physionet.org/files/mitdb/1.0.0/",
                        network=True)
\end{lstlisting}

Remote files can only be read. Each is downloaded in full the first time it
is opened, and kept for the rest of the process; a download that fails is
not kept, so the next time the file is opened it is requested again (a
server that was unavailable may have recovered). Redirects are followed, and
an HTTP status of 400 or more is a failure. A download gives up if
connecting takes 30 seconds, or if no data arrive for 30 seconds; there is
no limit on its total time. Each language downloads with its own means:
\begin{itemize}
\item Python: its standard \code{urllib};
\item C++ and Rust: the system \cmd{curl} program, which they run without a
  shell (a missing \cmd{curl} is reported when a URL is opened);
\item Julia: its standard \code{Downloads};
\item R and MATLAB/Octave: \cmd{curl} when it is installed, and otherwise
  R's \code{url()} or MATLAB's \code{webread} and \code{websave};
\item Kotlin: \code{java.net.HttpURLConnection} on the JVM and Android, and
  \code{NSURLSession} on iOS.
\end{itemize}

Only \code{http://} and \code{https://} URLs can be read. A name that begins
with another URL scheme (\code{ftp://}, \code{ftps://}, \code{sftp://},
\code{file://}, \ldots) is refused, with or without network access, with a
message the first time:
\begin{lstlisting}[style=shell]
WFDB-x: can't open ftp://example.org/db/100.hea: only http:// and https:// URLs
can be read (mount the server as a local folder, or download the files)
\end{lstlisting}
(on one line). The same holds for an input directory with such a scheme,
a file named in a header, and the name of an output file; the files that a
tool reads or creates itself are refused as URLs are
(\cref{sec:tool-files}), and an output directory with a URL scheme is
refused (\cref{sec:directories}). Such a name is never read as a local file
name (the directory \file{ftp:} is never looked in). A scheme is a letter
followed by at least one letter, digit, \code{+}, \code{-} or \code{.},
then \code{://}, so a one-letter prefix, as in \code{c://x} (like a drive
letter), is not one. The records of a server reached by FTP or SFTP can be
read once the server is mounted as a local folder, by their local paths, or
after the files are downloaded:
\begin{lstlisting}[style=shell]
sshfs user@example.org:/data/db ~/mnt/db       # SFTP (or: rclone mount)
rdsamp -r ~/mnt/db/100 -t 10
\end{lstlisting}
In legacy mode, every name containing \code{://} is refused, without a
message, as by the original library built without network support.

A header read from a URL can be read but not changed by the metadata
functions, except into a local directory (\cref{sec:metadata-file}).

\section{The standard input and output}
\label{sec:stdio}

The record name \code{-} stands for the standard input, when a file is read,
and for the standard output, when one is written, whatever the type of the
file: a header, an annotation file or a signal file. A signal file named
\code{-} in a header is the standard input (or output). The tools follow the
C tools in this, so that, for example, \cmd{wrsamp} reads its text input from
the standard input unless \opt{-i} is given. The metadata functions refuse
the record name \code{-}.

\section{Environment variables}
\label{sec:environment}

\Cref{tab:environment} lists every environment variable that \wfdbx{} reads.
The library reads each one when it needs it.

\begin{xltabular}{\linewidth}{@{}P{3.9cm}L@{}}
\caption{The environment variables used by \wfdbx{}.}
\label{tab:environment}\\
\toprule
{\tabhead Variable} & {\tabhead Meaning} \\
\midrule
\endfirsthead
\toprule
{\tabhead Variable} & {\tabhead Meaning} \\
\midrule
\endhead
\bottomrule
\endfoot
\code{WFDBX\_INPUT\_DIR} & The input directory (\cref{sec:directories}):
  a local directory, a \code{\textasciitilde} path or a URL. Not used in
  legacy mode. \\
\code{WFDBX\_OUTPUT\_DIR} & The output directory: a local directory or a
  \code{\textasciitilde} path. \\
\code{WFDBX\_NETWORK} & Network access is on if this is a nonzero number
  (\code{WFDBX\_NETWORK=1}). Not used in legacy mode. \\
\code{WFDB} & The WFDB path, used only in legacy mode. Outside legacy mode
  it is ignored, and if it is set the library says so once:
  \code{WFDB-x: the WFDB environment variable is ignored; set
  WFDBX\_INPUT\_DIR instead (or use legacy mode)}. In legacy mode,
  \code{setwfdb} sets it (and sets \code{WFDBCAL}, \code{WFDBANNSORT} and
  \code{WFDBGVMODE} to their defaults if they are not set), as the C library
  does, so that the programs it starts inherit them. \\
\code{WFDBANNSORT} & If it is \code{0}, an output annotation file that is
  not in canonical order is not sorted when it is closed; the library then
  prints the advice \code{Use the command: sortann \ldots} (\cref{sec:sorting}).
  The \code{sort} option of the high-level \code{write\_annotations} decides
  for its call, whatever this variable says. (\code{WFDBNOSORT}, named by
  an old manual page of \cmd{sortann}, is not used, as in C.) \\
\code{WFDBCAL} & The name of the calibration file that \code{calopen} reads
  when it is given no name; \file{wfdbcal} if it is not set
  (\cref{sec:calibration}). \\
\code{WFDBGVMODE} & The \code{getvec} mode used when \code{setgvmode} is
  given a negative mode, and by \cmd{sigamp} when its options do not set the
  mode: \code{0} for low resolution, \code{1} for high resolution, plus
  \code{2} to replace invalid samples by the previous valid ones
  (\cref{sec:frames}). \\
\code{WFDB\_FLAC\_CHECK\_MD5} & If it is set to a nonempty value, the MD5
  signature of each FLAC signal file that is read is verified when the file
  is closed; a mismatch gives \code{isigclose: warning: incorrect MD5 hash
  in FLAC input}. \\
\code{HOME}, \code{USERPROFILE} & The home directory, for names that begin
  with \code{\textasciitilde} (\code{USERPROFILE} on MS-Windows). \\
\code{PATH} & Where the system \cmd{curl} program is found (C++, Rust, and
  R and MATLAB/Octave), and, in legacy mode, the \cmd{sortann} program that
  the library runs to sort an annotation file (C++ and Rust look first in
  the directory of the running program). \\
\code{TZ} & On POSIX systems, the time zone of the local time, which the
  library uses where the C library does: for the current date and time
  (\code{setbasetime} with an empty string, and \cmd{mit2edf} for a record
  without a start time). \\
\code{WFDBX\_DEBUG} & (R and Kotlin tools.) If it is set, a tool that stops
  on an internal error (reported as \code{wfdbx: error: \ldots}) also prints
  the call stack. \\
\code{WFDBX\_STDIN} & (R and MATLAB/Octave.) The name of a file that is
  read as the standard input: in R, for the tests; in MATLAB/Octave, also
  set by the MATLAB driver \file{matlab/bin/wfdbx\_matlab.sh}, since MATLAB
  cannot read its standard input. \\
\code{WFDBX\_STDOUT}, \code{WFDBX\_STDERR} & (MATLAB/Octave.) Files to which
  the standard output and error are written, set by the MATLAB driver,
  which copies them to the real standard output and error when MATLAB
  exits. \\
\end{xltabular}

No environment variable turns legacy mode on: it is turned on for a tool by
its \opt{--legacy} option, and in a program by \code{wfdbx\_set\_legacy(1)}
or the \code{legacy} option of a high-level function (\cref{sec:legacy}).
The other variables that the C library reads are not used by \wfdbx{}: those
of its libcurl code (\code{WFDB\_PAGESIZE}, \code{WFDBPASSWORD},
\code{CURL\_CA\_BUNDLE}, \code{WFDB\_NET\_DEBUG}) and those that tune
libFLAC's encoder (\code{WFDB\_FLAC\_COMPRESSION\_LEVEL},
\code{WFDB\_FLAC\_BLOCK\_SIZE}, and the other \code{WFDB\_FLAC\_\ldots}
variables).

The conformance suite runs every case with \code{LC\_ALL=C} and \code{TZ=UTC},
and without the variables \code{WFDB}, \code{WFDBCAL}, \code{WFDBANNSORT},
\code{WFDBGVMODE}, \code{WFDBX\_INPUT\_DIR}, \code{WFDBX\_OUTPUT\_DIR} and
\code{WFDBX\_NETWORK}, so that its results do not depend on them.

% Part I, chapter 4: Differences from the original package.
\chapter{Differences from the original package}
\label{chap:behavior}

\wfdbx{} reproduces the WFDB Software Package \wfdbversion{} exactly, with
the deliberate exceptions described in this chapter. It has no built-in
signal parameters; it finds files only where they are named, and opens URLs
only when network access is on; it sorts out-of-order annotation files
itself, without losing annotations; it never writes a signal file over a
signal file it is reading; its tools report write errors in their exit
status; its \code{getinfo} finds the record line of a header as the header
parser does; and it gives a defined result for input on which the C package
crashes. It also adds features that the C package does not have.

These exceptions are the same in all seven implementations, and are tested
like everything else: the reference results of the conformance suite come
from the C package with the patch \file{tests/reference/divergences.patch}
applied, which is the precise definition of the differences. The original
behavior remains available in every case, on request, in \emph{legacy mode}
(\cref{sec:legacy}). The file \file{docs/DIVERGENCES.md} gives the reasons
for each difference in more detail, and \cref{sec:corrections-1-2-0,sec:corrections-1-2-2} list
the corrections of versions 1.2.0 and 1.2.2.

\section{No built-in sampling frequency, gain, ADC resolution or units}
\label{sec:no-defaults}

The C package falls back on built-in values for four parameters of a
signal whenever a record or the user does not give them
(\cref{tab:c-defaults}). These values describe the equipment with which the
first PhysioNet databases were recorded, and do not describe modern
recordings. A silently assumed value corrupts every time, interval,
resampled signal and physical value computed from it, and nothing in the
output says that it was assumed. \wfdbx{} therefore never assumes any of
these four parameters: it uses the values given in the record, or values
given explicitly by the user, and reports an error when an operation needs a
value that is not available. Where an output must still show missing units,
it shows \code{nd} (``not defined'', a notation that the C package already
used for empty units) rather than \code{mV}.

\begin{table}[htbp]
\caption{The values that the C package assumes.}
\label{tab:c-defaults}
\begin{tabularx}{\linewidth}{@{}P{3cm}L@{}}
\toprule
{\tabhead Parameter} & {\tabhead Built-in value in C} \\
\midrule
sampling frequency & 250\,Hz (\code{WFDB\_DEFFREQ}); 500\,Hz in
  \cmd{parsescp} \\
gain & 200\,adu/mV (\code{WFDB\_DEFGAIN}); 400\,adu/mV for AHA files; 12.5 or
  6400\,adu/mV for WAV files \\
ADC resolution & 12 bits (\code{WFDB\_DEFRES}), or the storage width of the
  format \\
units & millivolts \\
\bottomrule
\end{tabularx}
\end{table}

\paragraph{Library.}
\begin{itemize}
\item A header that does not give a sampling frequency leaves it
  unspecified: \code{sampfreq} returns 0 (C: 250), unless a frequency was set
  with \code{setsampfreq} before the record was opened (as in C).
\item A header that does not give an ADC resolution leaves it unspecified
  (0); C inferred it from the format or assumed 12 bits. Missing units remain
  unspecified (\code{NULL}), as in C, but are never taken to be millivolts.
\item The time conversions that need the frequency fail when it is
  unspecified: \code{strtim} reports \code{strtim: sampling frequency is not
  specified} and returns 0 for the time strings other than sample numbers
  (\code{s\ldots}), \code{e}, \code{i} and \code{o} (and counter values when
  there is no counter frequency); \code{timstr} and \code{mstimstr} report
  the same error and return the time as a sample number (\code{s} followed
  by the number). C used 1 sample per second. The new function
  \code{wfdbx\_time\_needs\_freq(string)} tells whether \code{strtim} needs
  the frequency for a string (\cref{sec:times}).
\item \code{setifreq} fails if the input frequency is unspecified
  (\code{setifreq: sampling frequency is not specified}).
\item \code{aduphys}, \code{physadu}, \code{adumuv} and \code{muvadu} fail
  for an uncalibrated signal (gain 0) or an invalid signal number: they
  report the error and return NaN (\code{aduphys}) or
  \code{WFDB\_INVALID\_SAMPLE} (the others). C assumed a gain of 200, and a
  baseline of 0 for invalid signal numbers.
\item The constants \code{WFDB\_DEFFREQ}, \code{WFDB\_DEFGAIN} and
  \code{WFDB\_DEFRES} remain defined, for compatibility, but the library does
  not use them (except in legacy mode).
\end{itemize}

\paragraph{Tools.} A tool that writes these parameters into a record
requires them; a tool that needs a parameter missing from its input stops
with an error only when the requested operation needs it:
\begin{itemize}
\item \cmd{wrsamp} requires \opt{-F}, \opt{-G} and the new option
  \opt{-R BITS} when it writes a record (\opt{-o}), and the new option
  \opt{-U UNITS} for calibrated signals unless the input gives the units;
  \opt{-G 0} declares the signals uncalibrated. \opt{-G}, \opt{-R} and
  \opt{-U} take one value per signal (the last applies to the rest).
\item \cmd{wav2mit} has the new options \opt{-G GAIN} (required) and
  \opt{-U UNITS} (required unless the gain is 0); \cmd{ad2m} and
  \cmd{ahaecg2mit} the new options \opt{-F FREQ}, \opt{-G GAIN},
  \opt{-R BITS} and \opt{-U UNITS}, all required (\opt{-U} unless the gain
  is 0).
\item \cmd{a2m}, \cmd{rr2ann}, \cmd{rdann}, \cmd{ann2rr}, \cmd{sortann} and
  \cmd{m2a} take \opt{-F FREQ} (new except in \cmd{rr2ann}), used when the
  record has no header or its header has no frequency, and required then
  when times are converted.
\item \cmd{rdsamp} needs the frequency for time arguments, time columns and
  XML output, and the gain of each selected signal for physical units
  (\opt{-p}); \cmd{sigamp}, \cmd{snip}, \cmd{wfdb2mat}, \cmd{mit2wav},
  \cmd{wfdbtime}, \cmd{xform} and \cmd{mit2edf} need the frequency (and
  \cmd{mit2edf} the gain and ADC resolution of every signal), and so does
  \cmd{time2sec} with \opt{-r} (since 1.2.2; C converted the time at
  250\,Hz, and \wfdbx{} until 1.2.1 printed a meaningless number,
  \code{-9223372036854775808}, with the exit status 0).
\item \cmd{wfdbdesc} prints \code{not specified} for an unspecified
  frequency and ADC resolution, \code{uncalibrated} for a gain of 0 (C:
  \code{uncalibrated; assume 200}), and \code{adu/nd} for unspecified units;
  \cmd{sampfreq} prints 0 for a record whose header gives no frequency (C:
  250); \cmd{parsescp} writes the frequency given in the SCP file (C always
  wrote 500\,Hz).
\end{itemize}
The messages are \code{TOOL: sampling frequency of record RECORD is not
specified} (followed by \code{ (use -F)} for the tools that have that
option), \code{TOOL: signal N of record RECORD is uncalibrated},
\code{TOOL: ADC resolution of signal N of record RECORD is not specified},
and \code{TOOL: the \ldots{} must be specified (-X)} for a missing required
option. \Cref{part:reference} gives the options of each tool.

\paragraph{High-level API.} \code{write\_record} requires the sampling
frequency (greater than 0), the gains, the ADC resolutions, and the units of
the calibrated signals; a gain of 0 declares a signal uncalibrated, which is
allowed only for samples in ADC units. \code{read\_record} reports an
unspecified frequency as 0, and fails for time strings that need it; with
physical units, it fails if a selected signal is uncalibrated.
\code{read\_annotations} takes an optional sampling frequency, used, like
the \opt{-F} option of the tools, when the record has no header or its
header gives none.

\paragraph{Reproducing the results of the original package.} The results of
the original tools are obtained by giving the values they assumed,
explicitly. For example, \cmd{wrsamp} writes a record as the original did
with \opt{-F 250 -G 200 -U mV} and \opt{-R} with the storage width of the
output format (\opt{-R 12} for format 212, \opt{-R 16} for formats 16, 61,
160 and 516, \opt{-R 8} for 8, 80 and 508, \opt{-R 10} for 310 and 311,
\opt{-R 24} for 24 and 524, \opt{-R 32} for 32); \cmd{wav2mit} with
\opt{-G 12.5 -U mV} for 8-bit files and \opt{-G 6400 -U mV} for 16-bit
files; \cmd{ad2m} and \cmd{ahaecg2mit} with \opt{-F 250 -G 400 -R 12 -U mV};
the annotation tools with \opt{-F 250} for a record without a frequency; and
a library program calls \code{setsampfreq(250)} before it opens a record
whose header gives no frequency. One difference remains by construction: a
header written with \opt{-U mV} says \code{200/mV} where the C tools wrote
\code{200}.

\paragraph{What does not change.} Values defined by a data format, not by an
assumption about the signal (the record lengths and time offsets of the AHA
database formats, and \cmd{md2a}'s default length); values read from the
input (the sampling frequency and resolution of a WAV file, the parameters
of an EDF file); and analysis settings that are documented options, such as
the measurement windows of \cmd{sigamp}. \cmd{parsescp} writes a gain of
200\,adu/mV because it rescales the SCP samples to units of 5\,$\mu$V, so
that value is exact, not assumed.

\section{Files are found only where they are named}
\label{sec:named-files}

The C library found a file by searching the WFDB path, a list of places
tried in order until the file was found, by default the current directory,
\file{/usr/local/database} and, in builds with network support, PhysioNet.
That was convenient, but:
\begin{itemize}
\item a search opens the first copy it finds, so a stale or partial copy in
  a directory earlier on the path is read instead of the intended one, and
  a signal file missing next to its header is replaced by a file of the same
  name found anywhere else on the path;
\item whether a record is read from disk or downloaded, and from where,
  depends on the environment, not on the program, and a missing local file
  turns into a network request instead of an error;
\item for every record it read, the library also asked the last entry of the
  path, PhysioNet, for an optional \file{.info} file when there was no local
  one, with the full local path of an absolute record name in the request;
\item a header could make the library read any local file
  (\code{../../.ssh/id\_rsa}) or contact a third-party server, through the
  names of its signal files;
\item the default URL is no longer where PhysioNet serves its files.
\end{itemize}
\wfdbx{} reads and writes every file at the one location that its name
designates, as \cref{chap:files} describes: under the input or output
directory for a relative name, as given for an absolute path, a
\code{\textasciitilde} path or a URL, and relative to the header for the
files that a header names, which must stay in the header's directory. URLs
are opened only when network access is on, and only \code{http://} and
\code{https://} URLs can be read. A name containing \code{..} is
accepted. \code{setwfdb} and \code{getwfdb} set and return the input
directory, and the \code{WFDB} environment variable is used only in legacy
mode. Reading PhysioNet records must therefore be asked for explicitly
(\cref{sec:network}); in return, every program reads exactly the files it
names, the same on every machine.

The files that a tool opens itself, by the name given on its command line,
rather than through the library (the \opt{-i} file of \cmd{wrsamp}, the WAV
file of \cmd{mit2wav}, and the others of \cref{tab:tool-files}), follow the
same rules: the C tools opened them relative to the current directory, so
that, with an input or output directory set, a tool read its records from
one place and its own files from another. In \wfdbx{} such a file is found,
or created, where the library finds, or creates, a file of that name
(\cref{sec:tool-files}); a URL is not accepted for these files. In legacy
mode, the files that a tool reads itself are opened as named, and the
output directory, if one is set, is a plain prefix of the names of the
files that it creates itself.

\section{Out-of-order annotation files are sorted in the library}
\label{sec:sorting}

When an annotation file written by the library is not in canonical order
(\cref{sec:annotations}), the C library sorted it when the file was closed by
running the \cmd{sortann} program through the shell, unless the environment
variable \code{WFDBANNSORT} was 0. The result then depended on which
\cmd{sortann} was found on the \code{PATH}, and a record or annotator name
with a space or a shell character broke the command.

\wfdbx{} sorts the file itself, in the same process. The file it writes is
byte for byte the one that \cmd{sortann} writes, in the same order (by time,
then annotation number, then channel; a later annotation with the same
time, number and channel replaces the earlier one), with the same treatment
of \code{NOTQRS} annotations (they are dropped) and of the time resolution,
and the same encoding. The messages and status of the library are the same:
\code{Rearranging annotations for output annotator ANN ...done!}, or the
advice \code{Use the command: sortann -r REC -a ANN} when sorting is turned
off or fails; only the messages that the separate \cmd{sortann} printed
itself are gone. \code{WFDBANNSORT=0} turns sorting off, and the
\code{sort} option of the high-level \code{write\_annotations} decides for
its call alone (sort, or leave the file as written), whatever
\code{WFDBANNSORT} says.

\paragraph{Sorting does not lose annotations.} The original \cmd{sortann}
silently loses annotations: when an annotation at time 0, with annotation
number 0 and channel 0, arrives after later ones, it is dropped, and so are
the annotations later inserted directly behind it (input sample times 5, 0,
3 give a file that contains only 5), because the search for the insertion
point mistakes the empty head entry of its list for an annotation. \wfdbx{}
puts such an annotation at the beginning of the list, where it belongs (5,
0, 3 give 0, 3, 5). Nothing else about the sorting changes. This applies to
the \cmd{sortann} tool and to the sorting done by the library.

\begin{note}
In legacy mode, the library runs \cmd{sortann} as the original does (with
its own \cmd{sortann} tool, started with \opt{--legacy}, without a shell),
and both lose these annotations as the original does. Legacy mode should not
be used to sort annotation files unless reproducing the original results,
losses included, is the point: nothing in the output says that annotations
were dropped.
\end{note}

\section{An output signal file is never one of the input signal files}
\label{sec:output-input}

The files of records can be shared: the segments of a multi-segment record
are other records, so that \cmd{xform -i vlay -o fmt16} reads
\file{fmt16.dat} (a segment of \code{vlay}) and writes \file{fmt16.dat}
(the output record \code{fmt16}). The original opens the output file, which
truncates it, while still reading it: the input is destroyed, and what is
read depends on how much of it the C library had read ahead, so that the
result differs between platforms.

\wfdbx{} refuses to create an output signal file that is the same file as a
signal file of the record open for reading (the same file, not the same
name: links and different spellings of a path are recognized), with the
message \code{osigfopen: output file NAME is also an input file}
(\code{osigopen:} when the file is opened by \code{osigopen}); nothing is
written. The check is in the library, so every tool and every program that
uses the library functions is protected. The high-level \code{write\_record}
reads nothing: it closes any open record first, and simply replaces
existing files of the same name. The check covers signal files only: for
example, \cmd{snip -i 100s -n 100s} still writes the header of its output
record over that of its input record, in C and in \wfdbx{} alike; the output
record needs a name of its own.

\section{Write errors give a nonzero exit status}
\label{sec:write-errors}

The original tools report a file that cannot be written (a full disk, for
example) only through the library's messages, such as \code{osigclose:
write error in signal group 0} or \code{oannclose: write error on annotation
file atr}, and then exit with status 0; they say nothing at all when the
standard output cannot be written, or when a file that a tool writes itself
fails. The tools of \wfdbx{} follow the convention of the standard Unix
tools: when an output file or the standard output cannot be written, the
exit status is 1. The library's messages are unchanged; a standard output
that cannot be written is reported as \code{PROGRAM: write error on standard
output}, and a file that \cmd{snip}, \cmd{mit2edf} or \cmd{parsescp} writes
itself as \code{PROGRAM: write error on FILE}. A tool that exits with
another nonzero status for another reason keeps that status, and the tools
keep running after a write error, as the original ones do.

The library counts the output files that it could not write, flush or close;
\code{wfdbx\_write\_errors()} returns that count, and
\code{wfdbx\_note\_write\_error()} adds one, for a program that writes a file
itself.

\section{The record line of a header is found as the header parser finds it}
\label{sec:record-line-divergence}

The info strings of a header are its lines that begin with \code{\#} after
the record line (\cref{sec:info-strings}). The C package has two rules for
finding the record line. The header parser (\code{readheader}) takes the
first line whose first token (separated by spaces, tabs, carriage returns
and line feeds) does not begin with \code{\#}, so it skips empty and blank
lines as well as comment lines; but \code{getinfo} takes the first line that
does not begin with \code{\#}. For a header whose first line is empty or
blank, such as
\begin{lstlisting}[style=text]
(an empty line)
# Age: 40
rec 1 250 10
rec.dat 16 200 12 0 0 0 0 ECG
# Age: 63
\end{lstlisting}
the original \code{getinfo} takes the empty line for the record line, and
returns \code{Age: 40}, a comment before the record line, as an info string,
as well as \code{Age: 63}: \cmd{wfdbdesc} prints both as notes, and the
\code{comments} of the record hold both.

Since version 1.2.2, the \code{getinfo} of \wfdbx{} finds the record line as
\code{readheader} does, so that the info strings are the lines that the
header parser leaves after the record line: \code{Age: 63} only. A header
with no record line has no info strings. The info strings after the record
line are found as before: every line whose first character is \code{\#} (a
comment between two signal lines included). This applies to \code{getinfo}
itself, to the \code{comments} of \code{read\_header} and
\code{read\_record}, and to the notes of \cmd{wfdbdesc}. The metadata
functions (\code{read\_metadata} and the others, and the \code{metadata} of
\code{read\_header}) follow the same rule in every mode, legacy mode
included (\cref{chap:metadata}). In legacy mode, \code{getinfo} follows the
original rule. The conformance cases \code{edge-hdr-leading-*} check both (in
the legacy suite, against the unmodified package).

\section{Input on which the C package crashes}
\label{sec:crashes}

Some malformed or unusual input makes the C package crash, run out of
memory, loop forever, or give a result that depends on the machine (reads of
uninitialized memory). A crash cannot be reproduced, so \wfdbx{} gives each
such input one defined, safe result, the same in all seven implementations,
and legacy mode gives the same result. For example:
\begin{itemize}
\item a header with a negative or impossibly large number of signals
  (\code{n -1 250 10}), for which C tries to allocate memory for 4,294,967,295
  signals, gives \code{init: unexpected EOF in header file for record n}, as
  C does for a header with fewer signal lines than it declares;
\item an EDF file whose signals have 0 samples per data record, on which C
  crashes with a floating-point exception, reads as a record with no
  samples;
\item a FLAC signal file with 3 or more samples per frame whose FLAC block
  size is not a multiple of them, for which C writes past the end of a
  buffer, is read correctly;
\item \cmd{wfdb2mat} with a one-character record name, for which C
  overwrites the header of its input record, names its output record as for
  longer names (\code{mm} for \code{m});
\item \cmd{edf2mit} with a file shorter than the 256-byte fixed part of the
  EDF header prints \code{edf2mit: incomplete EDF header in FILE} and exits
  with status 2, creating no files;
\item \cmd{signame -s N}, where \code{N} is the number of signals, prints
  \code{[INVALID]}, as its manual page says;
\item \code{osigopen(record, NULL, 0)}, for which C reads before the start of
  an array, closes the output signals already open (unless the record name
  begins with \code{+}), reads the header (an error if it cannot be read),
  opens nothing and returns 0.
\end{itemize}
The complete list, with the C behavior and the \wfdbx{} choice for each, is
the table ``Undefined behaviors (not reproduced)'' of
\file{docs/KNOWN\_C\_BEHAVIORS.md}. The conformance suite checks most of
these choices, with reference results recorded from \wfdbx{} instead of C.

\section{Additions}
\label{sec:additions}

These are new features, not changes of the original behavior, so they have
no legacy counterpart.
\begin{itemize}
\item \textbf{The high-level API}: in every language, a few functions that
  read and write whole records and annotation files as arrays, in one call,
  with named options and errors reported in the language's way
  (\cref{part:languages}).
\item \textbf{An output directory} (\code{wfdbx\_set\_output\_dir},
  \code{WFDBX\_OUTPUT\_DIR}, and the \code{dir} option of the high-level
  writers), in which output files are created, and an input directory
  (\code{wfdbx\_set\_input\_dir}, \code{WFDBX\_INPUT\_DIR}); see
  \cref{chap:files}.
\item \textbf{A record checker}, \code{check\_record(record, annotators)},
  which checks a record's header, signal files and annotation files, and
  reports each problem it finds with a severity (error, warning or
  information): missing or inconsistent signal specifications, ADC
  resolutions and ranges that the storage format cannot hold, samples outside
  the ADC range, checksum and length mismatches, problematic names,
  inconsistent segments of multi-segment records, and annotations that are
  out of order, undefined or after the end of the signals. Its report,
  including the text of every message, is identical in all seven
  implementations. It looks for problems similar to those found by Benjamin
  Moody's separate C program \cmd{wfdbcheck}, but was written independently.
\item \textbf{Fast reading and writing.} The high-level functions read and
  write whole blocks of samples, and \code{write\_annotations} encodes a whole
  MIT annotation file at once. The multi-frame readers added to the library
  functions (\code{getframes} in Python, R and MATLAB/Octave, \code{getvecs}
  in C++, Rust, Julia and Kotlin) decode whole blocks as well, including
  multi-frequency and multi-segment records. The samples and files are
  exactly those of the frame-by-frame functions.
\item \textbf{Signal names without surrounding spaces.} Some headers end a
  signal's description with spaces (\code{"ABP "}). The high-level readers
  give the names without white space at either end, and \code{read\_record}
  finds a channel by that name; the library functions and the tools keep the
  descriptions as the C library does.
\item \textbf{The header and the start time as numbers.} \code{read\_header}
  returns the header of a record without opening its signal files: the
  sampling and counter frequencies, the base counter value, the number of
  samples, the signal specifications, the info strings, the metadata, and the
  start time and date. \code{read\_header} and \code{read\_record} give the
  start time as text (\code{base\_time}), in milliseconds after midnight
  (\code{base\_time\_ms}, which keeps the milliseconds of a start time such
  as \code{14:00:03.500}), and the start date as year, month and day
  (\code{base\_date}).
\item \textbf{Write errors and frame checks in the high-level writers.}
  \code{write\_record} and \code{write\_annotations} raise an error
  (\code{write error on record REC}, \code{write error on annotator ANN for
  record REC}) when a file cannot be written completely, or when
  \code{putann} reports an error (version 1.2.2). \code{write\_record}
  also refuses, before creating any file, frame sizes that are not integers
  of at least 1 (\code{samps\_per\_frame must hold integers >= 1}) and
  signals whose number of columns is not the number of samples in a frame
  (\code{signals has N columns, but samps\_per\_frame adds up to M}).
\item \textbf{The high-level writers check every argument before touching
  any file, and leave the existing files intact after a failure} (version
  1.2.2). \code{write\_record} refuses an array of signals that is not 1-D or
  2-D (\code{signals must be a 1-D or 2-D array}), baselines, ADC zeros and
  ADC resolutions that are not whole numbers within the range of a 32-bit
  integer (\code{baselines must be whole numbers within the range of a
  32-bit integer}), and a record name, signal names, units or
  comments that cannot be written as header text (in Python, a character
  that is not Latin-1). \code{write\_annotations} refuses an array of times
  (or of codes, subtypes, channels or numbers) that is not 1-D
  (\code{samples must be a 1-D array}), annotation codes outside 0 to 58
  (\code{annotation code 61 is reserved (59 to 63 are the control codes of
  the annotation file format)}, \code{annotation code 64 is out of range (0
  to 58)}: the codes 59 to 63 are the control codes \code{SKIP}, \code{NUM},
  \code{SUB}, \code{CHN} and \code{AUX} of the MIT format, and an annotation
  written with one of them corrupts the file), and the times $-2^{63}$ and
  $2^{63}-1$ (\code{annotation time T cannot be written (annotation times
  must be integers from -9223372036854775807 to 9223372036854775806)}),
  which \code{putann} refuses (\cref{sec:known-c}); every other time can be
  written. Both writers then write their files to temporary files in the
  directories of the destinations (\file{.NAME.XXXXXXXX.tmp}), which are
  renamed over the destinations once all of them are complete, and removed
  after a failure, so that a failure never leaves a truncated header, signal
  file or annotation file in place of an existing one. A destination that is
  not a regular file, such as a device, is written directly; for a symbolic
  link, the file it leads to is replaced, and the link is kept. The files
  written are unchanged, byte for byte, and keep the permissions of the files
  they replace. (In legacy mode, an annotation file that must be sorted is
  renamed over its destination before the \cmd{sortann} program sorts it
  there.) The library functions (\code{putann}, \code{putvec},
  \code{newheader}, \ldots) and the tools keep the behavior of C.
\item \textbf{Read errors in the high-level readers.} \code{read\_record}
  raises an error (\code{can't read signal data of record REC (getvec status
  S)}) when a signal file ends before the length its header gives, or cannot
  be read, instead of returning the samples read so far. A read failure is
  tracked separately from the checksum: when a signal file ends in the frame
  at which the checksum of a signal is verified, for which \code{getvec}
  returns $-4$ (checksum error) instead of $-3$, the high-level readers
  still raise the read error (\code{getvec status -3}) and never return an
  incomplete frame, and \code{check\_record} reports the file as shorter
  than its header says, without a checksum error (version 1.2.2).
  \code{read\_record} allocates at most $2^{26}$
  samples in advance, whatever length the header gives, and grows its
  result as the samples are read, so that a header that gives far more
  samples than its files hold ends in the read error, not in an attempt to
  allocate memory for them. \code{read\_annotations} checks
  \code{sampfrom} and \code{sampto} as \code{read\_record} does (whole
  numbers, \code{sampfrom} not negative), before the file is opened, and
  closes the file whatever happens. A checksum error alone (every sample
  read, but their sum differs from the header's checksum) is a warning
  instead (\cref{sec:errors}).
\item \textbf{Header metadata} (version 1.2.0): the functions
  \code{read\_metadata}, \code{get\_metadata}, \code{set\_metadata},
  \code{remove\_metadata} and \code{update\_metadata} read and change the
  \code{Key: value} info strings of a header without changing anything else
  in it, \code{read\_header} returns them, and \code{write\_record} writes
  them (\cref{chap:metadata}). The info strings are those that
  \code{getinfo} reads, after the record line found as the header parser
  finds it, in every mode (\cref{sec:record-line-divergence}). The library
  functions \code{putinfo} and \code{setinfo}, and the tools, are unchanged,
  and the C library reads the changed headers as before.
\item \textbf{The high-level functions leave the annotation time resolution
  and the getvec mode as they were} (version 1.2.0). Opening an EDF record
  sets the time resolution of output annotators (\code{setafreq}) to the
  record's sampling frequency when none is set, and adds
  \code{WFDB\_HIGHRES} to the getvec mode (\code{setgvmode}), as in C, and
  \code{wfdbquit} resets neither. \code{read\_record}, \code{read\_header},
  \code{read\_annotations}, \code{write\_record}, \code{write\_annotations}
  and \code{check\_record} restore both previous values (\code{getafreq},
  \code{getgvmode}) when they return, also after an error, so that
  \code{read\_header("100s.edf")} changes nothing for a later
  \code{write\_annotations} or \code{getvec}; the \code{fs} argument of
  \code{write\_annotations} applies to its own file only. The library
  functions keep the behavior of C (\cref{sec:known-c}).
\item \textbf{Library functions} for these settings: \code{wfdbx\_set\_legacy},
  \code{wfdbx\_get\_legacy} and \code{wfdbx\_legacy\_args} (legacy mode),
  \code{wfdbx\_time\_needs\_freq}, \code{wfdbx\_write\_errors} and
  \code{wfdbx\_note\_write\_error}, the functions of the input and output
  directories and of network access, and \code{wfdbx\_input\_path} and
  \code{wfdbx\_output\_path} (version 1.2.0), which give the path under
  which a program opens a file that it reads or creates itself
  (\cref{sec:tool-files}).
\end{itemize}
The high-level functions behave identically in all seven implementations,
in legacy mode too, since they are not part of the original package.

\section{Corrections of version 1.2.0}
\label{sec:corrections-1-2-0}

Version 1.2.0 changes the behavior of WFDB-x in three points, which were
errors of the earlier versions:
\begin{itemize}
\item \textbf{The files that a tool reads or creates itself follow the input
  and output directories} (\cref{sec:tool-files}). The \opt{-i} file of
  \cmd{wrsamp}, \cmd{edf2mit}, \cmd{wav2mit}, \cmd{a2m} and \cmd{ad2m}, the
  input files of \cmd{ahaecg2mit} and the \opt{-S} script of \cmd{xform} are
  read from the input directory; the WAV file of \cmd{mit2wav}, the EDF file
  of \cmd{mit2edf}, the master header of \cmd{snip -m} and the \file{.des},
  \file{.ecg}, \file{.key} and \file{.txt} files of \cmd{parsescp} are
  created in the output directory; \cmd{wfdbdesc} measures a signal file of
  unknown length in the header's directory. They were opened relative to the
  current directory: with \code{WFDBX\_OUTPUT\_DIR} set, \cmd{mit2wav} left a
  WAV header without samples in the current directory and the samples
  without a header in the output directory. In legacy mode, input files are
  opened as named, and the output directory is a prefix of every output
  name. New library functions: \code{wfdbx\_input\_path} and
  \code{wfdbx\_output\_path} (\cref{sec:library-additions}).
\item \textbf{\code{osigopen(record, NULL, 0)}} returns 0 in every language
  (\cref{sec:crashes}); the Python implementation raised an
  \code{IndexError}.
\item \textbf{The high-level functions leave the annotation time resolution
  and the getvec mode as they were} (\cref{sec:additions}). After
  \code{read\_header("100s.edf")}, or any high-level function that opens an
  EDF record, a later \code{write\_annotations} in the same process no longer
  records the EDF record's sampling frequency as the time resolution of its
  file, and the getvec mode no longer stays in high resolution. The library
  functions keep the behavior of C.
\end{itemize}

\section{Corrections of version 1.2.2}
\label{sec:corrections-1-2-2}

Version 1.2.2 corrects the following points, in all seven implementations.
The library functions and the tools keep the behavior of C, except for
\code{getinfo} and the URL schemes.
\begin{itemize}
\item \textbf{Truncated records.} A signal file that ends before the length
  that its header gives, or that cannot be read, is a read error of
  \code{read\_record} (\code{can't read signal data of record REC (getvec
  status -3)}) even when the checksum of a signal is verified in the
  incomplete frame. \code{getvec} returns the checksum status ($-4$) for that
  frame, which replaced the end of file ($-3$), so \code{read\_record}
  returned the incomplete frame as data: a format-16 file holding the
  samples 1 and 2, whose header gives 3 samples, was read as 1, 2, 257, with
  a checksum warning. The high-level readers now track a read failure
  separately from the checksum, and never return an incomplete frame;
  \code{check\_record} reports such a file as shorter than its header says.
  A checksum error alone is still a warning (\cref{sec:additions}).
\item \textbf{\code{read\_record} allocates no more than it can read.} Its
  result was allocated from the length that the header gives (a header
  giving $10^{12}$ samples for a file of 4 bytes attempted to allocate
  terabytes); at most $2^{26}$ samples are now allocated in advance.
\item \textbf{The high-level writers check everything before touching any
  file, and leave the existing files intact after a failure}
  (\cref{sec:additions}). \code{write\_record} refuses, before any file is
  opened, arrays that are not 1-D or 2-D, baselines, ADC zeros and ADC
  resolutions that are not whole numbers within the range of a 32-bit
  integer, and text that cannot be written in a header (in Python,
  \code{units="}$\mu$\code{V"} left the existing header empty).
  \code{write\_annotations} refuses arrays of times that are not 1-D (a
  2-D array truncated the existing file), the annotation codes 59 to 63 and
  any code outside 0 to 58, and the times $-2^{63}$ and $2^{63}-1$, which
  were dropped while success was reported. An error that \code{putann}
  reports is an error of \code{write\_annotations}. Both writers write
  through temporary files that replace the destinations only when all of
  them are complete, keeping their permissions and symbolic links.
\item \textbf{\code{read\_annotations} checks \code{sampfrom} and
  \code{sampto}} as \code{read\_record} does (a fraction, NaN or an
  infinite value is an error, before the file is opened: \code{sampfrom
  must be a whole number}, \code{sampfrom must not be negative}), and closes
  the annotation file whatever happens.
\item \textbf{The record line of a header} is found by \code{getinfo} as
  the header parser finds it, except in legacy mode, and by the metadata
  functions in every mode (\cref{sec:record-line-divergence}): for a header
  whose first line is empty or blank, the comments before the record line
  were returned as info strings (in \code{comments} and in the notes of
  \cmd{wfdbdesc}).
\item \textbf{A remote file that could not be read is requested again}
  (\cref{sec:network}). A download that failed (a server error, a timeout, a
  file not yet published) was remembered until the process ended, so the
  file could not be read even after the server recovered; only the files
  that have been read are now kept.
\item \textbf{Only \code{http://} and \code{https://} URLs}
  (\cref{sec:network}). A name that begins with another URL scheme
  (\code{ftp://}, \code{sftp://}, \code{file://}, \ldots) was opened as a
  local path (\code{ftp://h/db/100} as the directory \file{ftp:}, then
  \file{h/db/100}); it is now refused in normal mode, with or without
  network access, with the message \code{WFDB-x: can't open NAME: only
  http:// and https:// URLs can be read (mount the server as a local folder,
  or download the files)}, and so are an input directory, a file named in a
  header, the files that a tool reads or creates itself, and an output
  directory with such a scheme. Legacy mode is unchanged.
\item \textbf{\cmd{time2sec -r} needs the sampling frequency of the record}:
  a record whose header gives none, or cannot be read, is an error (exit
  status 2: \code{time2sec: sampling frequency of record RECORD is not
  specified}). It printed \code{-9223372036854775808} with the exit status 0
  (a division by the unspecified frequency, 0); the C tool converted the time
  at 250\,Hz, as legacy mode still does (\cref{sec:no-defaults}).
\item \textbf{\code{write\_annotations} requires a label for each
  annotation}, in every language: with one annotation or more and neither
  symbols nor codes, it is an argument error (\code{symbols or codes must be
  given}; in R, \code{symbol or code must be given}), before any file is
  touched. C++, Rust and Julia wrote such annotations with the code 0,
  which is not an annotation code; Python, R and Kotlin refused even an
  empty list, which now writes an empty file, as in the other languages.
\item \textbf{\code{write\_record} removes a trailing \file{.hea}} from the
  record name, as the readers do: \code{write\_record("x.hea", \ldots)}
  writes \file{x.hea} and \file{x.dat}. Python failed (\code{can't create
  signal file for record x.hea}), and the other languages named the signal
  file \file{x.hea.dat}.
\item \textbf{Metadata changes} keep the permissions of the header, whatever
  its file name, and their temporary file is never more open than the header
  (\cref{sec:metadata-writing}); corrections in R, MATLAB/Octave and
  Kotlin.
\item \textbf{Licenses.} The file \file{NOTICE} (the licenses of
  \wfdbx{}, and the BSD-3-Clause notice of libFLAC, whose stream decoder the
  FLAC decoders translate) is included in the package of every language,
  with the license files (\cref{sec:license}). The R package declares
  \code{License: file LICENSE}, and in MATLAB/Octave the dispatcher of the
  \cmd{wfdbx} command (\file{+wfdbx/+internal/cli\_main.m} and
  \file{app\_run.m}), which only the tools use, is under the GPL, as the
  launchers of the other languages are.
\item \textbf{C++ package.} The installation provides
  \file{wfdbxConfigVersion.cmake}, so that
  \code{find\_package(wfdbx 1.2.2 CONFIG REQUIRED)} checks the version, and
  the \code{wfdbx} target requires C++17 of the programs that link to it.
\end{itemize}

\section{Legacy mode}
\label{sec:legacy}

Legacy mode makes \wfdbx{} behave exactly as the original package, for
reproducing old results. The default is the behavior described in the rest
of this manual. It is turned on in three ways:
\begin{itemize}
\item \textbf{Tools}: \opt{--legacy} as the \emph{first} argument of any tool
  (\cmd{rdsamp --legacy -r 100s -t 10}; in Python \cmd{wfdbx rdsamp --legacy
  -r 100s -t 10}). It is removed before the options are parsed, and the tool
  then behaves exactly as the original one: same options (the options added
  by \wfdbx{}, such as \opt{-F} in \cmd{rdann} or \opt{-R} in \cmd{wrsamp},
  are unrecognized), outputs, messages, help texts and exit statuses.
  \opt{--legacy} anywhere else is not recognized.
\item \textbf{Library}: \code{wfdbx\_set\_legacy(1)} turns legacy mode on,
  and \code{wfdbx\_set\_legacy(0)} off; each returns the previous setting,
  and \code{wfdbx\_get\_legacy()} returns the current one. The setting applies
  to the whole library. \code{wfdbx\_legacy\_args} is the function that the
  tools call to recognize and remove \opt{--legacy}.
\item \textbf{High-level API}: the \code{legacy} option of \code{read\_record},
  \code{read\_header}, \code{read\_annotations}, \code{write\_record},
  \code{write\_annotations} and the related functions turns the library's
  legacy mode on for the duration of the call, and restores the previous
  setting afterwards. The functions also follow the original rules when the
  library is in legacy mode.
\end{itemize}
No environment variable turns legacy mode on.

In legacy mode, every library difference of this chapter behaves as in the
original: a header without a frequency means 250\,Hz; a missing ADC
resolution is inferred from the format or taken as 12 bits; the time
conversions use 1 sample per second when no frequency is known;
\code{setifreq} reports \code{no open input record}; the unit conversions
assume a gain of 200 (and a baseline of 0 for invalid signal numbers); files
are found through the WFDB search path, without network access
(\cref{sec:resolve}); out-of-order annotation files are sorted by running
\cmd{sortann}, which loses annotations as the original does; an output file
may be one of the input files; \code{getinfo} takes the first line that does
not begin with \code{\#} for the record line; the whole record name is
checked when a header is written, and headers name records and files as
given. The input directory and \code{WFDBX\_NETWORK} are not used, no remote
file is read (every name containing \code{://} is refused), and the tools
open the files that they read themselves as named; the output directory is
used, also for the files that the tools create themselves. The tools exit
with the original status after a write error, and print no new messages. In \code{write\_record}, legacy mode
gives a gain of 200 and makes the units optional; its default ADC
resolution is the storage width of the format (8 bits for formats 80 and
508, 10 for 310 and 311, 12 for 212, 24 for 24 and 524, 32 for 32) and 16
bits for every other format, which differs from \cmd{wrsamp --legacy} only
for format 8 (12 bits there, as in the original \cmd{wrsamp}).

Legacy mode is tested against the unmodified C package by the legacy suite
of the conformance tests (886 cases, each run with \opt{--legacy}), and the
library's legacy values by the unit tests.

\begin{note}
Legacy mode is kept for compatibility only. The values it assumes are
exactly the silent defaults that \wfdbx{} removed, so results produced with
it carry the same risks: a time, interval, resampled signal or physical value
computed from an assumed frequency, gain, resolution or unit is wrong if the
assumption is, and nothing in the output says so.
\end{note}

\section{Errors and warnings}
\label{sec:errors}

\paragraph{The library functions} report errors as the C library does, in
every language: a function returns an error status (a negative number, a
null pointer, \ldots, as its C counterpart does), and the library prints a
message on the standard error output, such as \code{init: can't open header
for record 100}. \code{wfdbquiet()} stops the messages from being printed,
and \code{wfdbverbose()} prints them again; \code{wfdberror()} returns the
last message (before any error, \code{WFDB library version 10.7.0 (Jun 10
2022).}), and \code{wfdb\_error(format, \ldots)} reports a message as the
library does. The messages are those of the C library, byte for byte, so
that the tools print what the C tools print.

\paragraph{The high-level functions} report errors in the usual way of each
language, with a message that says which operation failed, such as
\code{can't read record 100s}; the library's own message, if any, is printed
too, unless \code{wfdbquiet()} was called, and returned by
\code{wfdberror()}.
\begin{itemize}
\item Python raises \code{OSError} (\code{IOError}) when a file cannot be
  read or written, \code{ValueError} for an invalid argument,
  \code{TypeError} for an argument of the wrong type, and
  \code{FileNotFoundError} when an output directory (\code{dir}) does not
  exist.
\item C++ throws \code{wfdbx::Error}, or, in the metadata functions (and
  in \code{write\_record}, for invalid metadata), one of its subclasses
  \code{ArgumentError} and \code{FileError} (in the namespace
  \code{wfdbx}).
\item Rust returns a \code{Result} whose error type is
  \code{wfdbx::record::Error}.
\item Julia throws a \code{WfdbError}, or, for some invalid arguments of
  the metadata functions and of \code{check\_record}, an
  \code{ArgumentError}.
\item R signals an R error (in the metadata functions, of class
  \code{wfdbx\_argument\_error} or \code{wfdbx\_file\_error}).
\item MATLAB/Octave throws an error, with an identifier such as
  \code{wfdbx:read\_record}.
\item Kotlin throws a \code{WfdbException}, or, for an invalid argument of
  the metadata functions, the writers and \code{readAnnotations}, its
  subclass \code{WfdbArgumentException}.
\end{itemize}
The chapters of \cref{part:languages} list the errors of each function.

\paragraph{The checksum warning.} When a high-level reader has read every
sample of a signal but their sum differs from the checksum of the header,
it returns every frame, the last one included, with the warning
\begin{lstlisting}[style=text]
checksum error in record REC: the signal data do not match the checksums in its header
\end{lstlisting}
in each language's usual way: Python \code{warnings} (a
\code{RuntimeWarning}), R \code{warning()}, Julia \code{@warn},
MATLAB/Octave \code{warning()} (identifier \code{wfdbx:checksum}), and the
standard error output in C++, Rust and Kotlin (also silenced by
\code{wfdbquiet()}; in Kotlin, an \code{onWarning} callback can receive it
instead). The \code{checksum\_warning} option (\code{checksumWarning} in
Kotlin) turns it off, and the record's \code{checksum\_ok} field is false in
either case, so that a program can test for it. The library function
\code{getvec} reports such an error as in C: \code{getvec: checksum error in
signal N}, and the status $-4$ for the last frame; the tools, such as
\cmd{rdsamp}, therefore stop before printing that frame, as the C tools do.
A signal file that ends early is a read error of the high-level readers,
not a checksum warning, even when the checksum is verified in its last,
incomplete frame (\cref{sec:additions}).

\section{Known C behaviors that WFDB-x reproduces}
\label{sec:known-c}

\wfdbx{} reproduces the C package exactly, including behaviors that may look
surprising; where the documentation of the C package and its code disagree,
\wfdbx{} follows the code. The file \file{docs/KNOWN\_C\_BEHAVIORS.md} lists
them, and says for each whether the official documentation describes it.
Among them:
\begin{itemize}
\item \textbf{Checksums.} \code{getvec} returns $-4$ for the frame at which a
  checksum mismatch is found, so that programs that loop while
  \code{getvec(\ldots) >= 0}, such as \cmd{rdsamp}, stop before printing the
  last frame.
\item \textbf{\code{getinfo} and the header parser disagree on where the
  record line is.} For a header whose first line is empty or blank, the C
  function \code{getinfo} takes that line for the record line, and returns
  the comments between it and the real record line as info strings. This
  is reproduced in legacy mode only: \wfdbx{}'s \code{getinfo} finds the
  record line as the header parser does (\cref{sec:record-line-divergence}).
\item \textbf{\code{putann} writes reserved annotation codes, and drops
  the extreme times.} \code{putann} does not check the annotation code: an
  annotation written with one of the codes 59 to 63, the control codes of
  the MIT format, is read back as a control word, which corrupts the rest of
  the file, and a code above 63 loses its high bits. An annotation at the
  time \code{WFDB\_TIME\_MIN} or \code{WFDB\_TIME\_MAX} ($-2^{63}$ or
  $2^{63}-1$) is refused with \code{putann: time overflow in annotation file
  N} and the status $-1$, and nothing is written for it. The library
  function \code{putann} of \wfdbx{} does the same; the high-level
  \code{write\_annotations} refuses such codes and times before writing
  anything (\cref{sec:additions}).
\item \textbf{Modification labels.} \code{annopen} consumes the \code{NOTE}
  annotations at time 0 (subtype 0) at the beginning of an annotation file,
  so \code{getann} does not return them; a file whose table redefines
  mnemonics makes every annotation file written later in the same process
  repeat the table.
\item \textbf{The time resolution note.} A program that calls
  \code{setafreq} with a positive value (as \cmd{wrann} does) writes a
  \code{\#\# time resolution:} note at the beginning of its output files.
\item \textbf{Duplicates at time 0.} An annotation at time 0 that follows
  another at time 0 counts as in order, so that annotations written twice at
  a time-0 location are not sorted, and both are kept.
\item \textbf{State that survives \code{wfdbquit}}: the counter frequency
  and base counter value, the output annotation time resolution, the buffer
  sizes, and the mnemonic and description tables are not reset.
\item \textbf{Opening an EDF record} sets the time resolution of output
  annotators to the record's sampling frequency, when \code{setafreq} has not
  set one, and adds \code{WFDB\_HIGHRES} to the getvec mode; \code{wfdbquit}
  resets neither, so that an annotation file written later by the same
  process begins with a time-resolution note (\code{\#\# time resolution:
  360} after \code{100s.edf}), and later records are read in high
  resolution. The library functions do the same; the high-level functions
  restore both settings (\cref{sec:additions}).
\item \textbf{Tools.} \cmd{sortann} keeps only the last of several
  annotations with the same time, \code{num} and \code{chan}; \cmd{wrann}
  turns an unrecognized mnemonic into a \code{NOTE} whose aux string starts
  with the mnemonic, and keeps the subtype, \code{chan} and \code{num} of the
  previous line when a line omits them; \cmd{rr2ann} repeats the previous
  interval for an empty line; \cmd{revise} rejects an old-style header of
  exactly 32 signals; \cmd{edf2mit} adds stray bytes to a 16-character signal
  label; \cmd{ahaecg2mit} writes the annotations of \file{.ecg} input one
  sample early; \cmd{mit2edf} uses the current date and time for a record
  that has no start time.
\item \textbf{Documentation and code disagree} on a few points, and \wfdbx{}
  follows the code: \code{wfdbquit} does not reset the counter frequency;
  the variable that turns sorting off is \code{WFDBANNSORT}, not
  \code{WFDBNOSORT}; \cmd{wrsamp}'s \opt{-c} has no effect and \opt{-l} is
  obsolete; its default field separators are space, tab, CR, LF and comma;
  \cmd{edf2mit} and \cmd{wav2mit} name their default output records after
  the input file name.
\item \textbf{Build-specific values.} \code{wfdbversion}, \code{wfdbldflags},
  \code{wfdbcflags} and \code{wfdbdefwfdb} return the values of a release
  build of 10.7.0 installed in \file{/usr/local}: \code{"10.7.0"},
  \code{""}, \code{""}, and the default WFDB path of the C package,
  \code{". /usr/local/\allowbreak database
  http://\allowbreak physionet.org/\allowbreak physiobank/\allowbreak database"};
  \code{wfdbdefwfdbcal} returns \code{"wfdbcal"}.
\end{itemize}

\paragraph{Behaviors that depend on the platform.} The reference results were
produced on x86-64 Linux with glibc, where \code{char} is signed, \code{long}
is 64 bits, and \code{WFDB\_Time} is 64 bits. \wfdbx{} reproduces that
platform's behavior where C's would differ elsewhere; for example,
floating-point operations are never contracted into fused multiply-adds
(which some compilers do by default, and which can change results such as
the baselines that \cmd{edf2mit} computes). One case remains: when an output
file is also an input file in legacy mode (\cref{sec:output-input}), what is
read depends on how far the runtime reads ahead, as in C; the R
implementation on macOS reads further ahead than glibc, and fails one legacy
conformance case for that reason.

% Part I, chapter 5: Header metadata.
\chapter{Header metadata}
\label{chap:metadata}

Many datasets store facts about a recording in the info strings of its
header (\cref{sec:info-strings}), as \code{Key: value} lines, for example:
\begin{lstlisting}[style=text]
# Age: 63
# Sex: F
# Dx: 164889003
\end{lstlisting}
These entries are the \emph{metadata} of the record. Five functions of the
high-level API, added in version 1.2.0, read and change them without
changing anything else in the header: \code{read\_metadata},
\code{get\_metadata}, \code{set\_metadata}, \code{remove\_metadata} and
\code{update\_metadata}. \code{read\_header} also returns the entries, and
\code{write\_record} writes them. The C package has no counterpart. The
library functions \code{putinfo} and \code{setinfo}, and the tools, are
unchanged, and so is \code{getinfo}, except for the rule by which it finds
the record line (\cref{sec:info-strings}), which the metadata functions
share; the C library reads the headers that these functions change exactly
as before (\file{tests/metadata/check\_c.py} checks this with the original
programs). The seven implementations follow the
rules of this chapter identically: they pass the same cases
(\file{tests/metadata/}) and write byte-identical headers.

\section{Entries}
\label{sec:metadata-entries}

A header is split into lines after each line feed (LF). Each line has a line
ending: CR LF, LF, or none (the last line of a file that does not end with a
line feed). The \emph{content} of a line is the line without its ending.

\paragraph{Info strings.} The record line is found as the header parser of
the library (\code{readheader}) finds it: it is the first line whose
content, up to a NUL character if there is one, has a token (separated by
spaces, tabs, carriage returns and line feeds) that does not begin with
\code{\#}; in other words, the first line that holds something other than
white space and does not begin, after white space, with \code{\#}. The
lines before it (comments, and empty or blank lines) are not info strings,
and are never read or changed. Every line after the record line whose
content begins with \code{\#} is an info string, as \code{getinfo} reads a
header (so a comment between two signal lines is an info string too). A
header with no record line has no info strings, and changing its metadata is
an argument error. This rule is the same in
every mode, legacy mode included (since version 1.2.2; before, the record
line was the first line whose content does not begin with \code{\#}, the
rule of the C function \code{getinfo}, which takes an empty first line for
the record line).

\paragraph{Entries and free comments.} An info string is an \emph{entry} if
the text after its \code{\#} contains a colon. The \emph{key} is the text
before the first colon, and the \emph{value} is the text after it; spaces,
tabs, vertical tabs, form feeds and carriage returns are removed at both ends
of each. If the key is then empty, or contains \code{\#}, the info string is
a \emph{free comment}. Free comments, and info strings without a colon, are
never returned and never changed. So \code{\#Age:45}, \code{\# Age: 45} and
\code{\#~~Age~:~45} all give the key \code{Age} and the value \code{45}, and
\cref{tab:metadata-reading} shows a few other cases.

\begin{table}[htbp]
\caption{How info strings are read.}
\label{tab:metadata-reading}
\begin{tabularx}{\linewidth}{@{}P{5.6cm}L@{}}
\toprule
{\tabhead Info string} & {\tabhead Entry (key, value)} \\
\midrule
\code{\#Age: 63} & (\code{Age}, \code{63}) \\
\code{\# Sex:F} & (\code{Sex}, \code{F}) \\
\code{\#~~Dx : 164889003, 59118001~} & (\code{Dx},
  \code{164889003, 59118001}) \\
\code{\#Notes:} & (\code{Notes}, the empty string) \\
\code{\#Time: 12:30:00} & (\code{Time}, \code{12:30:00}): the key ends at the
  first colon \\
\code{\#\textless{}tab\textgreater{}Tab key\textless{}tab\textgreater{}:%
  \textless{}tab\textgreater{}value with tabs} & (\code{Tab key},
  \code{value with tabs}) \\
\code{\# A free comment without a colon} & none: a free comment \\
\code{\#: empty key} & none: the key is empty \\
\code{\#\# Hash: not an entry} & none: the key (\code{\# Hash}) contains
  \code{\#} \\
\code{\#Produced by xform from record 100, beginning at 0:0} &
  (\code{Produced by xform from record 100, beginning at 0}, \code{0}) \\
\bottomrule
\end{tabularx}
\end{table}

The last line of \cref{tab:metadata-reading} is an info string of the sample
record \code{100s}: free text that happens to contain a colon is read as an
entry. Such an entry is harmless: it is changed only if its key is named.

\paragraph{Keys} are compared without regard to letter case, for the letters
A--Z and a--z only (no other character is changed: \code{\'E} and
\code{\'e} differ). So \code{get\_metadata(record, "age")} finds the entry
\code{Age}. Entries are returned in file order.

\paragraph{Duplicate keys.} A header should not hold two entries whose keys
are equal, ignoring letter case. If one does (a file written by other
software), the readers (\code{read\_metadata}, \code{get\_metadata},
\code{read\_header}) return the first of them, and give one warning for each
such key, naming the line numbers of its entries (from 1), such as
\code{header of record dups has more than one entry for key age (lines 4,
6); the first is used}; \code{set\_metadata} and \code{update\_metadata}
refuse to change that key (an argument error), and \code{remove\_metadata}
removes all of them, which repairs the header.

\section{Arguments}
\label{sec:metadata-arguments}

Before a key or a value given to a function is checked, the spaces, tabs,
vertical tabs, form feeds and carriage returns at both ends are removed.
\begin{itemize}
\item A \textbf{key} must then not be empty, and must not contain \code{:},
  \code{\#}, a carriage return, a line feed or a NUL character. It may
  contain spaces, as in \code{Date of birth}.
\item A \textbf{value} is text, which may be a whole sentence. It must not
  contain a carriage return, a line feed or a NUL character; it may be
  empty.
\item Text is one byte per character (as in the rest of \wfdbx{}: in
  Python, Latin-1 strings, or bytes as given).
\item In \code{update\_metadata}, no two keys may be equal, ignoring letter
  case, and a value of none (\code{None}, \code{nothing}, \code{NULL},
  \code{null}, or a missing value, depending on the language) means
  ``remove''.
\item The record name \code{-} (the standard input) is an argument error,
  and so is \code{get\_metadata} with an invalid key.
\end{itemize}
Every argument is checked before the header is read, and every change
before anything is written: after an error, the header is unchanged.

\section{Writing}
\label{sec:metadata-writing}

\paragraph{The written form.} An entry is written as \code{\# Key: value}:
\code{\#}, a space, the key as the caller spelled it, \code{:}, and, if the
value is not empty, a space and the value. An entry with an empty value is
therefore written \code{\# Key:}, with no trailing space.

\paragraph{The changes.}
\begin{itemize}
\item \code{set\_metadata}: if the header has one entry for the key, its line
  is replaced, at the same position, keeping its line ending, even if the
  value is the same (the line is then rewritten in the written form, with the
  caller's spelling of the key). Otherwise a new line is inserted.
\item \code{remove\_metadata}: every line of an entry for the key is deleted,
  with its line ending. The function returns the number of lines deleted (0
  if there is none).
\item \code{update\_metadata}: the changes are applied in the order given,
  each as \code{set\_metadata} or \code{remove\_metadata} would apply it, and
  written at once.
\end{itemize}

\paragraph{Where new lines go.} A new line is inserted after the last info
string of the header, or, if the header has none, after its last line whose
content is not blank (not made only of spaces, tabs, vertical tabs, form
feeds and carriage returns). Its line ending is that of the record line, or,
if the record line has none, that of the first line that has one, or CR LF
(the line ending of the headers written by the WFDB library). If the line
before it has no line ending, that line gets the same ending.

\paragraph{What is kept.} Every other byte of the header is kept: the
comments and the empty or blank lines before the record line, the record
line, the signal lines and their checksums, the free comments, and the line
endings. If nothing changes (and \code{dir} is not given), nothing is
written.

\paragraph{Atomic replacement.} The new header is written to a temporary
file in the directory of the destination, which is never more open than the
destination (in Python, it is created readable and writable by its owner
only). The temporary file then gets the permissions of the destination (or,
for a new file in \code{dir}, those of a new file) and replaces it (by
renaming it). An error never leaves a partly written header, and no other
file is left in the directory.

\section{Which file}
\label{sec:metadata-file}

The header is opened as the library opens the header of a record, after
removing a trailing \file{.hea} from the name: under the input directory, or
as given for an absolute path, a \code{\textasciitilde} path or a URL (a URL
only if network access is on); see \cref{sec:record-paths}. The functions
take the options of the high-level readers for this, an input directory
(\code{input\_dir}) and a network permission (\code{network}), for the
duration of the call. A header that cannot be opened is a file error.

A change is written to the file that was read, unless the call gives an
output directory \code{dir}: the changed header is then written to the file
of the same name in that directory, which must exist (else a file error),
and the header that was read is unchanged. The signal files are not copied.
A header read from a URL cannot be changed without \code{dir} (an argument
error).

\paragraph{Multi-segment records.} Only the named header is read or
changed: for a multi-segment record, that is the master header, unless the
name of a segment is given. (The info strings of a master header are read by
\code{getinfo} like those of any other header.)

\paragraph{Info files.} Only the header file (\file{.hea}) is read, never
the info file (\file{.info}) whose comment lines \code{getinfo} also
returns.

\section{read\_header and write\_record}
\label{sec:metadata-header}

\code{read\_header} returns the entries of a header, as \code{read\_metadata}
returns them (from the \file{.hea} file only), as the \code{metadata} of its
result, beside its \code{comments}, which are all the info strings.

\code{write\_record} writes the entries given in its \code{metadata}
argument (a mapping from keys to values, or a list of key and value pairs,
depending on the language) after its \code{comments}, in the written form
\code{\# Key: value}, through \code{putinfo}. Its keys must be unique,
ignoring letter case, and none may be the key of an entry in
\code{comments}; a value of none is an argument error. This is checked
before any file is created.

\section{An example}
\label{sec:metadata-example}

The header \file{tests/metadata/headers/basic.hea} (with CR LF line endings)
is:
\begin{lstlisting}[style=text]
# Converted on: 2026-11-01 (a comment before the record line)
basic 2 360 21600
100s.dat 212 200 11 1024 995 21537 0 MLII
100s.dat 212 200 11 1024 1011 -3962 0 V5
#Age: 63
# Sex:F
#  Dx : 164889003, 59118001
# A free comment without a colon
#Notes:
# Recorded by: Holter monitor, model X
\end{lstlisting}
\code{read\_metadata("basic")} returns the five entries, in file order:
(\code{Age}, \code{63}), (\code{Sex}, \code{F}), (\code{Dx},
\code{164889003, 59118001}), (\code{Notes}, the empty string) and
(\code{Recorded by}, \code{Holter monitor, model X}). The first line is
before the record line, and is not an info string; the free comment is not
an entry. In Python, the changes
\begin{lstlisting}[language=Python]
wfdbx.update_metadata("basic", [("age", "64"), ("Sex", None),
                                ("Weight", "72"), ("Height", "170")])
\end{lstlisting}
replace the entry \code{Age} (at its position, with the caller's spelling
\code{age}), remove the entry \code{Sex}, and add two new entries after the
last info string, with the line ending of the record line (CR LF). The other
lines are unchanged, byte for byte, including the unusual spacing of the
\code{Dx} line. The header becomes
\file{tests/metadata/expected/update\_basic.hea}:
\begin{lstlisting}[style=text]
# Converted on: 2026-11-01 (a comment before the record line)
basic 2 360 21600
100s.dat 212 200 11 1024 995 21537 0 MLII
100s.dat 212 200 11 1024 1011 -3962 0 V5
# age: 64
#  Dx : 164889003, 59118001
# A free comment without a colon
#Notes:
# Recorded by: Holter monitor, model X
# Weight: 72
# Height: 170
\end{lstlisting}

A header whose first line is empty, such as
\file{tests/metadata/headers/leading.hea}:
\begin{lstlisting}[style=text]
(an empty line)
# Age: preamble
leading 2 360 21600
100s.dat 212 200 11 1024 995 21537 0 MLII
100s.dat 212 200 11 1024 1011 -3962 0 V5
# Age: actual
\end{lstlisting}
has one entry, (\code{Age}, \code{actual}): its record line is the third
line, as for the header parser, and \code{\# Age: preamble}, before it, is
not an info string. \code{set\_metadata("leading", "Sex", "F")} adds
\code{\# Sex: F} after \code{\# Age: actual}, and leaves the first two lines
as they are. In Python:
\begin{lstlisting}[language=Python,style=example]
import wfdbx
with open("leading.hea", "w") as f:
    f.write("\n# Age: preamble\nleading 2 360 21600\n"
            "100s.dat 212 200 11 1024 995 21537 0 MLII\n"
            "100s.dat 212 200 11 1024 1011 -3962 0 V5\n# Age: actual\n")
print(wfdbx.read_metadata("leading"))           # [('Age', 'actual')]
print(wfdbx.read_header("leading").comments)    # [' Age: actual']
\end{lstlisting}

The following Python example writes a record with metadata, from the first
second of record \code{100s}, and changes its metadata:
\begin{lstlisting}[language=Python,style=example]
import wfdbx
rec = wfdbx.read_record("100s", sampto=360)
wfdbx.write_record("demo", rec.signals, fs=rec.fs, fmt=212, gains=rec.gains,
                   baselines=rec.baselines, adczeros=1024, adcres=11,
                   units="mV", names=rec.names,
                   metadata={"Age": "69", "Sex": "M"})
print(wfdbx.read_metadata("demo"))   # [('Age', '69'), ('Sex', 'M')]
wfdbx.set_metadata("demo", "age", "70")
wfdbx.update_metadata("demo", {"Sex": None, "Dx": "Normal sinus rhythm"})
print(wfdbx.get_metadata("demo", "AGE"))         # 70
print(wfdbx.read_metadata("demo"))   # [('age', '70'), ('Dx', 'Normal sinus rhythm')]
print(wfdbx.remove_metadata("demo", "dx"))       # 1
\end{lstlisting}
After the call of \code{update\_metadata}, the header \file{demo.hea} is:
\begin{lstlisting}[style=text]
demo 2 360 360
demo.dat 212 200/mV 11 1024 995 20844 0 MLII
demo.dat 212 200/mV 11 1024 1011 28153 0 V5
# age: 70
# Dx: Normal sinus rhythm
\end{lstlisting}

\section{The functions in each language}
\label{sec:metadata-functions}

\Cref{tab:metadata-functions} lists the five functions, with their names in
each language. They take the record name and, as options, the input
directory and the network permission for reading the header, and, for the
functions that change it, the output directory \code{dir}. The chapters of
\cref{part:languages} give their exact signatures, results and errors
(\cref{chap:python,chap:cpp,chap:rust,chap:julia,chap:r,chap:matlab,chap:kotlin}).

\begin{table}[htbp]
\caption{The header metadata functions.}
\label{tab:metadata-functions}
\begin{tabularx}{\linewidth}{@{}LP{3.2cm}P{3.0cm}P{2.6cm}@{}}
\toprule
{\tabhead What it does} & {\tabhead Python, C++, Rust, Julia, R} &
  {\tabhead MATLAB/Octave} & {\tabhead Kotlin} \\
\midrule
Read all the entries, in file order, as (key, value) pairs &
  \code{read\_metadata} & \code{wfdbx.read\_metadata} &
  \code{readMetadata} \\
Return the value of one entry, or none & \code{get\_metadata} &
  \code{wfdbx.get\_metadata} & \code{getMetadata} \\
Change the value of an entry, or add it & \code{set\_metadata} &
  \code{wfdbx.set\_metadata} & \code{setMetadata} \\
Remove the entries for a key; return the number of lines removed &
  \code{remove\_metadata} & \code{wfdbx.remove\_metadata} &
  \code{removeMetadata} \\
Make several changes in one write & \code{update\_metadata} &
  \code{wfdbx.update\_metadata} & \code{updateMetadata} \\
\bottomrule
\end{tabularx}
\end{table}

The functions are called with the prefix of each language: \code{wfdbx.}
in Python and MATLAB/Octave, \code{wfdbx::} in C++ and Rust, \code{WFDBX.}
in Julia, and none in R and Kotlin. In Rust, each function also has a form
with options, \code{read\_metadata\_with}, \code{get\_metadata\_with}, and so
on. The errors are those of the high-level functions of each language
(\cref{sec:errors}). An argument error and a file error are told apart: in
Python, \code{ValueError} (\code{TypeError} for an argument of the wrong
type) and \code{OSError} (\code{FileNotFoundError} for a missing
\code{dir}); in C++, \code{wfdbx::ArgumentError} and
\code{wfdbx::FileError}; in Julia, \code{ArgumentError} and
\code{WfdbError}; in R, the condition classes
\code{wfdbx\_argument\_error} and \code{wfdbx\_file\_error}; in Kotlin,
\code{WfdbArgumentException} and \code{WfdbException}.

\begin{phibox}
Metadata often holds information about patients, such as ages, dates,
diagnoses or names. \wfdbx{} stores it as plain text, exactly as given, and
does not de-identify it, in the header or anywhere else. Removing identifying
data from headers (their metadata, and the other info strings) before records
are shared is the responsibility of the user.
\end{phibox}


% ---- Part II: The languages --------------------------------------------------------------
\part{The seven languages}
\label{part:languages}
% The Python chapter of the WFDB-x user's manual. It is the template of the chapters of
% the other languages, which follow the same sections in the same order. Each \apientry
% of the section "The high-level API" is checked by check_manual.py against the code
% (python/src/wfdbx): its signature listing must equal the name of the function followed
% by str(inspect.signature(...)), its params must name every parameter and no other, and
% every listing of style "example" is run (in a copy of tests/data).
\chapter{Python}
\label{chap:python}
\chapterlanguage{python}

The Python implementation of WFDB-x is the package \code{wfdbx}, published on PyPI as
\code{wfdb-x}. It is the reference implementation: the six other languages are
translations of it, and its results are those that they must reproduce. It is written in
Python alone, on top of NumPy, and never calls the C library.

It offers WFDB-x in three ways:
\begin{itemize}
\item \textbf{The high-level API} (\cref{sec:python-api}): functions that read and write
  whole records, headers, annotation files and header metadata as NumPy arrays and Python
  objects, in one call, with keyword arguments, and that raise exceptions on errors. They
  read signal files in large blocks, and are the fastest way to read a record. They are an
  addition of WFDB-x: the C package has no equivalent. They are the recommended interface
  for new code, analysis and notebooks.
\item \textbf{The WFDB library functions} under their C names (\code{isigopen},
  \code{getvec}, \code{getann}, \code{putann}, \ldots), with the behavior of the C library
  (\cref{sec:python-library}). They are documented once, for all the languages, in
  \cref{chap:library}.
\item \textbf{The command-line tools} (\code{rdsamp}, \code{rdann}, \code{wrsamp},
  \code{xform} and 29 others), with the options, output and exit status of the C tools,
  run through the \code{wfdbx} command (\cref{sec:python-tools}). They are documented in
  \cref{chap:tools}.
\end{itemize}
The high-level functions return exactly the samples and annotations that the library
functions (and the C library) give.

% ==========================================================================================
\section{Installation and requirements}
\label{sec:python-install}

\subsection{Requirements}

WFDB-x for Python needs Python 3.9 or later and NumPy 1.17 or later, its only dependency.
The package is pure Python: it has no compiled code and needs no C compiler, no WFDB
library and no libFLAC (the FLAC formats 508, 516 and 524 are read and written by
WFDB-x's own codec). It runs on Linux, macOS and Windows.

Two sets of optional dependencies are declared:
\begin{itemize}
\item \code{examples}: matplotlib and Jupyter, to run the example notebooks;
\item \code{test}: pytest, nbformat, nbclient and ipykernel, to run the tests
  (\cref{sec:python-testing}).
\end{itemize}

\subsection{Installing from PyPI}

\begin{lstlisting}[style=shell]
pip install wfdb-x                # the package and the wfdbx command
pip install "wfdb-x[examples]"    # also matplotlib and Jupyter, for the notebooks
\end{lstlisting}

The package is imported as \code{wfdbx} (\code{import wfdbx}); the installation also
creates the command \cmd{wfdbx}, which runs the command-line tools. The installed version
is \code{wfdbx.\_\_version\_\_}, or the output of \cmd{wfdbx -{}-version}:

\begin{lstlisting}[style=shell]
python -c "import wfdbx; print(wfdbx.__version__)"    # 1.2.2
wfdbx --version                                       # WFDB-x 1.2.2
\end{lstlisting}

The wheel and the source distribution include the licenses, \file{LICENSE} (the library,
LGPL-2.1-or-later) and \file{LICENSE-GPL} (the command-line tools, GPL-3.0-or-later), and
the file \file{NOTICE}, which says which parts are under which license and gives the
BSD-3-Clause notice of libFLAC, from which the FLAC decoder (\file{wfdbx/\_flac.py}) is
translated (\cref{sec:license}). The license expression of the package is
\code{LGPL-2.1-or-later AND GPL-3.0-or-later AND BSD-3-Clause}.

\subsection{Installing from the repository}

From the root of a copy of the repository, the package is in the directory
\file{python}:

\begin{lstlisting}[style=shell]
pip install ./python                  # install the package
pip install -e "./python[test]"       # or: editable, with the test dependencies
\end{lstlisting}

Without installation, the package can also be used from its source directory by adding
\file{python/src} to \code{PYTHONPATH}; the command-line tools are then run as
\cmd{python -m wfdbx.cli} (\cref{sec:python-tools}).

\subsection{The sample record and the example notebooks}
\label{sec:python-resources}

The package includes the sample record \code{100s}, the sample record of the WFDB
Software Package: the first minute of record 100 of the MIT-BIH Arrhythmia Database, two
ECG signals (\code{MLII} and \code{V5}) sampled at 360\,Hz, with the reference beat
annotations (annotator \code{atr}). Its files (\file{100s.hea}, \file{100s.dat},
\file{100s.atr}) are in the directory returned by \code{wfdbx.sample\_data\_dir()}, so
that the examples and the notebooks run without a network connection. Work that uses the
MIT-BIH Arrhythmia Database cites it as described on PhysioNet.

The package also includes seven Jupyter notebooks (\cref{tab:python-notebooks}), in the
directory returned by \code{wfdbx.examples\_dir()}. Each runs without a network
connection (\code{06\_physionet\_records} reads record 100 from PhysioNet when it can,
and the sample record otherwise). The function \code{copy\_examples}, or the command
\cmd{wfdbx -{}-examples DIR}, copies them to a directory where they can be opened and run:

\begin{lstlisting}[style=shell]
wfdbx --examples my-examples      # without DIR: ./wfdbx-examples
cd my-examples
jupyter lab
\end{lstlisting}

\begin{table}[htbp]
\centering
\caption{The example notebooks of the Python package}
\label{tab:python-notebooks}
\begin{tabularx}{\linewidth}{@{}lL@{}}
\toprule
\tabhead Notebook & \tabhead Contents \\
\midrule
\code{01\_getting\_started} & reading headers and signals; ADC and physical units \\
\code{02\_annotations\_and\_heart\_rate} & reading annotations; RR intervals and heart
  rate \\
\code{03\_writing\_records} & writing records and annotations \\
\code{04\_command\_line\_tools} & running the tools \\
\code{05\_c\_library\_api} & the WFDB library functions \\
\code{06\_physionet\_records} & reading a record from PhysioNet \\
\code{07\_legacy\_mode} & the original package's behavior, on request \\
\bottomrule
\end{tabularx}
\end{table}

% ==========================================================================================
\section{Quick start}
\label{sec:python-quickstart}

This example reads the header, ten seconds of the signals and the annotations of the
sample record \code{100s}, from the package's data directory (\code{input\_dir}), and
computes the mean heart rate. It then writes the first ten seconds as a new record,
\code{first10}, in the current directory, with two metadata entries in its header.

\begin{lstlisting}[language=Python,style=example]
import numpy as np
import wfdbx

sample = wfdbx.sample_data_dir()        # the directory of the sample record 100s

# The header only: no samples are read.
hdr = wfdbx.read_header("100s", input_dir=sample)
print(hdr.nsig, hdr.fs, hdr.nsamples)   # 2 360.0 21600
print([s.desc for s in hdr.siginfo])    # ['MLII', 'V5']

# Seconds 10 to 20 of both signals, in physical units (float64): one row per sample,
# one column per signal.
rec = wfdbx.read_record("100s", sampfrom="0:10", sampto="0:20", physical=True,
                        input_dir=sample)
print(rec.signals.shape, rec.signals.dtype)    # (3600, 2) float64
print(rec.signals[0])                          # [-0.39  -0.275]

# The beat annotations, and the mean heart rate.
ann = wfdbx.read_annotations("100s", "atr", input_dir=sample)
print(len(ann), ann.symbols[:3], ann.samples[:3])    # 75 ['+', 'N', 'N'] [ 18  77 370]
beats = ann.samples[[bool(wfdbx.isqrs(c)) for c in ann.codes]]
rr = np.diff(beats) / ann.fs                         # RR intervals, in seconds
print(len(beats), "beats,", round(60 / rr.mean(), 1), "beats per minute")   # 74 beats, 73.9 ...

# Write the first 10 seconds, in ADC units, as the record "first10" (first10.hea and
# first10.dat, in the current directory), with the signal specifications of 100s.
x = wfdbx.read_record("100s", sampto="0:10", input_dir=sample)
wfdbx.write_record("first10", x.signals, fs=x.fs, fmt=16, gains=x.gains,
                   baselines=x.baselines, adczeros=[s.adczero for s in x.siginfo],
                   adcres=[s.adcres for s in x.siginfo], units="mV", names=x.names,
                   metadata={"Source": "MIT-BIH record 100"})

# Header metadata: "# Key: value" lines of the header.
wfdbx.set_metadata("first10", "Age", "69")
print(wfdbx.read_metadata("first10"))   # [('Source', 'MIT-BIH record 100'), ('Age', '69')]
print(wfdbx.check_record("first10"))    # no problems found
\end{lstlisting}

The header \file{100s.hea} gives no units: the example writes \code{"mV"} explicitly,
since WFDB-x never assumes units, a gain, an ADC resolution or a sampling frequency that
a record does not give (\cref{chap:behavior}).

% ==========================================================================================
\section{The high-level API}
\label{sec:python-api}

The high-level API is made of the functions of this section, all available at the top
level of the package (\code{wfdbx.read\_record}, \ldots), and of the classes of their
results (\cref{sec:python-results}). They are defined in the modules
\code{wfdbx.highlevel}, \code{wfdbx.metadata}, \code{wfdbx.check} and
\code{wfdbx.resources}; \code{help()} gives the documentation of each.

\subsection{Conventions}
\label{sec:python-conventions}

\paragraph{Record names and directories.} A record is named as in the other parts of
WFDB-x (\cref{chap:records,sec:record-paths}), without the \file{.hea} extension (a
trailing \code{".hea"} is removed, by the readers and the writers). A record name is one
of:
\begin{itemize}
\item a full path, absolute (\code{"/data/mitdb/100"}) or in the home directory
  (\code{"\textasciitilde/Documents/100s"}), used as given;
\item a relative path (\code{"100s"}, \code{"data/100s"}), read from the input directory
  if one is set, otherwise from the current directory, and written to the output
  directory if one is set, otherwise to the current directory (missing directories are
  created);
\item a URL (\code{"https://physionet.org/files/mitdb/1.0.0/100"}), read only if network
  access is on; a URL cannot be written.
\end{itemize}
The input and output directories are set with \code{wfdbx\_set\_input\_dir} and
\code{wfdbx\_set\_output\_dir} (\cref{sec:python-settings}) or with the environment
variables \code{WFDBX\_INPUT\_DIR} and \code{WFDBX\_OUTPUT\_DIR}. The functions that read
take the optional keyword arguments \code{input\_dir} and \code{network}, and those that
write take \code{dir}: they apply to that one call, and the previous settings are
restored when it returns, also after an error. \code{input\_dir} applies to relative
names only, not to full paths or URLs. A writer given a name with a path creates the
files in that directory, and the header names the record and its signal file without the
path (\code{rec}, \code{rec.dat}), so that the record opens under any path that leads to
it; in legacy mode, as in the original, the header gives the name with its path. With
\code{dir}, the name must be a plain name (\code{"rec"}): the files are created in
\code{dir}, and the record opens from that directory under its plain name. The signal
files that a header names are read from the header's directory.

\paragraph{Times.} The arguments \code{sampfrom} and \code{sampto} are sample numbers
(whole numbers) or WFDB time strings, as \code{strtim} reads them: \code{"0:10"} (10
seconds), \code{"1:02:03.5"}, \code{"s3600"} (sample 3600), \code{"e"} (the end of the
record), or a time of day in brackets, \code{"[14:00:00]"}. A time in seconds needs the
sampling frequency; when the record does not give it, such a time raises
\code{ValueError}.

\paragraph{Per-signal and per-annotation arguments.} The arguments that describe each
signal (\code{fmt}, \code{gains}, \code{units}, \ldots) or each annotation
(\code{subtypes}, \code{chans}, \ldots) take one value for all, or a list (or tuple, or
NumPy array) with exactly one value per signal or per annotation.

\paragraph{Signal parameters.} No sampling frequency, gain, ADC resolution or units is
ever assumed. Where one is needed and the record does not give it, it must be given as an
argument; otherwise the function raises \code{ValueError}. A record whose header gives no
sampling frequency has \code{fs = 0}.

\paragraph{Legacy mode.} With \code{legacy=True}, a function behaves as in the original
WFDB package, with its assumed signal parameters and its search for files
(\cref{chap:behavior}); the library's legacy mode is on for the duration of the call
only. With \code{legacy=False} (the default), the current setting of
\code{wfdbx\_set\_legacy} applies, which is off unless it was turned on.

\paragraph{State.} The functions that read and write records and annotation files close
anything opened with the WFDB library functions (they call \code{wfdbquit()}) before
they start, and close their own files before they return, whatever the outcome.
\code{read\_record} sets the getvec mode for its own reading only. \code{read\_record},
\code{read\_header}, \code{read\_annotations}, \code{write\_record},
\code{write\_annotations} and \code{check\_record} leave the time resolution of output
annotators (\code{getafreq()}) and the getvec mode (\code{getgvmode()}) as they were, also
after reading an EDF record, whose opening sets both (as in C; the library functions keep
that behavior). The metadata functions open no signals and leave the library's state as
it is.

\paragraph{The examples} of this section read the test records of WFDB-x (among them
\code{100s}) from the current directory, and write records of their own there. To run
them elsewhere, copy the sample record to the current directory first, for example with
\code{shutil.copytree(wfdbx.sample\_data\_dir(), ".", dirs\_exist\_ok=True)}.

% ------------------------------------------------------------------------------------------
\subsection{Reading records}
\label{sec:python-reading}

\apientry{read\_record}{Read the signals of a record}
\begin{lstlisting}[language=Python,style=signature]
read_record(record, channels=None, sampfrom=0, sampto=None, physical=False, highres=False, pad=False, *, input_dir=None, network=None, legacy=False, checksum_warning=True)
\end{lstlisting}
Reads the signals of a record, or part of it, into a NumPy array with one row per sample
and one column per signal, together with the description of the record. The samples are
those that \code{getvec} returns (\cref{chap:library}), read in large blocks. Invalid
samples are \code{wfdbx.WFDB\_INVALID\_SAMPLE} ($-32768$) in ADC units, and NaN in
physical units. When the read reaches the end of the record, the samples are compared
with the checksums of the header (\cref{sec:python-warnings}). The result is allocated
for at most $2^{26}$ samples in advance, whatever length the header gives, and grows as
the samples are read.
\begin{params}
\item[record] Record name (a string), without the \file{.hea} extension (a trailing
  \code{".hea"} is also accepted): a full path, absolute (\code{"/data/mitdb/100"}) or in
  the home directory (\code{"\textasciitilde/Documents/100s"}), used as given; a relative
  path (\code{"100s"}, \code{"data/100s"}), read from the input directory if one is set,
  otherwise from the current directory; or a URL
  (\code{"https://physionet.org/files/mitdb/1.0.0/100"}), read only if network access is
  on.
\item[channels] The signals to read, in the order in which they are returned: a list of
  signal numbers (\code{int}, from 0) or signal descriptions (\code{str}, as in
  \code{Record.names}, such as \code{"MLII"}), which may be mixed. Default \code{None}:
  all the signals, in the order of the header.
\item[sampfrom] The first sample to read: a sample number (a whole number $\ge 0$) or a
  time string. Default \code{0}: the beginning of the record.
\item[sampto] The sample after the last one to read: a sample number (a whole number) or
  a time string. A value at or before \code{sampfrom} reads no samples. Default
  \code{None}: the end of the record.
\item[physical] \code{False} (the default) to return ADC units, as \code{int32};
  \code{True} to return physical units, as \code{float64}, computed as
  $(\mathit{adc} - \mathit{baseline}) / \mathit{gain}$. Every signal read must then have a
  gain.
\item[highres] \code{False} (the default) to read a multi-frequency record at its frame
  frequency, with the samples of the faster signals averaged to one per frame;
  \code{True} to read every sample, at the highest sampling frequency of the record (the
  samples of the slower signals are repeated). \code{Record.fs} is then that frequency.
\item[pad] \code{True} to replace each invalid sample by the previous valid sample of the
  same signal. Default \code{False}.
\item[input\_dir] The input directory for this call only (optional): a local directory, a
  \code{\textasciitilde} path or the URL of a directory, under which relative record names
  are read. It does not apply to full paths or URLs. Default \code{None}: the current
  setting (\code{wfdbx\_set\_input\_dir}).
\item[network] \code{True} to allow reading URLs for this call only, \code{False} to forbid
  it. Default \code{None}: the current setting (\code{wfdbx\_set\_network}).
\item[legacy] \code{True} to behave as the original WFDB package (legacy mode, for this
  call only). Default \code{False}.
\item[checksum\_warning] \code{True} (the default) to issue a \code{RuntimeWarning} when
  the samples do not match the checksums of the header; \code{False} for no warning. The
  samples are returned either way, and \code{checksum\_ok} of the result is \code{False}.
\end{params}
\returns{A \code{Record} (\cref{tab:python-record}).}
\raises{\code{IOError} (that is, \code{OSError}) if the header or a signal file cannot be
read, if the record has no signals, if \code{sampfrom} is beyond the end of the record,
or if a signal file is shorter than its header says or cannot be read (``can't read
signal data of record REC (getvec status -3)''), even if the checksums do not match
either: no incomplete frame is ever returned; \code{ValueError} if \code{sampfrom} is
negative, if \code{sampfrom} or \code{sampto} is not a whole number (a fraction, NaN or an infinite
value), if a channel is not in the record, if a time string needs the sampling frequency
and the header gives none, or if \code{physical=True} and a signal read has no gain.
Warns with a \code{RuntimeWarning} on a checksum error (see \code{checksum\_warning}).}
\begin{lstlisting}[language=Python,style=example]
import wfdbx

# Seconds 10 to 20 of lead MLII, in physical units.
rec = wfdbx.read_record("100s", channels=["MLII"], sampfrom="0:10", sampto="0:20",
                        physical=True)
print(rec.signals.shape, rec.signals.dtype)   # (3600, 1) float64
print(rec.names, rec.fs, rec.sampfrom)        # ['MLII'] 360.0 3600

# The first 3 samples of both signals, in ADC units, V5 first.
rec = wfdbx.read_record("100s", channels=[1, 0], sampto=3)
print(rec.signals.tolist())     # [[1011, 995], [1011, 995], [1011, 995]]
print(rec.names, rec.gains, rec.baselines, rec.units)
# ['V5', 'MLII'] [200.0, 200.0] [1024, 1024] [None, None]

# The whole record: the checksums are verified.
rec = wfdbx.read_record("100s")
print(rec.signals.shape, rec.checksum_ok)     # (21600, 2) True
\end{lstlisting}
A multi-frequency record (here the test record \code{mfreq}, whose signals have 1, 2
and 4 samples per frame at a frame frequency of 100\,Hz) is read at its frame frequency,
or at its highest sampling frequency:
\begin{lstlisting}[language=Python,style=example]
import wfdbx

rec = wfdbx.read_record("mfreq")
print(rec.fs, rec.signals.shape)       # 100.0 (600, 3)
rec = wfdbx.read_record("mfreq", highres=True)
print(rec.fs, rec.signals.shape)       # 400.0 (2400, 3)
print([s.spf for s in rec.siginfo])    # [1, 2, 4]
\end{lstlisting}

\apientry{read\_header}{Read the header of a record}
\begin{lstlisting}[language=Python,style=signature]
read_header(record, *, input_dir=None, network=None, legacy=False)
\end{lstlisting}
Reads the header of a record, without opening its signal files. The result describes the
record and each of its signals, and holds the info strings and the metadata entries of the
header. For a multi-segment record, the master header is read, and \code{nsamples} is the
length of the whole record.
\begin{params}
\item[record] Record name (a string), as for \code{read\_record}: a full path (absolute or
  \code{\textasciitilde}), used as given; a relative path, read from the input directory
  if one is set, otherwise from the current directory; or a URL, read only if network
  access is on.
\item[input\_dir] The input directory for this call only (optional; as for
  \code{read\_record}), under which relative record names are read; it does not apply to
  full paths or URLs. Default \code{None}: the current setting.
\item[network] \code{True} to allow reading URLs for this call only, \code{False} to forbid
  it. Default \code{None}: the current setting.
\item[legacy] \code{True} to read the header as the original WFDB package does. Default
  \code{False}.
\end{params}
\returns{A \code{Header} (\cref{tab:python-header}).}
\raises{\code{IOError} if the header cannot be found or read.}
\begin{lstlisting}[language=Python,style=example]
import wfdbx

hdr = wfdbx.read_header("100s")
print(hdr.name, hdr.fs, hdr.nsig, hdr.nsamples)    # 100s 360.0 2 21600
for s in hdr.siginfo:
    print(s.desc, s.fname, s.fmt, s.gain, s.adcres, s.units)
# MLII 100s.dat 212 200.0 11 None
# V5 100s.dat 212 200.0 11 None
print(hdr.comments[0])          # " 69 M 1085 1629 x1": the first info string
print(hdr.base_time)            # None: the header gives no start time
\end{lstlisting}

\apientry{read\_annotations}{Read an annotation file}
\begin{lstlisting}[language=Python,style=signature]
read_annotations(record, annotator, sampfrom=0, sampto=None, fs=None, *, input_dir=None, network=None, legacy=False)
\end{lstlisting}
Reads the annotations of an annotation file, in file order, into NumPy arrays and lists.
The sampling frequency, needed to convert times in seconds, is that of the record's
header, or, for a record without a header or whose header gives none, the argument
\code{fs}.
\begin{params}
\item[record] Record name (a string), as for \code{read\_record}: a full path (absolute or
  \code{\textasciitilde}), used as given; a relative path, read from the input directory
  if one is set, otherwise from the current directory; or a URL, read only if network
  access is on. Its header, if there is one, gives the sampling frequency.
\item[annotator] Annotator name: the suffix of the annotation file. \code{"atr"} reads the
  file \file{RECORD.atr}.
\item[sampfrom] The first time to read: a sample number (a whole number $\ge 0$) or a
  time string. Default \code{0}: the beginning of the file.
\item[sampto] The time at which to stop: annotations at or after it are not read. A
  sample number (a whole number) or a time string. Default \code{None}: the end of the
  file.
\item[fs] The sampling frequency, in Hz (a number $> 0$), to use when the record has no
  header or its header gives none; it is ignored otherwise. Default \code{None}: none.
\item[input\_dir] The input directory for this call only (optional; as for
  \code{read\_record}), under which relative record names are read; it does not apply to
  full paths or URLs. Default \code{None}: the current setting.
\item[network] \code{True} to allow reading URLs for this call only, \code{False} to forbid
  it. Default \code{None}: the current setting.
\item[legacy] \code{True} to behave as the original WFDB package (which assumes 250\,Hz
  for a record without a frequency). Default \code{False}.
\end{params}
\returns{An \code{Annotations} object (\cref{tab:python-annotations}). Its \code{fs} is
the frequency in which the times are counted: the one written in the annotation file, if
there is one, or else the sampling frequency; it is 0 if none of them is known.}
\raises{\code{IOError} if the annotation file cannot be opened or read, or if it ends
before its end-of-file marker (``can't read annotation file ANN of record REC (getann
status -3)''); \code{ValueError} if \code{fs} is given and is not $> 0$, if
\code{sampfrom} is negative (``sampfrom must not be negative''), if \code{sampfrom} or
\code{sampto} is not a whole number (a fraction, NaN or an infinite value: ``sampfrom must
be a whole number''), or if a time string needs the sampling frequency and none is known.
The arguments are checked before the annotation file is opened (a time string is converted
with the sampling frequency of the record's header, which is read for it), and the
annotation file is closed whatever happens.}
\begin{lstlisting}[language=Python,style=example]
import wfdbx

ann = wfdbx.read_annotations("100s", "atr", sampfrom="0:10", sampto="0:30")
print(len(ann), ann.fs)                   # 24 360.0
print(ann.samples[:3], ann.symbols[:3])   # [3862 4170 4466] ['N', 'N', 'N']

# The whole file: the first annotation is a rhythm change, with an aux string.
ann = wfdbx.read_annotations("100s", "atr")
print(ann.symbols[0], ann.codes[0] == wfdbx.RHYTHM, ann.aux[0])   # + True (N
print(sorted(set(ann.symbols)))           # ['+', 'A', 'N']
\end{lstlisting}
For a record without a header, \code{fs} gives the frequency in which times in seconds
are converted (the library reports on the standard error that it cannot open the
header):
\begin{lstlisting}[language=Python,style=example]
import wfdbx

wfdbx.write_annotations("beats", "qrs", [100, 460, 820], symbols="N")   # no header
ann = wfdbx.read_annotations("beats", "qrs", sampto="0:02", fs=360)
print(ann.samples, ann.fs)                # [100 460] 360.0
\end{lstlisting}

% ------------------------------------------------------------------------------------------
\subsection{Writing records}
\label{sec:python-writing}

\apientry{write\_record}{Write a header and a signal file}
\begin{lstlisting}[language=Python,style=signature]
write_record(record, signals, fs, fmt=16, gains=None, baselines=None, adczeros=0, adcres=None, units=None, names=None, base_time=None, comments=None, physical=False, samps_per_frame=None, *, metadata=None, dir=None, legacy=False)
\end{lstlisting}
Writes signals as a new record: the header \file{RECORD.hea} and one signal file,
\file{RECORD.dat}, which holds all the signals (so they share one format). Existing files
of these names are replaced: both files are written to temporary files in the
destination directory (\file{.NAME.XXXXXXXX.tmp}, readable and writable by the owner
only), which replace the existing files, with their permissions, only when both are
complete; after an error, the temporary files are removed and the existing files are left
as they were. For a symbolic link, the file it leads to is replaced; a destination that
is not a regular file (a device) is written directly. The sampling frequency, the gain and
ADC resolution of each signal, and the units of each calibrated signal must be given:
there are no defaults. The per-signal arguments (\code{fmt}, \code{gains}, \code{baselines},
\code{adczeros}, \code{adcres}, \code{units}, \code{names}) take one value for all the
signals, or a list with one value per signal. All the arguments are checked before any
file is created or changed.
\begin{params}
\item[record] Record name (a non-empty string), without the \file{.hea} extension (a
  trailing \code{".hea"} is removed, as by the readers): a plain name (\code{"rec"}) or a relative path (\code{"data/rec"}), created in the output
  directory if one is set, otherwise in the current directory, missing directories
  included; or a full path, absolute or in the home directory
  (\code{"\textasciitilde/Documents/rec"}), used as given. With a path, the files are
  created in that directory, and the header names the record and its signal file without
  the path (\code{rec}, \code{rec.dat}), so that the record opens under any path that
  leads to it (in legacy mode, as in the original, the header gives the name with its
  path). A URL cannot be written.
\item[signals] The samples: an array (or anything \code{numpy.asarray} accepts) with one
  row per sample and one column per signal (2-D); a 1-D array is one signal. In ADC units
  (the default), the values are integers, or floats with whole values, within the range of
  a 32-bit integer; with \code{physical=True}, they are in physical units.
\item[fs] The sampling frequency, in Hz: a finite number $> 0$. Required.
\item[fmt] The storage format, such as \code{16}, \code{212}, \code{24} or \code{516}
  (FLAC); one value, or one per signal, all equal (\cref{chap:formats}). Default
  \code{16}.
\item[gains] The gain of each signal, in ADC units per physical unit: a finite number;
  \code{0} marks an uncalibrated signal, which is allowed only for samples given in ADC
  units. Negative gains are allowed. Required (default \code{None}: not given).
\item[baselines] The ADC value that corresponds to 0 physical units (a whole number
  within the range of a 32-bit integer). Default \code{None}: the ADC zero of the signal.
\item[adczeros] The ADC value at the middle of the ADC range (a whole number within the
  range of a 32-bit integer). Default \code{0}.
\item[adcres] The ADC resolution, in bits (a whole number $\ge 1$). Required (default
  \code{None}: not given).
\item[units] The physical units, such as \code{"mV"} (Latin-1 text, one byte per
  character, as all text in \code{wfdbx}). Required for each signal whose gain is not 0.
  Default \code{None}: not given.
\item[names] The signal descriptions, such as \code{"MLII"} (Latin-1 text). Default
  \code{None}: no descriptions.
\item[base\_time] The start time, with an optional date:
  \code{"hh:mm:ss"} or \code{"hh:mm:ss dd/mm/yyyy"}, with optional decimals of the
  seconds. Default \code{None}: no start time.
\item[comments] A list of info strings (Latin-1 text), written to the header as comment lines
  (\code{"\#"} followed by the string). Default \code{None}: none.
\item[physical] \code{True} if \code{signals} are in physical units. They are converted
  to ADC units as $\mathrm{round}(x \cdot \mathit{gain}) + \mathit{baseline}$ (halves
  rounded away from zero); NaN becomes an invalid sample, and infinite values and values
  beyond the range of 32-bit ADC units are errors. Default \code{False}: ADC units.
\item[samps\_per\_frame] For a multi-frequency record, the number of samples per frame of
  each signal (integers $\ge 1$). \code{signals} then has one row per frame and
  \code{sum(samps\_per\_frame)} columns: the samples of signal 0, then those of signal 1,
  and so on. \code{fs} is then the frame frequency. Default \code{None}: one sample per
  frame for every signal.
\item[metadata] Metadata entries (\cref{chap:metadata}) to write to the header after the
  comments, as \code{\# Key: value}: a mapping from key to value (strings), or a sequence
  of (key, value) pairs. Keys and values follow the rules of \code{set\_metadata}; no two
  keys may differ only in letter case, nor be the key of an entry in \code{comments}.
  Default \code{None}: none.
\item[dir] The directory (which must exist) in which to create the files, for this call
  only. \code{record} must then be a plain name, without a directory, and the record opens
  from that directory under its plain name. Default \code{None}: the output directory.
\item[legacy] \code{True} to behave as the original WFDB package: the gain defaults to
  200, and the ADC resolution to one based on the format. Default \code{False}.
\end{params}
\returns{\code{None}.}
\raises{\code{ValueError} if \code{record} is empty; if \code{signals} is not a 1-D or 2-D
array (``signals must be a 1-D or 2-D array''); if it has no signals
(``at least one signal is required''); if \code{fs}, a gain, an ADC resolution or the
units of a calibrated signal is missing or invalid (\code{fs} and the gains must be
finite, and \code{fs} $> 0$); if a baseline, an ADC zero or an ADC resolution is not a
whole number within the range of a 32-bit integer (``baselines must be whole numbers
within the range of a 32-bit integer''); if the record name, a signal name, the units or
a comment has a character that is not Latin-1 (``units 'XX' has a character that is not
Latin-1 (one byte per character)''); if \code{base\_time} cannot be parsed; if a list of
per-signal values does not have one value per signal; if a signal in physical units has a
gain of 0; if a sample is not valid (see \code{signals} and \code{physical}); if the
signals do not all have the same format; if a \code{samps\_per\_frame} value is not an
integer $\ge 1$, or the number of columns of \code{signals} is not their sum; if an entry
of \code{metadata} is not valid; or if \code{record} contains a directory and \code{dir}
is given. \code{TypeError} if \code{metadata} is neither a mapping nor a sequence of
pairs, or a key or value is not a string. \code{FileNotFoundError} if \code{dir} does not
exist. \code{IOError} if the signal file or the header cannot be created, or cannot be
written completely, for example on a full disk (``write error on record REC'').}
\begin{lstlisting}[language=Python,style=example]
import numpy as np
import wfdbx

# Two signals of random numbers, in millivolts, at 360 Hz: the record "rand"
# (rand.hea and rand.dat).
x = np.random.default_rng(1).standard_normal((3600, 2))
wfdbx.write_record("rand", x, fs=360, fmt=16, physical=True, gains=1000, adcres=16,
                   units="mV", names=["A", "B"], base_time="14:00:03 21/06/2024",
                   comments=["random numbers"], metadata=[("Generator", "numpy")])
rec = wfdbx.read_record("rand", physical=True)
print(rec.signals.shape, rec.base_time)       # (3600, 2) [14:00:03 21/06/2024]
print(np.abs(rec.signals - x).max() <= 0.0005)        # True: rounded to 1/1000 mV
print(wfdbx.read_header("rand").comments)     # ['random numbers', ' Generator: numpy']

# An uncalibrated signal (gain 0, no units), in ADC units, in the directory "out".
import os
os.makedirs("out", exist_ok=True)
wfdbx.write_record("counts", np.arange(1000), fs=100, gains=0, adcres=12, dir="out")
print(wfdbx.read_record("out/counts").signals[:3, 0])    # [0 1 2]
\end{lstlisting}
A multi-frequency record: signal \code{fast} has 2 samples per frame, signal
\code{slow} one; each row of \code{signals} is one frame, with the two samples of
\code{fast} and then the sample of \code{slow}:
\begin{lstlisting}[language=Python,style=example]
import numpy as np
import wfdbx

frames = 100                                    # 1 second at 100 frames per second
fast = np.arange(2 * frames).reshape(frames, 2)     # 200 samples per second
slow = np.full((frames, 1), 500)                    # 100 samples per second
wfdbx.write_record("mf", np.hstack([fast, slow]), fs=100, gains=200, adcres=12,
                   units="mV", names=["fast", "slow"], samps_per_frame=[2, 1])
rec = wfdbx.read_record("mf", highres=True)
print(rec.fs, rec.signals[:3].tolist())    # 200.0 [[0, 500], [1, 500], [2, 500]]
\end{lstlisting}

\apientry{write\_annotations}{Write an annotation file}
\begin{lstlisting}[language=Python,style=signature]
write_annotations(record, annotator, samples, symbols=None, codes=None, subtypes=0, chans=0, nums=0, aux=None, fs=None, *, sort=None, dir=None, legacy=False)
\end{lstlisting}
Writes an annotation file, \file{RECORD.ANNOTATOR}, in the MIT format. An existing file of
that name is replaced: the new file is written to a temporary file in the same directory
(readable and writable by the owner only), which replaces it, with its permissions, only
when it is complete (and, if needed, sorted); after an error, the temporary file is
removed and the existing file is left as it was. The record need not have a header. The
annotation types are given as \code{symbols} or as \code{codes} (\cref{chap:formats}). Each per-annotation
argument is one value for all the annotations, or a list with exactly one value per
annotation. If the annotations are not in time order, the file is sorted when it is
written, unless \code{sort=False}. All the arguments are checked before the file is
created or changed.
\begin{params}
\item[record] Record name (a string), without the \file{.hea} extension (a trailing
  \code{".hea"} is also accepted): a plain name or a relative path (\code{"data/100s"}),
  created in the output directory if one is set, otherwise in the current directory; or a
  full path, absolute or in the home directory (\code{"\textasciitilde/Documents/100s"}),
  used as given. The file is \file{RECORD.ANNOTATOR} at that place
  (\file{data/100s.qrs}). A URL cannot be written.
\item[annotator] Annotator name: the suffix of the file, made of letters, digits,
  \code{\_}, \code{-} and \code{\textasciitilde}.
\item[samples] The annotation times, in samples: a 1-D array or sequence of integers, one
  per annotation, from $-2^{63}+1$ to $2^{63}-2$ (the smallest and largest 64-bit
  integers, \code{WFDB\_TIME\_MIN} and \code{WFDB\_TIME\_MAX}, cannot be written).
\item[symbols] The mnemonics of the annotation types, such as \code{"N"} or \code{"V"}:
  one per annotation, or one string for all. An unknown mnemonic becomes code 0
  (\code{NOTQRS}). Default \code{None}: the types are given by \code{codes}.
\item[codes] The annotation codes (\code{wfdbx.NORMAL}, \code{wfdbx.PVC}, \ldots), from
  0 to 58 (59 to 63 are the control codes of the file format, \cref{sec:formats-mit-annotations}),
  instead of \code{symbols}: one per annotation, or one value for all. If both are given
  (as \code{read\_annotations} returns them), each symbol must be that of its code (an
  empty symbol is not checked). Default \code{None}.
\item[subtypes] The subtype of each annotation, stored as a signed byte (other values wrap
  around to $-128 \ldots 127$, as in C). Default \code{0}.
\item[chans] The channel number of each annotation, stored as an unsigned byte (other
  values wrap around to $0 \ldots 255$). Default \code{0}.
\item[nums] The annotator number of each annotation, stored as a signed byte. Default
  \code{0}.
\item[aux] The auxiliary text of each annotation, with \code{None} where there is none:
  a list, or one text for all. Texts longer than 255 bytes are cut. Default \code{None}:
  none.
\item[fs] If given (a finite number $> 0$), the frequency in which the times are
  counted, written to the file; it applies to this file only. Default \code{None}: none is
  written.
\item[sort] \code{True} to always sort the file, \code{False} to keep the given order
  (without the library's advice to run \code{sortann}), for this call only. Default
  \code{None}: the environment variable \code{WFDBANNSORT} decides (sorting is on unless
  it is \code{0}).
\item[dir] The directory (which must exist) in which to create the file, for this call
  only. \code{record} must then be a plain name, without a directory: the file is
  \file{DIR/RECORD.ANNOTATOR}. Default \code{None}: the output directory.
\item[legacy] \code{True} to behave as the original WFDB package. Default \code{False}.
\end{params}
\returns{\code{None}.}
\raises{\code{ValueError} if neither \code{symbols} nor \code{codes} is given for one
annotation or more (\code{symbols or codes must be given}; with no annotations, an empty
file is written); if both are
given and disagree (``symbols and codes disagree at annotation K''); if \code{samples}
(or a sequence of codes, subtypes, channels or annotator numbers) is not 1-D (``samples
must be a 1-D array''); if a list of per-annotation values does not have one value per
annotation; if a time is not an integer
within the range of a 64-bit integer (NaN, an infinite value, a fraction or a huge
value), or is $-2^{63}$ or $2^{63}-1$ (``annotation time T cannot be written''); if a
code is not from 0 to 58 (``annotation code 61 is reserved (59 to 63 are the control
codes of the annotation file format)'', ``annotation code 64 is out of range (0 to
58)''); if a sample, code, subtype, channel or annotator number is NaN or \code{None}; if
an aux text has a character that is not Latin-1; if \code{fs} is not finite and $> 0$; or if \code{record} contains a directory and
\code{dir} is given. \code{FileNotFoundError} if \code{dir} does not exist.
\code{IOError} if the annotation file cannot be created, or cannot be written completely,
or the library (\code{putann}) reports an error while writing it (``write error on
annotator ANN for record REC'').}
\begin{lstlisting}[language=Python,style=example]
import wfdbx

# Three beats and a rhythm change, by symbol, with an aux string for the rhythm.
wfdbx.write_annotations("100s", "test", [18, 77, 370, 662],
                        symbols=["+", "N", "V", "N"], aux=["(N", None, None, None])
ann = wfdbx.read_annotations("100s", "test")
print(ann.symbols, ann.aux)     # ['+', 'N', 'V', 'N'] ['(N', None, None, None]

# A copy of the reference annotations, as read_annotations returns them.
a = wfdbx.read_annotations("100s", "atr")
wfdbx.write_annotations("100s", "copy", a.samples, symbols=a.symbols, codes=a.codes,
                        subtypes=a.subtypes, chans=a.chans, nums=a.nums, aux=a.aux)
b = wfdbx.read_annotations("100s", "copy")
print((b.samples == a.samples).all(), b.symbols == a.symbols, b.aux == a.aux)
# True True True

# Annotations out of time order are sorted when the file is written (the library
# reports it on the standard error).
wfdbx.write_annotations("100s", "unsorted", [500, 100], codes=wfdbx.NORMAL)
print(wfdbx.read_annotations("100s", "unsorted").samples)      # [100 500]

# A reserved code is refused before anything is written: the file is unchanged.
try:
    wfdbx.write_annotations("100s", "unsorted", [100, 200], codes=[wfdbx.NORMAL, 61])
except ValueError as e:
    print("ValueError:", e)   # ValueError: annotation code 61 is reserved (59 to 63 ...)
print(wfdbx.read_annotations("100s", "unsorted").samples)      # [100 500]
\end{lstlisting}

% ------------------------------------------------------------------------------------------
\subsection{Checking records}
\label{sec:python-checking}

\apientry{check\_record}{Check a record and its annotation files}
\begin{lstlisting}[language=Python,style=signature]
check_record(record, annotators=(), *, input_dir=None, network=None)
\end{lstlisting}
Checks the header of a record, its signal specifications, every sample of its signal
files, each segment of a multi-segment record, and the annotation files named, and
reports each problem found. A problem in the record never raises an exception: it is
reported, with a severity (\code{"error"}: the record is damaged or does not follow the
WFDB format, and programs will read wrong data or fail; \code{"warning"}: the record can
be read, but something is missing or unusual, and may give wrong results; \code{"info"}:
a fact worth knowing, such as invalid samples). The problems, their codes and their
messages are the same in the seven languages (\cref{tab:python-issue-codes}). The checks
run in normal mode, even if legacy mode is on, and the library prints no messages during
the call. They are similar to those of Benjamin Moody's \code{wfdbcheck}.
\begin{params}
\item[record] Record name (a string), as for \code{read\_record}: a full path (absolute or
  \code{\textasciitilde}), used as given; a relative path, read from the input directory
  if one is set, otherwise from the current directory; or a URL, read only if network
  access is on.
\item[annotators] The annotator names whose files are checked, such as \code{["atr"]}; one
  name may be given as a string. Default \code{()}: none.
\item[input\_dir] The input directory for this call only (optional; as for
  \code{read\_record}), under which relative record names are read; it does not apply to
  full paths or URLs. Default \code{None}: the current setting.
\item[network] \code{True} to allow reading URLs for this call only, \code{False} to forbid
  it. Default \code{None}: the current setting.
\end{params}
\returns{A \code{CheckReport} (\cref{tab:python-checkreport}): \code{report.ok} is
\code{True} if no problem is an error, \code{report.issues} lists the problems
(\code{Issue} objects, \cref{tab:python-issue}), and \code{print(report)} prints one line
per problem, or \code{no problems found}.}
\raises{\code{TypeError} if \code{record} or an annotator name is not a string;
\code{ValueError} if \code{record} is empty. A record that cannot be read is reported
(code \code{header\_unreadable}), not raised.}
\begin{lstlisting}[language=Python,style=example]
import numpy as np
import wfdbx

report = wfdbx.check_record("100s", ["atr"])
print(report.ok)          # True: no errors
print(report)
# warning: signal 0: the units are not specified
# warning: signal 1: the units are not specified

# A damaged record: one sample of a 12-bit signal is overwritten with 2048.
wfdbx.write_record("damaged", np.arange(100), fs=100, gains=100, adcres=12, units="mV",
                   names="ramp")
with open("damaged.dat", "r+b") as f:
    f.seek(20)                    # sample 10 (2 bytes per sample in format 16)
    f.write(b"\x00\x08")          # 0x0800 = 2048, little-endian
report = wfdbx.check_record("damaged")
print(report.ok)          # False
for issue in report.issues:
    print(issue.severity, issue.code, issue.where)
# error sample_range signal 0
# error checksum signal 0
print(report.lines()[0])
# error: signal 0: 1 sample is outside the ADC range (-2048 to 2047): sample 10 (value 2048)
\end{lstlisting}

% ------------------------------------------------------------------------------------------
\subsection{Header metadata}
\label{sec:python-metadata}

These functions read and change the metadata of a header: its info strings of the form
\code{\# Key: value}, such as \code{\# Age: 63} (\cref{chap:metadata}). In short: keys are
compared without regard to letter case; a key must not be empty and must not contain
\code{:}, \code{\#}, a line break or a NUL character, and a value must not contain a line
break or a NUL character; spaces at either end of keys and values are removed; text is
one byte per character (Latin-1), and \code{bytes} are also accepted, as given. A change
rewrites the header in one step, through a temporary file, keeping every other byte of
it; all the arguments are checked first, so an error leaves the header unchanged. Only
the \file{.hea} file is read (not a \file{.info} file); for a multi-segment record, it is
the master header, unless a segment is named. Metadata often holds information about
patients: WFDB-x stores it as plain text, exactly as given, and does not de-identify it.

\begin{phibox}
A header can hold identifying data about a patient (a name, a date of birth, a medical
record number). The metadata functions read and write such data as given. Records that
are shared must be de-identified first.
\end{phibox}

\apientry{read\_metadata}{Read the metadata of a record}
\begin{lstlisting}[language=Python,style=signature]
read_metadata(record, *, input_dir=None, network=None)
\end{lstlisting}
Reads the metadata entries of the header of a record, in file order. An info string is an
entry if it contains a colon: the key is the text before the first colon, the value the
rest of the line, both without spaces at either end. Info strings without a colon are free
comments, which are not returned.
\begin{params}
\item[record] Record name (a string), without the \file{.hea} extension (a trailing
  \code{".hea"} is also accepted): a full path, absolute (\code{"/data/mitdb/100"}) or in
  the home directory (\code{"\textasciitilde/Documents/100s"}), used as given; a relative
  path (\code{"100s"}, \code{"data/100s"}), read from the input directory if one is set,
  otherwise from the current directory; or a URL
  (\code{"https://physionet.org/files/mitdb/1.0.0/100"}), read only if network access is
  on.
\item[input\_dir] The input directory for this call only (optional): a local directory, a
  \code{\textasciitilde} path or a URL, under which relative record names are read. It
  does not apply to full paths or URLs. Default \code{None}: the current setting.
\item[network] \code{True} to allow reading URLs for this call only, \code{False} to
  forbid it. Default \code{None}: the current setting.
\end{params}
\returns{A list of (key, value) pairs of strings, in file order. A key that appears more
than once (ignoring letter case) is returned once, with its first value, and a
\code{UserWarning} is issued (\cref{sec:python-warnings}).}
\raises{\code{IOError} if the header cannot be found or read; \code{ValueError} if
\code{record} is empty or \code{"-"} (the standard input).}
\begin{lstlisting}[language=Python,style=example]
import numpy as np
import wfdbx

wfdbx.write_record("pat", np.zeros(360, dtype=int), fs=360, gains=200, adcres=12,
                   units="mV", names="ECG", comments=["recorded at rest"],
                   metadata={"Age": "63", "Sex": "F"})
print(wfdbx.read_metadata("pat"))      # [('Age', '63'), ('Sex', 'F')]
print(wfdbx.read_header("pat").comments)
# ['recorded at rest', ' Age: 63', ' Sex: F']: the free comment is not an entry
\end{lstlisting}

\apientry{get\_metadata}{Return the value of a metadata entry}
\begin{lstlisting}[language=Python,style=signature]
get_metadata(record, key, *, input_dir=None, network=None)
\end{lstlisting}
Returns the value of one metadata entry of the header of a record.
\begin{params}
\item[record] Record name (as for \code{read\_metadata}).
\item[key] The key, such as \code{"Age"}. Letter case is ignored, and spaces at either end
  are removed.
\item[input\_dir] The input directory for this call only (optional; as for
  \code{read\_metadata}): it applies to relative record names only. Default \code{None}:
  the current setting.
\item[network] \code{True} to allow reading URLs for this call only, \code{False} to
  forbid it. Default \code{None}: the current setting.
\end{params}
\returns{The value, a string (\code{""} for an entry with an empty value), or
\code{None} if the header has no entry for \code{key}. If it has more than one, the first
is returned, and a \code{UserWarning} is issued.}
\raises{\code{IOError} if the header cannot be found or read; \code{ValueError} if
\code{record} is empty or \code{"-"}, or if \code{key} is not a valid key;
\code{TypeError} if \code{key} is not a string (or \code{bytes}).}
\begin{lstlisting}[language=Python,style=example]
import numpy as np
import wfdbx

wfdbx.write_record("pat", np.zeros(360, dtype=int), fs=360, gains=200, adcres=12,
                   units="mV", metadata={"Age": "63", "Sex": "F"})
print(wfdbx.get_metadata("pat", "age"))      # 63
print(wfdbx.get_metadata("pat", "Height"))   # None
\end{lstlisting}

\apientry{set\_metadata}{Set the value of a metadata entry}
\begin{lstlisting}[language=Python,style=signature]
set_metadata(record, key, value, *, input_dir=None, network=None, dir=None)
\end{lstlisting}
Sets the value of a metadata entry in the header of a record. If the header has an entry
for \code{key} (ignoring letter case), its line is replaced by \code{\# Key: value}, at
the same position, with the spelling of \code{key} given here; otherwise the entry is
added after the last info string of the header (or, if it has none, after its last line
that is not blank). The entry is written as \code{\# Key: value}, or \code{\# Key:} if
the value is empty. Every other line is kept, byte for byte, including its line ending.
\begin{params}
\item[record] Record name (as for \code{read\_metadata}). The header is changed where it
  was read, unless \code{dir} is given; a header read from a URL can only be written
  elsewhere, with \code{dir}.
\item[key] The key, such as \code{"Age"} or \code{"Date of birth"}. Spaces at either end
  are removed. It must not be empty, and must not contain \code{:}, \code{\#}, a line
  break or a NUL character.
\item[value] The value: a string, which may be a sentence. Spaces at either end are
  removed. It must not contain a line break or a NUL character; it may be empty.
\item[input\_dir] The input directory in which the header is read, for this call only
  (optional; as for \code{read\_metadata}): it applies to relative record names only.
  Default \code{None}: the current setting.
\item[network] \code{True} to allow reading the header from a URL, for this call only,
  \code{False} to forbid it. Default \code{None}: the current setting.
\item[dir] A directory, which must exist, in which to write the changed header, under the
  same file name, instead of changing the header itself; the signal files are not copied.
  Default \code{None}: the header itself is changed.
\end{params}
\returns{\code{None}.}
\raises{\code{IOError} if the header cannot be read, or the changed header cannot be
written; \code{ValueError} if \code{key} or \code{value} is not valid, if the header has
more than one entry for \code{key}, if it has no record line, if it was read from a URL
and \code{dir} is not given, or if \code{record} is empty or \code{"-"};
\code{TypeError} if \code{key} or \code{value} is not a string (or \code{bytes});
\code{FileNotFoundError} if \code{dir} does not exist.}
\begin{lstlisting}[language=Python,style=example]
import os
import numpy as np
import wfdbx

wfdbx.write_record("pat", np.zeros(360, dtype=int), fs=360, gains=200, adcres=12,
                   units="mV", metadata={"Age": "63", "Sex": "F"})
wfdbx.set_metadata("pat", "AGE", "64")              # replaces "# Age: 63"
wfdbx.set_metadata("pat", "Diagnosis", "sinus rhythm, no ectopy")   # a new entry
print(wfdbx.read_metadata("pat"))
# [('AGE', '64'), ('Sex', 'F'), ('Diagnosis', 'sinus rhythm, no ectopy')]

# Write the changed header to another directory; pat.hea is left unchanged.
os.makedirs("changed", exist_ok=True)
wfdbx.set_metadata("pat", "Sex", "M", dir="changed")
print(wfdbx.get_metadata("pat", "Sex"), wfdbx.get_metadata("changed/pat", "Sex"))   # F M
\end{lstlisting}

\apientry{remove\_metadata}{Remove a metadata entry}
\begin{lstlisting}[language=Python,style=signature]
remove_metadata(record, key, *, input_dir=None, network=None, dir=None)
\end{lstlisting}
Removes every entry for \code{key} (ignoring letter case) from the header of a record, so
it also repairs a header that has more than one. Other lines are kept, byte for byte.
Nothing is written if the header has no entry for \code{key}, unless \code{dir} is given.
\begin{params}
\item[record] Record name (as for \code{set\_metadata}).
\item[key] The key (as for \code{set\_metadata}).
\item[input\_dir] The input directory in which the header is read, for this call only
  (optional; as for \code{read\_metadata}): it applies to relative record names only.
  Default \code{None}: the current setting.
\item[network] \code{True} to allow reading the header from a URL, for this call only,
  \code{False} to forbid it. Default \code{None}: the current setting.
\item[dir] A directory, which must exist, in which to write the changed header instead of
  changing the header itself. Default \code{None}: the header itself is changed.
\end{params}
\returns{The number of lines removed (an \code{int}; 0 if the header has no entry for
\code{key}).}
\raises{As \code{set\_metadata}, except that more than one entry for \code{key} is not an
error.}
\begin{lstlisting}[language=Python,style=example]
import numpy as np
import wfdbx

wfdbx.write_record("pat", np.zeros(360, dtype=int), fs=360, gains=200, adcres=12,
                   units="mV", metadata={"Age": "63", "Sex": "F"})
print(wfdbx.remove_metadata("pat", "sex"))      # 1
print(wfdbx.remove_metadata("pat", "Sex"))      # 0: there is no such entry now
print(wfdbx.read_metadata("pat"))               # [('Age', '63')]
\end{lstlisting}

\apientry{update\_metadata}{Make several changes to the metadata}
\begin{lstlisting}[language=Python,style=signature]
update_metadata(record, changes, *, input_dir=None, network=None, dir=None)
\end{lstlisting}
Makes several changes to the metadata of a record, in one write of the header. Each
change is applied in order, as \code{set\_metadata} does, or, for a value of
\code{None}, as \code{remove\_metadata} does. All the changes are checked first: if one is
not valid, none is made.
\begin{params}
\item[record] Record name (as for \code{set\_metadata}).
\item[changes] A mapping from key to value, or a sequence of (key, value) pairs; a value
  of \code{None} removes the entry. No two keys may differ only in letter case.
\item[input\_dir] The input directory in which the header is read, for this call only
  (optional; as for \code{read\_metadata}): it applies to relative record names only.
  Default \code{None}: the current setting.
\item[network] \code{True} to allow reading the header from a URL, for this call only,
  \code{False} to forbid it. Default \code{None}: the current setting.
\item[dir] A directory, which must exist, in which to write the changed header instead of
  changing the header itself. Default \code{None}: the header itself is changed.
\end{params}
\returns{\code{None}.}
\raises{As \code{set\_metadata}; also \code{ValueError} if two keys of \code{changes}
differ only in letter case, and \code{TypeError} if \code{changes} is neither a mapping
nor a sequence of pairs.}
\begin{lstlisting}[language=Python,style=example]
import numpy as np
import wfdbx

wfdbx.write_record("pat", np.zeros(360, dtype=int), fs=360, gains=200, adcres=12,
                   units="mV", metadata={"Age": "63", "Sex": "F"})
wfdbx.update_metadata("pat", {"Age": "64", "Sex": None, "Dx": "164889003"})
print(wfdbx.read_metadata("pat"))      # [('Age', '64'), ('Dx', '164889003')]
try:
    wfdbx.update_metadata("pat", [("Dx", "426783006"), ("dx", None)])
except ValueError as e:
    print(e)        # metadata key 'dx' is given more than once (...)
print(wfdbx.get_metadata("pat", "Dx"))  # 164889003: nothing was changed
\end{lstlisting}

% ------------------------------------------------------------------------------------------
\subsection{The sample record and the example notebooks}
\label{sec:python-resource-functions}

\apientry{sample\_data\_dir}{The directory of the sample record}
\begin{lstlisting}[language=Python,style=signature]
sample_data_dir()
\end{lstlisting}
Returns the directory, within the installed package, that holds the sample record
\code{100s} (\cref{sec:python-resources}): its header, signal file and reference
annotations.
The function has no parameters.
\returns{The absolute path of the directory, a string.}
\raises{Nothing.}
\begin{lstlisting}[language=Python,style=example]
import os
import wfdbx

d = wfdbx.sample_data_dir()
print(sorted(f for f in os.listdir(d) if f.startswith("100s")))
# ['100s.atr', '100s.dat', '100s.hea']
rec = wfdbx.read_record("100s", sampto=10, input_dir=d)
print(rec.signals.shape)        # (10, 2)
\end{lstlisting}

\apientry{examples\_dir}{The directory of the example notebooks}
\begin{lstlisting}[language=Python,style=signature]
examples_dir()
\end{lstlisting}
Returns the directory, within the installed package, that holds the example notebooks
(\file{*.ipynb}, \cref{tab:python-notebooks}). The notebooks are meant to be copied
elsewhere (\code{copy\_examples}) before they are run.
The function has no parameters.
\returns{The absolute path of the directory, a string.}
\raises{Nothing.}
\begin{lstlisting}[language=Python,style=example]
import os
import wfdbx

print(sorted(os.listdir(wfdbx.examples_dir()))[0])     # 01_getting_started.ipynb
\end{lstlisting}

\apientry{copy\_examples}{Copy the example notebooks}
\begin{lstlisting}[language=Python,style=signature]
copy_examples(dest)
\end{lstlisting}
Copies the example notebooks to a directory, created if needed, where they can be opened
and run. Existing files of the same names are replaced. The command \cmd{wfdbx -{}-examples
DIR} does the same (\cref{sec:python-tools}).
\begin{params}
\item[dest] The directory into which the notebooks are copied (a path).
\end{params}
\returns{The list of the paths of the copied notebooks (\code{os.path.join(dest, name)}),
in the order of their names.}
\raises{\code{OSError} if the directory cannot be created or a notebook cannot be
written.}
\begin{lstlisting}[language=Python,style=example]
import os
import wfdbx

paths = wfdbx.copy_examples("notebooks")
print(len(paths), os.path.basename(paths[0]))    # 7 01_getting_started.ipynb
\end{lstlisting}

% ------------------------------------------------------------------------------------------
\subsection{The result classes}
\label{sec:python-results}

The readers return objects of the classes \code{Record}, \code{Header} and
\code{Annotations}, whose attributes are listed in
\cref{tab:python-record,tab:python-header,tab:python-annotations}; \code{check\_record}
returns a \code{CheckReport}, which holds \code{Issue} objects
(\cref{tab:python-checkreport,tab:python-issue}). The signals of a record are described
by \code{Siginfo} objects (\cref{tab:python-siginfo}), which are those of the WFDB library
(\cref{sec:python-library}). All these classes are available at the top level of the
package (\code{wfdbx.Record}, \ldots). The attributes of \code{Record}, \code{Header} and
\code{Annotations} are plain Python attributes, which may be changed.

\begin{xltabular}{\linewidth}{@{}lL@{}}
\caption{The attributes of \code{Record}, the result of \code{read\_record}}
\label{tab:python-record} \\
\toprule
\tabhead Attribute & \tabhead Contents \\
\midrule
\endfirsthead
\toprule
\tabhead Attribute & \tabhead Contents \\
\midrule
\endhead
\bottomrule
\endfoot
\code{name} & the record name, as given to \code{read\_record} \\
\code{fs} & the sampling frequency in Hz (a \code{float}), or 0 if the header gives none;
  with \code{highres=True}, the highest sampling frequency of a multi-frequency record \\
\code{signals} & a NumPy array with one row per sample and one column per signal read:
  \code{int32} ADC units, or \code{float64} physical units with \code{physical=True};
  invalid samples are \code{WFDB\_INVALID\_SAMPLE}, or NaN in physical units \\
\code{siginfo} & a list of \code{Siginfo} objects, one per signal read, in the order of
  the columns; \code{desc} has no white space at either end, and \code{adcres} is 0 if
  the header does not give it \\
\code{names} & the description of each signal, such as \code{"MLII"}, without the white
  space that a header may put at either end of it \\
\code{units} & the physical units of each signal, or \code{None} where the header gives
  none \\
\code{gains} & the gain of each signal, in ADC units per physical unit (\code{float}); 0
  means that the signal is uncalibrated \\
\code{baselines} & the ADC value of each signal that corresponds to 0 physical units \\
\code{base\_time} & the start time and date as text, to the second, such as
  \code{"[14:00:03 21/06/2024]"} (\code{"[14:00:03]"} without a date), or \code{None} if
  the header gives none \\
\code{base\_time\_ms} & the start time in milliseconds after midnight, such as
  \code{50403500} for 14:00:03.500, or \code{None} if \code{base\_time} is \code{None} \\
\code{base\_date} & the start date as a tuple \code{(year, month, day)}, or \code{None}
  if the header gives none \\
\code{sampfrom} & the number of the first sample returned \\
\code{comments} & the list of the info strings of the header (without the \code{\#}): its
  lines that begin with \code{\#} after the record line, found as the header parser finds
  it (\cref{sec:info-strings}) \\
\code{checksum\_ok} & \code{False} if the samples do not match the checksums of the
  header; checksums are checked only when a read reaches the end of the record, so a
  shorter read gives \code{True} \\
\end{xltabular}

\begin{xltabular}{\linewidth}{@{}lL@{}}
\caption{The attributes of \code{Header}, the result of \code{read\_header}}
\label{tab:python-header} \\
\toprule
\tabhead Attribute & \tabhead Contents \\
\midrule
\endfirsthead
\toprule
\tabhead Attribute & \tabhead Contents \\
\midrule
\endhead
\bottomrule
\endfoot
\code{name} & the record name, as given to \code{read\_header} \\
\code{fs} & the sampling frequency in Hz, or 0 if the header gives none \\
\code{counter\_freq} & the counter frequency, as \code{getcfreq()} returns it (the
  sampling frequency if the header gives none) \\
\code{base\_counter} & the base counter value, as \code{getbasecount()} returns it \\
\code{nsamples} & the number of samples per signal, or 0 if the header does not give it \\
\code{nsig} & the number of signals \\
\code{siginfo} & a list of \code{Siginfo} objects, one per signal; \code{desc} has no white
  space at either end, and \code{adcres} is 0 if the header does not give it \\
\code{base\_time}, \code{base\_time\_ms}, \code{base\_date} & the start time and date, as
  in \code{Record} \\
\code{comments} & the list of the info strings of the header (without the \code{\#}), as
  in \code{Record} \\
\code{metadata} & the list of the (key, value) pairs of the metadata entries of the
  header, in file order, as \code{read\_metadata} returns them \\
\end{xltabular}

\begin{xltabular}{\linewidth}{@{}lL@{}}
\caption{The attributes of \code{Annotations}, the result of \code{read\_annotations}.
  \code{len(ann)} is the number of annotations, which are in file order.}
\label{tab:python-annotations} \\
\toprule
\tabhead Attribute & \tabhead Contents \\
\midrule
\endfirsthead
\toprule
\tabhead Attribute & \tabhead Contents \\
\midrule
\endhead
\bottomrule
\endfoot
\code{record} & the record name, as given \\
\code{annotator} & the annotator name, such as \code{"atr"} \\
\code{samples} & the annotation times, in samples (an \code{int64} array) \\
\code{codes} & the annotation codes (an \code{int64} array), such as \code{wfdbx.NORMAL}
  or \code{wfdbx.PVC} \\
\code{symbols} & the list of the mnemonics of the codes, such as \code{"N"}, \code{"V"}
  or \code{"+"} \\
\code{subtypes} & the subtype of each annotation (an \code{int64} array) \\
\code{chans} & the channel number of each annotation (an \code{int64} array) \\
\code{nums} & the annotator number of each annotation (an \code{int64} array) \\
\code{aux} & the list of the auxiliary texts of the annotations (strings), with
  \code{None} where there is none \\
\code{fs} & the frequency in which \code{samples} are counted, normally the sampling
  frequency of the record; 0 if neither the record, nor the file, nor the argument
  \code{fs} of \code{read\_annotations} gives one \\
\end{xltabular}

\begin{table}[htbp]
\centering
\caption{The attributes and methods of \code{CheckReport}, the result of
  \code{check\_record}}
\label{tab:python-checkreport}
\begin{tabularx}{\linewidth}{@{}lL@{}}
\toprule
\tabhead Attribute & \tabhead Contents \\
\midrule
\code{record} & the record name, as given \\
\code{ok} & \code{True} if no issue has the severity \code{"error"} \\
\code{issues} & the list of the \code{Issue} objects, in the order of the checks \\
\code{lines()} & the text form of the report, as a list of lines: one line per issue
  (\code{str(issue)}), or the one line \code{no problems found} \\
\code{str(report)} & the lines of \code{lines()}, joined by line feeds (what
  \code{print(report)} prints) \\
\bottomrule
\end{tabularx}
\end{table}

\begin{table}[htbp]
\centering
\caption{The attributes of \code{Issue}, a problem found by \code{check\_record}}
\label{tab:python-issue}
\begin{tabularx}{\linewidth}{@{}lL@{}}
\toprule
\tabhead Attribute & \tabhead Contents \\
\midrule
\code{severity} & \code{"error"}, \code{"warning"} or \code{"info"} \\
\code{code} & a stable identifier of the kind of problem, such as \code{"checksum"} or
  \code{"gain\_missing"} (\cref{tab:python-issue-codes}) \\
\code{where} & where the problem is: \code{""} (the record), \code{"signal N"},
  \code{"segment SEG"}, \code{"segment SEG, signal N"} or \code{"annotator ANN"} \\
\code{message} & the description of the problem \\
\code{str(issue)} & \code{"SEVERITY: WHERE: MESSAGE"}, or \code{"SEVERITY: MESSAGE"}
  when \code{where} is empty \\
\bottomrule
\end{tabularx}
\end{table}
Two \code{Issue} objects are equal (\code{==}) when their four attributes are.

\begin{xltabular}{\linewidth}{@{}llL@{}}
\caption{The codes of the issues reported by \code{check\_record}}
\label{tab:python-issue-codes} \\
\toprule
\tabhead Code & \tabhead Severity & \tabhead Problem \\
\midrule
\endfirsthead
\toprule
\tabhead Code & \tabhead Severity & \tabhead Problem \\
\midrule
\endhead
\bottomrule
\endfoot
\multicolumn{3}{@{}l}{\textit{The header}} \\
\code{header\_unreadable} & error & the header cannot be read \\
\code{header\_line\_long} & error & a line of the header is longer than 255 characters \\
\code{header\_not\_ascii} & warning & a line of the header contains non-ASCII characters \\
\code{header\_name} & warning & the header names a record other than its file name \\
\code{fs\_missing} & warning & the sampling frequency is not specified \\
\code{length\_missing} & warning & the record length is not specified \\
\code{length\_large} & info & the record has more than 2147483647 samples \\
\code{no\_signals} & info & the record has no signals \\
\multicolumn{3}{@{}l}{\textit{The signal specifications}} \\
\code{desc\_missing} & warning & the signal has no description \\
\code{desc\_number} & warning & the description is only a number \\
\code{desc\_control} & error & the description contains control characters \\
\code{desc\_duplicate} & warning & several signals have the same description \\
\code{gain\_missing} & warning & the signal is uncalibrated (no gain is given) \\
\code{gain\_negative} & info & the gain is negative \\
\code{units\_missing} & warning & the units are not specified \\
\code{units\_control} & error & the units contain control characters \\
\code{adcres\_missing} & warning & the ADC resolution is not specified \\
\code{adcres\_too\_large} & error & the ADC resolution exceeds what the format can store \\
\code{adc\_range} & error & the ADC range does not fit in the format \\
\code{format\_difference} & warning & format 8 (first differences) cannot follow steep
  slopes \\
\code{filename} & warning & the signal file name contains characters other than letters,
  digits, \code{\_}, \code{-} and \code{.} \\
\multicolumn{3}{@{}l}{\textit{The signal files}} \\
\code{file\_unreadable} & error & a signal file cannot be read \\
\code{file\_large} & warning & a signal file has more than 2147483647 bytes \\
\code{read\_error} & error & the signal data cannot be read after some frames \\
\code{sample\_range} & error & samples are outside the ADC range \\
\code{invalid\_samples} & info & samples are invalid (missing data) \\
\code{checksum} & error & the checksum of a signal does not match its samples \\
\code{length\_short} & error & the signal files end before the length given in the
  header \\
\code{length\_long} & warning & the signal files hold more than the length given in the
  header \\
\code{length\_unspecified\_frames} & info & the number of frames that the signal files
  hold, when the header does not give the length \\
\multicolumn{3}{@{}l}{\textit{The segments of a multi-segment record}} \\
\code{segment\_unreadable} & error & a segment cannot be read \\
\code{nested} & error & a segment is itself a multi-segment record \\
\code{segment\_fs} & error & the sampling frequency of a segment differs from the
  record's \\
\code{segment\_lengths} & warning & the segment lengths do not add up to the length in
  the master header \\
\code{layout\_signals} & error & a segment has more signals than the record \\
\code{layout\_mismatch} & error & a signal of a segment differs from the record's
  signal \\
\code{layout\_missing\_signal} & error & a signal of a segment is not in the layout \\
\multicolumn{3}{@{}l}{\textit{The annotation files}} \\
\code{annotator\_name} & warning & the annotator name consists only of digits \\
\code{annotations\_unreadable} & error & the annotation file cannot be read \\
\code{annotations\_truncated} & error & the annotation file ends unexpectedly \\
\code{annotations\_order} & error & annotations are out of order \\
\code{annotation\_code} & error & annotations have an invalid code \\
\code{annotation\_null} & warning & annotations have code 0 (\code{NOTQRS}) \\
\code{annotation\_undefined} & warning & annotations have a code without a defined
  mnemonic \\
\code{annotations\_after\_end} & warning & annotations are after the end of the
  record \\
\end{xltabular}

\begin{xltabular}{\linewidth}{@{}lL@{}}
\caption{The attributes of \code{Siginfo}, the specification of a signal}
\label{tab:python-siginfo} \\
\toprule
\tabhead Attribute & \tabhead Contents \\
\midrule
\endfirsthead
\toprule
\tabhead Attribute & \tabhead Contents \\
\midrule
\endhead
\bottomrule
\endfoot
\code{fname} & the name of the signal file, as the header gives it \\
\code{desc} & the signal description, such as \code{"MLII"}, or \code{None} \\
\code{units} & the physical units, such as \code{"mV"}, or \code{None} if not specified \\
\code{gain} & the gain, in ADC units per physical unit (a \code{float}; 0: uncalibrated) \\
\code{initval} & the value of the first sample \\
\code{group} & the signal group number (the signals of one group share one file) \\
\code{fmt} & the storage format, such as \code{16}, \code{212} or \code{516} \\
\code{spf} & the number of samples per frame (more than 1 for the faster signals of a
  multi-frequency record) \\
\code{bsize} & the block size, in bytes (0 for ordinary files) \\
\code{adcres} & the ADC resolution, in bits (0: not specified) \\
\code{adczero} & the ADC output for an input of 0\,V \\
\code{baseline} & the ADC value that corresponds to 0 physical units \\
\code{nsamp} & the number of samples (0 if unknown) \\
\code{cksum} & the 16-bit checksum of the samples \\
\end{xltabular}

% ------------------------------------------------------------------------------------------
\subsection{Warnings and errors}
\label{sec:python-warnings}

\paragraph{Exceptions.} The high-level functions report errors by exceptions of the
standard classes (\cref{tab:python-exceptions}). Each function checks its arguments
before it creates or changes a file, so an invalid argument leaves the files as they
were; \code{write\_record} and \code{write\_annotations} also write through temporary
files, so that a failure while writing (a full disk, for example) leaves the existing
files as they were too. The library may also print a message on the standard error when
it cannot open a file (as in \code{init: can't open header for record REC}); \code{wfdbquiet()} turns these
messages off (\cref{sec:python-library}).

\begin{table}[htbp]
\centering
\caption{The exceptions of the high-level API}
\label{tab:python-exceptions}
\begin{tabularx}{\linewidth}{@{}lL@{}}
\toprule
\tabhead Exception & \tabhead Raised when \\
\midrule
\code{IOError} (\code{OSError}) & a header, signal file or annotation file cannot be
  found, read, created or written completely (for example on a full disk); a signal or
  annotation file ends before its header or its end-of-file marker says \\
\code{FileNotFoundError} & the directory given as \code{dir} does not exist (a subclass of
  \code{OSError}) \\
\code{ValueError} & an argument is invalid, or the operation needs a value that the
  record does not give (a time in seconds without a sampling frequency, physical units
  without a gain) \\
\code{TypeError} & an argument has the wrong type (a record or annotator name that is
  not a string in \code{check\_record}; a metadata key or value that is not a string;
  metadata that is neither a mapping nor a sequence of pairs) \\
\bottomrule
\end{tabularx}
\end{table}

\paragraph{The checksum warning.} When \code{read\_record} reads to the end of a record
whose samples do not match the checksums of its header, it returns all the samples and
issues a \code{RuntimeWarning}: ``checksum error in record REC: the signal data do not
match the checksums in its header''. \code{checksum\_warning=False} turns the warning off;
\code{checksum\_ok} of the result is \code{False} either way. A signal file that ends
early, or cannot be read, is an \code{IOError}, not a checksum warning, even when the
checksum is verified in its last, incomplete frame: no incomplete frame is returned.

\paragraph{The duplicate-key warning.} When a header holds two metadata entries whose keys
differ only in letter case (written by other software), \code{read\_metadata},
\code{get\_metadata} and \code{read\_header} use the first and issue a \code{UserWarning}:
``header of record REC has more than one entry for key KEY (lines L1, L2); the first is
used''. \code{set\_metadata} and \code{update\_metadata} then refuse to change that key
(\code{ValueError}), and \code{remove\_metadata} removes all of its entries.

Both warnings are issued with the module \code{warnings} of Python, so they can be
filtered, recorded or turned into errors (\code{warnings.simplefilter("error",
RuntimeWarning)}) as any other warning:

\begin{lstlisting}[language=Python,style=example]
import warnings
import numpy as np
import wfdbx

# A record whose last sample is changed after it was written: the checksum fails.
wfdbx.write_record("ck", np.arange(100), fs=100, gains=100, adcres=12, units="mV")
with open("ck.dat", "r+b") as f:
    f.seek(-2, 2)
    f.write(b"\x00\x00")
with warnings.catch_warnings(record=True) as caught:
    warnings.simplefilter("always")
    rec = wfdbx.read_record("ck")
print(rec.signals.shape, rec.checksum_ok)    # (100, 1) False: all the samples
print(caught[0].category.__name__)           # RuntimeWarning
rec = wfdbx.read_record("ck", checksum_warning=False)    # no warning

# Two entries for one key: the first is used, with a warning.
wfdbx.write_record("dup", np.zeros(10, dtype=int), fs=100, gains=0, adcres=12,
                   comments=[" Age: 63", " AGE: 64"])
with warnings.catch_warnings(record=True) as caught:
    warnings.simplefilter("always")
    print(wfdbx.get_metadata("dup", "age"))      # 63
print(caught[0].category.__name__)               # UserWarning
try:
    wfdbx.set_metadata("dup", "Age", "65")
except ValueError as e:
    print("ValueError:", e)
print(wfdbx.remove_metadata("dup", "age"))       # 2
\end{lstlisting}

\begin{lstlisting}[language=Python,style=example]
import numpy as np
import wfdbx

try:
    wfdbx.read_record("nosuchrecord")
except OSError as e:                  # IOError is OSError
    print("OSError:", e)              # OSError: can't read record nosuchrecord

wfdbx.write_record("raw", np.arange(10), fs=100, gains=0, adcres=12)   # uncalibrated
try:
    wfdbx.read_record("raw", physical=True)
except ValueError as e:
    print("ValueError:", e)           # ValueError: signal 0 of record raw is uncalibrated
\end{lstlisting}

% ==========================================================================================
\section{The WFDB library functions in Python}
\label{sec:python-library}

Every function of the WFDB library is available at the top level of the package under its
C name, with the same arguments, in the same order, and the same return values, so that
programs and documentation written for the C library carry over directly. The functions,
their arguments and their results are documented once, for all the languages, in
\cref{chap:library}; this section gives the rules by which the C declarations become
Python functions. \code{help(wfdbx.NAME)} gives the documentation of each function,
including what WFDB-x changes in it.

\subsection{Conventions}

\paragraph{Names.} The functions, the constants (\code{WFDB\_READ},
\code{WFDB\_INVALID\_SAMPLE}, \code{WFDB\_HIGHRES}, \ldots) and the annotation codes
(\code{NORMAL}, \code{PVC}, \ldots, \cref{chap:formats}) keep their C names. The macros of
\file{<wfdb/ecgmap.h>} are functions, under both their macro names and the names of the
functions behind them: \code{isqrs} and \code{wfdb\_isqrs}, \code{map1} and
\code{wfdb\_map1}, and so on. The structures are classes
(\cref{sec:python-structures}).

\paragraph{NULL is \code{None}.} A \code{NULL} pointer argument or result is
\code{None}: \code{sampfreq(None)} returns the frequency of the open record,
\code{getinfo(None)} continues with the info strings of the last record, and
\code{getinfo} returns \code{None} after the last one. Some trailing arguments have
defaults (\code{nann=None} and \code{nsig=None}: all the entries of the list):
\begin{lstlisting}[language=Python]
isigopen(record, siarray=None, nsig=0)
osigopen(record, siarray=None, nsig=0)
osigfopen(siarray, nsig=None)
setheader(record, siarray, nsig=None)
annopen(record, aiarray, nann=None)
\end{lstlisting}

\paragraph{Arrays and output arguments.} Where a C function fills an array that the caller
passes, the Python function fills a list that the caller passes:
\begin{itemize}
\item \code{isigopen} and \code{osigopen} empty the list \code{siarray} and append one
  \code{Siginfo} object (a copy) per signal; \code{siarray} may be \code{None}. With
  \code{nsig} $\le 0$, \code{isigopen} reads the header only, and returns the number of
  signals of the record; \code{osigopen(record, None, 0)} closes the open output signals,
  reads the header, opens nothing and returns 0.
\item \code{getvec}, \code{getframe} and \code{putvec} take a list (or a NumPy array) of
  integers, one per signal (\code{getframe} and \code{putvec}: one per sample of a
  frame); \code{getvec} and \code{getframe} replace its first entries, and it must have
  room for them (\code{[0] * nsig}), or \code{IndexError} is raised.
\item A C function that fills a structure through a pointer replaces the attributes of
  the object passed: \code{getann(n, annot)} fills an \code{Annotation},
  \code{getcal(desc, units, cal)} a \code{Calinfo}. \code{putann(n, annot)} and
  \code{putcal(cal)} read them.
\item An array of structures passed in is a list of objects: \code{annopen(record,
  [Anninfo("atr", WFDB\_READ)])}, \code{osigfopen([Siginfo(...), ...])}.
\item \code{getseginfo()} returns a list of \code{Seginfo} objects (empty if the open
  record is not a multi-segment record), instead of a pointer to an array.
\end{itemize}

\paragraph{Strings} are \code{str}. Text and binary data are handled as ISO-8859-1
(Latin-1), one character per byte, so every byte sequence round-trips unchanged. The
\code{aux} attribute of an annotation is \code{None}, or \code{bytes} whose first byte is
the length of the text that follows, as in C; \code{Annotation.aux\_string()} returns the
text (up to its first NUL byte), or \code{None}.

\paragraph{Return values and errors.} The functions return what the C functions return:
an \code{int} status (negative on errors), a number of signals, a \code{float}, a string
or \code{None}. Errors are reported by these values and by a message on the standard
error, as in C, never by exceptions (only a misuse of Python, such as a vector that is too
short, raises one). \code{wfdberror()} returns the last message; \code{wfdbquiet()} turns
the messages off and \code{wfdbverbose()} on again.

\paragraph{State.} As in C, the library keeps its state in the process (module-level
state): the open record, its input signals, the output signals and the annotators remain
open between calls, until \code{wfdbquit()} closes them. \code{wfdbquit()} leaves the
input and output directories and the network setting as they are, and so, as in C, the
time resolution of output annotators (\code{getafreq()}), which opening an EDF record sets
to its sampling frequency when none is set, and the getvec mode (\code{getgvmode()}), to
which it adds \code{WFDB\_HIGHRES}. The high-level functions of \cref{sec:python-api} call
\code{wfdbquit()}, which closes anything opened with the library functions, and leave
those two settings as they were.

\paragraph{Additions.} \code{getframes(n)} reads up to \code{n} frames at once, as
\code{n} calls of \code{getvec} would, but decoding whole blocks of the signal files. It
returns a tuple \code{(array, stat)}: an \code{int64} NumPy array with one row per frame
read and one column per signal, and the \code{getvec} status of the first frame that
could not be read, or the number of signals if all \code{n} frames were read.

\subsection{Module layout}

All the functions, classes and constants are imported by \code{import wfdbx}. They are
defined in these modules of the package:

\begin{table}[htbp]
\centering
\caption{The modules of the package \code{wfdbx}}
\label{tab:python-modules}
\begin{tabularx}{\linewidth}{@{}lL@{}}
\toprule
\tabhead Module & \tabhead Contents \\
\midrule
\code{wfdbx.\_wfdbio} & file lookup, the input and output directories, network access,
  error messages (\file{lib/wfdbio.c}); the settings functions \\
\code{wfdbx.\_signal} & headers, signal files in every format, time and date
  conversions, unit conversions (\file{lib/signal.c}); legacy mode \\
\code{wfdbx.\_annot} & annotation files, in the MIT and AHA formats
  (\file{lib/annot.c}) \\
\code{wfdbx.\_calib} & the calibration list (\file{lib/calib.c}) \\
\code{wfdbx.\_wfdbinit} & \code{wfdbinit}, \code{wfdbquit}, \code{wfdbflush}
  (\file{lib/wfdbinit.c}) \\
\code{wfdbx.\_ecgcodes} & the annotation codes and the mapping functions
  (\file{lib/ecgcodes.h}, \file{lib/ecgmap.h}) \\
\code{wfdbx.\_flac} & the FLAC codec of the formats 508, 516 and 524 \\
\code{wfdbx.\_cstd} & the exact behavior of the C standard library functions that the WFDB
  library relies on (\code{strtol}, \code{sscanf}, \code{printf}, \ldots) \\
\code{wfdbx.highlevel} & \code{read\_record}, \code{read\_header},
  \code{read\_annotations}, \code{write\_record}, \code{write\_annotations} and their
  result classes \\
\code{wfdbx.metadata} & the header metadata functions \\
\code{wfdbx.check} & \code{check\_record}, \code{CheckReport}, \code{Issue} \\
\code{wfdbx.resources} & \code{sample\_data\_dir}, \code{examples\_dir},
  \code{copy\_examples} \\
\code{wfdbx.cli} & the command \code{wfdbx} \\
\code{wfdbx.apps} & the command-line tools, one module per tool
  (\code{wfdbx.apps.rdsamp}, \ldots) \\
\bottomrule
\end{tabularx}
\end{table}

The modules whose names begin with an underscore are internal: their functions are used
through the top level of the package (\code{wfdbx.isigopen}, not
\code{wfdbx.\_signal.isigopen}).

\subsection{The structures}
\label{sec:python-structures}

The structures of the C library are classes whose attributes have the names of the
fields of the C structures (\cref{tab:python-structures}). Their constructors take the
attributes as arguments, in this order, all with defaults:
\begin{lstlisting}[language=Python]
Siginfo(fname=None, desc=None, units=None, gain=0.0, initval=0, group=0, fmt=0, spf=1,
        bsize=0, adcres=0, adczero=0, baseline=0, nsamp=0, cksum=0)
Seginfo(recname="", nsamp=0, samp0=0)
Anninfo(name=None, stat=WFDB_READ)
Annotation(time=0, anntyp=0, subtyp=0, chan=0, num=0, aux=None)
Calinfo(low=0.0, high=0.0, scale=0.0, sigtype=None, units=None, caltype=0)
\end{lstlisting}
\code{Siginfo}, \code{Annotation} and
\code{Calinfo} have a method \code{copy()}; \code{Annotation} also has
\code{assign(other)}, which copies the attributes of another annotation, and
\code{aux\_string()}.

\begin{xltabular}{\linewidth}{@{}lL@{}}
\caption{The classes of the C structures}
\label{tab:python-structures} \\
\toprule
\tabhead Class (C structure) & \tabhead Attributes \\
\midrule
\endfirsthead
\toprule
\tabhead Class (C structure) & \tabhead Attributes \\
\midrule
\endhead
\bottomrule
\endfoot
\code{Siginfo} (\code{WFDB\_Siginfo}) & \code{fname}, \code{desc}, \code{units},
  \code{gain}, \code{initval}, \code{group}, \code{fmt}, \code{spf}, \code{bsize},
  \code{adcres}, \code{adczero}, \code{baseline}, \code{nsamp}, \code{cksum}
  (\cref{tab:python-siginfo}) \\
\code{Seginfo} (\code{WFDB\_Seginfo}) & \code{recname} (the name of the segment),
  \code{nsamp} (its number of samples), \code{samp0} (the number of its first sample in
  the record) \\
\code{Anninfo} (\code{WFDB\_Anninfo}) & \code{name} (the annotator name, the suffix of the
  file), \code{stat} (\code{WFDB\_READ} or \code{WFDB\_WRITE} for the MIT format,
  \code{WFDB\_AHA\_READ} or \code{WFDB\_AHA\_WRITE} for the AHA format) \\
\code{Annotation} (\code{WFDB\_Annotation}) & \code{time} (in sample intervals),
  \code{anntyp} (the annotation code), \code{subtyp}, \code{chan}, \code{num},
  \code{aux} (\code{None}, or \code{bytes} whose first byte is the length of the text
  that follows) \\
\code{Calinfo} (\code{WFDB\_Calinfo}) & \code{low}, \code{high} (the levels of the
  calibration pulse, in physical units), \code{scale} (physical units per cm),
  \code{sigtype} (matched against the beginning of a signal description), \code{units},
  \code{caltype} (\code{WFDB\_AC\_COUPLED} or \code{WFDB\_DC\_COUPLED}, or-ed with
  \code{WFDB\_CAL\_SQUARE}, \code{WFDB\_CAL\_SINE} or \code{WFDB\_CAL\_SAWTOOTH}) \\
\end{xltabular}

\subsection{The settings functions}
\label{sec:python-settings}

The functions of \cref{tab:python-settings} are additions of WFDB-x to the library: they
set where files are read and written, network access and legacy mode, for the high-level
functions, the library functions and the tools alike (\cref{chap:files,chap:behavior}),
and give the path under which a program opens a file that it reads or creates itself, as
the tools do (\code{wfdbx\_input\_path}, \code{wfdbx\_output\_path};
\cref{sec:tool-files}). Each \code{set} function returns the previous setting, so that it
can be restored. The
C function \code{setwfdb(p)} sets the input directory too (\code{p} must name at most one
directory), and \code{getwfdb()} returns it (\code{"."} if none is set); in legacy mode
they set and return the WFDB path, as in C.

\begin{xltabular}{\linewidth}{@{}lL@{}}
\caption{The settings functions}
\label{tab:python-settings} \\
\toprule
\tabhead Function & \tabhead Effect and result \\
\midrule
\endfirsthead
\toprule
\tabhead Function & \tabhead Effect and result \\
\midrule
\endhead
\bottomrule
\endfoot
\code{wfdbx\_set\_input\_dir(d)} & sets the input directory: a local directory, a
  \code{\textasciitilde} path or a URL; \code{None} or \code{""} restores the default
  (\code{WFDBX\_INPUT\_DIR}, or else the current directory). Returns the previous input
  directory, or \code{None} if none was set \\
\code{wfdbx\_get\_input\_dir()} & returns the input directory, or \code{None} if none is
  set \\
\code{wfdbx\_set\_output\_dir(d)} & sets the output directory: a local directory or a
  \code{\textasciitilde} path (a URL is refused, with a message); \code{None} or
  \code{""} restores the default (\code{WFDBX\_OUTPUT\_DIR}, or else the current
  directory). Returns the previous output directory, or \code{None} (the current one if
  \code{d} is refused) \\
\code{wfdbx\_get\_output\_dir()} & returns the output directory, or \code{None} if none
  is set \\
\code{wfdbx\_set\_network(on)} & turns network access on (\code{True} or 1) or off
  (\code{False} or 0); \code{None} or a negative number restores the default (on if
  \code{WFDBX\_NETWORK} is a nonzero number). Returns the previous setting, 1 or 0 \\
\code{wfdbx\_get\_network()} & returns 1 if network access is on, 0 otherwise \\
\code{wfdbx\_set\_legacy(on)} & turns legacy mode on (\code{True} or 1) or off, for the
  whole library. Returns the previous setting, 1 or 0 \\
\code{wfdbx\_get\_legacy()} & returns 1 if legacy mode is on, 0 otherwise \\
\code{wfdbx\_legacy\_args(argv)} & if \code{argv[1]} (\code{argv}: the argument list,
  program name first) is \code{"-{}-legacy"}, removes it from the list, turns legacy mode on
  and returns 1; otherwise returns 0 \\
\code{wfdbx\_time\_needs\_freq(string)} & returns 1 if \code{strtim(string)} needs a
  sampling frequency, 0 otherwise \\
\code{wfdbx\_write\_errors()} & returns the number of write errors since the program
  started (output files that could not be written, flushed or closed) \\
\code{wfdbx\_note\_write\_error()} & counts a write error, for a program that writes a
  file itself \\
\code{wfdbx\_input\_path(name)} & returns the path under which a program opens a file
  that it reads itself: in normal mode, a relative name under the input directory, a
  \code{\textasciitilde} name expanded, an absolute name as given; in legacy mode, the
  name as given. \code{None}, \code{""} and \code{"-"} are returned as given; \code{None}
  for a URL, or a name that becomes one under the input directory \\
\code{wfdbx\_output\_path(name)} & returns the path under which a program creates a file
  that it writes itself: in normal mode, a relative name under the output directory (whose
  missing directories are then created), a \code{\textasciitilde} name expanded, an
  absolute name as given; in legacy mode, any name under the output directory, as a plain
  prefix. \code{None}, \code{""} and \code{"-"} are returned as given; \code{None} for a
  URL \\
\end{xltabular}

\subsection{Examples}

Reading the samples and the annotations of a record, as a C program does (the record
\code{100s} is read from the package's data directory):

\begin{lstlisting}[language=Python,style=example]
import wfdbx

previous = wfdbx.wfdbx_set_input_dir(wfdbx.sample_data_dir())

si = []                                  # receives one Siginfo per signal
nsig = wfdbx.isigopen("100s", si, 2)     # open 2 signals
print(nsig, [(s.desc, s.fmt, s.gain) for s in si])   # 2 [('MLII', 212, 200.0), ...]
print(wfdbx.sampfreq(None))              # 360.0: the frequency of the open record

v = [0] * nsig                           # receives one sample per signal
wfdbx.isigsettime(wfdbx.strtim("0:30"))  # skip to 30 seconds
wfdbx.getvec(v)
print(v, [wfdbx.aduphys(k, v[k]) for k in range(nsig)])

frames, stat = wfdbx.getframes(360)      # the next second, as a NumPy array
print(frames.shape, stat)                # (360, 2) 2

a = wfdbx.Annotation()
wfdbx.annopen("100s", [wfdbx.Anninfo("atr", wfdbx.WFDB_READ)])
n = 0
while wfdbx.getann(0, a) == 0:           # 0: an annotation was read; -1: the end
    if n < 2:
        print(wfdbx.mstimstr(a.time), wfdbx.annstr(a.anntyp), a.aux_string())
    n += 1
print(n, "annotations")                  # 75 annotations

wfdbx.wfdbquiet()                        # no messages on the standard error
print(wfdbx.isigopen("nosuchrecord", None, 0), wfdbx.wfdberror().strip())
wfdbx.wfdbverbose()

wfdbx.wfdbquit()                         # close everything
wfdbx.wfdbx_set_input_dir(previous)      # restore the input directory
\end{lstlisting}

Writing a record and an annotation file, sample by sample and annotation by annotation:

\begin{lstlisting}[language=Python,style=example]
import wfdbx

out = [wfdbx.Siginfo(fname="ramp.dat", desc="ramp", units="mV", gain=200.0, fmt=16,
                     adcres=12)]
print(wfdbx.osigfopen(out))         # 1: one output signal
wfdbx.setsampfreq(250.0)
for k in range(500):
    wfdbx.putvec([k])
print(wfdbx.newheader("ramp"))      # 0: ramp.hea is written
wfdbx.annopen("ramp", [wfdbx.Anninfo("qrs", wfdbx.WFDB_WRITE)])
for t in (100, 350):
    wfdbx.putann(0, wfdbx.Annotation(time=t, anntyp=wfdbx.NORMAL))
wfdbx.wfdbquit()                    # completes ramp.dat and ramp.qrs

print(open("ramp.hea").read().splitlines()[0])        # ramp 1 250 500
print(wfdbx.read_annotations("ramp", "qrs").samples)  # [100 350]
\end{lstlisting}

% ==========================================================================================
\section{The command-line tools in Python}
\label{sec:python-tools}

The 33 command-line tools (\code{rdsamp}, \code{rdann}, \code{wrsamp}, \code{xform},
\ldots) are run through the command \cmd{wfdbx}, installed with the package, followed by
the name of the tool and its arguments. The tools, their options, their output and their
exit status are those of the C tools, and are documented in \cref{chap:tools}.

\begin{lstlisting}[style=shell]
export WFDBX_INPUT_DIR=$(python -c "import wfdbx; print(wfdbx.sample_data_dir())")
wfdbx rdsamp -r 100s -t 10 -p     # the first 10 seconds of 100s, in physical units
wfdbx rdann -r 100s -a atr        # its annotations
wfdbx rdsamp --legacy -r 100s     # legacy mode: --legacy as the first argument
\end{lstlisting}

The command \cmd{wfdbx} also takes these options, instead of a tool:

\begin{table}[htbp]
\centering
\caption{The options of the command \code{wfdbx}}
\label{tab:python-wfdbx-options}
\begin{tabularx}{\linewidth}{@{}lL@{}}
\toprule
\tabhead Command & \tabhead Effect \\
\midrule
\cmd{wfdbx TOOL [ARGUMENTS ...]} & runs the tool \code{TOOL}, and exits with its
  status \\
\cmd{wfdbx -{}-list} & lists the names of the tools, one per line \\
\cmd{wfdbx -{}-examples [DIR]} & copies the example notebooks to the directory
  \code{DIR} (default \file{wfdbx-examples}), and prints their paths \\
\cmd{wfdbx -{}-version} & prints the version, \code{WFDB-x 1.2.2} \\
\cmd{wfdbx -{}-help}, \cmd{wfdbx -h} & prints the usage \\
\bottomrule
\end{tabularx}
\end{table}

Without arguments, \cmd{wfdbx} prints the usage and exits with status 1; an unknown tool
is an error (status 1). When an output file or the standard output cannot be written (a
full disk, for example), a tool reports it and exits with status 1.

\paragraph{Without the \cmd{wfdbx} command.} The package has no module
\code{wfdbx.\_\_main\_\_}, so \cmd{python -m wfdbx} is not supported. The command is the
module \code{wfdbx.cli}, and each tool is a module of \code{wfdbx.apps}, so that, with
\file{python/src} on \code{PYTHONPATH} for example:

\begin{lstlisting}[style=shell]
python -m wfdbx.cli rdsamp -r 100s -t 10     # as wfdbx rdsamp -r 100s -t 10
python -m wfdbx.apps.rdsamp -r 100s -t 10    # the tool's module itself
\end{lstlisting}

From Python, \code{wfdbx.cli.main(argv)} runs the command with the argument list
\code{argv} (without the program name), and the function \code{main(argv)} of a tool's
module runs the tool with \code{argv}, the complete argument list, program name first, as
in C. Both return the exit status; the tool writes to the standard output and error.

\begin{lstlisting}[language=Python,style=example]
import wfdbx.cli
from wfdbx.apps import rdsamp

status = wfdbx.cli.main(["rdsamp", "-r", "100s", "-t", "s2"])   # the first 2 samples
print(status)                                                   # 0
status = rdsamp.main(["rdsamp", "-r", "100s", "-f", "s10", "-t", "s11", "-p"])
\end{lstlisting}

\paragraph{Environment variables.} The tools, like the rest of the package, use the
environment variables of \cref{tab:python-environment}, also for the files that they read
or create themselves (\cref{sec:tool-files}): the \opt{-i} file of \cmd{wrsamp},
\cmd{edf2mit}, \cmd{wav2mit}, \cmd{a2m} and \cmd{ad2m}, the input files of
\cmd{ahaecg2mit} and the \opt{-S} script of \cmd{xform} are read from the input directory,
and the WAV file of \cmd{mit2wav}, the EDF file of \cmd{mit2edf}, the master header of
\cmd{snip -m} and the \file{.des}, \file{.ecg}, \file{.key} and \file{.txt} files of
\cmd{parsescp} are created in the output directory. Such a file cannot be read from a URL.
In legacy mode, the files that a tool reads itself are opened as named, and the output
directory, if set, is a plain prefix of every output name.

\begin{table}[htbp]
\centering
\caption{The environment variables used by the Python package}
\label{tab:python-environment}
\begin{tabularx}{\linewidth}{@{}lL@{}}
\toprule
\tabhead Variable & \tabhead Meaning \\
\midrule
\code{WFDBX\_INPUT\_DIR} & the input directory (a directory, a \code{\textasciitilde}
  path or a URL), when \code{wfdbx\_set\_input\_dir} does not set one \\
\code{WFDBX\_OUTPUT\_DIR} & the output directory, when \code{wfdbx\_set\_output\_dir} does
  not set one \\
\code{WFDBX\_NETWORK} & network access is on if it is a nonzero number, when
  \code{wfdbx\_set\_network} does not set it \\
\code{WFDBANNSORT} & \code{0} to keep annotation files in the order written (no sorting
  when they are closed) \\
\code{WFDBGVMODE} & the default getvec mode (\code{setgvmode} with a negative mode) \\
\code{WFDBCAL} & the calibration file read by \code{calopen(None)} (default
  \file{wfdbcal}) \\
\code{WFDB} & the WFDB path, used in legacy mode only (outside legacy mode, it is
  ignored, with a message) \\
\code{WFDB\_FLAC\_CHECK\_MD5} & if set, the MD5 signature of FLAC signal files is
  verified when they are read \\
\bottomrule
\end{tabularx}
\end{table}

% ==========================================================================================
\section{Testing}
\label{sec:python-testing}

The tests of the Python package are in \file{python/tests}. From the root of the
repository:

\begin{lstlisting}[style=shell]
pip install -e "./python[test]"       # the package and pytest, nbformat, nbclient, ipykernel
pytest python/tests
\end{lstlisting}

This runs the unit tests of the library, the high-level functions, the metadata functions
(with the cases shared by the seven languages, in \file{tests/metadata}) and
\code{check\_record}; it runs the example notebooks (when nbformat, nbclient and ipykernel
are installed); and it runs the cross-language conformance suite, which compares the
output of the tools with that of the original C programs, case by case (1155 cases, a few
minutes). The tests run against \file{python/src} when they are run from the repository,
in a temporary copy of the test records of \file{tests/data}. The tests that write to
\file{/dev/full} (write errors) are skipped where it does not exist, on macOS and Windows.

The environment variable \code{WFDBX\_CONFORMANCE} selects the conformance cases that
\cmd{pytest} runs: \code{0} skips the suite, and a list of tools, such as
\code{WFDBX\_CONFORMANCE=rdsamp,rdann}, runs the cases of those tools only.

\begin{lstlisting}[style=shell]
WFDBX_CONFORMANCE=0 pytest python/tests            # the unit tests only
WFDBX_CONFORMANCE=rdsamp,rdann pytest python/tests/test_conformance.py
\end{lstlisting}

The conformance suite can also be run on its own, with more options (the number of
parallel jobs, the legacy suite of 886 cases, which runs each tool with \opt{-{}-legacy},
verbose output, selected cases):

\begin{lstlisting}[style=shell]
PYTHONPATH=$PWD/python/src python3 tests/conformance/harness.py check --impl python -j4
PYTHONPATH=$PWD/python/src python3 tests/conformance/harness.py check --impl python --legacy
\end{lstlisting}

% The C++ chapter of the WFDB-x user's manual. It follows the sections of the Python
% chapter (chapters/python.tex), in the same order. Each \apientry of the section "The
% high-level API" is a function declared at namespace scope in
% cpp/include/wfdbx/record.hpp; its signature listing is copied from that declaration,
% and its params name every parameter of the declaration, in order. The examples of style
% "example" were compiled and run (each as the body of main(), after the lines given in
% the conventions of the section, in a copy of tests/data); check_manual.py does not run
% them.
\chapter{C++}
\label{chap:cpp}
\chapterlanguage{cpp}

The C++ implementation of WFDB-x is the library \code{libwfdbx}, with its headers in
\file{cpp/include/wfdbx} and everything in namespace \code{wfdbx}, and the 33
command-line tools, built as separate programs. It is a complete C++17 implementation of
the WFDB Software Package: it needs only the C++17 standard library, and never calls the
C library.

It offers WFDB-x in three ways:
\begin{itemize}
\item \textbf{The high-level API} (\cref{sec:cpp-api}): functions, declared in
  \file{<wfdbx/record.hpp>}, that read and write whole records, headers, annotation files
  and header metadata as \code{std::vector}s and structures, in one call, with options
  given as fields of a structure, and that throw \code{wfdbx::Error} on errors. They read
  signal files in large blocks, and are the fastest way to read a record. They are an
  addition of WFDB-x: the C package has no equivalent. They are the recommended interface
  for new code and analysis.
\item \textbf{The WFDB library functions} under their C names (\code{isigopen},
  \code{getvec}, \code{getann}, \code{putann}, \ldots), declared in
  \file{<wfdbx/wfdb.hpp>}, with the behavior of the C library (\cref{sec:cpp-library}).
  They are documented once, for all the languages, in \cref{chap:library}.
\item \textbf{The command-line tools} (\code{rdsamp}, \code{rdann}, \code{wrsamp},
  \code{xform} and 29 others), with the options, output and exit status of the C tools,
  one program per tool, run by its name (\cref{sec:cpp-tools}). They are documented in
  \cref{chap:tools}.
\end{itemize}
The high-level functions return exactly the samples and annotations that the library
functions (and the C library) give.

% ==========================================================================================
\section{Installation and requirements}
\label{sec:cpp-install}

\subsection{Requirements}

WFDB-x for C++ builds with CMake 3.14 or later and a C++17 compiler whose standard
library has \code{std::filesystem}: GCC 9, Clang 9 or MSVC 2019 and later (macOS 10.15 or
later), on Linux, macOS and Windows. It has no other dependency: no WFDB library and no
libFLAC (the FLAC formats 508, 516 and 524 are read and written by WFDB-x's own codec).
Network access (\cref{sec:network}) runs the program \cmd{curl}, which must then be
installed. The unit tests of \code{check\_record} also run \cmd{python3}
(\cref{sec:cpp-testing}).

\subsection{Building and installing}

From the root of a copy of the repository, these commands build the library, the tools
and the unit tests (a Release build), and install them:

\begin{lstlisting}[style=shell]
cmake -S cpp -B cpp/build                        # a Release build
cmake --build cpp/build -j8
cmake --install cpp/build --prefix /usr/local    # optional
\end{lstlisting}

The library is \file{cpp/build/libwfdbx.a}, and the tools are in \file{cpp/build/bin/}.
The install step copies the library, the headers (\file{include/wfdbx/}), the tools and a
CMake package file (\file{lib/cmake/wfdbx/wfdbxConfig.cmake}, with the target
\code{wfdbx::wfdbx}). These options of the build change what is built:

\begin{table}[htbp]
\centering
\caption{The options of the CMake build}
\label{tab:cpp-cmake-options}
\begin{tabularx}{\linewidth}{@{}llL@{}}
\toprule
\tabhead Option & \tabhead Default & \tabhead Effect \\
\midrule
\code{WFDBX\_BUILD\_TOOLS} & \code{ON} & build the command-line tools \\
\code{WFDBX\_BUILD\_TESTS} & \code{ON} & build the unit tests \\
\code{BUILD\_SHARED\_LIBS} & \code{OFF} & build a shared library \\
\code{CMAKE\_BUILD\_TYPE} & \code{Release} & the build type, when none is given \\
\bottomrule
\end{tabularx}
\end{table}

The version of the library is that of the CMake project, \code{1.2.0}. The tools, like
the C tools, have no \opt{-{}-version} option.

\subsection{Using the library from CMake}

An installed library is found with \code{find\_package}, which defines the imported
target \code{wfdbx::wfdbx} (the library and its include directory):

\begin{lstlisting}
find_package(wfdbx REQUIRED)
target_link_libraries(myprogram PRIVATE wfdbx::wfdbx)
\end{lstlisting}

A copy of the directory \file{cpp} (for example in a subdirectory of a project) can also
be built as part of a project with \code{add\_subdirectory}, which defines the target
\code{wfdbx} and its alias \code{wfdbx::wfdbx}. The options of
\cref{tab:cpp-cmake-options} then select what is built; for the library alone:

\begin{lstlisting}
set(WFDBX_BUILD_TOOLS OFF CACHE BOOL "Build the command-line tools" FORCE)
set(WFDBX_BUILD_TESTS OFF CACHE BOOL "Build the unit tests" FORCE)
add_subdirectory(path/to/wfdb-x/cpp ${CMAKE_CURRENT_BINARY_DIR}/wfdbx EXCLUDE_FROM_ALL)
target_link_libraries(myprogram PRIVATE wfdbx::wfdbx)
\end{lstlisting}

Without CMake, a program is compiled against the headers and linked with the library,
installed (\cmd{-lwfdbx}) or in the build directory:

\begin{lstlisting}[style=shell]
c++ -std=c++17 myprogram.cpp -lwfdbx                                  # installed
c++ -std=c++17 -Icpp/include myprogram.cpp cpp/build/libwfdbx.a -o myprogram
\end{lstlisting}

\subsection{The sample record}
\label{sec:cpp-resources}

The repository includes the sample record \code{100s}, the sample record of the WFDB
Software Package: the first minute of record 100 of the MIT-BIH Arrhythmia Database, two
ECG signals (\code{MLII} and \code{V5}) sampled at 360\,Hz, with the reference beat
annotations (annotator \code{atr}). Its files (\file{100s.hea}, \file{100s.dat},
\file{100s.atr}) are in the directory \file{tests/data} of a copy of the repository, with
the other test records of WFDB-x. The C++ library does not install them: programs name
that directory as their input directory (\code{wfdbx\_set\_input\_dir("tests/data")}).
Work that uses the MIT-BIH Arrhythmia Database cites it as described on PhysioNet.

% ==========================================================================================
\section{Quick start}
\label{sec:cpp-quickstart}

This program reads the header, ten seconds of the signals and the annotations of the
sample record \code{100s}, from the directory \file{tests/data} of the repository (the
input directory of each call), and
computes the mean heart rate. It then writes the first ten seconds as a new record,
\code{first10}, in the current directory, with a metadata entry in its header.

\begin{lstlisting}[language=C++,style=example]
#include <cstdio>
#include <vector>
#include <wfdbx/ecgmap.hpp>
#include <wfdbx/record.hpp>

int main() {
    wfdbx::InputOptions sample;
    sample.input_dir = "tests/data";             // the sample record 100s

    // The header only: no samples are read.
    wfdbx::Header hdr = wfdbx::read_header("100s", false, sample);
    std::printf("%d %g %lld\n", hdr.nsig, hdr.fs, hdr.nsamples);   // 2 360 21600
    std::printf("%s %s\n", hdr.siginfo[0].desc->c_str(),
                hdr.siginfo[1].desc->c_str());                     // MLII V5

    // Seconds 10 to 20 of both signals, frame by frame, in ADC units.
    wfdbx::ReadOptions opt;
    opt.input_dir = "tests/data";
    opt.from = "0:10";
    opt.to = "0:20";
    wfdbx::Record rec = wfdbx::read_record("100s", opt);
    std::printf("%zu %d\n", rec.nsamples(), rec.nsig);             // 3600 2
    std::printf("%d %g\n", rec.at(0, 0), rec.physical(0, 0));      // 946 -0.39

    // The beat annotations, and the mean heart rate.
    wfdbx::Annotations ann = wfdbx::read_annotations("100s", "atr", {}, {},
                                                     std::nullopt, false, sample);
    std::vector<long long> beats;
    for (size_t i = 0; i < ann.size(); i++)
        if (wfdbx::isqrs(ann.codes[i])) beats.push_back(ann.samples[i]);
    double rr = (beats.back() - beats.front()) / ann.fs / (beats.size() - 1);
    std::printf("%zu beats, %.1f beats per minute\n", beats.size(), 60 / rr);
    // 74 beats, 73.9 beats per minute

    // Write the first 10 seconds as the record "first10" (first10.hea and first10.dat,
    // in the current directory), with the signal specifications of 100s.
    opt.from = 0;
    opt.to = "0:10";
    wfdbx::Record x = wfdbx::read_record("100s", opt);
    wfdbx::WriteOptions w;
    w.gains = {200};
    w.baselines = {1024};
    w.adczeros = {1024};
    w.adcres = {11};
    w.units = {"mV"};
    w.names = x.names();
    w.metadata = {{"Source", "MIT-BIH record 100"}};
    wfdbx::write_record("first10", x.samples, x.nsig, x.fs, w);

    // Header metadata: "# Key: value" lines of the header.
    wfdbx::set_metadata("first10", "Age", "69");
    for (const auto &e : wfdbx::read_metadata("first10"))
        std::printf("%s = %s\n", e.first.c_str(), e.second.c_str());
    // Source = MIT-BIH record 100
    // Age = 69
    std::printf("%s\n", wfdbx::check_record("first10").to_string().c_str());
    // no problems found
}
\end{lstlisting}

Compiled and run from the root of the repository:

\begin{lstlisting}[style=shell]
c++ -std=c++17 -Icpp/include quick.cpp cpp/build/libwfdbx.a -o quick
./quick
\end{lstlisting}

The header \file{100s.hea} gives no units: the program writes \code{"mV"} explicitly,
since WFDB-x never assumes units, a gain, an ADC resolution or a sampling frequency that
a record does not give (\cref{chap:behavior}).

% ==========================================================================================
\section{The high-level API}
\label{sec:cpp-api}

The high-level API is made of the functions of this section, declared in
\file{<wfdbx/record.hpp>} in namespace \code{wfdbx}, of the structures of their options
(\cref{sec:cpp-options}) and of their results (\cref{sec:cpp-results}), and of the
exceptions \code{Error}, \code{ArgumentError} and \code{FileError}
(\cref{sec:cpp-warnings}). The comments of the header document each of them.

\subsection{Conventions}
\label{sec:cpp-conventions}

\paragraph{Record names and directories.} A record is named as in the other parts of
WFDB-x (\cref{chap:records,chap:files}), without the \file{.hea} extension (the readers
also accept a trailing \code{".hea"}; \code{write\_record} does not). A record name is
one of: a full path, absolute (\code{"/data/mitdb/100"}) or in the home directory
(\code{"\textasciitilde/Documents/100s"}), used as given; a relative path
(\code{"100s"}, \code{"data/100s"}), read from the input directory if one is set,
otherwise from the current directory, and written to the output directory if one is set,
otherwise to the current directory (missing directories are created); or a URL
(\code{"https://physionet.org/files/mitdb/1.0.0/100"}), read only if network access is
on, and which cannot be written. The input and output directories are set with
\code{wfdbx\_set\_input\_dir} and \code{wfdbx\_set\_output\_dir}
(\cref{sec:cpp-settings}) or with the environment variables \code{WFDBX\_INPUT\_DIR} and
\code{WFDBX\_OUTPUT\_DIR}. A writer given a name with a path creates the files in that
directory, and the header names the record and its signal file without the path
(\code{rec}, \code{rec.dat}), so that the record opens under any path that leads to it;
in legacy mode, as in the original, the header gives the name with its path. The readers
take the fields \code{input\_dir} and \code{network} of an \code{InputOptions} (or of a
structure derived from it), and the writers a \code{dir}: they apply to that one call,
and the previous settings are restored when it returns, also when it throws. The
\code{input\_dir} option is optional, and applies to relative names only, not to full
paths or URLs. With \code{dir}, the record name must be a plain name (\code{"rec"}): the
files are created in \code{dir}, and the record opens from that directory under its plain
name. The signal files that a header names are read from the header's directory.

\paragraph{Options.} The options of a function are the fields of a structure
(\code{ReadOptions}, \code{InputOptions}, \code{WriteOptions}, \code{MetadataOptions}),
given as its last argument, with a default value \code{\{\}}: a field that is not set
keeps its default. An option that may be absent is a \code{std::optional}.

\paragraph{Times.} The fields \code{from} and \code{to} of \code{ReadOptions}, and the
arguments \code{from} and \code{to} of \code{read\_annotations}, are \code{TimeSpec}
values: a sample number or a WFDB time string, as \code{strtim} reads it (\code{"0:10"},
10 seconds; \code{"1:02:03.5"}; \code{"s3600"}, sample 3600; \code{"e"}, the end of the
record; or a time of day in brackets, \code{"[14:00:00]"}). A \code{TimeSpec} converts
from integers and strings, so \code{opt.from = 3600;} and \code{opt.from = "0:10";} both
work; a default-constructed \code{TimeSpec} is empty (no time is given). A time in seconds
needs the sampling frequency; when the record does not give it, such a time throws
\code{Error}.

\paragraph{Samples.} Samples are \code{WFDB\_Sample} values (\code{int}), in ADC units,
stored frame by frame in one \code{std::vector}: \code{samples[i * nsig + s]} is sample
\code{i} of signal \code{s}. Invalid samples are \code{WFDB\_INVALID\_SAMPLE}
($-32768$), and NaN in physical units.

\paragraph{Per-signal and per-annotation values.} The per-signal fields of
\code{WriteOptions} (\code{gains}, \code{units}, \ldots) hold one value for all the
signals, or one value per signal; the per-annotation fields of \code{Annotations}
(\code{symbols}, \code{chans}, \ldots) are empty or hold one value per annotation. Any
other length throws \code{Error} before anything is written.

\paragraph{Signal parameters.} No sampling frequency, gain, ADC resolution or units is
ever assumed. Where one is needed and the record does not give it, it must be given;
otherwise the function throws \code{Error}. A record whose header gives no sampling
frequency has \code{fs = 0}.

\paragraph{Text.} Strings are \code{std::string}, handled as bytes: every byte sequence
round-trips unchanged, and UTF-8 text is stored as its bytes.

\paragraph{Legacy mode.} With \code{legacy} set to \code{true}, a function behaves as in
the original WFDB package, with its assumed signal parameters and its search for files
(\cref{chap:behavior}); the library's legacy mode is on for the duration of the call
only. With \code{false} (the default), the current setting of \code{wfdbx\_set\_legacy}
applies, which is off unless it was turned on.

\paragraph{State.} The functions that read and write records and annotation files call
\code{wfdbquit()} before and after their work, also when they throw, which closes
anything opened with the WFDB library functions. The library's settings (legacy mode,
getvec mode, time resolution of output annotators, input and output directories) are the
same after a call as before it. Opening an EDF record sets the time resolution of output
annotators (\code{getafreq()}) to its sampling frequency when none is set, and adds
\code{WFDB\_HIGHRES} to the getvec mode (\code{getgvmode()}), as in C, and
\code{wfdbquit()} resets neither: the functions that open a record restore both when they
return, also when they throw, so that a later \code{write\_annotations} does not record
that frequency; the C-style functions keep the behavior of C. The metadata functions do
not call \code{wfdbquit()}: they open only the header, and leave the library's open files
as they are. The functions are not thread-safe. Results are returned by value.

\paragraph{The examples} of this section are the body of a function \code{main}, compiled
after these lines, and run in a copy of the test records of \file{tests/data} (among them
\code{100s}), which is the current directory and holds the records that the examples
write:
\begin{lstlisting}[language=C++]
#include <algorithm>
#include <cmath>
#include <cstdio>
#include <filesystem>
#include <random>
#include <vector>
#include <wfdbx/ecgcodes.hpp>
#include <wfdbx/record.hpp>
using namespace wfdbx;
\end{lstlisting}

% ------------------------------------------------------------------------------------------
\subsection{Reading records}
\label{sec:cpp-reading}

\apientry{read\_record}{Read the signals of a record}
\begin{lstlisting}[language=C++,style=signature]
Record read_record(const std::string &record, const ReadOptions &opt = {});
\end{lstlisting}
Reads the signals of a record from \code{opt.from} up to \code{opt.to} (exclusive), in
ADC units, frame by frame, together with the description of the record. The samples are
those that \code{getvec} returns (\cref{chap:library}), read in large blocks with
\code{getvecs}. If the record ends before \code{opt.to}, the samples up to its end are
returned. When the read reaches the end of the record, the samples are compared with the
checksums of the header (\cref{sec:cpp-warnings}). The getvec mode (\code{getgvmode()})
and the time resolution of output annotators (\code{getafreq()}) are the same after the
call as before it (opening an EDF record sets both).
\begin{params}
\item[record] Record name, without the \file{.hea} extension (a trailing \code{".hea"} is
  also accepted): a full path, absolute (\code{"/data/mitdb/100"}) or in the home
  directory (\code{"\textasciitilde/Documents/100s"}), used as given; a relative path
  (\code{"100s"}, \code{"data/100s"}), read from the input directory if one is set,
  otherwise from the current directory; or a URL
  (\code{"https://physionet.org/files/mitdb/1.0.0/100"}), read only if network access is
  on.
\item[opt] The options (\cref{tab:cpp-readoptions}): the signals to read
  (\code{channels}), the interval (\code{from}, \code{to}), \code{highres}, \code{pad},
  \code{checksum\_warning}, \code{legacy}, and the input directory, under which relative
  record names are read, and network access of the call (\code{input\_dir},
  \code{network}). Default \code{\{\}}: every signal of the
  whole record, with the current settings.
\end{params}
\returns{A \code{Record} (\cref{tab:cpp-record}), with the samples in ADC units.}
\raises{\code{Error} if the record cannot be read (``can't read record REC'') or has no
signals; if a requested signal does not exist (``record REC has no signal C''); if
\code{opt.from} is negative (``read\_record: from must not be negative'') or cannot be
reached (``can't seek to sample N of record REC''); if a time string needs the sampling
frequency and the header gives none (``sampling frequency of record REC is not
specified''); or if a signal file is shorter than its header says, or cannot be read or
decoded (``can't read signal data of record REC (getvec status S)''). Data that do not
match the checksums of the header are returned, with a warning (see
\code{checksum\_warning}).}
\begin{lstlisting}[language=C++,style=example]
// Seconds 10 to 20 of lead MLII, in physical units.
ReadOptions opt;
opt.channels = {"MLII"};
opt.from = "0:10";
opt.to = "0:20";
Record rec = read_record("100s", opt);
std::vector<double> mlii = rec.physical_signal(0);
std::printf("%zu %g %lld %s\n", mlii.size(), rec.fs, rec.sampfrom,
            rec.names()[0].c_str());                   // 3600 360 3600 MLII

// The first 3 samples of both signals, in ADC units, V5 first.
ReadOptions first;
first.channels = {"1", "0"};
first.to = 3;
rec = read_record("100s", first);
for (size_t i = 0; i < rec.nsamples(); i++)
    std::printf("%d %d\n", rec.at(i, 0), rec.at(i, 1));   // 1011 995, 3 times
std::printf("%g %d\n", rec.siginfo[0].gain, rec.siginfo[0].baseline);   // 200 1024

// The whole record: the checksums are verified.
rec = read_record("100s");
std::printf("%zu %d\n", rec.nsamples(), rec.checksum_ok);   // 21600 1
\end{lstlisting}
A multi-frequency record (here the test record \code{mfreq}, whose signals have 1, 2
and 4 samples per frame at a frame frequency of 100\,Hz) is read at its frame frequency,
or at its highest sampling frequency:
\begin{lstlisting}[language=C++,style=example]
Record rec = read_record("mfreq");
std::printf("%g %zu\n", rec.fs, rec.nsamples());       // 100 600
ReadOptions opt;
opt.highres = true;
rec = read_record("mfreq", opt);
std::printf("%g %zu\n", rec.fs, rec.nsamples());       // 400 2400
for (const WFDB_Siginfo &s : rec.siginfo) std::printf("%d ", s.spf);   // 1 2 4
std::printf("\n");
\end{lstlisting}

\apientry{read\_header}{Read the header of a record}
\begin{lstlisting}[language=C++,style=signature]
Header read_header(const std::string &record, bool legacy = false,
                   const InputOptions &input = {});
\end{lstlisting}
Reads the header of a record, without opening its signal files. The result describes the
record and each of its signals, and holds the info strings and the metadata entries of the
header (from the \file{.hea} file only, as \code{read\_metadata} reads them, with its
warning for a key given more than once). For an EDF record, which has no \file{.hea}
file, and for a header read from the standard input (\code{"-"}), the metadata is empty,
without an error. For a multi-segment record, the master header is read, and
\code{nsamples} is the length of the whole record. The time resolution of output
annotators (\code{getafreq()}) and the getvec mode (\code{getgvmode()}) are the same after
the call as before it (opening an EDF record sets both).
\begin{params}
\item[record] Record name, without the \file{.hea} extension (a trailing \code{".hea"} is
  also accepted): a full path, absolute (\code{"/data/mitdb/100"}) or in the home
  directory (\code{"\textasciitilde/Documents/100s"}), used as given; a relative path
  (\code{"100s"}, \code{"data/100s"}), read from the input directory if one is set,
  otherwise from the current directory; or a URL
  (\code{"https://physionet.org/files/mitdb/1.0.0/100"}), read only if network access is
  on.
\item[legacy] \code{true} to read the header as the original WFDB package does. Default
  \code{false}.
\item[input] The input directory for this call only (optional), under which relative
  record names are read (it does not apply to full paths or URLs), and network access
  (\cref{tab:cpp-inputoptions}). Default \code{\{\}}: the current settings.
\end{params}
\returns{A \code{Header} (\cref{tab:cpp-header}).}
\raises{\code{Error} if the header cannot be read (``can't read header for record
REC'').}
\begin{lstlisting}[language=C++,style=example]
Header hdr = read_header("100s");
std::printf("%s %g %d %lld\n", hdr.name.c_str(), hdr.fs, hdr.nsig, hdr.nsamples);
// 100s 360 2 21600
for (const WFDB_Siginfo &s : hdr.siginfo)
    std::printf("%s %s %d %g %d %s\n", s.desc->c_str(), s.fname->c_str(), s.fmt,
                s.gain, s.adcres, s.units ? s.units->c_str() : "(none)");
// MLII 100s.dat 212 200 11 (none)
// V5 100s.dat 212 200 11 (none)
std::printf("[%s]\n", hdr.comments[0].c_str());   // [ 69 M 1085 1629 x1]
std::printf("%d\n", hdr.base_time.has_value());   // 0: the header gives no start time
\end{lstlisting}

\apientry{read\_annotations}{Read an annotation file}
\begin{lstlisting}[language=C++,style=signature]
Annotations read_annotations(const std::string &record,
                             const std::string &annotator,
                             const TimeSpec &from = {},
                             const TimeSpec &to = {},
                             std::optional<WFDB_Frequency> fs = std::nullopt,
                             bool legacy = false,
                             const InputOptions &input = {});
\end{lstlisting}
Reads the annotations of an annotation file, in file order, as parallel vectors. The
sampling frequency, needed to convert times in seconds, is that of the record's header,
or, for a record without a header or whose header gives none, the argument \code{fs}, as
the option \opt{-F} of \code{rdann}. The time resolution of output annotators
(\code{getafreq()}) and the getvec mode (\code{getgvmode()}) are the same after the call
as before it (opening an EDF record sets both).
\begin{params}
\item[record] Record name, without the \file{.hea} extension (a trailing \code{".hea"} is
  also accepted): a full path, absolute (\code{"/data/mitdb/100"}) or in the home
  directory (\code{"\textasciitilde/Documents/100s"}), used as given; a relative path
  (\code{"100s"}, \code{"data/100s"}), read from the input directory if one is set,
  otherwise from the current directory; or a URL
  (\code{"https://physionet.org/files/mitdb/1.0.0/100"}), read only if network access is
  on. Its header, if there is one, gives the sampling
  frequency.
\item[annotator] Annotator name: the suffix of the annotation file. \code{"atr"} reads the
  file \file{RECORD.atr}.
\item[from] The first time to read: a sample number or a time string. Default \code{\{\}}:
  the beginning of the file.
\item[to] The time at which to stop: annotations at or after it are not read. Default
  \code{\{\}}: the end of the file.
\item[fs] The sampling frequency, in Hz ($> 0$), to use when the record has no header or
  its header gives none; it is ignored otherwise. Default \code{std::nullopt}: none.
\item[legacy] \code{true} to behave as the original WFDB package (which assumes 250\,Hz
  for a record without a frequency). Default \code{false}.
\item[input] The input directory for this call only (optional), under which relative
  record names are read (it does not apply to full paths or URLs), and network access
  (\cref{tab:cpp-inputoptions}). Default \code{\{\}}: the current settings.
\end{params}
\returns{An \code{Annotations} (\cref{tab:cpp-annotations}). Its \code{fs} is the
frequency in which the times are counted: the one written in the annotation file, if
there is one, or else the sampling frequency; it is 0 if none of them is known.}
\raises{\code{Error} if \code{fs} is given and is not $> 0$ (``read\_annotations: fs must
be > 0''); if the annotation file cannot be opened (``can't read annotator ANN for record
REC''); if a time string needs the sampling frequency and none is known (``sampling
frequency of record REC is not specified''); or if the file ends early or cannot be read
before \code{to} is reached (``can't read annotation file ANN of record REC (getann status
S)'', where S is the status that \code{getann} returned, such as $-3$ for an unexpected
end of file).}
\begin{lstlisting}[language=C++,style=example]
Annotations ann = read_annotations("100s", "atr", "0:10", "0:30");
std::printf("%zu %g\n", ann.size(), ann.fs);                   // 24 360
std::printf("%lld %s\n", ann.samples[0], ann.symbols[0].c_str());   // 3862 N

// The whole file: the first annotation is a rhythm change, with an aux string.
ann = read_annotations("100s", "atr");
std::printf("%s %d %s\n", ann.symbols[0].c_str(), ann.codes[0] == RHYTHM,
            ann.aux[0]->c_str());                              // + 1 (N
\end{lstlisting}
For a record without a header, \code{fs} gives the frequency in which times in seconds
are converted (the library reports on the standard error that it cannot open the
header):
\begin{lstlisting}[language=C++,style=example]
Annotations beats;
beats.samples = {100, 460, 820};
beats.symbols = {"N", "N", "N"};
write_annotations("beats", "qrs", beats);             // no header
Annotations ann = read_annotations("beats", "qrs", {}, "0:02", 360.0);
std::printf("%zu %g\n", ann.size(), ann.fs);          // 2 360
\end{lstlisting}

% ------------------------------------------------------------------------------------------
\subsection{Writing records}
\label{sec:cpp-writing}

\apientry{write\_record}{Write a header and a signal file}
\begin{lstlisting}[language=C++,style=signature]
void write_record(const std::string &record,
                  const std::vector<WFDB_Sample> &samples, int nsig,
                  WFDB_Frequency fs, const WriteOptions &opt = {});
\end{lstlisting}
Writes samples in ADC units as a new record: the header \file{RECORD.hea} and one signal
file, \file{RECORD.dat}, which holds all the signals (so they share one format). Existing
files of these names are overwritten. The sampling frequency, the gain and ADC resolution
of each signal, and the units of each calibrated signal must be given: there are no
defaults. Every check except those of the writing itself comes before any file is
created. The time resolution of output annotators (\code{getafreq()}) and the getvec
mode (\code{getgvmode()}) are left as they were.
\begin{params}
\item[record] Record name (not empty), without the \file{.hea} extension: a plain name
  (\code{"rec"}) or a relative path (\code{"data/rec"}), created in the output directory
  if one is set, otherwise in the current directory, missing directories included; or a
  full path, absolute or in the home directory (\code{"\textasciitilde/Documents/rec"}),
  used as given. With a path, the files are created in that directory, and the header
  names the record and its signal file without the path (\code{rec}, \code{rec.dat}), so
  that the record opens under any path that leads to it (in legacy mode, as in the
  original, the header gives the name with its path). A URL cannot be written.
\item[samples] The samples, in ADC units, frame by frame: \code{nsig} values per frame,
  or, for a multi-frequency record, the sum of \code{opt.samps\_per\_frame} (the samples
  of signal 0, then those of signal 1, and so on). \code{WFDB\_INVALID\_SAMPLE} marks
  invalid samples. The number of samples must be a multiple of the frame size.
\item[nsig] The number of signals ($> 0$).
\item[fs] The sampling frequency, in Hz: a finite number $> 0$ (for a multi-frequency
  record, the frame frequency).
\item[opt] The options (\cref{tab:cpp-writeoptions}): the format, the signal
  specifications (\code{gains}, \code{adcres} and \code{units} are required), the start
  time, the comments and metadata, the directory \code{dir} (which must exist) in which
  to create the files for this call only, with a plain record name, and \code{legacy}.
\end{params}
\returns{Nothing.}
\raises{\code{Error} if a \code{samps\_per\_frame} value is less than 1
(``write\_record: samps\_per\_frame must hold integers >= 1'', checked first); if the
record name is empty; if \code{nsig} is not $> 0$ (``write\_record: at least one signal
is required''); if a per-signal field of \code{opt} has another length than 0, 1 or
\code{nsig}; if \code{fs} is not finite or not $> 0$; if a gain is not finite
(``write\_record: gains must be finite''); if the gains, the ADC resolutions or the units
of a calibrated signal are missing (``write\_record: gains must be specified'',
``write\_record: ADC resolution(s) must be specified'', ``write\_record: units must be
specified''); if \code{samples.size()} is not a multiple of the frame size; if
\code{opt.base\_time} is not a valid time (``write\_record: incorrect base time 'T'''); if
the files cannot be created; or if they cannot be written completely, for example on a
full disk (``write error on record REC''). With \code{opt.dir}, also if the directory does
not exist (``directory D does not exist'') or the record name contains one (``record name
must not contain a directory when dir is given''). \code{ArgumentError} (a subclass of
\code{Error}) if an entry of \code{opt.metadata} is not valid (``write\_record: metadata
key 'K' is also given in comments'', for example).}
\begin{lstlisting}[language=C++,style=example]
// A ramp of 1000 samples of one signal, in millivolts, at 100 Hz: the record "ramp"
// (ramp.hea and ramp.dat).
std::vector<WFDB_Sample> x(1000);
for (int i = 0; i < 1000; i++) x[i] = i;
WriteOptions w;
w.gains = {200};
w.adcres = {12};
w.units = {"mV"};
w.names = {"ramp"};
w.base_time = "14:00:03 21/06/2024";
w.comments = {"a ramp"};
w.metadata = {{"Generator", "loop"}};
write_record("ramp", x, 1, 100, w);
Record rec = read_record("ramp");
std::printf("%zu %s %g\n", rec.nsamples(), rec.base_time->c_str(),
            rec.physical(10, 0));                      // 1000 [14:00:03 21/06/2024] 0.05
for (const std::string &c : read_header("ramp").comments)
    std::printf("[%s]\n", c.c_str());
// [a ramp]
// [ Generator: loop]
\end{lstlisting}
A multi-frequency record: signal \code{fast} has 2 samples per frame, signal
\code{slow} one; each frame holds the two samples of \code{fast} and then the sample of
\code{slow}:
\begin{lstlisting}[language=C++,style=example]
std::vector<WFDB_Sample> frames;
for (int i = 0; i < 100; i++) {          // 1 second at 100 frames per second
    frames.push_back(2 * i);             // fast: 200 samples per second
    frames.push_back(2 * i + 1);
    frames.push_back(500);               // slow: 100 samples per second
}
WriteOptions w;
w.gains = {200};
w.adcres = {12};
w.units = {"mV"};
w.names = {"fast", "slow"};
w.samps_per_frame = {2, 1};
write_record("mf", frames, 2, 100, w);
ReadOptions opt;
opt.highres = true;
Record rec = read_record("mf", opt);
std::printf("%g %d %d %d %d\n", rec.fs, rec.at(0, 0), rec.at(0, 1), rec.at(1, 0),
            rec.at(1, 1));                         // 200 0 500 1 500
\end{lstlisting}

\apientry{write\_record\_physical}{Write a record from samples in physical units}
\begin{lstlisting}[language=C++,style=signature]
void write_record_physical(const std::string &record,
                           const std::vector<double> &samples, int nsig,
                           WFDB_Frequency fs, const WriteOptions &opt = {});
\end{lstlisting}
As \code{write\_record}, with the samples in physical units. Each sample $v$ is stored as
$\mathrm{round}(v \cdot \mathit{gain}) + \mathit{baseline}$ (halves rounded away from
zero), and NaN is stored as an invalid sample. Every signal must have a gain other than
0.
\begin{params}
\item[record] Record name, as for \code{write\_record}.
\item[samples] The samples, in physical units, frame by frame (as for
  \code{write\_record}).
\item[nsig] The number of signals ($> 0$).
\item[fs] The sampling frequency, in Hz: a finite number $> 0$.
\item[opt] The options, as for \code{write\_record} (\cref{tab:cpp-writeoptions}).
\end{params}
\returns{Nothing.}
\raises{As \code{write\_record}; also \code{Error} if a gain is 0 (``write\_record: signal
N is uncalibrated''), or if a sample is infinite or its ADC value is outside the range of
a 32-bit integer (``write\_record: sample I is out of range''). Nothing is written then.}
\begin{lstlisting}[language=C++,style=example]
// Two signals of random numbers, in millivolts, at 360 Hz: the record "rand".
std::mt19937 gen(1);
std::normal_distribution<double> randn;
std::vector<double> x(3600 * 2);                    // 3600 frames of 2 signals
for (double &v : x) v = randn(gen);
WriteOptions w;
w.fmt = 16;
w.gains = {1000};
w.adcres = {16};
w.units = {"mV"};
w.names = {"A", "B"};
write_record_physical("rand", x, 2, 360, w);
Record rec = read_record("rand");
double err = 0;
for (size_t i = 0; i < x.size(); i++)
    err = std::max(err, std::abs(rec.physical(i / 2, i % 2) - x[i]));
std::printf("%zu %d\n", rec.nsamples(), err <= 0.0005);   // 3600 1: rounded to 1/1000 mV
\end{lstlisting}

\apientry{write\_annotations}{Write an annotation file}
\begin{lstlisting}[language=C++,style=signature]
void write_annotations(const std::string &record, const std::string &annotator,
                       const Annotations &ann, WFDB_Frequency fs = 0,
                       const std::string &dir = {},
                       std::optional<bool> sort = std::nullopt);
\end{lstlisting}
Writes an annotation file, \file{RECORD.ANNOTATOR}, in the MIT format; an existing file of
that name is overwritten. The record need not have a header. The type of each annotation
is its \code{symbols} entry (\code{"N"}, \code{"V"}, \ldots; an unknown mnemonic gives
\code{NOTQRS}), or its \code{codes} entry (\cref{chap:formats}). Both may be given (as
\code{read\_annotations} returns them) if they agree: the code of each non-empty symbol
(\code{strann}) is the code given for that annotation; an empty symbol leaves the code.
\code{subtypes}, \code{chans} and \code{nums} are stored in 8 bits, as by \code{putann}
(300 is stored as 44, as in C), and \code{aux} texts are cut to 255 bytes. If the
annotations are not in time order, the file is sorted when it is written. The arguments
are checked before the file is created. The time resolution of output annotators
(\code{getafreq()}) and the getvec mode (\code{getgvmode()}) are left as they were:
\code{fs} applies to this file only.
\begin{params}
\item[record] Record name, without the \file{.hea} extension (a trailing \code{".hea"} is
  also accepted): a plain name or a relative path (\code{"data/100s"}) is created in the
  output directory if one is set, otherwise in the current directory; a full path,
  absolute or in the home directory (\code{"\textasciitilde/Documents/100s"}), is used as
  given. The file is \file{RECORD.ANNOTATOR} at that place (\file{data/100s.qrs}). A URL
  cannot be written.
\item[annotator] Annotator name: the suffix of the file, made of letters, digits,
  \code{\_}, \code{-} and \code{\textasciitilde}.
\item[ann] The annotations (\cref{tab:cpp-annotations}): \code{samples}, and
  \code{symbols} or \code{codes}; \code{subtypes}, \code{chans}, \code{nums} and
  \code{aux} are optional. Each of these vectors is empty (\code{subtypes}, \code{chans}
  and \code{nums} are then 0, and there is no aux) or has one entry per annotation.
  \code{ann.record}, \code{ann.annotator} and \code{ann.fs} are not used.
\item[fs] 0 (none, the default) or a finite frequency $> 0$, in which the times are
  counted, recorded in the file.
\item[dir] If not empty, the directory (which must exist) in which to create the file, for
  this call only; \code{record} must then be a plain name, without a directory: the file
  is \file{DIR/RECORD.ANNOTATOR}. Default \code{\{\}}: as \code{record} says.
\item[sort] \code{true} to always sort the file, \code{false} to keep the given order
  (without the library's advice to run \code{sortann}), for this call only. Default
  \code{std::nullopt}: the environment variable \code{WFDBANNSORT} decides (sorting is on
  unless it is \code{0}).
\end{params}
\returns{Nothing.}
\raises{\code{Error} if a per-annotation field has another length than 0 or the number of
annotations (``write\_annotations: symbols must have one entry per annotation (N), or
none'', for example); if \code{symbols} and \code{codes} are both given and disagree
(``write\_annotations: symbols and codes disagree at annotation K'', K counted from 0);
if \code{fs} is not 0 and not a finite frequency $> 0$; if the file cannot be created
(``can't create annotator ANN for record REC''); if it cannot be written completely
(``write error on annotator ANN for record REC''); or, with \code{dir}, as
\code{write\_record} does.}
\begin{lstlisting}[language=C++,style=example]
// Three beats and a rhythm change, by symbol, with an aux string for the rhythm.
Annotations a;
a.samples = {18, 77, 370, 662};
a.symbols = {"+", "N", "V", "N"};
a.aux = {std::string("(N"), std::nullopt, std::nullopt, std::nullopt};
write_annotations("100s", "test", a);
Annotations b = read_annotations("100s", "test");
std::printf("%s %s %s\n", b.symbols[0].c_str(), b.symbols[2].c_str(),
            b.aux[0]->c_str());                          // + V (N

// A copy of the reference annotations, as read_annotations returns them.
Annotations atr = read_annotations("100s", "atr");
write_annotations("100s", "copy", atr);
Annotations c = read_annotations("100s", "copy");
std::printf("%d\n", c.samples == atr.samples && c.aux == atr.aux);   // 1

// Annotations out of time order are sorted when the file is written (the library
// reports it on the standard error).
Annotations u;
u.samples = {500, 100};
u.codes = {NORMAL, NORMAL};
write_annotations("100s", "unsorted", u);
Annotations s = read_annotations("100s", "unsorted");
std::printf("%lld %lld\n", s.samples[0], s.samples[1]);    // 100 500
\end{lstlisting}

% ------------------------------------------------------------------------------------------
\subsection{Checking records}
\label{sec:cpp-checking}

\apientry{check\_record}{Check a record and its annotation files}
\begin{lstlisting}[language=C++,style=signature]
CheckReport check_record(const std::string &record,
                         const std::vector<std::string> &annotators = {},
                         const InputOptions &input = {});
\end{lstlisting}
Checks the header of a record, its signal specifications, every sample of its signal
files, each segment of a multi-segment record, and the annotation files named, and
reports each problem found. A problem in the record never throws: it is reported, with a
severity (\code{"error"}: the record is damaged or does not follow the WFDB format, and
programs will read wrong data or fail; \code{"warning"}: the record can be read, but
something is missing or unusual, and may give wrong results; \code{"info"}: a fact worth
knowing, such as invalid samples). The problems, their codes and their messages are the
same in the seven languages (\cref{tab:python-issue-codes}). The checks run in normal
mode, even if legacy mode is on, and the library prints no messages during the call (the
previous settings are restored afterwards, the time resolution of output annotators and
the getvec mode among them). They are similar to those of Benjamin Moody's
\code{wfdbcheck}.
\begin{params}
\item[record] Record name, without the \file{.hea} extension (a trailing \code{".hea"} is
  also accepted): a full path, absolute (\code{"/data/mitdb/100"}) or in the home
  directory (\code{"\textasciitilde/Documents/100s"}), used as given; a relative path
  (\code{"100s"}, \code{"data/100s"}), read from the input directory if one is set,
  otherwise from the current directory; or a URL
  (\code{"https://physionet.org/files/mitdb/1.0.0/100"}), read only if network access is
  on (as for \code{read\_record}).
\item[annotators] The annotator names whose files are checked, such as
  \code{\{"atr"\}}. Default \code{\{\}}: none.
\item[input] The input directory for this call only (optional), under which relative
  record names are read (it does not apply to full paths or URLs), and network access
  (\cref{tab:cpp-inputoptions}). Default \code{\{\}}: the current settings.
\end{params}
\returns{A \code{CheckReport} (\cref{tab:cpp-checkreport}): \code{report.ok} is
\code{true} if no problem is an error, \code{report.issues} lists the problems
(\code{Issue} values, \cref{tab:cpp-issue}), and \code{report.to\_string()} gives one line
per problem, or \code{no problems found}.}
\raises{\code{Error} if \code{record} is empty (``record must not be empty''). A record
that cannot be read is reported (code \code{header\_unreadable}), not thrown.}
\begin{lstlisting}[language=C++,style=example]
CheckReport report = check_record("100s", {"atr"});
std::printf("%d\n", report.ok);                    // 1: no errors
std::printf("%s\n", report.to_string().c_str());
// warning: signal 0: the units are not specified
// warning: signal 1: the units are not specified

// A damaged record: one sample of a 12-bit signal is overwritten with 2048.
std::vector<WFDB_Sample> x(100);
for (int i = 0; i < 100; i++) x[i] = i;
WriteOptions w;
w.gains = {100};
w.adcres = {12};
w.units = {"mV"};
w.names = {"ramp"};
write_record("damaged", x, 1, 100, w);
std::FILE *f = std::fopen("damaged.dat", "r+b");
std::fseek(f, 20, SEEK_SET);           // sample 10 (2 bytes per sample in format 16)
std::fputc(0x00, f);                   // 0x0800 = 2048, little-endian
std::fputc(0x08, f);
std::fclose(f);
report = check_record("damaged");
std::printf("%d\n", report.ok);                    // 0
for (const Issue &i : report.issues)
    std::printf("%s %s %s\n", i.severity.c_str(), i.code.c_str(), i.where.c_str());
// error sample_range signal 0
// error checksum signal 0
std::printf("%s\n", report.lines()[0].c_str());
// error: signal 0: 1 sample is outside the ADC range (-2048 to 2047): sample 10 (value 2048)
\end{lstlisting}

% ------------------------------------------------------------------------------------------
\subsection{Header metadata}
\label{sec:cpp-metadata}

These functions read and change the metadata of a header: its info strings of the form
\code{\# Key: value}, such as \code{\# Age: 63} (\cref{chap:metadata}). In short: keys are
compared without regard to letter case (A-Z and a-z only); a key must not be empty and
must not contain \code{:}, \code{\#}, CR, LF or NUL, and a value must not contain CR, LF
or NUL; spaces, tabs, VT, FF and CR at either end of keys and values are removed; keys and
values are bytes (\code{std::string}). A change rewrites the header in one step, through a
temporary file in the directory of the destination, which then replaces it, keeping its
permissions and every other byte of it; all the arguments are checked first, so an error
leaves the header unchanged, and nothing is written if nothing changes (unless
\code{dir} is given). Only the \file{.hea} file is read (not a \file{.info} file); for a
multi-segment record, it is the master header, unless a segment is named. The results
and arguments are of the types \code{Metadata} and \code{MetadataChanges}:
\begin{lstlisting}[language=C++]
using Metadata = std::vector<std::pair<std::string, std::string>>;
using MetadataChanges =
    std::vector<std::pair<std::string, std::optional<std::string>>>;
\end{lstlisting}
Errors throw \code{ArgumentError} (an invalid argument) or \code{FileError} (a file that
cannot be read or written), both subclasses of \code{Error} (\cref{sec:cpp-warnings}).
These functions do not call \code{wfdbquit()}. Metadata often holds information about
patients: WFDB-x stores it as plain text, exactly as given, and does not de-identify it.

\begin{phibox}
A header can hold identifying data about a patient (a name, a date of birth, a medical
record number). The metadata functions read and write such data as given. Records that
are shared must be de-identified first.
\end{phibox}

\apientry{read\_metadata}{Read the metadata of a record}
\begin{lstlisting}[language=C++,style=signature]
Metadata read_metadata(const std::string &record,
                       const InputOptions &input = {});
\end{lstlisting}
Reads the metadata entries of the header of a record, in file order. An info string is an
entry if it contains a colon: the key is the text before the first colon, the value the
rest of the line, both without spaces at either end. Info strings without a colon, or
whose key is empty or contains \code{\#}, are free comments, which are not returned.
\begin{params}
\item[record] Record name, without the \file{.hea} extension (a trailing \code{".hea"} is
  also accepted): a full path, absolute (\code{"/data/mitdb/100"}) or in the home
  directory (\code{"\textasciitilde/Documents/100s"}), used as given; a relative path
  (\code{"100s"}, \code{"data/100s"}), read from the input directory if one is set,
  otherwise from the current directory; or a URL
  (\code{"https://physionet.org/files/mitdb/1.0.0/100"}), read only if network access is
  on.
\item[input] The input directory for this call only (optional), under which relative
  record names are read (it does not apply to full paths or URLs), and network access
  (\cref{tab:cpp-inputoptions}). Default \code{\{\}}: the current settings.
\end{params}
\returns{The entries, as (key, value) pairs, in file order. A key that appears more than
once (ignoring letter case) is returned once, with its first value, with a warning
(\cref{sec:cpp-warnings}).}
\raises{\code{FileError} if the header cannot be found or read (``can't read header for
record REC''); \code{ArgumentError} if \code{record} is empty (``record name must not be
empty'') or \code{"-"} (the standard input).}
\begin{lstlisting}[language=C++,style=example]
WriteOptions w;
w.gains = {200};
w.adcres = {12};
w.units = {"mV"};
w.names = {"ECG"};
w.comments = {"recorded at rest"};
w.metadata = {{"Age", "63"}, {"Sex", "F"}};
write_record("pat", std::vector<WFDB_Sample>(360, 0), 1, 360, w);
for (const auto &e : read_metadata("pat"))
    std::printf("%s = %s\n", e.first.c_str(), e.second.c_str());
// Age = 63
// Sex = F
for (const std::string &c : read_header("pat").comments)
    std::printf("[%s]\n", c.c_str());
// [recorded at rest]: the free comment is not an entry
// [ Age: 63]
// [ Sex: F]
\end{lstlisting}

\apientry{get\_metadata}{Return the value of a metadata entry}
\begin{lstlisting}[language=C++,style=signature]
std::optional<std::string> get_metadata(const std::string &record,
                                        const std::string &key,
                                        const InputOptions &input = {});
\end{lstlisting}
Returns the value of one metadata entry of the header of a record.
\begin{params}
\item[record] Record name (as for \code{read\_metadata}).
\item[key] The key, such as \code{"Age"}. Letter case is ignored, and spaces at either end
  are removed; it must be a valid key (see \code{set\_metadata}).
\item[input] The input directory and network access for this call only. Default
  \code{\{\}}: the current settings.
\end{params}
\returns{The value (\code{""} for an entry with an empty value), or \code{std::nullopt} if
the header has no entry for \code{key}. If it has more than one, the first is returned,
with the warning of \code{read\_metadata}.}
\raises{As \code{read\_metadata}; also \code{ArgumentError} if \code{key} is not valid.}
\begin{lstlisting}[language=C++,style=example]
WriteOptions w;
w.gains = {200};
w.adcres = {12};
w.units = {"mV"};
w.metadata = {{"Age", "63"}, {"Notes", ""}};
write_record("pat", std::vector<WFDB_Sample>(360, 0), 1, 360, w);
std::printf("%s\n", get_metadata("pat", "age")->c_str());       // 63
std::printf("[%s]\n", get_metadata("pat", " NOTES ")->c_str()); // []: an empty value
std::printf("%d\n", get_metadata("pat", "Weight").has_value()); // 0: no such entry
\end{lstlisting}

\apientry{set\_metadata}{Set the value of a metadata entry}
\begin{lstlisting}[language=C++,style=signature]
void set_metadata(const std::string &record, const std::string &key,
                  const std::string &value, const MetadataOptions &opt = {});
\end{lstlisting}
Sets the value of a metadata entry in the header of a record. If the header has an entry
for \code{key} (ignoring letter case), its line is replaced by \code{\# Key: value}, at
the same position and with its line ending, with the spelling of \code{key} given here
(even if the value is the same). Otherwise the entry is added after the last info string
of the header (or, if it has none, after its last line that is not blank), with the line
ending of the record line. Every other byte of the header is kept.
\begin{params}
\item[record] Record name (as for \code{read\_metadata}). The header is changed where it
  was read, unless \code{opt.dir} is given; a header read from a URL can only be written
  elsewhere, with \code{opt.dir}.
\item[key] The key, such as \code{"Age"} or \code{"Date of birth"}. Spaces at either end
  are removed; it must not then be empty, and must not contain \code{:}, \code{\#}, CR,
  LF or NUL.
\item[value] The value, which may be a sentence. Spaces at either end are removed; it
  must not contain CR, LF or NUL, and may be empty (it is then written \code{\# Key:}).
\item[opt] The input directory and network access for reading the header, and the
  directory \code{dir} in which to write the changed header, for this call only
  (\cref{tab:cpp-metadataoptions}). Default \code{\{\}}: the current settings, and the
  header is changed where it was read.
\end{params}
\returns{Nothing.}
\raises{\code{ArgumentError} if \code{key} or \code{value} is not valid; if the header has
more than one entry for \code{key} (``header of record REC has more than one entry for key
KEY (lines L1, L2); remove them first''); if it has no record line; if it was read from a
URL and \code{opt.dir} is not given; or if \code{record} is empty or \code{"-"}.
\code{FileError} if the header cannot be read; if the changed header cannot be written
(``can't write header for record REC''); or if \code{opt.dir} does not exist
(``directory D does not exist'').}
\begin{lstlisting}[language=C++,style=example]
WriteOptions w;
w.gains = {200};
w.adcres = {12};
w.units = {"mV"};
w.metadata = {{"Age", "63"}};
write_record("pat", std::vector<WFDB_Sample>(360, 0), 1, 360, w);
set_metadata("pat", "AGE", "64");                  // replaces the line: "# AGE: 64"
set_metadata("pat", "Date of birth", "1961-02-03");    // adds "# Date of birth: 1961-02-03"
for (const auto &e : read_metadata("pat"))
    std::printf("%s = %s\n", e.first.c_str(), e.second.c_str());
// AGE = 64
// Date of birth = 1961-02-03
try {
    set_metadata("pat", "a:b", "1");
} catch (const ArgumentError &e) {
    std::printf("%s\n", e.what());      // metadata key 'a:b' must not contain ':'
}
\end{lstlisting}

\apientry{remove\_metadata}{Remove a metadata entry}
\begin{lstlisting}[language=C++,style=signature]
int remove_metadata(const std::string &record, const std::string &key,
                    const MetadataOptions &opt = {});
\end{lstlisting}
Removes a metadata entry from the header of a record. Every entry for \code{key}
(ignoring letter case) is removed, with its line ending, so this also repairs a header
that has more than one. Every other byte is kept. Nothing is written if the header has no
entry for \code{key} (unless \code{opt.dir} is given).
\begin{params}
\item[record] Record name (as for \code{set\_metadata}).
\item[key] The key (as for \code{set\_metadata}).
\item[opt] The input directory, network access and output directory \code{dir} of this
  call (as for \code{set\_metadata}). Default \code{\{\}}.
\end{params}
\returns{The number of lines removed (0 if the header has no entry for \code{key}).}
\raises{As \code{set\_metadata}.}
\begin{lstlisting}[language=C++,style=example]
WriteOptions w;
w.gains = {200};
w.adcres = {12};
w.units = {"mV"};
w.metadata = {{"Age", "63"}, {"Sex", "F"}};
write_record("pat", std::vector<WFDB_Sample>(360, 0), 1, 360, w);
std::printf("%d\n", remove_metadata("pat", "sex"));       // 1
std::printf("%d\n", remove_metadata("pat", "Weight"));    // 0: nothing is written
std::printf("%zu\n", read_metadata("pat").size());        // 1
\end{lstlisting}

\apientry{update\_metadata}{Make several changes to the metadata}
\begin{lstlisting}[language=C++,style=signature]
void update_metadata(const std::string &record, const MetadataChanges &changes,
                     const MetadataOptions &opt = {});
\end{lstlisting}
Makes several changes to the metadata of a record in one write. Each (key, value) of
\code{changes} is applied in order, as \code{set\_metadata} does, or, if the value is
\code{std::nullopt}, as \code{remove\_metadata} does. All the changes are checked first:
if one is not valid, none is made.
\begin{params}
\item[record] Record name (as for \code{set\_metadata}).
\item[changes] The changes, as (key, value) pairs; a value of \code{std::nullopt} removes
  the entry. No two keys may differ only in letter case. Empty \code{changes} change
  nothing (and write nothing, unless \code{opt.dir} is given).
\item[opt] The input directory, network access and output directory \code{dir} of this
  call (as for \code{set\_metadata}). Default \code{\{\}}.
\end{params}
\returns{Nothing.}
\raises{As \code{set\_metadata}; also \code{ArgumentError} if two keys differ only in
letter case (``metadata key 'K' is given more than once (keys are compared without regard
to letter case)'').}
\begin{lstlisting}[language=C++,style=example]
WriteOptions w;
w.gains = {200};
w.adcres = {12};
w.units = {"mV"};
w.metadata = {{"Age", "63"}, {"Dx", "164889003"}};
write_record("pat", std::vector<WFDB_Sample>(360, 0), 1, 360, w);
update_metadata("pat", {{"Age", "64"}, {"Dx", std::nullopt}, {"Sex", "F"}});
for (const auto &e : read_metadata("pat"))
    std::printf("%s = %s\n", e.first.c_str(), e.second.c_str());
// Age = 64
// Sex = F

// The header "pat" in the directory "in" is unchanged; the changed copy is written to
// the directory "out" (which must exist).
std::filesystem::create_directories("in");
std::filesystem::create_directories("out");
std::filesystem::copy_file("pat.hea", "in/pat.hea");
MetadataOptions o;
o.input_dir = "in";
o.dir = "out";
update_metadata("pat", {{"Age", "65"}}, o);
std::printf("%s %s\n", get_metadata("in/pat", "Age")->c_str(),
            get_metadata("out/pat", "Age")->c_str());      // 64 65
\end{lstlisting}

% ------------------------------------------------------------------------------------------
\subsection{The options}
\label{sec:cpp-options}

The options of the high-level functions are the fields of the structures of
\cref{tab:cpp-inputoptions,tab:cpp-readoptions,tab:cpp-writeoptions,tab:cpp-metadataoptions}.
\code{ReadOptions} and \code{MetadataOptions} are derived from \code{InputOptions}, and
have its fields too. A field that is not set keeps its default.

\begin{xltabular}{\linewidth}{@{}lL@{}}
\caption{The fields of \code{InputOptions}: where \code{read\_header},
  \code{read\_annotations}, \code{check\_record}, \code{read\_metadata} and
  \code{get\_metadata} find files, for one call}
\label{tab:cpp-inputoptions} \\
\toprule
\tabhead Field & \tabhead Meaning \\
\midrule
\endfirsthead
\toprule
\tabhead Field & \tabhead Meaning \\
\midrule
\endhead
\bottomrule
\endfoot
\code{std::optional<std::string> input\_dir} & the input directory for this call only
  (optional; see \code{wfdbx\_set\_input\_dir}): a local directory, a
  \code{\textasciitilde} path or the URL of a directory, under which relative record names
  are read. It does not apply to full paths or URLs. \code{""} means the default
  (\code{WFDBX\_INPUT\_DIR}, or the current directory). Not set: the current setting \\
\code{std::optional<bool> network} & \code{true} to allow opening URLs for this call only,
  \code{false} to forbid it (see \code{wfdbx\_set\_network}). Not set: the current
  setting \\
\end{xltabular}

\begin{xltabular}{\linewidth}{@{}lL@{}}
\caption{The fields of \code{ReadOptions}, the options of \code{read\_record} (with
  \code{input\_dir} and \code{network} of \code{InputOptions})}
\label{tab:cpp-readoptions} \\
\toprule
\tabhead Field & \tabhead Meaning (default) \\
\midrule
\endfirsthead
\toprule
\tabhead Field & \tabhead Meaning (default) \\
\midrule
\endhead
\bottomrule
\endfoot
\code{std::vector<std::string> channels} & the signals to return, in this order, by number
  (\code{"0"}, \code{"1"}, \ldots) or by name (\code{"V5"}, as in \code{Record::names()},
  then as \code{findsig} finds it). Default empty: all the signals \\
\code{TimeSpec from} & the first sample: a sample number ($\ge 0$) or a time string.
  Default empty: 0 \\
\code{TimeSpec to} & the sample after the last one: a sample number or a time string.
  Default empty: the end of the record \\
\code{bool highres} & \code{true} to read a multi-frequency record at its highest sampling
  frequency, repeating the samples of the slower signals; \code{false} averages the
  samples of the faster signals to one per frame. Default \code{false} \\
\code{bool pad} & \code{true} to replace each invalid sample by the previous valid sample
  of the same signal. Default \code{false} \\
\code{bool checksum\_warning} & \code{true} to print a warning on the standard error when
  the samples do not match the checksums of the header (\cref{sec:cpp-warnings}); the
  samples are returned either way. Default \code{true} \\
\code{bool legacy} & behave as the original WFDB package, for this call only. Default
  \code{false} \\
\end{xltabular}

\begin{xltabular}{\linewidth}{@{}lL@{}}
\caption{The fields of \code{WriteOptions}, the options of \code{write\_record} and
  \code{write\_record\_physical}. Each per-signal vector holds one value for all the
  signals, or one value per signal (or none, for those that have a default)}
\label{tab:cpp-writeoptions} \\
\toprule
\tabhead Field & \tabhead Meaning (default) \\
\midrule
\endfirsthead
\toprule
\tabhead Field & \tabhead Meaning (default) \\
\midrule
\endhead
\bottomrule
\endfoot
\code{int fmt} & the storage format, the same for all the signals, such as \code{16},
  \code{212}, \code{24} or \code{516} (FLAC) (\cref{chap:formats}). Default \code{16} \\
\code{std::vector<WFDB\_Gain> gains} & the gains, in ADC units per physical unit:
  required, and finite. A gain of 0 marks a signal as uncalibrated, which only
  \code{write\_record} (ADC units) allows. Negative gains are allowed \\
\code{std::vector<int> baselines} & the ADC values that correspond to 0 physical units.
  Default: the \code{adczeros} values \\
\code{std::vector<int> adczeros} & the ADC output for an input of 0\,V. Default
  \code{\{0\}} \\
\code{std::vector<int> adcres} & the ADC resolutions, in bits: required, each at least
  1 \\
\code{std::vector<std::optional<std::string>> units} & the physical units, such as
  \code{"mV"}: required for every signal whose gain is not 0; optional for uncalibrated
  signals \\
\code{std::vector<std::string> names} & the signal descriptions. Default: \code{""} \\
\code{std::vector<int> samps\_per\_frame} & the number of samples per frame of each
  signal, for a multi-frequency record (each $\ge 1$). Default: 1 \\
\code{std::optional<std::string> base\_time} & the start time and date, such as
  \code{"12:00:00 01/01/2020"} (see \code{setbasetime}). Default: none \\
\code{std::vector<std::string> comments} & the info strings, written to the header as
  comment lines (\code{"\#"} followed by the string). Default: none \\
\code{Metadata metadata} & metadata entries (\cref{chap:metadata}) written to the header
  after the comments, as \code{\# Key: value} (through \code{putinfo}), in this order. The
  keys and values follow the rules of \code{set\_metadata}; no two keys may differ only
  in letter case, and none may be the key of an entry in \code{comments}
  (\code{ArgumentError}, before any file is created). Default: none \\
\code{std::string dir} & if not empty, the directory (which must exist) in which the
  files are created, for this call only; the record name must then be a plain name,
  without a directory, and the record opens from that directory under its plain name.
  Default \code{""}: the output directory \\
\code{bool legacy} & behave as the original WFDB package, for this call only: the gain
  defaults to 200, and the ADC resolution to one based on the format. Default
  \code{false} \\
\end{xltabular}

\begin{xltabular}{\linewidth}{@{}lL@{}}
\caption{The fields of \code{MetadataOptions}, the options of \code{set\_metadata},
  \code{remove\_metadata} and \code{update\_metadata} (with \code{input\_dir} and
  \code{network} of \code{InputOptions}, which say where the header is read from)}
\label{tab:cpp-metadataoptions} \\
\toprule
\tabhead Field & \tabhead Meaning \\
\midrule
\endfirsthead
\toprule
\tabhead Field & \tabhead Meaning \\
\midrule
\endhead
\bottomrule
\endfoot
\code{std::string dir} & if not empty, the directory (which must exist; a
  \code{\textasciitilde} path is expanded) in which the changed header is written, under
  the file name of the header that was read, instead of changing that header. The header
  is then written even if nothing changes. The signal files are not copied. Default
  \code{""}: the header is changed where it was read \\
\end{xltabular}

\begin{xltabular}{\linewidth}{@{}lL@{}}
\caption{The members of \code{TimeSpec}, a sample number or a time string}
\label{tab:cpp-timespec} \\
\toprule
\tabhead Member & \tabhead Meaning \\
\midrule
\endfirsthead
\toprule
\tabhead Member & \tabhead Meaning \\
\midrule
\endhead
\bottomrule
\endfoot
\code{TimeSpec()} & an empty value: no time is given \\
\code{TimeSpec(WFDB\_Time t)}, \code{TimeSpec(int t)}, \code{TimeSpec(long t)} & the
  sample number \code{t} \\
\code{TimeSpec(const char *s)}, \code{TimeSpec(std::string s)} & the time string
  \code{s} \\
\code{bool empty() const} & \code{true} if no time was given \\
\code{bool needs\_freq() const} & \code{true} if the time is a string that needs a sampling
  frequency to be converted (see \code{wfdbx\_time\_needs\_freq}) \\
\code{WFDB\_Time resolve() const} & the sample number (0 if empty); a time string is
  converted with \code{strtim} for the record open at the time, and a time of day gives
  the positive sample number it corresponds to \\
\end{xltabular}

% ------------------------------------------------------------------------------------------
\subsection{The result types}
\label{sec:cpp-results}

The readers return structures of the types \code{Record}, \code{Header} and
\code{Annotations}, whose members are listed in
\cref{tab:cpp-record,tab:cpp-header,tab:cpp-annotations}; \code{check\_record} returns a
\code{CheckReport}, which holds \code{Issue} values
(\cref{tab:cpp-checkreport,tab:cpp-issue}). The signals of a record are described by
\code{WFDB\_Siginfo} structures (\cref{tab:cpp-siginfo}), which are those of the WFDB
library (\cref{sec:cpp-library}). Their fields are public, and may be changed.

\begin{xltabular}{\linewidth}{@{}lL@{}}
\caption{The members of \code{Record}, the result of \code{read\_record}}
\label{tab:cpp-record} \\
\toprule
\tabhead Member & \tabhead Contents \\
\midrule
\endfirsthead
\toprule
\tabhead Member & \tabhead Contents \\
\midrule
\endhead
\bottomrule
\endfoot
\code{std::string name} & the record name, as given to \code{read\_record} \\
\code{WFDB\_Frequency fs} & the sampling frequency of the returned samples, in Hz, or 0 if
  the header gives none; with \code{highres}, the highest sampling frequency of a
  multi-frequency record \\
\code{int nsig} & the number of signals returned \\
\code{std::vector<int> signals} & the number, in the record, of each signal returned \\
\code{std::vector<WFDB\_Sample> samples} & the samples in ADC units, frame by frame:
  \code{samples[i * nsig + s]} is sample \code{i} of signal \code{s}; invalid samples are
  \code{WFDB\_INVALID\_SAMPLE} \\
\code{std::vector<WFDB\_Siginfo> siginfo} & one per signal returned, in the order of the
  signals; \code{desc} has no white space at either end \\
\code{std::optional<std::string> base\_time} & the start time and date as text, to the
  second, such as \code{"[08:45:00 01/01/2001]"}; \code{std::nullopt} if the header gives
  none \\
\code{std::optional<long long> base\_time\_ms} & the start time in milliseconds after
  midnight (\code{"14:00:03.500"} gives \code{50403500}); given exactly when
  \code{base\_time} is \\
\code{std::optional<std::array<int, 3>> base\_date} & the start date as
  \code{\{year, month, day\}}; \code{std::nullopt} if the header gives no date \\
\code{WFDB\_Time sampfrom} & the number of the first sample returned \\
\code{std::vector<std::string> comments} & the info strings of the record: the comments of
  the header, and the lines of its \file{.info} file if it has one (see \code{getinfo}) \\
\code{bool checksum\_ok} & \code{false} if the samples do not match the checksums of the
  header; checksums are checked only when a read reaches the end of the record, so a
  shorter read gives \code{true} \\
\code{bool legacy} & \code{true} if the record was read in legacy mode; \code{physical()}
  then uses a gain of 200 for uncalibrated signals \\
\code{size\_t nsamples() const} & the number of samples per signal \\
\code{WFDB\_Sample at(size\_t i, int s) const} & sample \code{i} of signal \code{s}, in ADC
  units (not bounds-checked) \\
\code{double physical(size\_t i, int s) const} & sample \code{i} of signal \code{s} in
  physical units, $(\mathit{value} - \mathit{baseline}) / \mathit{gain}$; NaN if the
  sample is invalid. Throws \code{Error} (``signal N of record REC is uncalibrated'') if
  the gain of the signal is 0 \\
\code{std::vector<double> physical\_signal(int s) const} & all the samples of signal
  \code{s} in physical units \\
\code{std::vector<std::string> names() const} & the signal descriptions (\code{""} if
  absent), without the white space that a header may put at either end of them \\
\end{xltabular}

\begin{xltabular}{\linewidth}{@{}lL@{}}
\caption{The members of \code{Header}, the result of \code{read\_header}}
\label{tab:cpp-header} \\
\toprule
\tabhead Member & \tabhead Contents \\
\midrule
\endfirsthead
\toprule
\tabhead Member & \tabhead Contents \\
\midrule
\endhead
\bottomrule
\endfoot
\code{std::string name} & the record name, as given to \code{read\_header} \\
\code{WFDB\_Frequency fs} & the sampling frequency in Hz, as \code{sampfreq()} gives it; 0
  if the header gives none \\
\code{WFDB\_Frequency counter\_freq} & the counter frequency (\code{getcfreq()}) \\
\code{double base\_counter} & the base counter value (\code{getbasecount()}) \\
\code{WFDB\_Time nsamples} & the number of samples per signal, from the record line of the
  header (0 if the header does not give it) \\
\code{int nsig} & the number of signals \\
\code{std::vector<WFDB\_Siginfo> siginfo} & the signal specifications, one per signal;
  \code{desc} has no white space at either end \\
\code{base\_time}, \code{base\_time\_ms}, \code{base\_date} & the start time and date, as
  in \code{Record} \\
\code{std::vector<std::string> comments} & the info strings of the record, as in
  \code{Record} \\
\code{Metadata metadata} & the metadata entries of the header, in file order, as
  \code{read\_metadata} returns them: from the \file{.hea} file only; empty, without an
  error, for an EDF record and for a header read from the standard input \\
\end{xltabular}

\begin{xltabular}{\linewidth}{@{}lL@{}}
\caption{The members of \code{Annotations}, the result of \code{read\_annotations} and the
  argument of \code{write\_annotations}: parallel vectors, annotation \code{i} being
  \code{samples[i]}, \code{codes[i]}, \code{symbols[i]}, \ldots, in file order}
\label{tab:cpp-annotations} \\
\toprule
\tabhead Member & \tabhead Contents \\
\midrule
\endfirsthead
\toprule
\tabhead Member & \tabhead Contents \\
\midrule
\endhead
\bottomrule
\endfoot
\code{std::string record}, \code{annotator} & the record and annotator names, such as
  \code{"100s"} and \code{"atr"} \\
\code{WFDB\_Frequency fs} & the frequency in which \code{samples} are counted, normally the
  sampling frequency of the record; 0 if unknown \\
\code{std::vector<WFDB\_Time> samples} & the annotation times, in samples \\
\code{std::vector<int> codes} & the annotation codes, such as \code{NORMAL} or \code{PVC} \\
\code{std::vector<std::string> symbols} & the mnemonics of the codes, as \code{annstr}
  returns them, such as \code{"N"}, \code{"V"} or \code{"+"} \\
\code{std::vector<int> subtypes}, \code{chans}, \code{nums} & the subtype, channel number
  and annotator number of each annotation \\
\code{std::vector<std::optional<std::string>> aux} & the auxiliary text of each
  annotation, up to its first NUL byte, or \code{std::nullopt} where there is none \\
\code{size\_t size() const} & the number of annotations \\
\end{xltabular}

\begin{table}[htbp]
\centering
\caption{The members of \code{CheckReport}, the result of \code{check\_record}}
\label{tab:cpp-checkreport}
\begin{tabularx}{\linewidth}{@{}lL@{}}
\toprule
\tabhead Member & \tabhead Contents \\
\midrule
\code{std::string record} & the record name, as given \\
\code{bool ok} & \code{true} if no issue has the severity \code{"error"} \\
\code{std::vector<Issue> issues} & the problems, in the order of the checks \\
\code{std::vector<std::string> lines() const} & the text form of the report: one line per
  issue (\code{Issue::to\_string()}), or the one line \code{no problems found} \\
\code{std::string to\_string() const} & the lines of \code{lines()}, joined by line feeds
  (no final line feed) \\
\bottomrule
\end{tabularx}
\end{table}

\begin{table}[htbp]
\centering
\caption{The members of \code{Issue}, a problem found by \code{check\_record}}
\label{tab:cpp-issue}
\begin{tabularx}{\linewidth}{@{}lL@{}}
\toprule
\tabhead Member & \tabhead Contents \\
\midrule
\code{std::string severity} & \code{"error"}, \code{"warning"} or \code{"info"} \\
\code{std::string code} & a stable identifier of the kind of problem, such as
  \code{"checksum"} or \code{"gain\_missing"} (\cref{tab:python-issue-codes}) \\
\code{std::string where} & where the problem is: \code{""} (the record),
  \code{"signal N"}, \code{"segment SEG"}, \code{"segment SEG, signal N"} or
  \code{"annotator ANN"} \\
\code{std::string message} & the description of the problem \\
\code{std::string to\_string() const} & \code{"SEVERITY: WHERE: MESSAGE"}, or
  \code{"SEVERITY: MESSAGE"} when \code{where} is empty \\
\bottomrule
\end{tabularx}
\end{table}

\begin{xltabular}{\linewidth}{@{}lL@{}}
\caption{The fields of \code{WFDB\_Siginfo} (also named \code{Siginfo}), the specification
  of a signal}
\label{tab:cpp-siginfo} \\
\toprule
\tabhead Field & \tabhead Contents \\
\midrule
\endfirsthead
\toprule
\tabhead Field & \tabhead Contents \\
\midrule
\endhead
\bottomrule
\endfoot
\code{std::optional<std::string> fname} & the name of the signal file, as the header gives
  it \\
\code{std::optional<std::string> desc} & the signal description, such as \code{"MLII"} \\
\code{std::optional<std::string> units} & the physical units, such as \code{"mV"};
  \code{std::nullopt} if not specified \\
\code{WFDB\_Gain gain} & the gain, in ADC units per physical unit (0: uncalibrated) \\
\code{WFDB\_Sample initval} & the value of the first sample \\
\code{WFDB\_Group group} & the signal group number (the signals of one group share one
  file) \\
\code{int fmt} & the storage format, such as \code{16}, \code{212} or \code{516} \\
\code{int spf} & the number of samples per frame (more than 1 for the faster signals of a
  multi-frequency record) \\
\code{int bsize} & the block size, in bytes (0 for ordinary files) \\
\code{int adcres} & the ADC resolution, in bits (0: not specified) \\
\code{int adczero} & the ADC output for an input of 0\,V \\
\code{int baseline} & the ADC value that corresponds to 0 physical units \\
\code{long nsamp} & the number of samples (0: not specified) \\
\code{int cksum} & the 16-bit checksum of the samples \\
\end{xltabular}

% ------------------------------------------------------------------------------------------
\subsection{Warnings and errors}
\label{sec:cpp-warnings}

\paragraph{Exceptions.} The high-level functions report errors by throwing exceptions of
the classes of \cref{tab:cpp-exceptions}, all declared in \file{<wfdbx/record.hpp>}.
\code{e.what()} returns the message, which describes the failed operation, such as
``can't read record REC''. Each function checks its arguments before it creates or
changes a file, so an invalid argument leaves the files as they were. The library may
also print a message of its own on the standard error when it cannot open a file (as in
\code{init: can't open header for record REC}); \code{wfdberror()} returns the last such
message, and \code{wfdbquiet()} turns these messages off (\cref{sec:cpp-library}).

\begin{table}[htbp]
\centering
\caption{The exceptions of the high-level API}
\label{tab:cpp-exceptions}
\begin{tabularx}{\linewidth}{@{}lL@{}}
\toprule
\tabhead Exception & \tabhead Thrown when \\
\midrule
\code{Error} (derived from \code{std::runtime\_error}) & a record, header or annotation
  file cannot be found, read, created or written completely (for example on a full
  disk); a signal or annotation file ends before its header or its end-of-file marker
  says; an argument is invalid; or the operation needs a value that the record does not
  give (a time in seconds without a sampling frequency, physical units without a gain) \\
\code{ArgumentError} (derived from \code{Error}) & an argument of a metadata function is
  invalid: an invalid key or value, a record name that is empty or \code{"-"}, a key that
  the header has more than once (\code{set\_metadata}, \code{update\_metadata}), a header
  without a record line, or a header read from a URL to be changed without \code{dir};
  also an invalid \code{WriteOptions::metadata} in \code{write\_record} \\
\code{FileError} (derived from \code{Error}) & a metadata function cannot read or write a
  file: a header that cannot be found or read, a changed header that cannot be written,
  or a \code{dir} that does not exist \\
\bottomrule
\end{tabularx}
\end{table}

\paragraph{Warnings.} Warnings are printed on the standard error with the library's
messages, as \code{wfdb\_error} prints them: \code{wfdbquiet()} silences them, and
\code{wfdberror()} returns the last one. There is no callback.

\paragraph{The checksum warning.} When \code{read\_record} reads to the end of a record
whose samples do not match the checksums of its header, it returns all the samples and
prints the warning ``WFDB-x: checksum error in record REC: the signal data do not match
the checksums in its header''. \code{checksum\_warning = false} turns the warning off;
\code{checksum\_ok} of the result is \code{false} either way.

\paragraph{The duplicate-key warning.} When a header holds two metadata entries whose keys
differ only in letter case (written by other software), \code{read\_metadata},
\code{get\_metadata} and \code{read\_header} use the first and print the warning ``WFDB-x:
header of record REC has more than one entry for key KEY (lines L1, L2); the first is
used''. \code{set\_metadata} and \code{update\_metadata} then refuse to change that key
(\code{ArgumentError}), and \code{remove\_metadata} removes all of its entries.

\begin{lstlisting}[language=C++,style=example]
// A record whose last sample is changed after it was written: the checksum fails.
std::vector<WFDB_Sample> x(100);
for (int i = 0; i < 100; i++) x[i] = i;
WriteOptions w;
w.gains = {100};
w.adcres = {12};
w.units = {"mV"};
write_record("ck", x, 1, 100, w);
std::FILE *f = std::fopen("ck.dat", "r+b");
std::fseek(f, -2, SEEK_END);
std::fputc(0, f);
std::fputc(0, f);
std::fclose(f);
Record rec = read_record("ck");      // prints the checksum warning
std::printf("%zu %d\n", rec.nsamples(), rec.checksum_ok);   // 100 0: all the samples
ReadOptions quiet;
quiet.checksum_warning = false;
rec = read_record("ck", quiet);      // no warning

// Two entries for one key: the first is used, with a warning.
WriteOptions d = w;
d.gains = {0};
d.units = {};
d.comments = {" Age: 63", " AGE: 64"};
write_record("dup", std::vector<WFDB_Sample>(10, 0), 1, 100, d);
std::printf("%s\n", get_metadata("dup", "age")->c_str());   // 63, with a warning
try {
    set_metadata("dup", "Age", "65");
} catch (const ArgumentError &e) {
    std::printf("ArgumentError: %s\n", e.what());
    // ArgumentError: header of record dup has more than one entry for key Age
    // (lines 3, 4); remove them first
}
std::printf("%d\n", remove_metadata("dup", "age"));          // 2

try {
    read_record("nosuchrecord");
} catch (const Error &e) {
    std::printf("%s\n", e.what());   // can't read record nosuchrecord
}
try {
    read_metadata("nosuchrecord");
} catch (const FileError &e) {
    std::printf("%s\n", e.what());   // can't read header for record nosuchrecord
}
\end{lstlisting}

% ==========================================================================================
\section{The WFDB library functions in C++}
\label{sec:cpp-library}

Every function of the WFDB library is declared in \file{<wfdbx/wfdb.hpp>}, in namespace
\code{wfdbx}, under its C name, with the same arguments, in the same order, and the same
return values, so that programs and documentation written for the C library carry over
directly. The functions, their arguments and their results are documented once, for all
the languages, in \cref{chap:library}; this section gives the rules by which the C
declarations become C++ functions (see also \cref{sec:library-cpp}). The header
summarizes each function, and marks what WFDB-x changes in it.

\subsection{Conventions}

\paragraph{Names and headers.} The functions, the types (\code{WFDB\_Sample},
\code{WFDB\_Time}, \ldots), the structures and the constants (\code{WFDB\_READ},
\code{WFDB\_INVALID\_SAMPLE}, \code{WFDB\_HIGHRES}, \ldots) keep their C names, in
namespace \code{wfdbx}. \file{<wfdbx/ecgcodes.hpp>} defines the annotation codes
(\code{NORMAL}, \code{PVC}, \ldots, \cref{chap:formats}), and \file{<wfdbx/ecgmap.hpp>}
the macros of \file{<wfdb/ecgmap.h>} as inline functions (\code{isann}, \code{isqrs},
\code{map1}, \ldots, which call \code{wfdb\_isann}, \code{wfdb\_isqrs},
\code{wfdb\_map1}, \ldots). \file{<wfdbx/record.hpp>} includes \file{<wfdbx/wfdb.hpp>}.

\paragraph{Arguments.} String arguments are \code{const char *}, and \code{nullptr} is
accepted where C accepts \code{NULL}: \code{sampfreq(nullptr)} returns the frequency of
the open record, and \code{getinfo(nullptr)} continues with the info strings of the last
record. Arrays and structures are allocated by the caller and passed by pointer, as in
C: \code{isigopen} fills an array of \code{WFDB\_Siginfo}, \code{getvec} an array of
\code{WFDB\_Sample}, \code{getann} a \code{WFDB\_Annotation}. \code{getseginfo} takes a
\code{const WFDB\_Seginfo **}.

\paragraph{Strings} are returned by value, as \code{std::string}, or as
\code{std::optional<std::string>} where the C function may return \code{NULL}
(\code{annstr}, \code{ecgstr}, \code{anndesc}, \code{getinfo}, \code{wfdbfile}, \ldots);
unlike C's pointers to static buffers, they stay valid after later calls. Text is handled
as bytes.

\paragraph{Return values and errors.} The functions return what the C functions return:
an \code{int} status (negative on errors), a number of signals, a \code{double}, a string
or \code{std::nullopt}. Errors are reported by these values and by a message on the
standard error, as in C, never by exceptions. \code{wfdberror()} returns the last message;
\code{wfdbquiet()} turns the messages off and \code{wfdbverbose()} on again.

\paragraph{State.} As in C, the library keeps its state in the process: the open record,
its input signals, the output signals and the annotators remain open between calls, until
\code{wfdbquit()} closes them. \code{wfdbquit()} leaves the input and output directories
and the network setting as they are. The library is not thread-safe. The high-level
functions of \cref{sec:cpp-api} (except the metadata functions) call \code{wfdbquit()},
which closes anything opened with the library functions.

\paragraph{Additions.} \code{getvecs(buffer, n)} reads up to \code{n} sample vectors at
once, with exactly the result of \code{n} calls of \code{getvec}, but decoding whole
blocks of the signal files; \code{getvecs(buffer, n, \&status)} also gives the status of
the last \code{getvec} call:
\begin{lstlisting}[language=C++]
long long getvecs(WFDB_Sample *out, long long n, int *status = nullptr);
\end{lstlisting}

\subsection{The structures}
\label{sec:cpp-structures}

The structures of the C library have the names and the fields of the C structures,
with C's zero values as defaults; they are also available under the shorter names
\code{Siginfo}, \code{Anninfo}, \code{Annotation}, \code{Calinfo} and \code{Seginfo}
(\cref{tab:cpp-structures}).

\begin{xltabular}{\linewidth}{@{}lL@{}}
\caption{The C structures in C++}
\label{tab:cpp-structures} \\
\toprule
\tabhead Structure & \tabhead Fields \\
\midrule
\endfirsthead
\toprule
\tabhead Structure & \tabhead Fields \\
\midrule
\endhead
\bottomrule
\endfoot
\code{WFDB\_Siginfo} & \code{fname}, \code{desc}, \code{units} (\code{std::optional<std::string>}),
  \code{gain}, \code{initval}, \code{group}, \code{fmt}, \code{spf}, \code{bsize},
  \code{adcres}, \code{adczero}, \code{baseline}, \code{nsamp}, \code{cksum}
  (\cref{tab:cpp-siginfo}) \\
\code{WFDB\_Seginfo} & \code{recname} (the name of the segment, a \code{char} array),
  \code{nsamp} (its number of samples), \code{samp0} (the number of its first sample in
  the record) \\
\code{WFDB\_Anninfo} & \code{name} (the annotator name, a \code{std::string}), \code{stat}
  (\code{WFDB\_READ} or \code{WFDB\_WRITE} for the MIT format, \code{WFDB\_AHA\_READ} or
  \code{WFDB\_AHA\_WRITE} for the AHA format) \\
\code{WFDB\_Annotation} & \code{time} (\code{WFDB\_Time}, in sample intervals),
  \code{anntyp} (\code{char}), \code{subtyp} (\code{signed char}), \code{chan}
  (\code{unsigned char}), \code{num} (\code{signed char}), \code{aux}
  (\code{std::optional<std::string>}: absent, or a string whose first byte is the length
  of the text that follows, as in C) \\
\code{WFDB\_Calinfo} & \code{low}, \code{high} (the levels of the calibration pulse, in
  physical units), \code{scale} (physical units per cm), \code{sigtype} and \code{units}
  (\code{std::optional<std::string>}), \code{caltype} \\
\end{xltabular}

\code{WFDB\_Time} is \code{long long} (64 bits), \code{WFDB\_Sample} is \code{int},
\code{WFDB\_Date} is \code{long}, \code{WFDB\_Frequency} and \code{WFDB\_Gain} are
\code{double}, and \code{WFDB\_Group}, \code{WFDB\_Signal} and \code{WFDB\_Annotator} are
\code{unsigned int}.

\subsection{The settings functions}
\label{sec:cpp-settings}

The functions of \cref{tab:cpp-settings} are additions of WFDB-x to the library, declared
in \file{<wfdbx/wfdb.hpp>}: they set where files are read and written, network access and
legacy mode, for the high-level functions, the library functions and the tools alike
(\cref{chap:files,chap:behavior}). Each \code{set} function returns the previous setting,
so that it can be restored. The C function \code{setwfdb(p)} sets the input directory too
(\code{p} must name at most one directory), and \code{getwfdb()} returns it (\code{"."} if
none is set); in legacy mode they set and return the WFDB path, as in C.

\begin{xltabular}{\linewidth}{@{}lL@{}}
\caption{The settings functions}
\label{tab:cpp-settings} \\
\toprule
\tabhead Function & \tabhead Effect and result \\
\midrule
\endfirsthead
\toprule
\tabhead Function & \tabhead Effect and result \\
\midrule
\endhead
\bottomrule
\endfoot
\code{std::optional<std::string> wfdbx\_set\_input\_dir(const char *dir)} & sets the input
  directory: a local directory, a \code{\textasciitilde} path or a URL; \code{nullptr} or
  \code{""} restores the default (\code{WFDBX\_INPUT\_DIR}, or else the current
  directory). Returns the previous input directory, or \code{std::nullopt} if none was
  set \\
\code{std::optional<std::string> wfdbx\_get\_input\_dir()} & returns the input directory,
  or \code{std::nullopt} if none is set \\
\code{std::optional<std::string> wfdbx\_set\_output\_dir(const char *dir)} & sets the
  output directory: a local directory or a \code{\textasciitilde} path (a URL is refused,
  with a message); \code{nullptr} or \code{""} restores the default
  (\code{WFDBX\_OUTPUT\_DIR}, or else the current directory). Returns the previous output
  directory, or \code{std::nullopt} \\
\code{std::optional<std::string> wfdbx\_get\_output\_dir()} & returns the output directory,
  or \code{std::nullopt} if none is set \\
\code{int wfdbx\_set\_network(int on)} & turns network access on (1) or off (0); a
  negative number restores the default (on if \code{WFDBX\_NETWORK} is a nonzero number).
  Returns the previous setting, 1 or 0 \\
\code{int wfdbx\_get\_network()} & returns 1 if network access is on, 0 otherwise \\
\code{int wfdbx\_set\_legacy(int on)} & turns legacy mode on (1) or off (0), for the whole
  library. Returns the previous setting \\
\code{int wfdbx\_get\_legacy()} & returns 1 if legacy mode is on, 0 otherwise \\
\code{int wfdbx\_legacy\_args(int \&argc, char **\&argv)} & if \code{argv[1]} is
  \code{"-{}-legacy"}, removes it from the arguments, turns legacy mode on and returns 1;
  otherwise returns 0. It also installs the exit check of the tools (an exit status of 0
  becomes 1 after a write error, outside legacy mode) \\
\code{int wfdbx\_time\_needs\_freq(const char *time\_string)} & returns 1 if
  \code{strtim(time\_string)} needs a sampling frequency, 0 otherwise \\
\code{int wfdbx\_write\_errors()} & returns the number of write errors since the program
  started (output files that could not be written, flushed or closed) \\
\code{void wfdbx\_note\_write\_error()} & counts a write error, for a program that writes
  a file itself \\
\code{std::optional<std::string> wfdbx\_input\_path(const char *name)} & returns the path
  under which a program opens a file that it reads itself, where \code{wfdb\_open} finds a
  file of that name: in normal mode, a relative name under the input directory, a
  \code{\textasciitilde} name expanded, an absolute name as given; in legacy mode, the name
  as given. \code{""} and \code{"-"} are returned as given; \code{std::nullopt} for
  \code{nullptr} or a URL (also once the input directory is prefixed) \\
\code{std::optional<std::string> wfdbx\_output\_path(const char *name)} & returns the path
  under which a program creates a file that it writes itself, where \code{wfdb\_open}
  creates a file of that name: in normal mode, a relative name under the output directory,
  a \code{\textasciitilde} name expanded, an absolute name as given; in legacy mode, any
  name under the output directory, as a plain prefix. The missing directories are created
  when the output directory is prefixed. \code{""} and \code{"-"} are returned as given;
  \code{std::nullopt} for \code{nullptr} or a URL \\
\end{xltabular}

\subsection{Examples}

Reading the samples and the annotations of a record, as a C program does (the record
\code{100s} is read from the current directory):

\begin{lstlisting}[language=C++,style=example]
#include <cstdio>
#include <vector>
#include <wfdbx/ecgcodes.hpp>
#include <wfdbx/wfdb.hpp>
using namespace wfdbx;

int main() {
    WFDB_Siginfo si[2];                       // receives one WFDB_Siginfo per signal
    int nsig = isigopen("100s", si, 2);       // open 2 signals
    std::printf("%d %s %d %g\n", nsig, si[0].desc->c_str(), si[0].fmt, si[0].gain);
    // 2 MLII 212 200
    std::printf("%g\n", sampfreq(nullptr));   // 360: the frequency of the open record

    WFDB_Sample v[2];                         // receives one sample per signal
    isigsettime(strtim("0:30"));              // skip to 30 seconds
    getvec(v);
    std::printf("%d %d %g\n", v[0], v[1], aduphys(0, v[0]));    // 947 968 -0.385

    std::vector<WFDB_Sample> frames(360 * 2); // the next second, in one call
    int status;
    long long n = getvecs(frames.data(), 360, &status);
    std::printf("%lld %d\n", n, status);      // 360 2

    WFDB_Anninfo ai{"atr", WFDB_READ};
    WFDB_Annotation a;
    annopen("100s", &ai, 1);
    int count = 0;
    while (getann(0, &a) == 0) {              // 0: an annotation was read; -1: the end
        if (count < 2)
            std::printf("%s %s\n", mstimstr(a.time).c_str(), annstr(a.anntyp)->c_str());
        count++;
    }
    std::printf("%d annotations\n", count);   // 75 annotations

    wfdbquiet();                              // no messages on the standard error
    int stat = isigopen("nosuchrecord", nullptr, 0);
    std::printf("%d %s", stat, wfdberror().c_str());
    // -1 init: can't open header for record nosuchrecord
    wfdbverbose();
    wfdbquit();                               // close everything
}
\end{lstlisting}

Writing a record and an annotation file, sample by sample and annotation by annotation:

\begin{lstlisting}[language=C++,style=example]
#include <cstdio>
#include <wfdbx/ecgcodes.hpp>
#include <wfdbx/wfdb.hpp>
using namespace wfdbx;

int main() {
    WFDB_Siginfo out;
    out.fname = "ramp.dat";
    out.desc = "ramp";
    out.units = "mV";
    out.gain = 200;
    out.fmt = 16;
    out.adcres = 12;
    std::printf("%d\n", osigfopen(&out, 1));    // 1: one output signal
    setsampfreq(250);
    for (WFDB_Sample k = 0; k < 500; k++) putvec(&k);
    std::printf("%d\n", newheader("ramp"));     // 0: ramp.hea is written
    WFDB_Anninfo ai{"qrs", WFDB_WRITE};
    annopen("ramp", &ai, 1);
    for (WFDB_Time t : {100, 350}) {
        WFDB_Annotation a;
        a.time = t;
        a.anntyp = NORMAL;
        putann(0, &a);
    }
    wfdbquit();                                 // completes ramp.dat and ramp.qrs
}
\end{lstlisting}

% ==========================================================================================
\section{The command-line tools in C++}
\label{sec:cpp-tools}

The 33 command-line tools (\code{rdsamp}, \code{rdann}, \code{wrsamp}, \code{xform},
\ldots) are separate programs, built in \file{cpp/build/bin/} (or the \file{bin} directory
of the installation), one per file of \file{cpp/apps}, and run by their names. Each
follows its C program line by line. The tools, their options, their output and their exit
status are those of the C tools, and are documented in \cref{chap:tools}.

\begin{lstlisting}[style=shell]
export PATH=$PWD/cpp/build/bin:$PATH
export WFDBX_INPUT_DIR=tests/data     # run from the root of the repository
rdsamp -r 100s -t 10 -p     # the first 10 seconds of 100s, in physical units
rdann -r 100s -a atr        # its annotations
rdsamp --legacy -r 100s     # legacy mode: --legacy as the first argument
\end{lstlisting}

There is no launcher: a tool is run by its own name, and it has no \opt{-{}-version}
option, like the C tools. When an output file or the standard output cannot be written (a
full disk, for example), a tool reports it (``PROGRAM: write error on standard output'')
and exits with status 1; with \opt{-{}-legacy}, it keeps the original exit status.

\paragraph{Environment variables.} The tools, like the rest of the library, use the
environment variables of \cref{tab:cpp-environment}, also for the files that they read or
create themselves: the \opt{-i} file of \code{wrsamp}, \code{edf2mit}, \code{wav2mit},
\code{a2m} and \code{ad2m}, the input files of \code{ahaecg2mit} and the \opt{-S} script
of \code{xform} are read from the input directory, and the WAV file of \code{mit2wav},
the EDF file of \code{mit2edf}, the master header of \code{snip -m} and the \file{.des},
\file{.ecg}, \file{.key} and \file{.txt} files of \code{parsescp} are created in the
output directory, as the records they read and write are (with
\code{wfdbx\_input\_path} and \code{wfdbx\_output\_path}, \cref{tab:cpp-settings}). Such
a file cannot be read from a URL.

\begin{table}[htbp]
\centering
\caption{The environment variables used by the C++ library}
\label{tab:cpp-environment}
\begin{tabularx}{\linewidth}{@{}lL@{}}
\toprule
\tabhead Variable & \tabhead Meaning \\
\midrule
\code{WFDBX\_INPUT\_DIR} & the input directory (a directory, a \code{\textasciitilde}
  path or a URL), when \code{wfdbx\_set\_input\_dir} does not set one \\
\code{WFDBX\_OUTPUT\_DIR} & the output directory, when \code{wfdbx\_set\_output\_dir} does
  not set one \\
\code{WFDBX\_NETWORK} & network access is on if it is a nonzero number, when
  \code{wfdbx\_set\_network} does not set it \\
\code{WFDBANNSORT} & \code{0} to keep annotation files in the order written (no sorting
  when they are closed) \\
\code{WFDBGVMODE} & the default getvec mode (\code{setgvmode} with a negative mode) \\
\code{WFDBCAL} & the calibration file read by \code{calopen(nullptr)} (default
  \file{wfdbcal}) \\
\code{WFDB} & the WFDB path, used in legacy mode only (outside legacy mode, it is
  ignored, with a message) \\
\code{HOME} (\code{USERPROFILE} on Windows) & the home directory, for
  \code{\textasciitilde} paths \\
\bottomrule
\end{tabularx}
\end{table}

Tools that write signal parameters require them as options: \code{wrsamp} (with
\opt{-o}) requires \opt{-F} (frequency), \opt{-G} (gain), \opt{-R} (ADC resolution) and,
for calibrated signals, \opt{-U} (units); \code{wav2mit} requires \opt{-G}, and \opt{-U}
unless the gain is 0; \code{ad2m} and \code{ahaecg2mit} require \opt{-F}, \opt{-G},
\opt{-R}, and \opt{-U} unless the gain is 0; \code{a2m}, \code{rdann}, \code{ann2rr},
\code{sortann} and \code{m2a} require \opt{-F} when the record's header gives no
frequency (\cref{chap:tools}).

% ==========================================================================================
\section{Testing}
\label{sec:cpp-testing}

The unit tests of the C++ implementation are in \file{cpp/tests}, built with the library
(unless \code{WFDBX\_BUILD\_TESTS} is \code{OFF}) as the program
\file{cpp/build/wfdbx\_tests}. From the root of the repository, after building:

\begin{lstlisting}[style=shell]
ctest --test-dir cpp/build                          # the unit tests
cpp/build/wfdbx_tests                               # the same, one line per test
cpp/build/wfdbx_tests metadata_shared_cases         # selected tests, by name
\end{lstlisting}

The unit tests cover the WFDB library functions, the high-level functions, the metadata
functions (with the cases shared by the seven languages, in \file{tests/metadata}),
\code{check\_record} and the FLAC codec. They run in temporary copies of the test records
of \file{tests/data} (the environment variable \code{WFDBX\_TEST\_DATA} names another
directory). The tests of network access use a local HTTP server on 127.0.0.1 only, and
are skipped without \cmd{curl}. The tests of \code{check\_record} write the corpus of
\file{python/tests/check\_corpus.py} with \cmd{python3} (or the program named by the
environment variable \code{WFDBX\_PYTHON}) and the Python package of \file{python/src},
and compare the reports with its \file{expected.txt}; without Python they are skipped.
\cmd{cpp/build/wfdbx\_check\_corpus DIR} prints the reports for a corpus in \file{DIR}.

The cross-language conformance suite runs every tool and compares its results, byte for
byte, with those of the original C programs (1135 cases), and, with \opt{-{}-legacy}, in
legacy mode (883 cases). It runs the tools of \file{cpp/build/bin}:

\begin{lstlisting}[style=shell]
PYTHONPATH=$PWD/python/src python3 tests/conformance/harness.py check --impl cpp -j8
PYTHONPATH=$PWD/python/src python3 tests/conformance/harness.py check --impl cpp -j8 --legacy
\end{lstlisting}

% The Rust chapter of the WFDB-x user's manual. It follows the sections of the Python
% chapter (chapters/python.tex), in the same order. Each \apientry of the section "The
% high-level API" is a public free function of the crate root defined in
% rust/src/record.rs, rust/src/check.rs or rust/src/metadata.rs (each also a method of
% Wfdb, with the same parameters after &mut self); its signature listing is copied from
% that definition, and its params name every parameter of the definition, in order. The
% examples of style "example" were compiled and run (each as the body of main(), after
% the lines given in the conventions of the section, in a copy of tests/data);
% check_manual.py does not run them.
\chapter{Rust}
\label{chap:rust}
\chapterlanguage{rust}

The Rust implementation of WFDB-x is the crate \code{wfdbx}, in the directory \file{rust}
of the repository: a library and the 33 command-line tools, each a separate program of
the crate. It is a complete Rust implementation of the WFDB Software Package: it uses only
the Rust standard library (no other crate), contains no \code{unsafe} code, and never
calls the C library.

It offers WFDB-x in three ways:
\begin{itemize}
\item \textbf{The high-level API} (\cref{sec:rust-api}): functions that read and write
  whole records, headers, annotation files and header metadata as \code{Vec}s and
  structures, in one call, with options given as fields of a structure, and that return a
  \code{Result} whose error type is \code{wfdbx::record::Error}. They read signal files in
  large blocks, and are the fastest way to read a record. They are an addition of WFDB-x:
  the C package has no equivalent. They are the recommended interface for new code and
  analysis.
\item \textbf{The WFDB library functions} under their C names (\code{isigopen},
  \code{getvec}, \code{getann}, \code{putann}, \ldots), with the behavior of the C library
  (\cref{sec:rust-library}). They are documented once, for all the languages, in
  \cref{chap:library}.
\item \textbf{The command-line tools} (\code{rdsamp}, \code{rdann}, \code{wrsamp},
  \code{xform} and 29 others), with the options, output and exit status of the C tools,
  one program per tool, run by its name (\cref{sec:rust-tools}). They are documented in
  \cref{chap:tools}.
\end{itemize}
The high-level functions return exactly the samples and annotations that the library
functions (and the C library) give.

% ==========================================================================================
\section{Installation and requirements}
\label{sec:rust-install}

\subsection{Requirements}

WFDB-x for Rust needs Rust 1.70 or later (the stable toolchain, edition 2021; the crate
declares \code{rust-version = "1.70"}) and Cargo, on Linux, macOS or Windows. It has no
dependency beyond the Rust standard library: no WFDB library and no libFLAC (the FLAC
formats 508, 516 and 524 are read and written by WFDB-x's own codec). Network access
(\cref{sec:network}) runs the program \cmd{curl}, which must then be installed. The tests
of \code{check\_record} also run \cmd{python3}, with NumPy (\cref{sec:rust-testing}).

\subsection{Using the library in a Cargo project}

A Cargo project uses the crate as a dependency, from the repository or from a copy of it,
in its \file{Cargo.toml}:

\begin{lstlisting}[language=TOML]
[dependencies]
wfdbx = { git = "https://github.com/alphanumericslab/wfdb-x" }
# or, from a copy of the repository:
wfdbx = { path = "../wfdb-x/rust" }
\end{lstlisting}

The crate is \code{wfdbx}, version \code{1.2.1}: its library under the GNU Lesser
General Public License, version 2.1 or later, and its tools (the binaries of
\file{src/bin}) under the GNU General Public License, version 3 or later
(\cref{sec:license}). Programs that use it are built in release mode for speed
(\cmd{cargo build -{}-release}, \cmd{cargo run -{}-release}).

\subsection{Building and installing the tools}

From the root of a copy of the repository, \cmd{cargo install} builds the tools in
release mode and installs them in \file{\textasciitilde/.cargo/bin}; \cmd{cargo build
-{}-release}, in the directory \file{rust}, builds the library and the tools in
\file{rust/target/release/} without installing them:

\begin{lstlisting}[style=shell]
cargo install --path rust                                # the 33 tools, in ~/.cargo/bin
cargo build --release --manifest-path rust/Cargo.toml    # in rust/target/release/
cargo build --release --manifest-path rust/Cargo.toml --bin rdsamp   # one tool
\end{lstlisting}

The tools, like the C tools, have no \opt{-{}-version} option. \cmd{cargo doc -{}-open}, in
the directory \file{rust}, builds the documentation of the crate (its doc comments, with
examples) and opens it.

\subsection{The sample record}
\label{sec:rust-resources}

The repository includes the sample record \code{100s}, the sample record of the WFDB
Software Package: the first minute of record 100 of the MIT-BIH Arrhythmia Database, two
ECG signals (\code{MLII} and \code{V5}) sampled at 360\,Hz, with the reference beat
annotations (annotator \code{atr}). Its files (\file{100s.hea}, \file{100s.dat},
\file{100s.atr}) are in the directory \file{tests/data} of a copy of the repository, with
the other test records of WFDB-x. The crate does not install them: programs name that
directory as their input directory (\code{wfdbx\_set\_input\_dir("tests/data")}, or the
option \code{input\_dir}). Work that uses the MIT-BIH Arrhythmia Database cites it as
described on PhysioNet.

% ==========================================================================================
\section{Quick start}
\label{sec:rust-quickstart}

This program reads the header, ten seconds of the signals and the annotations of the
sample record \code{100s}, from the directory \file{tests/data} of the repository (the
input directory of each call), and computes the mean heart rate. It then writes the first
ten seconds as a new record, \code{first10}, in the current directory, with a metadata
entry in its header.

\begin{lstlisting}[language=Rust,style=example]
use wfdbx::{check_record, isqrs, read_annotations_with, read_header_with, read_metadata,
            read_record, set_metadata, write_record, ReadOptions, WriteOptions};

fn main() -> Result<(), Box<dyn std::error::Error>> {
    // The sample record 100s, in the directory tests/data of the repository.
    let sample = ReadOptions { input_dir: Some("tests/data".into()), ..Default::default() };

    // The header only: no samples are read.
    let hdr = read_header_with("100s", &sample)?;
    println!("{} {} {}", hdr.nsig, hdr.fs, hdr.nsamples); // 2 360 21600

    // Seconds 10 to 20 of both signals, frame by frame, in ADC units.
    let opt = ReadOptions {
        from: Some("0:10".into()),
        to: Some("0:20".into()),
        ..sample.clone()
    };
    let rec = read_record("100s", &opt)?;
    println!("{} {}", rec.nsamples(), rec.nsig); // 3600 2
    println!("{} {}", rec.at(0, 0), rec.physical(0, 0)?); // 946 -0.39

    // The beat annotations, and the mean heart rate.
    let ann = read_annotations_with("100s", "atr", None, &sample)?;
    let beats: Vec<i64> = (0..ann.len())
        .filter(|&i| isqrs(ann.codes[i]) != 0)
        .map(|i| ann.samples[i])
        .collect();
    let rr = (beats[beats.len() - 1] - beats[0]) as f64 / ann.fs / (beats.len() - 1) as f64;
    println!("{} beats, {:.1} beats per minute", beats.len(), 60.0 / rr);
    // 74 beats, 73.9 beats per minute

    // Write the first 10 seconds as the record "first10" (first10.hea and first10.dat,
    // in the current directory), with the signal specifications of 100s.
    let opt = ReadOptions { to: Some("0:10".into()), ..sample.clone() };
    let x = read_record("100s", &opt)?;
    let w = WriteOptions {
        gains: vec![200.0],
        baselines: vec![1024],
        adczeros: vec![1024],
        adcres: vec![11],
        units: vec![Some("mV".into())],
        names: x.names(),
        metadata: vec![("Source".into(), "MIT-BIH record 100".into())],
        ..Default::default()
    };
    write_record("first10", &x.samples, x.nsig, x.fs, &w)?;

    // Header metadata: "# Key: value" lines of the header (keys and values are bytes).
    set_metadata("first10", "Age", "69")?;
    for (key, value) in read_metadata("first10")? {
        println!("{} = {}", String::from_utf8_lossy(&key), String::from_utf8_lossy(&value));
    }
    // Source = MIT-BIH record 100
    // Age = 69
    println!("{}", check_record("first10", &[])?); // no problems found
    Ok(())
}
\end{lstlisting}

Run from the root of the repository, as the program \file{src/main.rs} of a Cargo project
that depends on \code{wfdbx} (\cmd{cargo run -{}-release}). The header \file{100s.hea}
gives no units: the program writes \code{"mV"} explicitly, since WFDB-x never assumes
units, a gain, an ADC resolution or a sampling frequency that a record does not give
(\cref{chap:behavior}).

% ==========================================================================================
\section{The high-level API}
\label{sec:rust-api}

The high-level API is made of the functions of this section, defined in the modules
\code{wfdbx::record}, \code{wfdbx::check} and \code{wfdbx::metadata} and available at the
root of the crate (\code{wfdbx::read\_record}, \ldots), of the structures of their options
(\cref{sec:rust-options}) and of their results (\cref{sec:rust-results}), and of the error
type \code{wfdbx::record::Error} (\cref{sec:rust-warnings}). The doc comments of the
crate (\cmd{cargo doc}) document each of them.

\subsection{Conventions}
\label{sec:rust-conventions}

\paragraph{Free functions and methods.} The state of the library (the open record, the
annotators, the input and output directories, network access, legacy mode) is held by a
value of the type \code{Wfdb}. Each high-level function exists twice: as a method of
\code{Wfdb}, such as \code{w.read\_record("100s", \&opt)}, which uses the state of
\code{w}, and as a free function of the crate with the same name and the same
parameters, such as \code{wfdbx::read\_record("100s", \&opt)}, which uses the default
\code{Wfdb} value of the calling thread (one per thread). The entries of this section
document the free functions; the methods take \code{\&mut self} first, and then the same
parameters, with the same types and results. Separate \code{Wfdb} values share no state
(except what belongs to the process: the count of write errors, the cache of fetched URLs
and the messages printed once per process), so several
records can be read at once, in one thread or in several; a \code{Wfdb} is \code{Send},
and is used by one thread at a time.

\begin{lstlisting}[language=Rust,style=example]
use wfdbx::Wfdb;

fn main() -> Result<(), wfdbx::record::Error> {
    let mut w = Wfdb::new(); // a library state of its own
    w.wfdbx_set_input_dir("tests/data");
    let a = w.read_record("100s", &Default::default())?; // on w
    let b = wfdbx::read_record("tests/data/100s", &Default::default())?; // on the default value
    println!("{}", a.samples == b.samples); // true
    Ok(())
}
\end{lstlisting}

\paragraph{Options.} The options of a function are the fields of a structure
(\code{ReadOptions}, \code{WriteOptions}, \code{AnnWriteOptions}, \code{MetadataOptions},
\cref{sec:rust-options}), passed by reference as the last argument. Each implements
\code{Default}, so that the fields that are not given keep their defaults:
\code{ReadOptions \{ highres: true, ..Default::default() \}}. Several functions have a
short form with no options or with a few arguments (\code{read\_header},
\code{read\_annotations}, \code{write\_annotations}, \code{check\_record}, the metadata
functions), and a form whose name ends in \code{\_with}, which takes the options
structure. An option that may be absent is an \code{Option}; \code{None} means the
default or the current setting.

\paragraph{Record names and directories.} A record is named as in the other parts of
WFDB-x (\cref{chap:records,sec:record-paths}), without the \file{.hea} extension (the
readers, and \code{write\_annotations}, also accept a trailing \code{".hea"};
\code{write\_record} does not). A record name is one of:
\begin{itemize}
\item a full path, absolute (\code{"/data/mitdb/100"}) or in the home directory
  (\code{"\textasciitilde/Documents/100s"}), used as given;
\item a relative path (\code{"100s"}, \code{"data/100s"}), read from the input directory
  if one is set, otherwise from the current directory, and written to the output
  directory if one is set, otherwise to the current directory (missing directories are
  created);
\item a URL (\code{"https://physionet.org/files/mitdb/1.0.0/100"}), read only if network
  access is on; a URL cannot be written.
\end{itemize}
The input and output directories are set with \code{wfdbx\_set\_input\_dir} and
\code{wfdbx\_set\_output\_dir} (\cref{sec:rust-settings}) or with the environment
variables \code{WFDBX\_INPUT\_DIR} and \code{WFDBX\_OUTPUT\_DIR}. The readers take the
optional fields \code{input\_dir} and \code{network} of a \code{ReadOptions} (or a
\code{MetadataOptions}), and the writers a \code{dir}: they apply to that one call, and
the previous settings are restored when it returns, also when it fails.
\code{input\_dir} applies to relative names only, not to full paths or URLs. A writer
given a name with a path creates the files in that directory, and the header names the
record and its signal file without the path (\code{rec}, \code{rec.dat}), so that the
record opens under any path that leads to it; in legacy mode, as in the original, the
header gives the name with its path. With \code{dir}, the name must be a plain name
(\code{"rec"}): the files are created in \code{dir}, and the record opens from that
directory under its plain name. The signal files that a header names are read from the
header's directory.

\paragraph{Times.} The fields \code{from} and \code{to} of \code{ReadOptions}, and the
arguments \code{from} and \code{to} of \code{read\_annotations}, are
\code{Option<TimeSpec>} values (\cref{tab:rust-timespec}): a sample number or a WFDB time
string, as \code{strtim} reads it (\code{"0:10"}, 10 seconds; \code{"1:02:03.5"};
\code{"s3600"}, sample 3600; \code{"e"}, the end of the record; or a time of day in
brackets, \code{"[14:00:00]"}). A \code{TimeSpec} converts from integers and strings, so
\code{Some(3600.into())} and \code{Some("0:10".into())} both work; \code{None} means that
no time is given. A time in seconds needs the sampling frequency; when the record does
not give it, such a time is an error.

\paragraph{Samples.} Samples are \code{i32} values (C's \code{WFDB\_Sample}), in ADC
units, stored frame by frame in one \code{Vec<i32>}: \code{samples[i * nsig + s]} is
sample \code{i} of signal \code{s}. Invalid samples are \code{WFDB\_INVALID\_SAMPLE}
($-32768$), and NaN in physical units. Sample numbers and times are \code{i64}.

\paragraph{Per-signal and per-annotation values.} The per-signal fields of
\code{WriteOptions} (\code{gains}, \code{units}, \ldots) hold one value for all the
signals, or one value per signal; the per-annotation fields of \code{Annotations}
(\code{symbols}, \code{chans}, \ldots) are empty or hold one value per annotation. Any
other length is an error, reported before anything is written.

\paragraph{Signal parameters.} No sampling frequency, gain, ADC resolution or units is
ever assumed. Where one is needed and the record does not give it, it must be given;
otherwise the function returns an error. A record whose header gives no sampling
frequency has \code{fs} 0, a signal without an ADC resolution has \code{adcres} 0, and a
signal without units has \code{units} \code{None}.

\paragraph{Text.} Text read from files (signal descriptions, units, info strings,
metadata, aux strings, mnemonics) is given as bytes, \code{Vec<u8>}, since WFDB files are
not UTF-8: every byte sequence round-trips unchanged. Arguments accept \code{\&str},
\code{\&[u8]}, \code{b"..."} or \code{Vec<u8>} (\code{impl AsRef<[u8]>}), and a
\code{Vec<u8>} field is set from a \code{\&str} with \code{.into()}.
\code{String::from\_utf8\_lossy} converts bytes for printing. The fields
\code{channels} and \code{input\_dir} of the options and \code{base\_time} of the results
are \code{String}s, and the directories \code{dir} are \code{PathBuf}s.

\paragraph{Legacy mode.} With \code{legacy} set to \code{true}, a function behaves as in
the original WFDB package, with its assumed signal parameters and its search for files
(\cref{chap:behavior}); the library's legacy mode is on for the duration of the call
only. With \code{false} (the default), the current setting of \code{wfdbx\_set\_legacy}
applies, which is off unless it was turned on.

\paragraph{State.} The functions that read and write records and annotation files call
\code{wfdbquit()} on the \code{Wfdb} value they use, before and after their work, also
when they fail, which closes anything opened with the WFDB library functions on that
value. The library's settings (legacy mode, getvec mode, input and output directories,
network access) are the same after a call as before it. \code{read\_record},
\code{read\_header}, \code{read\_annotations}, \code{write\_record},
\code{write\_record\_physical}, \code{write\_annotations}, their \code{\_with} forms and
\code{check\_record} also leave the time resolution of output annotators
(\code{getafreq()}) and the getvec mode (\code{getgvmode()}) as they were, also after
reading an EDF record, whose opening sets both (as in C; the library functions keep that
behavior). The metadata functions do not
call \code{wfdbquit()}: they open only the header, and leave the library's open files as
they are. Results are returned by value.

\paragraph{The examples} of this section are the body of the function
\code{fn main() -> Result<(), Box<dyn std::error::Error>{}>}, followed by \code{Ok(())},
compiled after these lines, and run with the free functions in a copy of the test
records of \file{tests/data} (among them \code{100s}), which is the current directory and
holds the records that the examples write:
\begin{lstlisting}[language=Rust]
use wfdbx::*;

/// Text (bytes) as a String, for printing.
fn s(b: &[u8]) -> String {
    String::from_utf8_lossy(b).into_owned()
}
\end{lstlisting}

% ------------------------------------------------------------------------------------------
\subsection{Reading records}
\label{sec:rust-reading}

\apientry{read\_record}{Read the signals of a record}
\begin{lstlisting}[language=Rust,style=signature]
pub fn read_record(record: impl AsRef<[u8]>, opt: &ReadOptions) -> Result<Record, Error>
\end{lstlisting}
Reads the signals of a record from \code{opt.from} up to \code{opt.to} (exclusive), in
ADC units, frame by frame, together with the description of the record. The samples are
those that \code{getvec} returns (\cref{chap:library}), read in large blocks with
\code{getvecs}. If the record ends before \code{opt.to}, the samples up to its end are
returned. When the read reaches the end of the record, the samples are compared with the
checksums of the header (\cref{sec:rust-warnings}). The getvec mode (\code{getgvmode()})
and the time resolution of output annotators (\code{getafreq()}) are the same after the
call as before it.
\begin{params}
\item[record] Record name, without the \file{.hea} extension (a trailing \code{".hea"} is
  also accepted): a full path, absolute (\code{"/data/mitdb/100"}) or in the home
  directory (\code{"\textasciitilde/Documents/100s"}), used as given; a relative path
  (\code{"100s"}, \code{"data/100s"}), read from the input directory if one is set,
  otherwise from the current directory; or a URL
  (\code{"https://physionet.org/files/mitdb/1.0.0/100"}), read only if network access is
  on.
\item[opt] The options (\cref{tab:rust-readoptions}): the signals to read
  (\code{channels}), the interval (\code{from}, \code{to}), \code{highres}, \code{pad},
  \code{checksum\_warning}, \code{legacy}, and the input directory (under which relative
  record names are read; it does not apply to full paths or URLs) and network access of
  the call (\code{input\_dir}, \code{network}). \code{\&ReadOptions::default()}: every
  signal of the whole record, with the current settings.
\end{params}
\returns{A \code{Record} (\cref{tab:rust-record}), with the samples in ADC units.}
\raises{An \code{Error} if the record cannot be read or has no signals (``can't read record
REC''), or its signals cannot be opened (``can't open signals of record REC''); if a
requested signal does not exist (``record REC has no signal C''); if \code{opt.from} is a
negative sample number (``read\_record: from must not be negative'') or cannot be reached
(``can't seek to sample N of record REC''); if a time string needs the sampling frequency
and the header gives none (``sampling frequency of record REC is not specified''); or if a
signal file is shorter than its header says, or cannot be read or decoded (``can't read
signal data of record REC (getvec status S)''). Data that do not match the checksums of
the header are returned, with a warning (see \code{checksum\_warning}).}
\begin{lstlisting}[language=Rust,style=example]
// Seconds 10 to 20 of lead MLII, in physical units.
let opt = ReadOptions {
    channels: vec!["MLII".into()],
    from: Some("0:10".into()),
    to: Some("0:20".into()),
    ..Default::default()
};
let rec = read_record("100s", &opt)?;
let mlii: Vec<f64> = rec.physical_signal(0)?;
println!("{} {} {} {}", mlii.len(), rec.fs, rec.sampfrom, s(&rec.names()[0]));
// 3600 360 3600 MLII

// The first 3 samples of both signals, in ADC units, V5 first.
let first = ReadOptions {
    channels: vec!["1".into(), "0".into()],
    to: Some(3.into()),
    ..Default::default()
};
let rec = read_record("100s", &first)?;
for i in 0..rec.nsamples() {
    println!("{} {}", rec.at(i, 0), rec.at(i, 1)); // 1011 995, 3 times
}
println!("{} {}", rec.siginfo[0].gain, rec.siginfo[0].baseline); // 200 1024

// The whole record: the checksums are verified.
let rec = read_record("100s", &ReadOptions::default())?;
println!("{} {}", rec.nsamples(), rec.checksum_ok); // 21600 true
\end{lstlisting}
A multi-frequency record (here the test record \code{mfreq}, whose signals have 1, 2
and 4 samples per frame at a frame frequency of 100\,Hz) is read at its frame frequency,
or at its highest sampling frequency:
\begin{lstlisting}[language=Rust,style=example]
let rec = read_record("mfreq", &ReadOptions::default())?;
println!("{} {}", rec.fs, rec.nsamples()); // 100 600
let opt = ReadOptions { highres: true, ..Default::default() };
let rec = read_record("mfreq", &opt)?;
println!("{} {}", rec.fs, rec.nsamples()); // 400 2400
let spf: Vec<i32> = rec.siginfo.iter().map(|si| si.spf).collect();
println!("{spf:?}"); // [1, 2, 4]
\end{lstlisting}

\apientry{read\_header}{Read the header of a record}
\begin{lstlisting}[language=Rust,style=signature]
pub fn read_header(record: impl AsRef<[u8]>, legacy: bool) -> Result<Header, Error>
\end{lstlisting}
Reads the header of a record, without opening its signal files. The result describes the
record and each of its signals, and holds the info strings and the metadata entries of the
header (from the \file{.hea} file only, as \code{read\_metadata} reads them, with its
warning for a key given more than once). For an EDF record, which has no \file{.hea}
file, and for a header read from the standard input (\code{"-"}), the metadata is empty,
without an error. For a multi-segment record, the master header is read, and
\code{nsamples} is the length of the whole record.
\begin{params}
\item[record] Record name, as for \code{read\_record}: a full path (absolute or
  \code{\textasciitilde}), used as given; a relative path, read from the input directory
  if one is set, otherwise from the current directory; or a URL, read only if network
  access is on.
\item[legacy] \code{true} to read the header as the original WFDB package does; normally
  \code{false}.
\end{params}
\returns{A \code{Header} (\cref{tab:rust-header}).}
\raises{An \code{Error} if the header cannot be read (``can't read header for record
REC'').}
\begin{lstlisting}[language=Rust,style=example]
let hdr = read_header("100s", false)?;
println!("{} {} {} {}", s(&hdr.name), hdr.fs, hdr.nsig, hdr.nsamples); // 100s 360 2 21600
for si in &hdr.siginfo {
    let desc = si.desc.as_deref().unwrap_or(b"");
    let fname = si.fname.as_deref().unwrap_or(b"");
    println!("{} {} {} {} {} {:?}", s(desc), s(fname), si.fmt, si.gain, si.adcres, si.units);
}
// MLII 100s.dat 212 200 11 None
// V5 100s.dat 212 200 11 None
println!("[{}]", s(&hdr.comments[0])); // [ 69 M 1085 1629 x1]
println!("{}", hdr.base_time.is_none()); // true: the header gives no start time
let (key, value) = &hdr.metadata[0];
println!("[{}] [{}]", s(key), s(value));
// [Produced by xform from record 100, beginning at 0] [0]
\end{lstlisting}

\apientry{read\_header\_with}{Read the header of a record, with options}
\begin{lstlisting}[language=Rust,style=signature]
pub fn read_header_with(record: impl AsRef<[u8]>, opt: &ReadOptions) -> Result<Header, Error>
\end{lstlisting}
As \code{read\_header}, with the options \code{input\_dir}, \code{network} and
\code{legacy} of \code{opt}, for this call only. The other fields of \code{opt} are not
used.
\begin{params}
\item[record] Record name, as for \code{read\_header}.
\item[opt] The options (\cref{tab:rust-readoptions}); only \code{input\_dir} (optional:
  the input directory under which relative record names are read; it does not apply to
  full paths or URLs), \code{network} and \code{legacy} are used.
\end{params}
\returns{A \code{Header} (\cref{tab:rust-header}).}
\raises{As \code{read\_header}.}
\begin{lstlisting}[language=Rust,style=example]
// A copy of the header 100s.hea in the directory "in".
std::fs::create_dir_all("in")?;
std::fs::copy("100s.hea", "in/100s.hea")?;
let opt = ReadOptions { input_dir: Some("in".into()), ..Default::default() };
let hdr = read_header_with("100s", &opt)?; // reads in/100s.hea
println!("{} {}", hdr.nsig, hdr.fs); // 2 360
println!("{}", read_header("100s", false).is_ok()); // true: the input directory is restored
\end{lstlisting}

\apientry{read\_annotations}{Read an annotation file}
\begin{lstlisting}[language=Rust,style=signature]
pub fn read_annotations(
    record: impl AsRef<[u8]>,
    annotator: impl AsRef<[u8]>,
    from: Option<TimeSpec>,
    to: Option<TimeSpec>,
    fs: Option<f64>,
    legacy: bool,
) -> Result<Annotations, Error>
\end{lstlisting}
Reads the annotations of an annotation file, in file order, as parallel vectors. The
sampling frequency, needed to convert times in seconds, is that of the record's header,
or, for a record without a header or whose header gives none, the argument \code{fs}, as
the option \opt{-F} of \code{rdann}.
\begin{params}
\item[record] Record name, as for \code{read\_record}: a full path (absolute or
  \code{\textasciitilde}), used as given; a relative path, read from the input directory
  if one is set, otherwise from the current directory; or a URL, read only if network
  access is on Its header, if there is one, gives the sampling frequency.
\item[annotator] Annotator name: the suffix of the annotation file. \code{"atr"} reads the
  file \file{RECORD.atr}.
\item[from] The first time to read: a sample number or a time string. \code{None}: the
  beginning of the file.
\item[to] The time at which to stop: annotations at or after it are not read.
  \code{None}: the end of the file.
\item[fs] The sampling frequency, in Hz ($> 0$), to use when the record has no header or
  its header gives none; it is ignored otherwise. \code{None}: none.
\item[legacy] \code{true} to behave as the original WFDB package (which assumes 250\,Hz
  for a record without a frequency); normally \code{false}.
\end{params}
\returns{An \code{Annotations} (\cref{tab:rust-annotations}). Its \code{fs} is the
frequency in which the times are counted: the one written in the annotation file, if
there is one, or else the sampling frequency; it is 0 if none of them is known.}
\raises{An \code{Error} if \code{fs} is given and is not $> 0$ (``fs must be > 0''); if
the annotation file cannot be opened (``can't read annotator ANN for record REC''); if a
time string needs the sampling frequency and none is known (``sampling frequency of
record REC is not specified''); if \code{from} cannot be reached (``can't seek in
annotator ANN''); or if the file ends early or cannot be read before \code{to} is reached
(``can't read annotation file ANN of record REC (getann status S)'', where S is the
status that \code{getann} returned, such as $-3$ for an unexpected end of file).}
\begin{lstlisting}[language=Rust,style=example]
let ann = read_annotations("100s", "atr", Some("0:10".into()), Some("0:30".into()), None, false)?;
println!("{} {}", ann.len(), ann.fs); // 24 360
println!("{} {}", ann.samples[0], s(&ann.symbols[0])); // 3862 N

// The whole file: the first annotation is a rhythm change, with an aux string.
let ann = read_annotations("100s", "atr", None, None, None, false)?;
let aux = ann.aux[0].as_deref().unwrap_or(b"");
println!("{} {} {}", s(&ann.symbols[0]), ann.codes[0] == RHYTHM, s(aux)); // + true (N
\end{lstlisting}
For a record without a header, \code{fs} gives the frequency in which times in seconds
are converted (the library reports on the standard error that it cannot open the
header):
\begin{lstlisting}[language=Rust,style=example]
let beats = Annotations {
    samples: vec![100, 460, 820],
    symbols: vec!["N".into(), "N".into(), "N".into()],
    ..Default::default()
};
write_annotations("beats", "qrs", &beats, 0.0, None)?; // no header
let ann = read_annotations("beats", "qrs", None, Some("0:02".into()), Some(360.0), false)?;
println!("{} {}", ann.len(), ann.fs); // 2 360
\end{lstlisting}

\apientry{read\_annotations\_with}{Read an annotation file, with options}
\begin{lstlisting}[language=Rust,style=signature]
pub fn read_annotations_with(
    record: impl AsRef<[u8]>,
    annotator: impl AsRef<[u8]>,
    fs: Option<f64>,
    opt: &ReadOptions,
) -> Result<Annotations, Error>
\end{lstlisting}
As \code{read\_annotations}, with the options \code{from}, \code{to}, \code{input\_dir},
\code{network} and \code{legacy} of \code{opt} (the last three for this call only). The
other fields of \code{opt} are not used.
\begin{params}
\item[record] Record name, as for \code{read\_annotations}.
\item[annotator] Annotator name, as for \code{read\_annotations}.
\item[fs] The sampling frequency to use when the record has no header or its header
  gives none, as for \code{read\_annotations}.
\item[opt] The options (\cref{tab:rust-readoptions}); only \code{from}, \code{to},
  \code{input\_dir} (optional: the input directory under which relative record names are
  read; it does not apply to full paths or URLs), \code{network} and \code{legacy} are
  used.
\end{params}
\returns{An \code{Annotations} (\cref{tab:rust-annotations}).}
\raises{As \code{read\_annotations}.}
\begin{lstlisting}[language=Rust,style=example]
// Seconds 10 to 30 of the annotations of 100s, read from the directory "in".
std::fs::create_dir_all("in")?;
for f in ["100s.hea", "100s.atr"] {
    std::fs::copy(f, format!("in/{f}"))?;
}
let opt = ReadOptions {
    from: Some("0:10".into()),
    to: Some("0:30".into()),
    input_dir: Some("in".into()),
    ..Default::default()
};
let ann = read_annotations_with("100s", "atr", None, &opt)?;
println!("{} {}", ann.len(), ann.samples[0]); // 24 3862
\end{lstlisting}

% ------------------------------------------------------------------------------------------
\subsection{Writing records}
\label{sec:rust-writing}

\apientry{write\_record}{Write a header and a signal file}
\begin{lstlisting}[language=Rust,style=signature]
pub fn write_record(
    record: impl AsRef<[u8]>,
    samples: &[i32],
    nsig: usize,
    fs: f64,
    opt: &WriteOptions,
) -> Result<(), Error>
\end{lstlisting}
Writes samples in ADC units as a new record: the header \file{RECORD.hea} and one signal
file, \file{RECORD.dat}, which holds all the signals (so they share one format). Existing
files of these names are overwritten. The sampling frequency, the gain and ADC resolution
of each signal, and the units of each calibrated signal must be given: there are no
defaults. Every check except those of the writing itself comes before any file is
created; whatever happens, the output signals are closed before the function returns,
and the getvec mode and the time resolution of output annotators are left as they were.
\begin{params}
\item[record] Record name (not empty), without the \file{.hea} extension: a plain name
  (\code{"rec"}) or a relative path (\code{"data/rec"}), created in the output directory
  if one is set, otherwise in the current directory, missing directories included; or a
  full path, absolute or in the home directory (\code{"\textasciitilde/Documents/rec"}),
  used as given. With a path, the files are created in that directory, and the header
  names the record and its signal file without the path (\code{rec}, \code{rec.dat}), so
  that the record opens under any path that leads to it (in legacy mode, as in the
  original, the header gives the name with its path). A URL cannot be written.
\item[samples] The samples, in ADC units, frame by frame: \code{nsig} values per frame,
  or, for a multi-frequency record, the sum of \code{opt.samps\_per\_frame} (the samples
  of signal 0, then those of signal 1, and so on). \code{WFDB\_INVALID\_SAMPLE} marks
  invalid samples. The number of samples must be a multiple of the frame size.
\item[nsig] The number of signals ($> 0$).
\item[fs] The sampling frequency, in Hz: a finite number $> 0$ (for a multi-frequency
  record, the frame frequency).
\item[opt] The options (\cref{tab:rust-writeoptions}): the format, the signal
  specifications (\code{gains}, \code{adcres} and \code{units} are required), the start
  time, the comments and metadata, the directory \code{dir} (which must exist) in which to
  create the files for this call only, \code{record} being then a plain name, and
  \code{legacy}.
\end{params}
\returns{\code{Ok(())}.}
\raises{An \code{Error} if a \code{samps\_per\_frame} value is less than 1
(``write\_record: samps\_per\_frame must hold integers >= 1'', checked first); with
\code{opt.dir}, if the directory does not exist (``directory D does not exist'') or the
record name contains one (``record name must not contain a directory when dir is
given''); if the record name is empty (``write\_record: the record name must not be
empty''); if \code{nsig} is 0 (``write\_record: at least one signal is required''); if a
per-signal field of \code{opt} has more than one value but not one per signal
(``write\_record: OPTION must have one value per signal''); if \code{fs} is not $> 0$ or
not finite (``write\_record: the sampling frequency must be > 0'', ``\ldots must be
finite''); if a gain is not finite (``write\_record: gains must be finite''); if an entry
of \code{opt.metadata} is not valid (``write\_record: metadata key 'K' is also given in
comments'', for example); if the gains, the ADC resolutions or the units of a calibrated
signal are missing (``write\_record: gains must be specified'', ``write\_record: ADC
resolution(s) must be specified'', ``write\_record: units must be specified''); if
\code{samples.len()} is not a multiple of the frame size (``write\_record: the number of
samples is not a multiple of the frame size''); if \code{opt.base\_time} is not a valid
time (``write\_record: incorrect base time 'T'''); if the files cannot be created
(``can't create signal file for record REC'', ``can't create header for record REC''); or
if they cannot be written completely, for example on a full disk (``write error on record
REC'').}
\begin{lstlisting}[language=Rust,style=example]
// A ramp of 1000 samples of one signal, in millivolts, at 100 Hz: the record "ramp"
// (ramp.hea and ramp.dat).
let x: Vec<i32> = (0..1000).collect();
let opt = WriteOptions {
    gains: vec![200.0],
    adcres: vec![12],
    units: vec![Some("mV".into())],
    names: vec!["ramp".into()],
    base_time: Some("14:00:03 21/06/2024".into()),
    comments: vec!["a ramp".into()],
    metadata: vec![("Generator".into(), "loop".into())],
    ..Default::default()
};
write_record("ramp", &x, 1, 100.0, &opt)?;
let rec = read_record("ramp", &ReadOptions::default())?;
println!("{} {} {}", rec.nsamples(), rec.base_time.as_deref().unwrap_or(""), rec.physical(10, 0)?);
// 1000 [14:00:03 21/06/2024] 0.05
for c in &read_header("ramp", false)?.comments {
    println!("[{}]", s(c));
}
// [a ramp]
// [ Generator: loop]
\end{lstlisting}
A multi-frequency record: signal \code{fast} has 2 samples per frame, signal
\code{slow} one; each frame holds the two samples of \code{fast} and then the sample of
\code{slow}:
\begin{lstlisting}[language=Rust,style=example]
let mut frames = Vec::new();
for i in 0..100 {
    // 1 second at 100 frames per second
    frames.push(2 * i); // fast: 200 samples per second
    frames.push(2 * i + 1);
    frames.push(500); // slow: 100 samples per second
}
let opt = WriteOptions {
    gains: vec![200.0],
    adcres: vec![12],
    units: vec![Some("mV".into())],
    names: vec!["fast".into(), "slow".into()],
    samps_per_frame: vec![2, 1],
    ..Default::default()
};
write_record("mf", &frames, 2, 100.0, &opt)?;
let hi = ReadOptions { highres: true, ..Default::default() };
let rec = read_record("mf", &hi)?;
println!("{} {} {} {} {}", rec.fs, rec.at(0, 0), rec.at(0, 1), rec.at(1, 0), rec.at(1, 1));
// 200 0 500 1 500
\end{lstlisting}

\apientry{write\_record\_physical}{Write a record from samples in physical units}
\begin{lstlisting}[language=Rust,style=signature]
pub fn write_record_physical(
    record: impl AsRef<[u8]>,
    samples: &[f64],
    nsig: usize,
    fs: f64,
    opt: &WriteOptions,
) -> Result<(), Error>
\end{lstlisting}
As \code{write\_record}, with the samples in physical units. Each sample $v$ is stored as
$\mathrm{round}(v \cdot \mathit{gain}) + \mathit{baseline}$ (halves rounded away from
zero), and NaN is stored as an invalid sample. Every signal must have a gain other than
0.
\begin{params}
\item[record] Record name, as for \code{write\_record}.
\item[samples] The samples, in physical units, frame by frame (as for
  \code{write\_record}).
\item[nsig] The number of signals ($> 0$).
\item[fs] The sampling frequency, in Hz: a finite number $> 0$.
\item[opt] The options, as for \code{write\_record} (\cref{tab:rust-writeoptions}).
\end{params}
\returns{\code{Ok(())}.}
\raises{As \code{write\_record}; also an \code{Error} if a gain is 0 (``write\_record:
signal N is uncalibrated''), or if a sample is infinite or its ADC value is outside the
range of \code{i32} (``write\_record: sample I is infinite or beyond the range of 32-bit
ADC units''). Nothing is written then.}
\begin{lstlisting}[language=Rust,style=example]
// Two sine waves, in millivolts, at 360 Hz, frame by frame: the record "sine".
let fs = 360.0;
let mut x = Vec::new();
for i in 0..3600 {
    let t = 2.0 * std::f64::consts::PI * i as f64 / fs;
    x.push(t.sin()); // signal 0
    x.push(0.5 * t.cos()); // signal 1
}
let opt = WriteOptions {
    gains: vec![1000.0],
    adcres: vec![16],
    units: vec![Some("mV".into())],
    names: vec!["A".into(), "B".into()],
    ..Default::default()
};
write_record_physical("sine", &x, 2, fs, &opt)?;
let rec = read_record("sine", &ReadOptions::default())?;
let mut err: f64 = 0.0;
for (i, v) in x.iter().enumerate() {
    err = err.max((rec.physical(i / 2, i % 2)? - v).abs());
}
println!("{} {}", rec.nsamples(), err <= 0.0005); // 3600 true: rounded to 1/1000 mV
\end{lstlisting}

\apientry{write\_annotations}{Write an annotation file}
\begin{lstlisting}[language=Rust,style=signature]
pub fn write_annotations(
    record: impl AsRef<[u8]>,
    annotator: impl AsRef<[u8]>,
    ann: &Annotations,
    fs: f64,
    dir: Option<&Path>,
) -> Result<(), Error>
\end{lstlisting}
Writes an annotation file, \file{RECORD.ANNOTATOR}, in the MIT format; an existing file of
that name is overwritten. The record need not have a header. The type of each annotation
is its \code{symbols} entry (\code{b"N"}, \code{b"V"}, \ldots; an unknown mnemonic gives
\code{NOTQRS}), or, if \code{symbols} is empty or the entry is empty, its \code{codes}
entry (\cref{chap:formats}). Both may be given (as \code{read\_annotations} returns them)
if they agree: the code of each non-empty symbol (\code{strann}) is the code given for
that annotation. \code{subtypes}, \code{chans} and \code{nums} are stored in one byte
each, as by \code{putann} (other values wrap around: the subtype and number to $-128$ to
127, the channel to 0 to 255), and \code{aux} texts are cut to 255 bytes. If the
annotations are not in time order, the file is sorted when it is written, unless the
environment variable \code{WFDBANNSORT} is \code{0} (\code{write\_annotations\_with}
decides this for one call). The arguments are checked before the file is created;
whatever happens, the annotator is closed before the function returns, and the getvec
mode and the time resolution of output annotators (\code{getafreq()}) are left as they
were (\code{fs} applies to this file only).
\begin{params}
\item[record] Record name, without the \file{.hea} extension (a trailing \code{".hea"} is
  also accepted): a plain name or a relative path (\code{"data/100s"}), created in the
  output directory if one is set, otherwise in the current directory; or a full path,
  absolute or in the home directory (\code{"\textasciitilde/Documents/100s"}), used as
  given. The file is \file{RECORD.ANNOTATOR} at that place (\file{data/100s.qrs}). A URL
  cannot be written.
\item[annotator] Annotator name: the suffix of the file. \code{"qrs"} writes
  \file{RECORD.qrs}.
\item[ann] The annotations (\cref{tab:rust-annotations}): \code{samples}, and
  \code{symbols} or \code{codes}; \code{subtypes}, \code{chans}, \code{nums} and
  \code{aux} are optional. Each of these vectors is empty (\code{subtypes}, \code{chans}
  and \code{nums} are then 0, and there is no aux) or has one entry per annotation.
  \code{ann.record}, \code{ann.annotator} and \code{ann.fs} are not used.
\item[fs] 0 (none) or a finite frequency $> 0$, in which the times are counted, recorded
  in the file; it applies to this file only.
\item[dir] The directory (which must exist) in which to create the file, for this call
  only. \code{record} must then be a plain name, without a directory: the file is
  \file{DIR/RECORD.ANNOTATOR}. \code{None}: the output directory.
\end{params}
\returns{\code{Ok(())}.}
\raises{An \code{Error} if a per-annotation field has another length than 0 or the number
of annotations (``write\_annotations: symbols must have one entry per annotation (N), or
none'', for example); if \code{symbols} and \code{codes} are both given and disagree
(``write\_annotations: symbols and codes disagree at annotation K'', K counted from 0);
if \code{fs} is not 0 and not a finite frequency $> 0$ (``write\_annotations: fs must be
finite and > 0 (or 0: none)''); if the file cannot be created (``can't create annotator
ANN for record REC''); if it cannot be written completely (``write error on annotator ANN
for record REC''); or, with \code{dir}, if the directory does not exist (``directory D
does not exist'') or \code{record} contains a directory (``record name must not contain a
directory when dir is given'').}
\begin{lstlisting}[language=Rust,style=example]
// Three beats and a rhythm change, by symbol, with an aux string for the rhythm.
let a = Annotations {
    samples: vec![18, 77, 370, 662],
    symbols: vec!["+".into(), "N".into(), "V".into(), "N".into()],
    aux: vec![Some("(N".into()), None, None, None],
    ..Default::default()
};
write_annotations("100s", "test", &a, 0.0, None)?;
let b = read_annotations("100s", "test", None, None, None, false)?;
let aux = b.aux[0].as_deref().unwrap_or(b"");
println!("{} {} {}", s(&b.symbols[0]), s(&b.symbols[2]), s(aux)); // + V (N

// A copy of the reference annotations, as read_annotations returns them.
let atr = read_annotations("100s", "atr", None, None, None, false)?;
write_annotations("100s", "copy", &atr, 0.0, None)?;
let c = read_annotations("100s", "copy", None, None, None, false)?;
println!("{}", c.samples == atr.samples && c.aux == atr.aux); // true
\end{lstlisting}

\apientry{write\_annotations\_with}{Write an annotation file, with options}
\begin{lstlisting}[language=Rust,style=signature]
pub fn write_annotations_with(
    record: impl AsRef<[u8]>,
    annotator: impl AsRef<[u8]>,
    ann: &Annotations,
    fs: f64,
    opt: &AnnWriteOptions,
) -> Result<(), Error>
\end{lstlisting}
As \code{write\_annotations}, with the output directory and the sorting of the file given
by \code{opt}, for this call only.
\begin{params}
\item[record] Record name, as for \code{write\_annotations}.
\item[annotator] Annotator name, as for \code{write\_annotations}.
\item[ann] The annotations, as for \code{write\_annotations}.
\item[fs] 0 or the frequency recorded in the file, as for \code{write\_annotations}.
\item[opt] The options (\cref{tab:rust-annwriteoptions}): \code{dir}, as the argument
  \code{dir} of \code{write\_annotations}, and \code{sort}: \code{Some(true)} always sorts
  annotations that are not in time order, \code{Some(false)} keeps the given order
  (without the library's advice to run \code{sortann}), and \code{None} lets
  \code{WFDBANNSORT} decide (sorting is on unless it is \code{0}). The environment is not
  changed.
\end{params}
\returns{\code{Ok(())}.}
\raises{As \code{write\_annotations}.}
\begin{lstlisting}[language=Rust,style=example]
let u = Annotations {
    samples: vec![500, 100],
    codes: vec![NORMAL, NORMAL],
    ..Default::default()
};
// Sorted when the file is written (the library reports it on the standard error).
write_annotations_with("100s", "sorted", &u, 0.0, &AnnWriteOptions::default())?;
let a = read_annotations("100s", "sorted", None, None, None, false)?;
println!("{:?}", a.samples); // [100, 500]
// Kept in the order given, in the existing directory "out".
std::fs::create_dir_all("out")?;
let opt = AnnWriteOptions { dir: Some("out".into()), sort: Some(false) };
write_annotations_with("100s", "unsorted", &u, 0.0, &opt)?;
println!("{}", std::path::Path::new("out/100s.unsorted").exists()); // true
\end{lstlisting}

% ------------------------------------------------------------------------------------------
\subsection{Checking records}
\label{sec:rust-checking}

\apientry{check\_record}{Check a record and its annotation files}
\begin{lstlisting}[language=Rust,style=signature]
pub fn check_record(record: impl AsRef<[u8]>, annotators: &[&str]) -> Result<CheckReport, Error>
\end{lstlisting}
Checks the header of a record, its signal specifications, every sample of its signal
files, each segment of a multi-segment record, and the annotation files named, and
reports each problem found. A problem in the record is never an error of the function: it
is reported, with a severity (\code{Severity::Error}: the record is damaged or does not
follow the WFDB format, and programs will read wrong data or fail;
\code{Severity::Warning}: the record can be read, but something is missing or unusual,
and may give wrong results; \code{Severity::Info}: a fact worth knowing, such as invalid
samples). The problems, their codes and their messages are the same in the seven
languages (\cref{tab:python-issue-codes}). The checks run in normal mode, even if legacy
mode is on, and the library prints no messages during the call (the previous settings
are restored afterwards, the getvec mode and the time resolution of output annotators
among them). They are similar to those of Benjamin Moody's \code{wfdbcheck}.
\begin{params}
\item[record] Record name, as for \code{read\_record}: a full path (absolute or
  \code{\textasciitilde}), used as given; a relative path, read from the input directory
  if one is set, otherwise from the current directory; or a URL, read only if network
  access is on.
\item[annotators] The annotator names whose files are checked, such as \code{\&["atr"]};
  \code{\&[]}: none.
\end{params}
\returns{A \code{CheckReport} (\cref{tab:rust-checkreport}): \code{report.ok} is
\code{true} if no problem is an error, \code{report.issues} lists the problems
(\code{Issue} values, \cref{tab:rust-issue}), and \code{report.lines()} (or
\code{Display}) gives one line per problem, or \code{no problems found}.}
\raises{An \code{Error} if \code{record} is empty (``record must not be empty''). A record
that cannot be read is reported (code \code{header\_unreadable}), not returned as an
error.}
\begin{lstlisting}[language=Rust,style=example]
let report = check_record("100s", &["atr"])?;
println!("{}", report.ok); // true: no errors
println!("{report}");
// warning: signal 0: the units are not specified
// warning: signal 1: the units are not specified

// A damaged record: one sample of a 12-bit signal is overwritten with 2048.
let x: Vec<i32> = (0..100).collect();
let opt = WriteOptions {
    gains: vec![100.0],
    adcres: vec![12],
    units: vec![Some("mV".into())],
    names: vec!["ramp".into()],
    ..Default::default()
};
write_record("damaged", &x, 1, 100.0, &opt)?;
let mut dat = std::fs::read("damaged.dat")?;
dat[20] = 0x00; // sample 10 (2 bytes per sample in format 16):
dat[21] = 0x08; // 0x0800 = 2048, little-endian
std::fs::write("damaged.dat", &dat)?;
let report = check_record("damaged", &[])?;
println!("{}", report.ok); // false
for i in &report.issues {
    println!("{} {} {}", i.severity, i.code, i.r#where);
}
// error sample_range signal 0
// error checksum signal 0
println!("{}", report.lines()[0]);
// error: signal 0: 1 sample is outside the ADC range (-2048 to 2047): sample 10 (value 2048)
\end{lstlisting}

\apientry{check\_record\_with}{Check a record, with options}
\begin{lstlisting}[language=Rust,style=signature]
pub fn check_record_with(
    record: impl AsRef<[u8]>,
    annotators: &[&str],
    opt: &ReadOptions,
) -> Result<CheckReport, Error>
\end{lstlisting}
As \code{check\_record}, with the options \code{input\_dir} and \code{network} of
\code{opt}, for this call only. The other fields of \code{opt} are not used.
\begin{params}
\item[record] Record name, as for \code{check\_record}.
\item[annotators] The annotator names whose files are checked, as for
  \code{check\_record}.
\item[opt] The options (\cref{tab:rust-readoptions}); only \code{input\_dir} (optional:
  the input directory under which relative record names are read; it does not apply to
  full paths or URLs) and \code{network} are used.
\end{params}
\returns{A \code{CheckReport} (\cref{tab:rust-checkreport}).}
\raises{As \code{check\_record}.}
\begin{lstlisting}[language=Rust,style=example]
std::fs::create_dir_all("in")?;
for f in ["100s.hea", "100s.dat", "100s.atr"] {
    std::fs::copy(f, format!("in/{f}"))?;
}
let opt = ReadOptions { input_dir: Some("in".into()), ..Default::default() };
let report = check_record_with("100s", &["atr"], &opt)?;
println!("{} {}", report.ok, report.issues.len()); // true 2
\end{lstlisting}

% ------------------------------------------------------------------------------------------
\subsection{Header metadata}
\label{sec:rust-metadata}

These functions read and change the metadata of a header: its info strings of the form
\code{\# Key: value}, such as \code{\# Age: 63} (\cref{chap:metadata}). In short: keys are
compared without regard to letter case (A-Z and a-z only); a key must not be empty and
must not contain \code{:}, \code{\#}, CR, LF or NUL, and a value must not contain CR, LF
or NUL; spaces, tabs, VT, FF and CR at either end of keys and values are removed; keys and
values are bytes (arguments \code{impl AsRef<[u8]>}, results \code{Vec<u8>}). A change
rewrites the header in one step, through a temporary file in the directory of the
destination, which then replaces it, keeping its permissions (on Unix) and every other
byte of it; the temporary file is removed on any failure. All the arguments are checked
first, so an error leaves the header unchanged, and nothing is written if nothing changes
(unless \code{dir} is given). Only the \file{.hea} file is read (not a \file{.info}
file); for a multi-segment record, it is the master header, unless a segment is named.
The header is opened as the library opens a header (\code{wfdb\_open}), so in legacy mode
it is searched for in the WFDB path. The entries are returned as a \code{Metadata}:
\begin{lstlisting}[language=Rust]
pub type Metadata = Vec<(Vec<u8>, Vec<u8>)>;
\end{lstlisting}
Each function has a form whose name ends in \code{\_with}, which takes a
\code{MetadataOptions} (\cref{tab:rust-metadataoptions}): the input directory and network
access for reading the header, and, for the functions that change it, the directory
\code{dir} in which to write it. These functions do not call \code{wfdbquit()}. Metadata
often holds information about patients: WFDB-x stores it as plain text, exactly as given,
and does not de-identify it.

\begin{phibox}
A header can hold identifying data about a patient (a name, a date of birth, a medical
record number). The metadata functions read and write such data as given. Records that
are shared must be de-identified first.
\end{phibox}

\apientry{read\_metadata}{Read the metadata of a record}
\begin{lstlisting}[language=Rust,style=signature]
pub fn read_metadata(record: impl AsRef<[u8]>) -> Result<Metadata, Error>
\end{lstlisting}
Reads the metadata entries of the header of a record, in file order. An info string is an
entry if it contains a colon: the key is the text before the first colon, the value the
rest of the line, both without spaces at either end. Info strings without a colon, or
whose key is empty or contains \code{\#}, are free comments, which are not returned.
\begin{params}
\item[record] Record name, without the \file{.hea} extension (a trailing \code{".hea"} is
  also accepted): a full path, absolute (\code{"/data/mitdb/100"}) or in the home
  directory (\code{"\textasciitilde/Documents/100s"}), used as given; a relative path
  (\code{"100s"}, \code{"data/100s"}), read from the input directory if one is set,
  otherwise from the current directory; or a URL
  (\code{"https://physionet.org/files/mitdb/1.0.0/100"}), read only if network access is
  on.
\end{params}
\returns{The entries, as (key, value) pairs, in file order. A key that appears more than
once (ignoring letter case) is returned once, with its first value, with a warning
(\cref{sec:rust-warnings}).}
\raises{An \code{Error} if the header cannot be found or read (``can't read header for
record REC''), or if \code{record} is empty (``record name must not be empty'') or
\code{"-"} (``the metadata of a header read from the standard input cannot be read or
changed'').}
\begin{lstlisting}[language=Rust,style=example]
let opt = WriteOptions {
    gains: vec![200.0],
    adcres: vec![12],
    units: vec![Some("mV".into())],
    names: vec!["ECG".into()],
    comments: vec!["recorded at rest".into()],
    metadata: vec![("Age".into(), "63".into()), ("Sex".into(), "F".into())],
    ..Default::default()
};
write_record("pat", &[0; 360], 1, 360.0, &opt)?;
for (key, value) in read_metadata("pat")? {
    println!("{} = {}", s(&key), s(&value));
}
// Age = 63
// Sex = F
for c in read_header("pat", false)?.comments {
    println!("[{}]", s(&c));
}
// [recorded at rest]: the free comment is not an entry
// [ Age: 63]
// [ Sex: F]
\end{lstlisting}

\apientry{read\_metadata\_with}{Read the metadata of a record, with options}
\begin{lstlisting}[language=Rust,style=signature]
pub fn read_metadata_with(
    record: impl AsRef<[u8]>,
    opt: &MetadataOptions,
) -> Result<Metadata, Error>
\end{lstlisting}
As \code{read\_metadata}, with the input directory and network access of \code{opt}, for
this call only. \code{opt.dir} is not used.
\begin{params}
\item[record] Record name, as for \code{read\_metadata}.
\item[opt] The options (\cref{tab:rust-metadataoptions}); only \code{input\_dir} and
  \code{network} are used.
\end{params}
\returns{The entries, as for \code{read\_metadata}.}
\raises{As \code{read\_metadata}.}
\begin{lstlisting}[language=Rust,style=example]
std::fs::create_dir_all("in")?;
std::fs::copy("100s.hea", "in/100s.hea")?;
let opt = MetadataOptions { input_dir: Some("in".into()), ..Default::default() };
let md = read_metadata_with("100s", &opt)?; // reads in/100s.hea
println!("{} {}", md.len(), s(&md[0].1)); // 1 0
\end{lstlisting}

\apientry{get\_metadata}{Return the value of a metadata entry}
\begin{lstlisting}[language=Rust,style=signature]
pub fn get_metadata(
    record: impl AsRef<[u8]>,
    key: impl AsRef<[u8]>,
) -> Result<Option<Vec<u8>>, Error>
\end{lstlisting}
Returns the value of one metadata entry of the header of a record.
\begin{params}
\item[record] Record name (as for \code{read\_metadata}).
\item[key] The key, such as \code{"Age"}. Letter case is ignored, and spaces at either end
  are removed; it must be a valid key (see \code{set\_metadata}), which is checked before
  the header is read.
\end{params}
\returns{\code{Some(value)} (empty for an entry with an empty value), or \code{None} if the
header has no entry for \code{key}. If it has more than one, the first is returned, with
the warning of \code{read\_metadata}.}
\raises{As \code{read\_metadata}; also an \code{Error} if \code{key} is not valid
(``metadata key 'a\#b' must not contain '\#''', for example).}
\begin{lstlisting}[language=Rust,style=example]
let opt = WriteOptions {
    gains: vec![200.0],
    adcres: vec![12],
    units: vec![Some("mV".into())],
    metadata: vec![("Age".into(), "63".into()), ("Notes".into(), "".into())],
    ..Default::default()
};
write_record("pat", &[0; 360], 1, 360.0, &opt)?;
println!("{:?}", get_metadata("pat", "age")?.map(|v| s(&v))); // Some("63")
println!("{:?}", get_metadata("pat", " NOTES ")?.map(|v| s(&v))); // Some(""): an empty value
println!("{:?}", get_metadata("pat", "Weight")?); // None: no such entry
\end{lstlisting}

\apientry{get\_metadata\_with}{Return the value of a metadata entry, with options}
\begin{lstlisting}[language=Rust,style=signature]
pub fn get_metadata_with(
    record: impl AsRef<[u8]>,
    key: impl AsRef<[u8]>,
    opt: &MetadataOptions,
) -> Result<Option<Vec<u8>>, Error>
\end{lstlisting}
As \code{get\_metadata}, with the input directory and network access of \code{opt}, for
this call only. \code{opt.dir} is not used.
\begin{params}
\item[record] Record name (as for \code{read\_metadata}).
\item[key] The key (as for \code{get\_metadata}).
\item[opt] The options (\cref{tab:rust-metadataoptions}); only \code{input\_dir} and
  \code{network} are used.
\end{params}
\returns{As \code{get\_metadata}.}
\raises{As \code{get\_metadata}.}
\begin{lstlisting}[language=Rust,style=example]
let opt = WriteOptions {
    gains: vec![200.0],
    adcres: vec![12],
    units: vec![Some("mV".into())],
    metadata: vec![("Age".into(), "63".into())],
    ..Default::default()
};
write_record("pat", &[0; 360], 1, 360.0, &opt)?;
std::fs::create_dir_all("in")?;
std::fs::copy("pat.hea", "in/pat.hea")?;
let opt = MetadataOptions { input_dir: Some("in".into()), ..Default::default() };
let age = get_metadata_with("pat", "Age", &opt)?.unwrap_or_default();
println!("{}", s(&age)); // 63
\end{lstlisting}

\apientry{set\_metadata}{Set the value of a metadata entry}
\begin{lstlisting}[language=Rust,style=signature]
pub fn set_metadata(
    record: impl AsRef<[u8]>,
    key: impl AsRef<[u8]>,
    value: impl AsRef<[u8]>,
) -> Result<(), Error>
\end{lstlisting}
Sets the value of a metadata entry in the header of a record. If the header has an entry
for \code{key} (ignoring letter case), its line is replaced by \code{\# Key: value}, at
the same position and with its line ending, with the spelling of \code{key} given here
(even if the value is the same). Otherwise the entry is added after the last info string
of the header (or, if it has none, after its last line that is not blank), with the line
ending of the record line. Every other byte of the header is kept.
\begin{params}
\item[record] Record name (as for \code{read\_metadata}). The header is changed where it
  was read, unless \code{dir} is given (\code{set\_metadata\_with}); a header read from a
  URL can only be written elsewhere, with \code{set\_metadata\_with} and \code{dir}.
\item[key] The key, such as \code{"Age"} or \code{"Date of birth"}. Spaces at either end
  are removed; it must not then be empty, and must not contain \code{:}, \code{\#}, CR,
  LF or NUL.
\item[value] The value, which may be a sentence. Spaces at either end are removed; it
  must not contain CR, LF or NUL, and may be empty (it is then written \code{\# Key:}).
\end{params}
\returns{\code{Ok(())}.}
\raises{An \code{Error} if \code{key} or \code{value} is not valid (``metadata key must not
be empty'', ``metadata key 'KEY' must not contain C'', ``value of metadata key 'KEY' must
not contain C'', where C is \code{':'}, \code{'\#'}, ``a line break'' or ``a NUL
character''); if \code{record} is empty or \code{"-"}; if the header cannot be read
(``can't read header for record REC''); if it has more than one entry for \code{key}
(``header of record REC has more than one entry for key KEY (lines L1, L2); remove them
first''); if it has no record line (``header of record REC has no record line''); if it
was read from a URL (``header of record REC was read from a URL, and cannot be changed
(give dir to write the changed header to a local directory)''); or if the changed header
cannot be written (``can't write header for record REC'').}
\begin{lstlisting}[language=Rust,style=example]
let opt = WriteOptions {
    gains: vec![200.0],
    adcres: vec![12],
    units: vec![Some("mV".into())],
    metadata: vec![("Age".into(), "63".into())],
    ..Default::default()
};
write_record("pat", &[0; 360], 1, 360.0, &opt)?;
set_metadata("pat", "AGE", "64")?; // replaces the line: "# AGE: 64"
set_metadata("pat", "Date of birth", "1961-02-03")?; // adds "# Date of birth: 1961-02-03"
for (key, value) in read_metadata("pat")? {
    println!("{} = {}", s(&key), s(&value));
}
// AGE = 64
// Date of birth = 1961-02-03
if let Err(e) = set_metadata("pat", "a:b", "1") {
    println!("{e}"); // metadata key 'a:b' must not contain ':'
}
\end{lstlisting}

\apientry{set\_metadata\_with}{Set the value of a metadata entry, with options}
\begin{lstlisting}[language=Rust,style=signature]
pub fn set_metadata_with(
    record: impl AsRef<[u8]>,
    key: impl AsRef<[u8]>,
    value: impl AsRef<[u8]>,
    opt: &MetadataOptions,
) -> Result<(), Error>
\end{lstlisting}
As \code{set\_metadata}, with the input directory and network access for reading the
header, and the directory \code{dir} in which to write the changed header, given by
\code{opt}, for this call only. With \code{dir}, the header that was read is unchanged,
and the changed header is written to the file of the same name in \code{dir}, even if
nothing changes.
\begin{params}
\item[record] Record name (as for \code{set\_metadata}). A header read from a URL can be
  changed only with \code{opt.dir}.
\item[key] The key (as for \code{set\_metadata}).
\item[value] The value (as for \code{set\_metadata}).
\item[opt] The options (\cref{tab:rust-metadataoptions}): \code{input\_dir},
  \code{network} and \code{dir}.
\end{params}
\returns{\code{Ok(())}.}
\raises{As \code{set\_metadata}; also an \code{Error} if \code{opt.dir} does not exist
(``directory D does not exist'').}
\begin{lstlisting}[language=Rust,style=example]
// The header of 100s (in the current directory) is unchanged; the changed copy is
// written to the directory "copy" (which must exist).
std::fs::create_dir_all("copy")?;
let opt = MetadataOptions { dir: Some("copy".into()), ..Default::default() };
set_metadata_with("100s", "Age", "69", &opt)?;
let age = get_metadata("copy/100s", "Age")?.unwrap_or_default();
println!("{} {:?}", s(&age), get_metadata("100s", "Age")?); // 69 None
\end{lstlisting}

\apientry{remove\_metadata}{Remove a metadata entry}
\begin{lstlisting}[language=Rust,style=signature]
pub fn remove_metadata(record: impl AsRef<[u8]>, key: impl AsRef<[u8]>) -> Result<usize, Error>
\end{lstlisting}
Removes a metadata entry from the header of a record. Every entry for \code{key}
(ignoring letter case) is removed, with its line ending, so this also repairs a header
that has more than one. Every other byte is kept. Nothing is written if the header has no
entry for \code{key}.
\begin{params}
\item[record] Record name (as for \code{set\_metadata}).
\item[key] The key (as for \code{set\_metadata}).
\end{params}
\returns{The number of lines removed (0 if the header has no entry for \code{key}).}
\raises{As \code{set\_metadata}, except that more than one entry for \code{key} is not an
error.}
\begin{lstlisting}[language=Rust,style=example]
let opt = WriteOptions {
    gains: vec![200.0],
    adcres: vec![12],
    units: vec![Some("mV".into())],
    metadata: vec![("Age".into(), "63".into()), ("Sex".into(), "F".into())],
    ..Default::default()
};
write_record("pat", &[0; 360], 1, 360.0, &opt)?;
println!("{}", remove_metadata("pat", "sex")?); // 1
println!("{}", remove_metadata("pat", "Weight")?); // 0: nothing is written
println!("{}", read_metadata("pat")?.len()); // 1
\end{lstlisting}

\apientry{remove\_metadata\_with}{Remove a metadata entry, with options}
\begin{lstlisting}[language=Rust,style=signature]
pub fn remove_metadata_with(
    record: impl AsRef<[u8]>,
    key: impl AsRef<[u8]>,
    opt: &MetadataOptions,
) -> Result<usize, Error>
\end{lstlisting}
As \code{remove\_metadata}, with the options of \code{opt}, for this call only (as for
\code{set\_metadata\_with}). With \code{dir}, the header is written to \code{dir} even if
nothing is removed.
\begin{params}
\item[record] Record name (as for \code{set\_metadata}).
\item[key] The key (as for \code{set\_metadata}).
\item[opt] The options (\cref{tab:rust-metadataoptions}): \code{input\_dir},
  \code{network} and \code{dir}.
\end{params}
\returns{The number of lines removed.}
\raises{As \code{remove\_metadata}; also an \code{Error} if \code{opt.dir} does not exist.}
\begin{lstlisting}[language=Rust,style=example]
let opt = WriteOptions {
    gains: vec![200.0],
    adcres: vec![12],
    units: vec![Some("mV".into())],
    metadata: vec![("Age".into(), "63".into()), ("Sex".into(), "F".into())],
    ..Default::default()
};
write_record("pat", &[0; 360], 1, 360.0, &opt)?;
std::fs::create_dir_all("out")?;
let opt = MetadataOptions { dir: Some("out".into()), ..Default::default() };
println!("{}", remove_metadata_with("pat", "Sex", &opt)?); // 1
println!("{} {}", read_metadata("pat")?.len(), read_metadata("out/pat")?.len()); // 2 1
\end{lstlisting}

\apientry{update\_metadata}{Make several changes to the metadata}
\begin{lstlisting}[language=Rust,style=signature]
pub fn update_metadata<K: AsRef<[u8]>, V: AsRef<[u8]>>(
    record: impl AsRef<[u8]>,
    changes: &[(K, Option<V>)],
) -> Result<(), Error>
\end{lstlisting}
Makes several changes to the metadata of a record in one write. Each (key, value) of
\code{changes} is applied in order, as \code{set\_metadata} does for \code{Some(value)}, or
as \code{remove\_metadata} does for \code{None}. All the changes are checked first: if one
is not valid, none is made.
\begin{params}
\item[record] Record name (as for \code{set\_metadata}).
\item[changes] The changes, as (key, value) pairs; a value of \code{None} removes the
  entry. No two keys may differ only in letter case. A list of removals only needs the
  type of its values: \code{\&[("Dx", None::<\&str>)]}. Empty \code{changes} change
  nothing (and write nothing).
\end{params}
\returns{\code{Ok(())}.}
\raises{As \code{set\_metadata}; also an \code{Error} if two keys differ only in letter
case (``metadata key 'K' is given more than once (keys are compared without regard to
letter case)'').}
\begin{lstlisting}[language=Rust,style=example]
let opt = WriteOptions {
    gains: vec![200.0],
    adcres: vec![12],
    units: vec![Some("mV".into())],
    metadata: vec![("Age".into(), "63".into()), ("Dx".into(), "164889003".into())],
    ..Default::default()
};
write_record("pat", &[0; 360], 1, 360.0, &opt)?;
update_metadata("pat", &[("Age", Some("64")), ("Dx", None), ("Sex", Some("F"))])?;
for (key, value) in read_metadata("pat")? {
    println!("{} = {}", s(&key), s(&value));
}
// Age = 64
// Sex = F
update_metadata("pat", &[("Sex", None::<&str>)])?; // removals only: the type of the values
println!("{}", read_metadata("pat")?.len()); // 1
\end{lstlisting}

\apientry{update\_metadata\_with}{Make several changes to the metadata, with options}
\begin{lstlisting}[language=Rust,style=signature]
pub fn update_metadata_with<K: AsRef<[u8]>, V: AsRef<[u8]>>(
    record: impl AsRef<[u8]>,
    changes: &[(K, Option<V>)],
    opt: &MetadataOptions,
) -> Result<(), Error>
\end{lstlisting}
As \code{update\_metadata}, with the options of \code{opt}, for this call only (as for
\code{set\_metadata\_with}). With \code{dir}, the header is written to \code{dir} even if
\code{changes} is empty.
\begin{params}
\item[record] Record name (as for \code{set\_metadata}).
\item[changes] The changes (as for \code{update\_metadata}).
\item[opt] The options (\cref{tab:rust-metadataoptions}): \code{input\_dir},
  \code{network} and \code{dir}.
\end{params}
\returns{\code{Ok(())}.}
\raises{As \code{update\_metadata}; also an \code{Error} if \code{opt.dir} does not
exist.}
\begin{lstlisting}[language=Rust,style=example]
let opt = WriteOptions {
    gains: vec![200.0],
    adcres: vec![12],
    units: vec![Some("mV".into())],
    metadata: vec![("Age".into(), "64".into())],
    ..Default::default()
};
write_record("pat", &[0; 360], 1, 360.0, &opt)?;
// The header "pat" in the directory "in" is unchanged; the changed copy is written to
// the directory "out" (which must exist).
std::fs::create_dir_all("in")?;
std::fs::create_dir_all("out")?;
std::fs::copy("pat.hea", "in/pat.hea")?;
let opt = MetadataOptions {
    input_dir: Some("in".into()),
    dir: Some("out".into()),
    ..Default::default()
};
update_metadata_with("pat", &[("Age", Some("65"))], &opt)?;
let a = get_metadata("in/pat", "Age")?.unwrap_or_default();
let b = get_metadata("out/pat", "Age")?.unwrap_or_default();
println!("{} {}", s(&a), s(&b)); // 64 65
\end{lstlisting}

% ------------------------------------------------------------------------------------------
\subsection{The options}
\label{sec:rust-options}

The options of the high-level functions are the fields of the structures of
\cref{tab:rust-readoptions,tab:rust-writeoptions,tab:rust-annwriteoptions,tab:rust-metadataoptions},
all public, and all with a \code{Default} implementation; the times are \code{TimeSpec}
values (\cref{tab:rust-timespec}). \code{ReadOptions}, \code{WriteOptions},
\code{AnnWriteOptions} and \code{TimeSpec} are defined in \code{wfdbx::record},
\code{MetadataOptions} in \code{wfdbx::metadata}; all are available at the root of the
crate.

\begin{xltabular}{\linewidth}{@{}lL@{}}
\caption{The fields of \code{ReadOptions}, the options of \code{read\_record} (all),
  \code{read\_header\_with} (\code{input\_dir}, \code{network}, \code{legacy}),
  \code{read\_annotations\_with} (\code{from}, \code{to}, \code{input\_dir},
  \code{network}, \code{legacy}) and \code{check\_record\_with} (\code{input\_dir},
  \code{network})}
\label{tab:rust-readoptions} \\
\toprule
\tabhead Field & \tabhead Meaning (default) \\
\midrule
\endfirsthead
\toprule
\tabhead Field & \tabhead Meaning (default) \\
\midrule
\endhead
\bottomrule
\endfoot
\code{channels: Vec<String>} & the signals to return, in this order, by number
  (\code{"0"}, \code{"1"}, \ldots) or by name (\code{"V5"}, as in \code{Record::names()},
  then as \code{findsig} finds it). Default empty: all the signals \\
\code{from: Option<TimeSpec>} & the first sample: a sample number ($\ge 0$) or a time
  string. Default \code{None}: 0 \\
\code{to: Option<TimeSpec>} & the sample after the last one: a sample number or a time
  string. Default \code{None}: the end of the record \\
\code{highres: bool} & \code{true} to read a multi-frequency record at its highest
  sampling frequency, repeating the samples of the slower signals; \code{false} averages
  the samples of the faster signals to one per frame. Default \code{false} \\
\code{pad: bool} & \code{true} to replace each invalid sample by the previous valid sample
  of the same signal. Default \code{false} \\
\code{legacy: bool} & behave as the original WFDB package, for this call only. Default
  \code{false} \\
\code{input\_dir: Option<String>} & the input directory for this call only (optional;
  see \code{wfdbx\_set\_input\_dir}): a local directory, a \code{\textasciitilde} path or
  the URL of a directory, under which relative record names are read. It does not apply
  to full paths or URLs. \code{""} means the default (\code{WFDBX\_INPUT\_DIR}, or the
  current directory). Default \code{None}: the current setting \\
\code{network: Option<bool>} & \code{Some(true)} to allow opening URLs for this call
  only, \code{Some(false)} to forbid it (see \code{wfdbx\_set\_network}). Default
  \code{None}: the current setting \\
\code{checksum\_warning: bool} & \code{true} to print a warning on the standard error when
  the samples do not match the checksums of the header (\cref{sec:rust-warnings}); the
  samples are returned either way. Default \code{true} \\
\end{xltabular}

\begin{xltabular}{\linewidth}{@{}lL@{}}
\caption{The fields of \code{WriteOptions}, the options of \code{write\_record} and
  \code{write\_record\_physical}. Each per-signal vector holds one value for all the
  signals, or one value per signal (or none, for those that have a default)}
\label{tab:rust-writeoptions} \\
\toprule
\tabhead Field & \tabhead Meaning (default) \\
\midrule
\endfirsthead
\toprule
\tabhead Field & \tabhead Meaning (default) \\
\midrule
\endhead
\bottomrule
\endfoot
\code{fmt: i32} & the storage format, the same for all the signals, such as \code{16},
  \code{212}, \code{24} or \code{516} (FLAC) (\cref{chap:formats}). Default \code{16} \\
\code{gains: Vec<f64>} & the gains, in ADC units per physical unit: required, and finite.
  A gain of 0 marks a signal as uncalibrated, which only \code{write\_record} (ADC units)
  allows. Negative gains are allowed \\
\code{baselines: Vec<i32>} & the ADC values that correspond to 0 physical units.
  Default: the \code{adczeros} values \\
\code{adczeros: Vec<i32>} & the ADC output for an input of 0\,V. Default
  \code{vec![0]} \\
\code{adcres: Vec<i32>} & the ADC resolutions, in bits: required, each at least 1 \\
\code{units: Vec<Option<Vec<u8>{}>{}>} & the physical units, such as
  \code{Some("mV".into())}: required for every signal whose gain is not 0; \code{None}
  for an uncalibrated signal \\
\code{names: Vec<Vec<u8>{}>} & the signal descriptions. Default: empty descriptions \\
\code{samps\_per\_frame: Vec<i32>} & the number of samples per frame of each signal, for a
  multi-frequency record (each $\ge 1$). Default: 1 \\
\code{base\_time: Option<Vec<u8>{}>} & the start time and date, such as
  \code{"12:00:00 01/01/2020"} (see \code{setbasetime}). Default \code{None}: none \\
\code{comments: Vec<Vec<u8>{}>} & the info strings, written to the header as comment lines
  (\code{"\#"} followed by the string). Default: none \\
\code{metadata: Vec<(Vec<u8>, Vec<u8>)>} & metadata entries (\cref{chap:metadata})
  written to the header after the comments, as \code{\# Key: value} (through
  \code{putinfo}), in this order. The keys and values follow the rules of
  \code{set\_metadata}; no two keys may differ only in letter case, and none may be the
  key of an entry in \code{comments} (an error, before any file is created). Default:
  none \\
\code{legacy: bool} & behave as the original WFDB package, for this call only: the gain
  defaults to 200, and the ADC resolution to one based on the format. Default
  \code{false} \\
\code{dir: Option<PathBuf>} & the directory (which must exist) in which the files are
  created, for this call only; the record name must then be a plain name, without a
  directory, and the record opens from that directory under its plain name. Default
  \code{None} (or an empty path): the output directory \\
\end{xltabular}

\begin{table}[htbp]
\centering
\caption{The fields of \code{AnnWriteOptions}, the options of
  \code{write\_annotations\_with}}
\label{tab:rust-annwriteoptions}
\begin{tabularx}{\linewidth}{@{}lL@{}}
\toprule
\tabhead Field & \tabhead Meaning (default) \\
\midrule
\code{dir: Option<PathBuf>} & the directory (which must exist) in which to create the
  file, for this call only; the record name must then be a plain name, without a
  directory: the file is \file{DIR/RECORD.ANNOTATOR}. Default \code{None}: the output
  directory \\
\code{sort: Option<bool>} & \code{Some(true)} always sorts annotations that are not in
  time order, \code{Some(false)} keeps the given order, for this call only. Default
  \code{None}: \code{WFDBANNSORT} decides (sorting is on unless it is \code{0}) \\
\bottomrule
\end{tabularx}
\end{table}

\begin{table}[htbp]
\centering
\caption{The fields of \code{MetadataOptions}, the options of the metadata functions whose
  names end in \code{\_with}}
\label{tab:rust-metadataoptions}
\begin{tabularx}{\linewidth}{@{}lL@{}}
\toprule
\tabhead Field & \tabhead Meaning (default) \\
\midrule
\code{input\_dir: Option<String>} & the input directory for reading the header, for this
  call only (optional), as in \code{ReadOptions}: it applies to relative record names
  only. Default \code{None}: the current setting \\
\code{network: Option<bool>} & network access for this call only, as in
  \code{ReadOptions}. Default \code{None}: the current setting \\
\code{dir: Option<PathBuf>} & for \code{set\_metadata\_with},
  \code{remove\_metadata\_with} and \code{update\_metadata\_with}: the directory (which
  must exist; a \code{\textasciitilde} path is expanded) in which the changed header is
  written, under the file name of the header that was read, instead of changing that
  header. The header is then written even if nothing changes. The signal files are not
  copied, and the output directory is not used. Not used by \code{read\_metadata\_with}
  and \code{get\_metadata\_with}. Default \code{None} (or an empty path): the header is
  changed where it was read \\
\bottomrule
\end{tabularx}
\end{table}

\begin{xltabular}{\linewidth}{@{}lL@{}}
\caption{\code{TimeSpec}, a sample number or a time string}
\label{tab:rust-timespec} \\
\toprule
\tabhead Item & \tabhead Meaning \\
\midrule
\endfirsthead
\toprule
\tabhead Item & \tabhead Meaning \\
\midrule
\endhead
\bottomrule
\endfoot
\code{TimeSpec::Sample(i64)} & a sample number \\
\code{TimeSpec::Str(Vec<u8>)} & a time string, converted with \code{strtim} for the open
  record \\
\code{From<i64>}, \code{From<i32>}, \code{From<u32>}, \code{From<usize>} & the sample
  number: \code{3600.into()} \\
\code{From<\&str>}, \code{From<String>}, \code{From<\&[u8]>} & the time string:
  \code{"0:10".into()} \\
\code{fn resolve(\&self, w: \&mut Wfdb) -> i64} & the sample number; a time string is
  converted with \code{strtim} for the record open in \code{w}, and a time of day gives
  the elapsed time from the start of the record (the absolute value of what
  \code{strtim} returns) \\
\code{fn needs\_freq(\&self, w: \&mut Wfdb) -> bool} & \code{true} if the time is a
  string that needs a sampling frequency to be converted (see
  \code{wfdbx\_time\_needs\_freq}) \\
\end{xltabular}

% ------------------------------------------------------------------------------------------
\subsection{The result types}
\label{sec:rust-results}

The readers return structures of the types \code{Record}, \code{Header} and
\code{Annotations}, whose fields and methods are listed in
\cref{tab:rust-record,tab:rust-header,tab:rust-annotations}; \code{check\_record} returns
a \code{CheckReport}, which holds \code{Issue} values
(\cref{tab:rust-checkreport,tab:rust-issue}). The signals of a record are described by
\code{Siginfo} structures (\cref{tab:rust-siginfo}), which are those of the WFDB library
(\cref{sec:rust-library}). Their fields are public, and may be changed. \code{Record},
\code{Header} and \code{Annotations} implement \code{Clone}, \code{Debug}, \code{Default}
and \code{PartialEq}; \code{CheckReport} and \code{Issue} implement \code{Clone},
\code{Debug}, \code{PartialEq} and \code{Eq}.

\begin{xltabular}{\linewidth}{@{}lL@{}}
\caption{The fields and methods of \code{Record}, the result of \code{read\_record}}
\label{tab:rust-record} \\
\toprule
\tabhead Member & \tabhead Contents \\
\midrule
\endfirsthead
\toprule
\tabhead Member & \tabhead Contents \\
\midrule
\endhead
\bottomrule
\endfoot
\code{name: Vec<u8>} & the record name, as given to \code{read\_record} \\
\code{fs: f64} & the sampling frequency of the returned samples, in Hz, or 0 if the header
  gives none; with \code{highres}, the highest sampling frequency of a multi-frequency
  record \\
\code{nsig: usize} & the number of signals returned \\
\code{signals: Vec<usize>} & the number, in the record, of each signal returned \\
\code{samples: Vec<i32>} & the samples in ADC units, frame by frame:
  \code{samples[i * nsig + s]} is sample \code{i} of signal \code{s}; invalid samples are
  \code{WFDB\_INVALID\_SAMPLE} \\
\code{siginfo: Vec<Siginfo>} & one per signal returned, in the order of the signals;
  \code{desc} has no white space at either end \\
\code{base\_time: Option<String>} & the start time and date as text, to the second, such
  as \code{"[14:00:03 21/06/2024]"}; \code{None} if the header gives none \\
\code{base\_time\_ms: Option<i64>} & the start time in milliseconds after midnight
  (\code{"14:00:03.500"} gives \code{50403500}); \code{Some} exactly when
  \code{base\_time} is \\
\code{base\_date: Option<[i32; 3]>} & the start date as \code{[year, month, day]};
  \code{None} if the header gives no date \\
\code{sampfrom: i64} & the number of the first sample returned \\
\code{comments: Vec<Vec<u8>{}>} & the info strings of the record: the comments of the
  header, and the lines of its \file{.info} file if it has one (see \code{getinfo}) \\
\code{checksum\_ok: bool} & \code{false} if the samples do not match the checksums of the
  header; checksums are checked only when a read reaches the end of the record, so a
  shorter read gives \code{true} \\
\code{legacy: bool} & \code{true} if the record was read in legacy mode;
  \code{physical()} then uses a gain of 200 for uncalibrated signals \\
\code{fn nsamples(\&self) -> usize} & the number of samples per signal (the number of
  frames) \\
\code{fn at(\&self, i: usize, s: usize) -> i32} & sample \code{i} of signal \code{s}, in
  ADC units; panics if \code{i} or \code{s} is out of range \\
\code{fn physical(\&self, i: usize, s: usize)} & a \code{Result<f64, Error>}: sample
  \code{i} of signal \code{s} in physical units, $(\mathit{value} - \mathit{baseline}) /
  \mathit{gain}$; NaN if the sample is invalid. An \code{Error} (``signal N of record REC
  is uncalibrated'') if the gain of the signal is 0 \\
\code{fn physical\_signal(\&self, s: usize)} & a \code{Result<Vec<f64>, Error>}: all the
  samples of signal \code{s} in physical units (the same error) \\
\code{fn names(\&self) -> Vec<Vec<u8>{}>} & the signal descriptions (empty if absent),
  without the white space that a header may put at either end of them \\
\end{xltabular}

\begin{xltabular}{\linewidth}{@{}lL@{}}
\caption{The fields of \code{Header}, the result of \code{read\_header} and
  \code{read\_header\_with}}
\label{tab:rust-header} \\
\toprule
\tabhead Field & \tabhead Contents \\
\midrule
\endfirsthead
\toprule
\tabhead Field & \tabhead Contents \\
\midrule
\endhead
\bottomrule
\endfoot
\code{name: Vec<u8>} & the record name, as given \\
\code{fs: f64} & the sampling frequency in Hz, as \code{sampfreq} gives it; 0 if the
  header gives none \\
\code{counter\_freq: f64} & the counter frequency (\code{getcfreq}) \\
\code{base\_counter: f64} & the base counter value (\code{getbasecount}) \\
\code{nsamples: i64} & the number of samples per signal, from the record line of the
  header (0 if the header does not give it) \\
\code{nsig: usize} & the number of signals \\
\code{siginfo: Vec<Siginfo>} & the signal specifications, one per signal; \code{desc} has
  no white space at either end \\
\code{base\_time}, \code{base\_time\_ms}, \code{base\_date} & the start time and date, as
  in \code{Record} \\
\code{comments: Vec<Vec<u8>{}>} & the info strings of the record, as in \code{Record} \\
\code{metadata: Vec<(Vec<u8>, Vec<u8>)>} & the metadata entries of the header, in file
  order, as \code{read\_metadata} returns them: from the \file{.hea} file only; empty,
  without an error, for an EDF record and for a header read from the standard input \\
\end{xltabular}

\begin{xltabular}{\linewidth}{@{}lL@{}}
\caption{The fields and methods of \code{Annotations}, the result of
  \code{read\_annotations} and the argument of \code{write\_annotations}: parallel
  vectors, annotation \code{i} being \code{samples[i]}, \code{codes[i]},
  \code{symbols[i]}, \ldots, in file order}
\label{tab:rust-annotations} \\
\toprule
\tabhead Member & \tabhead Contents \\
\midrule
\endfirsthead
\toprule
\tabhead Member & \tabhead Contents \\
\midrule
\endhead
\bottomrule
\endfoot
\code{record: Vec<u8>}, \code{annotator: Vec<u8>} & the record and annotator names, such
  as \code{b"100s"} and \code{b"atr"} \\
\code{fs: f64} & the frequency in which \code{samples} are counted, normally the sampling
  frequency of the record; 0 if unknown. \code{write\_annotations} does not use it \\
\code{samples: Vec<i64>} & the annotation times, in samples \\
\code{codes: Vec<i32>} & the annotation codes, such as \code{NORMAL} or \code{PVC} \\
\code{symbols: Vec<Vec<u8>{}>} & the mnemonics of the codes, as \code{annstr} returns them,
  such as \code{b"N"}, \code{b"V"} or \code{b"+"} \\
\code{subtypes: Vec<i32>}, \code{chans: Vec<i32>}, \code{nums: Vec<i32>} & the subtype,
  channel number and annotator number of each annotation \\
\code{aux: Vec<Option<Vec<u8>{}>{}>} & the auxiliary text of each annotation, without its
  length byte and up to its first NUL byte, or \code{None} where there is none \\
\code{fn len(\&self) -> usize} & the number of annotations \\
\code{fn is\_empty(\&self) -> bool} & \code{true} if there are no annotations \\
\end{xltabular}

\begin{table}[htbp]
\centering
\caption{The fields and methods of \code{CheckReport}, the result of \code{check\_record}}
\label{tab:rust-checkreport}
\begin{tabularx}{\linewidth}{@{}lL@{}}
\toprule
\tabhead Member & \tabhead Contents \\
\midrule
\code{record: Vec<u8>} & the record name, as given \\
\code{ok: bool} & \code{true} if no issue has the severity \code{Severity::Error} \\
\code{issues: Vec<Issue>} & the problems, in the order of the checks \\
\code{fn lines(\&self) -> Vec<String>} & the text form of the report: one line per issue
  (the \code{Display} form of the \code{Issue}), or the one line \code{no problems
  found} \\
\code{Display} & the lines of \code{lines()}, joined by line feeds (no final line feed) \\
\bottomrule
\end{tabularx}
\end{table}

\begin{table}[htbp]
\centering
\caption{The fields of \code{Issue}, a problem found by \code{check\_record}, and of
  \code{Severity}}
\label{tab:rust-issue}
\begin{tabularx}{\linewidth}{@{}lL@{}}
\toprule
\tabhead Member & \tabhead Contents \\
\midrule
\code{severity: Severity} & \code{Severity::Error}, \code{Severity::Warning} or
  \code{Severity::Info}; \code{as\_str()} and \code{Display} give \code{"error"},
  \code{"warning"} or \code{"info"} \\
\code{code: \&'static str} & a stable identifier of the kind of problem, such as
  \code{"checksum"} or \code{"gain\_missing"} (\cref{tab:python-issue-codes}) \\
\code{r\#where: String} & where the problem is: \code{""} (the record),
  \code{"signal N"}, \code{"segment SEG"}, \code{"segment SEG, signal N"} or
  \code{"annotator ANN"} (the field is named \code{where}, a keyword of Rust, written
  \code{r\#where}) \\
\code{message: String} & the description of the problem \\
\code{Display} & \code{SEVERITY: WHERE: MESSAGE}, or \code{SEVERITY: MESSAGE} when
  \code{where} is empty \\
\bottomrule
\end{tabularx}
\end{table}

\begin{xltabular}{\linewidth}{@{}lL@{}}
\caption{The fields of \code{Siginfo} (C's \code{WFDB\_Siginfo}), the specification of a
  signal}
\label{tab:rust-siginfo} \\
\toprule
\tabhead Field & \tabhead Contents \\
\midrule
\endfirsthead
\toprule
\tabhead Field & \tabhead Contents \\
\midrule
\endhead
\bottomrule
\endfoot
\code{fname: Option<Vec<u8>{}>} & the name of the signal file, as the header gives it \\
\code{desc: Option<Vec<u8>{}>} & the signal description, such as \code{b"MLII"} \\
\code{units: Option<Vec<u8>{}>} & the physical units, such as \code{b"mV"}; \code{None} if
  not specified \\
\code{gain: f64} & the gain, in ADC units per physical unit (0: uncalibrated) \\
\code{initval: i32} & the value of the first sample \\
\code{group: u32} & the signal group number (the signals of one group share one file) \\
\code{fmt: i32} & the storage format, such as \code{16}, \code{212} or \code{516} \\
\code{spf: i32} & the number of samples per frame (more than 1 for the faster signals of a
  multi-frequency record) \\
\code{bsize: i32} & the block size, in bytes (0 for ordinary files) \\
\code{adcres: i32} & the ADC resolution, in bits (0: not specified) \\
\code{adczero: i32} & the ADC output for an input of 0\,V \\
\code{baseline: i32} & the ADC value that corresponds to 0 physical units \\
\code{nsamp: i64} & the number of samples (0: not specified) \\
\code{cksum: i32} & the 16-bit checksum of the samples \\
\end{xltabular}

% ------------------------------------------------------------------------------------------
\subsection{Warnings and errors}
\label{sec:rust-warnings}

\paragraph{Errors.} Each high-level function returns a \code{Result}; its error type is
\code{wfdbx::record::Error}, the one public item of the high-level API that is not
available at the root of the crate. \code{e.message()} and \code{Display} give the
message, which describes the failed operation, such as ``can't read record REC''; the
type implements \code{std::error::Error}, \code{Clone}, \code{Debug} and \code{PartialEq},
so it converts to \code{Box<dyn std::error::Error>} with the \code{?} operator. There is
one error type: the messages tell the cases apart (\cref{tab:rust-errors}). Each function
checks its arguments before it creates or changes a file, so an invalid argument leaves
the files as they were. The library may also print a message of its own on the standard
error when it cannot open a file (as in \code{init: can't open header for record REC});
\code{wfdberror()} returns the last such message, and \code{wfdbquiet()} turns these
messages off (\cref{sec:rust-library}). \code{Record::at}, like the indexing of the
result vectors, panics on an index out of range.

\begin{table}[htbp]
\centering
\caption{The kinds of errors of the high-level API}
\label{tab:rust-errors}
\begin{tabularx}{\linewidth}{@{}lL@{}}
\toprule
\tabhead Kind & \tabhead Messages \\
\midrule
files & a record, header or annotation file cannot be found, read, created or written
  completely (for example on a full disk): ``can't read record REC'', ``can't read header
  for record REC'', ``can't write header for record REC'', ``write error on record REC'',
  ``directory D does not exist'', \ldots; a signal or annotation file ends before its
  header or its end-of-file marker says (``\ldots(getvec status S)'', ``\ldots(getann
  status S)'') \\
arguments & an invalid argument or option, reported before anything is written (most
  messages of \code{write\_record} and \code{write\_annotations} begin with the function's
  name); for the metadata functions, an invalid key or value, a record name that is empty
  or \code{"-"}, a key that the header has more than once, a header without a record
  line, or a header read from a URL to be changed without \code{dir} \\
missing values & the operation needs a value that the record does not give: a time in
  seconds without a sampling frequency (``sampling frequency of record REC is not
  specified''), physical units without a gain (``signal N of record REC is
  uncalibrated'') \\
\bottomrule
\end{tabularx}
\end{table}

\paragraph{Warnings.} Warnings are printed on the standard error with the library's
messages, through \code{wfdb\_error}: \code{wfdbquiet()} silences them, and
\code{wfdberror()} returns the last one. There is no callback.

\paragraph{Settings.} Whether a function returns a result or an error, the library's
settings are those it had before the call: the input and output directories, network
access and legacy mode (set for the call by the options), and, for \code{read\_record},
\code{read\_header}, \code{read\_annotations}, \code{write\_record},
\code{write\_record\_physical}, \code{write\_annotations}, their \code{\_with} forms and
\code{check\_record}, the getvec mode (\code{getgvmode()}) and the time resolution of
output annotators (\code{getafreq()}), which opening an EDF record sets (the
\code{WFDB\_HIGHRES} bit, and the record's sampling frequency when no resolution is set)
and \code{wfdbquit()} does not reset. A later \code{write\_annotations} therefore records
no frequency that an EDF record read before it left behind; its own \code{fs} applies to
its file only.

\paragraph{The checksum warning.} When \code{read\_record} reads to the end of a record
whose samples do not match the checksums of its header, it returns all the samples and
prints the warning ``WFDB-x: checksum error in record REC: the signal data do not match
the checksums in its header''. \code{checksum\_warning: false} turns the warning off;
\code{checksum\_ok} of the result is \code{false} either way.

\paragraph{The duplicate-key warning.} When a header holds two metadata entries whose keys
differ only in letter case (written by other software), \code{read\_metadata},
\code{get\_metadata} and \code{read\_header} use the first and print the warning ``WFDB-x:
header of record REC has more than one entry for key KEY (lines L1, L2); the first is
used'' (KEY spelled as in its first entry). \code{set\_metadata} and
\code{update\_metadata} then refuse to change that key (an \code{Error}), and
\code{remove\_metadata} removes all of its entries.

\begin{lstlisting}[language=Rust,style=example]
// A record whose last sample is changed after it was written: the checksum fails.
let x: Vec<i32> = (0..100).collect();
let opt = WriteOptions {
    gains: vec![100.0],
    adcres: vec![12],
    units: vec![Some("mV".into())],
    ..Default::default()
};
write_record("ck", &x, 1, 100.0, &opt)?;
let mut dat = std::fs::read("ck.dat")?;
let n = dat.len();
dat[n - 2] = 0;
dat[n - 1] = 0;
std::fs::write("ck.dat", &dat)?;
let rec = read_record("ck", &ReadOptions::default())?; // prints the checksum warning
println!("{} {}", rec.nsamples(), rec.checksum_ok); // 100 false: all the samples
let quiet = ReadOptions { checksum_warning: false, ..Default::default() };
let rec = read_record("ck", &quiet)?; // no warning

// Two entries for one key: the first is used, with a warning.
let dup = WriteOptions {
    gains: vec![0.0],
    adcres: vec![12],
    comments: vec![" Age: 63".into(), " AGE: 64".into()],
    ..Default::default()
};
write_record("dup", &[0; 10], 1, 100.0, &dup)?;
let age = get_metadata("dup", "age")?.unwrap_or_default(); // with a warning
println!("{}", s(&age)); // 63
if let Err(e) = set_metadata("dup", "Age", "65") {
    println!("{e}");
    // header of record dup has more than one entry for key Age (lines 3, 4); remove
    // them first
}
println!("{}", remove_metadata("dup", "age")?); // 2

match read_record("nosuchrecord", &ReadOptions::default()) {
    Err(e) => println!("{}", e.message()), // can't read record nosuchrecord
    Ok(_) => {}
}
\end{lstlisting}

% ==========================================================================================
\section{The WFDB library functions in Rust}
\label{sec:rust-library}

Every function of the WFDB library is a method of \code{Wfdb} and a free function of the
crate, under its C name, with the same arguments, in the same order, and the same return
values, so that programs and documentation written for the C library carry over
directly. The functions, their arguments and their results are documented once, for all
the languages, in \cref{chap:library}; this section gives the rules by which the C
declarations become Rust functions (see also \cref{sec:library-rust}). The doc comment of
each function summarizes it, and says what WFDB-x changes in it.

\subsection{Conventions}

\paragraph{The library state.} The state of the C library (the open record, its signals,
the annotators, the WFDB path, the time base, \ldots) is held by a \code{Wfdb} value
(\code{Wfdb::new()}), and each C function is a method of it with the C name:
\code{w.isigopen(...)}, \code{w.getvec(...)}. Each method is also a free function with
the same name and arguments (\code{wfdbx::isigopen(...)}), which uses the default
\code{Wfdb} value of the calling thread; a free function must not be called from inside a
closure given to \code{wfdbx::with\_default}, which borrows that value. Separate
\code{Wfdb} values share no state, unlike the C library. The functions that use no state
are free functions only (\code{wfdbldflags}, \code{wfdbcflags}, \code{wfdbx\_write\_errors},
\code{isann}, \ldots; \cref{sec:library-rust}).

\paragraph{Names and modules.} The functions, the constants (\code{WFDB\_READ},
\code{WFDB\_INVALID\_SAMPLE}, \code{WFDB\_HIGHRES}, \ldots), the structures and the
annotation codes (\code{NORMAL}, \code{PVC}, \ldots, \cref{chap:formats}) keep their C
names, and are available at the root of the crate (\code{use wfdbx::*;}). They are
defined in the modules of \cref{tab:rust-modules}. The macros of
\file{<wfdb/ecgmap.h>} are functions (\code{isann}, \code{isqrs}, \code{map1}, \ldots,
and \code{wfdb\_isann}, \code{wfdb\_isqrs}, \ldots).

\begin{table}[htbp]
\centering
\caption{The modules of the crate \code{wfdbx}}
\label{tab:rust-modules}
\begin{tabularx}{\linewidth}{@{}lL@{}}
\toprule
\tabhead Module & \tabhead Contents \\
\midrule
\code{record} & the high-level functions of reading and writing, their option and result
  types, and \code{Error} \\
\code{check} & \code{check\_record}, \code{check\_record\_with}, \code{CheckReport},
  \code{Issue}, \code{Severity} \\
\code{metadata} & the metadata functions, \code{MetadataOptions}, \code{Metadata} \\
\code{signal} & signals and headers (port of \file{lib/signal.c}) \\
\code{annot} & annotations (\file{lib/annot.c}) \\
\code{calib} & the calibration list (\file{lib/calib.c}) \\
\code{wfdbio} & file names, the input and output directories, network access, error
  messages, \code{WfdbFile} (\file{lib/wfdbio.c}) \\
\code{wfdbinit} & \code{wfdbinit}, \code{wfdbquit}, \code{wfdbflush}
  (\file{lib/wfdbinit.c}) \\
\code{ecgcodes} & the annotation codes and the functions of \file{<wfdb/ecgmap.h>} \\
\code{types} & the types, constants and structures of \file{<wfdb/wfdb.h>} \\
\code{wfdb} & the type \code{Wfdb}, \code{with\_default}, \code{OptBytes}, \code{NULL} \\
\code{cstd} & the C standard functions that the tools use, such as a \code{printf} that
  formats numbers exactly as the GNU C library does (\code{cformat!}, \code{cprintf!}) \\
\bottomrule
\end{tabularx}
\end{table}

\paragraph{Arguments.} String arguments are byte strings: they accept \code{\&str},
\code{\&[u8]} or \code{b"..."}, and a byte string ends at its first NUL byte, as a C
string does. Where C accepts \code{NULL}, they are \code{impl OptBytes}, which also
accepts \code{Option<\&[u8]>}, \code{Option<\&str>} and \code{\&Option<Vec<u8>{}>};
\code{wfdbx::NULL} is a typed \code{None}: \code{sampfreq(NULL)} returns the frequency of
the open record, and \code{getinfo(NULL)} continues with the info strings of the last
record. Where C fills an array, the caller passes a mutable slice
(\code{isigopen(record, \&mut si, n)}, \code{getvec(\&mut v)}); an empty slice stands for
C's \code{NULL} (\code{isigopen(record, \&mut [], 0)}); \code{osigopen(record, \&mut [], 0)}
closes the open output signals (unless the record name begins with \code{+}), reads the
header, opens nothing and returns 0. Where C fills a structure, the
caller passes a mutable reference (\code{getann(0, \&mut annot)}). \code{getseginfo}
takes an \code{\&mut Option<Vec<Seginfo>{}>}, which receives a copy of the segment table.

\paragraph{Strings} are returned by value, as \code{Vec<u8>}, or as
\code{Option<Vec<u8>{}>} where the C function may return \code{NULL} (\code{anndesc},
\code{getinfo}, \code{wfdbfile}, \ldots); unlike C's pointers to static buffers, they stay
valid after later calls. Text is handled as bytes.

\paragraph{Return values and errors.} The functions return what the C functions return:
an \code{i32} status (negative on errors), a number of signals, an \code{f64}, a string
or \code{None}. Errors are reported by these values and by a message on the standard
error, as in C, never by \code{Result} or panics. \code{wfdberror()} returns the last
message; \code{wfdbquiet()} turns the messages off and \code{wfdbverbose()} on again.
\code{wfdb\_error(msg)} takes a message already formatted, with the macro
\code{cformat!} (which formats as C's \code{printf}), and \code{wfdb\_error\_fmt(fmt,
args)} formats it.

\paragraph{State.} As in C, a \code{Wfdb} value keeps its state between calls: the open
record, its input signals, the output signals and the annotators remain open until
\code{wfdbquit()} closes them. \code{wfdbquit()} leaves the input and output directories
and the network setting as they are, and so, as in C, the time resolution of output
annotators (\code{getafreq()}), which opening an EDF record sets to its sampling frequency
when none is set, and the getvec mode (\code{getgvmode()}), to which it adds
\code{WFDB\_HIGHRES}. Dropping a \code{Wfdb} value flushes its output
buffers, as C's \code{exit} does, but does not close its records (no annotation end
marks, no header updates): a program calls \code{wfdbquit()} before it ends. The
high-level functions of \cref{sec:rust-api} (except the metadata functions) call
\code{wfdbquit()} on the \code{Wfdb} value they use, which closes anything opened with the
library functions on that value, and leave those two settings as they were.

\paragraph{Additions.} \code{getvecs(buffer, n)} reads up to \code{n} sample vectors at
once, with exactly the result of \code{n} calls of \code{getvec}, but decoding whole
blocks of the signal files; \code{getvecs\_status(buffer, n)} also gives the status of
the last \code{getvec} call (\cref{sec:library-additions}):
\begin{lstlisting}[language=Rust]
pub fn getvecs(out: &mut [i32], n: i64) -> i64
pub fn getvecs_status(out: &mut [i32], n: i64) -> (i64, i32)
\end{lstlisting}
\code{WfdbFile} is C's \code{FILE *} (or \code{WFDB\_FILE}): \code{wfdb\_open} and
\code{wfdb\_fopen} return an \code{Option<WfdbFile>}, whose methods (\code{getc},
\code{read}, \code{fgets}, \code{write}, \code{seek}, \code{close}, \ldots) are those of
the C standard library.

\subsection{The structures}
\label{sec:rust-structures}

The structures of the C library have the names of the C structures without the prefix
\code{WFDB\_}, and the C field names, with C's zero values as \code{Default}
(\cref{tab:rust-structures}).

\begin{xltabular}{\linewidth}{@{}lL@{}}
\caption{The C structures in Rust}
\label{tab:rust-structures} \\
\toprule
\tabhead Structure & \tabhead Fields \\
\midrule
\endfirsthead
\toprule
\tabhead Structure & \tabhead Fields \\
\midrule
\endhead
\bottomrule
\endfoot
\code{Siginfo} & \code{fname}, \code{desc}, \code{units} (\code{Option<Vec<u8>{}>}),
  \code{gain}, \code{initval}, \code{group}, \code{fmt}, \code{spf}, \code{bsize},
  \code{adcres}, \code{adczero}, \code{baseline}, \code{nsamp}, \code{cksum}
  (\cref{tab:rust-siginfo}) \\
\code{Seginfo} & \code{recname} (the name of the segment, \code{Vec<u8>}), \code{nsamp}
  (its number of samples, \code{i64}), \code{samp0} (the number of its first sample in the
  record, \code{i64}) \\
\code{Anninfo} & \code{name} (the annotator name, \code{Vec<u8>}), \code{stat}
  (\code{i32}: \code{WFDB\_READ} or \code{WFDB\_WRITE} for the MIT format,
  \code{WFDB\_AHA\_READ} or \code{WFDB\_AHA\_WRITE} for the AHA format) \\
\code{Annotation} & \code{time} (\code{i64}, in sample intervals), \code{anntyp}
  (\code{i8}), \code{subtyp} (\code{i8}), \code{chan} (\code{u8}), \code{num}
  (\code{i8}), \code{aux} (\code{Option<Vec<u8>{}>}: absent, or bytes whose first byte is
  the length of the text that follows, as in C) \\
\code{Calinfo} & \code{low}, \code{high} (the levels of the calibration pulse, in physical
  units, \code{f64}), \code{scale} (physical units per cm), \code{sigtype} and
  \code{units} (\code{Option<Vec<u8>{}>}), \code{caltype} (\code{i32}) \\
\end{xltabular}

The integer types follow the C reference platform: \code{WFDB\_Sample} and C's
\code{int} are \code{i32}, \code{WFDB\_Time}, \code{WFDB\_Date} and C's \code{long} are
\code{i64}, \code{WFDB\_Frequency} and \code{WFDB\_Gain} are \code{f64}, C's \code{char}
is \code{i8}, and \code{WFDB\_Signal}, \code{WFDB\_Group} and \code{WFDB\_Annotator} are
\code{u32}. C's wrap-around arithmetic is reproduced.

\subsection{The settings functions}
\label{sec:rust-settings}

The functions of \cref{tab:rust-settings} are additions of WFDB-x to the library: they
set where files are read and written, network access and legacy mode, for the
high-level functions, the library functions and the tools alike
(\cref{chap:files,chap:behavior}), and give the path under which a program opens a file
that it reads or creates itself, as the tools do (\code{wfdbx\_input\_path},
\code{wfdbx\_output\_path}; \cref{sec:tool-files}). Like the other library functions, each is a method of
\code{Wfdb} and a free function, except \code{wfdbx\_write\_errors} and
\code{wfdbx\_note\_write\_error} (free functions only: the count belongs to the process).
Each \code{set} function returns the previous setting, so that it can be restored. The C
function \code{setwfdb(p)} sets the input directory too (\code{p} must name at most one
directory), and \code{getwfdb()} returns it (\code{"."} if none is set); in legacy mode
they set and return the WFDB path, as in C.

\begin{xltabular}{\linewidth}{@{}lL@{}}
\caption{The settings functions (the signatures of the free functions)}
\label{tab:rust-settings} \\
\toprule
\tabhead Function & \tabhead Effect and result \\
\midrule
\endfirsthead
\toprule
\tabhead Function & \tabhead Effect and result \\
\midrule
\endhead
\bottomrule
\endfoot
\code{wfdbx\_set\_input\_dir(dir: impl OptBytes) -> Option<Vec<u8>{}>} & sets the input
  directory: a local directory, a \code{\textasciitilde} path or a URL; \code{NULL} or
  \code{""} restores the default (\code{WFDBX\_INPUT\_DIR}, or else the current
  directory). Returns the previous input directory, or \code{None} if none was set \\
\code{wfdbx\_get\_input\_dir() -> Option<Vec<u8>{}>} & returns the input directory, or
  \code{None} if none is set \\
\code{wfdbx\_set\_output\_dir(dir: impl OptBytes) -> Option<Vec<u8>{}>} & sets the output
  directory: a local directory or a \code{\textasciitilde} path (a URL is refused, with a
  message, and the setting is unchanged); \code{NULL} or \code{""} restores the default
  (\code{WFDBX\_OUTPUT\_DIR}, or else the current directory). Returns the previous output
  directory, or \code{None} \\
\code{wfdbx\_get\_output\_dir() -> Option<Vec<u8>{}>} & returns the output directory, or
  \code{None} if none is set \\
\code{wfdbx\_set\_network(on: i32) -> i32} & turns network access on ($> 0$) or off (0);
  a negative number restores the default (on if \code{WFDBX\_NETWORK} is a nonzero
  number). Returns the previous setting, 1 or 0 \\
\code{wfdbx\_get\_network() -> i32} & returns 1 if network access is on, 0 otherwise \\
\code{wfdbx\_set\_legacy(on: i32) -> i32} & turns legacy mode on (nonzero) or off (0).
  Returns the previous setting \\
\code{wfdbx\_get\_legacy() -> i32} & returns 1 if legacy mode is on, 0 otherwise \\
\code{wfdbx\_legacy\_args(argv: \&mut Vec<Vec<u8>{}>) -> i32} & if \code{argv[1]} is
  \code{"-{}-legacy"}, removes it from the arguments, turns legacy mode on and returns 1;
  otherwise returns 0. It also enables the exit check of the tools in
  \code{cstd::exit} (an exit status of 0 becomes 1 after a write error, outside legacy
  mode) \\
\code{wfdbx\_time\_needs\_freq(string: impl AsRef<[u8]>) -> i32} & returns 1 if
  \code{strtim(string)} needs a sampling frequency, 0 otherwise \\
\code{wfdbx\_write\_errors() -> i32} & returns the number of write errors since the
  program started (output files that could not be written, flushed or closed) \\
\code{wfdbx\_note\_write\_error()} & counts a write error, for a program that writes a file
  itself \\
\code{wfdbx\_input\_path(name: impl OptBytes) -> Option<Vec<u8>{}>} & returns the path
  under which a program opens a file that it reads itself: in normal mode, a relative
  name under the input directory, a \code{\textasciitilde} name expanded, an absolute name
  as given; in legacy mode, the name as given. \code{NULL}, \code{""} and \code{"-"} are
  returned as given; \code{None} for a URL, or a name that becomes one under the input
  directory \\
\code{wfdbx\_output\_path(name: impl OptBytes) -> Option<Vec<u8>{}>} & returns the path
  under which a program creates a file that it writes itself: in normal mode, a relative
  name under the output directory (whose missing directories are then created), a
  \code{\textasciitilde} name expanded, an absolute name as given; in legacy mode, any
  name under the output directory, as a plain prefix. \code{NULL}, \code{""} and
  \code{"-"} are returned as given; \code{None} for a URL \\
\end{xltabular}

\subsection{Examples}

Reading the samples and the annotations of a record, as a C program does (the record
\code{100s} is read from the current directory):

\begin{lstlisting}[language=Rust,style=example]
use wfdbx::*;

fn main() {
    let mut si = vec![Siginfo::default(); 2]; // receives one Siginfo per signal
    let nsig = isigopen("100s", &mut si, 2); // open 2 signals
    let desc = si[0].desc.clone().unwrap_or_default();
    println!("{} {} {} {}", nsig, String::from_utf8_lossy(&desc), si[0].fmt, si[0].gain);
    // 2 MLII 212 200
    println!("{}", sampfreq(NULL)); // 360: the frequency of the open record

    let mut v = [0i32; 2]; // receives one sample per signal
    isigsettime(strtim("0:30")); // skip to 30 seconds
    getvec(&mut v);
    println!("{} {} {}", v[0], v[1], aduphys(0, v[0])); // 947 968 -0.385

    let mut frames = vec![0i32; 360 * 2]; // the next second, in one call
    let (n, status) = getvecs_status(&mut frames, 360);
    println!("{n} {status}"); // 360 2

    let ai = Anninfo { name: b"atr".to_vec(), stat: WFDB_READ };
    annopen("100s", &[ai], 1);
    let mut a = Annotation::default();
    let mut count = 0;
    while getann(0, &mut a) == 0 {
        // 0: an annotation was read; -1: the end
        if count < 2 {
            let t = mstimstr(a.time);
            println!("{} {}", String::from_utf8_lossy(&t), String::from_utf8_lossy(&annstr(a.anntyp as i32)));
        }
        count += 1;
    }
    println!("{count} annotations"); // 75 annotations

    wfdbquiet(); // no messages on the standard error
    let stat = isigopen("nosuchrecord", &mut [], 0);
    print!("{} {}", stat, String::from_utf8_lossy(&wfdberror()));
    // -1 init: can't open header for record nosuchrecord
    wfdbverbose();
    wfdbquit(); // close everything
}
\end{lstlisting}

Writing a record and an annotation file, sample by sample and annotation by annotation:

\begin{lstlisting}[language=Rust,style=example]
use wfdbx::*;

fn main() {
    let out = Siginfo {
        fname: Some(b"ramp.dat".to_vec()),
        desc: Some(b"ramp".to_vec()),
        units: Some(b"mV".to_vec()),
        gain: 200.0,
        fmt: 16,
        adcres: 12,
        ..Default::default()
    };
    println!("{}", osigfopen(&[out], 1)); // 1: one output signal
    setsampfreq(250.0);
    for k in 0..500 {
        putvec(&[k]);
    }
    println!("{}", newheader("ramp")); // 0: ramp.hea is written
    let ai = Anninfo { name: b"qrs".to_vec(), stat: WFDB_WRITE };
    annopen("ramp", &[ai], 1);
    for t in [100, 350] {
        let a = Annotation { time: t, anntyp: NORMAL as i8, ..Default::default() };
        putann(0, &a);
    }
    wfdbquit(); // completes ramp.dat and ramp.qrs
}
\end{lstlisting}

% ==========================================================================================
\section{The command-line tools in Rust}
\label{sec:rust-tools}

The 33 command-line tools (\code{rdsamp}, \code{rdann}, \code{wrsamp}, \code{xform},
\ldots) are separate programs of the crate, one per file of \file{rust/src/bin}, built in
\file{rust/target/release/} (or installed in \file{\textasciitilde/.cargo/bin} by
\cmd{cargo install}), and run by their names. Each follows its C program line by line.
The tools, their options, their output and their exit status are those of the C tools,
and are documented in \cref{chap:tools}.

\begin{lstlisting}[style=shell]
export PATH=$PWD/rust/target/release:$PATH
export WFDBX_INPUT_DIR=tests/data     # run from the root of the repository
rdsamp -r 100s -t 10 -p     # the first 10 seconds of 100s, in physical units
rdann -r 100s -a atr        # its annotations
rdsamp --legacy -r 100s     # legacy mode: --legacy as the first argument
\end{lstlisting}

There is no launcher: a tool is run by its own name, and it has no \opt{-{}-version}
option, like the C tools. The tools have the same names as the original WFDB tools: if
both are installed, the one that comes first in the \code{PATH} runs. When an output file
or the standard output cannot be written (a full disk, for example), a tool reports it
(``PROGRAM: write error on standard output'') and exits with status 1; with
\opt{-{}-legacy}, it keeps the original exit status. A tool whose standard output is a
closed pipe ends with status 141, as a C program ended by \code{SIGPIPE}.

\paragraph{Environment variables.} The tools, like the rest of the library, use the
environment variables of \cref{tab:rust-environment}, also for the files that they read
or create themselves (\cref{sec:tool-files}): the \opt{-i} file of \cmd{wrsamp},
\cmd{edf2mit}, \cmd{wav2mit}, \cmd{a2m} and \cmd{ad2m}, the input files of
\cmd{ahaecg2mit} and the \opt{-S} script of \cmd{xform} are read from the input directory,
and the WAV file of \cmd{mit2wav}, the EDF file of \cmd{mit2edf}, the master header of
\cmd{snip -m} and the \file{.des}, \file{.ecg}, \file{.key} and \file{.txt} files of
\cmd{parsescp} are created in the output directory. Such a file cannot be read from a URL.
In legacy mode, the files that a tool reads itself are opened as named, and the output
directory, if set, is a plain prefix of every output name.

\begin{table}[htbp]
\centering
\caption{The environment variables used by the Rust library}
\label{tab:rust-environment}
\begin{tabularx}{\linewidth}{@{}lL@{}}
\toprule
\tabhead Variable & \tabhead Meaning \\
\midrule
\code{WFDBX\_INPUT\_DIR} & the input directory (a directory, a \code{\textasciitilde}
  path or a URL), when \code{wfdbx\_set\_input\_dir} does not set one \\
\code{WFDBX\_OUTPUT\_DIR} & the output directory, when \code{wfdbx\_set\_output\_dir} does
  not set one \\
\code{WFDBX\_NETWORK} & network access is on if it is a nonzero number, when
  \code{wfdbx\_set\_network} does not set it \\
\code{WFDBANNSORT} & \code{0} to keep annotation files in the order written (no sorting
  when they are closed) \\
\code{WFDBGVMODE} & the default getvec mode (\code{setgvmode} with a negative mode) \\
\code{WFDBCAL} & the calibration file read by \code{calopen(NULL)} (default
  \file{wfdbcal}) \\
\code{WFDB} & the WFDB path, used in legacy mode only (outside legacy mode, it is
  ignored, with a message) \\
\code{HOME} (\code{USERPROFILE} on Windows) & the home directory, for
  \code{\textasciitilde} paths \\
\bottomrule
\end{tabularx}
\end{table}

Tools that write signal parameters require them as options: \code{wrsamp} (with
\opt{-o}) requires \opt{-F} (frequency), \opt{-G} (gain), \opt{-R} (ADC resolution) and,
for calibrated signals, \opt{-U} (units); \code{wav2mit} requires \opt{-G}, and \opt{-U}
unless the gain is 0; \code{ad2m} and \code{ahaecg2mit} require \opt{-F}, \opt{-G},
\opt{-R}, and \opt{-U} unless the gain is 0; \code{a2m}, \code{rdann}, \code{ann2rr},
\code{sortann} and \code{m2a} require \opt{-F} when the record's header gives no
frequency (\cref{chap:tools}).

% ==========================================================================================
\section{Testing}
\label{sec:rust-testing}

The tests of the Rust implementation are the unit tests of the modules, the integration
tests of \file{rust/tests} and the examples of the doc comments (which read the sample
record \code{100s} from \file{tests/data}). From the directory \file{rust}:

\begin{lstlisting}[style=shell]
cargo test --release                       # all the tests, and the doc examples
cargo test --release --test metadata       # one file of rust/tests
cargo test --release shared_cases          # selected tests, by name
cargo clippy --release --all-targets       # lints
\end{lstlisting}

The integration tests cover the WFDB library functions (\file{signals.rs},
\file{annotations.rs}, \file{wfdbio.rs}, \file{bulk.rs}, the C standard functions in
\file{cstd.rs}), finding files and the settings (\file{files.rs}), network access
(\file{network.rs}), the high-level functions, the metadata functions (\file{metadata.rs},
with the cases shared by the seven languages, in \file{tests/metadata}) and
\code{check\_record} (\file{check.rs}). They run in temporary copies of the test records
of \file{tests/data}. The tests of network access use a local HTTP server on 127.0.0.1
only, and need \cmd{curl}. The tests of \code{check\_record} write the corpus of
\file{python/tests/check\_corpus.py} with \cmd{python3} (or the program named by the
environment variable \code{PYTHON}), NumPy and the Python package of \file{python/src},
and compare the reports with its \file{expected.txt}.
\cmd{cargo run -{}-release -{}-example check\_cases DIR} prints the reports for a corpus in
\file{DIR}.

The cross-language conformance suite runs every tool and compares its results, byte for
byte, with those of the original C programs (1135 cases), and, with \opt{-{}-legacy}, in
legacy mode (883 cases). It runs the tools of \file{rust/target/release}:

\begin{lstlisting}[style=shell]
cargo build --release --manifest-path rust/Cargo.toml
PYTHONPATH=$PWD/python/src python3 tests/conformance/harness.py check --impl rust -j8
PYTHONPATH=$PWD/python/src python3 tests/conformance/harness.py check --impl rust -j8 --legacy
\end{lstlisting}

% The Julia chapter of the WFDB-x user's manual. It follows the sections of the Python
% chapter (chapters/python.tex), in the same order. Each \apientry of the section "The
% high-level API" is a function exported by the module WFDBX and defined in
% julia/src/highlevel.jl, julia/src/check.jl or julia/src/metadata.jl; its signature
% listing is copied from the method definition that takes the library state w::Wfdb first
% (each such function also has a method without w, with the same other arguments), and its
% params name every argument and keyword argument of the definition except w, in order.
% The examples of style "example" were run (each in a new module, after "using WFDBX", in
% a fresh copy of tests/data that is the current directory); check_manual.py does not run
% them.
\chapter{Julia}
\label{chap:julia}
\chapterlanguage{julia}

The Julia implementation of WFDB-x is the package \code{WFDBX} (module \code{WFDBX}), in
the directory \file{julia} of the repository: a library and the 33 command-line tools,
run through one driver script. It is a complete Julia implementation of the WFDB Software
Package: it uses only Julia and its standard libraries (no other package), and never calls
the C library.

It offers WFDB-x in three ways:
\begin{itemize}
\item \textbf{The high-level API} (\cref{sec:julia-api}): functions that read and write
  whole records, headers, annotation files and header metadata as Julia matrices, vectors
  and structures, in one call, with keyword arguments, and that throw a \code{WfdbError}
  (or, for some invalid arguments, an \code{ArgumentError}) on errors. They read signal
  files in large blocks, and are the fastest way to read a record. They are an addition
  of WFDB-x: the C package has no equivalent. They are the recommended interface for new
  code, analysis and notebooks.
\item \textbf{The WFDB library functions} under their C names (\code{isigopen},
  \code{getvec}, \code{getann}, \code{putann}, \ldots), with the behavior of the C library
  (\cref{sec:julia-library}). They are documented once, for all the languages, in
  \cref{chap:library}.
\item \textbf{The command-line tools} (\code{rdsamp}, \code{rdann}, \code{wrsamp},
  \code{xform} and 29 others), with the options, output and exit status of the C tools,
  run through the driver script \file{julia/bin/wfdbx.jl} (\cref{sec:julia-tools}). They
  are documented in \cref{chap:tools}.
\end{itemize}
The high-level functions return exactly the samples and annotations that the library
functions (and the C library) give.

% ==========================================================================================
\section{Installation and requirements}
\label{sec:julia-install}

\subsection{Requirements}

WFDB-x for Julia needs Julia 1.9 or later, on Linux, macOS or Windows (the Windows code
paths are not tested). It has no dependency beyond Julia's standard libraries: no WFDB
library and no libFLAC (the FLAC formats 508, 516 and 524 are read and written by
WFDB-x's own codec). Network access (\cref{sec:network}) uses the standard library
\code{Downloads}, which is loaded only when the first URL is read. The unit tests use the
standard libraries \code{Test} and \code{Sockets}, and the tests of \code{check\_record}
also run \cmd{python3}, with NumPy (\cref{sec:julia-testing}).

\subsection{Installing the package}

The package is installed from a copy of the repository, with the package manager:

\begin{lstlisting}[language=Julia]
using Pkg
Pkg.develop(path="/path/to/wfdb-x/julia")
\end{lstlisting}

In the package REPL, \code{] dev /path/to/wfdb-x/julia} does the same. The package is
\code{WFDBX}, version \code{1.2.0}, under the GNU General Public License, version 3 or
later, and is loaded with \code{using WFDBX}. The first \code{using WFDBX} precompiles
the package, running a workload of every part of the library and of every tool, so that
later sessions and tool runs load it without compiling anything. Changing a source file
of the package makes the next session precompile it again.

\subsection{The command-line tools}

The tools are run through the driver script \file{julia/bin/wfdbx.jl}: its first argument
names the tool, and the other arguments go to the tool. A small shell script, saved as
\file{wfdbx} in a directory of the \code{PATH}, makes the command shorter
(\cref{sec:julia-tools}):

\begin{lstlisting}[style=shell]
#!/bin/sh
# wfdbx: run a WFDB-x tool, as in `wfdbx rdsamp -r 100s -t 10`
exec julia --startup-file=no --project=/path/to/wfdb-x/julia \
    /path/to/wfdb-x/julia/bin/wfdbx.jl "$@"
\end{lstlisting}

\cmd{wfdbx -{}-version} then prints the version of WFDB-x, \code{WFDB-x 1.2.0}. The
docstrings of the package document each function (\code{?read\_record} in the REPL), with
examples.

\subsection{The sample record}
\label{sec:julia-resources}

The repository includes the sample record \code{100s}, the sample record of the WFDB
Software Package: the first minute of record 100 of the MIT-BIH Arrhythmia Database, two
ECG signals (\code{MLII} and \code{V5}) sampled at 360\,Hz, with the reference beat
annotations (annotator \code{atr}). Its files (\file{100s.hea}, \file{100s.dat},
\file{100s.atr}) are in the directory \file{tests/data} of a copy of the repository, with
the other test records of WFDB-x. The package does not install them: programs name that
directory as their input directory (\code{wfdbx\_set\_input\_dir("/path/to/wfdb-x/tests/data")},
or the keyword argument \code{input\_dir}). Work that uses the MIT-BIH Arrhythmia
Database cites it as described on PhysioNet.

% ==========================================================================================
\section{Quick start}
\label{sec:julia-quickstart}

This program reads the header, ten seconds of the signals and the annotations of the
sample record \code{100s}, from the directory \file{tests/data} of the repository (the
input directory of each call), and computes the mean heart rate. It then writes the first ten seconds
as a new record, \code{first10}, in the current directory, with a metadata entry in its
header.

\begin{lstlisting}[language=Julia,style=example]
using WFDBX

# The sample record 100s, in the directory tests/data of the repository.
const DATA = "tests/data"

# The header only: no samples are read.
hdr = read_header("100s"; input_dir=DATA)
println(hdr.nsig, " ", hdr.fs, " ", hdr.nsamples)     # 2 360.0 21600

# Seconds 10 to 20 of both signals, in ADC units (one column per signal).
rec = read_record("100s"; from="0:10", to="0:20", input_dir=DATA)
println(size(rec.samples))                          # (3600, 2)
x = physical(rec)                                   # in physical units
println(rec.samples[1, 1], " ", x[1, 1])            # 946 -0.39

# The beat annotations, and the mean heart rate.
ann = read_annotations("100s", "atr"; input_dir=DATA)
beats = ann.samples[[isqrs(c) != 0 for c in ann.codes]]
rr = (beats[end] - beats[1]) / ann.fs / (length(beats) - 1)
println(length(beats), " beats, ", round(60 / rr; digits=1), " beats per minute")
# 74 beats, 73.9 beats per minute

# Write the first 10 seconds as the record "first10" (first10.hea and first10.dat, in
# the current directory, the output directory), with the signal specifications of 100s.
r10 = read_record("100s"; to="0:10", input_dir=DATA)
write_record("first10", r10.samples, r10.fs; gains=200.0, baselines=1024, adczeros=1024,
             adcres=11, units="mV", names=signal_names(r10),
             metadata=["Source" => "MIT-BIH record 100"])

# Header metadata: "# Key: value" lines of the header (keys and values are bytes).
set_metadata("first10", "Age", "69")
for (key, value) in read_metadata("first10")
    println(String(key), " = ", String(value))
end
# Source = MIT-BIH record 100
# Age = 69
println(check_record("first10"))                    # no problems found
\end{lstlisting}

The program runs from the root of the repository. The input directory is given to each
call that reads \code{100s} (\code{input\_dir}), so the record \code{first10}, written to
the current directory, is then read from it under its plain name. The header
\file{100s.hea} gives no units: the program writes \code{"mV"} explicitly, since WFDB-x
never assumes units, a gain, an ADC resolution or a sampling frequency that a record does
not give (\cref{chap:behavior}).

% ==========================================================================================
\section{The high-level API}
\label{sec:julia-api}

The high-level API is made of the functions of this section, defined in the files
\file{highlevel.jl}, \file{check.jl} and \file{metadata.jl} of \file{julia/src} and
exported by the module \code{WFDBX}, of the types of their results
(\cref{sec:julia-results}), and of the exception type \code{WfdbError}
(\cref{sec:julia-warnings}). The docstrings of the package document each of them.

\subsection{Conventions}
\label{sec:julia-conventions}

\paragraph{The library state: \code{[w::Wfdb]}.} The state of the library (the open
record, the annotators, the input and output directories, network access, legacy mode) is
held by an object of the abstract type \code{Wfdb}; \code{Wfdb()} creates a new one, and
\code{lib()} returns the default object of the process. Each high-level function that uses
the state (all but \code{physical}, \code{signal\_names} and \code{lines}) has two
methods: one that takes a \code{Wfdb} object \code{w} as its first argument, such as
\code{read\_record(w, "100s")}, which uses the state of \code{w}, and one without it, with
the same other arguments, such as \code{read\_record("100s")}, which uses \code{lib()}.
The docstrings write this optional first argument \code{[w::Wfdb]}. The entries of this
section show the method that takes \code{w}, as it is defined; the method without it is
defined as \code{read\_record(record::ByteStr; kw...) = read\_record(lib(), record;
kw...)}. Separate \code{Wfdb} objects share no state (except what belongs to the process:
the count of write errors and the standard streams), so several records can be read at
once.

\begin{lstlisting}[language=Julia,style=example]
w = Wfdb()                              # a library state of its own
wfdbx_set_input_dir(w, "in")            # (the input directory of w only)
mkpath("in"); cp("100s.hea", "in/100s.hea"); cp("100s.dat", "in/100s.dat")
a = read_record(w, "100s")              # reads in/100s.hea, on w
b = read_record("100s")                 # reads ./100s.hea, on lib()
println(a.samples == b.samples)         # true
\end{lstlisting}

\paragraph{Keyword arguments.} The options of a function are keyword arguments, after a
semicolon, such as \code{read\_record("100s"; from="0:10", highres=true)}; each has a
default, given in the signature. An option whose default is \code{nothing} means ``not
given'': the default behavior or the current setting.

\paragraph{Record names and directories.} A record is named as in the other parts of
WFDB-x (\cref{chap:records,chap:files}). A full path, a \code{\textasciitilde} path or a
URL is used as given; any other name, such as \code{"100s"} or \code{"mitdb/100"}, is read
from the input directory and written to the output directory. Both are the current
directory unless they are set, with \code{wfdbx\_set\_input\_dir} and
\code{wfdbx\_set\_output\_dir} (\cref{sec:julia-settings}) or with the environment
variables \code{WFDBX\_INPUT\_DIR} and \code{WFDBX\_OUTPUT\_DIR}. So a record written under
a plain name is written to the output directory and read again from the input directory:
the two are the same unless one of them is set. URLs are read only when network access is
on. The readers take the keyword arguments \code{input\_dir} and \code{network}, and the
writers \code{dir}: they apply to that one call, and the previous settings are restored
when it returns, also when it throws. The signal files that a header names are read from
the header's directory.

A record name is given without the \file{.hea} extension (the readers, and
\code{write\_annotations}, also accept a trailing \code{.hea}; \code{write\_record} does
not). It is a full path, absolute (\code{"/data/mitdb/100"}) or in the home directory
(\code{"\textasciitilde/Documents/100s"}), used as given; a relative path
(\code{"100s"}, \code{"data/100s"}), read from the input directory if one is set,
otherwise from the current directory, and written to the output directory if one is set,
otherwise to the current directory (missing directories are created); or a URL
(\code{"https://physionet.org/files/mitdb/1.0.0/100"}), read only if network access is
on (a URL cannot be written). The \code{input\_dir} option applies to relative names
only, not to full paths or URLs. A writer given a name with a path creates the files in
that directory, and the header names the record and its signal file without the path
(\code{rec}, \code{rec.dat}), so that the record opens under any path that leads to it;
in legacy mode, as in the original, the header gives the name with its path. With the
\code{dir} option, the name must be a plain name: the files are created in \code{dir},
and the record opens from that directory under its plain name.

\paragraph{Times.} The keyword arguments \code{from} and \code{to} of \code{read\_record}
and \code{read\_annotations} are sample numbers (whole numbers, such as \code{3600} or
\code{3600.0}) or WFDB time strings, as \code{strtim} reads them (\code{"0:10"}, 10
seconds; \code{"1:02:03.5"}; \code{"s3600"}, sample 3600; \code{"e"}, the end of the
record; or a time of day in brackets, \code{"[14:00:00]"}). A time in seconds needs the
sampling frequency; when the record does not give it, such a time is an error.

\paragraph{Samples.} Samples are \code{Int32} values (C's \code{WFDB\_Sample}), in ADC
units, in a \code{Matrix\{Int32\}} with one row per sample and one column per signal:
\code{samples[i, s]} is sample \code{i} of signal \code{s} (both counted from 1).
Invalid samples are \code{WFDB\_INVALID\_SAMPLE} ($-32768$), and \code{NaN} in physical
units. Sample numbers and times are \code{Int64}, counted from 0 as in C. Signal numbers
are counted from 0 where the API gives them as numbers (\code{channels}, and the
\code{channels} field of a \code{Record}), as in C and the tools.

\paragraph{Per-signal and per-annotation values.} The per-signal keyword arguments of
\code{write\_record} (\code{gains}, \code{units}, \ldots) take one value for all the
signals, or a vector with one value per signal; the per-annotation fields of
\code{Annotations} (\code{symbols}, \code{chans}, \ldots) are empty or hold one value per
annotation. Any other length is an error, reported before anything is written.

\paragraph{Signal parameters.} No sampling frequency, gain, ADC resolution or units is
ever assumed. Where one is needed and the record does not give it, it must be given;
otherwise the function throws a \code{WfdbError}. A record whose header gives no sampling
frequency has \code{fs} 0, a signal without an ADC resolution has \code{adcres} 0, and a
signal without units has \code{units} \code{nothing}.

\paragraph{Text.} Text read from files (signal descriptions, units, info strings,
metadata, aux strings, mnemonics) is given as bytes, \code{Vector\{UInt8\}}, since WFDB
files are not UTF-8: every byte sequence round-trips unchanged. \code{String(x)} (or
\code{String.(v)} for a vector) converts bytes for printing; \code{String(x)} takes the
vector it converts, so \code{String(copy(x))} keeps \code{x}. Arguments accept
\code{String}s or bytes (the type \code{ByteStr}, \code{Union\{AbstractString,
AbstractVector\{UInt8\}\}}; \code{OptByteStr} also accepts \code{nothing}); a
\code{String} is taken as its bytes (UTF-8), as given. The names of records
(\code{Record.name}, \code{Header.name}, \ldots) and the time strings are \code{String}s.

\paragraph{Legacy mode.} With \code{legacy=true}, a function behaves as in the original
WFDB package, with its assumed signal parameters and its search for files
(\cref{chap:behavior}); the library's legacy mode is on for the duration of the call
only. With \code{false} (the default), the current setting of \code{wfdbx\_set\_legacy}
applies, which is off unless it was turned on.

\paragraph{State.} The functions that read and write records and annotation files call
\code{wfdbquit()} on the \code{Wfdb} object they use, before and after their work, also
when they throw, which closes anything opened with the WFDB library functions on that
object. The library's settings (legacy mode, getvec mode, time resolution of output
annotators, input and output directories, network access) are the same after a call as
before it: the high-level functions and \code{check\_record} save \code{getgvmode()} and
\code{getafreq()} when they start and restore them when they return, also when they
throw, since opening an EDF record adds \code{WFDB\_HIGHRES} to the getvec mode and sets
the time resolution to the record's sampling frequency (when none is set), as in C, and
\code{wfdbquit()} resets neither. The C-style functions keep the behavior of C. The metadata functions do not
call \code{wfdbquit()}: they open only the header, and leave the library's open files as
they are.

\paragraph{The examples} of this section run, each in a new module after \code{using
WFDBX}, in a copy of the test records of \file{tests/data} (among them \code{100s}),
which is the current directory and holds the records that the examples write; the input
and output directories are not set. The quick start (\cref{sec:julia-quickstart}) runs
instead from the root of the repository.

% ------------------------------------------------------------------------------------------
\subsection{Reading records}
\label{sec:julia-reading}

\apientry{read\_record}{Read the signals of a record}
\begin{lstlisting}[language=Julia,style=signature]
read_record(w::Wfdb, record::ByteStr; channels=nothing, from=0, to=nothing,
            highres::Bool=false, pad::Bool=false, input_dir::OptByteStr=nothing,
            network::Union{Nothing,Bool}=nothing, legacy::Bool=false,
            checksum_warning::Bool=true)
\end{lstlisting}
Reads the signals of a record from \code{from} up to \code{to} (exclusive), in ADC units,
together with the description of the record. The samples are those that \code{getvec}
returns (\cref{chap:library}), read in large blocks with \code{getvecs}. If the record
ends before \code{to}, the samples up to its end are returned. When the read reaches the
end of the record, the samples are compared with the checksums of the header
(\cref{sec:julia-warnings}). The getvec mode (\code{getgvmode()}) and the time resolution
of output annotators (\code{getafreq()}) are the same after the call as before it:
\code{highres} and \code{pad} apply to this call only.
\begin{params}
\item[record] Record name, without the \file{.hea} extension (a trailing \code{".hea"} is
  also accepted): a full path, absolute (\code{"/data/mitdb/100"}) or in the home
  directory (\code{"\textasciitilde/Documents/100s"}), used as given; a relative path
  (\code{"100s"}, \code{"data/100s"}), read from the input directory if one is set,
  otherwise from the current directory; or a URL
  (\code{"https://physionet.org/files/mitdb/1.0.0/100"}), read only if network access is
  on.
\item[channels] The signals to read, in the order given: a vector of signal numbers
  (\code{0} is the first signal) or names, such as \code{["V5"]} (as
  \code{signal\_names} gives them). \code{nothing}: every signal.
\item[from] The first sample: a sample number ($\ge 0$) or a time string. Default 0.
\item[to] The sample after the last one: a sample number or a time string.
  \code{nothing}: the end of the record.
\item[highres] \code{true} to read a multi-frequency record at its highest sampling
  frequency (\code{fs} of the result is then that frequency); \code{false}: at its frame
  frequency, as \code{getvec} does in low-resolution mode.
\item[pad] \code{true} to replace each invalid sample by the previous valid sample of
  the same signal (getvec mode \code{WFDB\_GVPAD}).
\item[input\_dir] The input directory for this call only (optional): a local directory,
  a \code{\textasciitilde} path or a URL, under which relative record names are read. It
  does not apply to full paths or URLs. \code{nothing}: the current setting.
\item[network] \code{true} to allow reading URLs for this call only, \code{false} to
  forbid it. \code{nothing}: the current setting.
\item[legacy] \code{true} to behave as the original WFDB package, for this call only.
\item[checksum\_warning] \code{true} (the default) to log a warning when the data do not
  match the checksums of the header; \code{false}: no warning.
\end{params}
\returns{A \code{Record} (\cref{tab:julia-record}), with the samples in ADC units.}
\raises{A \code{WfdbError} if the record cannot be read or has no signals (``can't read
record REC''), or its signals cannot be opened (``can't open signals of record REC''); if
a requested signal does not exist (``record REC has no signal C''); if \code{from} is
negative (``read\_record: from must not be negative''), if \code{from} or \code{to} is a
number that is not whole (``read\_record: from must be a whole number'') or beyond the
range of a 64-bit integer, or is neither a number nor a string (``read\_record: from must
be a sample number or a time string''); if \code{from} cannot be reached (``can't seek
to sample N of record REC''); if a time string needs the sampling frequency and the
header gives none (``sampling frequency of record REC is not specified''); or if a signal
file is shorter than its header says, or cannot be read or decoded (``can't read signal
data of record REC (getvec status S)''). Data that do not match the checksums of the
header are returned, with a warning (see \code{checksum\_warning}).}
\begin{lstlisting}[language=Julia,style=example]
# Seconds 10 to 20 of lead MLII, in physical units.
rec = read_record("100s"; channels=["MLII"], from="0:10", to="0:20")
mlii = physical(rec)[:, 1]
println(length(mlii), " ", rec.fs, " ", rec.sampfrom, " ", String(signal_names(rec)[1]))
# 3600 360.0 3600 MLII

# The first 3 samples of both signals, in ADC units, V5 first.
rec = read_record("100s"; channels=[1, 0], to=3)
println(rec.samples)                                # Int32[1011 995; 1011 995; 1011 995]
println(rec.channels, " ", rec.siginfo[1].gain, " ", rec.siginfo[1].baseline)
# [1, 0] 200.0 1024

# The whole record: the checksums are verified.
rec = read_record("100s")
println(size(rec.samples, 1), " ", rec.checksum_ok)  # 21600 true
\end{lstlisting}
A multi-frequency record (here the test record \code{mfreq}, whose signals have 1, 2
and 4 samples per frame at a frame frequency of 100\,Hz) is read at its frame frequency,
or at its highest sampling frequency:
\begin{lstlisting}[language=Julia,style=example]
rec = read_record("mfreq")
println(rec.fs, " ", size(rec.samples))             # 100.0 (600, 3)
rec = read_record("mfreq"; highres=true)
println(rec.fs, " ", size(rec.samples))             # 400.0 (2400, 3)
println([si.spf for si in rec.siginfo])             # Int32[1, 2, 4]
\end{lstlisting}

\apientry{read\_header}{Read the header of a record}
\begin{lstlisting}[language=Julia,style=signature]
read_header(w::Wfdb, record::ByteStr; input_dir::OptByteStr=nothing,
            network::Union{Nothing,Bool}=nothing, legacy::Bool=false)
\end{lstlisting}
Reads the header of a record, without opening its signal files. The result describes the
record and each of its signals, and holds the info strings and the metadata entries of the
header (from the \file{.hea} file only, as \code{read\_metadata} reads them, with its
warning for a key given more than once). For an EDF record, which has no \file{.hea}
file, and for a header read from the standard input (\code{"-"}), the metadata is empty,
without an error. For a multi-segment record, the master header is read, and
\code{nsamples} is the length of the whole record.
\begin{params}
\item[record] Record name, without the \file{.hea} extension (a trailing \code{".hea"} is
  also accepted): a full path, absolute (\code{"/data/mitdb/100"}) or in the home
  directory (\code{"\textasciitilde/Documents/100s"}), used as given; a relative path
  (\code{"100s"}, \code{"data/100s"}), read from the input directory if one is set,
  otherwise from the current directory; or a URL
  (\code{"https://physionet.org/files/mitdb/1.0.0/100"}), read only if network access is
  on.
\item[input\_dir] The input directory for this call only (optional), as for
  \code{read\_record}: relative record names are read under it. It does not apply to full
  paths or URLs.
\item[network] Network access for this call only, as for \code{read\_record}.
\item[legacy] \code{true} to read the header as the original WFDB package does, for this
  call only.
\end{params}
\returns{A \code{Header} (\cref{tab:julia-header}).}
\raises{A \code{WfdbError} if the header cannot be read (``can't read header for record
REC'').}
\begin{lstlisting}[language=Julia,style=example]
hdr = read_header("100s")
println(hdr.name, " ", hdr.fs, " ", hdr.nsig, " ", hdr.nsamples)   # 100s 360.0 2 21600
for si in hdr.siginfo
    println(String(si.desc), " ", String(si.fname), " ", si.fmt, " ", si.gain, " ",
            si.adcres, " ", si.units)
end
# MLII 100s.dat 212 200.0 11 nothing
# V5 100s.dat 212 200.0 11 nothing
println("[", String(hdr.comments[1]), "]")          # [ 69 M 1085 1629 x1]
println(hdr.base_time === nothing)                  # true: the header gives no start time
key, value = hdr.metadata[1]
println("[", String(key), "] [", String(value), "]")
# [Produced by xform from record 100, beginning at 0] [0]

# A copy of the header in the directory "in", read for this call only.
mkpath("in"); cp("100s.hea", "in/100s.hea")
println(read_header("100s"; input_dir="in").nsig)   # 2
\end{lstlisting}

\apientry{read\_annotations}{Read an annotation file}
\begin{lstlisting}[language=Julia,style=signature]
read_annotations(w::Wfdb, record::ByteStr, annotator::ByteStr; from=nothing,
                 to=nothing, fs::Union{Nothing,Real}=nothing,
                 input_dir::OptByteStr=nothing, network::Union{Nothing,Bool}=nothing,
                 legacy::Bool=false)
\end{lstlisting}
Reads the annotations of an annotation file, in file order, as parallel vectors. The
sampling frequency, needed to convert times in seconds, is that of the record's header,
or, for a record without a header or whose header gives none, the keyword argument
\code{fs}, as the option \opt{-F} of \code{rdann}.
\begin{params}
\item[record] Record name, without the \file{.hea} extension (a trailing \code{".hea"} is
  also accepted): a full path, absolute (\code{"/data/mitdb/100"}) or in the home
  directory (\code{"\textasciitilde/Documents/100s"}), used as given; a relative path
  (\code{"100s"}, \code{"data/100s"}), read from the input directory if one is set,
  otherwise from the current directory; or a URL
  (\code{"https://physionet.org/files/mitdb/1.0.0/100"}), read only if network access is
  on. Its header, if there is one, gives the sampling
  frequency.
\item[annotator] Annotator name: the suffix of the annotation file. \code{"atr"} reads the
  file \file{RECORD.atr}.
\item[from] The first time to read: a sample number or a time string. \code{nothing}: the
  beginning of the file.
\item[to] The time at which to stop: annotations at or after it are not read.
  \code{nothing}: the end of the file.
\item[fs] The sampling frequency, in Hz ($> 0$), to use when the record has no header or
  its header gives none. \code{nothing}: none.
\item[input\_dir] The input directory for this call only (optional), as for
  \code{read\_record}: relative record names are read under it. It does not apply to full
  paths or URLs.
\item[network] Network access for this call only, as for \code{read\_record}.
\item[legacy] \code{true} to behave as the original WFDB package (which assumes 250\,Hz
  for a record without a frequency), for this call only.
\end{params}
\returns{An \code{Annotations} (\cref{tab:julia-annotations}). Its \code{fs} is the
frequency in which the times are counted: the one written in the annotation file, if
there is one, or else the sampling frequency; it is 0 if none of them is known.}
\raises{A \code{WfdbError} if \code{fs} is given and is not $> 0$ (``read\_annotations:
fs must be > 0''); if the annotation file cannot be opened (``can't read annotator ANN for
record REC''); if a time string needs the sampling frequency and none is known
(``sampling frequency of record REC is not specified''); if \code{from} cannot be reached
(``can't seek in annotator ANN''); or if the file ends early or cannot be read before
\code{to} is reached (``can't read annotation file ANN of record REC (getann status S)'',
where S is the status that \code{getann} returned, such as $-3$ for an unexpected end of
file).}
\begin{lstlisting}[language=Julia,style=example]
ann = read_annotations("100s", "atr"; from="0:10", to="0:30")
println(length(ann), " ", ann.fs)                   # 24 360.0
println(ann.samples[1], " ", String(ann.symbols[1])) # 3862 N

# The whole file: the first annotation is a rhythm change, with an aux string.
ann = read_annotations("100s", "atr")
println(String(ann.symbols[1]), " ", ann.codes[1] == WFDBX.RHYTHM, " ", String(ann.aux[1]))
# + true (N

# A record without a header: fs gives the frequency of times in seconds (the library
# reports on the standard error that it cannot open the header).
write_annotations("beats", "qrs", Annotations(samples=[100, 460, 820],
                                              symbols=[b"N", b"N", b"N"]))
ann = read_annotations("beats", "qrs"; to="0:02", fs=360)
println(length(ann), " ", ann.fs)                   # 2 360.0
\end{lstlisting}

\apientry{physical}{Convert the samples of a record to physical units}
\begin{lstlisting}[language=Julia,style=signature]
physical(rec::Record; legacy::Bool=false)
\end{lstlisting}
Returns the samples of a \code{Record} in physical units, computed as $(\mathit{adc} -
\mathit{baseline}) / \mathit{gain}$ for each signal, with the gain and baseline of its
\code{siginfo}. Invalid samples become \code{NaN}. The result has the shape of
\code{rec.samples}. It uses no library state.
\begin{params}
\item[rec] The record, as \code{read\_record} returns it.
\item[legacy] \code{true} to behave as the original WFDB package, which converts the
  samples of an uncalibrated signal (gain 0) with a gain of 200. With \code{false}, the
  legacy mode of the default library object (\code{wfdbx\_set\_legacy}) decides.
\end{params}
\returns{A \code{Matrix\{Float64\}}, one row per sample and one column per signal.}
\raises{A \code{WfdbError} if the gain of a signal is 0 (``signal N of record REC is
uncalibrated'', where N is the number of the signal in the record, counted from 0),
except in legacy mode.}
\begin{lstlisting}[language=Julia,style=example]
rec = read_record("100s"; to=2)
println(physical(rec))          # [-0.145 -0.065; -0.145 -0.065]
\end{lstlisting}

\apientry{signal\_names}{Return the signal descriptions of a record}
\begin{lstlisting}[language=Julia,style=signature]
signal_names(rec::Record)
\end{lstlisting}
Returns the description of each signal of a \code{Record}, such as \code{"MLII"}, as
bytes, without the white space that a header may put at either end of it (empty for a
signal without a description). These are the names that \code{channels} of
\code{read\_record} accepts. It uses no library state.
\begin{params}
\item[rec] The record, as \code{read\_record} returns it.
\end{params}
\returns{A \code{Vector\{Vector\{UInt8\}\}}, one element per signal of \code{rec}.}
\raises{Nothing.}
\begin{lstlisting}[language=Julia,style=example]
rec = read_record("100s"; to=1)
println(String.(signal_names(rec)))     # ["MLII", "V5"]
\end{lstlisting}

% ------------------------------------------------------------------------------------------
\subsection{Writing records}
\label{sec:julia-writing}

\apientry{write\_record}{Write a header and a signal file}
\begin{lstlisting}[language=Julia,style=signature]
write_record(w::Wfdb, record::ByteStr, samples::AbstractMatrix{<:Integer}, fs::Real;
             fmt::Integer=16, gains=nothing, baselines=nothing, adczeros=0,
             adcres=nothing, units=nothing, names=nothing, samps_per_frame=nothing,
             base_time=nothing, comments=(), metadata=nothing, dir=nothing,
             legacy::Bool=false)
\end{lstlisting}
Writes samples in ADC units as a new record: the header \file{RECORD.hea} and one signal
file, \file{RECORD.dat}, which holds all the signals (so they share one format). Existing
files of these names are overwritten. The sampling frequency, the gain and ADC resolution
of each signal, and the units of each calibrated signal must be given: there are no
defaults. Every check except those of the writing itself comes before any file is
created, so an invalid argument writes nothing. The files are written as \code{putvec}
and \code{newheader} write them, followed by the info strings (\code{putinfo}).
\begin{params}
\item[record] Record name, without the \file{.hea} extension: a plain name
  (\code{"rec"}) or a relative path (\code{"data/rec"}), created in the output directory
  if one is set, otherwise in the current directory, missing directories included; or a
  full path, absolute or in the home directory (\code{"\textasciitilde/Documents/rec"}),
  used as given. With a path, the files are created in that directory, and the header
  names the record and its signal file without the path (\code{"rec"}, \code{"rec.dat"}),
  so that the record opens under any path that leads to it (in legacy mode, as in the
  original, the header gives the name with its path). A URL cannot be written. It must
  not be empty.
\item[samples] The samples in ADC units, one row per sample (per frame, with
  \code{samps\_per\_frame}) and one column per signal: integers within the range of
  \code{Int32}. \code{WFDB\_INVALID\_SAMPLE} is written as an invalid sample.
\item[fs] The sampling frequency in Hz: finite and $> 0$ (the frame frequency of a
  multi-frequency record).
\item[fmt] The storage format of all the signals, such as 16, 212, 24 or 516 (FLAC).
  Default 16.
\item[gains] The gain of each signal, in ADC units per physical unit (required): finite;
  0 marks an uncalibrated signal, and negative gains are allowed.
\item[baselines] The ADC value of 0 physical units of each signal. \code{nothing}: the
  signal's \code{adczero}.
\item[adczeros] The ADC value of 0\,V of each signal. Default 0.
\item[adcres] The ADC resolution of each signal, in bits, at least 1 (required).
\item[units] The physical units of each signal, such as \code{"mV"} (required for each
  calibrated signal; \code{nothing} for an uncalibrated one).
\item[names] The description of each signal. \code{nothing}: none.
\item[samps\_per\_frame] The number of samples per frame of each signal, integers $\ge 1$,
  for a multi-frequency record: each row of \code{samples} then holds
  \code{samps\_per\_frame[s]} consecutive columns for signal \code{s}. \code{nothing}: 1
  for every signal.
\item[base\_time] The start time and date, as \code{"hh:mm:ss dd/mm/yyyy"} (or a part of
  it, as \code{setbasetime} reads it). \code{nothing}: none.
\item[comments] The info strings to write at the end of the header, without their
  \code{\#}. Default: none.
\item[metadata] \code{Key: value} entries, written after the comments as
  \code{\# Key: value} (\code{\# Key:} for an empty value): a vector of \code{key =>
  value} pairs (or \code{(key, value)} tuples), written in that order, or a dictionary.
  Keys and values follow the rules of \code{set\_metadata}; no two keys may differ only in
  letter case, none may be the key of an entry in \code{comments}, and a value must not be
  \code{nothing}. \code{nothing}: none.
\item[dir] The directory (which must exist) in which to create the files, for this call
  only (the output directory of the call). \code{record} must then be a plain name,
  without a directory, and the record opens from that directory under its plain name.
  \code{nothing}: the output directory.
\item[legacy] \code{true} to behave as the original WFDB package, which assumes a gain of
  200, an ADC resolution based on the format and no units when they are not given, for
  this call only.
\end{params}
\returns{\code{nothing}.}
\raises{A \code{WfdbError}, before any file is created, if \code{record} is empty
(``write\_record: the record name must not be empty''); if there are no signals
(``write\_record: at least one signal is required''); if \code{gains}, \code{adcres} or
\code{units} is missing (``write\_record: gains must be specified'', ``write\_record: ADC
resolution(s) must be specified'', ``write\_record: units must be specified''); if a
\code{samps\_per\_frame} value is not an integer $\ge 1$ (``write\_record:
samps\_per\_frame must hold integers >= 1''), or the number of columns of \code{samples}
is not their sum (``write\_record: signals has N columns, but samps\_per\_frame adds up to
M''); if \code{fs} is not finite and $> 0$; if a gain is not finite (``write\_record:
gains must be finite''); if a vector of per-signal values has neither one value nor one
per signal (``write\_record: OPTION must have one value per signal''); if the number of
columns does not match the signals; if a sample lies outside the range of \code{Int32};
if \code{base\_time} is not a valid time (``write\_record: incorrect base time 'T'''); if
an entry of \code{metadata} is not valid (``write\_record: metadata key "a:b" must not
contain ':''', \ldots), or its key is also that of an entry in \code{comments}; or if
\code{dir} does not exist (``directory D does not exist'') or \code{record} contains a
directory while \code{dir} is given. A \code{WfdbError} also if a file cannot be created
(``can't create signal file for record REC''), or cannot be written completely, for
example on a full disk (``write error on record REC'').}
\begin{lstlisting}[language=Julia,style=example]
x = Int32[0 100; 10 110; 20 120; 30 130]            # 4 samples of 2 signals
write_record("two", x, 250; gains=[100, 200], adcres=12, units="mV",
             names=["a", "b"], base_time="12:00:00 01/01/2020",
             comments=["recorded at rest"], metadata=["Age" => "63"])
print(read("two.hea", String))
# two 2 250 4 12:00:00 01/01/2020
# two.dat 16 100/mV 12 0 0 60 0 a
# two.dat 16 200/mV 12 0 100 460 0 b
# #recorded at rest
# # Age: 63
println(read_record("two").samples == x)            # true

# A multi-frequency record: signal 1 has 2 samples per frame.
y = Int32[1 2 3; 4 5 6]                             # 2 frames: a, then b twice
write_record("mf", y, 100; gains=200, adcres=12, units="mV", samps_per_frame=[1, 2])
println(size(read_record("mf"; highres=true).samples))   # (4, 2)
\end{lstlisting}

\apientry{write\_record\_physical}{Write a record from samples in physical units}
\begin{lstlisting}[language=Julia,style=signature]
write_record_physical(w::Wfdb, record::ByteStr, samples::AbstractMatrix{<:Real}, fs::Real;
                      fmt::Integer=16, gains=nothing, baselines=nothing,
                      adczeros=0, adcres=nothing, units=nothing, names=nothing,
                      samps_per_frame=nothing, base_time=nothing, comments=(),
                      metadata=nothing, dir=nothing, legacy::Bool=false)
\end{lstlisting}
Writes a record as \code{write\_record} does, from samples in physical units. Each value
\code{v} is stored as $\mathrm{round}(v \cdot \mathit{gain}) + \mathit{baseline}$, with
halfway cases rounded away from zero; \code{NaN} is stored as an invalid sample. The gain
of every signal must be nonzero.
\begin{params}
\item[record] Record name, as for \code{write\_record}.
\item[samples] The samples in physical units, one row per sample (or frame) and one column
  per signal: real numbers; \code{NaN} for an invalid sample.
\item[fs] The sampling frequency in Hz, as for \code{write\_record}.
\item[fmt] The storage format, as for \code{write\_record}. Default 16.
\item[gains] The gain of each signal (required), finite and nonzero.
\item[baselines] The baseline of each signal, as for \code{write\_record}.
\item[adczeros] The ADC zero of each signal, as for \code{write\_record}.
\item[adcres] The ADC resolution of each signal (required), as for \code{write\_record}.
\item[units] The physical units of each signal (required), as for \code{write\_record}.
\item[names] The signal descriptions, as for \code{write\_record}.
\item[samps\_per\_frame] The samples per frame of each signal, as for
  \code{write\_record}.
\item[base\_time] The start time and date, as for \code{write\_record}.
\item[comments] The info strings, as for \code{write\_record}.
\item[metadata] The metadata entries, as for \code{write\_record}.
\item[dir] The directory in which to write the files, as for \code{write\_record}.
\item[legacy] \code{true} to behave as the original WFDB package (a gain of 0 is then
  taken as 200), for this call only.
\end{params}
\returns{\code{nothing}.}
\raises{A \code{WfdbError} in the cases of \code{write\_record}; if a gain is 0
(``write\_record: signal N is uncalibrated'', N counted from 0); or if a value is
infinite, or its ADC value lies outside the range of \code{Int32} (``write\_record:
signal N has a sample that is infinite or beyond the range of 32-bit ADC units'', N
counting the columns from 0). Nothing is written then.}
\begin{lstlisting}[language=Julia,style=example]
x = [0.0 1.0; 0.5 -1.0; NaN 0.25]                  # in mV; NaN: an invalid sample
write_record_physical("phys", x, 100; gains=200, adcres=12, units="mV", names=["a", "b"])
rec = read_record("phys")
println(rec.samples)            # Int32[0 200; 100 -200; -32768 50]
println(physical(rec))          # [0.0 1.0; 0.5 -1.0; NaN 0.25]
\end{lstlisting}

\apientry{write\_annotations}{Write an annotation file}
\begin{lstlisting}[language=Julia,style=signature]
write_annotations(w::Wfdb, record::ByteStr, annotator::ByteStr, ann::Annotations;
                  fs::Real=0.0, sort::Union{Nothing,Bool}=nothing, dir=nothing,
                  legacy::Bool=false)
\end{lstlisting}
Writes the annotations of an \code{Annotations} object to the annotation file
\file{RECORD.ANNOTATOR}, in the output directory or in \code{dir}, as \code{putann}
writes them. An existing file is overwritten. The type of each annotation is its mnemonic
in \code{ann.symbols} if that is not empty, and its code in \code{ann.codes} otherwise;
both may be given (as \code{read\_annotations} returns them) if they agree. As in C, the
subtype, channel and annotator number are written in one byte each (other values wrap
around), and auxiliary text longer than 255 bytes is truncated.
\begin{params}
\item[record] Record name, without the \file{.hea} extension (a trailing \code{".hea"}
  is also accepted): the record of the annotations (its header is not needed). A plain
  name or a relative path (\code{"data/100s"}) is created in the output directory if one
  is set, otherwise in the current directory; a full path, absolute or in the home
  directory (\code{"\textasciitilde/Documents/100s"}), is used as given. The file is
  \file{RECORD.ANNOTATOR} at that place (\file{data/100s.qrs}). A URL cannot be
  written.
\item[annotator] Annotator name: the suffix of the file, such as \code{"qrs"}.
\item[ann] The annotations (\cref{tab:julia-annotations}): \code{samples} (the times),
  and \code{symbols} or \code{codes}; \code{subtypes}, \code{chans}, \code{nums} and
  \code{aux} are optional. Each field is empty or has one element per annotation.
  \code{ann.record}, \code{ann.annotator} and \code{ann.fs} are not used.
\item[fs] 0 (the default), or a finite frequency $> 0$, which is written in the file as
  the frequency in which the annotation times are counted. With 0, readers use the
  record's sampling frequency. It applies to this file only: the time resolution of
  output annotators (\code{getafreq()}) is left as it was.
\item[sort] If the annotations are not in time order, the file is sorted when it is
  written: \code{true} always sorts, and \code{false} keeps the given order (without the
  library's advice to run \code{sortann}), for this call only. \code{nothing}: the
  environment variable \code{WFDBANNSORT} decides (sorting is on unless it is \code{0}).
\item[dir] The directory (which must exist) in which to create the file, for this call
  only. \code{record} must then be a plain name, without a directory: the file is
  \file{DIR/RECORD.ANNOTATOR}.
\item[legacy] \code{true} to behave as the original WFDB package, for this call only.
\end{params}
\returns{\code{nothing}.}
\raises{A \code{WfdbError}, before the file is created, if a field of \code{ann} is
neither empty nor has one element per annotation (``write\_annotations: FIELD must have
one entry per annotation (N), or none''); if \code{ann.symbols} and \code{ann.codes} are
both given and disagree (``write\_annotations: symbols and codes disagree at annotation
K'', K counted from 0); if \code{fs} is neither 0 nor a finite frequency $> 0$; or if
\code{dir} does not exist or \code{record} contains a directory while \code{dir} is given.
A \code{WfdbError} also if the file cannot be created (``can't create annotator ANN for
record REC'') or cannot be written completely (``write error on annotator ANN for record
REC'').}
\begin{lstlisting}[language=Julia,style=example]
ann = Annotations(samples=[100, 460, 820], symbols=[b"N", b"N", b"V"])
write_annotations("100s", "qrs", ann; fs=360)       # writes 100s.qrs
a = read_annotations("100s", "qrs")
println(a.samples, " ", String.(a.symbols))         # [100, 460, 820] ["N", "N", "V"]
println(a.codes == [WFDBX.NORMAL, WFDBX.NORMAL, WFDBX.PVC])   # true

# Out of order: sorted when the file is written (the library prints a message
# "Rearranging annotations for output annotator srt ...done!").
write_annotations("100s", "srt", Annotations(samples=[460, 100], codes=Int32[1, 1]); sort=true)
println(read_annotations("100s", "srt").samples)    # [100, 460]
\end{lstlisting}

% ------------------------------------------------------------------------------------------
\subsection{Checking records}
\label{sec:julia-checking}

\apientry{check\_record}{Check a record and its annotation files}
\begin{lstlisting}[language=Julia,style=signature]
check_record(w::Wfdb, record::ByteStr; annotators=String[],
             input_dir::OptByteStr=nothing, network::Union{Nothing,Bool}=nothing)
\end{lstlisting}
Checks the header of a record, its signal specifications, every sample of its signal
files, each segment of a multi-segment record, and the annotation files named, and
reports each problem found. A problem in the record is never an error of the function: it
is reported, with a severity (\code{"error"}: the record is damaged or does not follow
the WFDB format, and programs will read wrong data or fail; \code{"warning"}: the record
can be read, but something is missing or unusual, and may give wrong results;
\code{"info"}: a fact worth knowing, such as invalid samples). The problems, their codes
and their messages are the same in the seven languages (\cref{tab:python-issue-codes}).
The checks run in normal mode, even if legacy mode is on, and the library prints no
messages during the call (the previous settings are restored afterwards). They are
similar to those of Benjamin Moody's \code{wfdbcheck}.
\begin{params}
\item[record] Record name, without the \file{.hea} extension (a trailing \code{".hea"} is
  also accepted): a full path, absolute (\code{"/data/mitdb/100"}) or in the home
  directory (\code{"\textasciitilde/Documents/100s"}), used as given; a relative path
  (\code{"100s"}, \code{"data/100s"}), read from the input directory if one is set,
  otherwise from the current directory; or a URL
  (\code{"https://physionet.org/files/mitdb/1.0.0/100"}), read only if network access is
  on.
\item[annotators] The annotator names whose files are checked, such as \code{["atr"]}
  (one name may be given as a string). Default: none.
\item[input\_dir] The input directory for this call only (optional), as for
  \code{read\_record}: relative record names are read under it. It does not apply to full
  paths or URLs.
\item[network] Network access for this call only, as for \code{read\_record}.
\end{params}
\returns{A \code{CheckReport} (\cref{tab:julia-checkreport}): \code{report.ok} is
\code{true} if no problem is an error, \code{report.issues} lists the problems
(\code{Issue} objects, \cref{tab:julia-issue}), and \code{lines(report)} (or
\code{print(report)}) gives one line per problem, or \code{no problems found}.}
\raises{An \code{ArgumentError} if \code{record} is empty (``record must not be empty'')
or an annotator name is not a string (``annotator names must be strings''). A record that
cannot be read is reported (code \code{header\_unreadable}), not thrown as an error.}
\begin{lstlisting}[language=Julia,style=example]
report = check_record("100s"; annotators=["atr"])
println(report.ok)                                  # true: no errors
println(report)
# warning: signal 0: the units are not specified
# warning: signal 1: the units are not specified

# A damaged record: one sample of a 12-bit signal is overwritten with 2048.
write_record("damaged", reshape(Int32.(0:99), :, 1), 100; gains=100, adcres=12,
             units="mV", names=["ramp"])
dat = read("damaged.dat")
dat[21] = 0x00; dat[22] = 0x08      # sample 10 (2 bytes per sample in format 16): 2048
write("damaged.dat", dat)
report = check_record("damaged")
println(report.ok)                                  # false
for i in report.issues
    println(i.severity, " ", i.code, " ", i.where)
end
# error sample_range signal 0
# error checksum signal 0
println(lines(report)[1])
# error: signal 0: 1 sample is outside the ADC range (-2048 to 2047): sample 10 (value 2048)
\end{lstlisting}

\apientry{lines}{The text form of a check report}
\begin{lstlisting}[language=Julia,style=signature]
lines(r::CheckReport)
\end{lstlisting}
Returns the lines of a report of \code{check\_record}: one line per issue, in the order of
the issues (\code{"SEVERITY: WHERE: MESSAGE"}, without \code{"WHERE: "} for the record
itself), or the one line \code{"no problems found"}. \code{print(report)} and
\code{show(report)} write these lines, joined by line feeds. The lines are the same in the
seven languages.
\begin{params}
\item[r] The report, as \code{check\_record} returns it.
\end{params}
\returns{A \code{Vector\{String\}}, with at least one element.}
\raises{Nothing.}
\begin{lstlisting}[language=Julia,style=example]
for l in lines(check_record("100s"))
    println(l)
end
# warning: signal 0: the units are not specified
# warning: signal 1: the units are not specified
x = zeros(Int32, 10, 1)
write_record("clean", x, 100; gains=200, adcres=12, units="mV", names=["a"])
println(lines(check_record("clean")))               # ["no problems found"]
\end{lstlisting}

% ------------------------------------------------------------------------------------------
\subsection{Header metadata}
\label{sec:julia-metadata}

These functions read and change the metadata of a header: its info strings of the form
\code{\# Key: value}, such as \code{\# Age: 63} (\cref{chap:metadata}). In short: keys are
compared without regard to letter case (A-Z and a-z only); a key must not be empty and
must not contain \code{:}, \code{\#}, CR, LF or NUL, and a value must not contain CR, LF
or NUL; spaces, tabs, VT, FF and CR at either end of keys and values are removed; keys and
values are bytes (arguments \code{ByteStr}: a \code{String} is taken as its bytes;
results \code{Vector\{UInt8\}}). A change rewrites the header in one step, through a
temporary file in the directory of the destination, which then replaces it, keeping its
permissions and every other byte of it; the temporary file is removed on any failure. All
the arguments are checked first, so an error leaves the header unchanged, and nothing is
written if nothing changes (unless \code{dir} is given). Only the \file{.hea} file is read
(not a \file{.info} file); for a multi-segment record, it is the master header, unless a
segment is named. The header is opened as the library opens a header
(\code{wfdb\_open}), so in legacy mode it is searched for in the WFDB path. The entries
are returned as \code{MetadataEntry} pairs:
\begin{lstlisting}[language=Julia]
const MetadataEntry = Pair{Vector{UInt8},Vector{UInt8}}     # key => value
\end{lstlisting}
The functions take the keyword arguments \code{input\_dir} and \code{network} for reading
the header, and the functions that change it also \code{dir}, the directory in which to
write it. \code{dir} \code{nothing} or \code{""} changes the header that was read. These
functions do not call \code{wfdbquit()}. Invalid arguments throw an
\code{ArgumentError}, and a header that cannot be read or written a \code{WfdbError}
(\cref{sec:julia-warnings}). Metadata often holds information about patients: WFDB-x
stores it as plain text, exactly as given, and does not de-identify it.

\begin{phibox}
A header can hold identifying data about a patient (a name, a date of birth, a medical
record number). The metadata functions read and write such data as given. Records that
are shared must be de-identified first.
\end{phibox}

\apientry{read\_metadata}{Read the metadata of a record}
\begin{lstlisting}[language=Julia,style=signature]
read_metadata(w::Wfdb, record::ByteStr; input_dir::OptByteStr=nothing,
              network::Union{Nothing,Bool}=nothing)
\end{lstlisting}
Reads the metadata entries of the header of a record, in file order. An info string is an
entry if it contains a colon: the key is the text before the first colon, the value the
rest of the line, both without spaces at either end. Info strings without a colon, or
whose key is empty or contains \code{\#}, are free comments, which are not returned.
\begin{params}
\item[record] Record name, without the \file{.hea} extension (a trailing \code{".hea"} is
  also accepted): a full path, absolute (\code{"/data/mitdb/100"}) or in the home
  directory (\code{"\textasciitilde/Documents/100s"}), used as given; a relative path
  (\code{"100s"}, \code{"data/100s"}), read from the input directory if one is set,
  otherwise from the current directory; or a URL
  (\code{"https://physionet.org/files/mitdb/1.0.0/100"}), read only if network access is
  on.
\item[input\_dir] The input directory for this call only (optional), as for
  \code{read\_record}: relative record names are read under it. It does not apply to full
  paths or URLs.
\item[network] Network access for this call only, as for \code{read\_record}.
\end{params}
\returns{The entries, as \code{key => value} pairs of bytes, in file order. A key that
appears more than once (ignoring letter case) is returned once, with its first value, and
a warning is logged (\cref{sec:julia-warnings}).}
\raises{A \code{WfdbError} if the header cannot be found or read (``can't read header for
record REC''); an \code{ArgumentError} if \code{record} is empty (``record name must not be
empty'') or \code{"-"} (``the metadata of a header read from the standard input cannot be
read or changed'').}
\begin{lstlisting}[language=Julia,style=example]
write_record("pat", zeros(Int32, 360, 1), 360; gains=200, adcres=12, units="mV",
             names=["ECG"], comments=["recorded at rest"],
             metadata=["Age" => "63", "Sex" => "F"])
for (key, value) in read_metadata("pat")
    println(String(key), " = ", String(value))
end
# Age = 63
# Sex = F
for c in read_header("pat").comments
    println("[", String(c), "]")
end
# [recorded at rest]: the free comment is not an entry
# [ Age: 63]
# [ Sex: F]
mkpath("in"); cp("100s.hea", "in/100s.hea")
md = read_metadata("100s"; input_dir="in")          # reads in/100s.hea
println(length(md), " ", String(md[1].second))      # 1 0
\end{lstlisting}

\apientry{get\_metadata}{Return the value of a metadata entry}
\begin{lstlisting}[language=Julia,style=signature]
get_metadata(w::Wfdb, record::ByteStr, key::ByteStr; input_dir::OptByteStr=nothing,
             network::Union{Nothing,Bool}=nothing)
\end{lstlisting}
Returns the value of one metadata entry of a record. Letter case is ignored, and the
spaces at either end of \code{key} are removed: \code{"age"} finds \code{Age}. If the
header has more than one entry for \code{key}, the first is returned; as in
\code{read\_metadata}, a warning is logged for each key of the header that has more than
one entry.
\begin{params}
\item[record] Record name, as for \code{read\_metadata}.
\item[key] The key, such as \code{"Age"} (see \code{set\_metadata} for the valid keys).
\item[input\_dir] The input directory for this call only (optional), as for
  \code{read\_record}: relative record names are read under it. It does not apply to full
  paths or URLs.
\item[network] Network access for this call only, as for \code{read\_record}.
\end{params}
\returns{The value, as bytes (an empty vector for an entry with an empty value), or
\code{nothing} if the header has no entry for \code{key}.}
\raises{An \code{ArgumentError} if \code{key} is not a valid key (checked before the header
is read), or as \code{read\_metadata}.}
\begin{lstlisting}[language=Julia,style=example]
write_record("pat", zeros(Int32, 10, 1), 100; gains=200, adcres=12, units="mV",
             metadata=["Age" => "63", "Notes" => ""])
println(String(get_metadata("pat", "age")))         # 63
println(get_metadata("pat", "Notes"))               # UInt8[]
println(get_metadata("pat", "Weight") === nothing)  # true
\end{lstlisting}

\apientry{set\_metadata}{Set the value of a metadata entry}
\begin{lstlisting}[language=Julia,style=signature]
set_metadata(w::Wfdb, record::ByteStr, key::ByteStr, value::ByteStr;
             input_dir::OptByteStr=nothing, network::Union{Nothing,Bool}=nothing,
             dir::OptByteStr=nothing)
\end{lstlisting}
Sets the value of a metadata entry in the header of a record. If the header has an entry
for \code{key} (ignoring letter case), its line is replaced by \code{\# Key: value}
(\code{\# Key:} if the value is empty), at the same position and with the same line
ending, with the spelling of \code{key} given here. Otherwise the entry is added after the
last info string of the header (or, if it has none, after its last line that is not
blank), ending as the record line does. Every other byte of the header is kept.
\begin{params}
\item[record] Record name, as for \code{read\_metadata}. The header is changed where it
  was read, unless \code{dir} is given; a header read from a URL can only be written
  elsewhere, with \code{dir}.
\item[key] The key, such as \code{"Age"} or \code{"Date of birth"}. Spaces at either end
  are removed. It must not be empty, and must not contain \code{:}, \code{\#}, a line
  break or a NUL byte.
\item[value] The value, which may be a sentence. Spaces at either end are removed. It must
  not contain a line break or a NUL byte; it may be empty.
\item[input\_dir] The input directory for reading the header, for this call only.
\item[network] Network access for reading the header, for this call only.
\item[dir] An existing directory in which to write the changed header, under the same file
  name, instead of changing the header itself (the signal files are not copied).
  \code{nothing} or \code{""}: change the header that was read.
\end{params}
\returns{\code{nothing}.}
\raises{An \code{ArgumentError} if \code{record} is empty or \code{"-"}, if \code{key} or
\code{value} is not valid (``metadata key "a:b" must not contain ':''', ``value of metadata
key "Age" must not contain a line break'', \ldots), if the header has more than one entry
for \code{key} (``header of record REC has more than one entry for key KEY (lines N, M);
remove them first''), if it has no record line, or if it was read from a URL and
\code{dir} is not given. A \code{WfdbError} if the header cannot be read (``can't read
header for record REC''), if the changed header cannot be written (``can't write header
for record REC''), or if \code{dir} does not exist (``directory D does not exist'').}
\begin{lstlisting}[language=Julia,style=example]
write_record("pat", zeros(Int32, 10, 1), 100; gains=200, adcres=12, units="mV",
             metadata=["Age" => "63"])
set_metadata("pat", "age", "64")                    # replaces the line "# Age: 63"
set_metadata("pat", "Date of birth", "1960-01-01")  # adds a line
for l in eachline("pat.hea")
    startswith(l, "#") && println(l)
end
# # age: 64
# # Date of birth: 1960-01-01
mkpath("copy")
set_metadata("pat", "Age", "65"; dir="copy")        # writes copy/pat.hea
println(String(get_metadata("copy/pat", "Age")), " ", String(get_metadata("pat", "Age")))
# 65 64
\end{lstlisting}

\apientry{remove\_metadata}{Remove a metadata entry}
\begin{lstlisting}[language=Julia,style=signature]
remove_metadata(w::Wfdb, record::ByteStr, key::ByteStr; input_dir::OptByteStr=nothing,
                network::Union{Nothing,Bool}=nothing, dir::OptByteStr=nothing)
\end{lstlisting}
Removes a metadata entry from the header of a record. Every entry for \code{key} (ignoring
letter case) is removed, with its line ending, so this also repairs a header that has more
than one. Every other byte of the header is kept. Nothing is written if the header has no
entry for \code{key} (unless \code{dir} is given).
\begin{params}
\item[record] Record name, as for \code{set\_metadata}.
\item[key] The key, as for \code{set\_metadata}.
\item[input\_dir] The input directory for reading the header, for this call only.
\item[network] Network access for reading the header, for this call only.
\item[dir] The directory in which to write the changed header, as for
  \code{set\_metadata}.
\end{params}
\returns{The number of lines removed (0 if the header has no entry for \code{key}).}
\raises{As \code{set\_metadata} (a key that the header has more than once is not an error
here).}
\begin{lstlisting}[language=Julia,style=example]
write_record("pat", zeros(Int32, 10, 1), 100; gains=200, adcres=12, units="mV",
             metadata=["Age" => "63", "Sex" => "F"])
println(remove_metadata("pat", "SEX"))              # 1
println(remove_metadata("pat", "Weight"))           # 0: nothing is written
println(length(read_metadata("pat")))               # 1
\end{lstlisting}

\apientry{update\_metadata}{Make several changes to the metadata}
\begin{lstlisting}[language=Julia,style=signature]
update_metadata(w::Wfdb, record::ByteStr, changes; input_dir::OptByteStr=nothing,
                network::Union{Nothing,Bool}=nothing, dir::OptByteStr=nothing)
\end{lstlisting}
Makes several changes to the metadata of a record in one write. Each change is applied in
order, as \code{set\_metadata} does, or, if its value is \code{nothing}, as
\code{remove\_metadata} does. All the changes are checked first: if one is not valid, none
is made.
\begin{params}
\item[record] Record name, as for \code{set\_metadata}.
\item[changes] A vector (or tuple) of \code{key => value} pairs (or \code{(key, value)}
  tuples), or a dictionary; a value of \code{nothing} removes the entry. No two keys may
  differ only in letter case. The changes of a \code{Dict} are applied in its iteration
  order, which is not defined: a vector fixes the order in which new entries are added.
\item[input\_dir] The input directory for reading the header, for this call only.
\item[network] Network access for reading the header, for this call only.
\item[dir] The directory in which to write the changed header, as for
  \code{set\_metadata}.
\end{params}
\returns{\code{nothing}.}
\raises{As \code{set\_metadata}; an \code{ArgumentError} also if \code{changes} is not a
vector, a tuple or a dictionary of pairs (``metadata must be a dictionary, or a vector of
key => value pairs''), if a value is neither a string nor \code{nothing}, or if a key is
given more than once (``metadata key "AGE" is given more than once (keys are compared
without regard to letter case)'').}
\begin{lstlisting}[language=Julia,style=example]
write_record("pat", zeros(Int32, 10, 1), 100; gains=200, adcres=12, units="mV",
             metadata=["Age" => "63", "Dx" => "164889003"])
update_metadata("pat", ["Age" => "64", "Dx" => nothing, "Sex" => "M"])
println([String(k) => String(v) for (k, v) in read_metadata("pat")])
# ["Age" => "64", "Sex" => "M"]
\end{lstlisting}

% ------------------------------------------------------------------------------------------
\subsection{The result types}
\label{sec:julia-results}

The readers return objects of the types \code{Record}, \code{Header} and
\code{Annotations}, whose fields are listed in
\cref{tab:julia-record,tab:julia-header,tab:julia-annotations}; \code{check\_record}
returns a \code{CheckReport}, which holds \code{Issue} objects
(\cref{tab:julia-checkreport,tab:julia-issue}). The signals of a record are described by
\code{Siginfo} structures (\cref{tab:julia-siginfo}), which are those of the WFDB library
(\cref{sec:julia-library}). \code{Record}, \code{Header}, \code{CheckReport} and
\code{Issue} are immutable structures (the vectors and matrices they hold can be changed
in place); \code{Annotations} and \code{Siginfo} are mutable, with keyword constructors
whose defaults are empty vectors or C's zero values, such as \code{Annotations(samples=[100],
symbols=[b"N"])}. \code{length(ann)} is the number of annotations.

\begin{xltabular}{\linewidth}{@{}lL@{}}
\caption{The fields of \code{Record}, the result of \code{read\_record}}
\label{tab:julia-record} \\
\toprule
\tabhead Field & \tabhead Contents \\
\midrule
\endfirsthead
\toprule
\tabhead Field & \tabhead Contents \\
\midrule
\endhead
\bottomrule
\endfoot
\code{name::String} & the record name, as given to \code{read\_record} \\
\code{fs::Float64} & the sampling frequency of the returned samples, in Hz, or 0 if the
  header gives none; with \code{highres}, the highest sampling frequency of a
  multi-frequency record \\
\code{samples::Matrix\{Int32\}} & the samples in ADC units, one row per sample and one
  column per signal returned; invalid samples are \code{WFDB\_INVALID\_SAMPLE} \\
\code{siginfo::Vector\{Siginfo\}} & one per signal returned, in the order of the columns;
  \code{desc} has no white space at either end \\
\code{base\_time::Union\{Nothing,String\}} & the start time and date as text, to the
  second, such as \code{"[14:00:03 21/06/2024]"}; \code{nothing} if the header gives
  none \\
\code{base\_time\_ms::Union\{Nothing,Int64\}} & the start time in milliseconds after
  midnight (\code{"14:00:03.500"} gives \code{50403500}); \code{nothing} exactly when
  \code{base\_time} is \\
\code{base\_date::Union\{Nothing,NTuple\{3,Int\}\}} & the start date as \code{(year,
  month, day)}; \code{nothing} if the header gives no date \\
\code{sampfrom::Int64} & the number of the first sample returned \\
\code{comments::Vector\{Vector\{UInt8\}\}} & the info strings of the record: the comments
  of the header, and the lines of its \file{.info} file if it has one (see
  \code{getinfo}) \\
\code{channels::Vector\{Int\}} & the number, in the record, of each signal returned
  (\code{0} is the first signal) \\
\code{checksum\_ok::Bool} & \code{false} if the samples do not match the checksums of the
  header; checksums are checked only when a read reaches the end of the record, so a
  shorter read gives \code{true} \\
\end{xltabular}

\begin{xltabular}{\linewidth}{@{}lL@{}}
\caption{The fields of \code{Header}, the result of \code{read\_header}}
\label{tab:julia-header} \\
\toprule
\tabhead Field & \tabhead Contents \\
\midrule
\endfirsthead
\toprule
\tabhead Field & \tabhead Contents \\
\midrule
\endhead
\bottomrule
\endfoot
\code{name::String} & the record name, as given \\
\code{fs::Float64} & the sampling frequency in Hz, as \code{sampfreq} gives it; 0 if the
  header gives none \\
\code{counter\_freq::Float64} & the counter frequency (\code{getcfreq}) \\
\code{base\_counter::Float64} & the base counter value (\code{getbasecount}) \\
\code{nsamples::Int64} & the number of samples per signal, from the record line of the
  header (0 if the header does not give it) \\
\code{nsig::Int} & the number of signals \\
\code{siginfo::Vector\{Siginfo\}} & the signal specifications, one per signal; \code{desc}
  has no white space at either end \\
\code{base\_time}, \code{base\_time\_ms}, \code{base\_date} & the start time and date, as
  in \code{Record} \\
\code{comments::Vector\{Vector\{UInt8\}\}} & the info strings of the record, as in
  \code{Record} \\
\code{metadata::Vector\{MetadataEntry\}} & the metadata entries of the header, in file
  order, as \code{read\_metadata} returns them: from the \file{.hea} file only; empty,
  without an error, for an EDF record and for a header read from the standard input \\
\end{xltabular}

\begin{xltabular}{\linewidth}{@{}lL@{}}
\caption{The fields of \code{Annotations}, the result of \code{read\_annotations} and the
  argument of \code{write\_annotations}: parallel vectors, annotation \code{k} being
  \code{samples[k]}, \code{codes[k]}, \code{symbols[k]}, \ldots, in file order}
\label{tab:julia-annotations} \\
\toprule
\tabhead Field & \tabhead Contents \\
\midrule
\endfirsthead
\toprule
\tabhead Field & \tabhead Contents \\
\midrule
\endhead
\bottomrule
\endfoot
\code{record::String}, \code{annotator::String} & the record and annotator names that
  were read, such as \code{"100s"} and \code{"atr"}; not used by
  \code{write\_annotations} \\
\code{fs::Float64} & the frequency in which \code{samples} are counted, normally the
  sampling frequency of the record; 0 if unknown. \code{write\_annotations} does not use
  it \\
\code{samples::Vector\{Int64\}} & the annotation times, in samples \\
\code{codes::Vector\{Int32\}} & the annotation codes, such as \code{WFDBX.NORMAL} or
  \code{WFDBX.PVC} \\
\code{symbols::Vector\{Vector\{UInt8\}\}} & the mnemonics of the codes, as \code{annstr}
  returns them, such as \code{b"N"}, \code{b"V"} or \code{b"+"}; when writing, an empty
  mnemonic means ``use \code{codes}'' \\
\code{subtypes}, \code{chans}, \code{nums} (\code{Vector\{Int32\}}) & the subtype,
  channel number and annotator number of each annotation \\
\code{aux::Vector\{Union\{Nothing,Vector\{UInt8\}\}\}} & the auxiliary text of each
  annotation, without its length byte, or \code{nothing} where there is none \\
\end{xltabular}

\begin{table}[htbp]
\centering
\caption{The fields of \code{CheckReport}, the result of \code{check\_record}}
\label{tab:julia-checkreport}
\begin{tabularx}{\linewidth}{@{}lL@{}}
\toprule
\tabhead Member & \tabhead Contents \\
\midrule
\code{record::String} & the record name, as given \\
\code{ok::Bool} & \code{true} if no issue has the severity \code{"error"} \\
\code{issues::Vector\{Issue\}} & the problems, in the order of the checks \\
\code{lines(report)} & the text form of the report: one line per issue, or the one line
  \code{no problems found} \\
\code{print}, \code{show} & write the lines of \code{lines(report)}, joined by line feeds
  (no final line feed) \\
\bottomrule
\end{tabularx}
\end{table}

\begin{table}[htbp]
\centering
\caption{The fields of \code{Issue}, a problem found by \code{check\_record}}
\label{tab:julia-issue}
\begin{tabularx}{\linewidth}{@{}lL@{}}
\toprule
\tabhead Member & \tabhead Contents \\
\midrule
\code{severity::String} & \code{"error"}, \code{"warning"} or \code{"info"} \\
\code{code::String} & a stable identifier of the kind of problem, such as
  \code{"checksum"} or \code{"gain\_missing"} (\cref{tab:python-issue-codes}) \\
\code{where::String} & where the problem is: \code{""} (the record), \code{"signal N"},
  \code{"segment SEG"}, \code{"segment SEG, signal N"} or \code{"annotator ANN"} \\
\code{message::String} & the description of the problem \\
\code{string(issue)} & \code{SEVERITY: WHERE: MESSAGE}, or \code{SEVERITY: MESSAGE} when
  \code{where} is empty \\
\bottomrule
\end{tabularx}
\end{table}

\begin{xltabular}{\linewidth}{@{}lL@{}}
\caption{The fields of \code{Siginfo} (C's \code{WFDB\_Siginfo}), the specification of a
  signal}
\label{tab:julia-siginfo} \\
\toprule
\tabhead Field & \tabhead Contents \\
\midrule
\endfirsthead
\toprule
\tabhead Field & \tabhead Contents \\
\midrule
\endhead
\bottomrule
\endfoot
\code{fname::Union\{Nothing,Vector\{UInt8\}\}} & the name of the signal file, as the
  header gives it \\
\code{desc::Union\{Nothing,Vector\{UInt8\}\}} & the signal description, such as
  \code{b"MLII"} \\
\code{units::Union\{Nothing,Vector\{UInt8\}\}} & the physical units, such as \code{b"mV"};
  \code{nothing} if not specified \\
\code{gain::Float64} & the gain, in ADC units per physical unit (0: uncalibrated) \\
\code{initval::Int32} & the value of the first sample \\
\code{group::UInt32} & the signal group number (the signals of one group share one
  file) \\
\code{fmt::Int32} & the storage format, such as \code{16}, \code{212} or \code{516} \\
\code{spf::Int32} & the number of samples per frame (more than 1 for the faster signals of
  a multi-frequency record) \\
\code{bsize::Int32} & the block size, in bytes (0 for ordinary files) \\
\code{adcres::Int32} & the ADC resolution, in bits (0: not specified) \\
\code{adczero::Int32} & the ADC output for an input of 0\,V \\
\code{baseline::Int32} & the ADC value that corresponds to 0 physical units \\
\code{nsamp::Int64} & the number of samples (0: not specified) \\
\code{cksum::Int32} & the 16-bit checksum of the samples \\
\end{xltabular}

% ------------------------------------------------------------------------------------------
\subsection{Warnings and errors}
\label{sec:julia-warnings}

\paragraph{Exceptions.} The high-level functions report errors by exceptions
(\cref{tab:julia-exceptions}). \code{WfdbError} is the exception type of WFDB-x; its field
\code{msg} (a \code{String}) describes the failed operation, such as ``can't read record
REC'', and \code{showerror} prints \code{WfdbError: MESSAGE}. The metadata functions and
\code{check\_record} throw Julia's \code{ArgumentError} (whose field is \code{msg} too)
for an invalid argument, and a \code{MethodError} is thrown, as for any Julia function,
for an argument of a type that no method accepts (such as a number as the value of
\code{set\_metadata}). Each function checks its arguments before it creates or changes a
file, so an invalid argument leaves the files as they were. The library may also print a
message of its own on the standard error when it cannot open a file (as in \code{init:
can't open header for record REC}); \code{wfdberror()} returns the last such message, and
\code{wfdbquiet()} turns these messages off (\cref{sec:julia-library}).

\begin{table}[htbp]
\centering
\caption{The exceptions of the high-level API}
\label{tab:julia-exceptions}
\begin{tabularx}{\linewidth}{@{}lL@{}}
\toprule
\tabhead Exception & \tabhead Cases \\
\midrule
\code{WfdbError} & a record, header or annotation file cannot be found, read, created or
  written completely (for example on a full disk): ``can't read record REC'', ``can't read
  header for record REC'', ``can't write header for record REC'', ``write error on record
  REC'', ``directory D does not exist'', \ldots; a signal or annotation file ends before
  its header or its end-of-file marker says (``\ldots(getvec status S)'', ``\ldots(getann
  status S)''); an invalid argument of the readers and writers (most messages of
  \code{write\_record} and \code{write\_annotations} begin with the function's name,
  including those of an invalid \code{metadata} argument of \code{write\_record}); the
  operation needs a value that the record does not give: a time in seconds without a
  sampling frequency (``sampling frequency of record REC is not specified''), physical
  units without a gain (``signal N of record REC is uncalibrated'') \\
\code{ArgumentError} & an invalid argument of the metadata functions: an invalid key or
  value, a record name that is empty or \code{"-"}, a key given more than once in
  \code{changes}, a key that the header has more than once (to be set), a header without a
  record line, or a header read from a URL to be changed without \code{dir}; an empty
  record name or an annotator name that is not a string, in \code{check\_record} \\
\code{MethodError} & an argument of a type that no method accepts \\
\bottomrule
\end{tabularx}
\end{table}

\paragraph{Warnings} are logged with Julia's logging macro \code{@warn}, at the level
\code{Warn}, so the current logger decides where they go (the standard error by default);
\code{Test.@test\_logs} checks them, and a logger such as
\code{Base.CoreLogging.NullLogger()} silences them.

\paragraph{The checksum warning.} When \code{read\_record} reads to the end of a record
whose samples do not match the checksums of its header, it returns all the samples and
logs the warning ``checksum error in record REC: the signal data do not match the
checksums in its header''. \code{checksum\_warning=false} turns the warning off;
\code{checksum\_ok} of the result is \code{false} either way.

\paragraph{The duplicate-key warning.} When a header holds two metadata entries whose keys
differ only in letter case (written by other software), \code{read\_metadata},
\code{get\_metadata} and \code{read\_header} use the first and log the warning ``header of
record REC has more than one entry for key KEY (lines L1, L2); the first is used'' (KEY
spelled as in its first entry). \code{set\_metadata} and \code{update\_metadata} then
refuse to change that key (an \code{ArgumentError}), and \code{remove\_metadata} removes
all of its entries.

\paragraph{The library's settings after an error.} Whether a high-level function returns or
throws, the getvec mode (\code{getgvmode()}) and the time resolution of output annotators
(\code{getafreq()}) are those of before the call, as are legacy mode and the input and
output directories (\cref{sec:julia-conventions}). Opening an EDF record changes both, as
in C, and \code{wfdbquit()} resets neither: after \code{read\_header("100s.edf")}, a
\code{write\_annotations} without \code{fs} writes no time resolution, and the next
\code{read\_record} reads in the mode it reads in by default. The C-style functions
(\code{isigopen}, \code{wfdbquit}, \ldots) keep the behavior of C.

\begin{lstlisting}[language=Julia,style=example]
# A record whose last sample is changed after it was written: the checksum fails.
write_record("ck", reshape(Int32.(0:99), :, 1), 100; gains=100, adcres=12, units="mV")
dat = read("ck.dat"); dat[end-1:end] .= 0x00; write("ck.dat", dat)
rec = read_record("ck")         # logs the checksum warning
println(size(rec.samples, 1), " ", rec.checksum_ok)  # 100 false: all the samples
rec = read_record("ck"; checksum_warning=false)     # no warning

# Two entries for one key: the first is used, with a warning.
write_record("dup", zeros(Int32, 10, 1), 100; gains=0, adcres=12,
             comments=[" Age: 63", " AGE: 64"])
println(String(get_metadata("dup", "age")))         # 63 (with a warning)
try
    set_metadata("dup", "Age", "65")
catch e
    println(e.msg)
    # header of record dup has more than one entry for key Age (lines 3, 4); remove them first
end
println(remove_metadata("dup", "age"))              # 2

try
    read_record("nosuchrecord")     # (the library also prints a message)
catch e
    println(e isa WfdbError, " ", e.msg)            # true can't read record nosuchrecord
end
\end{lstlisting}

% ==========================================================================================
\section{The WFDB library functions in Julia}
\label{sec:julia-library}

Every function of the WFDB library is a function of the module \code{WFDBX}, under its C
name, with the same arguments, in the same order, and the same return values, so that
programs and documentation written for the C library carry over directly. The functions,
their arguments and their results are documented once, for all the languages, in
\cref{chap:library}; this section gives the rules by which the C declarations become
Julia functions (see also \cref{sec:library-julia}). The docstring of each function
(\code{?getvec}) summarizes it, and says what WFDB-x changes in it.

\subsection{Conventions}

\paragraph{The library state.} The state of the C library (the open record, its signals,
the annotators, the WFDB path, the time base, \ldots) is held by a \code{Wfdb} object, and
each C function that uses it takes one as an optional first argument:
\code{isigopen(w, "100s", si, 2)}, \code{getvec(w, v)}. Without it, the function uses the
default object, \code{lib()}. \code{Wfdb()} creates an independent one; separate objects
share no state, unlike the C library. The functions that use no state take none:
\code{wfdbldflags}, \code{wfdbcflags}, \code{wfdbdefwfdb}, \code{wfdbdefwfdbcal},
\code{isann}, \code{ammap}, \code{mamap} (and their \code{wfdb\_} forms), and
\code{wfdbx\_write\_errors} and \code{wfdbx\_note\_write\_error}, whose count belongs to
the process.

\paragraph{Names.} The functions, the constants (\code{WFDB\_READ},
\code{WFDB\_INVALID\_SAMPLE}, \code{WFDB\_HIGHRES}, \ldots) and the structures keep their C
names. \code{using WFDBX} makes most of them available; the others are called with the
module prefix: \code{WFDBX.wfdb\_error}, \code{WFDBX.wfdb\_me\_fatal}, the annotation codes
(\code{WFDBX.NORMAL}, \code{WFDBX.PVC}, \ldots, \cref{chap:formats}), and the function
forms of the macros of \file{<wfdb/ecgmap.h>} (\code{WFDBX.wfdb\_isann},
\code{WFDBX.wfdb\_isqrs}, \ldots; \code{isann}, \code{isqrs}, \code{map1}, \ldots{} are
exported).

\paragraph{Arguments.} String arguments are \code{String}s or bytes (\code{ByteStr}); a
byte string ends at its first NUL byte, as a C string does. C's \code{NULL} is
\code{nothing}: \code{sampfreq(nothing)} returns the frequency of the open record, and
\code{getinfo(nothing)} continues with the info strings of the last record. Where C fills
an array, the caller passes a vector: of \code{Siginfo} structures (\code{[Siginfo() for
\_ in 1:n]}), or of samples, which must be a \code{Vector\{Int32\}}; \code{nothing} stands
for C's \code{NULL} (\code{isigopen("100s", nothing, 0)}). Where C fills a structure, the
caller passes a mutable structure (\code{getann(0, a)} with \code{a = Annotation()}).
Where C fills an \code{int} or a pointer, the caller passes a \code{Ref}:
\code{getseginfo(r)} takes a \code{Ref}, which receives a vector of \code{Seginfo} (or
\code{nothing}).

\paragraph{Strings} are returned as bytes, \code{Vector\{UInt8\}}, or \code{nothing} where
the C function may return \code{NULL} (\code{anndesc}, \code{getinfo}, \code{wfdbfile},
\ldots), except the results of \code{timstr}, \code{mstimstr} and \code{datstr}, which are
\code{String}s; unlike C's pointers to static buffers, they stay valid after later calls.

\paragraph{Numbers.} The integer types are those of the C reference platform:
\code{WFDB\_Sample} and C's \code{int} are \code{Int32}, \code{WFDB\_Time},
\code{WFDB\_Date} and C's \code{long} are \code{Int64}, \code{WFDB\_Frequency} and
\code{WFDB\_Gain} are \code{Float64}, C's \code{char} is \code{Int8}, and
\code{WFDB\_Signal}, \code{WFDB\_Group} and \code{WFDB\_Annotator} are \code{UInt32}.
Julia's fixed-width arithmetic wraps around as C's does. Arrays are indexed from 1, but
the values that C counts from 0 (signal numbers, annotator numbers, sample numbers) keep
their C values.

\paragraph{Return values and errors.} The functions return what the C functions return:
an \code{Int32} status (negative on errors), a number of signals, a \code{Float64}, a
string or \code{nothing}. Errors are reported by these values and by a message on the
standard error, as in C, never by exceptions. \code{wfdberror()} returns the last message;
\code{wfdbquiet()} turns the messages off and \code{wfdbverbose()} on again.
\code{WFDBX.wfdb\_error(fmt, args...)} formats its message as C's \code{printf}.

\paragraph{State.} As in C, a \code{Wfdb} object keeps its state between calls: the open
record, its input signals, the output signals and the annotators remain open until
\code{wfdbquit()} closes them. \code{wfdbquit()} leaves the input and output directories
and the network setting as they are. The high-level functions of \cref{sec:julia-api}
(except the metadata functions) call \code{wfdbquit()} on the \code{Wfdb} object they use,
which closes anything opened with the library functions on that object.

\paragraph{Additions.} \code{getvecs(out, n)} reads up to \code{n} sample vectors at once
into \code{out}, a \code{Vector\{Int32\}} with room for \code{n} times the number of
values of one vector, with exactly the result of \code{n} calls of \code{getvec}, but
decoding whole blocks of the signal files; it returns the number of vectors read
($\ge 0$), or the negative status of \code{getvec} if the first call failed.
\code{getvecs\_status(out, n)} also gives the status of the last \code{getvec} call
(\cref{sec:library-additions}):
\begin{lstlisting}[language=Julia]
getvecs([w], out::Vector{Int32}, n) -> Int64
getvecs_status([w], out::Vector{Int32}, n) -> (Int64, Int32)
\end{lstlisting}
\code{WfdbFile} is C's \code{FILE *} (or \code{WFDB\_FILE}): \code{wfdb\_open} and
\code{wfdb\_fopen} return a \code{WfdbFile} or \code{nothing}, read and written with the
functions of the C standard library (\code{getc}, \code{fread!}, \code{fgets!},
\code{fwrite}, \code{fseek}, \code{fclose}, \ldots).

\subsection{The structures}
\label{sec:julia-structures}

The structures of the C library have the names of the C structures without the prefix
\code{WFDB\_}, and the C field names; they are mutable, with keyword constructors whose
defaults are C's zero values (\cref{tab:julia-structures}).

\begin{xltabular}{\linewidth}{@{}lL@{}}
\caption{The C structures in Julia}
\label{tab:julia-structures} \\
\toprule
\tabhead Structure & \tabhead Fields \\
\midrule
\endfirsthead
\toprule
\tabhead Structure & \tabhead Fields \\
\midrule
\endhead
\bottomrule
\endfoot
\code{Siginfo} & \code{fname}, \code{desc}, \code{units}
  (\code{Union\{Nothing,Vector\{UInt8\}\}}), \code{gain}, \code{initval}, \code{group},
  \code{fmt}, \code{spf}, \code{bsize}, \code{adcres}, \code{adczero}, \code{baseline},
  \code{nsamp}, \code{cksum} (\cref{tab:julia-siginfo}) \\
\code{Seginfo} & \code{recname} (the name of the segment, bytes), \code{nsamp} (its
  number of samples, \code{Int64}), \code{samp0} (the number of its first sample in the
  record, \code{Int64}) \\
\code{Anninfo} & \code{name} (the annotator name, bytes), \code{stat} (\code{Int32}:
  \code{WFDB\_READ} or \code{WFDB\_WRITE} for the MIT format, \code{WFDB\_AHA\_READ} or
  \code{WFDB\_AHA\_WRITE} for the AHA format); \code{Anninfo("atr", WFDB\_READ)} \\
\code{Annotation} & \code{time} (\code{Int64}, in sample intervals), \code{anntyp}
  (\code{Int8}), \code{subtyp} (\code{Int8}), \code{chan} (\code{UInt8}), \code{num}
  (\code{Int8}), \code{aux} (\code{nothing}, or bytes whose first byte is the length of
  the text that follows, as in C) \\
\code{Calinfo} & \code{low}, \code{high} (the levels of the calibration pulse, in physical
  units, \code{Float64}), \code{scale} (physical units per cm), \code{sigtype} and
  \code{units} (bytes or \code{nothing}), \code{caltype} (\code{Int32}) \\
\end{xltabular}

\subsection{The settings functions}
\label{sec:julia-settings}

The functions of \cref{tab:julia-settings} are additions of WFDB-x to the library: they
set where files are read and written, network access and legacy mode, for the
high-level functions, the library functions and the tools alike
(\cref{chap:files,chap:behavior}). Like the other library functions, each takes an
optional \code{Wfdb} object first, except \code{wfdbx\_write\_errors} and
\code{wfdbx\_note\_write\_error} (the count belongs to the process). Each \code{set}
function returns the previous setting, so that it can be restored. The C function
\code{setwfdb(p)} sets the input directory too (\code{p} must name at most one directory),
and \code{getwfdb()} returns it (\code{"."} if none is set); in legacy mode they set and
return the WFDB path, as in C.

\begin{xltabular}{\linewidth}{@{}lL@{}}
\caption{The settings functions (\code{[w]}: the optional \code{Wfdb} object)}
\label{tab:julia-settings} \\
\toprule
\tabhead Function & \tabhead Effect and result \\
\midrule
\endfirsthead
\toprule
\tabhead Function & \tabhead Effect and result \\
\midrule
\endhead
\bottomrule
\endfoot
\code{wfdbx\_set\_input\_dir([w], dir::OptByteStr)} & sets the input directory: a local
  directory, a \code{\textasciitilde} path or a URL; \code{nothing} or \code{""} restores
  the default (\code{WFDBX\_INPUT\_DIR}, or else the current directory). Returns the
  previous input directory as bytes, or \code{nothing} if none was set \\
\code{wfdbx\_get\_input\_dir([w])} & returns the input directory as bytes, or
  \code{nothing} if none is set \\
\code{wfdbx\_set\_output\_dir([w], dir::OptByteStr)} & sets the output directory: a local
  directory or a \code{\textasciitilde} path (a URL is refused, with a message, and the
  setting is unchanged); \code{nothing} or \code{""} restores the default
  (\code{WFDBX\_OUTPUT\_DIR}, or else the current directory). Returns the previous output
  directory, or \code{nothing} \\
\code{wfdbx\_get\_output\_dir([w])} & returns the output directory as bytes, or
  \code{nothing} if none is set \\
\code{wfdbx\_set\_network([w], on::Union\{Nothing,Bool,Integer\})} & turns network access
  on (\code{true} or $> 0$) or off (\code{false} or 0); \code{nothing} or a negative
  number restores the default (on if \code{WFDBX\_NETWORK} is a nonzero number). Returns
  the previous setting, 1 or 0 (\code{Int32}) \\
\code{wfdbx\_get\_network([w])} & returns 1 if network access is on, 0 otherwise \\
\code{wfdbx\_set\_legacy([w], on::Union\{Bool,Integer\})} & turns legacy mode on
  (\code{true} or nonzero) or off. Returns the previous setting, 1 or 0 \\
\code{wfdbx\_get\_legacy([w])} & returns 1 if legacy mode is on, 0 otherwise \\
\code{wfdbx\_legacy\_args([w], args::Vector\{String\})} & if \code{args[1]} (the
  arguments \emph{without} the program name, as Julia's \code{ARGS}) is
  \code{"-{}-legacy"}, removes it, turns legacy mode on and returns 1; otherwise returns
  0 \\
\code{wfdbx\_time\_needs\_freq([w], string::ByteStr)} & returns 1 if
  \code{strtim(string)} needs a sampling frequency, 0 otherwise \\
\code{wfdbx\_write\_errors()} & returns the number of write errors since the program
  started (output files that could not be written, flushed or closed); under
  \code{WFDBX.run\_tool}, since the tool started \\
\code{wfdbx\_note\_write\_error()} & counts a write error, for a program that writes a file
  itself \\
\code{wfdbx\_input\_path([w], name::OptByteStr)} & returns, as bytes, the path at which a
  program opens a file that it reads itself (not through the library): where
  \code{wfdb\_open} finds a file of that name (in normal mode, a relative name under the
  input directory, a \code{\textasciitilde} name expanded, an absolute name as given; in
  legacy mode, the name as given). \code{""} and \code{"-"} are returned as given;
  \code{nothing} for \code{nothing} or a URL (also one made by the input directory) \\
\code{wfdbx\_output\_path([w], name::OptByteStr)} & returns, as bytes, the path at which a
  program creates a file that it writes itself: where \code{wfdb\_open} creates a file of
  that name (in normal mode, a relative name under the output directory, a
  \code{\textasciitilde} name expanded, an absolute name as given; in legacy mode, any
  name under the output directory, used as a plain prefix). When the output directory has
  been prefixed, the missing directories of the path are created. \code{""} and
  \code{"-"} are returned as given; \code{nothing} for \code{nothing} or a URL \\
\end{xltabular}

\subsection{Examples}

Reading the samples and the annotations of a record, as a C program does (the record
\code{100s} is read from the current directory):

\begin{lstlisting}[language=Julia,style=example]
si = [Siginfo() for _ in 1:2]                       # receives one Siginfo per signal
nsig = isigopen("100s", si, 2)                      # open 2 signals
println(nsig, " ", String(si[1].desc), " ", si[1].fmt, " ", si[1].gain)  # 2 MLII 212 200.0
println(sampfreq(nothing))                          # 360.0: the frequency of the open record

v = zeros(Int32, 2)                                 # receives one sample per signal
isigsettime(strtim("0:30"))                         # skip to 30 seconds
getvec(v)
println(v[1], " ", v[2], " ", aduphys(0, v[1]))     # 947 968 -0.385

frames = zeros(Int32, 360 * 2)                      # the next second, in one call
println(getvecs_status(frames, 360))                # (360, 2)

annopen("100s", [Anninfo("atr", WFDB_READ)], 1)
a = Annotation()
n = 0
while getann(0, a) == 0                             # 0: an annotation was read; -1: the end
    global n += 1
    n <= 2 && println(mstimstr(a.time), " ", String(annstr(a.anntyp)))
end
#     0:00.050 +
#     0:00.214 N
println(n, " annotations")                          # 75 annotations

wfdbquiet()                                         # no messages on the standard error
stat = isigopen("nosuchrecord", nothing, 0)
print(stat, " ", String(wfdberror()))
# -1 init: can't open header for record nosuchrecord
wfdbverbose()
wfdbquit()                                          # close everything
\end{lstlisting}

Writing a record and an annotation file, sample by sample and annotation by annotation:

\begin{lstlisting}[language=Julia,style=example]
out = Siginfo(fname=b"ramp.dat", desc=b"ramp", units=b"mV", gain=200.0, fmt=16, adcres=12)
println(osigfopen([out], 1))                        # 1: one output signal
setsampfreq(250.0)
for k in 0:499
    putvec(Int32[k])
end
println(newheader("ramp"))                          # 0: ramp.hea is written
annopen("ramp", [Anninfo("qrs", WFDB_WRITE)], 1)
for t in (100, 350)
    putann(0, Annotation(time=t, anntyp=WFDBX.NORMAL))
end
wfdbquit()                                          # completes ramp.dat and ramp.qrs
println(length(read_annotations("ramp", "qrs")))    # 2
\end{lstlisting}

% ==========================================================================================
\section{The command-line tools in Julia}
\label{sec:julia-tools}

The 33 command-line tools (\code{rdsamp}, \code{rdann}, \code{wrsamp}, \code{xform},
\ldots) are functions of the package, one file per tool in \file{julia/src/apps}, run
through the driver script \file{julia/bin/wfdbx.jl}: its first argument names the tool,
and the other arguments go to the tool. Each follows its C program line by line. The
tools, their options, their output and their exit status are those of the C tools, and
are documented in \cref{chap:tools}. With the shell script \file{wfdbx} of
\cref{sec:julia-install}:

\begin{lstlisting}[style=shell]
export WFDBX_INPUT_DIR=/path/to/wfdb-x/tests/data
wfdbx rdsamp -r 100s -t 10 -p     # the first 10 seconds of 100s, in physical units
wfdbx rdann -r 100s -a atr        # its annotations
wfdbx rdsamp --legacy -r 100s     # legacy mode: --legacy as the first argument
wfdbx --version                   # WFDB-x 1.2.0
\end{lstlisting}

Without the script, the command is \cmd{julia -{}-startup-file=no
-{}-project=/path/to/wfdb-x/julia /path/to/wfdb-x/julia/bin/wfdbx.jl TOOL ARGUMENTS}.
Every run starts a new Julia process, which loads the precompiled package; most of the
start-up time of a tool is Julia's own. With no arguments, or an unknown tool name, the
driver prints a message on the standard error and exits with status 1. When an output
file or the standard output cannot be written (a full disk, for example), a tool reports
it (``PROGRAM: write error on standard output'') and exits with status 1; with
\opt{-{}-legacy}, it keeps the original exit status. A tool whose standard output is a
closed pipe ends with status 141, as a C program ended by \code{SIGPIPE}. The driver finds
the compiled package in the depot of the user's own home directory even when the
environment variable \code{HOME} is changed (\code{JULIA\_DEPOT\_PATH}, if set, is used as
given).

\paragraph{Inside a Julia session.} \code{WFDBX.run\_tool(tool, args; stdin=UInt8[],
dir=nothing, env=nothing)} runs a tool in the current process, with a fresh library state,
in the directory \code{dir}, with the environment variables of \code{env} set for the run
and the bytes \code{stdin} as its standard input, and returns its exit status and what it
wrote to the standard output and the standard error, as bytes. The caller's library state
and standard streams are restored when the tool returns. \code{WFDBX.tools()} returns the
names of the tools, and an unknown name is an \code{ArgumentError}.

\begin{lstlisting}[language=Julia,style=example]
status, out, err = WFDBX.run_tool("rdsamp", ["-r", "100s", "-t", "s3"]; dir=".")
println(status)                                     # 0
print(String(out))       # 3 lines: the sample number, then 995 and 1011 (tab-separated)
println(length(WFDBX.tools()))                      # 33
\end{lstlisting}

\paragraph{Environment variables.} The tools, like the rest of the library, use the
environment variables of \cref{tab:julia-environment}, also for the files that they read
or create themselves: the \opt{-i} file of \code{wrsamp}, \code{edf2mit},
\code{wav2mit}, \code{a2m} and \code{ad2m}, the input files of \code{ahaecg2mit} and the
\opt{-S} script of \code{xform} are read from the input directory (\code{wfdbdesc}
measures a signal file in the directory of its header), and the WAV file of
\code{mit2wav}, the EDF file of \code{mit2edf}, the master header of \code{snip -m} and
the \file{.des}, \file{.ecg}, \file{.key} and \file{.txt} files of \code{parsescp} are
created in the output directory, as \code{wfdbx\_input\_path} and
\code{wfdbx\_output\_path} locate them (\cref{sec:julia-settings}). Such a file cannot be
read from a URL. In legacy mode, the files that a tool reads itself are opened as named.

\begin{table}[htbp]
\centering
\caption{The environment variables used by the Julia library}
\label{tab:julia-environment}
\begin{tabularx}{\linewidth}{@{}lL@{}}
\toprule
\tabhead Variable & \tabhead Meaning \\
\midrule
\code{WFDBX\_INPUT\_DIR} & the input directory (a directory, a \code{\textasciitilde}
  path or a URL), when \code{wfdbx\_set\_input\_dir} does not set one \\
\code{WFDBX\_OUTPUT\_DIR} & the output directory, when \code{wfdbx\_set\_output\_dir} does
  not set one \\
\code{WFDBX\_NETWORK} & network access is on if it is a nonzero number, when
  \code{wfdbx\_set\_network} does not set it \\
\code{WFDBANNSORT} & \code{0} to keep annotation files in the order written (no sorting
  when they are closed) \\
\code{WFDBGVMODE} & the default getvec mode (\code{setgvmode} with a negative mode) \\
\code{WFDBCAL} & the calibration file read by \code{calopen(nothing)} (default
  \file{wfdbcal}) \\
\code{WFDB} & the WFDB path, used in legacy mode only (outside legacy mode, it is
  ignored, with a message) \\
\code{HOME} (\code{USERPROFILE} on Windows) & the home directory, for
  \code{\textasciitilde} paths \\
\bottomrule
\end{tabularx}
\end{table}

Tools that write signal parameters require them as options: \code{wrsamp} (with
\opt{-o}) requires \opt{-F} (frequency), \opt{-G} (gain), \opt{-R} (ADC resolution) and,
for calibrated signals, \opt{-U} (units); \code{wav2mit} requires \opt{-G}, and \opt{-U}
unless the gain is 0; \code{ad2m} and \code{ahaecg2mit} require \opt{-F}, \opt{-G},
\opt{-R}, and \opt{-U} unless the gain is 0; \code{a2m}, \code{rdann}, \code{ann2rr},
\code{sortann} and \code{m2a} require \opt{-F} when the record's header gives no
frequency (\cref{chap:tools}).

% ==========================================================================================
\section{Testing}
\label{sec:julia-testing}

The unit tests of the Julia implementation are in \file{julia/test}, and use the standard
library \code{Test}. From the root of the repository:

\begin{lstlisting}[style=shell]
julia --startup-file=no --project=julia julia/test/runtests.jl   # all the unit tests
julia --project=julia -e 'using Pkg; Pkg.test()'                  # the same, through Pkg
\end{lstlisting}

The tests cover the C standard functions (\code{printf}, \code{strtol}, \code{strtod},
\code{sscanf}, \ldots, against glibc), the files and the library functions
(\file{runtests.jl}, \file{signals.jl}, \file{test\_annotations.jl}), the FLAC codec
(\file{test\_flac.jl}), finding files, the settings and network access
(\file{test\_files.jl}), the high-level functions, the metadata functions
(\file{test\_metadata.jl}, with the cases shared by the seven languages, in
\file{tests/metadata}), \code{check\_record} (\file{test\_check.jl}) and the tools run
in-process. They run in temporary copies of the test records of \file{tests/data}. The
tests of network access use a local HTTP server on 127.0.0.1 only (written with the
standard library \code{Sockets}). The tests of \code{check\_record} write the corpus of
\file{python/tests/check\_corpus.py} with \cmd{python3}, NumPy and the Python package of
\file{python/src}, and compare the reports with its \file{expected.txt}; they are skipped
if that script cannot run.

The cross-language conformance suite runs every tool and compares its results, byte for
byte, with those of the original C programs (1135 cases), and, with \opt{-{}-legacy}, in
legacy mode (883 cases). It runs the tools through \file{julia/bin/wfdbx.jl}:

\begin{lstlisting}[style=shell]
PYTHONPATH=$PWD/python/src python3 tests/conformance/harness.py check --impl julia -j8
PYTHONPATH=$PWD/python/src python3 tests/conformance/harness.py check --impl julia -j8 --legacy
\end{lstlisting}

% The R chapter of the WFDB-x user's manual. It follows the sections of the Python chapter
% (python.tex), in the same order. Each \apientry of the section "The high-level API" (and
% of the section on the command-line tools) documents one function exported by r/NAMESPACE
% from r/R/highlevel.R, highlevel_annot.R, metadata.R, check.R and cli.R: its signature
% listing is the usage of the function, with the argument names and defaults of its
% definition in the code, and its params name every formal argument, in code order. The
% listings of style "example" were run with Rscript (the package installed from r/), in a
% copy of tests/data as the current directory, and print what their comments say.
\chapter{R}
\label{chap:r}
\chapterlanguage{r}

The R implementation of WFDB-x is the package \code{wfdbx}, in the directory \file{r} of
the repository. It is written in base R alone: it needs no other package, has no compiled
code, and never calls the C library. It reads and writes the same files as the original C
library, and its tools print the same output.

It offers WFDB-x in three ways:
\begin{itemize}
\item \textbf{The high-level API} (\cref{sec:r-api}): functions that read and write whole
  records, headers, annotation files and header metadata as R matrices, lists and data
  frames, in one call, with named arguments, and that stop with R errors on failure. They
  read signal files in large blocks with vectorized code, and are the fastest way to read
  a record. They are an addition of WFDB-x: the C package has no equivalent. They are the
  recommended interface for new code and analysis.
\item \textbf{The WFDB library functions} under their C names (\code{isigopen},
  \code{getvec}, \code{getann}, \code{putann}, \ldots), with the behavior of the C library
  (\cref{sec:r-library}). They are documented once, for all the languages, in
  \cref{chap:library}.
\item \textbf{The command-line tools} (\code{rdsamp}, \code{rdann}, \code{wrsamp},
  \code{xform} and 29 others), with the options, output and exit status of the C tools,
  run with \cmd{Rscript} through the driver script \file{wfdbx.R}, or from R with
  \code{wfdbx\_tool} (\cref{sec:r-tools}). They are documented in \cref{chap:tools}.
\end{itemize}
The high-level functions return exactly the samples and annotations that the library
functions (and the C library) give.

% ==========================================================================================
\section{Installation and requirements}
\label{sec:r-install}

\subsection{Requirements}

WFDB-x for R needs R 4.1 or later, and nothing else: the package uses base R only (its
one import, \code{tools}, is part of R). It has no compiled code and needs no C compiler,
no WFDB library and no libFLAC (the FLAC formats 508, 516 and 524 are read and written by
WFDB-x's own codec, in R). It runs on Linux and macOS. On Windows, the command-line tools
write their output as text, so NUL bytes in it are lost (\cref{sec:r-tools}).

Two programs are used if they are installed: \cmd{curl}, to download remote files (base
R's \code{url()} is used otherwise), and \cmd{stat}, to recognize an output file that is
also an input file under another path. The tests also use \cmd{python3}
(\cref{sec:r-testing}).

\subsection{Installing the package}

From the root of a copy of the repository, the package is in the directory \file{r}:

\begin{lstlisting}[style=shell]
R CMD INSTALL r                   # install the package wfdbx
\end{lstlisting}

or, from R, with the path of the directory \file{r}:

\begin{lstlisting}[language=R]
install.packages("r", repos = NULL, type = "source")
library(wfdbx)
\end{lstlisting}

The package is attached with \code{library(wfdbx)}. Its version is given by
\code{packageVersion("wfdbx")}, and by the driver script of the tools:

\begin{lstlisting}[style=shell]
Rscript -e 'cat(format(packageVersion("wfdbx")), "\n")'    # 1.2.1
Rscript r/inst/bin/wfdbx.R --version                        # WFDB-x 1.2.1
\end{lstlisting}

The help pages of the package (\code{?read\_record}, \code{?read\_metadata},
\code{?isigopen}, \ldots) document each function; \code{?wfdbx} gives an overview.

\subsection{The command-line tools}

The tools are run by the driver script \file{wfdbx.R} with \cmd{Rscript}. In a copy of the
repository, the script is \file{r/inst/bin/wfdbx.R}, and it runs the tools without the
package being installed (it then loads the sources of \file{r/R}); in an installed package,
its path is \code{system.file("bin", "wfdbx.R", package = "wfdbx")}
(\cref{sec:r-tools}).

\subsection{The sample record}
\label{sec:r-resources}

The repository includes the sample record \code{100s}, the sample record of the WFDB
Software Package: the first minute of record 100 of the MIT-BIH Arrhythmia Database, two
ECG signals (\code{MLII} and \code{V5}) sampled at 360\,Hz, with the reference beat
annotations (annotator \code{atr}). Its files (\file{100s.hea}, \file{100s.dat},
\file{100s.atr}) are in the directory \file{tests/data}, with the other test records of
WFDB-x; the R package does not install them. Work that uses the MIT-BIH Arrhythmia
Database cites it as described on PhysioNet.

% ==========================================================================================
\section{Quick start}
\label{sec:r-quickstart}

This example reads the header, ten seconds of the signals and the annotations of the
sample record \code{100s}, from the directory \file{tests/data} of a copy of the
repository (\code{input\_dir}; the path is that of the copy), and computes the mean heart
rate. It then writes the first ten seconds as a new record, \code{first10}, in the current
directory, with two metadata entries in its header.

\begin{lstlisting}[language=R,style=example]
library(wfdbx)

data <- "wfdb-x/tests/data"             # the directory of the sample record 100s

# The header only: no samples are read.
hdr <- read_header("100s", input_dir = data)
cat(hdr$nsig, hdr$fs, hdr$nsamples, "\n")           # 2 360 21600
print(sapply(hdr$siginfo, function(s) s$desc))     # "MLII" "V5"

# Seconds 10 to 20 of both signals, in physical units: one row per sample, one
# column per signal.
rec <- read_record("100s", sampfrom = "0:10", sampto = "0:20", physical = TRUE,
                   input_dir = data)
print(dim(rec$signals))                 # 3600 2
print(rec$signals[1, ])                 # -0.390 -0.275

# The beat annotations, and the mean heart rate.
ann <- read_annotations("100s", "atr", input_dir = data)
cat(nrow(ann), head(ann$symbol, 3), head(ann$sample, 3), "\n")   # 75 + N N 18 77 370
beats <- ann$sample[isqrs(ann$code) == 1]
rr <- diff(beats) / attr(ann, "fs")     # RR intervals, in seconds
cat(length(beats), "beats,", round(60 / mean(rr), 1), "beats per minute\n")   # 74 beats, 73.9 ...

# Write the first 10 seconds, in ADC units, as the record "first10" (first10.hea and
# first10.dat, in the current directory), with the signal specifications of 100s.
x <- read_record("100s", sampto = "0:10", input_dir = data)
write_record("first10", x$signals, fs = x$fs, fmt = 16, gains = x$gains,
             baselines = x$baselines,
             adczeros = sapply(x$siginfo, function(s) s$adczero),
             adcres = sapply(x$siginfo, function(s) s$adcres), units = "mV",
             names = x$names, metadata = list(Source = "MIT-BIH record 100"))

# Header metadata: "# Key: value" lines of the header.
set_metadata("first10", "Age", "69")
print(read_metadata("first10"))         # c(Source = "MIT-BIH record 100", Age = "69")
print(check_record("first10"))          # no problems found
\end{lstlisting}

The header \file{100s.hea} gives no units: the example writes \code{"mV"} explicitly,
since WFDB-x never assumes units, a gain, an ADC resolution or a sampling frequency that
a record does not give (\cref{chap:behavior}).

% ==========================================================================================
\section{The high-level API}
\label{sec:r-api}

The high-level API is made of the functions of this section and of the lists and data
frames that they return (\cref{sec:r-results}). They are defined in the files
\file{highlevel.R}, \file{highlevel\_annot.R}, \file{metadata.R} and \file{check.R} of
\file{r/R}, and documented by the help pages \code{?read\_record} (which also covers
\code{read\_header} and \code{write\_record}), \code{?read\_annotations} (also
\code{write\_annotations}), \code{?check\_record} and \code{?read\_metadata} (the five
metadata functions).

\subsection{Conventions}
\label{sec:r-conventions}

\paragraph{Record names and directories.} A record is named as in the other parts of
WFDB-x (\cref{chap:records,chap:files}), without the \code{.hea} extension (the readers
also accept a trailing \code{.hea}; \code{write\_record} does not). A record name is a
full path, absolute (\code{"/data/mitdb/100"}) or in the home directory
(\code{"\textasciitilde/Documents/100s"}), used as given; a relative path, such as \code{"100s"} or
\code{"data/100s"}, read from the input directory if one is set, otherwise from the
current directory, and written to the output directory if one is set, otherwise to the
current directory (missing directories are created); or a URL
(\code{"https://physionet.org/files/mitdb/1.0.0/100"}), read only if network access is
on, which cannot be written. A writer given a name with a path creates the files in that
directory, and the header names the record and its signal file without the path
(\code{rec}, \code{rec.dat}), so that the record opens under any path that leads to it;
in legacy mode, as in the original, the header gives the name with its path. The input
and output directories are the current directory unless they are set, with \code{wfdbx\_set\_input\_dir} and
\code{wfdbx\_set\_output\_dir} (\cref{sec:r-settings}) or with the environment variables
\code{WFDBX\_INPUT\_DIR} and \code{WFDBX\_OUTPUT\_DIR}. URLs are read only when network
access is on. The functions that read take the arguments \code{input\_dir} and
\code{network}, and those that write take \code{dir}: they apply to that one call, and the
previous settings are restored when it returns, also after an error. \code{NULL}, the
default of these arguments, keeps the current setting (for \code{dir}, \code{""} too).
\code{input\_dir} applies to relative names only, not to full paths or URLs. With
\code{dir}, the name must be a plain name (\code{"rec"}): the files are created in
\code{dir}, and the record opens from that directory under its plain name.
The signal files that a header names are read from the header's directory. A name with
another extension, such as \code{"rec.edf"}, is read as an EDF file.

\paragraph{Times.} The arguments \code{sampfrom} and \code{sampto} are sample numbers
(whole numbers) or WFDB time strings, as \code{strtim} reads them: \code{"0:10"} (10
seconds), \code{"1:02:03.5"}, \code{"s3600"} (sample 3600), \code{"e"} (the end of the
record), or a time of day in brackets, \code{"[14:00:00]"}. A time in seconds needs the
sampling frequency; when the record does not give it, such a time is an error.

\paragraph{Per-signal and per-annotation arguments.} The arguments that describe each
signal (\code{fmt}, \code{gains}, \code{units}, \ldots) or each annotation
(\code{subtype}, \code{chan}, \ldots) take one value for all, or a vector (or, for
\code{units} and \code{names}, a list) with exactly one value per signal or per
annotation; any other length is an error.

\paragraph{Numbers and text.} Numbers are doubles: samples, times and counts are
\code{double} vectors (exact for every 32-bit sample and for times up to $2^{53}$).
Signal numbers start at 0, as in WFDB. Text is handled as bytes: what the files hold is
returned as character strings of their bytes, never re-encoded, and character strings
given are written as their bytes.

\paragraph{Signal parameters.} No sampling frequency, gain, ADC resolution or units is
ever assumed. Where one is needed and the record does not give it, it must be given as an
argument; otherwise the function stops with an error. A record whose header gives no
sampling frequency has \code{fs = 0}.

\paragraph{Legacy mode.} With \code{legacy = TRUE}, a function behaves as in the original
WFDB package, with its assumed signal parameters and its search for files
(\cref{chap:behavior}); the library's legacy mode is on for the duration of the call
only. With \code{legacy = FALSE} (the default), the current setting of
\code{wfdbx\_set\_legacy} applies, which is off unless it was turned on.

\paragraph{State.} The functions that read and write records and annotation files close
anything opened with the WFDB library functions (they call \code{wfdbquit()}) before
they start, and close their own files before they return, whatever the outcome.
\code{read\_record} sets the getvec mode for its own reading only. \code{read\_record},
\code{read\_header}, \code{read\_annotations}, \code{write\_record},
\code{write\_annotations} and \code{check\_record} leave the getvec mode
(\code{getgvmode}) and the time resolution of output annotators (\code{getafreq}) as they
were, also after an error: opening an EDF record sets both, as in C (\code{WFDB\_HIGHRES}
is added to the getvec mode, and the time resolution becomes the record's sampling
frequency when none is set), and \code{wfdbquit()} resets neither; the \code{fs} of
\code{write\_annotations} applies to its file only. The metadata functions open no
signals and leave the library's state as it is.

\paragraph{The examples} of this section read the test records of WFDB-x (among them
\code{100s}) from the current directory, a copy of \file{tests/data}, and write records of
their own there. They were run with \cmd{Rscript}, and print what their comments say.

% ------------------------------------------------------------------------------------------
\subsection{Reading records}
\label{sec:r-reading}

\apientry{read\_record}{Read the signals of a record}
\begin{lstlisting}[language=R,style=signature]
read_record(record, channels = NULL, sampfrom = 0, sampto = NULL,
            physical = FALSE, highres = FALSE, pad = FALSE,
            input_dir = NULL, network = NULL,
            legacy = FALSE, checksum_warning = TRUE)
\end{lstlisting}
Reads the signals of a record, or part of it, into a numeric matrix with one row per
sample and one column per signal, together with the description of the record. The
samples are those that \code{getvec} returns (\cref{chap:library}), read in large blocks.
Invalid samples are \code{WFDB\_INVALID\_SAMPLE} ($-32768$) in ADC units, and \code{NA}
in physical units. When the read reaches the end of the record, the samples are compared
with the checksums of the header (\cref{sec:r-warnings}).
\begin{params}
\item[record] Record name (a character string), without the \code{.hea} extension (a
  trailing \code{".hea"} is also accepted): a full path, absolute
  (\code{"/data/mitdb/100"}) or in the home directory
  (\code{"\textasciitilde/Documents/100s"}), used as given; a relative path (\code{"100s"},
  \code{"data/100s"}), read from the input directory if one is set, otherwise from the
  current directory; or a URL (\code{"https://physionet.org/files/mitdb/1.0.0/100"}),
  read only if network access is on.
\item[channels] The signals to read, in the order in which they are returned: signal
  numbers (from 0) or signal descriptions (as in \code{names}, such as \code{"MLII"}).
  Default \code{NULL}: all the signals, in the order of the header.
\item[sampfrom] The first sample to read: a sample number (a whole number $\ge 0$) or a
  time string. Default \code{0}: the beginning of the record.
\item[sampto] The sample after the last one to read: a sample number (a whole number) or
  a time string. A value at or before \code{sampfrom} reads no samples. Default
  \code{NULL}: the end of the record.
\item[physical] \code{FALSE} (the default) to return ADC units; \code{TRUE} to return
  physical units, computed as $(\mathit{adc} - \mathit{baseline}) / \mathit{gain}$. Every
  signal read must then have a gain.
\item[highres] \code{FALSE} (the default) to read a multi-frequency record at its frame
  frequency, with the samples of the faster signals averaged to one per frame;
  \code{TRUE} to read every sample, at the highest sampling frequency of the record (the
  samples of the slower signals are repeated). \code{fs} of the result is then that
  frequency.
\item[pad] \code{TRUE} to replace each invalid sample by the previous valid sample of the
  same signal. Default \code{FALSE}.
\item[input\_dir] The input directory for this call only (optional): a local directory, a
  \code{\textasciitilde} path or a URL, under which relative record names are read. It does not
  apply to full paths or URLs. Default \code{NULL}: the current setting
  (\code{wfdbx\_set\_input\_dir}).
\item[network] \code{TRUE} to allow reading URLs for this call only, \code{FALSE} to forbid
  it. Default \code{NULL}: the current setting (\code{wfdbx\_set\_network}).
\item[legacy] \code{TRUE} to behave as the original WFDB package (legacy mode, for this
  call only). Default \code{FALSE}.
\item[checksum\_warning] \code{TRUE} (the default) to issue an R warning when the samples
  do not match the checksums of the header; \code{FALSE} for no warning. The samples are
  returned either way, and \code{checksum\_ok} of the result is \code{FALSE}.
\end{params}
\returns{A list (\cref{tab:r-record}).}
\raises{Stops with an error if the header or a signal file cannot be read (``can't read
record REC''), if the record has no signals, if \code{sampfrom} cannot be reached (``can't
seek to sample N of record REC''), or if a signal file is shorter than its header says
(``can't read signal data of record REC (getvec status -3)''); if \code{sampfrom} is
negative (``sampfrom must not be negative''), or \code{sampfrom} or \code{sampto} is not a
whole number (a fraction, \code{NA}, \code{NaN} or an infinite value: ``sampfrom must be
a whole number''); if a channel is not in the record (``record REC has no signal C''), if
a time string needs the sampling frequency and the header gives none (``sampling
frequency of record REC is not specified''), or if \code{physical = TRUE} and a signal
read has no gain (``signal K of record REC is uncalibrated''). Warns on a checksum error
(see \code{checksum\_warning}).}
\begin{lstlisting}[language=R,style=example]
library(wfdbx)

# Seconds 10 to 20 of lead MLII, in physical units.
rec <- read_record("100s", channels = "MLII", sampfrom = "0:10", sampto = "0:20",
                   physical = TRUE)
print(dim(rec$signals))                  # 3600 1
cat(rec$names, rec$fs, rec$sampfrom, "\n")   # MLII 360 3600

# The first 3 samples of both signals, in ADC units, V5 first.
rec <- read_record("100s", channels = c(1, 0), sampto = 3)
print(rec$signals)                       # three rows of 1011 995
print(rec$gains)                         # 200 200
print(rec$baselines)                     # 1024 1024
str(rec$units)                           # a list of two NULL: no units in the header

# The whole record: the checksums are verified.
rec <- read_record("100s")
cat(dim(rec$signals), rec$checksum_ok, "\n")   # 21600 2 TRUE
\end{lstlisting}
A multi-frequency record (here the test record \code{mfreq}, whose signals have 1, 2
and 4 samples per frame at a frame frequency of 100\,Hz) is read at its frame frequency,
or at its highest sampling frequency:
\begin{lstlisting}[language=R,style=example]
library(wfdbx)

rec <- read_record("mfreq")
cat(rec$fs, dim(rec$signals), "\n")      # 100 600 3
rec <- read_record("mfreq", highres = TRUE)
cat(rec$fs, dim(rec$signals), "\n")      # 400 2400 3
print(sapply(rec$siginfo, function(s) s$spf))   # 1 2 4
\end{lstlisting}

\apientry{read\_header}{Read the header of a record}
\begin{lstlisting}[language=R,style=signature]
read_header(record, input_dir = NULL, network = NULL,
            legacy = FALSE)
\end{lstlisting}
Reads the header of a record, without opening its signal files. The result describes the
record and each of its signals, and holds the info strings and the metadata entries of the
header. For a multi-segment record, the master header is read, and \code{nsamples} is the
length of the whole record. The metadata are read from the \file{.hea} file only; an EDF
file, which has no \file{.hea} file, and a header read from the standard input
(\code{"-"}) have none.
\begin{params}
\item[record] Record name (a character string), without the \code{.hea} extension (a
  trailing \code{".hea"} is also accepted): a full path, absolute
  (\code{"/data/mitdb/100"}) or in the home directory
  (\code{"\textasciitilde/Documents/100s"}), used as given; a relative path (\code{"100s"},
  \code{"data/100s"}), read from the input directory if one is set, otherwise from the
  current directory; or a URL (\code{"https://physionet.org/files/mitdb/1.0.0/100"}),
  read only if network access is on.
\item[input\_dir] The input directory for this call only (optional): a local directory, a
  \code{\textasciitilde} path or a URL, under which relative record names are read. It does not
  apply to full paths or URLs. Default \code{NULL}: the current setting
  (\code{wfdbx\_set\_input\_dir}).
\item[network] \code{TRUE} to allow reading URLs for this call only, \code{FALSE} to
  forbid it. Default \code{NULL}: the current setting.
\item[legacy] \code{TRUE} to read the header as the original WFDB package does. Default
  \code{FALSE}.
\end{params}
\returns{A list of class \code{"wfdbx\_header"} (\cref{tab:r-header}).}
\raises{Stops with an error if the header cannot be found or read (``can't read header
for record REC''). Warns, as \code{read\_metadata} does, for each key of which the header
has more than one entry.}
\begin{lstlisting}[language=R,style=example]
library(wfdbx)

hdr <- read_header("100s")
cat(hdr$name, hdr$fs, hdr$nsig, hdr$nsamples, "\n")   # 100s 360 2 21600
for (s in hdr$siginfo)
  cat(s$desc, s$fname, s$fmt, s$gain, s$adcres, is.null(s$units), "\n")
# MLII 100s.dat 212 200 11 TRUE
# V5 100s.dat 212 200 11 TRUE
print(hdr$comments[1])          # " 69 M 1085 1629 x1": the first info string
print(is.null(hdr$base_time))   # TRUE: the header gives no start time
print(hdr$metadata)
# c("Produced by xform from record 100, beginning at 0" = "0"): the last info string
# contains a colon, so it is an entry
\end{lstlisting}

\apientry{read\_annotations}{Read an annotation file}
\begin{lstlisting}[language=R,style=signature]
read_annotations(record, annotator, sampfrom = 0, sampto = NULL, fs = NULL,
                 input_dir = NULL, network = NULL,
                 legacy = FALSE)
\end{lstlisting}
Reads the annotations of an annotation file, in file order, into a data frame with one
row per annotation. The sampling frequency, needed to convert times in seconds, is that of
the record's header, or, for a record without a header or whose header gives none, the
argument \code{fs}.
\begin{params}
\item[record] Record name (a character string), without the \code{.hea} extension (a
  trailing \code{".hea"} is also accepted): a full path, absolute
  (\code{"/data/mitdb/100"}) or in the home directory
  (\code{"\textasciitilde/Documents/100s"}), used as given; a relative path (\code{"100s"},
  \code{"data/100s"}), read from the input directory if one is set, otherwise from the
  current directory; or a URL (\code{"https://physionet.org/files/mitdb/1.0.0/100"}),
  read only if network access is on. Its header, if there is one, gives
  the sampling frequency.
\item[annotator] Annotator name: the suffix of the annotation file. \code{"atr"} reads the
  file \file{RECORD.atr}.
\item[sampfrom] The first time to read: a sample number or a time string. Default
  \code{0}: the beginning of the file.
\item[sampto] The time at which to stop: annotations at or after it are not read. A
  sample number or a time string. Default \code{NULL}: the end of the file.
\item[fs] The sampling frequency, in Hz (a number $> 0$), to use when the record has no
  header or its header gives none; it is ignored otherwise. Default \code{NULL}: none.
\item[input\_dir] The input directory for this call only (optional): a local directory, a
  \code{\textasciitilde} path or a URL, under which relative record names are read. It does not
  apply to full paths or URLs. Default \code{NULL}: the current setting
  (\code{wfdbx\_set\_input\_dir}).
\item[network] \code{TRUE} to allow reading URLs for this call only, \code{FALSE} to
  forbid it. Default \code{NULL}: the current setting.
\item[legacy] \code{TRUE} to behave as the original WFDB package (which assumes 250\,Hz
  for a record without a frequency). Default \code{FALSE}.
\end{params}
\returns{A data frame (\cref{tab:r-annotations}), with the attributes \code{record},
\code{annotator} and \code{fs}. \code{attr(ann, "fs")} is the frequency in which the times
are counted: the one written in the annotation file, if there is one, or else the
sampling frequency; it is 0 if none of them is known.}
\raises{Stops with an error if the annotation file cannot be opened (``can't read
annotator ANN for record REC'') or read, or if it ends before its end-of-file marker
(``can't read annotation file ANN of record REC (getann status -3)''); if \code{fs} is
given and is not $> 0$ (``fs must be > 0''); or if a time string needs the sampling
frequency and none is known.}
\begin{lstlisting}[language=R,style=example]
library(wfdbx)

ann <- read_annotations("100s", "atr", sampfrom = "0:10", sampto = "0:30")
cat(nrow(ann), attr(ann, "fs"), "\n")    # 24 360
print(head(ann, 3))
#   sample symbol code subtype chan num  aux
# 1   3862      N    1       0    0   0 <NA>
# 2   4170      N    1       0    0   0 <NA>
# 3   4466      N    1       0    0   0 <NA>

# The whole file: the first annotation is a rhythm change, with an aux string.
ann <- read_annotations("100s", "atr")
cat(ann$symbol[1], ann$code[1] == RHYTHM, ann$aux[1], "\n")   # + TRUE (N
print(sort(unique(ann$symbol)))          # "+" "A" "N"
\end{lstlisting}
For a record without a header, \code{fs} gives the frequency in which times in seconds
are converted (the library reports on the standard error that it cannot open the
header):
\begin{lstlisting}[language=R,style=example]
library(wfdbx)

write_annotations("beats", "qrs", c(100, 460, 820), symbol = "N")   # no header
ann <- read_annotations("beats", "qrs", sampto = "0:02", fs = 360)
cat(ann$sample, attr(ann, "fs"), "\n")   # 100 460 360
\end{lstlisting}

% ------------------------------------------------------------------------------------------
\subsection{Writing records}
\label{sec:r-writing}

\apientry{write\_record}{Write a header and a signal file}
\begin{lstlisting}[language=R,style=signature]
write_record(record, signals, fs, fmt = 16, gains, baselines = NULL,
             adczeros = 0, adcres = NULL, units = NULL, names = NULL,
             base_time = NULL, comments = NULL, physical = FALSE,
             samps_per_frame = NULL, dir = NULL,
             legacy = FALSE, metadata = NULL)
\end{lstlisting}
Writes signals as a new record: the header \file{RECORD.hea} and one signal file,
\file{RECORD.dat}, which holds all the signals (so they share one format). Existing files
of these names are overwritten. The sampling frequency, the gain and ADC resolution of
each signal, and the units of each calibrated signal must be given: there are no
defaults. The per-signal arguments (\code{fmt}, \code{gains}, \code{baselines},
\code{adczeros}, \code{adcres}, \code{units}, \code{names}) take one value for all the
signals, or a vector (or list) with one value per signal. All the arguments, the samples
included, are checked before any file is created; whatever happens, the output signals
are closed (\code{wfdbquit}) before the function returns.
\begin{params}
\item[record] Record name (a non-empty character string), without the \code{.hea}
  extension: a plain name (\code{"rec"}) or a relative path (\code{"data/rec"}), created
  in the output directory if one is set, otherwise in the current directory, missing
  directories included; or a full path, absolute or in the home directory
  (\code{"\textasciitilde/Documents/rec"}), used as given. With a path, the files are created in that
  directory, and the header names the record and its signal file without the path
  (\code{rec}, \code{rec.dat}), so that the record opens under any path that leads to it
  (in legacy mode, as in the original, the header gives the name with its path). A URL
  cannot be written.
\item[signals] The samples: a numeric matrix with one row per sample and one column per
  signal; a vector is one signal. In ADC units (the default), the values are integers
  within the range of a 32-bit integer, and \code{NA} is written as an invalid sample;
  with \code{physical = TRUE}, they are in physical units.
\item[fs] The sampling frequency, in Hz: a finite number $> 0$. Required.
\item[fmt] The storage format, such as \code{16}, \code{212}, \code{24} or \code{516}
  (FLAC); one value, or one per signal, all equal (\cref{chap:formats}). Default
  \code{16}.
\item[gains] The gain of each signal, in ADC units per physical unit: a finite number;
  \code{0} marks an uncalibrated signal, which is allowed only for samples given in ADC
  units. Negative gains are allowed. Required (no default).
\item[baselines] The ADC value that corresponds to 0 physical units. Default
  \code{NULL}: the ADC zero of the signal.
\item[adczeros] The ADC value at the middle of the ADC range. Default \code{0}.
\item[adcres] The ADC resolution, in bits (an integer $\ge 1$). Required (default
  \code{NULL}: not given).
\item[units] The physical units, such as \code{"mV"}, as a character vector or a list.
  Required for each signal whose gain is not 0. Default \code{NULL}: not given.
\item[names] The signal descriptions, such as \code{"MLII"}, as a character vector or a
  list. Default \code{NULL}: no descriptions.
\item[base\_time] The start time, with an optional date: \code{"hh:mm:ss"} or
  \code{"hh:mm:ss dd/mm/yyyy"}, with optional decimals of the seconds. Default
  \code{NULL}: no start time.
\item[comments] A character vector of info strings, written to the header as comment
  lines (\code{"\#"} followed by the string). Default \code{NULL}: none.
\item[physical] \code{TRUE} if \code{signals} are in physical units. They are converted
  to ADC units as $\mathrm{round}(x \cdot \mathit{gain}) + \mathit{baseline}$ (halves
  rounded away from zero); \code{NA} becomes an invalid sample, and infinite values and
  values beyond the range of 32-bit ADC units are errors. Default \code{FALSE}: ADC units.
\item[samps\_per\_frame] For a multi-frequency record, the number of samples per frame of
  each signal (integers $\ge 1$). \code{signals} then has one row per frame and
  \code{sum(samps\_per\_frame)} columns: the samples of signal 0, then those of signal 1,
  and so on. \code{fs} is then the frame frequency. Default \code{NULL}: one sample per
  frame for every signal.
\item[dir] The directory (which must exist) in which to create the files, for this call
  only; \code{record} must then be a plain name, without a directory, and the record
  opens from that directory under its plain name. Default \code{NULL} (or \code{""}): the
  output directory.
\item[legacy] \code{TRUE} to behave as the original WFDB package: the gain defaults to
  200, and the ADC resolution to one based on the format. Default \code{FALSE}.
\item[metadata] Metadata entries (\cref{chap:metadata}) to write to the header after the
  comments, as \code{\# Key: value} (through \code{putinfo}): a named list or a named
  character vector, such as \code{list(Age = "63")} or the result of
  \code{read\_metadata}. Keys and values follow the rules of \code{set\_metadata}; a
  \code{NULL} or \code{NA} value is an error, no two keys may differ only in letter case,
  and none may be the key of an entry in \code{comments}. Default \code{NULL}: none.
\end{params}
\returns{\code{NULL}, invisibly.}
\raises{Stops with an error if \code{record} is empty (``record name must not be
empty''); if \code{signals} has no signals (``at least one signal is required''); if
\code{fs}, a gain, an ADC resolution or the units of a calibrated signal is missing or
invalid (``fs must be finite and > 0'', ``gains must be specified'', ``gains must be
finite'', ``ADC resolution(s) must be specified'', ``units must be specified''); if
\code{base\_time} cannot be parsed (``invalid base\_time 'TEXT'\,''); if a vector of
per-signal values does not have one value per signal (``NAME must have one value per
signal''); if a signal in physical units has a gain of 0 (``signal K is uncalibrated'');
if a sample is not valid (``samples in ADC units must be integers'', ``samples in ADC
units must be within the range of a 32-bit integer''); if the signals do not all have the
same format; if a \code{samps\_per\_frame} value is not an integer $\ge 1$ (``samps\_per\_frame
must hold integers >= 1''), or the number of columns of \code{signals} is not their sum
(``signals has N columns, but samps\_per\_frame adds up to M''); if \code{metadata} is not
valid (an error of class \code{wfdbx\_argument\_error}, \cref{sec:r-warnings}); if
\code{record} contains a directory and \code{dir} is given, or \code{dir} does not exist
(``directory DIR does not exist''); and if the signal file or the header cannot be created,
or cannot be written completely, for example on a full disk (``write error on record
REC'').}
\begin{lstlisting}[language=R,style=example]
library(wfdbx)

# Two signals of random numbers, in millivolts, at 360 Hz: the record "rand"
# (rand.hea and rand.dat).
set.seed(1)
x <- matrix(rnorm(7200), ncol = 2)
write_record("rand", x, fs = 360, fmt = 16, physical = TRUE, gains = 1000, adcres = 16,
             units = "mV", names = c("A", "B"), base_time = "14:00:03 21/06/2024",
             comments = "random numbers", metadata = list(Generator = "rnorm"))
rec <- read_record("rand", physical = TRUE)
cat(dim(rec$signals), rec$base_time, "\n")      # 3600 2 [14:00:03 21/06/2024]
print(max(abs(rec$signals - x)) <= 0.0005)      # TRUE: rounded to 1/1000 mV
print(read_header("rand")$comments)             # "random numbers" " Generator: rnorm"

# An uncalibrated signal (gain 0, no units), in ADC units, in the directory "out".
dir.create("out", showWarnings = FALSE)
write_record("counts", 0:999, fs = 100, gains = 0, adcres = 12, dir = "out")
print(read_record("out/counts")$signals[1:3, 1])   # 0 1 2
\end{lstlisting}
A multi-frequency record: signal \code{fast} has 2 samples per frame, signal
\code{slow} one; each row of \code{signals} is one frame, with the two samples of
\code{fast} and then the sample of \code{slow}:
\begin{lstlisting}[language=R,style=example]
library(wfdbx)

frames <- 100                                   # 1 second at 100 frames per second
fast <- matrix(0:(2 * frames - 1), ncol = 2, byrow = TRUE)   # 200 samples per second
slow <- matrix(500, frames, 1)                  # 100 samples per second
write_record("mf", cbind(fast, slow), fs = 100, gains = 200, adcres = 12, units = "mV",
             names = c("fast", "slow"), samps_per_frame = c(2, 1))
rec <- read_record("mf", highres = TRUE)
cat(rec$fs, "\n")                               # 200
print(rec$signals[1:3, ])                       # rows 0 500, 1 500, 2 500
\end{lstlisting}

\apientry{write\_annotations}{Write an annotation file}
\begin{lstlisting}[language=R,style=signature]
write_annotations(record, annotator, sample, symbol = NULL, code = NULL,
                  subtype = 0, chan = 0, num = 0, aux = NULL, fs = NULL,
                  dir = NULL, legacy = FALSE, sort = NULL)
\end{lstlisting}
Writes an annotation file, \file{RECORD.ANNOTATOR}, in the MIT format; an existing file of
that name is overwritten. The record need not have a header. The annotation types are
given as \code{symbol} or as \code{code} (\cref{chap:formats}). Each per-annotation
argument is one value for all the annotations, or a vector with exactly one value per
annotation. If the annotations are not in time order, the file is sorted when it is
written, unless \code{sort = FALSE}. All the arguments are checked before the file is
created; whatever happens, the annotator is closed (\code{wfdbquit}) before the function
returns.
\begin{params}
\item[record] Record name (a character string), without the \code{.hea} extension (a
  trailing \code{".hea"} is also accepted): a plain name or a relative path
  (\code{"data/100s"}) is created in the output directory if one is set, otherwise in the
  current directory; a full path, absolute or in the home directory
  (\code{"\textasciitilde/Documents/100s"}), is used as given. The file is \file{RECORD.ANNOTATOR} at
  that place (\file{data/100s.qrs}). A URL cannot be written.
\item[annotator] Annotator name: the suffix of the file, made of letters, digits,
  \code{\_}, \code{-} and \code{\textasciitilde}.
\item[sample] The annotation times, in samples: integers within the range of a 64-bit
  integer, one per annotation. It may also be a data frame returned by
  \code{read\_annotations}, whose columns then give the other fields (\code{code},
  \code{subtype}, \code{chan}, \code{num} and \code{aux}) unless they are given.
\item[symbol] The mnemonics of the annotation types, such as \code{"N"} or \code{"V"}:
  one per annotation, or one string for all. An unknown mnemonic becomes code 0
  (\code{NOTQRS}). Default \code{NULL}: the types are given by \code{code}.
\item[code] The annotation codes (\code{NORMAL}, \code{PVC}, \ldots), instead of
  \code{symbol}: one per annotation, or one value for all. If both are given (as
  \code{read\_annotations} returns them), each symbol must be the mnemonic of its code (an
  empty symbol agrees with any code). Default \code{NULL}.
\item[subtype] The subtype of each annotation, stored as a signed byte (other values wrap
  around to $-128 \ldots 127$, as in C). Default \code{0}.
\item[chan] The channel number of each annotation, stored as an unsigned byte (other
  values wrap around to $0 \ldots 255$). Default \code{0}.
\item[num] The annotator number of each annotation, stored as a signed byte. Default
  \code{0}.
\item[aux] The auxiliary text of each annotation, with \code{NA} where there is none: a
  character vector, or one text for all. Texts longer than 255 bytes are cut. Default
  \code{NULL}: none.
\item[fs] If given (a finite number $> 0$), the frequency in which the times are
  counted, written to the file. Default \code{NULL}: none is written.
\item[dir] The directory (which must exist) in which to create the file, for this call
  only; \code{record} must then be a plain name, without a directory: the file is
  \file{DIR/RECORD.ANNOTATOR}. Default \code{NULL} (or \code{""}): the output directory.
\item[legacy] \code{TRUE} to behave as the original WFDB package (an out-of-order file is
  then sorted by running the \code{sortann} tool). Default \code{FALSE}.
\item[sort] \code{TRUE} to always sort the file, \code{FALSE} to keep the given order
  (without the library's advice to run \code{sortann}), for this call only. Default
  \code{NULL}: the environment variable \code{WFDBANNSORT} decides (sorting is on unless
  it is \code{0}).
\end{params}
\returns{The status of the last \code{putann} (0), invisibly.}
\raises{Stops with an error if neither \code{symbol} nor \code{code} is given (``symbol
or code must be given''); if both are given and disagree (``symbols and codes disagree at
annotation K'', K counting from 0); if a vector of per-annotation values does not have one
value per annotation (``NAME must have one value per annotation''); if a time is not an
integer within the range of a 64-bit integer; if a code, subtype, channel or annotator
number is \code{NA}, \code{NaN} or infinite (``NAME must not be NA, NaN or infinite''); if
\code{fs} is not finite and $> 0$; if \code{sort} is not \code{NULL}, \code{TRUE} or
\code{FALSE}; if \code{record} contains a directory and \code{dir} is given, or \code{dir}
does not exist; and if the annotation file cannot be created (``can't create annotator ANN
for record REC'') or cannot be written completely (``write error on annotator ANN for
record REC'').}
\begin{lstlisting}[language=R,style=example]
library(wfdbx)

# Three beats and a rhythm change, by symbol, with an aux string for the rhythm.
write_annotations("100s", "test", c(18, 77, 370, 662),
                  symbol = c("+", "N", "V", "N"), aux = c("(N", NA, NA, NA))
ann <- read_annotations("100s", "test")
print(ann$symbol)                        # "+" "N" "V" "N"
print(ann$aux)                           # "(N" NA NA NA

# A copy of the reference annotations: the data frame gives all the fields.
a <- read_annotations("100s", "atr")
write_annotations("100s", "copy", a)
b <- read_annotations("100s", "copy")
print(identical(b$sample, a$sample) && identical(b$symbol, a$symbol) &&
      identical(b$aux, a$aux))           # TRUE

# Annotations out of time order are sorted when the file is written (the library
# reports it on the standard error).
write_annotations("100s", "unsorted", c(500, 100), code = NORMAL)
print(read_annotations("100s", "unsorted")$sample)   # 100 500
\end{lstlisting}

% ------------------------------------------------------------------------------------------
\subsection{Checking records}
\label{sec:r-checking}

\apientry{check\_record}{Check a record and its annotation files}
\begin{lstlisting}[language=R,style=signature]
check_record(record, annotators = character(0), input_dir = NULL,
             network = NULL)
\end{lstlisting}
Checks the header of a record, its signal specifications, every sample of its signal
files, each segment of a multi-segment record, and the annotation files named, and
reports each problem found. A problem in the record never stops with an error: it is
reported, with a severity (\code{"error"}: the record is damaged or does not follow the
WFDB format, and programs will read wrong data or fail; \code{"warning"}: the record can
be read, but something is missing or unusual, and may give wrong results; \code{"info"}:
a fact worth knowing, such as invalid samples). The problems, their codes and their
messages are the same in the seven languages (\cref{tab:python-issue-codes}). The checks
run in normal mode, even if legacy mode is on, and the library prints no messages during
the call; both settings are restored afterwards. They are similar to those of Benjamin
Moody's \code{wfdbcheck}.
\begin{params}
\item[record] Record name (a character string), without the \code{.hea} extension (a
  trailing \code{".hea"} is also accepted): a full path, absolute
  (\code{"/data/mitdb/100"}) or in the home directory
  (\code{"\textasciitilde/Documents/100s"}), used as given; a relative path (\code{"100s"},
  \code{"data/100s"}), read from the input directory if one is set, otherwise from the
  current directory; or a URL (\code{"https://physionet.org/files/mitdb/1.0.0/100"}),
  read only if network access is on.
\item[annotators] A character vector of the annotator names whose files are checked, such
  as \code{"atr"}. Default \code{character(0)}: none.
\item[input\_dir] The input directory for this call only (optional): a local directory, a
  \code{\textasciitilde} path or a URL, under which relative record names are read. It does not
  apply to full paths or URLs. Default \code{NULL}: the current setting
  (\code{wfdbx\_set\_input\_dir}).
\item[network] \code{TRUE} to allow reading URLs for this call only, \code{FALSE} to
  forbid it. Default \code{NULL}: the current setting.
\end{params}
\returns{A list of class \code{"wfdbx\_check\_report"} (\cref{tab:r-checkreport}):
\code{report\$ok} is \code{TRUE} if no problem is an error, \code{report\$issues} is a data
frame with one row per problem, \code{format(report)} gives the text form, one line per
problem, and \code{print(report)} prints it (\code{no problems found} if there is none).}
\raises{Stops with an error if \code{record} is not a single string (``record must be a
string'') or is empty (``record must not be empty''), or if \code{annotators} is not a
character vector without \code{NA} (``annotator names must be strings''). A record that
cannot be read is reported (code \code{header\_unreadable}), not raised.}
\begin{lstlisting}[language=R,style=example]
library(wfdbx)

report <- check_record("100s", "atr")
print(report$ok)          # TRUE: no errors
report
# warning: signal 0: the units are not specified
# warning: signal 1: the units are not specified

# A damaged record: one sample of a 12-bit signal is overwritten with 2048.
write_record("damaged", 0:99, fs = 100, gains = 100, adcres = 12, units = "mV",
             names = "ramp")
con <- file("damaged.dat", "r+b")
invisible(seek(con, 20, rw = "write"))  # sample 10 (2 bytes per sample in format 16)
writeBin(as.raw(c(0x00, 0x08)), con)    # 0x0800 = 2048, little-endian
close(con)
report <- check_record("damaged")
print(report$ok)          # FALSE
print(report$issues[, c("severity", "code", "where")])
#   severity         code    where
# 1    error sample_range signal 0
# 2    error     checksum signal 0
print(format(report)[1])
# "error: signal 0: 1 sample is outside the ADC range (-2048 to 2047): sample 10 (value 2048)"
\end{lstlisting}

% ------------------------------------------------------------------------------------------
\subsection{Header metadata}
\label{sec:r-metadata}

These functions read and change the metadata of a header: its info strings of the form
\code{\# Key: value}, such as \code{\# Age: 63} (\cref{chap:metadata}). In short: keys are
compared without regard to letter case (A-Z and a-z only); a key must not be empty and
must not contain \code{:}, \code{\#}, a line break or a NUL byte, and a value must not
contain a line break or a NUL byte; spaces at either end of keys and values are removed.
Keys and values are handled as bytes: a character string is taken as its bytes, never
re-encoded, and a raw vector is also accepted, as given. A change rewrites the header in
one step (a temporary file in the same directory, renamed over the header, keeping its
permission bits), keeping every other byte of it; all the arguments are checked first, so
an error leaves the header unchanged, and nothing is written if nothing changes (unless
\code{dir} is given). Only the \file{.hea} file is read (not a \file{.info} file); for a
multi-segment record, it is the master header, unless a segment is named. The output
directory is not used: a change is written to the header that was read, or to \code{dir}.

The entries are a named character vector: the names are the keys and the elements the
values, in file order, with the spelling of the file, such as \code{c(Age = "63", Sex =
"F")}. Since reading returns each key once, \code{md[["Age"]]} gives a value (with the
file's spelling of the key; \code{get\_metadata} ignores letter case). The same form, or
a named list, is what \code{update\_metadata} and the argument \code{metadata} of
\code{write\_record} take, so the entries of one header can be written to another.

Metadata often holds information about patients: WFDB-x stores it as plain text, exactly
as given, and does not de-identify it.

\begin{phibox}
A header can hold identifying data about a patient (a name, a date of birth, a medical
record number). The metadata functions read and write such data as given. Records that
are shared must be de-identified first.
\end{phibox}

\apientry{read\_metadata}{Read the metadata of a record}
\begin{lstlisting}[language=R,style=signature]
read_metadata(record, input_dir = NULL, network = NULL)
\end{lstlisting}
Reads the metadata entries of the header of a record, in file order. An info string is an
entry if it contains a colon: the key is the text before the first colon, the value the
rest of the line, both without spaces at either end. Info strings without a colon are free
comments, which are not returned.
\begin{params}
\item[record] Record name (a character string), without the \code{.hea} extension (a
  trailing \code{".hea"} is also accepted): a full path, absolute
  (\code{"/data/mitdb/100"}) or in the home directory
  (\code{"\textasciitilde/Documents/100s"}), used as given; a relative path (\code{"100s"},
  \code{"data/100s"}), read from the input directory if one is set, otherwise from the
  current directory; or a URL (\code{"https://physionet.org/files/mitdb/1.0.0/100"}),
  read only if network access is on.
\item[input\_dir] The input directory for this call only (optional): a local directory, a
  \code{\textasciitilde} path or a URL, under which relative record names are read. It does not
  apply to full paths or URLs. Default \code{NULL}: the current setting
  (\code{wfdbx\_set\_input\_dir}).
\item[network] \code{TRUE} to allow reading URLs for this call only, \code{FALSE} to
  forbid it. Default \code{NULL}: the current setting.
\end{params}
\returns{A named character vector of the entries, in file order (an empty named character
vector if there are none). A key that appears more than once (ignoring letter case) is
returned once, with its first value, and an R warning is issued (\cref{sec:r-warnings}).}
\raises{An error of class \code{wfdbx\_file\_error} if the header cannot be found or read
(``can't read header for record REC''); of class \code{wfdbx\_argument\_error} if
\code{record} is empty or \code{"-"} (the standard input).}
\begin{lstlisting}[language=R,style=example]
library(wfdbx)

write_record("pat", rep(0, 360), fs = 360, gains = 200, adcres = 12, units = "mV",
             names = "ECG", comments = "recorded at rest",
             metadata = list(Age = "63", Sex = "F"))
print(read_metadata("pat"))             # c(Age = "63", Sex = "F")
print(read_header("pat")$comments)
# "recorded at rest" " Age: 63" " Sex: F": the free comment is not an entry
\end{lstlisting}

\apientry{get\_metadata}{Return the value of a metadata entry}
\begin{lstlisting}[language=R,style=signature]
get_metadata(record, key, input_dir = NULL, network = NULL)
\end{lstlisting}
Returns the value of one metadata entry of the header of a record.
\begin{params}
\item[record] Record name (as for \code{read\_metadata}).
\item[key] The key, such as \code{"Age"}: a character string or a raw vector. Letter case
  is ignored, and spaces at either end are removed.
\item[input\_dir] The input directory for this call only. Default \code{NULL}: the
  current setting.
\item[network] \code{TRUE} to allow reading URLs for this call only, \code{FALSE} to
  forbid it. Default \code{NULL}: the current setting.
\end{params}
\returns{The value, a character string (\code{""} for an entry with an empty value), or
\code{NULL} if the header has no entry for \code{key}. If it has more than one, the first
is returned, and an R warning is issued.}
\raises{An error of class \code{wfdbx\_file\_error} if the header cannot be found or
read; of class \code{wfdbx\_argument\_error} if \code{record} is empty or \code{"-"}, or
if \code{key} is not a valid key or not a string.}
\begin{lstlisting}[language=R,style=example]
library(wfdbx)

write_record("pat", rep(0, 360), fs = 360, gains = 200, adcres = 12, units = "mV",
             metadata = c(Age = "63", Sex = "F"))
print(get_metadata("pat", "age"))       # "63"
print(get_metadata("pat", "Height"))    # NULL
\end{lstlisting}

\apientry{set\_metadata}{Set the value of a metadata entry}
\begin{lstlisting}[language=R,style=signature]
set_metadata(record, key, value, input_dir = NULL, network = NULL, dir = NULL)
\end{lstlisting}
Sets the value of a metadata entry in the header of a record. If the header has an entry
for \code{key} (ignoring letter case), its line is replaced by \code{\# Key: value}, at
the same position, with the spelling of \code{key} given here; otherwise the entry is
added after the last info string of the header (or, if it has none, after its last line
that is not blank). The entry is written as \code{\# Key: value}, or \code{\# Key:} if
the value is empty. Every other line is kept, byte for byte, including its line ending.
\begin{params}
\item[record] Record name (as for \code{read\_metadata}). The header is changed where it
  was read, unless \code{dir} is given; a header read from a URL can only be written
  elsewhere, with \code{dir}.
\item[key] The key, such as \code{"Age"} or \code{"Date of birth"}: a character string
  or a raw vector. Spaces at either end are removed. It must not be empty, and must not
  contain \code{:}, \code{\#}, a line break or a NUL byte.
\item[value] The value: a character string, which may be a sentence, or a raw vector.
  Spaces at either end are removed. It must not contain a line break or a NUL byte; it
  may be empty.
\item[input\_dir] The input directory in which the header is read, for this call only.
  Default \code{NULL}: the current setting.
\item[network] \code{TRUE} to allow reading the header from a URL, for this call only,
  \code{FALSE} to forbid it. Default \code{NULL}: the current setting.
\item[dir] A directory, which must exist, in which to write the changed header, under the
  same file name, instead of changing the header itself; the signal files are not copied.
  The header is then written even if nothing changes. Default \code{NULL} (or \code{""}):
  the header itself is changed.
\end{params}
\returns{\code{NULL}, invisibly.}
\raises{An error of class \code{wfdbx\_file\_error} if the header cannot be read, if the
changed header cannot be written (``can't write header for record REC''), or if
\code{dir} does not exist (``directory DIR does not exist''); of class
\code{wfdbx\_argument\_error} if \code{key} or \code{value} is not valid or not a string,
if the header has more than one entry for \code{key} (``header of record REC has more than
one entry for key KEY (lines L1, L2); remove them first''), if it has no record line, if
it was read from a URL and \code{dir} is not given, or if \code{record} is empty or
\code{"-"}.}
\begin{lstlisting}[language=R,style=example]
library(wfdbx)

write_record("pat", rep(0, 360), fs = 360, gains = 200, adcres = 12, units = "mV",
             metadata = c(Age = "63", Sex = "F"))
set_metadata("pat", "AGE", "64")                                # replaces "# Age: 63"
set_metadata("pat", "Diagnosis", "sinus rhythm, no ectopy")     # a new entry
print(read_metadata("pat"))
# c(AGE = "64", Sex = "F", Diagnosis = "sinus rhythm, no ectopy")

# Write the changed header to another directory; pat.hea is left unchanged.
dir.create("changed", showWarnings = FALSE)
set_metadata("pat", "Sex", "M", dir = "changed")
cat(get_metadata("pat", "Sex"), get_metadata("changed/pat", "Sex"), "\n")   # F M
\end{lstlisting}

\apientry{remove\_metadata}{Remove a metadata entry}
\begin{lstlisting}[language=R,style=signature]
remove_metadata(record, key, input_dir = NULL, network = NULL, dir = NULL)
\end{lstlisting}
Removes every entry for \code{key} (ignoring letter case) from the header of a record, so
it also repairs a header that has more than one. Other lines are kept, byte for byte.
Nothing is written if the header has no entry for \code{key}, unless \code{dir} is given.
\begin{params}
\item[record] Record name (as for \code{set\_metadata}).
\item[key] The key (as for \code{set\_metadata}).
\item[input\_dir] The input directory in which the header is read, for this call only.
  Default \code{NULL}: the current setting.
\item[network] \code{TRUE} to allow reading the header from a URL, for this call only,
  \code{FALSE} to forbid it. Default \code{NULL}: the current setting.
\item[dir] A directory, which must exist, in which to write the changed header instead of
  changing the header itself. Default \code{NULL} (or \code{""}): the header itself is
  changed.
\end{params}
\returns{The number of lines removed (a number; 0 if the header has no entry for
\code{key}).}
\raises{As \code{set\_metadata}, except that more than one entry for \code{key} is not an
error.}
\begin{lstlisting}[language=R,style=example]
library(wfdbx)

write_record("pat", rep(0, 360), fs = 360, gains = 200, adcres = 12, units = "mV",
             metadata = c(Age = "63", Sex = "F"))
print(remove_metadata("pat", "sex"))    # 1
print(remove_metadata("pat", "Sex"))    # 0: there is no such entry now
print(read_metadata("pat"))             # c(Age = "63")
\end{lstlisting}

\apientry{update\_metadata}{Make several changes to the metadata}
\begin{lstlisting}[language=R,style=signature]
update_metadata(record, changes, input_dir = NULL, network = NULL, dir = NULL)
\end{lstlisting}
Makes several changes to the metadata of a record, in one write of the header. Each
change is applied in order, as \code{set\_metadata} does, or, for a value of \code{NULL}
or \code{NA}, as \code{remove\_metadata} does. All the changes are checked first: if one
is not valid, none is made.
\begin{params}
\item[record] Record name (as for \code{set\_metadata}).
\item[changes] A named list or a named character vector: the names are the keys, and each
  element a value (a character string or a raw vector), or \code{NULL} or \code{NA} to
  remove the entry, as in \code{list(Age = "64", Dx = NULL)} or \code{c(Age = "64", Dx =
  NA)}. No two keys may differ only in letter case. An empty list or vector makes no
  change.
\item[input\_dir] The input directory in which the header is read, for this call only.
  Default \code{NULL}: the current setting.
\item[network] \code{TRUE} to allow reading the header from a URL, for this call only,
  \code{FALSE} to forbid it. Default \code{NULL}: the current setting.
\item[dir] A directory, which must exist, in which to write the changed header instead of
  changing the header itself. Default \code{NULL} (or \code{""}): the header itself is
  changed.
\end{params}
\returns{\code{NULL}, invisibly.}
\raises{As \code{set\_metadata}; also an error of class \code{wfdbx\_argument\_error} if
two keys of \code{changes} differ only in letter case, or if \code{changes} is not a
named list or a named character vector (\code{NULL} included).}
\begin{lstlisting}[language=R,style=example]
library(wfdbx)

write_record("pat", rep(0, 360), fs = 360, gains = 200, adcres = 12, units = "mV",
             metadata = c(Age = "63", Sex = "F"))
update_metadata("pat", list(Age = "64", Sex = NULL, Dx = "164889003"))
print(read_metadata("pat"))             # c(Age = "64", Dx = "164889003")
tryCatch(update_metadata("pat", c(Dx = "426783006", dx = NA)),
         wfdbx_argument_error = function(e) print(conditionMessage(e)))
# "metadata key 'dx' is given more than once (keys are compared without regard to letter case)"
print(get_metadata("pat", "Dx"))        # "164889003": nothing was changed
\end{lstlisting}

% ------------------------------------------------------------------------------------------
\subsection{The results}
\label{sec:r-results}

The readers return plain R objects: \code{read\_record} a list
(\cref{tab:r-record}), \code{read\_header} a list of class \code{"wfdbx\_header"}
(\cref{tab:r-header}), \code{read\_annotations} a data frame (\cref{tab:r-annotations}),
\code{read\_metadata} a named character vector (\cref{sec:r-metadata}), and
\code{check\_record} a list of class \code{"wfdbx\_check\_report"}
(\cref{tab:r-checkreport}). The signals are described by \code{siginfo} lists
(\cref{tab:r-siginfo}), which are those of the WFDB library functions
(\cref{sec:r-library}). The elements of these lists are listed in the order in which they
appear.

\begin{xltabular}{\linewidth}{@{}lL@{}}
\caption{The elements of the list returned by \code{read\_record}}
\label{tab:r-record} \\
\toprule
\tabhead Element & \tabhead Contents \\
\midrule
\endfirsthead
\toprule
\tabhead Element & \tabhead Contents \\
\midrule
\endhead
\bottomrule
\endfoot
\code{name} & the record name, as given to \code{read\_record} \\
\code{fs} & the sampling frequency in Hz, or 0 if the header gives none; with
  \code{highres = TRUE}, the highest sampling frequency of a multi-frequency record \\
\code{signals} & a numeric matrix with one row per sample and one column per signal read:
  ADC units, or physical units with \code{physical = TRUE}; invalid samples are
  \code{WFDB\_INVALID\_SAMPLE}, or \code{NA} in physical units \\
\code{siginfo} & a list of \code{siginfo} lists, one per signal read, in the order of the
  columns; \code{desc} has no white space at either end, and \code{adcres} is 0 if the
  header does not give it \\
\code{names} & the description of each signal (a character vector), such as
  \code{"MLII"}, without the white space that a header may put at either end of it \\
\code{units} & the physical units of each signal, as a list, with \code{NULL} where the
  header gives none \\
\code{gains} & the gain of each signal, in ADC units per physical unit; 0 means that the
  signal is uncalibrated \\
\code{baselines} & the ADC value of each signal that corresponds to 0 physical units \\
\code{base\_time} & the start time and date as text, to the second, such as
  \code{"[14:00:03 21/06/2024]"} (\code{"[14:00:03]"} without a date), or \code{NULL} if
  the header gives none \\
\code{base\_time\_ms} & the start time in milliseconds after midnight, such as
  \code{50403500} for 14:00:03.500, or \code{NULL} if \code{base\_time} is \code{NULL} \\
\code{base\_date} & the start date as an integer vector \code{c(year, month, day)}, or
  \code{NULL} if the header gives none \\
\code{sampfrom} & the number of the first sample returned \\
\code{comments} & a character vector of the info strings of the header (without the
  \code{\#}) \\
\code{checksum\_ok} & \code{FALSE} if the samples do not match the checksums of the
  header; checksums are checked only when a read reaches the end of the record, so a
  shorter read gives \code{TRUE} \\
\end{xltabular}

\begin{xltabular}{\linewidth}{@{}lL@{}}
\caption{The elements of the list of class \code{"wfdbx\_header"} returned by
  \code{read\_header}}
\label{tab:r-header} \\
\toprule
\tabhead Element & \tabhead Contents \\
\midrule
\endfirsthead
\toprule
\tabhead Element & \tabhead Contents \\
\midrule
\endhead
\bottomrule
\endfoot
\code{name} & the record name, as given to \code{read\_header} \\
\code{fs} & the sampling frequency in Hz, or 0 if the header gives none \\
\code{counter\_freq} & the counter frequency, as \code{getcfreq()} returns it (the
  sampling frequency if the header gives none) \\
\code{base\_counter} & the base counter value, as \code{getbasecount()} returns it \\
\code{nsamples} & the number of samples per signal, or 0 if the header does not give it \\
\code{nsig} & the number of signals \\
\code{siginfo} & a list of \code{siginfo} lists, one per signal; \code{desc} has no white
  space at either end, and \code{adcres} is 0 if the header does not give it \\
\code{base\_time}, \code{base\_time\_ms}, \code{base\_date} & the start time and date, as
  in the result of \code{read\_record} \\
\code{comments} & a character vector of the info strings of the header (without the
  \code{\#}) \\
\code{metadata} & the metadata entries of the header, as \code{read\_metadata} returns
  them: a named character vector, in file order (empty for an EDF file and for the
  standard input) \\
\end{xltabular}

\begin{table}[htbp]
\centering
\caption{The columns of the data frame returned by \code{read\_annotations}, one row per
  annotation, in file order}
\label{tab:r-annotations}
\begin{tabularx}{\linewidth}{@{}lL@{}}
\toprule
\tabhead Column or attribute & \tabhead Contents \\
\midrule
\code{sample} & the annotation times, in samples \\
\code{symbol} & the mnemonics of the codes, such as \code{"N"}, \code{"V"} or \code{"+"}
  (as \code{annstr} gives them) \\
\code{code} & the annotation codes, such as \code{NORMAL} or \code{PVC} \\
\code{subtype} & the subtype of each annotation \\
\code{chan} & the channel number of each annotation \\
\code{num} & the annotator number of each annotation \\
\code{aux} & the auxiliary texts of the annotations, with \code{NA} where there is
  none \\
\code{attr(, "record")} & the record name, as given \\
\code{attr(, "annotator")} & the annotator name, such as \code{"atr"} \\
\code{attr(, "fs")} & the frequency in which \code{sample} is counted, normally the
  sampling frequency of the record; 0 if neither the record, nor the file, nor the
  argument \code{fs} of \code{read\_annotations} gives one \\
\bottomrule
\end{tabularx}
\end{table}

\begin{table}[htbp]
\centering
\caption{The list of class \code{"wfdbx\_check\_report"} returned by \code{check\_record},
  and its methods}
\label{tab:r-checkreport}
\begin{tabularx}{\linewidth}{@{}lL@{}}
\toprule
\tabhead Element or method & \tabhead Contents \\
\midrule
\code{record} & the record name, as given \\
\code{ok} & \code{TRUE} if no issue has the severity \code{"error"} \\
\code{issues} & a data frame with one row per problem, in the order of the checks, and
  the columns \code{severity} (\code{"error"}, \code{"warning"} or \code{"info"}),
  \code{code} (a stable identifier of the kind of problem, such as \code{"checksum"} or
  \code{"gain\_missing"}, \cref{tab:python-issue-codes}), \code{where} (\code{""} for the
  record, \code{"signal N"}, \code{"segment SEG"}, \code{"segment SEG, signal N"} or
  \code{"annotator ANN"}) and \code{message} \\
\code{format(report)} & the text form of the report, as a character vector: one line per
  issue, \code{"SEVERITY: WHERE: MESSAGE"} (\code{"SEVERITY: MESSAGE"} when \code{where}
  is empty), or the one line \code{no problems found} \\
\code{print(report)} & prints the lines of \code{format(report)}, and returns the report
  invisibly \\
\bottomrule
\end{tabularx}
\end{table}

\begin{xltabular}{\linewidth}{@{}lL@{}}
\caption{The elements of a \code{siginfo} list, the specification of a signal}
\label{tab:r-siginfo} \\
\toprule
\tabhead Element & \tabhead Contents \\
\midrule
\endfirsthead
\toprule
\tabhead Element & \tabhead Contents \\
\midrule
\endhead
\bottomrule
\endfoot
\code{fname} & the name of the signal file, as the header gives it \\
\code{desc} & the signal description, such as \code{"MLII"}, or \code{NULL} \\
\code{units} & the physical units, such as \code{"mV"}, or \code{NULL} if not specified \\
\code{gain} & the gain, in ADC units per physical unit (0: uncalibrated) \\
\code{initval} & the value of the first sample \\
\code{group} & the signal group number (the signals of one group share one file) \\
\code{fmt} & the storage format, such as \code{16}, \code{212} or \code{516} \\
\code{spf} & the number of samples per frame (more than 1 for the faster signals of a
  multi-frequency record) \\
\code{bsize} & the block size, in bytes (0 for ordinary files) \\
\code{adcres} & the ADC resolution, in bits (0: not specified) \\
\code{adczero} & the ADC output for an input of 0\,V \\
\code{baseline} & the ADC value that corresponds to 0 physical units \\
\code{nsamp} & the number of samples (0 if unknown) \\
\code{cksum} & the 16-bit checksum of the samples \\
\end{xltabular}

% ------------------------------------------------------------------------------------------
\subsection{Warnings and errors}
\label{sec:r-warnings}

\paragraph{Errors.} The high-level functions report errors as R errors (conditions of
class \code{"error"}, raised by \code{stop}), whose messages are given in the entries of
this section; \code{tryCatch} handles them. Each function checks its arguments before it
creates or changes a file, so an invalid argument leaves the files as they were. The
library may also print a message on the standard error when it cannot open a file (as in
\code{init: can't open header for record REC}); \code{wfdbquiet()} turns these messages
off (\cref{sec:r-library}).

The errors of the metadata functions, and those of the argument \code{metadata} of
\code{write\_record}, also have a class that tells their kind (\cref{tab:r-errors}), so
that a program can handle each kind on its own, with
\code{tryCatch(..., wfdbx\_file\_error = function(e) ...)}.

\begin{table}[htbp]
\centering
\caption{The classes of the errors of the metadata functions}
\label{tab:r-errors}
\begin{tabularx}{\linewidth}{@{}lL@{}}
\toprule
\tabhead Class & \tabhead Raised when \\
\midrule
\code{wfdbx\_file\_error} & a header cannot be found, read or written, or the directory
  given as \code{dir} does not exist (the errors that Python reports as \code{IOError}) \\
\code{wfdbx\_argument\_error} & an argument is invalid (a record name that is empty or
  \code{"-"}, a key or value that is not valid or not a string, changes or metadata that
  are not a named list or named character vector, two keys that differ only in letter
  case), or the header does not allow the change (more than one entry for the key, no
  record line, a header read from a URL without \code{dir}); Python's \code{ValueError}
  and \code{TypeError} \\
\bottomrule
\end{tabularx}
\end{table}
Both classes are also of class \code{"error"} (and \code{"condition"}), so
\code{tryCatch(..., error = )} catches them too.

\paragraph{The checksum warning.} When \code{read\_record} reads to the end of a record
whose samples do not match the checksums of its header, it returns all the samples and
issues an R warning: ``checksum error in record REC: the signal data do not match the
checksums in its header''. \code{checksum\_warning = FALSE} turns the warning off;
\code{checksum\_ok} of the result is \code{FALSE} either way.

\paragraph{The duplicate-key warning.} When a header holds two metadata entries whose keys
differ only in letter case (written by other software), \code{read\_metadata},
\code{get\_metadata} and \code{read\_header} use the first and issue an R warning:
``header of record REC has more than one entry for key KEY (lines L1, L2); the first is
used''. \code{set\_metadata} and \code{update\_metadata} then refuse to change that key
(an error of class \code{wfdbx\_argument\_error}), and \code{remove\_metadata} removes all
of its entries.

Both warnings are issued with \code{warning()}, without the call. As any other R warning,
they can be handled (\code{withCallingHandlers}), silenced (\code{suppressWarnings}) or
turned into errors (\code{options(warn = 2)}):

\begin{lstlisting}[language=R,style=example]
library(wfdbx)

# A record whose last sample is changed after it was written: the checksum fails.
write_record("ck", 0:99, fs = 100, gains = 100, adcres = 12, units = "mV")
con <- file("ck.dat", "r+b")
invisible(seek(con, 198, rw = "write"))
writeBin(as.raw(c(0, 0)), con)
close(con)
w <- NULL
rec <- withCallingHandlers(read_record("ck"), warning = function(c) {
  w <<- conditionMessage(c)
  invokeRestart("muffleWarning")
})
cat(dim(rec$signals), rec$checksum_ok, "\n")    # 100 1 FALSE: all the samples
print(w)
# "checksum error in record ck: the signal data do not match the checksums in its header"
rec <- read_record("ck", checksum_warning = FALSE)      # no warning

# Two entries for one key: the first is used, with a warning.
write_record("dup", rep(0, 10), fs = 100, gains = 0, adcres = 12,
             comments = c(" Age: 63", " AGE: 64"))
print(suppressWarnings(get_metadata("dup", "age")))     # "63"
tryCatch(set_metadata("dup", "Age", "65"),
         wfdbx_argument_error = function(e) print(conditionMessage(e)))
# "header of record dup has more than one entry for key Age (lines 3, 4); remove them first"
print(remove_metadata("dup", "age"))                    # 2
\end{lstlisting}

\begin{lstlisting}[language=R,style=example]
library(wfdbx)

tryCatch(read_record("nosuchrecord"), error = function(e) print(conditionMessage(e)))
# "can't read record nosuchrecord" (after the library's message on the standard error)

write_record("raw", 0:9, fs = 100, gains = 0, adcres = 12)      # uncalibrated
tryCatch(read_record("raw", physical = TRUE), error = function(e) print(conditionMessage(e)))
# "signal 0 of record raw is uncalibrated"

# The metadata functions: the class tells the kind of error.
e <- tryCatch(read_metadata("nosuchrecord"), error = function(e) e)
print(class(e))           # "wfdbx_file_error" "error" "condition"
e <- tryCatch(set_metadata("raw", "a:b", "1"), error = function(e) e)
print(class(e))           # "wfdbx_argument_error" "error" "condition"
print(conditionMessage(e))   # "metadata key 'a:b' must not contain ':'"
\end{lstlisting}

% ==========================================================================================
\section{The WFDB library functions in R}
\label{sec:r-library}

Every function of the WFDB library is available under its C name, with the same
arguments, in the same order, and the same meaning, so that programs and documentation
written for the C library carry over directly. The functions, their arguments and their
results are documented once, for all the languages, in \cref{chap:library}; this section
gives the rules by which the C declarations become R functions. The help pages
\code{?isigopen}, \code{?timstr}, \code{?annopen}, \code{?annstr}, \code{?calopen},
\code{?getwfdb}, \code{?ecgcodes}, \code{?isann} and \code{?siginfo} list the functions of
each part of the library, with the result of each.

\subsection{Conventions}

\paragraph{Names.} The functions, the constants (\code{WFDB\_READ},
\code{WFDB\_INVALID\_SAMPLE}, \code{WFDB\_HIGHRES}, \ldots) and the annotation codes
(\code{NORMAL}, \code{PVC}, \ldots, \cref{chap:formats}; \code{ecgcodes()} returns them as
a list) keep their C names. The macros of \file{<wfdb/ecgmap.h>} are functions, under both
their macro names and the names of the functions behind them: \code{isqrs} and
\code{wfdb\_isqrs}, \code{map1} and \code{wfdb\_map1}, and so on; they take vectors. The C
function \code{sample} is \code{wfdb\_sample}, so that it does not mask R's
\code{sample}.

\paragraph{Numbers} are doubles. Signal, group and annotator numbers start at 0, as in C
(signal \code{s} is element \code{s + 1} of an R vector). C's integer arithmetic (wrap
around, casts to \code{short} or \code{char}) is applied where C applies it.

\paragraph{NULL is \code{NULL}.} A \code{NULL} pointer argument or result is
\code{NULL}: \code{sampfreq(NULL)} returns the frequency of the open record,
\code{getinfo(NULL)} continues with the info strings of the last record, and
\code{getinfo} returns \code{NULL} after the last one. A list element that is
\code{NULL} is set with \code{x["f"] <- list(NULL)} (\code{x\$f <- NULL} removes it).

\paragraph{Structures} are lists with the C field names, made by the constructors
\code{siginfo()}, \code{anninfo()}, \code{annotation()}, \code{calinfo()} and
\code{seginfo()} (\cref{sec:r-structures}). An array of structures is a list of them:
\code{annopen(record, list(anninfo("atr", WFDB\_READ)))}, \code{osigfopen(list(siginfo(...),
...))}.

\paragraph{Arrays and output arguments.} R functions do not change their arguments: where
a C function fills an array or a structure that the caller passes, the R function returns
it instead, in a list with the C return value:
\begin{lstlisting}[language=R]
r <- isigopen(record, nsig)      # r$nsig: the C result; r$si: the siginfo lists
r <- osigopen(record, nsig)      # the same
g <- getvec()                    # g$stat: the C result; g$v: one sample per signal
g <- getframe()                  # g$stat; g$v: one sample per sample of a frame
g <- getann(an)                  # g$stat; g$annot: the annotation (NULL on errors)
g <- getcal(desc, units)         # g$stat; g$cal: the calinfo list (NULL if none)
r <- wfdbinit(record, aiarray, nsig)   # r$stat; r$si
s <- getseginfo()                # a list of seginfo lists (empty: not multi-segment)
\end{lstlisting}
With \code{nsig} $\le 0$, \code{isigopen} reads the header only: \code{r\$nsig} is then
the number of signals of the record, and \code{r\$si} the specifications of up to
\code{-nsig} of them. The arguments that C takes only to receive results are left out
(\code{getvec()} has none), and the lengths of arrays have defaults: \code{osigfopen(siarray,
nsig = length(siarray))}, \code{setheader(record, siarray, nsig = length(siarray))},
\code{setmsheader(record, segment\_name, nsegments = length(segment\_name))},
\code{annopen(record, aiarray, nann = NULL)}. \code{putvec(vector)}, \code{putann(an,
annot)}, \code{ungetann(an, annot)} and \code{putcal(cal)} take the values to write.

\paragraph{Strings} are character strings, handled as bytes (ISO-8859-1, never
re-encoded), so every byte sequence round-trips unchanged; functions that take a string
also take a raw vector. The \code{aux} element of an annotation is \code{NULL}, or a raw
vector whose first byte is the length of the text that follows, as in C;
\code{aux\_string(aux)} returns the text (up to its first NUL byte), or \code{NULL}.

\paragraph{Return values and errors.} The functions return what the C functions return:
a status (negative on errors), a number of signals, a frequency, a string or
\code{NULL}. Errors are reported by these values and by a message on the standard error,
as in C, never by R errors. \code{wfdberror()} returns the last message;
\code{wfdbquiet()} turns the messages off and \code{wfdbverbose()} on again;
\code{wfdb\_error(fmt, ...)} reports a message of a program's own in the same way.

\paragraph{State.} As in C, the library keeps its state in the R session: the open record,
its input signals, the output signals and the annotators remain open between calls,
until \code{wfdbquit()} closes them. \code{wfdbquit()} leaves the input and output
directories and the network setting as they are. The high-level functions of
\cref{sec:r-api} call \code{wfdbquit()}, which closes anything opened with the library
functions.

\paragraph{Additions.} \code{getframes(n, minstat = 0)} reads up to \code{n} sample vectors
at once (\code{n} may be \code{Inf}), as \code{n} calls of \code{getvec} would, but
decoding whole blocks of the signal files. It returns \code{list(v, stat)}: a matrix with
one row per vector read and one column per signal, and the \code{getvec} status of the
first vector that could not be read, or the number of signals if all \code{n} were read;
with \code{minstat}, it stops at the first vector for which \code{getvec} returns less
than \code{minstat}. \code{putframes(m)} writes the rows of a matrix, as \code{putvec} on
each row. \code{sample\_valid()} returns the validity flag of the last
\code{wfdb\_sample}.

\paragraph{Speed.} Reading sample by sample with \code{getvec} in an R loop is much slower
than reading in blocks: \code{read\_record} and \code{getframes} decode whole blocks with
vectorized code, and are the way to read large records (in chunks, with \code{sampfrom}
and \code{sampto}, for a record too large for memory).

\subsection{The structures}
\label{sec:r-structures}

The structures of the C library are lists whose elements have the names of the fields of
the C structures (\cref{tab:r-structures}). Their constructors take the elements as
arguments, in this order, all with defaults:
\begin{lstlisting}[language=R]
siginfo(fname = NULL, desc = NULL, units = NULL, gain = 0,
        initval = 0, group = 0, fmt = 0, spf = 1, bsize = 0,
        adcres = 0, adczero = 0, baseline = 0, nsamp = 0,
        cksum = 0)
anninfo(name = NULL, stat = WFDB_READ)
annotation(time = 0, anntyp = 0, subtyp = 0, chan = 0, num = 0,
           aux = NULL)
calinfo(low = 0, high = 0, scale = 0, sigtype = NULL,
        units = NULL, caltype = 0)
seginfo(recname = "", nsamp = 0, samp0 = 0)
\end{lstlisting}
For a list of \code{n} structures, \code{replicate(n, siginfo(), simplify = FALSE)}.

\begin{xltabular}{\linewidth}{@{}lL@{}}
\caption{The lists of the C structures}
\label{tab:r-structures} \\
\toprule
\tabhead Constructor (C structure) & \tabhead Elements \\
\midrule
\endfirsthead
\toprule
\tabhead Constructor (C structure) & \tabhead Elements \\
\midrule
\endhead
\bottomrule
\endfoot
\code{siginfo} (\code{WFDB\_Siginfo}) & \code{fname}, \code{desc}, \code{units},
  \code{gain}, \code{initval}, \code{group}, \code{fmt}, \code{spf}, \code{bsize},
  \code{adcres}, \code{adczero}, \code{baseline}, \code{nsamp}, \code{cksum}
  (\cref{tab:r-siginfo}) \\
\code{seginfo} (\code{WFDB\_Seginfo}) & \code{recname} (the name of the segment),
  \code{nsamp} (its number of samples), \code{samp0} (the number of its first sample in
  the record) \\
\code{anninfo} (\code{WFDB\_Anninfo}) & \code{name} (the annotator name, the suffix of the
  file), \code{stat} (\code{WFDB\_READ} or \code{WFDB\_WRITE} for the MIT format,
  \code{WFDB\_AHA\_READ} or \code{WFDB\_AHA\_WRITE} for the AHA format) \\
\code{annotation} (\code{WFDB\_Annotation}) & \code{time} (in sample intervals),
  \code{anntyp} (the annotation code), \code{subtyp}, \code{chan}, \code{num},
  \code{aux} (\code{NULL}, or a raw vector whose first byte is the length of the text
  that follows) \\
\code{calinfo} (\code{WFDB\_Calinfo}) & \code{low}, \code{high} (the levels of the
  calibration pulse, in physical units), \code{scale} (physical units per cm),
  \code{sigtype} (matched against the beginning of a signal description), \code{units},
  \code{caltype} (\code{WFDB\_AC\_COUPLED} or \code{WFDB\_DC\_COUPLED}, or-ed with
  \code{WFDB\_CAL\_SQUARE}, \code{WFDB\_CAL\_SINE} or \code{WFDB\_CAL\_SAWTOOTH}) \\
\end{xltabular}

\subsection{The settings functions}
\label{sec:r-settings}

The functions of \cref{tab:r-settings} are additions of WFDB-x to the library: they set
where files are read and written, network access and legacy mode, for the high-level
functions, the library functions and the tools alike (\cref{chap:files,chap:behavior}).
Each \code{set} function returns the previous setting, invisibly, so that it can be
restored. The C function \code{setwfdb(p)} sets the input directory too, and
\code{getwfdb()} returns it; in legacy mode they set and return the WFDB path, as in C.

\begin{xltabular}{\linewidth}{@{}lL@{}}
\caption{The settings functions}
\label{tab:r-settings} \\
\toprule
\tabhead Function & \tabhead Effect and result \\
\midrule
\endfirsthead
\toprule
\tabhead Function & \tabhead Effect and result \\
\midrule
\endhead
\bottomrule
\endfoot
\code{wfdbx\_set\_input\_dir(dir)} & sets the input directory: a local directory, a
  \code{\textasciitilde} path or a URL; \code{NULL} or \code{""} restores the default
  (\code{WFDBX\_INPUT\_DIR}, or else the current directory). Returns the previous input
  directory, or \code{NULL} if none was set \\
\code{wfdbx\_get\_input\_dir()} & returns the input directory, or \code{NULL} if none is
  set \\
\code{wfdbx\_set\_output\_dir(dir)} & sets the output directory: a local directory or a
  \code{\textasciitilde} path (a URL is refused, with a message, and the setting is
  unchanged); \code{NULL} or \code{""} restores the default (\code{WFDBX\_OUTPUT\_DIR}, or
  else the current directory). Returns the previous output directory, or \code{NULL} \\
\code{wfdbx\_get\_output\_dir()} & returns the output directory, or \code{NULL} if none
  is set \\
\code{wfdbx\_set\_network(on)} & turns network access on (\code{TRUE} or 1) or off
  (\code{FALSE} or 0); \code{NULL} or a negative number restores the default (on if
  \code{WFDBX\_NETWORK} is a nonzero number). Returns the previous setting, 1 or 0 \\
\code{wfdbx\_get\_network()} & returns 1 if network access is on, 0 otherwise \\
\code{wfdbx\_set\_legacy(on)} & turns legacy mode on (\code{TRUE}) or off, for the whole
  library. Returns the previous setting, \code{TRUE} or \code{FALSE} \\
\code{wfdbx\_get\_legacy()} & returns \code{TRUE} if legacy mode is on, \code{FALSE}
  otherwise \\
\code{wfdbx\_legacy\_args(argv)} & if \code{argv[2]} (\code{argv}: the argument vector,
  program name first) is \code{"--legacy"}, turns legacy mode on; returns
  \code{list(argv, legacy)}: the arguments without \code{"--legacy"}, and whether it was
  found \\
\code{wfdbx\_time\_needs\_freq(string)} & returns 1 if \code{strtim(string)} needs a
  sampling frequency, 0 otherwise \\
\code{wfdbx\_write\_errors()} & returns the number of write errors since the R session
  started (output files that could not be written, flushed or closed) \\
\code{wfdbx\_note\_write\_error()} & counts a write error, for a program that writes a
  file itself \\
\code{wfdbx\_input\_path(name)} & returns the path where a program finds a file that it
  reads itself, as \code{wfdb\_open} finds a file of that name: in normal mode, a
  relative name under the input directory, if one is set, a \code{\textasciitilde} name
  expanded, an absolute name as given; in legacy mode, the name as given. \code{""} and
  \code{"-"} are returned as given; \code{NULL} for \code{NULL} or a URL, or when the
  input directory makes the name a URL \\
\code{wfdbx\_output\_path(name)} & returns the path where a program creates a file that
  it writes itself, as \code{wfdb\_open} creates a file of that name: in normal mode, a
  relative name under the output directory, if one is set, a \code{\textasciitilde} name
  expanded, an absolute name as given; in legacy mode, any name under the output
  directory, if one is set, as a plain prefix. When the output directory is prefixed,
  the missing directories of the path are created. \code{""} and \code{"-"} are returned
  as given; \code{NULL} for \code{NULL} or a URL \\
\end{xltabular}

\subsection{Examples}

Reading the samples and the annotations of a record, as a C program does (the record
\code{100s} is read from the current directory, a copy of \file{tests/data}):

\begin{lstlisting}[language=R,style=example]
library(wfdbx)

r <- isigopen("100s", 2)                 # open 2 signals
nsig <- r$nsig
cat(nsig, sapply(r$si, function(s) paste(s$desc, s$fmt, s$gain)), "\n")
# 2 MLII 212 200 V5 212 200
print(sampfreq(NULL))                    # 360: the frequency of the open record

stat <- isigsettime(strtim("0:30"))      # skip to 30 seconds
g <- getvec()                            # list(stat, v)
cat(g$stat, g$v, sapply(0:(nsig - 1), function(k) aduphys(k, g$v[k + 1])), "\n")
# 2 947 968 -0.385 -0.28

f <- getframes(360)                      # the next second, as a matrix
cat(dim(f$v), f$stat, "\n")              # 360 2 2

stat <- annopen("100s", list(anninfo("atr", WFDB_READ)))
n <- 0
while ((g <- getann(0))$stat == 0) {     # 0: an annotation was read; -1: the end
  if (n < 2) cat(mstimstr(g$annot$time), annstr(g$annot$anntyp),
                 aux_string(g$annot$aux), "\n")
  n <- n + 1
}
cat(n, "annotations\n")                  # 75 annotations

wfdbquiet()                              # no messages on the standard error
cat(isigopen("nosuchrecord", 0)$nsig, wfdberror())
# -1 init: can't open header for record nosuchrecord
wfdbverbose()
stat <- wfdbquit()                       # close everything
\end{lstlisting}

Writing a record and an annotation file, sample by sample and annotation by annotation:

\begin{lstlisting}[language=R,style=example]
library(wfdbx)

out <- list(siginfo(fname = "ramp.dat", desc = "ramp", units = "mV", gain = 200,
                    fmt = 16, adcres = 12))
print(osigfopen(out))                    # 1: one output signal
stat <- setsampfreq(250)
for (k in 0:499) putvec(k)
print(newheader("ramp"))                 # 0: ramp.hea is written
stat <- annopen("ramp", list(anninfo("qrs", WFDB_WRITE)))
for (t in c(100, 350)) putann(0, annotation(time = t, anntyp = NORMAL))
stat <- wfdbquit()                       # completes ramp.dat and ramp.qrs

print(readLines("ramp.hea", n = 1))              # "ramp 1 250 500"
print(read_annotations("ramp", "qrs")$sample)    # 100 350
\end{lstlisting}

% ==========================================================================================
\section{The command-line tools in R}
\label{sec:r-tools}

The 33 command-line tools (\code{rdsamp}, \code{rdann}, \code{wrsamp}, \code{xform},
\ldots) are run by the driver script \file{wfdbx.R} with \cmd{Rscript}, followed by the
name of the tool and its arguments. The tools, their options, their output and their exit
status are those of the C tools, and are documented in \cref{chap:tools}. In a copy of
the repository the script is \file{r/inst/bin/wfdbx.R}, which runs without the package
being installed; in an installed package its path is \code{system.file("bin", "wfdbx.R",
package = "wfdbx")}. A shell alias makes it convenient:

\begin{lstlisting}[style=shell]
alias wfdbx='Rscript /path/to/wfdb-x/r/inst/bin/wfdbx.R'
export WFDBX_INPUT_DIR=/path/to/wfdb-x/tests/data   # the directory of 100s
wfdbx rdsamp -r 100s -t 10 -p     # the first 10 seconds of 100s, in physical units
wfdbx rdann -r 100s -a atr        # its annotations
wfdbx rdsamp --legacy -r 100s     # legacy mode: --legacy as the first argument
\end{lstlisting}

The driver script also takes these options, instead of a tool:

\begin{table}[htbp]
\centering
\caption{The options of the driver script \file{wfdbx.R}}
\label{tab:r-wfdbx-options}
\begin{tabularx}{\linewidth}{@{}lL@{}}
\toprule
\tabhead Command & \tabhead Effect \\
\midrule
\cmd{Rscript wfdbx.R TOOL [ARGUMENTS ...]} & runs the tool \code{TOOL}, and exits with
  its status \\
\cmd{Rscript wfdbx.R --list} & lists the names of the tools, one per line \\
\cmd{Rscript wfdbx.R --version} & prints the version, \code{WFDB-x 1.2.1} \\
\cmd{Rscript wfdbx.R --help}, \cmd{Rscript wfdbx.R -h} & prints the usage \\
\bottomrule
\end{tabularx}
\end{table}

Without arguments, the script prints the usage and exits with status 1; an unknown tool
is an error (``wfdbx: unknown tool 'TOOL' (try 'wfdbx -{}-list')'', status 1), and an R
error in a tool is reported as \code{wfdbx: error: MESSAGE} (status 1; with the
environment variable \code{WFDBX\_DEBUG} set, the call stack follows). When an output file
or the standard output cannot be written (a full disk, for example), a tool reports it and
exits with status 1 (legacy mode keeps the original status, 0). The files that a tool
reads or creates itself follow the input and output directories, as the records do: the
\opt{-i} file of \code{wrsamp}, \code{edf2mit}, \code{wav2mit}, \code{a2m} and
\code{ad2m}, the input files of \code{ahaecg2mit} and the \opt{-S} script of
\code{xform} are read from the input directory, and the WAV file of \code{mit2wav}, the
EDF file of \code{mit2edf}, the master header of \code{snip -m} and the \file{.des},
\file{.ecg}, \file{.key} and \file{.txt} files of \code{parsescp} are created in the
output directory (\code{wfdbx\_input\_path} and \code{wfdbx\_output\_path},
\cref{tab:r-settings}); such a file cannot be read from a URL. The script runs the tools
in the C locale, with R's warnings turned off, so that their output is that of the C
tools.

\paragraph{The start-up cache.} When run from a copy of the repository, the script keeps a
compiled copy of the sources in \file{\textasciitilde/.cache/R/wfdbx/} (or under
\code{\$XDG\_CACHE\_HOME}), so that a tool starts quickly; \code{WFDBX\_NO\_CACHE=1} turns
this off.

\paragraph{Windows.} The tools write their output in binary through \file{/dev/stdout} and
\file{/dev/stderr}. On Windows these do not exist, so the output is written as text and
NUL bytes are lost.

\paragraph{The tools from R.} \code{wfdbx\_tool} runs a tool in the R session, and
\code{wfdbx\_tools} lists them.

\apientry{wfdbx\_tool}{Run a command-line tool}
\begin{lstlisting}[language=R,style=signature]
wfdbx_tool(tool, args = character(0))
\end{lstlisting}
Runs one of the command-line tools with the given arguments, as the command line would,
in the R session. The output of the tool goes directly to the standard output and error
of the R process: \code{sink} and \code{capture.output} do not capture it, and a graphical
console may not show it. To capture the output, the tool is run in a separate process
with the driver script, as in \code{system2("Rscript", c(driver, "rdsamp", ...), stdout =
TRUE)}.
\begin{params}
\item[tool] The name of the tool, a character string (one of \code{wfdbx\_tools()}).
\item[args] The command-line arguments of the tool, as a character vector (one element
  per argument, without the program name). Default \code{character(0)}: none.
\end{params}
\returns{The exit status of the tool (0 on success), a number.}
\raises{Stops with an error if \code{tool} is not the name of a tool (``unknown tool
'TOOL'\,'').}
\begin{lstlisting}[language=R,style=example]
library(wfdbx)

status <- wfdbx_tool("rdsamp", c("-r", "100s", "-t", "s2"))     # the first 2 samples
# rdsamp prints two lines: the sample numbers 0 and 1, and the samples 995 1011
print(status)                            # 0
\end{lstlisting}

\apientry{wfdbx\_tools}{List the command-line tools}
\begin{lstlisting}[language=R,style=signature]
wfdbx_tools()
\end{lstlisting}
Returns the names of the command-line tools.
The function has no parameters.
\returns{A character vector of the names of the tools, in alphabetical order (C
collation).}
\raises{Nothing.}
\begin{lstlisting}[language=R,style=example]
library(wfdbx)

print(length(wfdbx_tools()))             # 33
print(head(wfdbx_tools(), 4))            # "a2m" "ad2m" "ahaecg2mit" "ann2rr"
\end{lstlisting}

\paragraph{Environment variables.} The tools, like the rest of the package, use the
environment variables of \cref{tab:r-environment}. They are read when they are first
needed in an R session.

\begin{table}[htbp]
\centering
\caption{The environment variables used by the R package}
\label{tab:r-environment}
\begin{tabularx}{\linewidth}{@{}lL@{}}
\toprule
\tabhead Variable & \tabhead Meaning \\
\midrule
\code{WFDBX\_INPUT\_DIR} & the input directory (a directory, a \code{\textasciitilde}
  path or a URL), when \code{wfdbx\_set\_input\_dir} does not set one \\
\code{WFDBX\_OUTPUT\_DIR} & the output directory, when \code{wfdbx\_set\_output\_dir} does
  not set one \\
\code{WFDBX\_NETWORK} & network access is on if it is a nonzero number, when
  \code{wfdbx\_set\_network} does not set it \\
\code{WFDBANNSORT} & \code{0} to keep annotation files in the order written (no sorting
  when they are closed) \\
\code{WFDBGVMODE} & the default getvec mode (\code{setgvmode} with a negative mode) \\
\code{WFDBCAL} & the calibration file read by \code{calopen(NULL)} (default
  \file{wfdbcal}) \\
\code{WFDB} & the WFDB path, used in legacy mode only (outside legacy mode, it is
  ignored, with a message) \\
\code{WFDB\_FLAC\_CHECK\_MD5} & if set, the MD5 signature of FLAC signal files is
  verified when they are read \\
\code{WFDBX\_DEBUG} & if set, the driver script prints the call stack of an R error in a
  tool \\
\code{WFDBX\_NO\_CACHE} & \code{1} to turn the start-up cache of the driver script off \\
\code{XDG\_CACHE\_HOME} & the directory of the start-up cache (default
  \file{\textasciitilde/.cache}) \\
\bottomrule
\end{tabularx}
\end{table}

Remote files are downloaded with the program \cmd{curl} if it is installed, and with R's
\code{url()} otherwise; each is read once and kept for the rest of the R session, and a
download gives up when it cannot connect, or receives nothing, for 30 seconds.

% ==========================================================================================
\section{Testing}
\label{sec:r-testing}

The unit tests of the R package are in \file{r/tests}: the runner \file{run\_tests.R} and
the files \file{unit/test\_*.R}, written with base R only. From the root of the
repository:

\begin{lstlisting}[style=shell]
Rscript r/tests/run_tests.R              # all the unit tests
Rscript r/tests/run_tests.R metadata     # the files whose name matches "metadata"
\end{lstlisting}

The runner loads the sources of \file{r/R} (or the installed package, under \cmd{R CMD
check}), runs each test file, and prints one line per file and one per failed check; it
exits with status 1 if a check fails. The tests check the library functions with expected
values of the C library, the high-level functions, the bulk paths against the
sample-by-sample paths, \code{printf} and \code{strtod} against the GNU C library, the
FLAC codec, \code{check\_record}, and the header metadata functions against the cases
shared by the seven languages (in \file{tests/metadata}: results, warnings and the
headers, byte for byte). Network access is tested against a local HTTP server
(127.0.0.1, run with \cmd{python3}; skipped without it). The tests work in temporary
directories, with copies of the test records of \file{tests/data}; the tests that need
\file{tests/data} are skipped when it is not available (in the installed package).

The package check builds the package and runs its examples and tests:

\begin{lstlisting}[style=shell]
R CMD build r && R CMD check --no-manual wfdbx_*.tar.gz
\end{lstlisting}

The cross-language conformance suite compares the output of the tools with that of the
original C programs, case by case: 1135 cases, and 883 in legacy mode, which runs each
tool with \opt{--legacy}. It runs the driver script \file{r/inst/bin/wfdbx.R}. From the
root of the repository:

\begin{lstlisting}[style=shell]
PYTHONPATH=$PWD/python/src python3 tests/conformance/harness.py check --impl r -j8
PYTHONPATH=$PWD/python/src python3 tests/conformance/harness.py check --impl r -j8 --legacy
\end{lstlisting}

% The MATLAB/Octave chapter of the WFDB-x user's manual. It follows the sections of the
% Python chapter (chapters/python.tex), in the same order. Each \apientry of the section
% "The high-level API" is a function of matlab/+wfdbx listed under "High-level functions"
% or "Header metadata" in matlab/+wfdbx/Contents.m; its params name the positional
% arguments of the function line, then the name-value options in the order of the
% structure of defaults that the function parses (wfdbx.internal.hl_options). The
% examples of style "example" were run with GNU Octave 8.4, in a copy of tests/data that
% is the current folder (the paths /path/to/wfdb-x replaced by those of the repository);
% check_manual.py does not run them.
\chapter{MATLAB and Octave}
\label{chap:matlab}
\chapterlanguage{matlab}

The MATLAB/Octave implementation of WFDB-x is the package \code{wfdbx}: the folder
\file{matlab/+wfdbx} of the repository, with its internal functions in
\file{+wfdbx/+internal} and the command-line tools in \file{+wfdbx/+apps}. It is plain
MATLAB code, which runs unchanged in GNU Octave and in MATLAB, needs no toolbox, and never
calls the C library. Every function is called with the name of the package,
\code{wfdbx.NAME}.

It offers WFDB-x in three ways:
\begin{itemize}
\item \textbf{The high-level API} (\cref{sec:matlab-api}): functions that read and write
  whole records, headers, annotation files and header metadata as matrices of doubles,
  structures and cell arrays, in one call, with name-value options, and that throw MATLAB
  errors (with an identifier, such as \code{wfdbx:read\_record}) on failure. They read
  signal files in large blocks, and are the fastest way to read a record. They are an
  addition of WFDB-x: the C package has no equivalent. They are the recommended interface
  for new code and analysis.
\item \textbf{The WFDB library functions} under their C names, called with the name of
  the package as a prefix: \code{wfdbx.isigopen}, \code{wfdbx.getvec},
  \code{wfdbx.getann}, \code{wfdbx.putann}, and so on. They have the
  behavior of the C library (\cref{sec:matlab-library}). They are documented once, for
  all the languages, in \cref{chap:library}.
\item \textbf{The command-line tools} (\code{rdsamp}, \code{rdann}, \code{wrsamp},
  \code{xform} and 29 others), with the options, output and exit status of the C tools,
  run from the shell through an Octave script or a MATLAB shell script, or inside a
  session as functions (\cref{sec:matlab-tools}). They are documented in
  \cref{chap:tools}.
\end{itemize}
The high-level functions return exactly the samples and annotations that the library
functions (and the C library) give.

% ==========================================================================================
\section{Installation and requirements}
\label{sec:matlab-install}

\subsection{Requirements}

WFDB-x for MATLAB and Octave needs GNU Octave 6 or later, or MATLAB R2019b or later, and
no toolbox or other package. It uses only the language features and functions that both
products provide (no \code{string} objects, no \code{arguments} blocks, no Octave-only
syntax), so the same files run in both. It has no compiled code, and needs no C compiler,
no WFDB library and no libFLAC: the FLAC formats 508, 516 and 524 are read and written by
WFDB-x's own codec.

The code is tested with GNU Octave 8.4 (\cref{sec:matlab-testing}). It has not yet been
run in MATLAB: it is written to the rules that both products share, which are checked by
reading the code, and its MATLAB command-line driver, \file{matlab/bin/wfdbx\_matlab.sh},
has been checked in Octave only.

Remote files (\cref{chap:files}) are downloaded with the program \cmd{curl} when it is
installed, and otherwise with \code{webread} (Octave) or \code{websave} (MATLAB).

\subsection{Installing}

The package is used from a copy of the repository: its folder \file{matlab}, which holds
\file{+wfdbx}, is added to the path:

\begin{lstlisting}[language=Matlab]
addpath('/path/to/wfdb-x/matlab')
savepath      % optional: keep the folder on the path in later sessions
\end{lstlisting}

The folder \file{matlab} also holds copies of the license files of the repository:
\file{LICENSE} (the GNU Lesser General Public License, version 2.1, of the library),
\file{LICENSE-GPL} (the GNU General Public License, version 3, of the command-line tools:
\file{+wfdbx/+apps}, the drivers of \file{bin} and the dispatcher
\file{+wfdbx/+internal/cli\_main.m} and \file{app\_run.m}; the internal functions that
the library also uses, such as \file{tool\_command.m}, which \code{wfdbx.oannclose} calls
to run \cmd{sortann} in legacy mode, stay under the LGPL)
and \file{NOTICE}, which summarizes the licenses and gives the BSD-3-Clause notice of
libFLAC, whose stream decoder the FLAC code of \file{+wfdbx/+internal} translates. The
folder is the MATLAB and Octave distribution of WFDB-x: a copy of it that is passed on
keeps these three files.

\code{help wfdbx.Contents} lists the functions, by group (in MATLAB, \code{help wfdbx}
also does), and \code{help wfdbx.NAME} describes each one. \file{+wfdbx/Contents.m} gives
the version of the package, \code{Version 1.2.2}, and the command-line driver prints it:

\begin{lstlisting}[style=shell]
octave-cli --no-gui --norc -q matlab/bin/wfdbx_cli.m --version    # WFDB-x 1.2.2
\end{lstlisting}

\subsection{The sample record}
\label{sec:matlab-resources}

The repository includes the sample record \code{100s} in \file{tests/data}: the first
minute of record 100 of the MIT-BIH Arrhythmia Database, two ECG signals (\code{MLII} and
\code{V5}) sampled at 360\,Hz, with the reference beat annotations (annotator
\code{atr}), in the files \file{100s.hea}, \file{100s.dat} and \file{100s.atr}. The
package has no function that returns this folder: a program names it, or sets it as the
input directory. Work that uses the MIT-BIH Arrhythmia Database cites it as described on
PhysioNet.

% ==========================================================================================
\section{Quick start}
\label{sec:matlab-quickstart}

This example reads the header, ten seconds of the signals and the annotations of the
sample record \code{100s}, from the folder \file{tests/data} of a copy of the repository
(\code{'input\_dir'}), and computes the mean heart rate. It then writes the first ten
seconds as a new record, \code{first10}, in the current folder, with two metadata entries
in its header.

\begin{lstlisting}[language=Matlab,style=example]
addpath('/path/to/wfdb-x/matlab');
data = '/path/to/wfdb-x/tests/data';     % the folder of the sample record 100s

% The header only: no samples are read.
hdr = wfdbx.read_header('100s', 'input_dir', data);
fprintf('%d %g %d\n', hdr.nsig, hdr.fs, hdr.nsamples);   % 2 360 21600
fprintf('%s\n', hdr.siginfo.desc);                       % MLII, then V5

% Seconds 10 to 20 of both signals, in physical units: one row per sample, one column
% per signal.
rec = wfdbx.read_record('100s', 'sampfrom', '0:10', 'sampto', '0:20', ...
    'physical', true, 'input_dir', data);
disp(size(rec.signals));                % 3600 2
fprintf('%g %g\n', rec.signals(1, :));  % -0.39 -0.275

% The beat annotations, and the mean heart rate.
ann = wfdbx.read_annotations('100s', 'atr', 'input_dir', data);
fprintf('%d %s %s %s\n', numel(ann.samples), ann.symbols{1:3});   % 75 + N N
beats = ann.samples(wfdbx.isqrs(ann.codes) == 1);
rr = diff(beats) / ann.fs;              % RR intervals, in seconds
fprintf('%d beats, %.1f beats per minute\n', numel(beats), 60 / mean(rr));
% 74 beats, 73.9 beats per minute

% Write the first 10 seconds, in ADC units, as the record "first10" (first10.hea and
% first10.dat, in the current folder), with the signal specifications of 100s.
x = wfdbx.read_record('100s', 'sampto', '0:10', 'input_dir', data);
wfdbx.write_record('first10', x.signals, x.fs, 'fmt', 16, 'gains', x.gains, ...
    'baselines', x.baselines, 'adczeros', [x.siginfo.adczero], ...
    'adcres', [x.siginfo.adcres], 'units', 'mV', 'names', x.names, ...
    'metadata', {'Source', 'MIT-BIH record 100'});

% Header metadata: "# Key: value" lines of the header.
wfdbx.set_metadata('first10', 'Age', '69');
m = wfdbx.read_metadata('first10')';    % (transposed, for fprintf)
fprintf('%s: %s\n', m{:});      % Source: MIT-BIH record 100, then Age: 69
wfdbx.check_record('first10');  % prints: no problems found
\end{lstlisting}

The header \file{100s.hea} gives no units: the example writes \code{'mV'} explicitly,
since WFDB-x never assumes units, a gain, an ADC resolution or a sampling frequency that
a record does not give (\cref{chap:behavior}).

% ==========================================================================================
\section{The high-level API}
\label{sec:matlab-api}

The high-level API is made of the functions of this section, in the package
\code{wfdbx} (\code{wfdbx.read\_record}, \ldots), and of the structures of their results
(\cref{sec:matlab-results}). Each is one file of \file{matlab/+wfdbx};
\code{help wfdbx.NAME} gives its documentation.

\subsection{Conventions}
\label{sec:matlab-conventions}

\paragraph{Arguments and options.} The first arguments of a function are positional (the
record, the annotator, the samples, \ldots); the others are options, given as name-value
pairs after them, in any order: \code{wfdbx.read\_record('100s', 'sampto', 1000,
'physical', true)}. Option names are case-insensitive. An option that is not given keeps
its default, as listed in the call syntax of each entry; \code{[]} means ``not given'' for
the options whose default is \code{[]}. An unknown option, or a name without a value, is
an error (\cref{sec:matlab-warnings}). Text arguments are \code{char} row vectors or
MATLAB strings (converted to \code{char}).

\paragraph{Record names and directories.} A record is named as in the other parts of
WFDB-x (\cref{chap:records,chap:files}), without the \code{.hea} extension of its header
(the readers and the writers, \code{wfdbx.write\_record} and
\code{wfdbx.write\_annotations}, also accept a trailing \code{.hea}, which is removed). A
full path, absolute (\code{'/data/mitdb/100'}) or
in the home directory (\code{'\textasciitilde/Documents/100s'}), is used as given; a
relative path, such as \code{'100s'} or \code{'data/100s'}, is read from the input
directory if one is set, otherwise from the current folder, and written to the output
directory if one is set, otherwise to the current folder (missing folders are created);
a URL (\code{'https://physionet.org/files/mitdb/1.0.0/100'}) is read only if network
access is on, and cannot be written. A writer given a name with a path creates the files
in that folder, and the header names the record and its signal file without the path
(\code{'rec'}, \code{'rec.dat'}), so that the record opens under any path that leads to
it; in legacy mode, as in the original package, the header gives the name with its path.
The input and output directories are set with \code{wfdbx.wfdbx\_set\_input\_dir} and
\code{wfdbx.wfdbx\_set\_output\_dir} (\cref{sec:matlab-settings}) or with the
environment variables \code{WFDBX\_INPUT\_DIR} and \code{WFDBX\_OUTPUT\_DIR}. The
functions that read take the options
\code{'input\_dir'} and \code{'network'}, and those that write take \code{'dir'}: they
apply to that one call, and the previous settings are restored when it returns, also after
an error. \code{'input\_dir'} is optional, and applies to relative names only, not to full
paths or URLs; with \code{'dir'}, the record name must be a plain name (\code{'rec'}),
without a folder. The signal files that a header names are read from the header's
folder.

\paragraph{Times.} The options \code{'sampfrom'} and \code{'sampto'} are sample numbers
(whole numbers) or WFDB time strings, as \code{strtim} reads them: \code{'0:10'} (10
seconds), \code{'1:02:03.5'}, \code{'s3600'} (sample 3600), \code{'e'} (the end of the
record), or a time of day in brackets, \code{'[14:00:00]'}. A time in seconds needs the
sampling frequency; when the record does not give it, such a time is an error.

\paragraph{Numbers and arrays.} Samples, times, counts and codes are doubles. Signal
numbers start at 0, as in C. The signals of a record are a matrix with one row per sample
and one column per signal; the annotations of a file are column vectors and cell columns,
one row per annotation.

\paragraph{Per-signal and per-annotation options.} The options that describe each signal
(\code{'fmt'}, \code{'gains'}, \code{'units'}, \ldots) or each annotation
(\code{'subtypes'}, \code{'chans'}, \ldots) take one value for all, or a vector (for
numbers) or cell array (for text) with exactly one value per signal or per annotation.

\paragraph{Signal parameters.} No sampling frequency, gain, ADC resolution or units is
ever assumed. Where one is needed and the record does not give it, it must be given as an
argument; otherwise the function throws an error. A record whose header gives no
sampling frequency has \code{fs} 0.

\paragraph{Text.} Text is held as bytes: each character of a \code{char} vector is one byte
of a file (codes 0 to 255), so every byte sequence round-trips unchanged, and UTF-8 text
is stored as its bytes. In Octave, a \code{char} array holds bytes already; in MATLAB, a
character above 255 cannot be written to a header. The metadata functions also accept a
\code{uint8} vector of bytes as a key or value.

\paragraph{Legacy mode.} With \code{'legacy', true}, a function behaves as in the
original WFDB package, with its assumed signal parameters and its search for files
(\cref{chap:behavior}); the library's legacy mode is on for the duration of the call
only. With \code{false} (the default), the current setting of
\code{wfdbx.wfdbx\_set\_legacy} applies, which is off unless it was turned on.

\paragraph{State.} The functions that read and write records and annotation files call
\code{wfdbx.wfdbquit()}, which closes anything opened with the WFDB library functions,
and close their own files before they return, whatever the outcome.
\code{wfdbx.read\_record} sets the getvec mode for its own reading only. The functions
\code{wfdbx.read\_record}, \code{wfdbx.read\_header}, \code{wfdbx.read\_annotations},
\code{wfdbx.write\_record}, \code{wfdbx.write\_annotations} and
\code{wfdbx.check\_record} leave the getvec mode (\code{wfdbx.getgvmode}) and the time
resolution of output annotators (\code{wfdbx.getafreq}) as they were when they return,
also after an error: opening an EDF record sets both (it adds \code{WFDB\_HIGHRES} to the
getvec mode, and sets the time resolution to the record's sampling frequency when none is
set), as in C, and \code{wfdbx.wfdbquit()} resets neither. The metadata
functions open no signals, and leave the library's open files as they are.

\paragraph{The examples} of this section read the test records of WFDB-x (among them
\code{100s}) from the current folder, and write records of their own there: they run in a
copy of the folder \file{tests/data} of the repository, made the current folder (with
\file{matlab} on the path).

% ------------------------------------------------------------------------------------------
\subsection{Reading records}
\label{sec:matlab-reading}

\apientry{wfdbx.read\_record}{Read the signals of a record}
\begin{lstlisting}[language=Matlab,style=signature]
rec = wfdbx.read_record(record, 'channels', [], 'sampfrom', 0, 'sampto', [], ...
    'physical', false, 'highres', false, 'pad', false, 'input_dir', [], ...
    'network', [], 'legacy', false, 'checksum_warning', true)
\end{lstlisting}
Reads the signals of a record, or part of it, into a matrix of doubles with one row per
sample and one column per signal, together with the description of the record. The
samples are those that \code{getvec} returns (\cref{chap:library}), read in large blocks.
Invalid samples are \code{wfdbx.WFDB\_INVALID\_SAMPLE()} ($-32768$) in ADC units, and NaN
in physical units. When the read reaches the end of the record, the samples are compared
with the checksums of the header (\cref{sec:matlab-warnings}).
\begin{params}
\item[record] Record name (a \code{char} vector or a MATLAB string), without the \code{.hea} extension (a trailing \code{.hea}
  is also accepted): a full path, absolute (\code{'/data/mitdb/100'}) or in the home
  directory (\code{'\textasciitilde/Documents/100s'}), used as given; a relative path
  (\code{'100s'}, \code{'data/100s'}), read from the input directory if one is set,
  otherwise from the current folder; or a URL
  (\code{'https://physionet.org/files/mitdb/1.0.0/100'}), read only if network access is
  on (\cref{sec:record-paths}).
\item[channels] The signals to read, in the order in which they are returned: a vector of
  signal numbers (from 0), or a cell array of signal numbers and signal descriptions (as
  in \code{rec.names}, such as \code{'MLII'}), which may be mixed. Default \code{[]}: all
  the signals, in the order of the header.
\item[sampfrom] The first sample to read: a sample number (a whole number $\ge 0$) or a
  time string. Default \code{0}: the beginning of the record.
\item[sampto] The sample after the last one to read: a sample number (a whole number) or
  a time string. A value at or before \code{sampfrom} reads no samples. Default
  \code{[]}: the end of the record.
\item[physical] \code{false} (the default) to return ADC units; \code{true} to return
  physical units, computed as $(\mathit{adc} - \mathit{baseline}) / \mathit{gain}$. Every
  signal read must then have a gain.
\item[highres] \code{false} (the default) to read a multi-frequency record at its frame
  frequency, with the samples of the faster signals averaged to one per frame;
  \code{true} to read every sample, at the highest sampling frequency of the record (the
  samples of the slower signals are repeated). \code{rec.fs} is then that frequency.
\item[pad] \code{true} to replace each invalid sample by the previous valid sample of the
  same signal. Default \code{false}.
\item[input\_dir] The input directory for this call only (optional): a local folder, a
  \code{\textasciitilde} path or the URL of a directory, under which relative record names
  are read (\code{''}: the default). It does not apply to full paths or URLs. Default
  \code{[]}: the current setting (\code{wfdbx.wfdbx\_set\_input\_dir}).
\item[network] \code{true} to allow reading URLs for this call only, \code{false} to
  forbid it. Default \code{[]}: the current setting (\code{wfdbx.wfdbx\_set\_network}).
\item[legacy] \code{true} to behave as the original WFDB package (legacy mode, for this
  call only). Default \code{false}.
\item[checksum\_warning] \code{true} (the default) to issue the warning
  \code{wfdbx:checksum} when the samples do not match the checksums of the header;
  \code{false} for no warning. The samples are returned either way, and
  \code{rec.checksum\_ok} is \code{false}.
\end{params}
The memory allocated does not depend on the length that the header gives, which may be
far more than the signal files hold: the samples are read in chunks of about $2^{22}$,
joined at the end.
\returns{A structure (\cref{tab:matlab-record}).}
\raises{Identifier \code{wfdbx:read\_record}: the record cannot be read (``can't read
record REC''), or its signals cannot be opened; a signal file is shorter than its header
says, or cannot be read (a decoder or I/O error): ``can't read signal data of record REC
(getvec status -3)'', even if the checksums do not match either (the checksum error is
found in the frame in which the file ended), so that no incomplete frame is ever
returned; a requested signal is not in the record; \code{sampfrom} is negative, or
\code{sampfrom} or \code{sampto} is not a whole number (a fraction, NaN or an infinite
value: ``read\_record: sampfrom must be a whole number''), or is neither a number nor a
time string (``\ldots{} must be a sample number or a time string''); a time string needs the
sampling frequency and the header gives none; \code{'physical'} is \code{true} and a
signal read is uncalibrated; or an option is not valid. Warning \code{wfdbx:checksum} on
a checksum error (see \code{checksum\_warning}).}
\begin{lstlisting}[language=Matlab,style=example]
% Seconds 10 to 20 of lead MLII, in physical units.
rec = wfdbx.read_record('100s', 'channels', {'MLII'}, 'sampfrom', '0:10', ...
    'sampto', '0:20', 'physical', true);
disp(size(rec.signals));                         % 3600 1
fprintf('%s %g %d\n', rec.names{1}, rec.fs, rec.sampfrom);   % MLII 360 3600

% The first 3 samples of both signals, in ADC units, V5 first.
rec = wfdbx.read_record('100s', 'channels', [1 0], 'sampto', 3);
disp(rec.signals);          % three rows of 1011 995
disp(rec.names);            % {'V5', 'MLII'}
disp([rec.gains; rec.baselines]);      % 200 200; 1024 1024

% The whole record: the checksums are verified.
rec = wfdbx.read_record('100s');
fprintf('%d %d %d\n', size(rec.signals), rec.checksum_ok);    % 21600 2 1
\end{lstlisting}
A multi-frequency record (here the test record \code{mfreq}, whose signals have 1, 2
and 4 samples per frame at a frame frequency of 100\,Hz) is read at its frame frequency,
or at its highest sampling frequency:
\begin{lstlisting}[language=Matlab,style=example]
rec = wfdbx.read_record('mfreq');
fprintf('%g %d %d\n', rec.fs, size(rec.signals));     % 100 600 3
rec = wfdbx.read_record('mfreq', 'highres', true);
fprintf('%g %d %d\n', rec.fs, size(rec.signals));     % 400 2400 3
disp([rec.siginfo.spf]);                              % 1 2 4
\end{lstlisting}

\apientry{wfdbx.read\_header}{Read the header of a record}
\begin{lstlisting}[language=Matlab,style=signature]
hdr = wfdbx.read_header(record, 'input_dir', [], 'network', [], 'legacy', false)
\end{lstlisting}
Reads the header of a record, without opening its signal files. The result describes the
record and each of its signals, and holds the info strings and the metadata entries of the
header. For a multi-segment record, the master header is read, and \code{nsamples} is the
length of the whole record. The metadata are read from the \file{.hea} file only; an EDF
record (which has no \file{.hea} file) and a header read from the standard input
(\code{'-'}) have none.
\begin{params}
\item[record] Record name, without the \code{.hea} extension (a trailing \code{.hea}
  is also accepted): a full path, absolute (\code{'/data/mitdb/100'}) or in the home
  directory (\code{'\textasciitilde/Documents/100s'}), used as given; a relative path
  (\code{'100s'}, \code{'data/100s'}), read from the input directory if one is set,
  otherwise from the current folder; or a URL
  (\code{'https://physionet.org/files/mitdb/1.0.0/100'}), read only if network access is
  on (\cref{sec:record-paths}).
\item[input\_dir] The input directory for this call only (optional), under which
  relative record names are read; it does not apply to full paths or URLs (as for
  \code{wfdbx.read\_record}). Default \code{[]}: the current setting.
\item[network] \code{true} to allow reading URLs for this call only, \code{false} to
  forbid it. Default \code{[]}: the current setting.
\item[legacy] \code{true} to read the header as the original WFDB package does. Default
  \code{false}.
\end{params}
\returns{A structure (\cref{tab:matlab-header}).}
\raises{Identifier \code{wfdbx:read\_header}: the header cannot be found or read
(``can't read header for record REC''), or an option is not valid. Warning
\code{wfdbx:duplicate\_key} for a metadata key that the header has more than once.}
\begin{lstlisting}[language=Matlab,style=example]
hdr = wfdbx.read_header('100s');
fprintf('%s %g %d %d\n', hdr.name, hdr.fs, hdr.nsig, hdr.nsamples);  % 100s 360 2 21600
for k = 1:hdr.nsig
    s = hdr.siginfo(k);
    fprintf('%s %s %d %g %d\n', s.desc, s.fname, s.fmt, s.gain, s.adcres);
end
% MLII 100s.dat 212 200 11
% V5 100s.dat 212 200 11
disp(hdr.comments{1});          % " 69 M 1085 1629 x1": the first info string
disp(isempty(hdr.base_time));   % 1: the header gives no start time
\end{lstlisting}

\apientry{wfdbx.read\_annotations}{Read an annotation file}
\begin{lstlisting}[language=Matlab,style=signature]
ann = wfdbx.read_annotations(record, annotator, 'sampfrom', 0, 'sampto', [], ...
    'fs', [], 'input_dir', [], 'network', [], 'legacy', false)
ann = wfdbx.read_annotations(record, annotator, sampfrom, sampto, ...)
\end{lstlisting}
Reads the annotations of an annotation file, in file order, into column vectors and cell
columns. The sampling frequency, needed to convert times in seconds, is that of the
record's header, or, for a record without a header or whose header gives none, the option
\code{'fs'}. \code{'sampfrom'} and \code{'sampto'} may also be given as the third and
fourth arguments, without their names.
\begin{params}
\item[record] Record name, without the \code{.hea} extension (a trailing \code{.hea}
  is also accepted): a full path, absolute (\code{'/data/mitdb/100'}) or in the home
  directory (\code{'\textasciitilde/Documents/100s'}), used as given; a relative path
  (\code{'100s'}, \code{'data/100s'}), read from the input directory if one is set,
  otherwise from the current folder; or a URL
  (\code{'https://physionet.org/files/mitdb/1.0.0/100'}), read only if network access is
  on (\cref{sec:record-paths}). Its header, if there is one, gives the
  sampling frequency.
\item[annotator] Annotator name: the suffix of the annotation file. \code{'atr'} reads the
  file \file{RECORD.atr}.
\item[sampfrom] The first time to read: a sample number (a whole number $\ge 0$) or a
  time string. Default \code{0}: the beginning of the file.
\item[sampto] The time at which to stop: annotations at or after it are not read. A
  sample number (a whole number) or a time string. Default \code{[]}: the end of the
  file.
\item[fs] The sampling frequency, in Hz (a number $> 0$), to use when the record has no
  header or its header gives none; it is ignored otherwise. Default \code{[]}: none.
\item[input\_dir] The input directory for this call only (optional), under which
  relative record names are read; it does not apply to full paths or URLs (as for
  \code{wfdbx.read\_record}). Default \code{[]}: the current setting.
\item[network] \code{true} to allow reading URLs for this call only, \code{false} to
  forbid it. Default \code{[]}: the current setting.
\item[legacy] \code{true} to behave as the original WFDB package (which assumes 250\,Hz
  for a record without a frequency). Default \code{false}.
\end{params}
\returns{A structure (\cref{tab:matlab-annotations}). Its \code{fs} is the frequency in
which the times are counted: the one written in the annotation file, if there is one, or
else the sampling frequency; it is 0 if none of them is known.}
\raises{Identifier \code{wfdbx:read\_annotations}: the annotation file cannot be opened
(``can't read annotator ANN for record REC''), or it ends unexpectedly or cannot be read
before \code{sampto} (``can't read annotation file ANN of record REC (getann status
S)''); \code{'fs'} is given and is not $> 0$; \code{sampfrom} or \code{sampto} is not
valid, as for \code{wfdbx.read\_record} (``read\_annotations: sampfrom must be a whole
number'', ``\ldots{} must be a sample number or a time string'', ``\ldots{} is beyond
the range of a 64-bit integer'', ``read\_annotations: sampfrom must not be negative''),
which is checked before any file is opened; a time string needs the sampling frequency
and none is known; or an option is not valid. Whatever happens, the annotation file is
closed (\code{wfdbx.wfdbquit}) before the function returns.}
\begin{lstlisting}[language=Matlab,style=example]
ann = wfdbx.read_annotations('100s', 'atr', 'sampfrom', '0:10', 'sampto', '0:30');
fprintf('%d %g\n', numel(ann.samples), ann.fs);       % 24 360
fprintf('%d %s\n', ann.samples(1), ann.symbols{1});   % 3862 N

% The whole file: the first annotation is a rhythm change, with an aux string.
ann = wfdbx.read_annotations('100s', 'atr');
fprintf('%s %d %s\n', ann.symbols{1}, ann.codes(1) == wfdbx.RHYTHM(), ann.aux{1});
% + 1 (N
disp(unique(ann.symbols)');     % {'+', 'A', 'N'}
\end{lstlisting}
For a record without a header, \code{'fs'} gives the frequency in which times in seconds
are converted (the library reports on the standard error that it cannot open the
header):
\begin{lstlisting}[language=Matlab,style=example]
wfdbx.write_annotations('beats', 'qrs', [100 460 820], 'symbols', 'N');   % no header
ann = wfdbx.read_annotations('beats', 'qrs', 'sampto', '0:02', 'fs', 360);
fprintf('%d %d %g\n', ann.samples, ann.fs);           % 100 460 360
\end{lstlisting}

% ------------------------------------------------------------------------------------------
\subsection{Writing records}
\label{sec:matlab-writing}

\apientry{wfdbx.write\_record}{Write a header and a signal file}
\begin{lstlisting}[language=Matlab,style=signature]
wfdbx.write_record(record, signals, fs, 'fmt', 16, 'gains', [], 'baselines', [], ...
    'adczeros', 0, 'adcres', [], 'units', [], 'names', [], 'base_time', [], ...
    'comments', [], 'physical', false, 'samps_per_frame', [], 'metadata', [], ...
    'dir', [], 'legacy', false)
\end{lstlisting}
Writes signals as a new record: the header \file{RECORD.hea} and one signal file,
\file{RECORD.dat}, which holds all the signals (so they share one format). Existing files
of these names are replaced: both files are written to temporary files in the
destination folder (\file{.NAME.XXXXXXXX.tmp}, created with the permissions of the files
they replace), which replace the existing files only when both are complete, so that an
error leaves the existing record as it was, and no temporary file is left. The sampling
frequency, the gain and ADC resolution of
each signal, and the units of each calibrated signal must be given: there are no
defaults. The per-signal options (\code{'fmt'}, \code{'gains'}, \code{'baselines'},
\code{'adczeros'}, \code{'adcres'}, \code{'units'}, \code{'names'}) take one value for all
the signals, or a vector or cell array with one value per signal. All the arguments are
checked before any file is created or changed.
\begin{params}
\item[record] Record name (not empty), without the \code{.hea} extension (a trailing
  \code{.hea} is removed, as by the readers): a plain name
  (\code{'rec'}) or a relative path (\code{'data/rec'}), created in the output directory
  if one is set, otherwise in the current folder, missing folders included; or a full
  path, absolute or in the home directory (\code{'\textasciitilde/Documents/rec'}), used
  as given. With a path, the files are created in that folder, and the header names the
  record and its signal file without the path (\code{'rec'}, \code{'rec.dat'}), so that
  the record opens under any path that leads to it (in legacy mode, as in the original,
  the header gives the name with its path). A URL cannot be written.
\item[signals] The samples: a matrix (a 1-D or 2-D array) with one row per sample and one
  column per signal.
  In ADC units (the default), the values are integers within the range of a 32-bit
  integer; with \code{'physical', true}, they are in physical units.
\item[fs] The sampling frequency, in Hz: a finite number $> 0$.
\item[fmt] The storage format: 8, 16, 61, 80, 160, 212, 310, 311, 24, 32, or 508, 516 or
  524 (FLAC); one value, or one per signal, all equal (\cref{chap:formats}). Default
  \code{16}.
\item[gains] The gain of each signal, in ADC units per physical unit: a finite number;
  \code{0} marks an uncalibrated signal, which is allowed only for samples given in ADC
  units. Negative gains are allowed. Required (default \code{[]}: not given).
\item[baselines] The ADC value that corresponds to 0 physical units (a whole number within
  the range of a 32-bit integer; in a cell array, \code{[]} is the ADC zero). Default
  \code{[]}: the ADC zero of the signal.
\item[adczeros] The ADC value at the middle of the ADC range (a whole number within the
  range of a 32-bit integer; in a cell array, \code{[]} is 0). Default \code{0}.
\item[adcres] The ADC resolution, in bits (a whole number $\ge 1$, within the range of a
  32-bit integer). Required (default \code{[]}: not given).
\item[units] The physical units, such as \code{'mV'}: one string for all the signals, or a
  cell array with one per signal. Required for each signal whose gain is not 0. Default
  \code{[]}: not given.
\item[names] The signal descriptions, such as \code{'MLII'}: one string, or a cell array
  with one per signal. Default \code{[]}: no descriptions.
\item[base\_time] The start time, with an optional date: \code{'hh:mm:ss'} or
  \code{'hh:mm:ss dd/mm/yyyy'}, with optional decimals of the seconds, as
  \code{wfdbx.setbasetime} reads it. Default \code{[]}: no start time.
\item[comments] A cell array of info strings (or one string), written to the header as
  comment lines (\code{\#} followed by the string). Default \code{[]}: none.
\item[physical] \code{true} if \code{signals} are in physical units. They are converted to
  ADC units as $\mathrm{round}(x \cdot \mathit{gain}) + \mathit{baseline}$ (halves rounded
  away from zero); NaN becomes an invalid sample, and infinite values and values beyond
  the range of 32-bit ADC units are errors. Default \code{false}: ADC units.
\item[samps\_per\_frame] For a multi-frequency record, the number of samples per frame of
  each signal (integers $\ge 1$). \code{signals} then has one row per frame and
  \code{sum(samps\_per\_frame)} columns: the samples of signal 0, then those of signal 1,
  and so on. \code{fs} is then the frame frequency. Default \code{[]}: one sample per
  frame for every signal.
\item[metadata] Metadata entries (\cref{chap:metadata}) to write to the header after the
  comments, as \code{\# Key: value} (\code{\# Key:} for an empty value): an N-by-2 cell
  array with one entry per row, \code{\{KEY, VALUE; ...\}}. Keys and values follow the
  rules of \code{wfdbx.set\_metadata}; no two keys may be equal without regard to letter
  case, nor be the key of an entry in \code{'comments'}, and a value must not be
  \code{[]}. Default \code{[]}: none.
\item[dir] The folder (which must exist) in which to create the files, for this call
  only; \code{record} must then be a plain name, without a folder, and the record opens
  from that folder under its plain name. Default \code{[]}: the output directory.
\item[legacy] \code{true} to behave as the original WFDB package: the gain defaults to
  200, and the ADC resolution to one based on the format. Default \code{false}.
\end{params}
\returns{Nothing.}
\raises{Identifier \code{wfdbx:write\_record}: \code{record} is empty; \code{record}
has a character beyond 255 (``record name 'X' has a character that is not Latin-1 (one
byte per character)''; the same for a signal name, the units and a comment: ``signal
name 'X' \ldots'', ``units 'X' \ldots'', ``comment 'X' \ldots''; text is one byte per
character, so this can happen only in MATLAB, whose characters are 16-bit);
\code{signals} has more than two dimensions (``signals must be a 1-D or 2-D array''); a
\code{'samps\_per\_frame'} value is not an integer $\ge 1$ (``samps\_per\_frame must hold
integers >= 1''), or the number of columns of \code{signals} is not their sum (``signals
has N columns, but samps\_per\_frame adds up to M''); \code{signals} has no signals (``at
least one signal is required''); \code{fs} is not finite and $> 0$; a gain is not finite;
the gains, the ADC resolutions or the required units are not given; an ADC resolution, ADC
zero or baseline is not a whole number within the range of a 32-bit integer (text, NaN,
12.5: ``adcres must be whole numbers within the range of a 32-bit integer'', and the same
for \code{adczeros} and \code{baselines}); a per-signal option
does not have one value or one value per signal; a sample in ADC units is not an integer
or is outside the range of a 32-bit integer, or a sample in physical units is infinite or
its ADC value is outside that range; \code{'physical'} is \code{true} and a gain is 0;
\code{'base\_time'} cannot be parsed (``invalid base\_time 'T'''); the signals do not all
have the same format; an entry of \code{'metadata'} is not valid; \code{record} contains
a folder and \code{'dir'} is given, or \code{'dir'} does not exist; an option is not
valid; or the files cannot be created, or cannot be written completely, for example on a
full disk (``write error on record REC'', after the library's message; the existing files
are then left as they were).}
\begin{lstlisting}[language=Matlab,style=example]
% Two signals of random numbers, in millivolts, at 360 Hz: the record "rand"
% (rand.hea and rand.dat).
x = randn(3600, 2);
wfdbx.write_record('rand', x, 360, 'fmt', 16, 'physical', true, 'gains', 1000, ...
    'adcres', 16, 'units', 'mV', 'names', {'A', 'B'}, ...
    'base_time', '14:00:03 21/06/2024', 'comments', {'random numbers'}, ...
    'metadata', {'Generator', 'randn'});
rec = wfdbx.read_record('rand', 'physical', true);
fprintf('%d %d %s\n', size(rec.signals), rec.base_time);   % 3600 2 [14:00:03 21/06/2024]
disp(max(abs(rec.signals(:) - x(:))) <= 0.0005);           % 1: rounded to 1/1000 mV
disp(rec.comments');            % {'random numbers', ' Generator: randn'}

% An uncalibrated signal (gain 0, no units), in ADC units, in the folder "out".
mkdir('out');
wfdbx.write_record('counts', (0:999)', 100, 'gains', 0, 'adcres', 12, 'dir', 'out');
rec = wfdbx.read_record('out/counts');
disp(rec.signals(1:3)');        % 0 1 2
\end{lstlisting}
A multi-frequency record: signal \code{fast} has 2 samples per frame, signal
\code{slow} one; each row of \code{signals} is one frame, with the two samples of
\code{fast} and then the sample of \code{slow}:
\begin{lstlisting}[language=Matlab,style=example]
frames = 100;                                   % 1 second at 100 frames per second
fast = reshape(0:2 * frames - 1, 2, frames)';   % 200 samples per second
slow = repmat(500, frames, 1);                  % 100 samples per second
wfdbx.write_record('mf', [fast, slow], 100, 'gains', 200, 'adcres', 12, ...
    'units', 'mV', 'names', {'fast', 'slow'}, 'samps_per_frame', [2 1]);
rec = wfdbx.read_record('mf', 'highres', true);
fprintf('%g\n', rec.fs);        % 200
disp(rec.signals(1:3, :));      % 0 500; 1 500; 2 500
\end{lstlisting}

\apientry{wfdbx.write\_annotations}{Write an annotation file}
\begin{lstlisting}[language=Matlab,style=signature]
wfdbx.write_annotations(record, annotator, samples, 'symbols', [], 'codes', [], ...
    'subtypes', 0, 'chans', 0, 'nums', 0, 'aux', [], 'fs', [], 'dir', [], ...
    'sort', [], 'legacy', false)
\end{lstlisting}
Writes an annotation file, \file{RECORD.ANNOTATOR}, in the MIT format. An existing file of
that name is replaced: the new file is written to a temporary file in the same folder
(created with the permissions of the file it replaces), which replaces it only when it is
complete (and sorted), so that an error leaves it as it was. The record need not have a
header. The annotation types are
given as \code{'symbols'} or as \code{'codes'} (\cref{chap:formats}). Each per-annotation
option is one value for all the annotations, or a list with exactly one value per
annotation. If the annotations are not in time order, the file is sorted when it is
written, unless \code{'sort'} is \code{false}. All the arguments are checked before the
file is created or changed.
\begin{params}
\item[record] Record name, without the \code{.hea} extension (a trailing \code{.hea} is
  also accepted): a plain name or a relative path (\code{'data/100s'}) is created in the
  output directory if one is set, otherwise in the current folder; a full path, absolute
  or in the home directory (\code{'\textasciitilde/Documents/100s'}), is used as given.
  The file is \file{RECORD.ANNOTATOR} at that place (\file{data/100s.qrs}). A URL cannot
  be written.
\item[annotator] Annotator name: the suffix of the file, made of letters, digits,
  \code{\_}, \code{-} and \code{\textasciitilde}.
\item[samples] The annotation times, in samples: a vector (a 1-D array) of integers from
  $-2^{63}+1$ to $2^{63}-2$, one per annotation (the smallest and the largest 64-bit
  integers are reserved by the file format; doubles are exact up to $2^{53}$, and
  \code{int64} values are accepted).
\item[symbols] The mnemonics of the annotation types, such as \code{'N'} or \code{'V'}: a
  cell array with one per annotation, a \code{char} vector of one-character mnemonics
  (\code{'NNV'}), or one mnemonic for all. An unknown mnemonic becomes code 0
  (\code{NOTQRS}). Default \code{[]}: the types are given by \code{'codes'}.
\item[codes] The annotation codes (\code{wfdbx.NORMAL()}, \code{wfdbx.PVC()}, \ldots; 0
  to 58, as 59 to 63 are the control codes of the file format), instead of
  \code{'symbols'}: one per annotation, or one value for all. If both are given
  (as \code{wfdbx.read\_annotations} returns them), each symbol must be that of its code
  (an empty symbol is not checked), and the codes are written. Default \code{[]}.
\item[subtypes] The subtype of each annotation, stored as a signed byte (other values wrap
  around to $-128 \ldots 127$, as in C). Default \code{0}.
\item[chans] The channel number of each annotation, stored as an unsigned byte (other
  values wrap around to $0 \ldots 255$). Default \code{0}.
\item[nums] The annotator number of each annotation, stored as a signed byte. Default
  \code{0}.
\item[aux] The auxiliary text of each annotation: a cell array with one per annotation,
  with \code{[]} where there is none, or one \code{char} vector for all. Default
  \code{[]}: none.
\item[fs] If given (a finite number $> 0$), the frequency in which the times are
  counted, written to the file. Default \code{[]}: none is written.
\item[dir] The folder (which must exist) in which to create the file, for this call
  only; \code{record} must then be a plain name, without a folder: the file is
  \file{DIR/RECORD.ANNOTATOR}. Default \code{[]}: the output directory.
\item[sort] \code{true} to always sort the file, \code{false} to keep the given order
  (without the library's advice to run \code{sortann}), for this call only. Default
  \code{[]}: the environment variable \code{WFDBANNSORT} decides (sorting is on unless it
  is \code{0}).
\item[legacy] \code{true} to behave as the original WFDB package. Default \code{false}.
\end{params}
\returns{Nothing.}
\raises{Identifier \code{wfdbx:write\_annotations}: \code{samples} is a matrix or an
array of more dimensions (``samples must be a 1-D array''; the same for \code{'codes'},
\code{'subtypes'}, \code{'chans'}, \code{'nums'}, and for \code{'symbols'} given with
\code{'codes'}); neither \code{'symbols'} nor \code{'codes'} is given for one annotation
or more (``symbols or codes must be given''; with no annotations, an empty file is
written); a per-annotation
option given as a list does not have one value per annotation (``\ldots{} must have one
value per annotation''); a time is not an integer within the range of a 64-bit integer
(NaN, an infinite value or a fraction); a code, subtype, channel or annotator number is
NaN or infinite; \code{'fs'} is not finite and $> 0$; a time is $-2^{63}$ or, as an
\code{int64}, $2^{63}-1$ (``annotation time T cannot be written (annotation times must be
integers from -9223372036854775807 to 9223372036854775806)''); a code is not from 0 to 58
(``annotation code C is reserved (59 to 63 are the control codes of the annotation file
format)'', ``annotation code C is out of range (0 to 58)''); both \code{'symbols'} and
\code{'codes'} are given and disagree (``symbols and codes disagree at annotation K'', K
counting from 0); an aux text has a character beyond 255 (``aux text 'X' has a character
that is not Latin-1 (one byte per character)'', in MATLAB only); \code{record} contains a
folder and \code{'dir'} is given, or \code{'dir'} does not exist; an option is not valid;
or the file cannot be created, or cannot be written completely, or the library reports an
error while writing it (``write error on annotator ANN for record REC''; the existing file
is then left as it was).}
\begin{lstlisting}[language=Matlab,style=example]
% Three beats and a rhythm change, by symbol, with an aux string for the rhythm.
wfdbx.write_annotations('100s', 'test', [18 77 370 662], ...
    'symbols', {'+', 'N', 'V', 'N'}, 'aux', {'(N', [], [], []});
ann = wfdbx.read_annotations('100s', 'test');
fprintf('%s', ann.symbols{:}); fprintf(' %s\n', ann.aux{1});    % +NVN (N

% A copy of the reference annotations, as read_annotations returns them.
a = wfdbx.read_annotations('100s', 'atr');
wfdbx.write_annotations('100s', 'copy', a.samples, 'symbols', a.symbols, ...
    'codes', a.codes, 'subtypes', a.subtypes, 'chans', a.chans, 'nums', a.nums, ...
    'aux', a.aux);
b = wfdbx.read_annotations('100s', 'copy');
disp([isequal(b.samples, a.samples), isequal(b.symbols, a.symbols), ...
    isequal(b.aux, a.aux)]);    % 1 1 1

% Annotations out of time order are sorted when the file is written (the library
% reports it on the standard error).
wfdbx.write_annotations('100s', 'unsorted', [500 100], 'codes', wfdbx.NORMAL());
ann = wfdbx.read_annotations('100s', 'unsorted');
disp(ann.samples');             % 100 500
\end{lstlisting}

% ------------------------------------------------------------------------------------------
\subsection{Checking records}
\label{sec:matlab-checking}

\apientry{wfdbx.check\_record}{Check a record and its annotation files}
\begin{lstlisting}[language=Matlab,style=signature]
report = wfdbx.check_record(record, 'annotators', {}, 'input_dir', [], 'network', [])
\end{lstlisting}
Checks the header of a record, its signal specifications, every sample of its signal
files, each segment of a multi-segment record, and the annotation files named, and
reports each problem found. A problem in the record never throws an error: it is
reported, with a severity (\code{'error'}: the record is damaged or does not follow the
WFDB format, and programs will read wrong data or fail; \code{'warning'}: the record can
be read, but something is missing or unusual, and may give wrong results; \code{'info'}:
a fact worth knowing, such as invalid samples). The problems, their codes and their
messages are the same in the seven languages (\cref{tab:python-issue-codes}). The checks
run in normal mode, even if legacy mode is on, and the library prints no messages during
the call. They are similar to those of Benjamin Moody's \code{wfdbcheck}. Called without
an output argument, the function displays the lines of the report.
\begin{params}
\item[record] Record name, without the \code{.hea} extension (a trailing \code{.hea}
  is also accepted): a full path, absolute (\code{'/data/mitdb/100'}) or in the home
  directory (\code{'\textasciitilde/Documents/100s'}), used as given; a relative path
  (\code{'100s'}, \code{'data/100s'}), read from the input directory if one is set,
  otherwise from the current folder; or a URL
  (\code{'https://physionet.org/files/mitdb/1.0.0/100'}), read only if network access is
  on (\cref{sec:record-paths}).
\item[annotators] The annotator names whose files are checked: a cell array, such as
  \code{\{'atr'\}}, or one name as a \code{char} vector. Default \code{\{\}}: none.
\item[input\_dir] The input directory for this call only (optional), under which
  relative record names are read; it does not apply to full paths or URLs (as for
  \code{wfdbx.read\_record}). Default \code{[]}: the current setting.
\item[network] \code{true} to allow reading URLs for this call only, \code{false} to
  forbid it. Default \code{[]}: the current setting.
\end{params}
\returns{A structure (\cref{tab:matlab-checkreport}): \code{report.ok} is \code{true} if
no problem is an error, \code{report.issues} is a structure array of the problems
(\cref{tab:matlab-issue}), and \code{report.lines} is a cell column with one line per
problem, or the one line \code{no problems found}.}
\raises{Identifier \code{wfdbx:check\_record}, only for invalid arguments:
\code{record} is not a string or is empty, an annotator name is not a string, or an
option is not valid. A record that cannot be read is reported (code
\code{header\_unreadable}), not thrown.}
\begin{lstlisting}[language=Matlab,style=example]
report = wfdbx.check_record('100s', 'annotators', {'atr'});
disp(report.ok);                % 1: no errors
fprintf('%s\n', report.lines{:});
% warning: signal 0: the units are not specified
% warning: signal 1: the units are not specified

% A damaged record: one sample of a 12-bit signal is overwritten with 2048.
wfdbx.write_record('damaged', (0:99)', 100, 'gains', 100, 'adcres', 12, ...
    'units', 'mV', 'names', 'ramp');
fid = fopen('damaged.dat', 'r+b');
fseek(fid, 20, 'bof');          % sample 10 (2 bytes per sample in format 16)
fwrite(fid, [0 8], 'uint8');    % 0x0800 = 2048, little-endian
fclose(fid);
report = wfdbx.check_record('damaged');
disp(report.ok);                % 0
for k = 1:numel(report.issues)
    i = report.issues(k);
    fprintf('%s %s %s\n', i.severity, i.code, i.where);
end
% error sample_range signal 0
% error checksum signal 0
disp(report.lines{1});
% error: signal 0: 1 sample is outside the ADC range (-2048 to 2047): sample 10 (value 2048)
\end{lstlisting}

% ------------------------------------------------------------------------------------------
\subsection{Header metadata}
\label{sec:matlab-metadata}

These functions read and change the metadata of a header: its info strings of the form
\code{\# Key: value}, such as \code{\# Age: 63} (\cref{chap:metadata}). In short: keys are
compared without regard to letter case (A--Z and a--z only); a key must not be empty and
must not contain \code{:}, \code{\#}, a line break or a NUL character, and a value must
not contain a line break or a NUL character; spaces, tabs, VT, FF and CR at either end of
keys and values are removed; text is one byte per character, and keys and values may be
\code{char} vectors, MATLAB strings or \code{uint8} vectors of bytes. The entries are an
N-by-2 cell array of \code{char} row vectors, one entry per row, \code{\{KEY, VALUE;
...\}}, in file order (\code{cell(0, 2)} if there is none); a structure cannot hold keys
with spaces, such as \code{Date of birth}, so none is used. The info strings are the
lines that begin with \code{\#} after the record line, which is found as the header
parser of the library finds it, in every mode: the first line that holds something
other than white space (spaces, tabs, CRs and LFs, up to a NUL character) and does not
begin, after white space, with \code{\#}. Comment lines, and empty or blank lines,
before the record line are never read or changed. A change rewrites the header
in one step, through a temporary file, keeping every other byte of it; all the arguments
are checked first, so an error leaves the header unchanged. Only the \file{.hea} file is
read (not a \file{.info} file); for a multi-segment record, it is the master header,
unless a segment is named. Metadata often holds information about patients: WFDB-x stores
it as plain text, exactly as given, and does not de-identify it.

Each function has two error identifiers (\cref{sec:matlab-warnings}): \code{wfdbx:NAME}
for an invalid argument, and \code{wfdbx:NAME:file} for a header that cannot be read or
written, or a \code{'dir'} that does not exist (NAME is the name of the function, such as
\code{set\_metadata}).

\begin{phibox}
A header can hold identifying data about a patient (a name, a date of birth, a medical
record number). The metadata functions read and write such data as given. Records that
are shared must be de-identified first.
\end{phibox}

\begin{note}[Permissions of a changed header]
The temporary file is renamed over the header (\code{rename} in Octave, \code{movefile}
in MATLAB), so the header keeps its permissions only if the temporary file has them. The
temporary file is created with the permissions of the header, whatever the name of the
header (parentheses, spaces and any other characters included), and no shell command is
run on the name: Octave, which has no \code{chmod} function, creates it under a
\code{umask} that leaves exactly the read and write permissions of the header, so that a
private (0600) header stays private, and is never readable by others while it is
rewritten; MATLAB gives it the read, write and execute permissions of the header with
\code{fileattrib}. In Octave, execute permissions (which a header does not normally have)
are not kept. The files that \code{wfdbx.write\_record} and \code{wfdbx.write\_annotations}
replace keep their permissions in the same way.
\end{note}

\apientry{wfdbx.read\_metadata}{Read the metadata of a record}
\begin{lstlisting}[language=Matlab,style=signature]
M = wfdbx.read_metadata(record, 'input_dir', [], 'network', [])
\end{lstlisting}
Reads the metadata entries of the header of a record, in file order. An info string is an
entry if it contains a colon: the key is the text before the first colon, the value the
rest of the line, both without spaces at either end. Info strings without a colon, and
those whose key is empty or contains \code{\#}, are free comments, which are not returned.
\begin{params}
\item[record] Record name, without the \code{.hea} extension (a trailing \code{.hea}
  is also accepted): a full path, absolute (\code{'/data/mitdb/100'}) or in the home
  directory (\code{'\textasciitilde/Documents/100s'}), used as given; a relative path
  (\code{'100s'}, \code{'data/100s'}), read from the input directory if one is set,
  otherwise from the current folder; or a URL
  (\code{'https://physionet.org/files/mitdb/1.0.0/100'}), read only if network access is
  on (\cref{sec:record-paths}).
\item[input\_dir] The input directory for this call only (optional), under which
  relative record names are read; it does not apply to full paths or URLs (as for
  \code{wfdbx.read\_record}). Default \code{[]}: the current setting.
\item[network] \code{true} to allow reading URLs for this call only, \code{false} to
  forbid it. Default \code{[]}: the current setting.
\end{params}
\returns{An N-by-2 cell array of \code{char} row vectors, \code{\{KEY, VALUE; ...\}}, in
file order. A key that appears more than once (ignoring letter case) is returned once,
with its first value, and the warning \code{wfdbx:duplicate\_key} is issued
(\cref{sec:matlab-warnings}).}
\raises{Identifier \code{wfdbx:read\_metadata}: \code{record} is empty or \code{'-'} (the
standard input), or an option is not valid. Identifier \code{wfdbx:read\_metadata:file}:
the header cannot be found or read (``can't read header for record REC'').}
\begin{lstlisting}[language=Matlab,style=example]
wfdbx.write_record('pat', zeros(360, 1), 360, 'gains', 200, 'adcres', 12, ...
    'units', 'mV', 'names', 'ECG', 'comments', {'recorded at rest'}, ...
    'metadata', {'Age', '63'; 'Sex', 'F'});
M = wfdbx.read_metadata('pat');
disp(size(M));                  % 2 2
M = M';
fprintf('%s: %s\n', M{:});      % Age: 63, then Sex: F
hdr = wfdbx.read_header('pat');
disp(hdr.comments');
% {'recorded at rest', ' Age: 63', ' Sex: F'}: the free comment is not an entry
\end{lstlisting}

\apientry{wfdbx.get\_metadata}{Return the value of a metadata entry}
\begin{lstlisting}[language=Matlab,style=signature]
value = wfdbx.get_metadata(record, key, 'input_dir', [], 'network', [])
\end{lstlisting}
Returns the value of one metadata entry of the header of a record.
\begin{params}
\item[record] Record name (as for \code{wfdbx.read\_metadata}).
\item[key] The key, such as \code{'Age'}. Letter case is ignored, and spaces at either end
  are removed.
\item[input\_dir] The input directory for this call only (optional), under which a
  relative record name is read (as for \code{wfdbx.read\_metadata}). Default \code{[]}:
  the current setting.
\item[network] \code{true} to allow reading URLs for this call only, \code{false} to
  forbid it. Default \code{[]}: the current setting.
\end{params}
\returns{The value, a \code{char} row vector (empty for an entry with an empty value), or
\code{[]} if the header has no entry for \code{key} (\code{ischar} tells them apart). If
it has more than one, the first is returned, and the warning \code{wfdbx:duplicate\_key}
is issued.}
\raises{Identifier \code{wfdbx:get\_metadata}: \code{record} is empty or \code{'-'},
\code{key} is missing, not a string (or \code{uint8} vector) or not a valid key, or an
option is not valid. Identifier \code{wfdbx:get\_metadata:file}: the header cannot be
found or read.}
\begin{lstlisting}[language=Matlab,style=example]
wfdbx.write_record('pat', zeros(360, 1), 360, 'gains', 200, 'adcres', 12, ...
    'units', 'mV', 'metadata', {'Age', '63'; 'Sex', 'F'});
disp(wfdbx.get_metadata('pat', 'age'));         % 63
v = wfdbx.get_metadata('pat', 'Height');
disp([isempty(v), ischar(v)]);                  % 1 0: there is no such entry
\end{lstlisting}

\apientry{wfdbx.set\_metadata}{Set the value of a metadata entry}
\begin{lstlisting}[language=Matlab,style=signature]
wfdbx.set_metadata(record, key, value, 'input_dir', [], 'network', [], 'dir', [])
\end{lstlisting}
Sets the value of a metadata entry in the header of a record. If the header has an entry
for \code{key} (ignoring letter case), its line is replaced by \code{\# Key: value}, at
the same position and with the same line ending, with the spelling of \code{key} given
here (even if the value is the same); otherwise the entry is added after the last info
string of the header (or, if it has none, after its last line that is not blank), with
the line ending of the record line. The entry is written as \code{\# Key: value}, or
\code{\# Key:} if the value is empty. Every other line is kept, byte for byte, including
its line ending. The header is changed where it was read (not in the output directory),
unless \code{'dir'} is given.
\begin{params}
\item[record] Record name (as for \code{wfdbx.read\_metadata}). The header is changed
  where it was read, unless \code{'dir'} is given; a header read from a URL can only be
  written elsewhere, with \code{'dir'}.
\item[key] The key, such as \code{'Age'} or \code{'Date of birth'}. Spaces at either end
  are removed. It must not be empty, and must not contain \code{:}, \code{\#}, a line
  break or a NUL character.
\item[value] The value: text, which may be a sentence. Spaces at either end are removed.
  It must not contain a line break or a NUL character; it may be empty (\code{''}).
\item[input\_dir] The input directory in which the header is read, for this call only.
  Default \code{[]}: the current setting.
\item[network] \code{true} to allow reading the header from a URL, for this call only,
  \code{false} to forbid it. Default \code{[]}: the current setting.
\item[dir] A folder, which must exist, in which to write the changed header, under the
  same file name, instead of changing the header itself (which is then unchanged); the
  signal files are not copied. An unchanged header is also written there. Default
  \code{[]} (or \code{''}): the header itself is changed.
\end{params}
\returns{Nothing.}
\raises{Identifier \code{wfdbx:set\_metadata}: \code{key} or \code{value} is missing, not
a string (or \code{uint8} vector) or not valid; the header has more than one entry for
\code{key} (``\ldots{} remove them first''); it has no record line; it was read from a URL
and \code{'dir'} is not given; \code{record} is empty or \code{'-'}; \code{'dir'} is not a
string; or an option is not valid. Identifier \code{wfdbx:set\_metadata:file}: the header
cannot be read, or the changed header cannot be written (``can't write header for record
REC''); or \code{'dir'} does not exist (``directory DIR does not exist'').}
\begin{lstlisting}[language=Matlab,style=example]
wfdbx.write_record('pat', zeros(360, 1), 360, 'gains', 200, 'adcres', 12, ...
    'units', 'mV', 'metadata', {'Age', '63'; 'Sex', 'F'});
wfdbx.set_metadata('pat', 'AGE', '64');                 % replaces "# Age: 63"
wfdbx.set_metadata('pat', 'Diagnosis', 'sinus rhythm, no ectopy');   % a new entry
M = wfdbx.read_metadata('pat')';
fprintf('%s: %s\n', M{:});
% AGE: 64
% Sex: F
% Diagnosis: sinus rhythm, no ectopy

% Write the changed header to another folder; pat.hea is left unchanged.
mkdir('changed');
wfdbx.set_metadata('pat', 'Sex', 'M', 'dir', 'changed');
fprintf('%s %s\n', wfdbx.get_metadata('pat', 'Sex'), ...
    wfdbx.get_metadata('changed/pat', 'Sex'));         % F M
\end{lstlisting}

\apientry{wfdbx.remove\_metadata}{Remove a metadata entry}
\begin{lstlisting}[language=Matlab,style=signature]
n = wfdbx.remove_metadata(record, key, 'input_dir', [], 'network', [], 'dir', [])
\end{lstlisting}
Removes every entry for \code{key} (ignoring letter case) from the header of a record,
line ending included, so it also repairs a header that has more than one. Other lines are
kept, byte for byte. Nothing is written if the header has no entry for \code{key}, unless
\code{'dir'} is given.
\begin{params}
\item[record] Record name (as for \code{wfdbx.set\_metadata}).
\item[key] The key (as for \code{wfdbx.set\_metadata}).
\item[input\_dir] The input directory in which the header is read, for this call only.
  Default \code{[]}: the current setting.
\item[network] \code{true} to allow reading the header from a URL, for this call only,
  \code{false} to forbid it. Default \code{[]}: the current setting.
\item[dir] A folder, which must exist, in which to write the changed header instead of
  changing the header itself. Default \code{[]}: the header itself is changed.
\end{params}
\returns{The number of lines removed (0 if the header has no entry for \code{key}).}
\raises{As \code{wfdbx.set\_metadata}, with the identifiers
\code{wfdbx:remove\_metadata} and \code{wfdbx:remove\_metadata:file}, except that more
than one entry for \code{key} is not an error.}
\begin{lstlisting}[language=Matlab,style=example]
wfdbx.write_record('pat', zeros(360, 1), 360, 'gains', 200, 'adcres', 12, ...
    'units', 'mV', 'metadata', {'Age', '63'; 'Sex', 'F'});
disp(wfdbx.remove_metadata('pat', 'sex'));      % 1
disp(wfdbx.remove_metadata('pat', 'Sex'));      % 0: there is no such entry now
disp(wfdbx.read_metadata('pat'));               % {'Age', '63'}
\end{lstlisting}

\apientry{wfdbx.update\_metadata}{Make several changes to the metadata}
\begin{lstlisting}[language=Matlab,style=signature]
wfdbx.update_metadata(record, changes, 'input_dir', [], 'network', [], 'dir', [])
\end{lstlisting}
Makes several changes to the metadata of a record, in one write of the header. Each
change is applied in order, as \code{wfdbx.set\_metadata} does, or, for a value of
\code{[]}, as \code{wfdbx.remove\_metadata} does. All the changes are checked first: if one
is not valid, none is made.
\begin{params}
\item[record] Record name (as for \code{wfdbx.set\_metadata}).
\item[changes] An N-by-2 cell array with one change per row, \code{\{KEY, VALUE; ...\}};
  a value of \code{[]} (an empty double) removes the entry, and \code{''} sets an empty
  value. No two keys may be equal without regard to letter case. An empty cell array
  (\code{\{\}} or \code{cell(0, 2)}) changes nothing.
\item[input\_dir] The input directory in which the header is read, for this call only.
  Default \code{[]}: the current setting.
\item[network] \code{true} to allow reading the header from a URL, for this call only,
  \code{false} to forbid it. Default \code{[]}: the current setting.
\item[dir] A folder, which must exist, in which to write the changed header instead of
  changing the header itself. Default \code{[]}: the header itself is changed.
\end{params}
\returns{Nothing.}
\raises{As \code{wfdbx.set\_metadata}, with the identifiers
\code{wfdbx:update\_metadata} and \code{wfdbx:update\_metadata:file}; also an argument
error if \code{changes} is not an N-by-2 cell array, or if two of its keys differ only in
letter case.}
\begin{lstlisting}[language=Matlab,style=example]
wfdbx.write_record('pat', zeros(360, 1), 360, 'gains', 200, 'adcres', 12, ...
    'units', 'mV', 'metadata', {'Age', '63'; 'Sex', 'F'});
wfdbx.update_metadata('pat', {'Age', '64'; 'Sex', []; 'Dx', '164889003'});
M = wfdbx.read_metadata('pat')';
fprintf('%s: %s\n', M{:});      % Age: 64, then Dx: 164889003
try
    wfdbx.update_metadata('pat', {'Dx', '426783006'; 'dx', []});
catch err
    disp(err.identifier);       % wfdbx:update_metadata
    disp(err.message);          % metadata key 'dx' is given more than once (...)
end
disp(wfdbx.get_metadata('pat', 'Dx'));          % 164889003: nothing was changed
\end{lstlisting}

% ------------------------------------------------------------------------------------------
\subsection{The result structures}
\label{sec:matlab-results}

The readers return structures, whose fields are listed in
\cref{tab:matlab-record,tab:matlab-header,tab:matlab-annotations};
\code{wfdbx.check\_record} returns a structure that holds a structure array of issues
(\cref{tab:matlab-checkreport,tab:matlab-issue}). The signals of a record are described by
\code{wfdbx.siginfo} structures (\cref{tab:matlab-siginfo}), which are those of the WFDB
library (\cref{sec:matlab-library}). The fields are ordinary structure fields, which may
be changed.

\begin{xltabular}{\linewidth}{@{}lL@{}}
\caption{The fields of the result of \code{wfdbx.read\_record}}
\label{tab:matlab-record} \\
\toprule
\tabhead Field & \tabhead Contents \\
\midrule
\endfirsthead
\toprule
\tabhead Field & \tabhead Contents \\
\midrule
\endhead
\bottomrule
\endfoot
\code{name} & the record name, as given to \code{wfdbx.read\_record} \\
\code{fs} & the sampling frequency in Hz, or 0 if the header gives none; with
  \code{'highres', true}, the highest sampling frequency of a multi-frequency record \\
\code{signals} & a matrix of doubles with one row per sample and one column per signal
  read: ADC units, or physical units with \code{'physical', true}; invalid samples are
  \code{wfdbx.WFDB\_INVALID\_SAMPLE()} ($-32768$), or NaN in physical units \\
\code{siginfo} & a 1-by-N \code{wfdbx.siginfo} structure array, one per signal read, in
  the order of the columns; \code{desc} has no white space at either end, and
  \code{adcres} is 0 if the header does not give it \\
\code{names} & a 1-by-N cell array of the signal descriptions, such as \code{'MLII'},
  without the white space that a header may put at either end of them \\
\code{units} & a 1-by-N cell array of the physical units, with \code{[]} where the header
  gives none \\
\code{gains} & a row vector of the gains, in ADC units per physical unit; 0 means that the
  signal is uncalibrated \\
\code{baselines} & a row vector of the ADC values that correspond to 0 physical units \\
\code{base\_time} & the start time and date as text, to the second, such as
  \code{'[14:00:03 21/06/2024]'} (\code{'[14:00:03]'} without a date), or \code{[]} if
  the header gives none \\
\code{base\_time\_ms} & the start time in milliseconds after midnight, such as
  \code{50403500} for 14:00:03.500, or \code{[]} if \code{base\_time} is \code{[]} \\
\code{base\_date} & the start date as \code{[year month day]}, or \code{[]} if the header
  gives none \\
\code{sampfrom} & the number of the first sample returned \\
\code{comments} & a cell column of the info strings of the header (without the
  \code{\#}): its lines that begin with \code{\#} after the record line, found as the
  header parser finds it (an empty or blank first line is not the record line; in legacy
  mode, as the C function \code{getinfo} finds it: the first line that does not begin
  with \code{\#}) \\
\code{checksum\_ok} & \code{false} if the samples do not match the checksums of the
  header; checksums are checked only when a read reaches the end of the record, so a
  shorter read gives \code{true} \\
\end{xltabular}

\begin{xltabular}{\linewidth}{@{}lL@{}}
\caption{The fields of the result of \code{wfdbx.read\_header}}
\label{tab:matlab-header} \\
\toprule
\tabhead Field & \tabhead Contents \\
\midrule
\endfirsthead
\toprule
\tabhead Field & \tabhead Contents \\
\midrule
\endhead
\bottomrule
\endfoot
\code{name} & the record name, as given to \code{wfdbx.read\_header} \\
\code{fs} & the sampling frequency in Hz, or 0 if the header gives none \\
\code{counter\_freq} & the counter frequency, as \code{wfdbx.getcfreq()} returns it (the
  sampling frequency if the header gives none) \\
\code{base\_counter} & the base counter value, as \code{wfdbx.getbasecount()} returns
  it \\
\code{nsamples} & the number of samples per signal, or 0 if the header does not give it \\
\code{nsig} & the number of signals \\
\code{siginfo} & a 1-by-N \code{wfdbx.siginfo} structure array, one per signal;
  \code{desc} has no white space at either end, and \code{adcres} is 0 if the header does
  not give it \\
\code{base\_time}, \code{base\_time\_ms}, \code{base\_date} & the start time and date, as
  for \code{wfdbx.read\_record} \\
\code{comments} & a cell column of the info strings of the header (without the
  \code{\#}), as for \code{wfdbx.read\_record} \\
\code{metadata} & the metadata entries of the header, an N-by-2 cell array
  \code{\{KEY, VALUE; ...\}} in file order, as \code{wfdbx.read\_metadata} returns them:
  from the \file{.hea} file only; \code{cell(0, 2)}, without an error, for an EDF record
  and for a header read from the standard input \\
\end{xltabular}

\begin{xltabular}{\linewidth}{@{}lL@{}}
\caption{The fields of the result of \code{wfdbx.read\_annotations}: one row per
  annotation, in file order}
\label{tab:matlab-annotations} \\
\toprule
\tabhead Field & \tabhead Contents \\
\midrule
\endfirsthead
\toprule
\tabhead Field & \tabhead Contents \\
\midrule
\endhead
\bottomrule
\endfoot
\code{record} & the record name, as given \\
\code{annotator} & the annotator name, such as \code{'atr'} \\
\code{samples} & a column of the annotation times, in samples \\
\code{codes} & a column of the annotation codes, such as \code{wfdbx.NORMAL()} or
  \code{wfdbx.PVC()} \\
\code{symbols} & a cell column of the mnemonics of the codes, such as \code{'N'},
  \code{'V'} or \code{'+'} \\
\code{subtypes} & a column of the subtypes \\
\code{chans} & a column of the channel numbers \\
\code{nums} & a column of the annotator numbers \\
\code{aux} & a cell column of the auxiliary texts (\code{char} vectors), with \code{[]}
  where there is none \\
\code{fs} & the frequency in which \code{samples} are counted, normally the sampling
  frequency of the record; 0 if neither the record, nor the file, nor the option
  \code{'fs'} of \code{wfdbx.read\_annotations} gives one \\
\end{xltabular}

\begin{table}[htbp]
\centering
\caption{The fields of the result of \code{wfdbx.check\_record}}
\label{tab:matlab-checkreport}
\begin{tabularx}{\linewidth}{@{}lL@{}}
\toprule
\tabhead Field & \tabhead Contents \\
\midrule
\code{record} & the record name, as given \\
\code{ok} & \code{true} if no issue has the severity \code{'error'} \\
\code{issues} & an N-by-1 structure array of the issues (\cref{tab:matlab-issue}), in the
  order of the checks \\
\code{lines} & the text form of the report, a cell column: one line per issue,
  \code{'SEVERITY: WHERE: MESSAGE'} (\code{'SEVERITY: MESSAGE'} when \code{where} is
  empty), or the one line \code{'no problems found'} \\
\bottomrule
\end{tabularx}
\end{table}

\begin{table}[htbp]
\centering
\caption{The fields of an issue, a problem found by \code{wfdbx.check\_record}}
\label{tab:matlab-issue}
\begin{tabularx}{\linewidth}{@{}lL@{}}
\toprule
\tabhead Field & \tabhead Contents \\
\midrule
\code{severity} & \code{'error'}, \code{'warning'} or \code{'info'} \\
\code{code} & a stable identifier of the kind of problem, such as \code{'checksum'} or
  \code{'gain\_missing'} (\cref{tab:python-issue-codes}) \\
\code{where} & where the problem is: \code{''} (the record), \code{'signal N'},
  \code{'segment SEG'}, \code{'segment SEG, signal N'} or \code{'annotator ANN'} \\
\code{message} & the description of the problem \\
\bottomrule
\end{tabularx}
\end{table}

\begin{xltabular}{\linewidth}{@{}lL@{}}
\caption{The fields of a \code{wfdbx.siginfo} structure, the specification of a signal}
\label{tab:matlab-siginfo} \\
\toprule
\tabhead Field & \tabhead Contents \\
\midrule
\endfirsthead
\toprule
\tabhead Field & \tabhead Contents \\
\midrule
\endhead
\bottomrule
\endfoot
\code{fname} & the name of the signal file, as the header gives it \\
\code{desc} & the signal description, such as \code{'MLII'}, or \code{[]} \\
\code{units} & the physical units, such as \code{'mV'}, or \code{[]} if not specified \\
\code{gain} & the gain, in ADC units per physical unit (0: uncalibrated) \\
\code{initval} & the value of the first sample \\
\code{group} & the signal group number, from 0 (the signals of one group share one
  file) \\
\code{fmt} & the storage format, such as \code{16}, \code{212} or \code{516} \\
\code{spf} & the number of samples per frame (more than 1 for the faster signals of a
  multi-frequency record) \\
\code{bsize} & the block size, in bytes (0 for ordinary files) \\
\code{adcres} & the ADC resolution, in bits (0: not specified) \\
\code{adczero} & the ADC output for an input of 0\,V \\
\code{baseline} & the ADC value that corresponds to 0 physical units \\
\code{nsamp} & the number of samples (0 if unknown) \\
\code{cksum} & the 16-bit checksum of the samples \\
\end{xltabular}

% ------------------------------------------------------------------------------------------
\subsection{Warnings and errors}
\label{sec:matlab-warnings}

\paragraph{Errors.} The high-level functions report errors with \code{error}, with a
message and an identifier (\cref{tab:matlab-errors}), so that \code{err.identifier} of a
\code{try}/\code{catch} tells which function failed, and, for the metadata functions,
whether an argument was invalid or a file could not be read or written. Each function
checks its arguments before it creates or changes a file, so an invalid argument leaves
the files as they were. After an error, as after a successful call, the settings of the
call (\code{'input\_dir'}, \code{'network'}, \code{'dir'}, \code{'legacy'}), the getvec
mode and the time resolution of output annotators are as they were before it
(\cref{sec:matlab-conventions}). The library may also print a message on the standard error when
it cannot open a file (as in \code{init: can't open header for record REC});
\code{wfdbx.wfdbquiet()} turns these messages off (\cref{sec:matlab-library}).

\begin{table}[htbp]
\centering
\caption{The error and warning identifiers of the high-level API}
\label{tab:matlab-errors}
\begin{tabularx}{\linewidth}{@{}lL@{}}
\toprule
\tabhead Identifier & \tabhead Issued when \\
\midrule
\code{wfdbx:NAME} & an error of the function \code{wfdbx.NAME}: for
  \code{read\_record}, \code{read\_header}, \code{read\_annotations},
  \code{write\_record}, \code{write\_annotations} and \code{check\_record}, any error (a
  file that cannot be found, read, created or written completely; an invalid option or
  argument; a value that the operation needs and the record does not give, such as a
  time in seconds without a sampling frequency); for the five metadata functions, an
  invalid argument: an invalid key or value, a record name that is empty or \code{'-'}, a
  key that the header has more than once (\code{set\_metadata}, \code{update\_metadata}),
  a header without a record line, a header read from a URL to be changed without
  \code{'dir'}, or an invalid option \\
\code{wfdbx:NAME:file} & a metadata function cannot read or write a file: a header that
  cannot be found or read, a changed header that cannot be written, or a \code{'dir'}
  that does not exist \\
\code{wfdbx:checksum} & (a warning) the samples read by \code{wfdbx.read\_record} do not
  match the checksums of the header \\
\code{wfdbx:duplicate\_key} & (a warning) a header has more than one metadata entry for a
  key \\
\bottomrule
\end{tabularx}
\end{table}

\paragraph{The checksum warning.} When \code{wfdbx.read\_record} reads to the end of a
record whose samples do not match the checksums of its header, it returns all the samples
and issues the warning \code{wfdbx:checksum}: ``checksum error in record REC: the signal
data do not match the checksums in its header''. \code{'checksum\_warning', false} turns
the warning off; \code{rec.checksum\_ok} is \code{false} either way. This applies only
when every sample could be read: a signal file that ends before the length that its
header gives is an error (``can't read signal data of record REC (getvec status -3)''),
even though the checksum, verified at the frame in which the file ended, does not match
either; \code{wfdbx.check\_record} reports it as \code{length\_short}, not as a
checksum error.

\paragraph{The duplicate-key warning.} When a header holds two metadata entries whose keys
differ only in letter case (written by other software), \code{wfdbx.read\_metadata},
\code{wfdbx.get\_metadata} and \code{wfdbx.read\_header} use the first and issue the
warning \code{wfdbx:duplicate\_key}: ``header of record REC has more than one entry for
key KEY (lines L1, L2); the first is used''. \code{wfdbx.set\_metadata} and
\code{wfdbx.update\_metadata} then refuse to change that key (an argument error), and
\code{wfdbx.remove\_metadata} removes all of its entries.

Both are issued with \code{warning}, so \code{warning('off', 'wfdbx:checksum')} turns one
off, \code{lastwarn} returns the last one, and \code{warning('error', ...)} (MATLAB) makes
it an error, as for any other warning:

\begin{lstlisting}[language=Matlab,style=example]
% A record whose last sample is changed after it was written: the checksum fails.
wfdbx.write_record('ck', (0:99)', 100, 'gains', 100, 'adcres', 12, 'units', 'mV');
fid = fopen('ck.dat', 'r+b');
fseek(fid, -2, 'eof');
fwrite(fid, [0 0], 'uint8');
fclose(fid);
lastwarn('');
rec = wfdbx.read_record('ck');                  % prints the checksum warning
fprintf('%d %d %d\n', size(rec.signals), rec.checksum_ok);   % 100 1 0: all the samples
[msg, id] = lastwarn();
disp(id);                                       % wfdbx:checksum
rec = wfdbx.read_record('ck', 'checksum_warning', false);    % no warning

% Two entries for one key: the first is used, with a warning.
wfdbx.write_record('dup', zeros(10, 1), 100, 'gains', 0, 'adcres', 12, ...
    'comments', {' Age: 63', ' AGE: 64'});
disp(wfdbx.get_metadata('dup', 'age'));         % 63, with the warning
[msg, id] = lastwarn();
disp(id);                                       % wfdbx:duplicate_key
try
    wfdbx.set_metadata('dup', 'Age', '65');
catch err
    fprintf('%s: %s\n', err.identifier, err.message);
    % wfdbx:set_metadata: header of record dup has more than one entry for key Age
    % (lines 3, 4); remove them first
end
disp(wfdbx.remove_metadata('dup', 'age'));      % 2
\end{lstlisting}

\begin{lstlisting}[language=Matlab,style=example]
try
    wfdbx.read_record('nosuchrecord');
catch err
    fprintf('%s: %s\n', err.identifier, err.message);
    % wfdbx:read_record: can't read record nosuchrecord
end
try
    wfdbx.read_metadata('nosuchrecord');
catch err
    fprintf('%s: %s\n', err.identifier, err.message);
    % wfdbx:read_metadata:file: can't read header for record nosuchrecord
end
wfdbx.write_record('raw', (0:9)', 100, 'gains', 0, 'adcres', 12);   % uncalibrated
try
    wfdbx.read_record('raw', 'physical', true);
catch err
    fprintf('%s: %s\n', err.identifier, err.message);
    % wfdbx:read_record: signal 0 of record raw is uncalibrated
end
\end{lstlisting}

% ==========================================================================================
\section{The WFDB library functions in MATLAB and Octave}
\label{sec:matlab-library}

Every function of the WFDB library is a function of the package \code{wfdbx} under its C
name, with the same arguments, in the same order, and the same return values, so that
programs and documentation written for the C library carry over directly. The functions,
their arguments and their results are documented once, for all the languages, in
\cref{chap:library}; this section gives the rules by which the C declarations become
MATLAB functions (see also \cref{sec:library-matlab}). \code{help wfdbx.NAME} gives the
documentation of each function, including what WFDB-x changes in it.

\subsection{Conventions}

\paragraph{Names.} The functions keep their C names, with the prefix of the package:
\code{wfdbx.isigopen}, \code{wfdbx.getvec}, \ldots{} The constants (\code{WFDB\_READ},
\code{WFDB\_INVALID\_SAMPLE}, \code{WFDB\_HIGHRES}, \ldots) and the annotation codes
(\code{NORMAL}, \code{PVC}, \ldots, \cref{chap:formats}) are functions without arguments:
\code{wfdbx.WFDB\_READ()}, \code{wfdbx.NORMAL()}; \code{C = wfdbx.ecgcodes()} returns all
the annotation codes as the fields of one structure (\code{C.NORMAL}, \code{C.PVC},
\ldots), which is faster in loops. The calibration constants (\code{WFDB\_AC\_COUPLED},
\ldots) and \code{WFDB\_FMT\_LIST} are not defined; their values are given in
\cref{tab:library-constants}. The macros of \file{<wfdb/ecgmap.h>} are functions, under
both their macro names and the names of the functions behind them:
\code{wfdbx.isqrs} and \code{wfdbx.wfdb\_isqrs}, \code{wfdbx.map1} and
\code{wfdbx.wfdb\_map1}, and so on.

\paragraph{NULL is \code{[]}.} A C \code{NULL} string, as argument or result, is
\code{[]} (an empty double), and an empty C string is \code{''}; \code{ischar} tells them
apart. \code{wfdbx.sampfreq([])} returns the frequency of the open record,
\code{wfdbx.getinfo([])} continues with the info strings of the last record, and
\code{wfdbx.getinfo} returns \code{[]} after the last one.

\paragraph{Arrays and output arguments.} Where a C function fills an array or a structure
that the caller passes, the MATLAB function returns it as an additional output value, and
omits it from the arguments:
\begin{lstlisting}[language=Matlab]
[nsig, si] = wfdbx.isigopen(record, nsig)     % si: a struct array of siginfo
[nsig, si] = wfdbx.osigopen(record, nsig)
[stat, si] = wfdbx.wfdbinit(record, aiarray, nann, nsig)
[stat, v] = wfdbx.getvec(v0)                  % v0 (optional): the vector C leaves unchanged
[stat, v] = wfdbx.getframe(v0)
[stat, annot] = wfdbx.getann(an, annot0)      % annot0 (optional), as v0
[stat, cal] = wfdbx.getcal(desc, units)
[stat, si] = wfdbx.setheader(record, si, nsig)
segs = wfdbx.getseginfo()                     % a struct array of seginfo
[r, argv] = wfdbx.wfdbx_legacy_args(argv)     % argv: a cell array
\end{lstlisting}
With \code{nsig} $\le 0$, \code{wfdbx.isigopen} reads the header only, and \code{si}
describes \code{-nsig} signals. Arrays of structures passed in are struct arrays:
\code{wfdbx.annopen(record, wfdbx.anninfo('atr', wfdbx.WFDB\_READ()))},
\code{wfdbx.osigfopen(si, nsig)}. \code{wfdbx.setmsheader(record, segnames)} takes a cell
array of segment names.

\paragraph{Numbers and strings.} Samples, times, counts and frequencies are doubles,
exact for every 32-bit sample and for times up to $2^{53}$. Strings are \code{char} row
vectors whose characters are the bytes of the C strings (codes 0 to 255); record names
may also be MATLAB strings. The \code{aux} field of an annotation is \code{[]}, or a
\code{uint8} row vector whose first element is the length of the text that follows, as in
C: \code{char(a.aux(2:1 + a.aux(1)))} is the text.

\paragraph{Return values and errors.} The functions return what the C functions return: a
status (negative on errors), a number of signals, a frequency, a string or \code{[]}.
Errors are reported by these values and by a message on the standard error, as in C,
never by MATLAB errors. \code{wfdbx.wfdberror()} returns the last message;
\code{wfdbx.wfdbquiet()} turns the messages off and \code{wfdbx.wfdbverbose()} on again;
\code{wfdbx.wfdb\_error(fmt, ...)} reports a message as the library does, formatting as
C's \code{printf}.

\paragraph{State.} As in C, the library keeps its state in the session (in
\code{persistent} variables): the open record, its input signals, the output signals and
the annotators remain open between calls, until \code{wfdbx.wfdbquit()} closes them.
\code{wfdbx.wfdbquit()} leaves the input and output directories, the network setting,
legacy mode, the getvec mode and the time resolution of output annotators (which opening
an EDF record sets, as in C) as they are; \code{clear functions} (or \code{clear all}) resets the whole
state. The high-level functions of \cref{sec:matlab-api} (except the metadata functions)
call \code{wfdbx.wfdbquit()}, which closes anything opened with the library functions.

\paragraph{Additions.} \code{[v, stat, last] = wfdbx.getframes(n, minstat)} reads up to
\code{n} sample vectors at once (\code{n} may be \code{Inf}), as \code{n} calls of
\code{wfdbx.getvec} would, but decoding whole blocks of the signal files: \code{v} is a
matrix with one row per vector and one column per signal, \code{stat} the status of the
first vector that could not be read, or the number of signals if all \code{n} were read,
and \code{last} the vector filled by the call that stopped the reading
(\cref{sec:library-additions}). \code{stat = wfdbx.putframes(M)} writes the rows of
\code{M}, as \code{wfdbx.putvec} for each row would. Calling \code{wfdbx.getvec} sample by
sample is much slower in MATLAB and Octave than these functions, and than the high-level
functions.

\subsection{The structures}
\label{sec:matlab-structures}

The structures of the C library are MATLAB structures with the names of the fields of
the C structures (\cref{tab:matlab-structures}), made with C's zero values by these
functions:
\begin{lstlisting}[language=Matlab]
si = wfdbx.siginfo(n)            % 1-by-n struct array (n = 1 by default); spf is 1
s = wfdbx.seginfo(n)             % 1-by-n struct array (n = 1 by default)
ai = wfdbx.anninfo(name, stat)   % one structure
a = wfdbx.annotation(n)          % 1-by-n struct array (n = 1 by default)
cal = wfdbx.calinfo()            % one structure
\end{lstlisting}

\begin{xltabular}{\linewidth}{@{}lL@{}}
\caption{The C structures in MATLAB and Octave}
\label{tab:matlab-structures} \\
\toprule
\tabhead Constructor (C structure) & \tabhead Fields \\
\midrule
\endfirsthead
\toprule
\tabhead Constructor (C structure) & \tabhead Fields \\
\midrule
\endhead
\bottomrule
\endfoot
\code{wfdbx.siginfo} (\code{WFDB\_Siginfo}) & \code{fname}, \code{desc}, \code{units}
  (\code{[]} or \code{char}), \code{gain}, \code{initval}, \code{group}, \code{fmt},
  \code{spf}, \code{bsize}, \code{adcres}, \code{adczero}, \code{baseline}, \code{nsamp},
  \code{cksum} (\cref{tab:matlab-siginfo}) \\
\code{wfdbx.seginfo} (\code{WFDB\_Seginfo}) & \code{recname} (the name of the segment),
  \code{nsamp} (its number of samples), \code{samp0} (the number of its first sample in
  the record) \\
\code{wfdbx.anninfo} (\code{WFDB\_Anninfo}) & \code{name} (the annotator name, the suffix
  of the file), \code{stat} (\code{wfdbx.WFDB\_READ()} or \code{wfdbx.WFDB\_WRITE()} for
  the MIT format, \code{wfdbx.WFDB\_AHA\_READ()} or \code{wfdbx.WFDB\_AHA\_WRITE()} for the
  AHA format) \\
\code{wfdbx.annotation} (\code{WFDB\_Annotation}) & \code{time} (in sample intervals),
  \code{anntyp} (the annotation code), \code{subtyp}, \code{chan}, \code{num}, \code{aux}
  (\code{[]}, or a \code{uint8} row vector whose first element is the length of the text
  that follows) \\
\code{wfdbx.calinfo} (\code{WFDB\_Calinfo}) & \code{low}, \code{high} (the levels of the
  calibration pulse, in physical units), \code{scale} (physical units per cm),
  \code{sigtype} (matched against the beginning of a signal description), \code{units},
  \code{caltype} \\
\end{xltabular}

\subsection{The settings functions}
\label{sec:matlab-settings}

The functions of \cref{tab:matlab-settings} are additions of WFDB-x to the library: they
set where files are read and written, network access and legacy mode, for the high-level
functions, the library functions and the tools alike (\cref{chap:files,chap:behavior}),
and give the path under which a program opens a file that it reads or creates itself, as
the tools do (\code{wfdbx.wfdbx\_input\_path}, \code{wfdbx.wfdbx\_output\_path};
\cref{sec:tool-files}). Each \code{set} function returns the previous setting, so that it
can be restored. The C function \code{wfdbx.setwfdb(p)} sets the input directory too (\code{p} must name at most
one directory), and \code{wfdbx.getwfdb()} returns it; in legacy mode they set and return
the WFDB path, as in C.

\begin{xltabular}{\linewidth}{@{}lL@{}}
\caption{The settings functions}
\label{tab:matlab-settings} \\
\toprule
\tabhead Function & \tabhead Effect and result \\
\midrule
\endfirsthead
\toprule
\tabhead Function & \tabhead Effect and result \\
\midrule
\endhead
\bottomrule
\endfoot
\code{previous = wfdbx.wfdbx\_set\_input\_dir(d)} & sets the input directory: a local
  folder, a \code{\textasciitilde} path or a URL; \code{[]} or \code{''} restores the
  default (\code{WFDBX\_INPUT\_DIR}, or else the current folder). Returns the previous
  input directory, or \code{[]} if none was set \\
\code{d = wfdbx.wfdbx\_get\_input\_dir()} & returns the input directory, or \code{[]} if
  none is set \\
\code{previous = wfdbx.wfdbx\_set\_output\_dir(d)} & sets the output directory: a local
  folder or a \code{\textasciitilde} path (a URL is refused, with a message, and the
  setting is unchanged); \code{[]} or \code{''} restores the default
  (\code{WFDBX\_OUTPUT\_DIR}, or else the current folder). Returns the previous output
  directory, or \code{[]} \\
\code{d = wfdbx.wfdbx\_get\_output\_dir()} & returns the output directory, or \code{[]}
  if none is set \\
\code{previous = wfdbx.wfdbx\_set\_network(on)} & turns network access on (\code{true} or
  1) or off (\code{false} or 0); \code{[]} or a negative number restores the default (on
  if \code{WFDBX\_NETWORK} is a nonzero number). Returns the previous setting, 1 or 0 \\
\code{on = wfdbx.wfdbx\_get\_network()} & returns 1 if network access is on, 0
  otherwise \\
\code{previous = wfdbx.wfdbx\_set\_legacy(on)} & turns legacy mode on (\code{true} or 1)
  or off, for the whole library. Returns the previous setting, 1 or 0 \\
\code{on = wfdbx.wfdbx\_get\_legacy()} & returns 1 if legacy mode is on, 0 otherwise \\
\code{[r, argv] = wfdbx.wfdbx\_legacy\_args(argv)} & if \code{argv\{2\}} (\code{argv}: a
  cell row of the arguments, program name first) is \code{'-{}-legacy'}, removes it,
  turns legacy mode on and returns \code{r} = 1; otherwise returns \code{argv} unchanged
  and \code{r} = 0 \\
\code{r = wfdbx.wfdbx\_time\_needs\_freq(string)} & returns 1 if
  \code{wfdbx.strtim(string)} needs a sampling frequency, 0 otherwise \\
\code{n = wfdbx.wfdbx\_write\_errors()} & returns the number of write errors since the
  program started (output files that could not be written, flushed or closed) \\
\code{wfdbx.wfdbx\_note\_write\_error()} & counts a write error, for a program that
  writes a file itself \\
\code{p = wfdbx.wfdbx\_input\_path(name)} & returns the path under which a program opens
  a file that it reads itself: in normal mode, a relative name under the input directory,
  a \code{\textasciitilde} name expanded, an absolute name as given; in legacy mode, the
  name as given. \code{[]}, \code{''} and \code{'-'} are returned as given; \code{[]} for
  a URL, or a name that becomes one under the input directory \\
\code{p = wfdbx.wfdbx\_output\_path(name)} & returns the path under which a program
  creates a file that it writes itself: in normal mode, a relative name under the output
  directory (whose missing folders are then created), a \code{\textasciitilde} name
  expanded, an absolute name as given; in legacy mode, any name under the output
  directory, as a plain prefix. \code{[]}, \code{''} and \code{'-'} are returned as given;
  \code{[]} for a URL \\
\end{xltabular}

\subsection{Examples}

Reading the samples and the annotations of a record, as a C program does (the record
\code{100s} is read from the current folder):

\begin{lstlisting}[language=Matlab,style=example]
[nsig, si] = wfdbx.isigopen('100s', 2);         % open 2 signals
fprintf('%d %s %d %g\n', nsig, si(1).desc, si(1).fmt, si(1).gain);   % 2 MLII 212 200
fprintf('%g\n', wfdbx.sampfreq([]));            % 360: the frequency of the open record

wfdbx.isigsettime(wfdbx.strtim('0:30'));        % skip to 30 seconds
[stat, v] = wfdbx.getvec();                     % one sample per signal
fprintf('%d %d %g\n', v(1), v(2), wfdbx.aduphys(0, v(1)));      % 947 968 -0.385

[frames, stat] = wfdbx.getframes(360);          % the next second, as a matrix
fprintf('%d %d %d\n', size(frames), stat);      % 360 2 2

wfdbx.annopen('100s', wfdbx.anninfo('atr', wfdbx.WFDB_READ()));
n = 0;
[stat, a] = wfdbx.getann(0);
while stat == 0                                 % 0: an annotation was read; -1: the end
    if n < 2
        fprintf('%s %s\n', wfdbx.mstimstr(a.time), wfdbx.annstr(a.anntyp));
    end
    n = n + 1;
    [stat, a] = wfdbx.getann(0);
end
fprintf('%d annotations\n', n);                 % 75 annotations

wfdbx.wfdbquiet();                              % no messages on the standard error
stat = wfdbx.isigopen('nosuchrecord', 0);
fprintf('%d %s', stat, wfdbx.wfdberror());
% -1 init: can't open header for record nosuchrecord
wfdbx.wfdbverbose();
wfdbx.wfdbquit();                               % close everything
\end{lstlisting}

Writing a record and an annotation file, sample by sample and annotation by annotation:

\begin{lstlisting}[language=Matlab,style=example]
out = wfdbx.siginfo(1);
out.fname = 'ramp.dat';
out.desc = 'ramp';
out.units = 'mV';
out.gain = 200;
out.fmt = 16;
out.adcres = 12;
disp(wfdbx.osigfopen(out, 1));      % 1: one output signal
wfdbx.setsampfreq(250);
for k = 0:499
    wfdbx.putvec(k);
end
disp(wfdbx.newheader('ramp'));      % 0: ramp.hea is written
wfdbx.annopen('ramp', wfdbx.anninfo('qrs', wfdbx.WFDB_WRITE()));
a = wfdbx.annotation();
a.anntyp = wfdbx.NORMAL();
for t = [100 350]
    a.time = t;
    wfdbx.putann(0, a);
end
wfdbx.wfdbquit();                   % completes ramp.dat and ramp.qrs

ann = wfdbx.read_annotations('ramp', 'qrs');
disp(ann.samples');                 % 100 350
\end{lstlisting}

% ==========================================================================================
\section{The command-line tools in MATLAB and Octave}
\label{sec:matlab-tools}

The 33 command-line tools (\code{rdsamp}, \code{rdann}, \code{wrsamp}, \code{xform},
\ldots) are functions of the package \code{wfdbx.apps}, one file per tool in
\file{+wfdbx/+apps}. The tools, their options, their output and their exit status are
those of the C tools, and are documented in \cref{chap:tools}; \code{help
wfdbx.apps.TOOL} shows the usage of a tool.

\paragraph{From the shell, with Octave.} The Octave script \file{matlab/bin/wfdbx\_cli.m}
runs a tool: its first argument names the tool, and the others are the tool's arguments.
It adds \file{matlab} to the path, and exits with the tool's status. An alias makes the
command shorter:

\begin{lstlisting}[style=shell]
alias wfdbx='octave-cli --no-gui --norc -q /path/to/wfdb-x/matlab/bin/wfdbx_cli.m'
export WFDBX_INPUT_DIR=/path/to/wfdb-x/tests/data
wfdbx rdsamp -r 100s -t 10 -p     # the first 10 seconds of 100s, in physical units
wfdbx rdann -r 100s -a atr        # its annotations
wfdbx rdsamp --legacy -r 100s     # legacy mode: --legacy as the first argument
\end{lstlisting}

The script also takes these arguments, instead of a tool:

\begin{table}[htbp]
\centering
\caption{The arguments of the script \code{wfdbx\_cli.m} (as \cmd{wfdbx})}
\label{tab:matlab-wfdbx-options}
\begin{tabularx}{\linewidth}{@{}lL@{}}
\toprule
\tabhead Command & \tabhead Effect \\
\midrule
\cmd{wfdbx TOOL [ARGUMENTS ...]} & runs the tool \code{TOOL}, and exits with its
  status \\
\cmd{wfdbx -{}-list} & lists the names of the tools, one per line \\
\cmd{wfdbx -{}-version} & prints the version, \code{WFDB-x 1.2.2} \\
\cmd{wfdbx -{}-help}, \cmd{wfdbx -h} & prints the usage \\
\bottomrule
\end{tabularx}
\end{table}

Without arguments, the script prints the usage and exits with status 1; an unknown tool
is an error (``wfdbx: unknown tool 'TOOL' (try 'wfdbx -{}-list')'', status 1), and a
MATLAB error in a tool is reported as \code{wfdbx: error: MESSAGE} (status 1). When an
output file cannot be written completely (a full disk, for example), or the standard
output cannot be written (``TOOL: write error on standard output''), a tool exits with
status 1; in legacy mode, it keeps the original status, as in C. Octave starts in a fraction of a second, so
the script suits everyday use.

\paragraph{From the shell, with MATLAB.} The shell script
\file{matlab/bin/wfdbx\_matlab.sh} runs a tool with \cmd{matlab -batch}, or the program
named by the environment variable \code{WFDBX\_MATLAB}. It saves its standard input to a
temporary file (named by the environment variable \code{WFDBX\_STDIN}) before it starts
MATLAB, so a script whose standard input stays open runs it with \cmd{< /dev/null}, and it
passes the arguments in hexadecimal, so that any bytes arrive unchanged. MATLAB starts
slowly, so the script is mainly useful for testing.

\begin{lstlisting}[style=shell]
export WFDBX_INPUT_DIR=/path/to/wfdb-x/tests/data
/path/to/wfdb-x/matlab/bin/wfdbx_matlab.sh rdsamp -r 100s -t 10 -p
\end{lstlisting}

\paragraph{In a session.} Each tool is a function that takes the arguments as a cell
array (without the program name), writes to the standard output and error, and returns
the exit status:
\begin{lstlisting}[language=Matlab]
status = wfdbx.apps.rdsamp({'-r', '100s', '-t', 's2'})    % the first 2 samples
\end{lstlisting}
A call of \code{exit} in a C tool ends the function, not the session.

\paragraph{Environment variables.} The tools, like the rest of the package, use the
environment variables of \cref{tab:matlab-environment}, also for the files that they read
or create themselves (\cref{sec:tool-files}): the \opt{-i} file of \cmd{wrsamp},
\cmd{edf2mit}, \cmd{wav2mit}, \cmd{a2m} and \cmd{ad2m}, the input files of
\cmd{ahaecg2mit} and the \opt{-S} script of \cmd{xform} are read from the input directory,
and the WAV file of \cmd{mit2wav}, the EDF file of \cmd{mit2edf}, the master header of
\cmd{snip -m} and the \file{.des}, \file{.ecg}, \file{.key} and \file{.txt} files of
\cmd{parsescp} are created in the output directory. Such a file cannot be read from a URL.
In legacy mode, the files that a tool reads itself are opened as named, and the output
directory, if set, is a plain prefix of every output name. In Octave and in MATLAB,
\code{getenv} returns \code{''} for an unset and an empty variable alike, so an empty
variable counts as unset.

\begin{table}[htbp]
\centering
\caption{The environment variables used by the MATLAB/Octave package}
\label{tab:matlab-environment}
\begin{tabularx}{\linewidth}{@{}lL@{}}
\toprule
\tabhead Variable & \tabhead Meaning \\
\midrule
\code{WFDBX\_INPUT\_DIR} & the input directory (a folder, a \code{\textasciitilde} path or
  a URL), when \code{wfdbx.wfdbx\_set\_input\_dir} does not set one \\
\code{WFDBX\_OUTPUT\_DIR} & the output directory, when \code{wfdbx.wfdbx\_set\_output\_dir}
  does not set one \\
\code{WFDBX\_NETWORK} & network access is on if it is a nonzero number, when
  \code{wfdbx.wfdbx\_set\_network} does not set it \\
\code{WFDBANNSORT} & \code{0} to keep annotation files in the order written (no sorting
  when they are closed) \\
\code{WFDBGVMODE} & the default getvec mode (\code{setgvmode} with a negative mode) \\
\code{WFDBCAL} & the calibration file read by \code{wfdbx.calopen([])} (default
  \file{wfdbcal}) \\
\code{WFDB} & the WFDB path, used in legacy mode only (outside legacy mode, it is
  ignored, with a message) \\
\code{WFDB\_FLAC\_CHECK\_MD5} & if set, the MD5 signature of FLAC signal files is
  verified when they are read \\
\code{WFDBX\_MATLAB} & the MATLAB command that \file{wfdbx\_matlab.sh} runs (default
  \cmd{matlab}) \\
\bottomrule
\end{tabularx}
\end{table}

Tools that write signal parameters require them as options: \code{wrsamp} (with
\opt{-o}) requires \opt{-F} (frequency), \opt{-G} (gain), \opt{-R} (ADC resolution) and,
for calibrated signals, \opt{-U} (units); \code{wav2mit} requires \opt{-G}, and \opt{-U}
unless the gain is 0; \code{ad2m} and \code{ahaecg2mit} require \opt{-F}, \opt{-G},
\opt{-R}, and \opt{-U} unless the gain is 0; \code{a2m}, \code{rdann}, \code{ann2rr},
\code{sortann} and \code{m2a} require \opt{-F} when the record's header gives no
frequency (\cref{chap:tools}).

% ==========================================================================================
\section{Testing}
\label{sec:matlab-testing}

The unit tests of the MATLAB/Octave implementation are in \file{matlab/tests}: one
function \file{test\_*.m} per group, run by \code{run\_tests}. From the root of the
repository:

\begin{lstlisting}[style=shell]
octave-cli --no-gui --norc -q --eval "cd matlab/tests; run_tests"           # Octave
matlab -batch "cd matlab/tests; run_tests"                                  # MATLAB
octave-cli --no-gui --norc -q --eval "cd matlab/tests; run_tests('metadata')"  # by name
\end{lstlisting}

\code{run\_tests} prints \code{ok NAME} or \code{FAIL NAME: MESSAGE} for each test, then
a summary, and ends with an error (so that the command exits with status 1) if a test
failed; \code{run\_tests(PATTERN)} runs the tests whose names match a regular expression,
and \code{n = run\_tests(...)} returns the number of failures. The tests cover the C
standard library emulation, the WFDB library functions (with expected values printed by
the C library), the high-level functions, the bulk readers and writers against the
per-sample functions, the FLAC codec, finding files, \code{wfdbx.check\_record}, and the
metadata functions, with the cases shared by the seven languages (in
\file{tests/metadata}: results, warnings, and the headers, byte for byte). They run in
temporary folders, with copies of the test records of \file{tests/data}. The tests of
network access use a local HTTP server on 127.0.0.1 only (a Python script, started with
\cmd{python3}; without it, they are skipped).

The cross-language conformance suite runs every tool through \file{matlab/bin/wfdbx\_cli.m}
and compares its results, byte for byte, with those of the original C programs (1155
cases), and, with \opt{-{}-legacy}, in legacy mode (886 cases). From the root of the
repository:

\begin{lstlisting}[style=shell]
PYTHONPATH=$PWD/python/src python3 tests/conformance/harness.py check --impl octave -j8
PYTHONPATH=$PWD/python/src python3 tests/conformance/harness.py check --impl octave -j8 --legacy
\end{lstlisting}

The implementation \code{matlab} of the harness runs the tools through
\file{matlab/bin/wfdbx\_matlab.sh} instead.

% The Kotlin chapter of the WFDB-x user's manual. It follows the sections of the Python
% chapter (python.tex), in the same order. Each \apientry of the section "The high-level
% API" documents one public high-level function of kotlin/src/commonMain/kotlin/wfdbx
% (Record.kt, Check.kt, Metadata.kt): its signature listing holds the declarations of the
% extension function of Wfdb (all its overloads), copied from the code without their
% bodies, and its params name every parameter of these declarations, in code order. Each
% function is also a top-level function with the same parameters (the default instance).
% The listings of style "example" were compiled and run against the library.
\chapter{Kotlin}
\label{chap:kotlin}
\chapterlanguage{kotlin}

The Kotlin implementation of WFDB-x is the Kotlin Multiplatform library \code{wfdbx}, with
the Gradle coordinates \code{io.github.alphanumericslab:wfdbx:1.2.0}, in the directory
\file{kotlin} of the repository. It is written in common Kotlin code, which runs on every
JVM, on Android and on iOS, with the Kotlin standard library as its only dependency, and
never calls the C library. It is meant for mobile applications and JVM programs that read
and write WFDB records on the device, with the behavior of the other implementations of
WFDB-x, without native libraries, JNI or another runtime.

It offers WFDB-x in three ways:
\begin{itemize}
\item \textbf{The high-level API} (\cref{sec:kotlin-api}): functions that read and write
  whole records, headers, annotation files and header metadata in one call, with named
  arguments and default values; the samples are an \code{IntArray} (and a
  \code{DoubleArray} in physical units), frame by frame, and errors throw a
  \code{WfdbException}. They read signal files in large blocks, and are the fastest way
  to read a record. They are an addition of WFDB-x: the C package has no equivalent. They
  are the recommended interface for new code and apps.
\item \textbf{The WFDB library functions} under their C names (\code{isigopen},
  \code{getvec}, \code{getann}, \code{putann}, \ldots), with the behavior of the C library
  (\cref{sec:kotlin-library}). They are documented once, for all the languages, in
  \cref{chap:library}.
\item \textbf{The command-line tools} (\code{rdsamp}, \code{rdann}, \code{wrsamp},
  \code{xform} and 29 others), with the options, output and exit status of the C tools,
  run on the JVM through the launcher \file{kotlin/bin/wfdbx} (\cref{sec:kotlin-tools}).
  They are documented in \cref{chap:tools}.
\end{itemize}
The high-level functions return exactly the samples and annotations that the library
functions (and the C library) give.

% ==========================================================================================
\section{Installation and requirements}
\label{sec:kotlin-install}

\subsection{Requirements}

The library is built with Gradle 8.14.3 (the Gradle wrapper of the repository,
\file{kotlin/gradlew}, which downloads itself) and Kotlin 2.1.21, and needs a JDK 11 or
later to build. It has no dependency besides the Kotlin standard library (and
\code{kotlin-test} for the tests); the FLAC formats 508, 516 and 524 are read and written
by WFDB-x's own codec. The targets are:
\begin{itemize}
\item \code{jvm()}: bytecode for Java 11, which runs on every current JVM and on Android
  (API level 26 or later; the library uses \code{java.nio.file});
\item \code{iosArm64()}, \code{iosSimulatorArm64()} and \code{iosX64()}: declared, with a
  small platform layer for iOS (POSIX files and \code{NSURLSession}); they can only be
  built on macOS;
\item \code{linuxX64()}, built only with the Gradle option \code{-Pwfdbx.linux=true},
  with the Kotlin/Native toolchain.
\end{itemize}
The native code (iOS and Linux) is not compiled by the tests of the repository, which
build the JVM code and check that the common code, which holds all of the library's
logic, uses only common APIs.

\subsection{Building and installing the library}

The library is built from a copy of the repository and installed in the local Maven
repository (\file{\textasciitilde/.m2}). In the directory \file{kotlin}:

\begin{lstlisting}[style=shell]
./gradlew publishToMavenLocal
\end{lstlisting}

A Gradle project (an Android app, a JVM program, or the common source set of a Kotlin
Multiplatform project) that lists \code{mavenLocal()} among its repositories then adds it
to its dependencies:

\begin{lstlisting}[language=Kotlin]
dependencies {
    implementation("io.github.alphanumericslab:wfdbx:1.2.0")
}
\end{lstlisting}

The build downloads the Kotlin Gradle plugin from Maven Central. Where Maven Central
cannot be reached, the Gradle property \code{mavenMirror} names another Maven repository,
for plugins and dependencies alike, on the command line or in
\file{\textasciitilde/.gradle/gradle.properties}:

\begin{lstlisting}[style=shell]
./gradlew -PmavenMirror=https://maven-central.storage-download.googleapis.com/maven2/ publishToMavenLocal
\end{lstlisting}

The Kotlin package is \code{wfdbx} (\code{import wfdbx.*}); the C standard library
emulation is in \code{wfdbx.cstd}, the tools in \code{wfdbx.apps}, and the command-line
driver in \code{wfdbx.cli} (\cref{tab:kotlin-packages}).

\subsection{The command-line tools}

The tools need the runnable jar \file{kotlin/build/libs/wfdbx-all.jar}, built in the
directory \file{kotlin} by:

\begin{lstlisting}[style=shell]
./gradlew fatJar
bin/wfdbx --version                # WFDB-x 1.2.0
bin/wfdbx --list                   # the tools
\end{lstlisting}

\subsection{The sample record}
\label{sec:kotlin-resources}

The repository includes the sample record \code{100s} in \file{tests/data}: the first
minute of record 100 of the MIT-BIH Arrhythmia Database, two ECG signals (\code{MLII} and
\code{V5}) sampled at 360\,Hz, with the reference beat annotations (annotator
\code{atr}). The library jar does not include it: an app copies the files it needs
(\file{100s.hea}, \file{100s.dat}, \file{100s.atr}) to its own storage. Work that uses the
MIT-BIH Arrhythmia Database cites it as described on PhysioNet.

\subsection{Android and iOS}

On Android and iOS, the current directory is not a useful place for files: the input and
output directories are set to a directory of the app, such as
\code{context.filesDir.path} on Android (\cref{sec:kotlin-settings}). Reading URLs on
Android needs the \code{INTERNET} permission; files are downloaded with
\code{java.net.HttpURLConnection} on the JVM and Android, and with \code{NSURLSession} on
iOS. Legacy-mode sorting of annotation files runs the \code{sortann} program, which an
iOS app cannot do (the sort of normal mode, done in the library, works everywhere).

% ==========================================================================================
\section{Quick start}
\label{sec:kotlin-quickstart}

This example reads the header, ten seconds of the signals and the annotations of the
sample record \code{100s}, and computes the mean heart rate. It then writes the first ten
seconds as a new record, \code{first10}, with a metadata entry in its header. It assumes
that the input and output directories are a copy of \file{tests/data}
(\code{wfdbxSetInputDir("tests/data")}, as in \file{kotlin/README.md}, reads the sample
record from the repository).

\begin{lstlisting}[language=Kotlin,style=example]
import wfdbx.*

// The header only: no samples are read.
val hdr = readHeader("100s")
println("${hdr.nsig} ${hdr.fs} ${hdr.nsamples}")       // 2 360.0 21600

// Seconds 10 to 20 of both signals, in ADC and physical units, frame by frame.
val rec = readRecord("100s", from = TimeSpec("0:10"), to = TimeSpec("0:20"), physical = true)
println("${rec.names} ${rec.nsamples}")                // [MLII, V5] 3600
println("${rec.physical!![0]} ${rec.physicalAt(0, 1)}")    // -0.39 -0.275

// The beat annotations, and the mean heart rate.
val ann = readAnnotations("100s", "atr")
println("${ann.size} ${ann.symbols.take(3)} ${ann.samples.take(3)}")   // 75 [+, N, N] [18, 77, 370]
val beats = (0 until ann.size).filter { isqrs(ann.codes[it]) != 0 }.map { ann.samples[it] }
val rr = (1 until beats.size).map { (beats[it] - beats[it - 1]) / ann.fs }
println("${beats.size} beats, ${"%.1f".format(60 / rr.average())} beats per minute")   // 74 beats, 73.9 ...

// Write the first 10 seconds, in ADC units, as the record "first10".
val x = readRecord("100s", to = TimeSpec("0:10"))
writeRecord("first10", x.samples, x.nsig, x.fs, fmt = 16, gains = x.gains.toDoubleArray(),
            baselines = x.baselines.toIntArray(), adczeros = x.siginfo.map { it.adczero }.toIntArray(),
            adcres = x.siginfo.map { it.adcres }.toIntArray(), units = listOf("mV"), names = x.names,
            metadata = listOf("Source" to "MIT-BIH record 100"))

// Header metadata: "# Key: value" lines of the header.
setMetadata("first10", "Age", "69")
println(readMetadata("first10"))     // [(Source, MIT-BIH record 100), (Age, 69)]
println(checkRecord("first10"))      // no problems found
\end{lstlisting}

The header \file{100s.hea} gives no units: the example writes \code{"mV"} explicitly,
since WFDB-x never assumes units, a gain, an ADC resolution or a sampling frequency that
a record does not give (\cref{chap:behavior}).

% ==========================================================================================
\section{The high-level API}
\label{sec:kotlin-api}

The high-level API is made of the functions of this section and of the classes of their
results (\cref{sec:kotlin-results}), all in the package \code{wfdbx}. They are defined in
the files \file{Record.kt}, \file{Check.kt} and \file{Metadata.kt} of
\file{kotlin/src/commonMain/kotlin/wfdbx}, whose KDoc comments document each.

\subsection{Conventions}
\label{sec:kotlin-conventions}

\paragraph{The library state: \code{Wfdb} objects and the default instance.} The state of
the library (the open record, the annotators, the input and output directories, network
access, legacy mode) is held by an object of the class \code{Wfdb}. Each high-level
function is an extension function of \code{Wfdb}, such as
\code{fun Wfdb.readRecord(...)}, called on an object: \code{w.readRecord("100s")}. Each is
also a top-level function with the same name and the same parameters, such as
\code{readRecord("100s")}, which runs it on the default instance of the calling thread,
\code{wfdbDefault()} (one per thread; \code{withDefault \{ w -> ... \}} runs a block with
it). The entries of this section give the declarations of the extension functions; the
top-level functions have the same parameters, defaults and results. Separate \code{Wfdb}
objects share no state (except the process-wide list of open output files and the cache
of fetched URLs), so several records can be read at once, in one thread or in several; a
\code{Wfdb} object is used by one thread at a time.

\begin{lstlisting}[language=Kotlin]
val w = Wfdb()                          // a library state of its own
w.wfdbxSetInputDir("tests/data")
val a = w.readRecord("100s")            // on w
val b = readRecord("100s", inputDir = "tests/data")   // on the default instance
\end{lstlisting}

\paragraph{Named arguments.} The options are parameters with default values, given as
named arguments (\code{physical = true}); there are no option structures. A \code{null}
default means ``not given'' (for \code{inputDir} and \code{network}: the current
setting).

\paragraph{Record names and directories.} A record is named as in the other parts of
WFDB-x (\cref{chap:records,chap:files}), without the \file{.hea} extension (the readers,
and \code{writeAnnotations}, also accept a trailing \code{.hea}; \code{writeRecord} does
not). A record name is a full path, absolute (\code{"/data/mitdb/100"}) or in the home
directory (\code{"\textasciitilde/Documents/100s"}), used as given; a relative path
(\code{"100s"}, \code{"data/100s"}), read from the input directory if one is set,
otherwise from the current directory, and written to the output directory if one is set,
otherwise to the current directory (missing directories are created); or a URL
(\code{"https://physionet.org/files/mitdb/1.0.0/100"}), read only if network access is
on, which cannot be written. A writer given a name with a path creates the files in that
directory, and the header names the record and its signal file without the path
(\code{rec}, \code{rec.dat}), so that the record opens under any path that leads to it;
in legacy mode, as in the original, the header gives the name with its path. The input
and output directories are the current directory unless they are set, with \code{wfdbxSetInputDir} and \code{wfdbxSetOutputDir}
(\cref{sec:kotlin-settings}) or with the environment variables \code{WFDBX\_INPUT\_DIR}
and \code{WFDBX\_OUTPUT\_DIR}. URLs are read only when network access is on. The
functions that read take the parameters \code{inputDir} and \code{network}, and those
that write take \code{dir}: they apply to that one call, and the previous settings are
restored when it returns, also after an error. \code{inputDir} is optional, and applies
to relative names only, not to full paths or URLs; with \code{dir}, the record name must
be a plain name (\code{"rec"}), and the record opens from that directory under it. The signal files that a header names are
read from the header's directory. Record, annotator and directory names are converted to
bytes with the platform's file name encoding.

\paragraph{Times.} The parameters \code{from} and \code{to} are \code{TimeSpec} values: a
sample number (a \code{TimeSpec.Sample}, made from an \code{Int} or a \code{Long}), or a
WFDB time string, as \code{strtim} reads it (a \code{TimeSpec.Text}):
\begin{lstlisting}[language=Kotlin]
TimeSpec(3600)            // sample 3600
TimeSpec("0:10")          // 10 seconds; also "1:02:03.5"
TimeSpec("s3600")         // sample 3600
TimeSpec("e")             // the end of the record
TimeSpec("[14:00:00]")    // a time of day
\end{lstlisting}
A time in seconds needs the sampling frequency; when the record does not give it, such a
time throws a \code{WfdbException}.

\paragraph{Samples, frame by frame.} The samples of a record are one array, frame by
frame, as the library returns them: \code{samples[i * nsig + s]} is sample \code{i} of
signal \code{s} (\code{rec.at(i, s)}), so that a record of millions of samples is one
array, without a copy per signal. A signal with \code{spf} samples per frame has
\code{spf} consecutive values in each frame.

\paragraph{Per-signal and per-annotation arguments.} The arguments that describe each
signal (\code{gains}, \code{units}, \ldots) take one value for all the signals, or exactly
one value per signal. Those that describe each annotation (\code{subtypes}, \code{chans},
\ldots) have exactly one value per annotation.

\paragraph{Text.} Text in the files (signal descriptions, units, info strings, metadata,
aux strings, mnemonics) is a \code{String} with one character per byte (ISO-8859-1), so
that every byte round-trips. The fields of \code{Siginfo} are byte strings
(\code{ByteArray}), as in C; \code{latin1()}, from \code{wfdbx.cstd}, converts a
\code{ByteArray} to a \code{String} and back.

\paragraph{Signal parameters.} No sampling frequency, gain, ADC resolution or units is
ever assumed. Where one is needed and the record does not give it, it must be given as an
argument; otherwise the function throws a \code{WfdbException}. A record whose header
gives no sampling frequency has \code{fs = 0}.

\paragraph{Legacy mode.} With \code{legacy = true}, a function behaves as in the original
WFDB package, with its assumed signal parameters and its search for files
(\cref{chap:behavior}); legacy mode is on, on that \code{Wfdb} object, for the duration
of the call only. With \code{legacy = false} (the default), the current setting of
\code{wfdbxSetLegacy} applies, which is off unless it was turned on.

\paragraph{State.} The functions that read and write records and annotation files close
anything opened with the WFDB library functions on the same \code{Wfdb} object (they call
\code{wfdbquit()}) before they start, and close their own files before they return,
whatever the outcome. They leave the time resolution of output annotators
(\code{getafreq()}) and the getvec mode (\code{getgvmode()}) as they were, also after an
error and after reading an EDF record, whose opening sets both, as in C (\code{wfdbquit}
resets neither): \code{readRecord} sets the getvec mode for its own reading only, and the
\code{fs} of \code{writeAnnotations} applies to that file only. The metadata functions
open no signals and leave the library's state as it is.

\paragraph{The examples} of this section use the top-level functions, with the input and
output directories set to a copy of \file{tests/data} (\code{wfdbxSetInputDir(d)} and
\code{wfdbxSetOutputDir(d)}), which holds the test records of WFDB-x, among them
\code{100s}; they write records of their own there. The listings of style \code{example}
were compiled and run with the library on the JVM, and print what their comments say.

% ------------------------------------------------------------------------------------------
\subsection{Reading records}
\label{sec:kotlin-reading}

\apientry{readRecord}{Read the signals of a record}
\begin{lstlisting}[language=Kotlin,style=signature]
fun Wfdb.readRecord(
    record: String,
    channels: List<String>? = null,
    from: TimeSpec? = null,
    to: TimeSpec? = null,
    physical: Boolean = false,
    highres: Boolean = false,
    pad: Boolean = false,
    checksumWarning: Boolean = true,
    onWarning: ((String) -> Unit)? = null,
    inputDir: String? = null,
    network: Boolean? = null,
    legacy: Boolean = false,
): Record
\end{lstlisting}
Reads the signals of a record, or part of it, into one \code{IntArray} of ADC units,
frame by frame, together with the description of the record. The samples are those that
\code{getvec} returns (\cref{chap:library}), read in large blocks. Invalid samples are
\code{WFDB\_INVALID\_SAMPLE} ($-32768$) in ADC units, and NaN in physical units. When the
read reaches the end of the record, the samples are compared with the checksums of the
header (\cref{sec:kotlin-warnings}). On any error once the record is open, its input
signals are closed, and the getvec mode and the settings changed for the call are
restored.
\begin{params}
\item[record] Record name, without the \file{.hea} extension (a trailing \code{".hea"} is
  also accepted): a full path, absolute (\code{"/data/mitdb/100"}) or in the home
  directory (\code{"\textasciitilde/Documents/100s"}), used as given; a relative path
  (\code{"100s"}, \code{"data/100s"}), read from the input directory if one is set,
  otherwise from the current directory; or a URL
  (\code{"https://physionet.org/files/mitdb/1.0.0/100"}), read only if network access is
  on.
\item[channels] The signals to read, in the order in which they are returned: a list of
  signal numbers (\code{"0"}, \code{"1"}, \ldots) or signal descriptions (as in
  \code{Record.names}, such as \code{"MLII"}), which may be mixed. A name is looked up
  among the descriptions first, then with \code{findsig}. Default \code{null}: all the
  signals, in the order of the header.
\item[from] The first sample to read: a \code{TimeSpec}, a sample number ($\ge 0$) or a
  time string. Default \code{null}: the beginning of the record.
\item[to] The sample after the last one to read: a \code{TimeSpec}. Default \code{null}:
  the end of the record.
\item[physical] \code{true} to return the samples in physical units too, in
  \code{Record.physical} (a \code{DoubleArray} in the layout of \code{samples}), computed
  as $(\mathit{adc} - \mathit{baseline}) / \mathit{gain}$. Every signal read must then have
  a gain. Default \code{false}.
\item[highres] \code{false} (the default) to read a multi-frequency record at its frame
  frequency, with the samples of the faster signals averaged to one per frame;
  \code{true} to read every sample, at the highest sampling frequency of the record (the
  samples of the slower signals are repeated). \code{Record.fs} is then that frequency.
\item[pad] \code{true} to replace each invalid sample by the previous valid sample of the
  same signal. Default \code{false}.
\item[checksumWarning] \code{true} (the default) to warn when the samples do not match
  the checksums of the header; \code{false} for no warning. The samples are returned
  either way, and \code{checksumOk} of the result is \code{false}.
\item[onWarning] A function that receives the warning (without the \code{WFDB-x: }
  prefix and the newline). Default \code{null}: the library prints \code{WFDB-x: }
  followed by the warning on the standard error, unless \code{wfdbquiet()} was called.
\item[inputDir] The input directory for this call only (optional): a local directory, a
  \code{\textasciitilde} path or a URL, under which relative record names are read. It
  does not apply to full paths or URLs. Default \code{null}: the current setting
  (\code{wfdbxSetInputDir}).
\item[network] \code{true} to allow reading URLs for this call only, \code{false} to
  forbid it. Default \code{null}: the current setting (\code{wfdbxSetNetwork}).
\item[legacy] \code{true} to behave as the original WFDB package (legacy mode, for this
  call only). Default \code{false}.
\end{params}
\returns{A \code{Record} (\cref{tab:kotlin-record}).}
\raises{\code{WfdbException} if the header or a signal file cannot be read (``can't read
record REC''), if a channel is not in the record (``record REC has no signal C''), if
\code{from} is a negative sample number (``readRecord: from must not be negative''), if a
time string needs the sampling frequency and the header gives none (``sampling frequency
of record REC is not specified''), if \code{physical} is \code{true} and a signal read has
no gain (``signal N of record REC is uncalibrated''), if \code{from} is beyond the end of
the record (``can't seek to sample N of record REC''), or if a signal file is shorter than
its header says (``can't read signal data of record REC (getvec status -3)''). Warns on a
checksum error (see \code{checksumWarning} and \code{onWarning}).}
\begin{lstlisting}[language=Kotlin,style=example]
import wfdbx.*

// Seconds 10 to 20 of lead MLII, in physical units.
var rec = readRecord("100s", channels = listOf("MLII"), from = TimeSpec("0:10"),
                     to = TimeSpec("0:20"), physical = true)
println("${rec.nsamples} ${rec.nsig} ${rec.names} ${rec.fs} ${rec.sampfrom}")   // 3600 1 [MLII] 360.0 3600

// The first 3 samples of both signals, in ADC units, V5 first.
rec = readRecord("100s", channels = listOf("1", "0"), to = TimeSpec(3))
println(rec.samples.toList())       // [1011, 995, 1011, 995, 1011, 995]
println("${rec.names} ${rec.gains} ${rec.baselines} ${rec.units}")
// [V5, MLII] [200.0, 200.0] [1024, 1024] [null, null]

// The whole record: the checksums are verified.
rec = readRecord("100s")
println("${rec.nsamples} ${rec.checksumOk}")      // 21600 true
\end{lstlisting}
A multi-frequency record (here the test record \code{mfreq}, whose signals have 1, 2
and 4 samples per frame at a frame frequency of 100\,Hz) is read at its frame frequency,
or at its highest sampling frequency:
\begin{lstlisting}[language=Kotlin,style=example]
import wfdbx.*

var rec = readRecord("mfreq")
println("${rec.fs} ${rec.nsamples} ${rec.nsig}")     // 100.0 600 3
rec = readRecord("mfreq", highres = true)
println("${rec.fs} ${rec.nsamples} ${rec.nsig}")     // 400.0 2400 3
println(rec.siginfo.map { it.spf })                  // [1, 2, 4]
\end{lstlisting}

\apientry{readHeader}{Read the header of a record}
\begin{lstlisting}[language=Kotlin,style=signature]
fun Wfdb.readHeader(
    record: String,
    inputDir: String? = null,
    network: Boolean? = null,
    legacy: Boolean = false,
    onWarning: ((String) -> Unit)? = null,
): Header
\end{lstlisting}
Reads the header of a record, without opening its signal files. The result describes the
record and each of its signals, and holds the info strings and the metadata entries of the
header (of the \file{.hea} file only, not of a \file{.info} file; none for an EDF file or a
header read from the standard input). For a multi-segment record, the master header is
read, and \code{nsamples} is the length of the whole record.
\begin{params}
\item[record] Record name, without the \file{.hea} extension (a trailing \code{".hea"} is
  also accepted): a full path, absolute (\code{"/data/mitdb/100"}) or in the home
  directory (\code{"\textasciitilde/Documents/100s"}), used as given; a relative path
  (\code{"100s"}, \code{"data/100s"}), read from the input directory if one is set,
  otherwise from the current directory; or a URL
  (\code{"https://physionet.org/files/mitdb/1.0.0/100"}), read only if network access is
  on.
\item[inputDir] The input directory for this call only (optional): a local directory, a
  \code{\textasciitilde} path or a URL, under which relative record names are read. It
  does not apply to full paths or URLs. Default \code{null}: the current setting
  (\code{wfdbxSetInputDir}).
\item[network] \code{true} to allow reading URLs for this call only, \code{false} to
  forbid it. Default \code{null}: the current setting.
\item[legacy] \code{true} to read the header as the original WFDB package does. Default
  \code{false}.
\item[onWarning] A function that receives the warnings about \code{Header.metadata} (a
  key given more than once, \cref{sec:kotlin-warnings}). Default \code{null}: the library
  prints \code{WFDB-x: } followed by the warning on the standard error, unless
  \code{wfdbquiet()} was called.
\end{params}
\returns{A \code{Header} (\cref{tab:kotlin-header}).}
\raises{\code{WfdbException} if the header cannot be found or read (``can't read header
for record REC'').}
\begin{lstlisting}[language=Kotlin,style=example]
import wfdbx.*
import wfdbx.cstd.latin1

val hdr = readHeader("100s")
println("${hdr.name} ${hdr.fs} ${hdr.nsig} ${hdr.nsamples}")    // 100s 360.0 2 21600
for (s in hdr.siginfo) {
    println("${s.desc?.latin1()} ${s.fname?.latin1()} ${s.fmt} ${s.gain} ${s.adcres} ${s.units?.latin1()}")
}
// MLII 100s.dat 212 200.0 11 null
// V5 100s.dat 212 200.0 11 null
println("\"${hdr.comments[0]}\"")   // " 69 M 1085 1629 x1": the first info string
println(hdr.baseTime)               // null: the header gives no start time
\end{lstlisting}

\apientry{readAnnotations}{Read an annotation file}
\begin{lstlisting}[language=Kotlin,style=signature]
fun Wfdb.readAnnotations(
    record: String,
    annotator: String,
    from: TimeSpec? = null,
    to: TimeSpec? = null,
    fs: Double? = null,
    inputDir: String? = null,
    network: Boolean? = null,
    legacy: Boolean = false,
): Annotations
\end{lstlisting}
Reads the annotations of an annotation file, in file order. The sampling frequency,
needed to convert times in seconds, is that of the record's header, or, for a record
without a header or whose header gives none, the argument \code{fs}.
\begin{params}
\item[record] Record name, without the \file{.hea} extension (a trailing \code{".hea"} is
  also accepted): a full path, absolute (\code{"/data/mitdb/100"}) or in the home
  directory (\code{"\textasciitilde/Documents/100s"}), used as given; a relative path
  (\code{"100s"}, \code{"data/100s"}), read from the input directory if one is set,
  otherwise from the current directory; or a URL
  (\code{"https://physionet.org/files/mitdb/1.0.0/100"}), read only if network access is
  on. Its header, if there is one, gives the sampling frequency.
\item[annotator] Annotator name: the suffix of the annotation file. \code{"atr"} reads the
  file \file{RECORD.atr}.
\item[from] The first time to read: a \code{TimeSpec}. Default \code{null}: the beginning
  of the file.
\item[to] The time at which to stop: annotations at or after it are not read. Default
  \code{null}: the end of the file.
\item[fs] The sampling frequency, in Hz ($> 0$), to use when the record has no header or
  its header gives none; it is ignored otherwise. Default \code{null}: none.
\item[inputDir] The input directory for this call only (optional): a local directory, a
  \code{\textasciitilde} path or a URL, under which relative record names are read. It
  does not apply to full paths or URLs. Default \code{null}: the current setting
  (\code{wfdbxSetInputDir}).
\item[network] \code{true} to allow reading URLs for this call only, \code{false} to
  forbid it. Default \code{null}: the current setting.
\item[legacy] \code{true} to behave as the original WFDB package (which assumes 250\,Hz
  for a record without a frequency). Default \code{false}.
\end{params}
\returns{An \code{Annotations} object (\cref{tab:kotlin-annotations}). Its \code{fs} is
the frequency in which the times are counted: the one written in the annotation file, if
there is one, or else the sampling frequency; it is 0 if none of them is known.}
\raises{\code{WfdbException} if \code{fs} is given and is not $> 0$ (``readAnnotations:
fs must be > 0''), if the annotation file cannot be opened (``can't read annotator ANN for
record REC''), if a time string needs the sampling frequency and none is known, or if the
file ends before its end-of-file marker (``can't read annotation file ANN of record REC
(getann status -3)'').}
\begin{lstlisting}[language=Kotlin,style=example]
import wfdbx.*

var ann = readAnnotations("100s", "atr", from = TimeSpec("0:10"), to = TimeSpec("0:30"))
println("${ann.size} ${ann.fs}")                              // 24 360.0
println("${ann.samples.take(3)} ${ann.symbols.take(3)}")      // [3862, 4170, 4466] [N, N, N]

// The whole file: the first annotation is a rhythm change, with an aux string.
ann = readAnnotations("100s", "atr")
println("${ann.symbols[0]} ${ann.codes[0] == RHYTHM} ${ann.aux[0]}")   // + true (N
println(ann.symbols.toSortedSet())                            // [+, A, N]
\end{lstlisting}
For a record without a header, \code{fs} gives the frequency in which times in seconds
are converted (the library reports on the standard error that it cannot open the
header):
\begin{lstlisting}[language=Kotlin,style=example]
import wfdbx.*

writeAnnotations("beats", "qrs", longArrayOf(100, 460, 820), symbols = listOf("N", "N", "N"))
val ann = readAnnotations("beats", "qrs", to = TimeSpec("0:02"), fs = 360.0)
println("${ann.samples.toList()} ${ann.fs}")        // [100, 460] 360.0
\end{lstlisting}

% ------------------------------------------------------------------------------------------
\subsection{Writing records}
\label{sec:kotlin-writing}

\apientry{writeRecord}{Write a header and a signal file}
\begin{lstlisting}[language=Kotlin,style=signature]
fun Wfdb.writeRecord(
    record: String,
    samples: IntArray,
    nsig: Int,
    fs: Double,
    fmt: Int = 16,
    gains: DoubleArray? = null,
    baselines: IntArray? = null,
    adczeros: IntArray? = null,
    adcres: IntArray? = null,
    units: List<String?>? = null,
    names: List<String>? = null,
    baseTime: String? = null,
    comments: List<String>? = null,
    sampsPerFrame: IntArray? = null,
    dir: String? = null,
    legacy: Boolean = false,
    metadata: List<Pair<String, String?>>? = null,
)

fun Wfdb.writeRecord(
    record: String,
    samples: DoubleArray,
    nsig: Int,
    fs: Double,
    fmt: Int = 16,
    gains: DoubleArray? = null,
    baselines: IntArray? = null,
    adczeros: IntArray? = null,
    adcres: IntArray? = null,
    units: List<String?>? = null,
    names: List<String>? = null,
    baseTime: String? = null,
    comments: List<String>? = null,
    sampsPerFrame: IntArray? = null,
    dir: String? = null,
    legacy: Boolean = false,
    metadata: List<Pair<String, String?>>? = null,
)
\end{lstlisting}
Writes samples as a new record: the header \file{RECORD.hea} and one signal file,
\file{RECORD.dat}, which holds all the signals. Existing files of these names are
overwritten. The first overload takes the samples in ADC units (an \code{IntArray}), the
second in physical units (a \code{DoubleArray}), converted to ADC units as
$\mathrm{round}(x \cdot \mathit{gain}) + \mathit{baseline}$, with halves rounded away from
zero; NaN becomes an invalid sample, and an infinite value, or one whose ADC value is
beyond the range of 32-bit integers, is an error. The sampling frequency, the gain and
ADC resolution of each signal, and the units of each calibrated signal must be given:
there are no defaults. The per-signal arguments (\code{gains}, \code{baselines},
\code{adczeros}, \code{adcres}, \code{units}, \code{names}, \code{sampsPerFrame}) take one
value for all the signals, or one value per signal. All the arguments are checked before
any file is created, and the output signals are closed whatever happens.
\begin{params}
\item[record] Record name (not empty), without the \file{.hea} extension: a plain name
  (\code{"rec"}) or a relative path (\code{"data/rec"}), created in the output directory
  if one is set, otherwise in the current directory, missing directories included; or a
  full path, absolute or in the home directory (\code{"\textasciitilde/Documents/rec"}),
  used as given. With a path, the files are created in that directory, and the header
  names the record and its signal file without the path (\code{rec}, \code{rec.dat}), so
  that the record opens under any path that leads to it (in legacy mode, as in the
  original, the header gives the name with its path). A URL cannot be written.
\item[samples] The samples, frame by frame: signal 0, signal 1, \ldots{} of frame 0, then
  those of frame 1, and so on (a signal with \code{spf} samples per frame has \code{spf}
  consecutive values in each frame). An \code{IntArray} of ADC units, or a
  \code{DoubleArray} of physical units. Its size must be a multiple of the frame size.
\item[nsig] The number of signals (at least 1).
\item[fs] The sampling frequency, in Hz (frames per second): a finite number $> 0$.
\item[fmt] The storage format of all the signals, such as \code{16}, \code{212},
  \code{24} or \code{516} (FLAC) (\cref{chap:formats}). Default \code{16}.
\item[gains] The gain of each signal, in ADC units per physical unit: a finite number;
  \code{0} marks an uncalibrated signal, which is allowed only for samples given in ADC
  units. Negative gains are allowed. Required (default \code{null}: not given).
\item[baselines] The ADC value that corresponds to 0 physical units. Default
  \code{null}: the ADC zero of the signal.
\item[adczeros] The ADC value at the middle of the ADC range. Default \code{null}: 0.
\item[adcres] The ADC resolution, in bits ($\ge 1$). Required (default \code{null}: not
  given).
\item[units] The physical units, such as \code{"mV"}; required for each signal whose gain
  is not 0 (\code{null} for an uncalibrated signal). Default \code{null}: not given.
\item[names] The signal descriptions, such as \code{"MLII"}. Default \code{null}: no
  descriptions.
\item[baseTime] The start time, with an optional date: \code{"hh:mm:ss"} or
  \code{"hh:mm:ss dd/mm/yyyy"}. Default \code{null}: no start time.
\item[comments] Info strings, written at the end of the header as comment lines
  (\code{"\#"} followed by the string). Default \code{null}: none.
\item[sampsPerFrame] For a multi-frequency record, the number of samples per frame of
  each signal (integers $\ge 1$); \code{fs} is then the frame frequency. Default
  \code{null}: one sample per frame for every signal.
\item[dir] A directory (which must exist) in which to create the files, for this call
  only. \code{record} must then be a plain name, without a directory, and the record opens
  from that directory under its plain name. Default \code{null} (or \code{""}): the output
  directory.
\item[legacy] \code{true} to behave as the original WFDB package: the gain defaults to
  200, the ADC resolution to one based on the format, and no units are required. Default
  \code{false}.
\item[metadata] Metadata entries (\cref{chap:metadata}) to write to the header after the
  comments, as \code{\# Key: value} (through \code{putinfo}), in the order given: a list
  of pairs, \code{listOf("Age" to "63")} (a map converts with \code{toList()}). Keys and
  values follow the rules of \code{setMetadata}; no two keys may be equal (ignoring
  letter case), none may be the key of an entry in \code{comments}, and a \code{null}
  value is not allowed. Default \code{null}: none.
\end{params}
\returns{Nothing (\code{Unit}).}
\raises{\code{WfdbException}, before any file is created: ``writeRecord: sampsPerFrame
must hold integers >= 1'' (checked first, also in legacy mode), ``writeRecord: the record
name must not be empty'', ``writeRecord: at least one signal is required'',
``writeRecord: NAME must have one value per signal'', ``writeRecord: the sampling
frequency must be > 0'' (not checked in legacy mode), ``writeRecord: the sampling
frequency must be finite'', ``writeRecord: gains must be specified'', ``writeRecord:
gains must be finite'', ``writeRecord: ADC resolution(s) must be specified'',
``writeRecord: units must be specified'', ``writeRecord: incorrect base time 'TIME'{}'',
``writeRecord: the number of samples is not a multiple of the frame size'',
``writeRecord: signal N is uncalibrated'' and ``writeRecord: sample I is infinite or
beyond the range of 32-bit ADC units'' (physical units), ``record name must not contain a
directory when dir is given'', ``directory DIR does not exist''; a
\code{WfdbArgumentException} if an entry of \code{metadata} is not valid (``writeRecord:
metadata key 'a:b' must not contain ':'{}'', ``writeRecord: value of metadata key 'Age'
must not be null'', ``writeRecord: metadata key 'Age' is also given in comments'',
\ldots). \code{WfdbException} if the signal file or the header cannot be created
(``can't create signal file for record REC'', ``can't create header for record REC''),
or cannot be written completely, for example on a full disk (``write error on record
REC'').}
\begin{lstlisting}[language=Kotlin,style=example]
import wfdbx.*
import kotlin.math.abs

// Two signals of random numbers, in millivolts, at 360 Hz: the record "rand"
// (rand.hea and rand.dat).
val rnd = kotlin.random.Random(1)
val x = DoubleArray(3600 * 2) { rnd.nextDouble(-3.0, 3.0) }
writeRecord("rand", x, nsig = 2, fs = 360.0, fmt = 16, gains = doubleArrayOf(1000.0),
            adcres = intArrayOf(16), units = listOf("mV"), names = listOf("A", "B"),
            baseTime = "14:00:03 21/06/2024", comments = listOf("random numbers"),
            metadata = listOf("Generator" to "kotlin.random"))
val rec = readRecord("rand", physical = true)
println("${rec.nsamples} ${rec.baseTime}")     // 3600 [14:00:03 21/06/2024]
println(x.indices.all { abs(rec.physical!![it] - x[it]) <= 0.0005 })   // true: rounded to 1/1000 mV
println(readHeader("rand").comments)           // [random numbers,  Generator: kotlin.random]

// An uncalibrated signal (gain 0, no units), in ADC units, in the directory "out".
writeRecord("out/counts", IntArray(1000) { it }, nsig = 1, fs = 100.0, gains = doubleArrayOf(0.0),
            adcres = intArrayOf(12))
println(readRecord("out/counts").samples.take(3))    // [0, 1, 2]
\end{lstlisting}
A multi-frequency record: signal \code{fast} has 2 samples per frame, signal
\code{slow} one; each frame holds the two samples of \code{fast} and then the sample of
\code{slow}:
\begin{lstlisting}[language=Kotlin,style=example]
import wfdbx.*

val frames = 100                                    // 1 second at 100 frames per second
val x = IntArray(3 * frames) { if (it % 3 == 2) 500 else 2 * (it / 3) + it % 3 }
writeRecord("mf", x, nsig = 2, fs = 100.0, gains = doubleArrayOf(200.0), adcres = intArrayOf(12),
            units = listOf("mV"), names = listOf("fast", "slow"), sampsPerFrame = intArrayOf(2, 1))
val rec = readRecord("mf", highres = true)
println("${rec.fs} ${rec.samples.take(6)}")        // 200.0 [0, 500, 1, 500, 2, 500]
\end{lstlisting}

\apientry{writeAnnotations}{Write an annotation file}
\begin{lstlisting}[language=Kotlin,style=signature]
fun Wfdb.writeAnnotations(
    record: String,
    annotator: String,
    samples: LongArray,
    symbols: List<String>? = null,
    codes: IntArray? = null,
    subtypes: IntArray? = null,
    chans: IntArray? = null,
    nums: IntArray? = null,
    aux: List<String?>? = null,
    fs: Double? = null,
    sort: Boolean? = null,
    dir: String? = null,
    legacy: Boolean = false,
)

fun Wfdb.writeAnnotations(
    record: String,
    annotator: String,
    a: Annotations,
    fs: Double? = null,
    sort: Boolean? = null,
    dir: String? = null,
    legacy: Boolean = false,
)
\end{lstlisting}
Writes an annotation file, \file{RECORD.ANNOTATOR}, in the MIT format; an existing file of
that name is overwritten. The record need not have a header. The annotation types are
given as \code{symbols} or as \code{codes} (\cref{chap:formats}). If the annotations are
not in time order, the file is sorted when it is written, unless \code{sort = false}. The
second overload writes the annotations of an \code{Annotations} object, as
\code{readAnnotations} returns them: its samples, codes, subtypes, channels, annotator
numbers and aux strings (its record, annotator and \code{fs} are not used). All the
arguments are checked before the file is created.
\begin{params}
\item[record] Record name, without the \file{.hea} extension (a trailing \code{".hea"} is
  also accepted): a plain name or a relative path (\code{"data/100s"}) is created in the
  output directory if one is set, otherwise in the current directory; a full path,
  absolute or in the home directory (\code{"\textasciitilde/Documents/100s"}), is used as
  given. The file is \file{RECORD.ANNOTATOR} at that place (\file{data/100s.qrs}). A URL
  cannot be written.
\item[annotator] Annotator name: the suffix of the file, made of letters, digits,
  \code{\_}, \code{-} and \code{\textasciitilde}.
\item[samples] The annotation times, in samples, one per annotation.
\item[symbols] The mnemonics of the annotation types, such as \code{"N"} or \code{"V"},
  one per annotation. An unknown mnemonic becomes code 0 (\code{NOTQRS}). Default
  \code{null}: the types are given by \code{codes}.
\item[codes] The annotation codes (\code{NORMAL}, \code{PVC}, \ldots), one per annotation,
  instead of \code{symbols}. If both are given (as \code{readAnnotations} returns them),
  each symbol must be that of its code (an empty symbol is not checked). Default
  \code{null}.
\item[subtypes] The subtype of each annotation, stored as a signed byte (other values wrap
  around to $-128 \ldots 127$, as in C). Default \code{null}: 0.
\item[chans] The channel number of each annotation, stored as an unsigned byte (other
  values wrap around to $0 \ldots 255$). Default \code{null}: 0.
\item[nums] The annotator number of each annotation, stored as a signed byte. Default
  \code{null}: 0.
\item[aux] The auxiliary text of each annotation, with \code{null} where there is none.
  Texts longer than 255 bytes are cut. Default \code{null}: none.
\item[fs] If given (a finite number $> 0$), the frequency in which the times are counted,
  written to the file. Default \code{null}: none is written.
\item[sort] \code{true} to always sort the file, \code{false} to keep the given order
  (without the library's advice to run \code{sortann}), for this call only. Default
  \code{null}: the environment variable \code{WFDBANNSORT} decides (sorting is on unless
  it is \code{0}).
\item[dir] A directory (which must exist) in which to create the file, for this call
  only. \code{record} must then be a plain name, without a directory: the file is
  \file{DIR/RECORD.ANNOTATOR}. Default \code{null}: the output directory.
\item[legacy] \code{true} to behave as the original WFDB package. Default \code{false}.
\item[a] The annotations to write (second overload): an \code{Annotations} object, as
  \code{readAnnotations} returns it.
\end{params}
\returns{Nothing (\code{Unit}).}
\raises{\code{WfdbException}, before the file is created, if neither \code{symbols} nor
\code{codes} is given (``symbols or codes must be given''), if both are given and
disagree (``writeAnnotations: symbols and codes disagree at annotation K'', K counting
from 0), if an array does not have one value per annotation (``writeAnnotations: NAME
must have one value per annotation''), if \code{fs} is not finite and $> 0$
(``writeAnnotations: fs must be finite and > 0''), or as for the \code{dir} of
\code{writeRecord}; if the file cannot be created (``can't create annotator ANN for record
REC''), or cannot be written completely (``write error on annotator ANN for record
REC'').}
\begin{lstlisting}[language=Kotlin,style=example]
import wfdbx.*

// Three beats and a rhythm change, by symbol, with an aux string for the rhythm.
writeAnnotations("100s", "test", longArrayOf(18, 77, 370, 662),
                 symbols = listOf("+", "N", "V", "N"), aux = listOf("(N", null, null, null))
var ann = readAnnotations("100s", "test")
println("${ann.symbols} ${ann.aux}")       // [+, N, V, N] [(N, null, null, null]

// A copy of the reference annotations, as readAnnotations returns them.
val a = readAnnotations("100s", "atr")
writeAnnotations("100s", "copy", a)
val b = readAnnotations("100s", "copy")
println("${b.samples.contentEquals(a.samples)} ${b.symbols == a.symbols} ${b.aux == a.aux}")
// true true true

// Annotations out of time order are sorted when the file is written (the library
// reports it on the standard error).
writeAnnotations("100s", "unsorted", longArrayOf(500, 100), codes = intArrayOf(NORMAL, NORMAL))
println(readAnnotations("100s", "unsorted").samples.toList())     // [100, 500]
\end{lstlisting}

% ------------------------------------------------------------------------------------------
\subsection{Checking records}
\label{sec:kotlin-checking}

\apientry{checkRecord}{Check a record and its annotation files}
\begin{lstlisting}[language=Kotlin,style=signature]
fun Wfdb.checkRecord(
    record: String,
    annotators: List<String> = emptyList(),
    inputDir: String? = null,
    network: Boolean? = null,
): CheckReport
\end{lstlisting}
Checks the header of a record, its signal specifications, every sample of its signal
files, each segment of a multi-segment record, and the annotation files named, and
reports each problem found. A problem in the record never throws an exception: it is
reported, with a severity (\code{Severity.ERROR}: the record is damaged or does not follow
the WFDB format, and programs will read wrong data or fail; \code{Severity.WARNING}: the
record can be read, but something is missing or unusual, and may give wrong results;
\code{Severity.INFO}: a fact worth knowing, such as invalid samples). The problems, their
codes and their messages are the same in the seven languages
(\cref{tab:python-issue-codes}). The checks run in normal mode, even if legacy mode is on,
and the library prints no messages during the call. They are similar to those of Benjamin
Moody's \code{wfdbcheck}.
\begin{params}
\item[record] Record name, without the \file{.hea} extension (a trailing \code{".hea"} is
  also accepted): a full path, absolute (\code{"/data/mitdb/100"}) or in the home
  directory (\code{"\textasciitilde/Documents/100s"}), used as given; a relative path
  (\code{"100s"}, \code{"data/100s"}), read from the input directory if one is set,
  otherwise from the current directory; or a URL
  (\code{"https://physionet.org/files/mitdb/1.0.0/100"}), read only if network access is
  on (as for \code{readRecord}).
\item[annotators] The annotator names whose files are checked, such as
  \code{listOf("atr")}. Default \code{emptyList()}: none.
\item[inputDir] The input directory for this call only (optional): a local directory, a
  \code{\textasciitilde} path or a URL, under which relative record names are read. It
  does not apply to full paths or URLs. Default \code{null}: the current setting
  (\code{wfdbxSetInputDir}).
\item[network] \code{true} to allow reading URLs for this call only, \code{false} to
  forbid it. Default \code{null}: the current setting.
\end{params}
\returns{A \code{CheckReport} (\cref{tab:kotlin-checkreport}): \code{report.ok} is
\code{true} if no problem is an error, \code{report.issues} lists the problems
(\code{Issue} objects, \cref{tab:kotlin-issue}), and \code{println(report)} prints one
line per problem, or \code{no problems found}.}
\raises{\code{WfdbException} if \code{record} is empty (``record must not be empty''). A
record that cannot be read is reported (code \code{header\_unreadable}), not thrown.}
\begin{lstlisting}[language=Kotlin,style=example]
import wfdbx.*

var report = checkRecord("100s", listOf("atr"))
println(report.ok)          // true: no errors
println(report)
// warning: signal 0: the units are not specified
// warning: signal 1: the units are not specified

// A damaged record: one sample of a 12-bit signal is overwritten with 2048 (JVM file
// access, in the output directory).
writeRecord("damaged", IntArray(100) { it }, nsig = 1, fs = 100.0, gains = doubleArrayOf(100.0),
            adcres = intArrayOf(12), units = listOf("mV"), names = listOf("ramp"))
java.io.RandomAccessFile(wfdbxGetOutputDir() + "/damaged.dat", "rw").use { f ->
    f.seek(20)                                  // sample 10 (2 bytes per sample in format 16)
    f.write(byteArrayOf(0x00, 0x08))            // 0x0800 = 2048, little-endian
}
report = checkRecord("damaged")
println(report.ok)          // false
for (issue in report.issues) println("${issue.severity} ${issue.code} ${issue.where}")
// error sample_range signal 0
// error checksum signal 0
println(report.lines()[0])
// error: signal 0: 1 sample is outside the ADC range (-2048 to 2047): sample 10 (value 2048)
\end{lstlisting}

% ------------------------------------------------------------------------------------------
\subsection{Header metadata}
\label{sec:kotlin-metadata}

These functions read and change the metadata of a header: its info strings of the form
\code{\# Key: value}, such as \code{\# Age: 63} (\cref{chap:metadata}). In short: keys are
compared without regard to letter case (A--Z and a--z only); a key must not be empty and
must not contain \code{:}, \code{\#}, a line break or a NUL character, and a value must
not contain a line break or a NUL character; spaces, tabs, VT, FF and CR at either end of
keys and values are removed; text is one character per byte (ISO-8859-1), and a character
beyond U+00FF is an argument error. A change rewrites the header in one step: the new
header is written to a temporary file in the header's directory, which then replaces the
header, keeping its permission bits where the platform has them; every other byte of the
header is kept, and all the arguments are checked first, so an error leaves the header
unchanged. If nothing changes (and \code{dir} is not given), nothing is written. Only the
\file{.hea} file is read (not a \file{.info} file); for a multi-segment record, it is the
master header, unless a segment is named. Invalid arguments throw a
\code{WfdbArgumentException}, and a header that cannot be read or written a
\code{WfdbException} (\cref{sec:kotlin-warnings}). Metadata often holds information about
patients: WFDB-x stores it as plain text, exactly as given, and does not de-identify it.

\begin{phibox}
A header can hold identifying data about a patient (a name, a date of birth, a medical
record number). The metadata functions read and write such data as given. Records that
are shared must be de-identified first.
\end{phibox}

\apientry{readMetadata}{Read the metadata of a record}
\begin{lstlisting}[language=Kotlin,style=signature]
fun Wfdb.readMetadata(
    record: String,
    onWarning: ((String) -> Unit)? = null,
    inputDir: String? = null,
    network: Boolean? = null,
): List<Pair<String, String>>
\end{lstlisting}
Reads the metadata entries of the header of a record, in file order. An info string is an
entry if it contains a colon: the key is the text before the first colon, the value the
rest of the line, both without spaces at either end. Info strings without a colon are free
comments, which are not returned.
\begin{params}
\item[record] Record name, without the \file{.hea} extension (a trailing \code{".hea"} is
  also accepted): a full path, absolute (\code{"/data/mitdb/100"}) or in the home
  directory (\code{"\textasciitilde/Documents/100s"}), used as given; a relative path
  (\code{"100s"}, \code{"data/100s"}), read from the input directory if one is set,
  otherwise from the current directory; or a URL
  (\code{"https://physionet.org/files/mitdb/1.0.0/100"}), read only if network access is
  on.
\item[onWarning] A function that receives the warnings (without the \code{WFDB-x: }
  prefix and the newline). Default \code{null}: the library prints \code{WFDB-x: }
  followed by the warning on the standard error, unless \code{wfdbquiet()} was called.
\item[inputDir] The input directory for this call only (optional): a local directory, a
  \code{\textasciitilde} path or a URL, under which relative record names are read. It
  does not apply to full paths or URLs. Default \code{null}: the current setting
  (\code{wfdbxSetInputDir}).
\item[network] \code{true} to allow reading URLs for this call only, \code{false} to
  forbid it. Default \code{null}: the current setting.
\end{params}
\returns{A list of (key, value) pairs, in file order. A key that appears more than once
(ignoring letter case) is returned once, with its first value, and a warning is issued
(\cref{sec:kotlin-warnings}).}
\raises{\code{WfdbException} if the header cannot be found or read (``can't read header
for record REC''); \code{WfdbArgumentException} if \code{record} is empty or \code{"-"}
(the standard input).}
\begin{lstlisting}[language=Kotlin,style=example]
import wfdbx.*

writeRecord("pat", IntArray(360), nsig = 1, fs = 360.0, gains = doubleArrayOf(200.0),
            adcres = intArrayOf(12), units = listOf("mV"), names = listOf("ECG"),
            comments = listOf("recorded at rest"), metadata = listOf("Age" to "63", "Sex" to "F"))
println(readMetadata("pat"))           // [(Age, 63), (Sex, F)]
println(readHeader("pat").comments)
// [recorded at rest,  Age: 63,  Sex: F]: the free comment is not an entry
\end{lstlisting}

\apientry{getMetadata}{Return the value of a metadata entry}
\begin{lstlisting}[language=Kotlin,style=signature]
fun Wfdb.getMetadata(
    record: String,
    key: String,
    onWarning: ((String) -> Unit)? = null,
    inputDir: String? = null,
    network: Boolean? = null,
): String?
\end{lstlisting}
Returns the value of one metadata entry of the header of a record. The key is checked
before the header is read.
\begin{params}
\item[record] Record name (as for \code{readMetadata}).
\item[key] The key, such as \code{"Age"}. Letter case is ignored, and spaces at either end
  are removed.
\item[onWarning] A function that receives the warnings (as for \code{readMetadata}).
  Default \code{null}: the standard error.
\item[inputDir] The input directory for this call only. Default \code{null}: the current
  setting.
\item[network] \code{true} to allow reading URLs for this call only, \code{false} to
  forbid it. Default \code{null}: the current setting.
\end{params}
\returns{The value (\code{""} for an entry with an empty value), or \code{null} if the
header has no entry for \code{key}. If it has more than one, the first is returned, with a
warning.}
\raises{\code{WfdbException} if the header cannot be found or read;
\code{WfdbArgumentException} if \code{record} is empty or \code{"-"}, or if \code{key} is
not a valid key (see \code{setMetadata}).}
\begin{lstlisting}[language=Kotlin,style=example]
import wfdbx.*

writeRecord("pat", IntArray(360), nsig = 1, fs = 360.0, gains = doubleArrayOf(200.0),
            adcres = intArrayOf(12), units = listOf("mV"), metadata = listOf("Age" to "63", "Sex" to "F"))
println(getMetadata("pat", "age"))       // 63
println(getMetadata("pat", "Height"))    // null
\end{lstlisting}

\apientry{setMetadata}{Set the value of a metadata entry}
\begin{lstlisting}[language=Kotlin,style=signature]
fun Wfdb.setMetadata(
    record: String,
    key: String,
    value: String,
    inputDir: String? = null,
    network: Boolean? = null,
    dir: String? = null,
)
\end{lstlisting}
Sets the value of a metadata entry in the header of a record. If the header has an entry
for \code{key} (ignoring letter case), its line is replaced by \code{\# Key: value}, at
the same position and with the same line ending, with the spelling of \code{key} given
here (even if the value is the same); otherwise the entry is added after the last info
string of the header (or, if it has none, after its last line that is not blank), ending
as the record line ends. The entry is written as \code{\# Key: value}, or \code{\# Key:}
if the value is empty. Every other byte of the header is kept.
\begin{params}
\item[record] Record name (as for \code{readMetadata}). The header is changed where it was
  read, unless \code{dir} is given; a header read from a URL can only be written
  elsewhere, with \code{dir}.
\item[key] The key, such as \code{"Age"} or \code{"Date of birth"}. Spaces at either end
  are removed. It must not be empty, and must not contain \code{:}, \code{\#}, a line
  break or a NUL character.
\item[value] The value, which may be a sentence. Spaces at either end are removed. It
  must not contain a line break or a NUL character; it may be empty.
\item[inputDir] The input directory in which the header is read, for this call only.
  Default \code{null}: the current setting.
\item[network] \code{true} to allow reading the header from a URL, for this call only,
  \code{false} to forbid it. Default \code{null}: the current setting.
\item[dir] A directory, which must exist, in which to write the changed header, under the
  same file name, instead of changing the header itself (the header is then written even
  if nothing changes); the signal files are not copied. Default \code{null}: the header
  itself is changed.
\end{params}
\returns{Nothing (\code{Unit}).}
\raises{\code{WfdbArgumentException} if \code{key} or \code{value} is not valid
(``metadata key 'a:b' must not contain ':'{}'', ``value of metadata key 'Age' must not
contain a line break'', \ldots), if the header has more than one entry for \code{key}
(``header of record REC has more than one entry for key KEY (lines L1, L2); remove them
first''), if it has no record line, if it was read from a URL and \code{dir} is not given,
or if \code{record} is empty or \code{"-"}; \code{WfdbException} if the header cannot be
read (``can't read header for record REC''), if \code{dir} does not exist (``directory
DIR does not exist''), or if the changed header cannot be written (``can't write header
for record REC'').}
\begin{lstlisting}[language=Kotlin,style=example]
import wfdbx.*

writeRecord("pat", IntArray(360), nsig = 1, fs = 360.0, gains = doubleArrayOf(200.0),
            adcres = intArrayOf(12), units = listOf("mV"), metadata = listOf("Age" to "63", "Sex" to "F"))
setMetadata("pat", "AGE", "64")                                // replaces "# Age: 63"
setMetadata("pat", "Diagnosis", "sinus rhythm, no ectopy")     // a new entry
println(readMetadata("pat"))
// [(AGE, 64), (Sex, F), (Diagnosis, sinus rhythm, no ectopy)]

// Write the changed header to another directory; pat.hea is left unchanged.
val changed = wfdbxGetOutputDir() + "/changed"
java.io.File(changed).mkdir()                                  // (JVM)
setMetadata("pat", "Sex", "M", dir = changed)
println("${getMetadata("pat", "Sex")} ${getMetadata("$changed/pat", "Sex")}")    // F M
\end{lstlisting}

\apientry{removeMetadata}{Remove a metadata entry}
\begin{lstlisting}[language=Kotlin,style=signature]
fun Wfdb.removeMetadata(
    record: String,
    key: String,
    inputDir: String? = null,
    network: Boolean? = null,
    dir: String? = null,
): Int
\end{lstlisting}
Removes every entry for \code{key} (ignoring letter case) from the header of a record,
with its line ending, so it also repairs a header that has more than one. Other lines are
kept, byte for byte. Nothing is written if the header has no entry for \code{key}, unless
\code{dir} is given.
\begin{params}
\item[record] Record name (as for \code{setMetadata}).
\item[key] The key (as for \code{setMetadata}).
\item[inputDir] The input directory in which the header is read, for this call only.
  Default \code{null}: the current setting.
\item[network] \code{true} to allow reading the header from a URL, for this call only,
  \code{false} to forbid it. Default \code{null}: the current setting.
\item[dir] A directory, which must exist, in which to write the changed header instead of
  changing the header itself. Default \code{null}: the header itself is changed.
\end{params}
\returns{The number of lines removed (0 if the header has no entry for \code{key}).}
\raises{As \code{setMetadata}, except that more than one entry for \code{key} is not an
error.}
\begin{lstlisting}[language=Kotlin,style=example]
import wfdbx.*

writeRecord("pat", IntArray(360), nsig = 1, fs = 360.0, gains = doubleArrayOf(200.0),
            adcres = intArrayOf(12), units = listOf("mV"), metadata = listOf("Age" to "63", "Sex" to "F"))
println(removeMetadata("pat", "sex"))     // 1
println(removeMetadata("pat", "Sex"))     // 0: there is no such entry now
println(readMetadata("pat"))              // [(Age, 63)]
\end{lstlisting}

\apientry{updateMetadata}{Make several changes to the metadata}
\begin{lstlisting}[language=Kotlin,style=signature]
fun Wfdb.updateMetadata(
    record: String,
    changes: List<Pair<String, String?>>,
    inputDir: String? = null,
    network: Boolean? = null,
    dir: String? = null,
)

fun Wfdb.updateMetadata(
    record: String,
    changes: Map<String, String?>,
    inputDir: String? = null,
    network: Boolean? = null,
    dir: String? = null,
)
\end{lstlisting}
Makes several changes to the metadata of a record, in one write of the header. Each
change is applied in order, as \code{setMetadata} does, or, for a \code{null} value, as
\code{removeMetadata} does. All the changes are checked first: if one is not valid, none
is made. The second overload takes the changes as a map, applied in its iteration order
(the order of insertion for the maps of \code{mapOf} and \code{linkedMapOf}).
\begin{params}
\item[record] Record name (as for \code{setMetadata}).
\item[changes] The changes: a list of (key, value) pairs, or a map from key to value; a
  \code{null} value removes the entry. No two keys may differ only in letter case.
\item[inputDir] The input directory in which the header is read, for this call only.
  Default \code{null}: the current setting.
\item[network] \code{true} to allow reading the header from a URL, for this call only,
  \code{false} to forbid it. Default \code{null}: the current setting.
\item[dir] A directory, which must exist, in which to write the changed header instead of
  changing the header itself. Default \code{null}: the header itself is changed.
\end{params}
\returns{Nothing (\code{Unit}).}
\raises{As \code{setMetadata}; also \code{WfdbArgumentException} if two keys of
\code{changes} differ only in letter case (``metadata key 'dx' is given more than once
(keys are compared without regard to letter case)'').}
\begin{lstlisting}[language=Kotlin,style=example]
import wfdbx.*

writeRecord("pat", IntArray(360), nsig = 1, fs = 360.0, gains = doubleArrayOf(200.0),
            adcres = intArrayOf(12), units = listOf("mV"), metadata = listOf("Age" to "63", "Sex" to "F"))
updateMetadata("pat", mapOf("Age" to "64", "Sex" to null, "Dx" to "164889003"))
println(readMetadata("pat"))         // [(Age, 64), (Dx, 164889003)]
try {
    updateMetadata("pat", listOf("Dx" to "426783006", "dx" to null))
} catch (e: WfdbArgumentException) {
    println(e.message)        // metadata key 'dx' is given more than once (...)
}
println(getMetadata("pat", "Dx"))    // 164889003: nothing was changed
\end{lstlisting}

% ------------------------------------------------------------------------------------------
\subsection{Settings}
\label{sec:kotlin-settings}

These functions set where files are read and written, network access and legacy mode, on
a \code{Wfdb} object (the top-level functions: on the default instance), for the
high-level functions and the library functions alike (\cref{chap:files,chap:behavior}).
Each \code{set} function returns the previous setting, so that it can be restored. They
are the high-level forms of the library functions \code{wfdbx\_set\_input\_dir},
\code{wfdbx\_set\_network}, \ldots{} (\cref{sec:kotlin-library}), which take and return
bytes and integers.

\apientry{wfdbxSetInputDir}{Set the input directory}
\begin{lstlisting}[language=Kotlin,style=signature]
fun Wfdb.wfdbxSetInputDir(dir: String?): String?
\end{lstlisting}
Sets the input directory, from which names that are not full paths are read.
\begin{params}
\item[dir] A local directory, a \code{\textasciitilde} path or a URL; \code{null} (or
  \code{""}) restores the default: \code{WFDBX\_INPUT\_DIR}, or else the current directory.
\end{params}
\returns{The previous input directory, or \code{null} if none was set.}
\raises{Nothing.}
\begin{lstlisting}[language=Kotlin]
val previous = wfdbxSetInputDir("/data/mitdb")
val rec = readRecord("100")             // reads /data/mitdb/100.hea and 100.dat
wfdbxSetInputDir(previous)
\end{lstlisting}

\apientry{wfdbxGetInputDir}{Return the input directory}
\begin{lstlisting}[language=Kotlin,style=signature]
fun Wfdb.wfdbxGetInputDir(): String?
\end{lstlisting}
Returns the input directory. The function has no parameters.
\returns{The input directory, or \code{null} if none is set (the default).}
\raises{Nothing.}
\begin{lstlisting}[language=Kotlin]
println(Wfdb().wfdbxGetInputDir())      // null, unless WFDBX_INPUT_DIR is set
\end{lstlisting}

\apientry{wfdbxSetOutputDir}{Set the output directory}
\begin{lstlisting}[language=Kotlin,style=signature]
fun Wfdb.wfdbxSetOutputDir(dir: String?): String?
\end{lstlisting}
Sets the output directory, in which files named without a full path are created.
\begin{params}
\item[dir] A local directory or a \code{\textasciitilde} path (a URL is refused, with a
  message); \code{null} (or \code{""}) restores the default: \code{WFDBX\_OUTPUT\_DIR}, or
  else the current directory.
\end{params}
\returns{The previous output directory, or \code{null} if none was set.}
\raises{Nothing.}
\begin{lstlisting}[language=Kotlin]
wfdbxSetOutputDir("out")
writeAnnotations("100", "qrs", longArrayOf(100, 460), symbols = listOf("N", "N"))   // out/100.qrs
\end{lstlisting}

\apientry{wfdbxGetOutputDir}{Return the output directory}
\begin{lstlisting}[language=Kotlin,style=signature]
fun Wfdb.wfdbxGetOutputDir(): String?
\end{lstlisting}
Returns the output directory. The function has no parameters.
\returns{The output directory, or \code{null} if none is set (the default).}
\raises{Nothing.}
\begin{lstlisting}[language=Kotlin]
val out = wfdbxGetOutputDir() ?: "."
\end{lstlisting}

\apientry{wfdbxSetNetwork}{Turn network access on or off}
\begin{lstlisting}[language=Kotlin,style=signature]
fun Wfdb.wfdbxSetNetwork(on: Boolean): Boolean
\end{lstlisting}
Turns network access (reading URLs) on or off. It is off unless the environment variable
\code{WFDBX\_NETWORK} is a nonzero number.
\begin{params}
\item[on] \code{true} to allow reading URLs, \code{false} to forbid it.
\end{params}
\returns{The previous setting.}
\raises{Nothing.}
\begin{lstlisting}[language=Kotlin]
wfdbxSetNetwork(true)
wfdbxSetInputDir("https://physionet.org/files/mitdb/1.0.0/")
val rec = readRecord("100", to = TimeSpec("0:05"))      // from PhysioNet
\end{lstlisting}

\apientry{wfdbxGetNetwork}{Return whether network access is on}
\begin{lstlisting}[language=Kotlin,style=signature]
fun Wfdb.wfdbxGetNetwork(): Boolean
\end{lstlisting}
Returns whether network access is on. The function has no parameters.
\returns{\code{true} if network access is on.}
\raises{Nothing.}
\begin{lstlisting}[language=Kotlin]
if (!wfdbxGetNetwork()) println("network access is off")
\end{lstlisting}

\apientry{wfdbxSetLegacy}{Turn legacy mode on or off}
\begin{lstlisting}[language=Kotlin,style=signature]
fun Wfdb.wfdbxSetLegacy(on: Boolean): Boolean
\end{lstlisting}
Turns legacy mode (the behavior of the original WFDB package, including its assumed
signal parameters and its search of the WFDB path, \cref{chap:behavior}) on or off, for
the whole \code{Wfdb} object. The parameter \code{legacy} of the high-level functions
turns it on for one call.
\begin{params}
\item[on] \code{true} to turn legacy mode on, \code{false} to turn it off.
\end{params}
\returns{The previous setting.}
\raises{Nothing.}
\begin{lstlisting}[language=Kotlin]
val previous = wfdbxSetLegacy(true)
val rec = readRecord("100s")            // as the original package reads it
wfdbxSetLegacy(previous)
\end{lstlisting}

\apientry{wfdbxGetLegacy}{Return whether legacy mode is on}
\begin{lstlisting}[language=Kotlin,style=signature]
fun Wfdb.wfdbxGetLegacy(): Boolean
\end{lstlisting}
Returns whether legacy mode is on. The function has no parameters.
\returns{\code{true} if legacy mode is on.}
\raises{Nothing.}
\begin{lstlisting}[language=Kotlin]
println(wfdbxGetLegacy())               // false, unless it was turned on
\end{lstlisting}

% ------------------------------------------------------------------------------------------
\subsection{The result classes}
\label{sec:kotlin-results}

The readers return objects of the classes \code{Record}, \code{Header} and
\code{Annotations}, whose properties are listed in
\cref{tab:kotlin-record,tab:kotlin-header,tab:kotlin-annotations}; \code{checkRecord}
returns a \code{CheckReport}, which holds \code{Issue} objects
(\cref{tab:kotlin-checkreport,tab:kotlin-issue}). The signals of a record are described
by \code{Siginfo} objects (\cref{tab:kotlin-siginfo}), which are those of the WFDB library
(\cref{sec:kotlin-library}). \code{Record}, \code{Header}, \code{Annotations},
\code{CheckReport}, \code{Issue} and \code{BaseDate} are data classes, whose properties
are read-only (\code{val}); \code{Record} and \code{Annotations} compare their arrays by
content in \code{equals}. A \code{BaseDate} has the properties \code{year}, \code{month}
(1--12) and \code{day} (1--31).

\begin{xltabular}{\linewidth}{@{}lL@{}}
\caption{The properties and methods of \code{Record}, the result of \code{readRecord}}
\label{tab:kotlin-record} \\
\toprule
\tabhead Property & \tabhead Contents \\
\midrule
\endfirsthead
\toprule
\tabhead Property & \tabhead Contents \\
\midrule
\endhead
\bottomrule
\endfoot
\code{name} & the record name, as given to \code{readRecord} (\code{String}) \\
\code{fs} & the sampling frequency in Hz (\code{Double}), or 0 if the header gives none;
  with \code{highres = true}, the highest sampling frequency of a multi-frequency
  record \\
\code{nsig} & the number of signals returned \\
\code{nsamples} & the number of samples per signal (the number of frames, or of sample
  vectors with \code{highres}): \code{samples.size / nsig} \\
\code{signalNumbers} & the number, in the record, of each signal returned
  (\code{IntArray}) \\
\code{samples} & the samples in ADC units, frame by frame (\code{IntArray}):
  \code{samples[i * nsig + s]} is sample \code{i} of signal \code{s}; invalid samples are
  \code{WFDB\_INVALID\_SAMPLE} \\
\code{physical} & with \code{physical = true}, the samples in physical units, in the
  layout of \code{samples} (\code{DoubleArray}; NaN for an invalid sample); otherwise
  \code{null} \\
\code{siginfo} & a \code{List<Siginfo>}, one per signal returned; \code{desc} has no white
  space at either end, and \code{adcres} is 0 if the header does not give it \\
\code{names} & the description of each signal, such as \code{"MLII"} (\code{""} if none),
  without the white space that a header may put at either end of it \\
\code{units} & the physical units of each signal, or \code{null} where the header gives
  none (\code{List<String?>}) \\
\code{gains} & the gain of each signal, in ADC units per physical unit
  (\code{List<Double>}); 0 means that the signal is uncalibrated \\
\code{baselines} & the ADC value of each signal that corresponds to 0 physical units
  (\code{List<Int>}) \\
\code{baseTime} & the start time and date as text, to the second, such as
  \code{"[14:00:03 21/06/2024]"}, or \code{null} if the header gives none \\
\code{baseTimeMs} & the start time in milliseconds after midnight, such as
  \code{50403500} for 14:00:03.500 (\code{Long}), or \code{null} if \code{baseTime} is
  \code{null} \\
\code{baseDate} & the start date as a \code{BaseDate(year, month, day)}, or \code{null}
  if the header gives none \\
\code{sampfrom} & the number of the first sample returned (\code{Long}) \\
\code{comments} & the info strings of the header (without the \code{\#}), and of the
  record's \file{.info} file, if any \\
\code{checksumOk} & \code{false} if the samples do not match the checksums of the
  header; checksums are checked only when a read reaches the end of the record, so a
  shorter read gives \code{true} \\
\code{legacy} & \code{true} if the record was read in legacy mode (a gain of 200 is then
  used for uncalibrated signals in physical units) \\
\code{at(i, s)} & sample \code{i} of signal \code{s}, in ADC units \\
\code{signal(s)} & all the samples of signal \code{s}, in ADC units (\code{IntArray});
  \code{IllegalArgumentException} (``no signal S'') if there is no signal \code{s} \\
\code{physicalAt(i, s)} & sample \code{i} of signal \code{s} in physical units,
  $(\mathit{adc} - \mathit{baseline}) / \mathit{gain}$ (NaN for an invalid sample);
  \code{WfdbException} (``signal N of record REC is uncalibrated'') if the signal has no
  gain \\
\code{physicalSignal(s)} & all the samples of signal \code{s} in physical units
  (\code{DoubleArray}), as \code{physicalAt} \\
\end{xltabular}

\begin{xltabular}{\linewidth}{@{}lL@{}}
\caption{The properties of \code{Header}, the result of \code{readHeader}}
\label{tab:kotlin-header} \\
\toprule
\tabhead Property & \tabhead Contents \\
\midrule
\endfirsthead
\toprule
\tabhead Property & \tabhead Contents \\
\midrule
\endhead
\bottomrule
\endfoot
\code{name} & the record name, as given to \code{readHeader} \\
\code{fs} & the sampling frequency in Hz, or 0 if the header gives none \\
\code{counterFreq} & the counter frequency, as \code{getcfreq()} returns it \\
\code{baseCounter} & the base counter value, as \code{getbasecount()} returns it \\
\code{nsamples} & the number of samples per signal (\code{Long}), or 0 if the header does
  not give it \\
\code{nsig} & the number of signals \\
\code{siginfo} & a \code{List<Siginfo>}, one per signal; \code{desc} has no white space at
  either end, and \code{adcres} is 0 if the header does not give it \\
\code{baseTime}, \code{baseTimeMs}, \code{baseDate} & the start time and date, as in
  \code{Record} \\
\code{comments} & the info strings of the header (without the \code{\#}) \\
\code{metadata} & the (key, value) pairs of the metadata entries of the header
  (\code{List<Pair<String, String>{}>}), in file order, as \code{readMetadata} returns
  them; empty for an EDF file or a header read from the standard input \\
\end{xltabular}

\begin{xltabular}{\linewidth}{@{}lL@{}}
\caption{The properties of \code{Annotations}, the result of \code{readAnnotations}.
  The annotations are in file order; each array or list has one entry per annotation.}
\label{tab:kotlin-annotations} \\
\toprule
\tabhead Property & \tabhead Contents \\
\midrule
\endfirsthead
\toprule
\tabhead Property & \tabhead Contents \\
\midrule
\endhead
\bottomrule
\endfoot
\code{record} & the record name, as given \\
\code{annotator} & the annotator name, such as \code{"atr"} \\
\code{fs} & the frequency in which \code{samples} are counted, normally the sampling
  frequency of the record; 0 if neither the record, nor the file, nor the argument
  \code{fs} of \code{readAnnotations} gives one \\
\code{samples} & the annotation times, in samples (\code{LongArray}) \\
\code{codes} & the annotation codes (\code{IntArray}), such as \code{NORMAL} or
  \code{PVC} \\
\code{symbols} & the mnemonics of the codes, such as \code{"N"}, \code{"V"} or \code{"+"}
  (\code{List<String>}) \\
\code{subtypes} & the subtype of each annotation (\code{IntArray}) \\
\code{chans} & the channel number of each annotation (\code{IntArray}) \\
\code{nums} & the annotator number of each annotation (\code{IntArray}) \\
\code{aux} & the auxiliary text of each annotation, or \code{null} where there is none
  (\code{List<String?>}) \\
\code{size} & the number of annotations \\
\end{xltabular}

\begin{table}[htbp]
\centering
\caption{The properties and methods of \code{CheckReport}, the result of
  \code{checkRecord}}
\label{tab:kotlin-checkreport}
\begin{tabularx}{\linewidth}{@{}lL@{}}
\toprule
\tabhead Property & \tabhead Contents \\
\midrule
\code{record} & the record name, as given \\
\code{ok} & \code{true} if no issue has the severity \code{Severity.ERROR} \\
\code{issues} & the \code{List<Issue>}, in the order of the checks \\
\code{lines()} & the text form of the report, as a list of lines: one line per issue
  (\code{issue.toString()}), or the one line \code{no problems found} \\
\code{toString()} & the lines of \code{lines()}, joined by line feeds (what
  \code{println(report)} prints) \\
\bottomrule
\end{tabularx}
\end{table}

\begin{table}[htbp]
\centering
\caption{The properties of \code{Issue}, a problem found by \code{checkRecord}}
\label{tab:kotlin-issue}
\begin{tabularx}{\linewidth}{@{}lL@{}}
\toprule
\tabhead Property & \tabhead Contents \\
\midrule
\code{severity} & a \code{Severity}: \code{ERROR}, \code{WARNING} or \code{INFO}; its
  \code{text} (and \code{toString()}) is \code{"error"}, \code{"warning"} or
  \code{"info"} \\
\code{code} & a stable identifier of the kind of problem, such as \code{"checksum"} or
  \code{"gain\_missing"} (\cref{tab:python-issue-codes}) \\
\code{where} & where the problem is: \code{""} (the record), \code{"signal N"},
  \code{"segment SEG"}, \code{"segment SEG, signal N"} or \code{"annotator ANN"} \\
\code{message} & the description of the problem \\
\code{toString()} & \code{"SEVERITY: WHERE: MESSAGE"}, or \code{"SEVERITY: MESSAGE"} when
  \code{where} is empty \\
\bottomrule
\end{tabularx}
\end{table}
Two \code{Issue} objects are equal (\code{==}) when their four properties are. The codes
of the issues are those of \cref{tab:python-issue-codes}.

\begin{xltabular}{\linewidth}{@{}lL@{}}
\caption{The properties of \code{Siginfo}, the specification of a signal (mutable, with
  C's zero values by default)}
\label{tab:kotlin-siginfo} \\
\toprule
\tabhead Property & \tabhead Contents \\
\midrule
\endfirsthead
\toprule
\tabhead Property & \tabhead Contents \\
\midrule
\endhead
\bottomrule
\endfoot
\code{fname} & the name of the signal file, as the header gives it (\code{ByteArray?}) \\
\code{desc} & the signal description, such as \code{"MLII"} (\code{ByteArray?}) \\
\code{units} & the physical units, such as \code{"mV"}, or \code{null} if not specified
  (\code{ByteArray?}) \\
\code{gain} & the gain, in ADC units per physical unit (\code{Double}; 0:
  uncalibrated) \\
\code{initval} & the value of the first sample \\
\code{group} & the signal group number (the signals of one group share one file) \\
\code{fmt} & the storage format, such as \code{16}, \code{212} or \code{516} \\
\code{spf} & the number of samples per frame (more than 1 for the faster signals of a
  multi-frequency record) \\
\code{bsize} & the block size, in bytes (0 for ordinary files) \\
\code{adcres} & the ADC resolution, in bits (0: not specified) \\
\code{adczero} & the ADC output for an input of 0\,V \\
\code{baseline} & the ADC value that corresponds to 0 physical units \\
\code{nsamp} & the number of samples (\code{Long}; 0 if unknown) \\
\code{cksum} & the 16-bit checksum of the samples \\
\end{xltabular}

% ------------------------------------------------------------------------------------------
\subsection{Warnings and errors}
\label{sec:kotlin-warnings}

\paragraph{Exceptions.} The high-level functions report errors by exceptions
(\cref{tab:kotlin-exceptions}), whose messages say what failed. Each function checks its
arguments before it creates or changes a file, so an invalid argument leaves the files as
they were. The library may also print a message on the standard error when it cannot open
a file (as in \code{init: can't open header for record REC}); \code{wfdbquiet()} turns
these messages off (\cref{sec:kotlin-library}).

\begin{xltabular}{\linewidth}{@{}lL@{}}
\caption{The exceptions of the high-level API}
\label{tab:kotlin-exceptions} \\
\toprule
\tabhead Exception & \tabhead Thrown when \\
\midrule
\endfirsthead
\toprule
\tabhead Exception & \tabhead Thrown when \\
\midrule
\endhead
\bottomrule
\endfoot
\code{WfdbException} & every error of a high-level function: a header, signal file or
  annotation file cannot be found, read, created or written completely (for example on a
  full disk); a signal or annotation file ends before its header or its end-of-file marker
  says; an argument is invalid, or the operation needs a value that the record does not
  give (a time in seconds without a sampling frequency, physical units without a gain);
  the directory given as \code{dir} does not exist \\
\code{WfdbArgumentException} & a subclass of \code{WfdbException}, thrown by the metadata
  functions (and for the \code{metadata} argument of \code{writeRecord}) for an invalid
  argument: an invalid key or value, a key given twice, a header with two entries for the
  key to set, a header without a record line, a header read from a URL without
  \code{dir}, an empty record name or \code{"-"} (Python's \code{ValueError}); a header
  that cannot be read or written, or a \code{dir} that does not exist, is a plain
  \code{WfdbException} (Python's \code{IOError}) \\
\end{xltabular}

\paragraph{The checksum warning.} When \code{readRecord} reads to the end of a record
whose samples do not match the checksums of its header, it returns all the samples and
warns: ``checksum error in record REC: the signal data do not match the checksums in its
header''. \code{checksumWarning = false} turns the warning off; \code{checksumOk} of the
result is \code{false} either way.

\paragraph{The duplicate-key warning.} When a header holds two metadata entries whose keys
differ only in letter case (written by other software), \code{readMetadata},
\code{getMetadata} and \code{readHeader} use the first and warn: ``header of record REC
has more than one entry for key KEY (lines L1, L2); the first is used''.
\code{setMetadata} and \code{updateMetadata} then refuse to change that key
(\code{WfdbArgumentException}), and \code{removeMetadata} removes all of its entries.

Both warnings go to the function given as \code{onWarning}, which receives the message;
by default, the library prints \code{WFDB-x: } followed by the message on the standard
error, through \code{wfdb\_error}, so that \code{wfdbquiet()} silences it and
\code{wfdberror()} returns it:

\begin{lstlisting}[language=Kotlin,style=example]
import wfdbx.*

// A record whose last sample is changed after it was written: the checksum fails.
writeRecord("ck", IntArray(100) { it }, nsig = 1, fs = 100.0, gains = doubleArrayOf(100.0),
            adcres = intArrayOf(12), units = listOf("mV"))
java.io.RandomAccessFile(wfdbxGetOutputDir() + "/ck.dat", "rw").use { f ->
    f.seek(f.length() - 2)
    f.write(byteArrayOf(0, 0))
}
val caught = ArrayList<String>()
var rec = readRecord("ck", onWarning = { caught.add(it) })
println("${rec.nsamples} ${rec.checksumOk}")    // 100 false: all the samples
println(caught[0])
// checksum error in record ck: the signal data do not match the checksums in its header
rec = readRecord("ck", checksumWarning = false)      // no warning

// Two entries for one key: the first is used, with a warning.
writeRecord("dup", IntArray(10), nsig = 1, fs = 100.0, gains = doubleArrayOf(0.0),
            adcres = intArrayOf(12), comments = listOf(" Age: 63", " AGE: 64"))
println(getMetadata("dup", "age", onWarning = { caught.add(it) }))     // 63
println(caught[1])
// header of record dup has more than one entry for key Age (lines 3, 4); the first is used
try {
    setMetadata("dup", "Age", "65")
} catch (e: WfdbArgumentException) {
    println("WfdbArgumentException: ${e.message}")
}
println(removeMetadata("dup", "age"))       // 2
\end{lstlisting}

\begin{lstlisting}[language=Kotlin,style=example]
import wfdbx.*

wfdbquiet()                           // no library messages on the standard error
try {
    readRecord("nosuchrecord")
} catch (e: WfdbException) {
    println("WfdbException: ${e.message}")    // WfdbException: can't read record nosuchrecord
}
wfdbverbose()

writeRecord("raw", IntArray(10) { it }, nsig = 1, fs = 100.0, gains = doubleArrayOf(0.0),
            adcres = intArrayOf(12))          // uncalibrated
try {
    readRecord("raw", physical = true)
} catch (e: WfdbException) {
    println("WfdbException: ${e.message}")    // WfdbException: signal 0 of record raw is uncalibrated
}
\end{lstlisting}

% ==========================================================================================
\section{The WFDB library functions in Kotlin}
\label{sec:kotlin-library}

Every function of the WFDB library is available under its C name, with the same
arguments, in the same order, and the same return values, so that programs and
documentation written for the C library carry over directly. The functions, their
arguments and their results are documented once, for all the languages, in
\cref{chap:library}; this section gives the rules by which the C declarations become
Kotlin functions. The KDoc comment of each function documents it, including what WFDB-x
changes in it.

\subsection{Conventions}

\paragraph{The library state.} All the state that the C library keeps in global variables
(the WFDB path, the open record, its signals, the annotators, the calibration list, the
settings) is held by a \code{Wfdb} object. Each C function is an extension function of
\code{Wfdb} with the C name and arguments, \code{fun Wfdb.isigopen(record, siarray,
nsig)}, called as \code{w.isigopen(...)}, and also a top-level function with the same name
and arguments, \code{isigopen(...)}, which uses the default instance of the calling thread
(\code{wfdbDefault()}), for line-by-line ports of C programs. Separate \code{Wfdb} objects
share no state, so several records can be open at once, in one thread or in several; a
\code{Wfdb} object is used by one thread at a time (it may be handed from one thread to
another). As in C, the open record, its input signals, the output signals and the
annotators remain open between calls, until \code{wfdbquit()} closes them;
\code{wfdbquit()} leaves the input and output directories and the network setting as they
are. There are no destructors: an output file is completed by \code{wfdbquit()} (or the
function that closes it), and \code{wfdbx.cstd.exit} flushes the open output files as C's
\code{exit} does, but writes no annotation end marks and no header updates.

\paragraph{Names.} The functions, the constants (\code{WFDB\_READ},
\code{WFDB\_INVALID\_SAMPLE}, \code{WFDB\_HIGHRES}, \ldots) and the annotation codes
(\code{NORMAL}, \code{PVC}, \ldots, \cref{chap:formats}) keep their C names. The functions
of \file{<wfdb/ecgmap.h>} are provided under their macro names (\code{isqrs},
\code{map1}, \code{map2}, \ldots).

\paragraph{Strings are bytes.} WFDB files are not UTF-8: every string of the C-style API
is a \code{ByteArray}, which, like a C string, ends at its first NUL byte, and C's
\code{NULL} is \code{null}. Where natural, a function also has an overload that takes a
\code{String}, one character per byte (ISO-8859-1): \code{"..".latin1()} and
\code{bytes.latin1()}, from \code{wfdbx.cstd}, convert. Where both overloads exist, a
\code{null} argument is written \code{null as ByteArray?}, as in
\code{sampfreq(null as ByteArray?)}. Functions that return a C string (\code{timstr},
\code{mstimstr}, \code{datstr}, \code{getinfo}, \ldots) return a new \code{ByteArray} (a
copy of C's static buffer); \code{getinfo} returns \code{null} after the last info
string.

\paragraph{Numbers.} C's \code{int} and \code{WFDB\_Sample} are \code{Int},
\code{long}, \code{WFDB\_Time} and \code{WFDB\_Date} are \code{Long}, and \code{double}
and \code{WFDB\_Frequency} are \code{Double}. \code{WFDB\_Signal} and \code{WFDB\_Group}
are \code{Int}; negative values are invalid signal numbers, as C's unsigned values are.
The \code{char} fields of the structures are \code{Int} properties that keep C's range:
\code{anntyp}, \code{subtyp} and \code{num} from $-128$ to 127, \code{chan} from 0 to
255.

\paragraph{Arrays and output arguments.} Where a C function fills an array that the caller
passes, the Kotlin function fills an array that the caller passes:
\begin{itemize}
\item \code{isigopen(record, siarray, nsig)} (also \code{osigopen} and \code{wfdbinit})
  takes an \code{Array<Siginfo>?}, whose elements it fills (allocated by the caller:
  \code{Array(n) \{ Siginfo() \}}); \code{null} is C's \code{NULL}. With \code{nsig}
  $\le 0$, \code{isigopen} reads the header only, and returns the number of signals of the
  record.
\item \code{getvec(v)}, \code{getframe(v)} and \code{putvec(v)} take an \code{IntArray};
  a shorter array receives what fits.
\item A C function that fills a structure through a pointer fills the object passed:
  \code{getann(n, annot)} fills an \code{Annotation}, \code{getcal(desc, units, cal)} a
  \code{Calinfo}. \code{putann(n, annot)} and \code{putcal(cal)} read them.
\item An array of structures passed in is an array of objects, as in
  \code{osigfopen(siarray, nsig)} and \code{annopen(record, aiarray, nann)}, where
  \code{aiarray} is, for example, \code{arrayOf(Anninfo("atr", WFDB\_READ))}.
\item \code{getseginfo(sarray)} takes a \code{MutableList<Seginfo>?}, which it clears and
  fills with copies of the segments; \code{setmsheader(record, names, n)} takes an
  \code{Array<ByteArray>} or an \code{Array<String>}.
\end{itemize}

\paragraph{Return values and errors.} The functions return what the C functions return:
an \code{Int} status (negative on errors), a number of signals, a \code{Double}, a byte
string or \code{null}. Errors are reported by these values and by a message on the
standard error, as in C, never by exceptions. \code{wfdberror()} returns the last message;
\code{wfdbquiet()} turns the messages off and \code{wfdbverbose()} on again.
\code{wfdbx\_write\_errors()} returns the number of write errors since the program
started (output files that could not be written, flushed or closed),
\code{wfdbx\_note\_write\_error()} counts one more, and
\code{wfdbx\_fclose\_output(file, pname, name)} closes a file that a program writes itself,
reporting ``PNAME: write error on NAME'' and counting it.
\code{wfdbx\_input\_path(name)} and \code{wfdbx\_output\_path(name)} return the path under
which a program opens a file that it reads or creates itself, where the library would
open a file of that name: under the input or output directory (the missing directories
of an output path are created), a \code{\textasciitilde} name expanded, a full path as
given; \code{null} for a URL, and \code{"-"} as given. In legacy mode, an input name is
used as given, and the output directory, if set, is a prefix of any output name.
\code{wfdbquit()} closes the open record and annotators, and keeps the directory and
network settings, the time resolution of output annotators (\code{getafreq()}) and the
getvec mode (\code{getgvmode()}), as in C; the high-level functions restore the last two
(\cref{sec:kotlin-conventions}).

\paragraph{Additions.} \code{getvecs(out, n)} reads up to \code{n} sample vectors at once
into the \code{IntArray} \code{out}, vector after vector, as \code{n} calls of
\code{getvec} would, but decoding whole blocks of the signal files; it returns the number
of vectors read, or the negative status of \code{getvec} if the first read failed.
\code{getvecs\_status(out, n)} returns a \code{GetvecsResult(count, status)}, with the
value returned by the last \code{getvec}. The C standard library functions that the WFDB
library relies on are in \code{wfdbx.cstd}: \code{cformat} and \code{sprintf} (to a
\code{ByteArray}), \code{printf} and \code{eprintf}, which format numbers exactly as the
GNU C library does, \code{strtol}, \code{strtod}, \code{sscanf}, \ldots

\paragraph{Settings.} The settings functions of WFDB-x (\code{wfdbx\_set\_input\_dir},
\code{wfdbx\_get\_input\_dir}, \code{wfdbx\_set\_output\_dir},
\code{wfdbx\_get\_output\_dir}, \code{wfdbx\_set\_network}, \code{wfdbx\_get\_network},
\code{wfdbx\_set\_legacy}, \code{wfdbx\_get\_legacy}, \code{wfdbx\_legacy\_args}) keep
their C names and types (byte strings, \code{Int} for the switches); the high-level
functions of \cref{sec:kotlin-settings} are their Kotlin forms.

\subsection{Package layout}

\begin{table}[htbp]
\centering
\caption{The packages of the Kotlin library}
\label{tab:kotlin-packages}
\begin{tabularx}{\linewidth}{@{}lL@{}}
\toprule
\tabhead Package or file & \tabhead Contents \\
\midrule
\code{wfdbx}: \file{Wfdb.kt} & the class \code{Wfdb}, \code{wfdbDefault()},
  \code{withDefault} \\
\file{Types.kt} & the constants and structures of \file{wfdb.h} \\
\file{WfdbFile.kt}, \file{Wfdbio.kt} & files, the input and output directories, network
  access, error messages (\file{lib/wfdbio.c}) \\
\file{Signal.kt}, \file{SignalTime.kt}, \file{Flac.kt} & headers, signal files in every
  format, time and date conversions (\file{lib/signal.c}), the FLAC codec \\
\file{Annot.kt}, \file{Calib.kt}, \file{Wfdbinit.kt}, \file{Ecgcodes.kt} & annotation files
  (\file{lib/annot.c}), the calibration list (\file{lib/calib.c}), \code{wfdbinit}, the
  annotation codes and \file{ecgmap.h} \\
\file{Record.kt}, \file{Check.kt}, \file{Metadata.kt} & the high-level API
  (\cref{sec:kotlin-api}) \\
\code{wfdbx.cstd} & the C standard library emulation (\code{printf}, \code{scanf},
  \code{strtol}, \code{rand}, the standard streams, \code{exit}, \code{latin1}) and the
  platform layer \\
\code{wfdbx.apps} & the command-line tools, one file per tool, as functions
  \code{fun Wfdb.rdsampMain(argv: Array<ByteArray>): Int}, \ldots \\
\code{wfdbx.cli} & the command-line driver \code{wfdbx.cli.MainKt} (JVM) \\
\bottomrule
\end{tabularx}
\end{table}

\subsection{The structures}
\label{sec:kotlin-structures}

The structures of the C library are mutable classes whose properties have the names of
the fields of the C structures and C's zero values by default
(\cref{tab:kotlin-structures}). \code{copy()} is C's structure assignment, and
\code{set(o)} copies another object into an existing one (what a C function filling
\code{*p} does); \code{equals} compares contents.
\begin{lstlisting}[language=Kotlin]
Siginfo(fname: ByteArray? = null, desc: ByteArray? = null, units: ByteArray? = null, gain: Double = 0.0,
        initval: Int = 0, group: Int = 0, fmt: Int = 0, spf: Int = 0, bsize: Int = 0, adcres: Int = 0,
        adczero: Int = 0, baseline: Int = 0, nsamp: Long = 0, cksum: Int = 0)
Seginfo(recname: ByteArray = ByteArray(0), nsamp: Long = 0, samp0: Long = 0)
Anninfo(name: ByteArray = ByteArray(0), stat: Int = 0)      // also Anninfo(name: String, stat: Int)
Annotation(time: Long = 0, anntyp: Int = 0, subtyp: Int = 0, chan: Int = 0, num: Int = 0, aux: ByteArray? = null)
Calinfo(low: Double = 0.0, high: Double = 0.0, scale: Double = 0.0, sigtype: ByteArray? = null,
        units: ByteArray? = null, caltype: Int = 0)
\end{lstlisting}

\begin{xltabular}{\linewidth}{@{}lL@{}}
\caption{The classes of the C structures}
\label{tab:kotlin-structures} \\
\toprule
\tabhead Class (C structure) & \tabhead Properties \\
\midrule
\endfirsthead
\toprule
\tabhead Class (C structure) & \tabhead Properties \\
\midrule
\endhead
\bottomrule
\endfoot
\code{Siginfo} (\code{WFDB\_Siginfo}) & \code{fname}, \code{desc}, \code{units},
  \code{gain}, \code{initval}, \code{group}, \code{fmt}, \code{spf}, \code{bsize},
  \code{adcres}, \code{adczero}, \code{baseline}, \code{nsamp}, \code{cksum}
  (\cref{tab:kotlin-siginfo}) \\
\code{Seginfo} (\code{WFDB\_Seginfo}) & \code{recname} (the name of the segment),
  \code{nsamp} (its number of samples), \code{samp0} (the number of its first sample in
  the record) \\
\code{Anninfo} (\code{WFDB\_Anninfo}) & \code{name} (the annotator name, the suffix of the
  file), \code{stat} (\code{WFDB\_READ} or \code{WFDB\_WRITE} for the MIT format,
  \code{WFDB\_AHA\_READ} or \code{WFDB\_AHA\_WRITE} for the AHA format) \\
\code{Annotation} (\code{WFDB\_Annotation}) & \code{time} (in sample intervals),
  \code{anntyp} (the annotation code), \code{subtyp}, \code{chan}, \code{num},
  \code{aux} (\code{null}, or bytes whose first byte is the length of the text that
  follows, as in C) \\
\code{Calinfo} (\code{WFDB\_Calinfo}) & \code{low}, \code{high} (the levels of the
  calibration pulse, in physical units), \code{scale} (physical units per cm),
  \code{sigtype} (matched against the beginning of a signal description), \code{units},
  \code{caltype} \\
\end{xltabular}

\subsection{Examples}

Reading the samples and the annotations of a record, as a C program does, on a
\code{Wfdb} object of its own (the input directory is that of the examples):

\begin{lstlisting}[language=Kotlin,style=example]
import wfdbx.*
import wfdbx.cstd.latin1

val w = Wfdb()                              // a library state of its own
w.wfdbx_set_input_dir(wfdbxGetInputDir()!!)
val si = Array(2) { Siginfo() }             // filled with the signal specifications
val nsig = w.isigopen("100s", si, 2)        // open 2 signals
println("$nsig ${si.map { "${it.desc?.latin1()} ${it.fmt} ${it.gain}" }}")   // 2 [MLII 212 200.0, V5 212 200.0]
println(w.sampfreq(null as ByteArray?))     // 360.0: the frequency of the open record

val v = IntArray(nsig)                      // receives one sample per signal
w.isigsettime(w.strtim("0:30"))             // skip to 30 seconds
w.getvec(v)
println("${v.toList()} ${(0 until nsig).map { w.aduphys(it, v[it]) }}")   // [947, 968] [-0.385, -0.28]

val frames = IntArray(360 * nsig)
println(w.getvecs(frames, 360))             // 360: the next second, frame by frame

val a = Annotation()
w.annopen("100s", arrayOf(Anninfo("atr", WFDB_READ)), 1)
var n = 0
while (w.getann(0, a) == 0) {               // 0: an annotation was read; -1: the end
    if (n < 2) println("${w.mstimstr(a.time).latin1().trim()} ${w.annstr(a.anntyp).latin1()}")   // 0:00.050 +, 0:00.214 N
    n++
}
println("$n annotations")                   // 75 annotations

w.wfdbquiet()                               // no messages on the standard error
println("${w.isigopen("nosuchrecord", null, 0)} ${w.wfdberror().latin1().trim()}")
// -1 init: can't open header for record nosuchrecord
w.wfdbverbose()
w.wfdbquit()                                // close everything
\end{lstlisting}

Writing a record and an annotation file, sample by sample and annotation by annotation:

\begin{lstlisting}[language=Kotlin,style=example]
import wfdbx.*
import wfdbx.cstd.latin1

val w = Wfdb()
w.wfdbx_set_output_dir(wfdbxGetOutputDir()!!)
val out = arrayOf(Siginfo(fname = "ramp.dat".latin1(), desc = "ramp".latin1(), units = "mV".latin1(),
                          gain = 200.0, fmt = 16, adcres = 12))
println(w.osigfopen(out, 1))                // 1: one output signal
w.setsampfreq(250.0)
for (k in 0 until 500) w.putvec(intArrayOf(k))
println(w.newheader("ramp"))                // 0: ramp.hea is written
w.annopen("ramp", arrayOf(Anninfo("qrs", WFDB_WRITE)), 1)
for (t in longArrayOf(100, 350)) w.putann(0, Annotation(time = t, anntyp = NORMAL))
w.wfdbquit()                                // completes ramp.dat and ramp.qrs

println(readAnnotations("ramp", "qrs").samples.toList())    // [100, 350]
\end{lstlisting}

% ==========================================================================================
\section{The command-line tools in Kotlin}
\label{sec:kotlin-tools}

The 33 command-line tools (\code{rdsamp}, \code{rdann}, \code{wrsamp}, \code{xform},
\ldots) run on the JVM, through the launcher \file{kotlin/bin/wfdbx} (a POSIX shell
script), followed by the name of the tool and its arguments. The tools, their options,
their output and their exit status are those of the C tools, and are documented in
\cref{chap:tools}. The launcher needs the jar built by \cmd{./gradlew fatJar}
(\cref{sec:kotlin-install}). In the directory \file{kotlin}:

\begin{lstlisting}[style=shell]
export WFDBX_INPUT_DIR=../tests/data
bin/wfdbx rdsamp -r 100s -t 10 -p      # the first 10 seconds of 100s, in physical units
bin/wfdbx rdann -r 100s -a atr         # its annotations
bin/wfdbx rdsamp --legacy -r 100s      # legacy mode: --legacy as the first argument
\end{lstlisting}

The command also takes these options, instead of a tool:

\begin{table}[htbp]
\centering
\caption{The options of the launcher \code{bin/wfdbx}}
\label{tab:kotlin-wfdbx-options}
\begin{tabularx}{\linewidth}{@{}lL@{}}
\toprule
\tabhead Command & \tabhead Effect \\
\midrule
\cmd{bin/wfdbx TOOL [ARGUMENTS ...]} & runs the tool \code{TOOL}, and exits with its
  status \\
\cmd{bin/wfdbx --list} & lists the names of the tools, one per line \\
\cmd{bin/wfdbx --version} & prints the version, \code{WFDB-x 1.2.0} \\
\cmd{bin/wfdbx --help}, \cmd{bin/wfdbx -h} & prints the usage \\
\bottomrule
\end{tabularx}
\end{table}

Without arguments, the command prints the usage and exits with status 1; an unknown tool
is an error (``wfdbx: unknown tool 'TOOL' (try 'wfdbx --list')'', status 1), and an
uncaught exception is reported as \code{wfdbx: error: MESSAGE} (status 1). When an output
file or the standard output cannot be written (a full disk, for example), a tool reports
it and exits with status 1 (legacy mode keeps the original status, 0).

The files that a tool reads or creates itself are found where the library finds its own
files: the \opt{-i} file of \code{wrsamp}, \code{edf2mit}, \code{wav2mit}, \code{a2m} and
\code{ad2m}, the input files of \code{ahaecg2mit}, the \opt{-S} script of \code{xform} and
the signal file that \code{wfdbdesc} measures are read under the input directory, and the
WAV file of \code{mit2wav}, the EDF file of \code{mit2edf}, the master header of
\cmd{snip -m} and the \file{.des}, \file{.ecg}, \file{.key} and \file{.txt} files of
\code{parsescp} are created under the output directory (with
\code{wfdbx\_input\_path} and \code{wfdbx\_output\_path}, \cref{sec:kotlin-library}). Such
a file cannot be read from a URL. In legacy mode, they are read as named, relative to the
current directory, and the output directory, if set, is a prefix of their names.

\paragraph{The launcher.} \file{bin/wfdbx} runs \cmd{java -cp build/libs/wfdbx-all.jar
wfdbx.cli.MainKt TOOL ...}, or the jar named by the environment variable
\code{WFDBX\_KOTLIN\_JAR}. It starts the JVM with options for a short start-up (its first
compiler only, the serial collector, and the class data sharing archive
\file{wfdbx-all.jsa} that \cmd{./gradlew fatJar} makes next to the jar), and moves
\code{JAVA\_TOOL\_OPTIONS} to the command line. A link to \file{bin/wfdbx} named after a
tool runs that tool. Without the launcher, the jar runs the tools directly:

\begin{lstlisting}[style=shell]
java -cp build/libs/wfdbx-all.jar wfdbx.cli.MainKt rdsamp -r 100s -t 10
\end{lstlisting}

\paragraph{The tools as functions.} Each tool is a function of the package
\code{wfdbx.apps}, such as
\begin{lstlisting}[language=Kotlin]
fun Wfdb.rdsampMain(argv: Array<ByteArray>): Int
\end{lstlisting}
where \code{argv} is C's \code{argv}, the program name first, and the result is the exit
status. A tool may end the process (with
\code{exit}, as C programs do), so an app uses the high-level functions or the library
functions instead.

\paragraph{Environment variables.} The tools, like the rest of the library, use the
environment variables of \cref{tab:kotlin-environment}. On the JVM, which cannot change
its own environment, \code{setenv} (in \code{wfdbx.cstd}) keeps values that the library
and the processes it starts see.

\begin{table}[htbp]
\centering
\caption{The environment variables used by the Kotlin library}
\label{tab:kotlin-environment}
\begin{tabularx}{\linewidth}{@{}lL@{}}
\toprule
\tabhead Variable & \tabhead Meaning \\
\midrule
\code{WFDBX\_INPUT\_DIR} & the input directory (a directory, a \code{\textasciitilde}
  path or a URL), when \code{wfdbxSetInputDir} does not set one \\
\code{WFDBX\_OUTPUT\_DIR} & the output directory, when \code{wfdbxSetOutputDir} does not
  set one \\
\code{WFDBX\_NETWORK} & network access is on if it is a nonzero number, when
  \code{wfdbxSetNetwork} does not set it \\
\code{WFDBANNSORT} & \code{0} to keep annotation files in the order written (no sorting
  when they are closed) \\
\code{WFDBGVMODE} & the default getvec mode (\code{setgvmode} with a negative mode) \\
\code{WFDBCAL} & the calibration file read by \code{calopen(null)} (default
  \file{wfdbcal}) \\
\code{WFDB} & the WFDB path, used in legacy mode only \\
\code{WFDB\_FLAC\_CHECK\_MD5} & if set, the MD5 signature of FLAC signal files is
  verified when they are read \\
\code{WFDBX\_KOTLIN\_JAR} & the jar that \file{bin/wfdbx} runs (default
  \file{build/libs/wfdbx-all.jar}) \\
\bottomrule
\end{tabularx}
\end{table}

% ==========================================================================================
\section{Testing}
\label{sec:kotlin-testing}

The unit tests of the Kotlin library are in \file{kotlin/src/commonTest} (run on every
target) and \file{kotlin/src/jvmTest} (JVM only). In the directory \file{kotlin}:

\begin{lstlisting}[style=shell]
./gradlew jvmTest                           # the unit tests, on the JVM
./gradlew compileCommonMainKotlinMetadata   # the common code uses only common APIs
./gradlew jvmTest -Pwfdbx.linux=true linuxX64Test   # also on Linux/Native, with its toolchain
\end{lstlisting}

The unit tests check the library functions with expected values printed by the C
library, the high-level functions, the bulk paths against the per-frame paths, the header
metadata functions against the cases shared by the seven languages (in
\file{tests/metadata}: results, warnings and the headers, byte for byte), the examples of
\file{kotlin/README.md}, and \code{checkRecord} against the reports of the Python
implementation on a corpus of damaged records (this test is skipped, with a message, when
\code{python3} with NumPy is not available). Network access is tested against a local
HTTP server (127.0.0.1). The tests copy \file{tests/data} to a scratch directory, and
never change the current directory.

The cross-language conformance suite compares the output of the tools with that of the
original C programs, case by case: 1135 cases, and 883 in legacy mode, which runs each
tool with \opt{--legacy}. It runs the launcher \file{kotlin/bin/wfdbx}, so the jar is
built first. From the root of the repository:

\begin{lstlisting}[style=shell]
(cd kotlin && ./gradlew fatJar)
PYTHONPATH=$PWD/python/src python3 tests/conformance/harness.py check --impl kotlin -j8
PYTHONPATH=$PWD/python/src python3 tests/conformance/harness.py check --impl kotlin -j8 --legacy
\end{lstlisting}


% ---- Part III: Reference -----------------------------------------------------------------
\part{Reference}
\label{part:reference}
% The WFDB library functions: the reference of the C-style API of WFDB-x.
% Part of the WFDB-x user's manual (see ../README.md for the conventions).
% check_manual.py checks that every library function of the Python
% implementation (the C-named functions and the wfdbx_* additions) has
% exactly one \apientry here, and that every entry is one of them.

\chapter{The WFDB library functions}
\label{chap:library}
% (long identifiers: allow more stretch in this chapter; reset at its end)
\setlength{\emergencystretch}{2.5em}

Besides its high-level functions, which read and write whole records in the
idiom of each language, WFDB-x provides the functions of the C library of the
WFDB Software Package \wfdbversion{} (\code{libwfdb}, declared in
\file{<wfdb/wfdb.h>}) in all seven languages. They keep the C names
(\code{isigopen}, \code{getvec}, \code{getann}, \code{strtim}, ...), the C
arguments and their order, the C return values and error codes, the messages
that the C library prints, and its global state: the open record, signals and
annotators, the sampling frequency and base time used by the time conversion
functions, the calibration list, the tables of annotation mnemonics. A program
written for the C library, or a passage of the \emph{WFDB Programmer's Guide},
carries over call by call; the command-line tools of WFDB-x
(\cref{chap:tools}) are written with these functions, as the C tools are with
\code{libwfdb}.

The behavior is that of the C library, edge cases included, with the
deliberate changes of WFDB-x described in \cref{chap:behavior}: no built-in
sampling frequency, gain, ADC resolution or units; files found only where
they are named (an input directory, an output directory, network access off
by default) instead of being searched for in the WFDB path; out-of-order
annotation files sorted in the same process, without losing annotations; an
output signal file never one of the input files; write errors counted. Legacy
mode (\code{wfdbx\_set\_legacy}) restores the original behavior. Each entry
below says what changes, if anything. WFDB-x adds a few functions of its own,
whose names begin with \code{wfdbx\_} (\cref{sec:library-additions}), and a
reader of many sample vectors at once (\code{getframes}, or \code{getvecs}).

Since the names, the arguments and the behavior are the same in the seven
languages, each function is documented once, in C notation: the first listing
of each entry is the prototype of the C library (for the WFDB-x additions,
the declaration of the C++ implementation, \file{cpp/include/wfdbx/wfdb.hpp},
which is in the same style). \Cref{sec:library-calling} shows how the C
conventions (arrays filled by a function, output arguments, \code{NULL},
strings, structures) appear in each language. The record \code{100s} used in
the examples is the sample record included with WFDB-x (the first minute of
record 100 of the MIT-BIH Arrhythmia Database, in \file{tests/data}).

\begin{note}
These functions are for programs that need the C library's model: reading
sample by sample, writing headers and annotation files piece by piece, or
porting an existing C program. To read or write a whole record or annotation
file, the high-level functions of each language (\code{read\_record},
\code{write\_annotations}, ...; \cref{chap:python,chap:cpp,chap:rust,chap:julia,chap:r,chap:matlab,chap:kotlin})
are simpler and faster.
\end{note}

% ==========================================================================
\section{Calling the library functions from each language}
\label{sec:library-calling}

In every language the functions have their C names and take their arguments
in the C order. What differs is how a language expresses C's pointers:
arrays that a function fills, output arguments, \code{NULL}, and strings.
\Cref{tab:library-languages} summarizes these conventions; the notes and the
example that follow give the details for each language. Each example opens
the record \code{100s} (from the current directory, the default input
directory) with \code{isigopen}, reads one sample of each signal with
\code{getvec}, and prints the description of signal 0 and the two samples
(\code{MLII 995 1011}).

The conventions common to all languages:

\begin{itemize}
\item Signal, group and annotator numbers start at 0, as in C, also in the
  languages whose arrays start at 1 (Julia, R, MATLAB/Octave).
\item Errors are reported as in C: by the return value, with a message on
  the standard error output, not by exceptions. \code{wfdberror} returns the
  last message, and \code{wfdbquiet} stops the messages from being printed.
\item The state of the library is kept between calls, as in C.
  \code{wfdbquit} closes the files and resets the record's state; the input
  and output directories and the network setting are kept.
\item Text is handled as bytes, one character per byte (ISO-8859-1 where a
  language's strings are Unicode), so that every byte of a header or an
  annotation file is kept.
\item An annotation's \code{aux} field is either absent (C's \code{NULL}) or a
  sequence of bytes whose first byte is the length of the data that follow,
  as in C.
\end{itemize}

\begin{xltabular}{\linewidth}{@{}P{1.6cm}LLLL@{}}
\caption{The C conventions in each language.}
\label{tab:library-languages}\\
\toprule
{\tabhead Language} & {\tabhead Calling} & {\tabhead Arrays and output
arguments} & {\tabhead \code{NULL}} & {\tabhead Strings returned}\\
\midrule
\endfirsthead
\toprule
{\tabhead Language} & {\tabhead Calling} & {\tabhead Arrays and output
arguments} & {\tabhead \code{NULL}} & {\tabhead Strings returned}\\
\midrule
\endhead
\bottomrule
\endfoot
Python & \code{wfdbx.isigopen(...)} (module \code{wfdbx}) & a list passed by
the caller is filled; an object (\code{Annotation}, \code{Calinfo}) passed by
the caller is filled & \code{None} & \code{str}; \code{aux} is \code{bytes}\\
C++ & \code{wfdbx::\allowbreak isigopen(...)} (namespace \code{wfdbx}, header
\file{<wfdbx/wfdb.hpp>}) & as in C: arrays and structures allocated by the
caller, passed by pointer & \code{nullptr} & \code{std::string}, or
\code{std::optional<\allowbreak std::string>} where C may return
\code{NULL}\\
Rust & \code{wfdbx::\allowbreak isigopen(...)}, or the method \code{w.isigopen(...)} of a
\code{Wfdb} value & a mutable slice or reference passed by the caller is
filled & \code{None} (\code{wfdbx::NULL}) & \code{Vec<u8>}, or
\code{Option<Vec<u8>>} where C may return \code{NULL}\\
Julia & \code{isigopen(...)} after \code{using WFDBX}, or
\code{WFDBX.\allowbreak isigopen(...)}; optionally a \code{Wfdb} object first & a vector
or mutable structure passed by the caller is filled; a \code{Ref} receives a
pointer result & \code{nothing} & \code{Vector\{UInt8\}}; \code{String} for
\code{timstr}, \code{mstimstr}, \code{datstr}\\
R & \code{isigopen(...)} (package \code{wfdbx}) & returned with the C result
in a list (\code{r\$nsig}, \code{r\$si}; \code{g\$stat}, \code{g\$v}) & \code{NULL}
& character strings; \code{aux} is a raw vector\\
MATLAB/ Octave & \code{wfdbx.isigopen(...)} (package \code{+wfdbx}) & returned as
extra output values (\code{[nsig, si] = ...}) & \code{[]} & \code{char} row
vectors; \code{aux} is \code{uint8}\\
Kotlin & \code{isigopen(...)} (package \code{wfdbx}), or the extension
function \code{w.isigopen(...)} of a \code{Wfdb} object & an array or object
passed by the caller is filled & \code{null} & \code{ByteArray}, or
\code{ByteArray?} where C may return \code{NULL}\\
\end{xltabular}

\subsection{Python}
\label{sec:library-python}

The functions, the structures and the constants are in the package
\code{wfdbx} (\code{import wfdbx}). The C structures are the classes
\code{Siginfo}, \code{Anninfo}, \code{Annotation}, \code{Calinfo} and
\code{Seginfo}, with the C field names; their constructors take the fields
as keyword arguments, with C's zero values, except that \code{Siginfo.spf} is
1 by default.

\begin{itemize}
\item Where C fills an array of structures (\code{isigopen},
  \code{osigopen}), the caller passes a list, which is emptied and receives
  one \code{Siginfo} per signal. Where C fills an array of samples
  (\code{getvec}, \code{getframe}), the caller passes a list with room for
  every value (\code{[0] * nsig}), whose first elements are replaced.
  \code{putvec} reads a list (or any sequence).
\item Where C fills a structure (\code{getann}, \code{getcal}), the caller
  passes an object (\code{Annotation()}, \code{Calinfo()}), whose fields are
  replaced.
\item \code{getseginfo()} takes no argument and returns the list of
  \code{Seginfo} objects (empty if the record is not a multi-segment
  record), instead of a count and a pointer.
\item The counts that only give the length of an array have defaults:
  \code{annopen(record, aiarray)}, \code{osigfopen(siarray)},
  \code{setheader(record, siarray)} and \code{setmsheader(record, names)}
  use the length of the list; \code{isigopen} and \code{osigopen} take
  \code{siarray=None, nsig=0}, \code{getinfo} and \code{calopen} take
  \code{None}.
\item Strings are \code{str}; an annotation's \code{aux} is \code{None} or
  \code{bytes}, and \code{Annotation.aux\_string()} returns its text.
  \code{wfdb\_error(fmt, *args)} formats as C's \code{printf}.
\item \code{getframes(n)} reads up to \code{n} sample vectors into a NumPy
  array (\cref{sec:library-additions}).
\item The functions of \file{<wfdb/ecgmap.h>} (\code{isann}, \code{isqrs},
  \code{map1}, ...) and the annotation codes (\code{NORMAL}, \code{PVC}, ...)
  are in the package; \code{isann} returns \code{True} or \code{False}
  (C: 1 or 0).
\end{itemize}

\begin{lstlisting}[language=Python,style=example]
import wfdbx

si = []
nsig = wfdbx.isigopen("100s", si, 2)    # si receives 2 Siginfo objects
v = [0] * 2
if nsig == 2 and wfdbx.getvec(v) > 0:   # v receives one sample per signal
    print(si[0].desc, v[0], v[1])       # MLII 995 1011
wfdbx.wfdbquit()
\end{lstlisting}

\subsection{C++}
\label{sec:library-cpp}

\file{<wfdbx/wfdb.hpp>} declares the functions, types, structures and
constants in namespace \code{wfdbx}; \file{<wfdbx/ecgcodes.hpp>} defines the
annotation codes, and \file{<wfdbx/ecgmap.hpp>} the macros of
\file{<wfdb/ecgmap.h>} as inline functions (\code{isann}, \code{isqrs}, ...,
which call \code{wfdb\_isann}, \code{wfdb\_isqrs}, ...).

\begin{itemize}
\item Arrays and structures are allocated by the caller and passed by
  pointer, as in C. String arguments are \code{const char *}, and
  \code{nullptr} is accepted where C accepts \code{NULL}.
\item Strings are returned by value, as \code{std::string}, or as
  \code{std::optional<std::string>} where the C function may return
  \code{NULL} (\code{annstr}, \code{ecgstr}, \code{anndesc}, \code{getinfo},
  \code{wfdbfile}); unlike C's pointers to static buffers, they stay valid
  after later calls.
\item The string fields of \code{WFDB\_Siginfo}, \code{WFDB\_Calinfo} and
  \code{WFDB\_Annotation} (including \code{aux}) are
  \code{std::optional<std::string>}, since C distinguishes \code{NULL} from
  \code{""}; \code{WFDB\_Anninfo::name} is a \code{std::string}. The
  structures are also available under the shorter names \code{Siginfo},
  \code{Anninfo}, \code{Annotation}, \code{Calinfo} and \code{Seginfo}.
\item \code{WFDB\_Time} is \code{long long} (64 bits); \code{getseginfo}
  takes a \code{const WFDB\_Seginfo **}.
\item \code{getvecs} reads many sample vectors at once
  (\cref{sec:library-additions}). The library is not thread-safe, as in C.
\end{itemize}

\begin{lstlisting}[language=C++,style=example]
#include <cstdio>
#include <wfdbx/wfdb.hpp>
using namespace wfdbx;

int main() {
    WFDB_Siginfo si[2];
    WFDB_Sample v[2];
    int nsig = isigopen("100s", si, 2);
    if (nsig == 2 && getvec(v) > 0)
        std::printf("%s %d %d\n", si[0].desc->c_str(), v[0], v[1]);
    wfdbquit();
    return 0;
}
\end{lstlisting}

\subsection{Rust}
\label{sec:library-rust}

The state of the library is held by a \code{Wfdb} value, and each function is
a method of it (\code{w.isigopen(...)}). Every function is also a free
function of the crate with the same name and arguments
(\code{wfdbx::\allowbreak isigopen(...)}), which uses a default \code{Wfdb} value, one per
thread; separate \code{Wfdb} values share no state, so that several records
can be open at once. The functions that use no state of their own are free
functions only: \code{wfdbldflags}, \code{wfdbcflags}, \code{wfdbdefwfdb},
\code{wfdbdefwfdbcal}, the count of write errors (\code{wfdbx\_write\_errors},
\code{wfdbx\_note\_write\_error}, which belongs to the process),
\code{isann}, \code{ammap} and \code{mamap} (and their \code{wfdb\_} forms),
and the macros of \file{<wfdb/ecgmap.h>} that use the mapping tables
(\code{isqrs}, \code{map1}, ...), which use those of the default value (the
methods \code{w.wfdb\_isqrs(...)}, \code{w.wfdb\_map1(...)}, ... use those of
\code{w}).

\begin{itemize}
\item Where C fills an array, the caller passes a mutable slice
  (\code{\&mut si}, \code{\&mut v}); where C fills a structure, a mutable
  reference (\code{\&mut annot}). C's \code{NULL} array, as in
  \code{isigopen(record, NULL, 0)}, is an empty slice.
\item String arguments accept \code{\&str}, \code{\&[u8]} or \code{b"..."};
  where C accepts \code{NULL}, they also accept \code{Option}s, and
  \code{wfdbx::NULL} is a typed \code{None}. Strings are returned as
  \code{Vec<u8>}, or \code{Option<Vec<u8>>} where C may return \code{NULL}
  (\code{anndesc}, \code{getinfo}, \code{wfdbfile}, ...).
\item C's \code{int} is \code{i32}, \code{long} and \code{WFDB\_Time} are
  \code{i64}, \code{WFDB\_Signal}, \code{WFDB\_Group} and
  \code{WFDB\_Annotator} are \code{u32}. The structures (\code{Siginfo},
  \code{Anninfo}, \code{Annotation}, \code{Calinfo}, \code{Seginfo}) have the
  C field names; \code{Default::default()} gives C's zero values; the
  \code{char} fields of \code{Annotation} are \code{i8} (\code{chan}:
  \code{u8}).
\item \code{getseginfo(\&mut segs)} takes an \code{\&mut Option<Vec<Seginfo>>},
  which receives a copy of the segment table (\code{None} if there is none).
\item \code{wfdb\_error(msg)} takes a message already formatted (with the
  crate's \code{cformat!} macro, which formats as C's \code{printf}), and
  \code{wfdb\_error\_fmt(fmt, args)} formats it.
\item \code{getvecs} and \code{getvecs\_status} read many sample vectors at
  once (\cref{sec:library-additions}); \code{wfdbx\_fclose\_output} closes a
  file that a program wrote itself and counts a write error
  (\cref{sec:library-additions}).
\end{itemize}

\begin{lstlisting}[language=Rust,style=example]
use wfdbx::Siginfo;

fn main() {
    let mut si = vec![Siginfo::default(); 2];
    let nsig = wfdbx::isigopen("100s", &mut si, 2);
    let mut v = vec![0i32; 2];
    if nsig == 2 && wfdbx::getvec(&mut v) > 0 {
        let desc = si[0].desc.clone().unwrap_or_default();
        println!("{} {} {}", String::from_utf8_lossy(&desc), v[0], v[1]);
    }
    wfdbx::wfdbquit();
}
\end{lstlisting}

\subsection{Julia}
\label{sec:library-julia}

The functions are in the module \code{WFDBX}. \code{using WFDBX} makes the
names of most of them available; those that it does not export are called
with the module prefix: \code{WFDBX.wfdb\_error}, \code{WFDBX.wfdb\_me\_fatal},
and the function forms of the macros of \file{<wfdb/ecgmap.h>},
\code{WFDBX.wfdb\_isann}, \code{WFDBX.wfdb\_isqrs}, ... (\code{isann},
\code{isqrs}, ... are exported). The state of the library is held by a
\code{Wfdb} object: every function that uses it takes one as an optional
first argument (\code{isigopen(w, "100s", si, 2)}) and otherwise uses the
default object, \code{lib()}; \code{Wfdb()} creates an independent one. The
functions that use no such state take none: \code{wfdbldflags},
\code{wfdbcflags}, \code{wfdbdefwfdb}, \code{wfdbdefwfdbcal}, \code{isann},
\code{ammap}, \code{mamap} (and their \code{wfdb\_} forms), and
\code{wfdbx\_write\_errors} and \code{wfdbx\_note\_write\_error}, whose
count belongs to the process.

\begin{itemize}
\item Where C fills an array, the caller passes a vector: of \code{Siginfo}
  structures (\code{[Siginfo() for \_ in 1:n]}), or of samples, which must
  be a \code{Vector\{Int32\}}. Where C fills a structure, the caller passes
  a mutable structure (\code{Annotation()}). \code{getseginfo(r)} takes a
  \code{Ref}, which receives a vector of \code{Seginfo} (or \code{nothing}).
\item C's \code{NULL} is \code{nothing}. Arguments may be \code{String}s or
  bytes. Text is returned as bytes (\code{Vector\{UInt8\}}; \code{String(x)}
  converts it), except the results of \code{timstr}, \code{mstimstr} and
  \code{datstr}, which are \code{String}s; the string fields of the
  structures are bytes, or \code{nothing}.
\item The integer types are those of C: \code{WFDB\_Sample} and \code{int} are
  \code{Int32}, \code{WFDB\_Time} and \code{long} are \code{Int64}; the
  \code{char} fields of \code{Annotation} are \code{Int8} (\code{chan}:
  \code{UInt8}).
\item \code{wfdb\_error(fmt, args...)} formats as C's \code{printf}.
  \code{wfdbx\_legacy\_args(args)} takes the arguments without the program
  name (\cref{sec:library-additions}).
\item \code{getvecs} and \code{getvecs\_status} read many sample vectors at
  once (\cref{sec:library-additions}).
\end{itemize}

\begin{lstlisting}[language=Julia,style=example]
using WFDBX

si = [Siginfo() for _ in 1:2]
nsig = isigopen("100s", si, 2)          # si receives the specifications
v = zeros(Int32, 2)
if nsig == 2 && getvec(v) > 0           # v receives one sample per signal
    println(String(si[1].desc), " ", v[1], " ", v[2])
end
wfdbquit()
\end{lstlisting}

\subsection{R}
\label{sec:library-r}

The functions are those of the package \code{wfdbx} (\code{library(wfdbx)}).
R has no output arguments, so a function that fills an array or a
structure in C returns it, together with the C return value, in a list; the
argument that C fills is omitted.

\begin{itemize}
\item \code{isigopen(record, nsig)} and \code{osigopen(record, nsig)} return
  \code{list(nsig, si)}: the C result and the list of \code{siginfo} lists;
  \code{wfdbinit(record, aiarray, nsig)} (without \code{nann}) returns
  \code{list(stat, si)}; \code{getvec()} and \code{getframe()} return
  \code{list(stat, v)}; \code{getann(an)} returns \code{list(stat, annot)};
  \code{getcal(desc, units)} returns \code{list(stat, cal)};
  \code{setheader} returns \code{list(stat, si)} (with the units and
  descriptions as \code{setheader} truncates them);
  \code{getseginfo()} returns the list of \code{seginfo} lists;
  \code{wfdbx\_legacy\_args(argv)} returns \code{list(argv, legacy)}.
\item The C structures are lists with the C field names, made by
  \code{siginfo()}, \code{anninfo()}, \code{annotation()}, \code{calinfo()}
  and \code{seginfo()} (\code{siginfo()} gives \code{spf} 1). Numbers are
  doubles; strings are character strings, and C's \code{NULL} is
  \code{NULL}; an annotation's \code{aux} is \code{NULL} or a raw vector,
  and \code{aux\_string(aux)} returns its text.
\item The C function \code{sample} is \code{wfdb\_sample}, so that it does not
  mask R's \code{sample}. The functions of \file{<wfdb/ecgmap.h>} accept
  vectors of codes. \code{wfdb\_error(fmt, ...)} formats as C's
  \code{printf}.
\item \code{getframes(n)} reads up to \code{n} sample vectors into a matrix
  (\cref{sec:library-additions}); \code{putframes(m)} writes the rows of a
  matrix with \code{putvec}.
\end{itemize}

\begin{lstlisting}[language=R,style=example]
library(wfdbx)

r <- isigopen("100s", 2)           # r$nsig, and r$si: the specifications
g <- getvec()                      # g$stat, and g$v: one sample per signal
if (r$nsig == 2 && g$stat > 0)
  cat(r$si[[1]]$desc, g$v, "\n")
wfdbquit()
\end{lstlisting}

\subsection{MATLAB/Octave}
\label{sec:library-matlab}

The functions are those of the package \file{+wfdbx}, called with the prefix
\code{wfdbx.} (\code{wfdbx.isigopen(...)}). An array or structure that C
fills is returned as an additional output value, and omitted from the
arguments.

\begin{itemize}
\item \code{[nsig, si] = wfdbx.isigopen(record, nsig)} and
  \code{[nsig, si] = wfdbx.osigopen(record, nsig)} return the C result and a
  struct array of \code{wfdbx.siginfo} structures; \code{[stat, si] =
  wfdbx.wfdbinit(record, aiarray, nann, nsig)}; \code{[stat, v] =
  wfdbx.getvec()} and \code{[stat, v] = wfdbx.getframe()} (an initial vector
  may be passed: its elements are kept where C leaves them unchanged);
  \code{[stat, annot] = wfdbx.getann(an)}; \code{[stat, cal] =
  wfdbx.getcal(desc, units)}; \code{[stat, si] = wfdbx.setheader(record, si,
  nsig)}; \code{segs = wfdbx.getseginfo()} returns a struct array;
  \code{[r, argv] = wfdbx.wfdbx\_legacy\_args(argv)}.
\item The C structures are MATLAB structures with the C field names, made by
  \code{wfdbx.siginfo(n)}, \code{wfdbx.anninfo(name, stat)},
  \code{wfdbx.annotation()}, \code{wfdbx.calinfo()} and \code{wfdbx.seginfo(n)}
  (\code{siginfo} gives \code{spf} 1). Samples, times and counts are
  doubles. Strings are \code{char} row vectors (record names can also be
  MATLAB strings), and C's \code{NULL} is \code{[]}; an annotation's
  \code{aux} is \code{[]} or a \code{uint8} row vector.
\item Constants are functions without arguments: \code{wfdbx.WFDB\_READ()},
  \code{wfdbx.NORMAL()}, ...; \code{wfdbx.ecgcodes()} returns all the
  annotation codes in one structure. The calibration constants
  (\code{WFDB\_AC\_COUPLED}, ...) and \code{WFDB\_FMT\_LIST} are not defined;
  their values are given in \cref{tab:library-constants}.
\item \code{wfdbx.wfdb\_error(fmt, ...)} formats as C's \code{printf}.
  \code{wfdbx.getframes(n)} reads up to \code{n} sample vectors into a matrix
  (\cref{sec:library-additions}), and \code{wfdbx.putframes(M)} writes the
  rows of a matrix.
\end{itemize}

\begin{lstlisting}[language=Matlab,style=example]
[nsig, si] = wfdbx.isigopen('100s', 2);   % si: the specifications
[stat, v] = wfdbx.getvec();               % v: one sample per signal
if nsig == 2 && stat > 0
    fprintf('%s %d %d\n', si(1).desc, v(1), v(2));
end
wfdbx.wfdbquit();
\end{lstlisting}

\subsection{Kotlin}
\label{sec:library-kotlin}

The functions are in the package \code{wfdbx} (\code{import wfdbx.*}). The
state of the library is held by a \code{Wfdb} object, and each function is an
extension function of it (\code{w.isigopen(...)}); every function is also a
top-level function with the same name and arguments
(\code{isigopen(...)}), which uses a default \code{Wfdb} object, one per
thread. Separate \code{Wfdb} objects share no state; one object must be used
by one thread at a time. As in Rust, \code{wfdbldflags}, \code{wfdbcflags},
\code{wfdbdefwfdb}, \code{wfdbdefwfdbcal}, \code{wfdbx\_write\_errors},
\code{wfdbx\_note\_write\_error}, \code{isann}, \code{ammap},
\code{mamap} (and their \code{wfdb\_} forms) and the macros of
\file{<wfdb/ecgmap.h>} that use the mapping tables (\code{isqrs},
\code{map1}, ..., on the default object; the extension functions are
\code{w.wfdb\_isqrs(...)}, ...) are top-level functions only.

\begin{itemize}
\item Where C fills an array, the caller passes an array
  (\code{Array<Siginfo>}, \code{IntArray}); where C fills a structure, an
  object (\code{Annotation()}). \code{getseginfo(list)} takes a
  \code{MutableList<Seginfo>?}, which receives the segments.
\item Strings are bytes (\code{ByteArray}), as in C; most functions also
  accept a \code{String} (one character per byte), and \code{latin1()} (in
  \code{wfdbx.cstd}) converts in either direction. C's \code{NULL} is
  \code{null}; results that C may return as \code{NULL} are
  \code{ByteArray?}.
\item C's \code{int} and \code{WFDB\_Sample} are \code{Int}, \code{long} and
  \code{WFDB\_Time} are \code{Long}, \code{double} is \code{Double}. The
  structures (\code{Siginfo}, \code{Anninfo}, \code{Annotation},
  \code{Calinfo}, \code{Seginfo}) have the C field names and C's zero values;
  the \code{char} fields of \code{Annotation} are \code{Int} properties that
  keep C's range (\code{anntyp}, \code{subtyp}, \code{num}: -128 to 127;
  \code{chan}: 0 to 255).
\item \code{wfdb\_error(msg)} takes a message already formatted (with
  \code{cformat} of \code{wfdbx.cstd}, which formats as C's \code{printf}).
\item \code{getvecs} and \code{getvecs\_status} read many sample vectors at
  once, and \code{wfdbx\_fclose\_output} closes a file that a program wrote
  itself (\cref{sec:library-additions}).
\end{itemize}

\begin{lstlisting}[language=Kotlin,style=example]
import wfdbx.*
import wfdbx.cstd.latin1

fun main() {
    val si = Array(2) { Siginfo() }
    val nsig = isigopen("100s", si, 2)
    val v = IntArray(2)
    if (nsig == 2 && getvec(v) > 0)
        println("${si[0].desc!!.latin1()} ${v[0]} ${v[1]}")
    wfdbquit()
}
\end{lstlisting}

\begin{note}[Functions that exist only in some languages]
Every function documented in this chapter exists in the seven languages
(\file{docs/INVENTORY.md}). The reader of many sample vectors has two names:
\code{getframes} in Python, R and MATLAB/Octave, and \code{getvecs} (with
\code{getvecs\_status}) in C++, Rust, Julia and Kotlin; it is documented under
\code{getframes}. A few helpers have no counterpart in the Python
implementation, and therefore no entry here: \code{putframes} (R,
MATLAB/Octave), \code{wfdbx\_fclose\_output} (Rust, Kotlin),
\code{wfdb\_error\_fmt} (Rust), \code{aux\_string} (R; a method of
\code{Annotation} in Python) and the structure constructors of R and
MATLAB/Octave. In R, \code{sample} is named \code{wfdb\_sample}.
\end{note}

% ==========================================================================
\section{Opening and closing records}
\label{sec:library-open}

A record is opened for reading by \code{isigopen} (its signals) and
\code{annopen} (its annotation files), or by \code{wfdbinit}, which calls
both; output signals are opened by \code{osigopen} or \code{osigfopen}, and
output annotation files by \code{annopen}. Programs that write files must call
\code{wfdbquit} (or close the files with the functions given in each entry)
before they exit.

A record name is found as \cref{chap:behavior} describes: a relative name
under the input directory (for reading) or the output directory (for
writing), an absolute name, a \code{\textasciitilde} name or a URL as given; a
name that begins with \code{+} keeps the files of the same kind already open,
and the record name \code{-} is the standard input or output. A trailing
\code{.hea} is removed from a record name.

\apientry{annopen}{Open annotation files}
\begin{lstlisting}[language=C,style=signature]
int annopen(char *record, const WFDB_Anninfo *aiarray, unsigned int nann);
\end{lstlisting}
Opens, for \code{record}, the \code{nann} annotation files described by
\code{aiarray}: for reading if their \code{stat} is \code{WFDB\_READ} or
\code{WFDB\_AHA\_READ}, for writing (created, or truncated) if it is
\code{WFDB\_WRITE} or \code{WFDB\_AHA\_WRITE}. The file of annotator
\code{atr} of record \code{100s} is \file{100s.atr}. Input and output
annotators are numbered separately, from 0, in the order of \code{aiarray}:
these are the numbers given to \code{getann} and \code{putann}. All the
annotation files already open are closed first (output files are completed,
and sorted if needed, as by \code{oannclose}), unless \code{record} begins with
\code{+}, which is then removed. All the \code{stat} values are checked before
any file is opened. An input file whose contents show the other format (MIT or
AHA) than its \code{stat} says is read in the format it has, with a warning.
The modification labels (mnemonics and descriptions defined at the beginning
of an input file) are applied to the tables of \code{annstr} and
\code{anndesc}; the time resolution recorded in an input file is used by
\code{getann} (see \code{getiaorigfreq}). An output annotator name may contain
only letters, digits, \code{\_}, \code{\textasciitilde}, \code{-} and
\code{/}.

With \code{nann} 0, \code{annopen} only closes the open annotation files
(\code{annopen("", NULL, 0)}) and reports whether that succeeded.
\begin{params}
\item[record] The record name.
\item[aiarray] An array of \code{nann} \code{WFDB\_Anninfo} structures
  (\cref{sec:library-structures}): the annotator name (\code{name}) and the
  mode (\code{stat}) of each file.
\item[nann] The number of annotators in \code{aiarray}.
\end{params}
\returns{0 on success; -3 if an input file cannot be opened; -4 if an output
file cannot be created or its annotator name is illegal; -5 if a \code{stat}
is not one of the four values. With \code{nann} 0: 0, -6 if an output file
closed by the call was out of order (whether or not it could be sorted), or
-7 if one could not be written completely.}
\raises{\code{annopen: can't read annotator ATR for record REC};
\code{annopen: can't write annotator ATR for record REC};
\code{annopen: illegal stat N for annotator ATR, record REC}; \code{init:
illegal character N in annotator name}.}

WFDB-x: input files are read only where their name says, output files are
created under the output directory, and out-of-order output files closed by
the call are sorted in the same process (\cref{chap:behavior}).

\apientry{isigopen}{Open the input signals of a record}
\begin{lstlisting}[language=C,style=signature]
int isigopen(char *record, WFDB_Siginfo *siarray, int nsig);
\end{lstlisting}
Reads the header of \code{record} and opens its signal files for reading.
The input signals already open are closed first, unless \code{record} begins
with \code{+}, which is then removed (the new signals are numbered after
those open). The header also sets the sampling and counter frequencies, the
base counter value, the base time and date, and the length of the record,
which the time conversion functions use.

With \code{nsig} greater than 0, at most \code{nsig} signals are opened, and
the structures \code{siarray[0]} to \code{siarray[n-1]} are filled with their
specifications (signal numbers are assigned in the order of the header).
Signals stored in the same file (a signal group) are opened together: a
group is skipped if its file cannot be opened, or if its signals do not fit
in the remaining entries (a record whose signals are all in one file is
therefore not opened at all if \code{nsig} is smaller than its number of
signals).

With \code{nsig} 0 or negative, no signals are opened (and the open input
signals are closed, as above): the header is read, up to \code{-nsig}
structures of \code{siarray} are filled from it (\code{siarray} may be
\code{NULL} when \code{nsig} is 0), and the number of signals named in the
header is returned. \code{isigopen(record, NULL, 0)} thus gives the number of
signals, to allocate \code{siarray} before opening them, and sets the time
base of the record for \code{strtim} and \code{timstr}.

A record name that contains a \code{.}, other than in a final \code{.hea},
names an EDF file, read as a record (\cref{chap:formats}). A multi-segment
record opens its segments as they are read (see \code{getseginfo}).
\begin{params}
\item[record] The record name.
\item[siarray] An array of at least \code{nsig} (or \code{-nsig})
  \code{WFDB\_Siginfo} structures, filled by the function.
\item[nsig] The maximum number of signals to open; 0 or a negative number to
  read the header only.
\end{params}
\returns{With \code{nsig} > 0: the number of signals opened (possibly fewer
than \code{nsig}), 0 if the record has no signals or none could be opened.
With \code{nsig} <= 0: the number of signals of the record. In either case:
-1 if the header cannot be read (probably an incorrect record name), -2 if it
is not a valid header, -3 if the internal limits are exceeded (the total
number of samples per frame, or the skew, too large), in which case all the
input signals are closed.}
\raises{\code{init: can't open header for record REC}; \code{isigopen:
record REC has no signals}; \code{isigopen: none of the signals for record
REC is readable}; the messages of the header parser (\code{init: ...}).}

WFDB-x (\cref{chap:behavior}): the header is read only where the record name
says, and the signal files and segments that it names only from the header's
directory (\code{init: file F named in the header of record REC is not in
the header's directory} for a name that leaves it); a header without a
sampling frequency, ADC resolution or units leaves them unspecified (0, 0,
\code{NULL}) instead of assuming 250 Hz, a resolution and millivolts; a header
with a negative or impossibly large number of signals gives \code{init:
unexpected EOF in header file for record REC} (C tries to allocate memory for
that count).

\apientry{osigopen}{Open output signals as described by a header}
\begin{lstlisting}[language=C,style=signature]
int osigopen(char *record, WFDB_Siginfo *siarray, unsigned int nsig);
\end{lstlisting}
Reads the header of the existing record \code{record} and opens for writing
(creating or truncating them) the signal files of its first \code{nsig}
signals, with the specifications of the header; \code{siarray} receives them
(with \code{nsamp} and \code{cksum} 0). The output signals already open are
closed first, unless \code{record} begins with \code{+}, which is then
removed. The samples are then written with \code{putvec}, and a header for
the new signals with \code{newheader}. A header of an existing database
record names that record's signal files, which would be overwritten: the
header used here is usually written for the purpose (a ``piped'' record such
as \code{8l}, whose signal files are \file{data0} and \file{data1}).
\begin{params}
\item[record] The record whose header specifies the output signals.
\item[siarray] An array of at least \code{nsig} \code{WFDB\_Siginfo}
  structures, filled by the function.
\item[nsig] The number of output signals to open.
\end{params}
\returns{The number of output signals opened (\code{nsig}, 0 when \code{nsig}
is 0); -1 if the header cannot be read; -2 if it is not valid, or if the
signals of a group differ in format or block size; -3 if the record has fewer
than \code{nsig} signals or a signal file cannot be created.}
\raises{\code{osigopen: record REC has fewer signals than needed};
\code{osigopen: can't open FILE}; \code{osigopen: error in specification of
signal N or M}.}

The Programmer's Guide allows \code{nsig} 0 with \code{siarray} \code{NULL},
to close the output signals; the C function then reads before the start of
its internal array. In WFDB-x, in every language, \code{osigopen(record, NULL,
0)} closes the output signals already open (unless \code{record} begins with
\code{+}), reads the header (and returns -1 if it cannot be read), opens
nothing, and returns 0 (\cref{sec:crashes}). \code{osigfopen(NULL, 0)} also
closes the output signals.

WFDB-x: a signal file is created under the output directory, and a signal
file that is one of the signal files of the open input record is refused
(\code{osigopen: output file FILE is also an input file}, -3; legacy mode does
not check); see \cref{chap:behavior}.

\apientry{osigfopen}{Open output signals as described by the caller}
\begin{lstlisting}[language=C,style=signature]
int osigfopen(const WFDB_Siginfo *siarray, unsigned int nsig);
\end{lstlisting}
Closes the open output signals, and opens \code{nsig} output signals as
\code{siarray} describes them: file name, description, units, gain, format,
samples per frame, block size, ADC resolution and zero, baseline
(\code{initval}, \code{nsamp} and \code{cksum} are not used). Signals written
to the same file (a multiplexed signal file) must be consecutive and have the
same \code{fname} and \code{group}, and the same format and block size; each
new file has the next group number, starting at 0. A \code{fname} of
\code{-} is the standard output; a signal of format 0 with no file name (or
\code{\textasciitilde}) is not written, and a signal of format 0 with a file
name is written in format 16. An \code{spf} smaller than 1 is taken as 1.

\code{osigfopen(NULL, 0)} closes the output signals (and the header opened by
\code{newheader}, \code{setheader} or \code{setmsheader}) and reports whether
they were written successfully.
\begin{params}
\item[siarray] An array of \code{nsig} \code{WFDB\_Siginfo} structures,
  filled by the caller.
\item[nsig] The number of output signals.
\end{params}
\returns{The number of output signals opened (\code{nsig}); -2 for an error in
a specification (\code{fname} and \code{desc} together longer than 200
characters, a negative \code{bsize}, a format not in
\code{WFDB\_FMT\_LIST}, a wrong group or file name, or signals of a group with
different formats or block sizes); -3 if a signal file cannot be created.
With \code{nsig} 0: 0, or -4 if a signal file or header closed by the call
could not be written completely.}
\raises{\code{osigfopen: error in specification of signal N};
\code{osigfopen: incorrect file name or group for signal N};
\code{osigfopen: can't open FILE}; \code{osigclose: write error in signal
group N}.}

WFDB-x: as for \code{osigopen}, files are created under the output directory,
and a signal file of the open input record is refused (\code{osigfopen:
output file FILE is also an input file}); see \cref{chap:behavior}.

\apientry{wfdbinit}{Open the annotation files and the input signals}
\begin{lstlisting}[language=C,style=signature]
int wfdbinit(char *record, const WFDB_Anninfo *aiarray, unsigned int nann,
             WFDB_Siginfo *siarray, unsigned int nsig);
\end{lstlisting}
Calls \code{annopen(record, aiarray, nann)} and, if it returns 0,
\code{isigopen(record, siarray, nsig)}. The Programmer's Guide advises
against using \code{wfdbinit} together with \code{setifreq}: \code{setifreq}
must be called after the signals are opened and before the annotation files
are, to apply to the annotation times.
\begin{params}
\item[record] The record name.
\item[aiarray] The annotators, as for \code{annopen}.
\item[nann] The number of annotators.
\item[siarray] An array of at least \code{nsig} \code{WFDB\_Siginfo}
  structures, as for \code{isigopen}.
\item[nsig] The maximum number of input signals to open.
\end{params}
\returns{The negative result of \code{annopen} if it fails; otherwise the
result of \code{isigopen}: the number of signals opened, 0 if none (no
message if \code{nsig} is 0), or -1, -2, -3.}

In R, \code{wfdbinit(record, aiarray, nsig)} has no \code{nann} argument, and
returns \code{list(stat, si)}; in MATLAB/Octave, \code{[stat, si] =
wfdbx.wfdbinit(record, aiarray, nann, nsig)}.

\apientry{wfdbquit}{Close all files and reset the library}
\begin{lstlisting}[language=C,style=signature]
void wfdbquit(void);
\end{lstlisting}
Closes all the annotation files (completing and, if needed, sorting the
output files, as \code{oannclose}), the info file opened by \code{setinfo},
the input and output signal files (writing out the buffered samples) and the
header being written, and releases the buffers of \code{sample} and the info
strings of \code{getinfo}. It resets the sampling frequencies, the base time
and date, the record length and the segments of the open record. It does not
write a header (\code{newheader} or \code{setheader} must be called first),
and, as in C, it does not reset the counter frequency and base counter value,
the output annotation time resolution (\code{setafreq}), the buffer sizes
(\code{setibsize}, \code{setobsize}), the calibration list, or the tables of
mnemonics, descriptions and code mappings (\cref{chap:behavior}). Programs
that write signals or annotations must call it (or the closing functions of
each kind of file) before they exit.

WFDB-x: the input directory, the output directory and the network setting
are kept; in legacy mode, the WFDB path is restored as by \code{resetwfdb}
(\cref{chap:behavior}).

\apientry{wfdbflush}{Write the buffered output}
\begin{lstlisting}[language=C,style=signature]
void wfdbflush(void);
\end{lstlisting}
Writes the buffered output annotations and samples to their files, so that
the files are up to date while they stay open.

% ==========================================================================
\section{Reading and writing signals}
\label{sec:library-signals}

The input signals are read one sample vector at a time, sequentially, with
\code{getvec} (one sample of each signal) or \code{getframe} (every sample of
a frame); \code{isigsettime}, \code{isgsettime} and \code{tnextvec} move to
another time, and \code{sample} gives random access. Samples are integers in
ADC units; a missing or invalid sample is \code{WFDB\_INVALID\_SAMPLE}
(-32768) unless the \code{WFDB\_GVPAD} mode is set. Times are in sample
intervals: of the frame, of the highest sampling frequency in
\code{WFDB\_HIGHRES} mode, or of the frequency set by \code{setifreq}.

\apientry{getvec}{Read a sample of each input signal}
\begin{lstlisting}[language=C,style=signature]
int getvec(WFDB_Sample *vector);
\end{lstlisting}
Reads the next sample of each input signal into \code{vector}: \code{vector[i]}
receives the sample of signal \code{i}. In \code{WFDB\_LOWRES} mode (the
default), a signal with several samples per frame gives the mean of the
samples of the frame (rounded to the nearest integer, halves upward; invalid
if one of them is); in \code{WFDB\_HIGHRES} mode, each call gives the next
sample at the highest sampling frequency, and the other signals are repeated
(see \code{setgvmode}). After \code{setifreq}, the signals are resampled at
the chosen frequency by linear interpolation. Skewed signals are corrected
for their skew. The checksum of a signal is verified when the number of
samples given in the header has been read, unless the signals were
positioned (\code{isigsettime}, \code{isgsettime}) before.
\begin{params}
\item[vector] An array with room for one sample per input signal, filled by
  the function.
\end{params}
\returns{The number of input signals, on success; 0 if no input signals are
open; -1 at the end of the data (the contents of \code{vector} are then not
valid); -3 at an unexpected end of a signal file (or a FLAC decoding error);
-4 on a checksum error, detected on the last frame (which is valid).}
\raises{\code{getvec: unexpected EOF in signal N}; \code{getvec: checksum
error in signal N}; \code{getvec: error opening segment SEG}.}

\apientry{getframe}{Read a frame of the input signals}
\begin{lstlisting}[language=C,style=signature]
int getframe(WFDB_Sample *vector);
\end{lstlisting}
Reads the next frame: every sample of every input signal for one frame
interval, the \code{spf} samples of signal 0 first, then those of signal 1,
and so on. For a record whose signals all have one sample per frame,
\code{getframe} and \code{getvec} (without \code{setifreq}) read the same
vectors. \code{setgvmode} and \code{setifreq} do not affect it.
\begin{params}
\item[vector] An array with room for the sum of the \code{spf} fields of the
  input signals, filled by the function.
\end{params}
\returns{As \code{getvec}.}

\apientry{putvec}{Write a sample to each output signal}
\begin{lstlisting}[language=C,style=signature]
int putvec(const WFDB_Sample *vector);
\end{lstlisting}
Writes a frame to the output signals: \code{vector} holds the \code{spf}
samples of output signal 0, then those of signal 1, and so on (one sample per
signal when every \code{spf} is 1). The number of samples and the checksum of
each signal are kept, for \code{newheader}. A sample equal to
\code{WFDB\_INVALID\_SAMPLE} is written as the invalid-sample value of the
format; a sample outside the range of the format is written as C writes it
(its low-order bits). In format 8 (first differences), a difference outside
-128 to 127 is limited to that range, and the following samples correct the
level as soon as the differences allow.
\begin{params}
\item[vector] The samples, given by the caller.
\end{params}
\returns{The number of output signals on success; 0 if a difference was too
large in format 8; -1 on a write error.}
\raises{\code{putvec: write error in signal N}.}

\apientry{isigsettime}{Skip to a time in all input signals}
\begin{lstlisting}[language=C,style=signature]
int isigsettime(WFDB_Time t);
\end{lstlisting}
Positions the input signals so that the next sample read is the one at time
\code{t}; only the magnitude of \code{t} counts, so that a value returned by
\code{strtim} for a time of day can be given as it is. Nothing is done if the
signals are already at \code{t} (without \code{setifreq}). A pipe, such as the
standard input, cannot be positioned backward.
\begin{params}
\item[t] The time, in sample intervals (see the introduction of this
  section).
\end{params}
\returns{0 on success (also if no input signals are open); -1 at the end of
the data or on an improper seek; -2 if a signal group is incorrect.}
\raises{\code{isigsettime: improper seek on signal group N}.}

\apientry{isgsettime}{Skip to a time in one signal group}
\begin{lstlisting}[language=C,style=signature]
int isgsettime(WFDB_Group sgroup, WFDB_Time t);
\end{lstlisting}
As \code{isigsettime}, for the signals of group \code{sgroup} only (the
signals stored in one file); the other groups keep their position. It serves
programs that open several records at once (\code{isigopen("+REC", ...)}).
\begin{params}
\item[sgroup] The signal group number.
\item[t] The time, in sample intervals.
\end{params}
\returns{0 on success; -1 at the end of the data or on an improper seek; -2
if \code{sgroup} is not a signal group.}
\raises{\code{isigsettime: improper seek on signal group N};
\code{isgsettime: incorrect signal group number N}.}

WFDB-x: a group number beyond the groups of the open record is reported as
incorrect (-2) also when the signals are already at \code{t}, where C reads
past the end of its array (\cref{chap:behavior}).

\apientry{tnextvec}{Find the next valid sample of a signal}
\begin{lstlisting}[language=C,style=signature]
WFDB_Time tnextvec(WFDB_Signal s, WFDB_Time t);
\end{lstlisting}
Finds the first time, at or after \code{t}, at which input signal \code{s}
has a valid sample (not \code{WFDB\_INVALID\_SAMPLE}), and positions the
input signals there, so that the next \code{getvec} reads it. It skips long
gaps (in a multi-segment record, segments without the signal) efficiently.
\begin{params}
\item[s] The input signal number.
\item[t] The time at which the search begins, in sample intervals.
\end{params}
\returns{The time of the next valid sample (\code{>= 0}); -1 at the end of
the data, on an improper seek, or if \code{s} is not an input signal (or the
negative status of \code{getvec}).}
\raises{\code{nextvect: illegal signal number N} (the C library's spelling).}

\apientry{sample}{Read a sample by its time}
\begin{lstlisting}[language=C,style=signature]
WFDB_Sample sample(WFDB_Signal s, WFDB_Time t);
\end{lstlisting}
Returns sample \code{t} of input signal \code{s}, in any order: the samples
are read with \code{getvec} (which moves the input signals) into a buffer of
the 4096 most recent sample vectors, from which they are returned. A negative
\code{t} is taken as 0. \code{sample\_valid} tells whether the value is
valid. \code{wfdbquit} releases the buffer. In R, this function is
\code{wfdb\_sample}.
\begin{params}
\item[s] The input signal number (0 to the number of input signals minus 1).
\item[t] The sample number.
\end{params}
\returns{The sample, in ADC units; after the end of the record, the last
sample read; \code{WFDB\_INVALID\_SAMPLE} if the signal has no valid sample at
\code{t}, if \code{s} is not an input signal, or if the signals cannot be
positioned.}

\apientry{sample\_valid}{Whether the last sample read by sample is valid}
\begin{lstlisting}[language=C,style=signature]
int sample_valid(void);
\end{lstlisting}
Tells whether the value most recently returned by \code{sample} is valid, so
that \code{while (sample\_valid())} reads to the end of the record and
\code{sample\_valid() > 0} selects the values usable in a computation.
\returns{1 if the value was valid; 0 if \code{t} was after the end of the
record (the value is padding); -1 if the value was
\code{WFDB\_INVALID\_SAMPLE} (the signal is not available at \code{t}, or the
signal number is not valid).}

\apientry{findsig}{Find an input signal by name or number}
\begin{lstlisting}[language=C,style=signature]
int findsig(const char *signame);
\end{lstlisting}
Returns the number of the input signal designated by \code{signame}: a string
of decimal digits that is the number of an open input signal is taken as that
number; otherwise the lowest-numbered input signal whose description
(\code{desc}) is exactly \code{signame} is returned. The descriptions are
compared as the header gives them, spaces included (the high-level readers
remove surrounding spaces from signal names; the library does not).
\begin{params}
\item[signame] A signal number or a signal description, such as \code{"MLII"}.
\end{params}
\returns{The signal number, or -1 if there is no such signal.}

\apientry{getspf}{Number of samples per signal per frame}
\begin{lstlisting}[language=C,style=signature]
int getspf(void);
\end{lstlisting}
Returns the number of sample vectors that \code{getvec} returns per frame:
1, except in \code{WFDB\_HIGHRES} mode for a record with oversampled signals,
where it is the largest number of samples per frame of the input signals (the
sampling frequency divided by the frame frequency, rounded).
\returns{The number of samples per signal per frame.}

\apientry{setgvmode}{Set the mode of getvec}
\begin{lstlisting}[language=C,style=signature]
void setgvmode(int mode);
\end{lstlisting}
Sets how \code{getvec} reads records whose signals have different numbers of
samples per frame, and how invalid samples are returned.
\code{WFDB\_LOWRES} (0, the default): one sample per signal per frame, the
oversampled signals being averaged over the frame. \code{WFDB\_HIGHRES}
(1): every sample of the most oversampled signal, the other signals being
repeated (zero-order hold); the sampling frequency returned by
\code{sampfreq(NULL)}, and the times of the time conversion functions and of
the annotations, are then in units of the high-resolution sample interval
(for annotation files opened after the call). \code{WFDB\_GVPAD} (2), added to
either: an invalid or missing sample is replaced by the previous valid sample
of its signal, in \code{getvec}, \code{getframe} and \code{sample}. A negative
\code{mode} restores the default, read from the environment variable
\code{WFDBGVMODE} (or \code{WFDB\_LOWRES} if it is not set).
\begin{params}
\item[mode] \code{WFDB\_LOWRES} or \code{WFDB\_HIGHRES}, optionally plus
  \code{WFDB\_GVPAD}; or a negative value.
\end{params}

\apientry{getgvmode}{The mode of getvec}
\begin{lstlisting}[language=C,style=signature]
int getgvmode(void);
\end{lstlisting}
Returns the mode set by \code{setgvmode}: the \code{WFDB\_HIGHRES} bit is
set in high-resolution mode, the \code{WFDB\_GVPAD} bit if invalid samples
are replaced.
\returns{The mode (0 to 3).}

\apientry{setifreq}{Resample the input signals}
\begin{lstlisting}[language=C,style=signature]
int setifreq(WFDB_Frequency freq);
\end{lstlisting}
Makes \code{getvec} return the input signals resampled at \code{freq} samples
per second, by linear interpolation between the two nearest samples
(\code{getframe} is not affected). It must be called after the signals are
opened. From then on, the time conversion functions, \code{isigsettime} and
\code{isgsettime} use sample intervals at \code{freq}, the annotators opened
afterwards return times at \code{freq}, and the time resolution of the output
annotation files opened afterwards is \code{freq} (it calls
\code{setafreq(freq)}). The resampling uses the greatest common divisor of the
two frequencies, treated as multiples of 0.005 Hz.
\begin{params}
\item[freq] The sampling frequency wanted, in Hz (> 0).
\end{params}
\returns{0 on success; -1 if no record is open, if \code{freq} is not
positive, or (WFDB-x) if the sampling frequency of the open record is not
specified.}
\raises{\code{setifreq: no open input record}; \code{setifreq: improper
frequency F (must be > 0)}; \code{setifreq: sampling frequency is not
specified}.}

WFDB-x: a record without a sampling frequency is an error
(\cref{chap:behavior}); in legacy mode, it is reported as \code{setifreq: no
open input record}, as in C.

\apientry{getifreq}{The frequency of the samples returned by getvec}
\begin{lstlisting}[language=C,style=signature]
WFDB_Frequency getifreq(void);
\end{lstlisting}
Returns the sampling frequency of the samples returned by \code{getvec}: the
one set by \code{setifreq}, or else that of the open record (as
\code{sampfreq(NULL)}).
\returns{The frequency, in Hz (0 if no record is open, or if its frequency is
not specified).}

\apientry{setibsize}{Set the size of the input buffers}
\begin{lstlisting}[language=C,style=signature]
int setibsize(int input_buffer_size);
\end{lstlisting}
Sets the size of the buffers in which the input signal files are read, for
the signals opened afterwards (a block size given in the header takes
precedence). It cannot be used while input signals are open.
\begin{params}
\item[input\_buffer\_size] The size in bytes; 0 for the default
  (\code{BUFSIZ}, 8192 bytes).
\end{params}
\returns{The new size; -1 if input signals are open; -2 if the size is
negative.}
\raises{\code{setibsize: can't change buffer size after isigopen};
\code{setibsize: illegal buffer size N}.}

\apientry{setobsize}{Set the size of the output buffers}
\begin{lstlisting}[language=C,style=signature]
int setobsize(int output_buffer_size);
\end{lstlisting}
As \code{setibsize}, for the output signal files opened afterwards by
\code{osigopen} and \code{osigfopen}; it cannot be used while output signals
are open.
\begin{params}
\item[output\_buffer\_size] The size in bytes; 0 for the default
  (\code{BUFSIZ}, 8192 bytes).
\end{params}
\returns{The new size; -1 if output signals are open; -2 if the size is
negative.}
\raises{\code{setobsize: can't change buffer size after osig[f]open};
\code{setobsize: illegal buffer size N}.}

% ==========================================================================
\section{Time and date conversion}
\label{sec:library-time}

Times are kept as \code{WFDB\_Time} values, in sample intervals. The
functions of this section convert them to and from text, using the sampling
frequency (or the frequency set by \code{setifreq}), the base time and date
and the counter frequency of the open record, set by \code{isigopen},
\code{wfdbinit} or \code{sampfreq} from a header, or by \code{setsampfreq},
\code{setbasetime} and \code{setcfreq}. A positive time is an elapsed time
(an interval from the beginning of the record); a time of day is represented
by the negated number of sample intervals from the beginning of the record
(0 or negative).

WFDB-x (\cref{chap:behavior}): when no sampling frequency is known, C took a
sample interval to be one second; WFDB-x reports \code{FUNCTION: sampling
frequency is not specified}, and \code{timstr} and \code{mstimstr} return the
time as a sample number (\code{s} followed by its magnitude), \code{strtim}
returns 0 for a string that needs the frequency (\code{wfdbx\_time\_needs\_freq}
tells which). Legacy mode keeps C's behavior.

\apientry{timstr}{Format a time}
\begin{lstlisting}[language=C,style=signature]
char *timstr(WFDB_Time t);
\end{lstlisting}
Converts \code{t} to text. If \code{t} is positive, or if neither a base time
nor a base date is defined, \code{|t|} is an elapsed time, written
\code{h:mm:ss} right-aligned in 8 characters (more from 100 hours on),
without the hours under one hour (\code{"    1:00"}, \code{" 1:02:05"}). Otherwise \code{-t} sample intervals
after the base time are written as a time of day in brackets, with all the
digits, followed by the date if a base date is defined, or else by the
number of days after the base time if it is not 0: \code{"[12:00:01
01/01/2020]"}, \code{"[12:00:00 1]"}. The time is rounded to the
millisecond, and the milliseconds are then dropped. A
base time of midnight (0) without a base date counts as undefined.
\begin{params}
\item[t] The time, in sample intervals.
\end{params}
\returns{The text. In C, it is in a static buffer (shared with
\code{mstimstr}) that the next call overwrites; in WFDB-x, it is a new string
in every language.}
\raises{WFDB-x: \code{timstr: sampling frequency is not specified}, and the
result \code{s} followed by \code{|t|} (\code{"s100"}), if the frequency is
not known.}

\apientry{mstimstr}{Format a time with milliseconds}
\begin{lstlisting}[language=C,style=signature]
char *mstimstr(WFDB_Time t);
\end{lstlisting}
As \code{timstr}, with the milliseconds (rounded): an elapsed time
\code{h:mm:ss.mmm} right-aligned in 12 characters (\code{"    0:01.000"}), a
time of day \code{"[hh:mm:ss.mmm dd/mm/yyyy]"}. \code{mstimstr(0)} gives the
base time and date of the record (there is no \code{getbasetime}).
\begin{params}
\item[t] The time, in sample intervals.
\end{params}
\returns{The text, as for \code{timstr}.}
\raises{WFDB-x: \code{mstimstr: sampling frequency is not specified}, as for
\code{timstr}.}

\apientry{strtim}{Parse a time}
\begin{lstlisting}[language=C,style=signature]
WFDB_Time strtim(const char *time_string);
\end{lstlisting}
Converts a time string to sample intervals. Leading spaces, tabs, carriage
returns and newlines are skipped; then, according to the first character:
\begin{description}
\item[\code{s}\emph{N}] sample number \emph{N} (\code{s12345});
\item[\code{c}\emph{X}] counter value \emph{X}, converted with the counter
  frequency and the base counter value (\code{c350.5}); without a counter
  frequency, \emph{X} is a sample number;
\item[\code{f}\emph{N}] frame number \emph{N}, converted to sample intervals;
\item[\code{e}] the end of the record (its length; 0 if the header does not
  give it), in high-resolution intervals in \code{WFDB\_HIGHRES} mode;
\item[\code{i}] the time of the next sample of the input signals;
\item[\code{o}] the time of the next sample of the output signals;
\item[\code{[}\emph{time} \emph{date}\code{]}] a time of day, \code{[hh:mm:ss]},
  followed optionally by a date \code{dd/mm/yyyy} or by a number of days after
  the base date (\code{[13:6:0]}, \code{[8:0:0 1]}, \code{[12:0:0 1/3/1992]}):
  the result is the negated number of sample intervals from the base time to
  that time (0 or negative); without the closing bracket, the result is 0;
\item[otherwise] an elapsed time \code{h:m:s}, \code{m:s} or \code{s}, each
  part with optional decimals (\code{2:14.875}, \code{143}, \code{4:02:01}),
  rounded to the nearest sample interval.
\end{description}
For records with sampling frequencies below 1 kHz,
\code{strtim(mstimstr(t))} is \code{t}.
\begin{params}
\item[time\_string] The time string.
\end{params}
\returns{A positive number of sample intervals for an elapsed time, a sample
number, a counter value, \code{e}, \code{i} or \code{o}; the negated number
of sample intervals for a time of day; 0 for a string that is not a time
(indistinguishable from a time of 0).}
\raises{WFDB-x: \code{strtim: sampling frequency is not specified}, and 0, if
the frequency is not known and the string needs it.}

\apientry{datstr}{Format a date}
\begin{lstlisting}[language=C,style=signature]
char *datstr(WFDB_Date d);
\end{lstlisting}
Converts a Julian date (a day number; these Julian dates begin at midnight
local time, not at noon) to text, \code{" dd/mm/yyyy"}, with a leading space,
which the Programmer's Guide does not mention (\code{datstr(1753639)} is
\code{" 14/07/1789"}). The Gregorian calendar applies from 15 October 1582;
years before 1 A.D. are negative (there is no year 0).
\begin{params}
\item[d] The Julian date.
\end{params}
\returns{The text (a static buffer in C; a new string in WFDB-x).}

\apientry{strdat}{Parse a date}
\begin{lstlisting}[language=C,style=signature]
WFDB_Date strdat(const char *date_string);
\end{lstlisting}
Converts a date \code{dd/mm/yyyy} to a Julian date. The year is taken as
written: \code{15/3/89} is in the year 89, not 1989. The number of days
between two dates is the difference of their Julian dates.
\begin{params}
\item[date\_string] The date.
\end{params}
\returns{The Julian date; 0 if the string does not contain two \code{/}, the
day is not 1 to 31, the month is not 1 to 12, or the year is 0.}

\apientry{sampfreq}{Sampling frequency of a record}
\begin{lstlisting}[language=C,style=signature]
WFDB_Frequency sampfreq(char *record);
\end{lstlisting}
With a record name, reads the header of \code{record} and returns its
sampling frequency; the frequency, base time and date and counter frequency
of the header then become those used by the time conversion functions (as
after \code{isigopen(record, NULL, 0)}, without closing the open signals).
With \code{NULL}, returns the sampling frequency in use. For a record with
oversampled signals, the result is the frame frequency in
\code{WFDB\_LOWRES} mode and the highest sampling frequency in
\code{WFDB\_HIGHRES} mode.
\begin{params}
\item[record] The record name, or \code{NULL}.
\end{params}
\returns{The sampling frequency in Hz; -1 if the header cannot be read; -2 if
it is not valid; WFDB-x: 0 if the header does not give a frequency (C: 250,
\code{WFDB\_DEFFREQ}; see \cref{chap:behavior}). \code{sampfreq(record)}
first sets the frequency to 0, so that a frequency set by \code{setsampfreq}
does not apply to it.}

\apientry{setsampfreq}{Set the sampling frequency}
\begin{lstlisting}[language=C,style=signature]
int setsampfreq(WFDB_Frequency sampling_frequency);
\end{lstlisting}
Sets the frame frequency used by the time conversion functions and written
by \code{newheader} and \code{setheader}: it is the sampling frequency of a
record whose signals have one sample per frame (in \code{WFDB\_HIGHRES} mode,
the sampling frequency becomes the frame frequency times the largest number
of samples per frame). A header read afterwards by \code{isigopen}, \code{osigopen} or
\code{wfdbinit} replaces it, unless the
header gives no frequency: \code{setsampfreq(250)} before opening a record
whose header gives none reproduces C's default (\cref{chap:behavior}).
\begin{params}
\item[sampling\_frequency] The frame frequency, in Hz (0 or more).
\end{params}
\returns{0 on success; -1 if the frequency is negative.}
\raises{\code{setsampfreq: sampling frequency must not be negative}.}

\apientry{setbasetime}{Set the base time and date}
\begin{lstlisting}[language=C,style=signature]
int setbasetime(char *time_string);
\end{lstlisting}
Sets the base time (the time of day of sample 0) and the base date, used by
the time conversion functions and written by \code{newheader} and
\code{setheader}. The string is a time \code{hh:mm:ss} (with optional
decimals), optionally followed by a space and a date \code{dd/mm/yyyy};
leading spaces are skipped. Only a space separates the date (not a tab, which
the Programmer's Guide also allows). An empty string or \code{NULL} sets the
current local time and date.
\begin{params}
\item[time\_string] The time and date, or \code{""} or \code{NULL}.
\end{params}
\returns{0 on success; -1 if neither a time nor a date can be read from the
string (and it does not begin with \code{[}).}
\raises{\code{setbasetime: incorrect time format, 'STRING'}.}

\apientry{getcfreq}{Counter frequency}
\begin{lstlisting}[language=C,style=signature]
WFDB_Frequency getcfreq(void);
\end{lstlisting}
Returns the counter frequency: the rate at which a counter (of a tape
recorder, or the pages of a chart) changes, as given by the header or by
\code{setcfreq}, which \code{strtim} uses for \code{c} times. If none is
defined (0 or negative), it returns the frame frequency.
\returns{The counter frequency, in Hz.}

\apientry{setcfreq}{Set the counter frequency}
\begin{lstlisting}[language=C,style=signature]
void setcfreq(WFDB_Frequency counter_frequency);
\end{lstlisting}
Sets the counter frequency, used by \code{strtim} and written by
\code{newheader} and \code{setheader}; 0 or a negative value means none. A
header read afterwards replaces it. \code{wfdbquit} does not reset it.
\begin{params}
\item[counter\_frequency] The counter frequency, in Hz.
\end{params}

\apientry{getbasecount}{Base counter value}
\begin{lstlisting}[language=C,style=signature]
double getbasecount(void);
\end{lstlisting}
Returns the counter value of sample 0, as given by the header or by
\code{setbasecount} (0 if neither gave one).
\returns{The base counter value.}

\apientry{setbasecount}{Set the base counter value}
\begin{lstlisting}[language=C,style=signature]
void setbasecount(double count);
\end{lstlisting}
Sets the counter value of sample 0, used by \code{strtim} and written by
\code{newheader} and \code{setheader}. A header read afterwards replaces it.
\code{wfdbquit} does not reset it.
\begin{params}
\item[count] The base counter value.
\end{params}

% ==========================================================================
\section{Amplitude conversion}
\label{sec:library-amplitude}

These functions convert between ADC units and physical units with the gain
(and, for \code{aduphys} and \code{physadu}, the baseline) of an open input
signal. \code{adumuv} and \code{muvadu} convert differences of level (the
baseline is not used), between ADC units and thousandths of the physical
unit (microvolts for a signal in millivolts).

WFDB-x (\cref{chap:behavior}): for an uncalibrated signal (gain 0) or a
number that is not that of an open input signal, C assumed a gain of 200
(\code{WFDB\_DEFGAIN}) and a baseline of 0; WFDB-x reports \code{FUNCTION:
signal N is uncalibrated} or \code{FUNCTION: invalid signal number N}, and
returns \code{WFDB\_INVALID\_SAMPLE} (\code{aduphys}: NaN). Legacy mode keeps
C's behavior.

\apientry{aduphys}{ADC units to physical units}
\begin{lstlisting}[language=C,style=signature]
double aduphys(WFDB_Signal s, WFDB_Sample a);
\end{lstlisting}
Returns \code{(a - baseline) / gain} for input signal \code{s}.
\begin{params}
\item[s] The input signal number.
\item[a] The sample value, in ADC units.
\end{params}
\returns{The value in physical units; WFDB-x: NaN on error.}

\apientry{physadu}{Physical units to ADC units}
\begin{lstlisting}[language=C,style=signature]
WFDB_Sample physadu(WFDB_Signal s, double v);
\end{lstlisting}
Returns \code{v * gain}, rounded to the nearest integer (halves away from
zero), plus the baseline of input signal \code{s}.
\begin{params}
\item[s] The input signal number.
\item[v] The value, in physical units.
\end{params}
\returns{The sample value in ADC units; WFDB-x: \code{WFDB\_INVALID\_SAMPLE}
on error.}

\apientry{adumuv}{ADC units to microvolts}
\begin{lstlisting}[language=C,style=signature]
int adumuv(WFDB_Signal s, WFDB_Sample a);
\end{lstlisting}
Returns \code{a * 1000 / gain} for input signal \code{s}, rounded to the
nearest integer (halves away from zero).
\begin{params}
\item[s] The input signal number.
\item[a] The difference of level, in ADC units.
\end{params}
\returns{The difference in thousandths of the physical unit; WFDB-x:
\code{WFDB\_INVALID\_SAMPLE} on error.}

\apientry{muvadu}{Microvolts to ADC units}
\begin{lstlisting}[language=C,style=signature]
WFDB_Sample muvadu(WFDB_Signal s, int microvolts);
\end{lstlisting}
Returns \code{microvolts * gain / 1000} for input signal \code{s}, rounded to
the nearest integer (halves away from zero).
\begin{params}
\item[s] The input signal number.
\item[microvolts] The difference of level, in thousandths of the physical
  unit.
\end{params}
\returns{The difference in ADC units; WFDB-x: \code{WFDB\_INVALID\_SAMPLE} on
error.}

% ==========================================================================
\section{Annotations}
\label{sec:library-annotations}

Annotation files are opened by \code{annopen} (or \code{wfdbinit}), read
sequentially with \code{getann}, and written with \code{putann}, in MIT
format (\code{WFDB\_READ}, \code{WFDB\_WRITE}) or AHA format
(\code{WFDB\_AHA\_READ}, \code{WFDB\_AHA\_WRITE}); \cref{chap:formats}
describes both formats and lists the annotation codes. An annotation is a
\code{WFDB\_Annotation} structure (\cref{sec:library-structures}).

\apientry{getann}{Read the next annotation}
\begin{lstlisting}[language=C,style=signature]
int getann(WFDB_Annotator a, WFDB_Annotation *annot);
\end{lstlisting}
Reads the next annotation of input annotator \code{a} into \code{annot} (an
annotation pushed back by \code{ungetann} first, if there is one). Its time
is in sample intervals of the input sampling frequency in effect when the
annotator was opened (that of \code{setifreq}, or of the record, in
high-resolution intervals in \code{WFDB\_HIGHRES} mode), converted from the
time resolution recorded in the file if it records one, or in the units set by
\code{setiafreq}. In C, \code{aux} points to a buffer of the annotator that
later calls overwrite.
\begin{params}
\item[a] The input annotator number (0 for the first input annotator of
  \code{annopen}).
\item[annot] The annotation structure, filled by the function.
\end{params}
\returns{0 on success; -1 at the end of the file (\code{*annot} is not
valid); -2 if \code{a} is not an open input annotator; -3 at an unexpected end
of the file.}
\raises{\code{getann: can't read annotator N}; \code{getann: unexpected EOF
in annotator ATR}; \code{getann: unexpected annot number in annotator ATR}.}

\apientry{ungetann}{Push an annotation back}
\begin{lstlisting}[language=C,style=signature]
int ungetann(WFDB_Annotator a, const WFDB_Annotation *annot);
\end{lstlisting}
Makes \code{*annot} the next annotation that \code{getann} returns for input
annotator \code{a}. It need not have been read from that annotator. Only one
annotation can be pushed back at a time for each annotator.
\begin{params}
\item[a] The input annotator number.
\item[annot] The annotation to push back.
\end{params}
\returns{0 on success; -1 if an annotation is already pushed back (this one
is not); -2 if \code{a} is not an input annotator.}
\raises{\code{ungetann: annotator N is not initialized}; \code{ungetann:
pushback buffer is full}.}

\apientry{putann}{Write an annotation}
\begin{lstlisting}[language=C,style=signature]
int putann(WFDB_Annotator a, const WFDB_Annotation *annot);
\end{lstlisting}
Writes \code{*annot} to output annotator \code{a}; all its fields must be
set. Annotations may be written in any order: an annotation is in order if
its time is later than that of the previous annotation, or its \code{chan} or
its \code{num} is greater, or both are at time 0; a file that received one out
of order is sorted when it is closed (see \code{oannclose}). The first
\code{putann} to a file writes, before the annotation, the time resolution set
by \code{setafreq} and the modification labels of the mnemonics and
descriptions changed by \code{setannstr} and \code{setanndesc}. In an AHA
file, at most 6 bytes of \code{aux} are written.
\begin{params}
\item[a] The output annotator number (0 for the first output annotator of
  \code{annopen}).
\item[annot] The annotation, filled by the caller.
\end{params}
\returns{0 on success; -1 on a write error, or for a time equal to
\code{WFDB\_TIME\_MIN} or \code{WFDB\_TIME\_MAX}; -2 if \code{a} is not an
open output annotator.}
\raises{\code{putann: can't write annotation file N}; \code{putann: write
error on annotation file ATR}; \code{putann: time overflow in annotation file
N}.}

\apientry{iannsettime}{Skip to a time in all input annotators}
\begin{lstlisting}[language=C,style=signature]
int iannsettime(WFDB_Time t);
\end{lstlisting}
Positions every input annotator at its first annotation at or after time
\code{|t|} (rewinding the file if needed, which a pipe cannot do), so that a
value of \code{strtim} for a time of day can be given as it is.
\code{WFDB\_TIME\_MIN} rewinds to the very beginning, before annotations at
negative times. A pushed-back annotation is discarded if the file is rewound.
\begin{params}
\item[t] The time, in the units of the annotators' times.
\end{params}
\returns{0 if at least one annotator is positioned on an annotation; -1 if
every annotator reached the end of its file (or none is open), or on an
improper seek. (The Programmer's Guide also lists -3, which the C code never
returns.)}
\raises{\code{iannsettime: improper seek}.}

\apientry{iannclose}{Close an input annotator}
\begin{lstlisting}[language=C,style=signature]
void iannclose(WFDB_Annotator a);
\end{lstlisting}
Closes the file of input annotator \code{a}. The input annotators that follow
are renumbered: annotator \code{a+1} becomes \code{a}, and so on. Nothing is
done if \code{a} is not open.
\begin{params}
\item[a] The input annotator number.
\end{params}

\apientry{oannclose}{Close an output annotator}
\begin{lstlisting}[language=C,style=signature]
void oannclose(WFDB_Annotator a);
\end{lstlisting}
Completes and closes the file of output annotator \code{a} (an end marker in
an MIT file, padding to a 1024-byte block in an AHA file); the output
annotators that follow are renumbered. If annotations were written out of
order, the file is sorted, unless the environment variable \code{WFDBANNSORT}
is 0: the result is that of \code{sortann -r. REC -a ATR}, with the messages
\code{Rearranging annotations for output annotator ATR ...done!}; if it is
not sorted, the advice \code{Use the command: sortann -r REC -a ATR} is
printed. The status is kept for \code{annopen(..., 0)}: -7 after a write
error, -6 after an out-of-order file.
\begin{params}
\item[a] The output annotator number.
\end{params}
\raises{\code{oannclose: write error on annotation file ATR};
\code{Annotations still need to be rearranged.}}

WFDB-x (\cref{chap:behavior}): the file is sorted in the same process (C ran
the program \code{sortann} through the shell), and an annotation at time 0,
number 0 and channel 0 that arrives after later ones is no longer lost, nor
are those written after it; legacy mode runs the \code{sortann} tool of the
same language, with C's result (without a shell; R and MATLAB/Octave quote
its arguments instead). A write error is counted
(\code{wfdbx\_write\_errors}). The high-level \code{write\_annotations} can
decide for its call whether the file is sorted.

\apientry{setafreq}{Set the time resolution of output annotation files}
\begin{lstlisting}[language=C,style=signature]
void setafreq(WFDB_Frequency f);
\end{lstlisting}
Sets the time resolution, in ticks per second, of the output annotation files
to which no annotation has been written yet (contrary to the Programmer's
Guide, this includes those already open): the resolution is recorded at the
beginning of the file, and the times given to \code{putann} are in ticks. 0
(the default) records none: the times are then in sample intervals of the
record. \code{setifreq} calls it. \code{wfdbquit} does not reset it.
\begin{params}
\item[f] The time resolution, in Hz; 0 for none.
\end{params}

\apientry{getafreq}{Time resolution of output annotation files}
\begin{lstlisting}[language=C,style=signature]
WFDB_Frequency getafreq(void);
\end{lstlisting}
Returns the value set by \code{setafreq} (or by \code{setifreq}).
\returns{The time resolution in Hz; 0 if none was set.}

\apientry{setiafreq}{Set the time units of an input annotator}
\begin{lstlisting}[language=C,style=signature]
void setiafreq(WFDB_Annotator a, WFDB_Frequency f);
\end{lstlisting}
Makes \code{getann} return the times of input annotator \code{a} in units of
\code{1/f} second (the times read and pushed back so far are converted too).
\code{annopen} sets it to the input sampling frequency. With \code{f} 0 or
negative, the times are those of the file, unconverted. Nothing is done if
\code{a} is not an open input annotator.
\begin{params}
\item[a] The input annotator number.
\item[f] The frequency, in Hz.
\end{params}

\apientry{getiafreq}{Time units of an input annotator}
\begin{lstlisting}[language=C,style=signature]
WFDB_Frequency getiafreq(WFDB_Annotator a);
\end{lstlisting}
Returns the frequency of the time units that \code{getann} uses for input
annotator \code{a} (see \code{setiafreq}).
\returns{The frequency in Hz; -2 if \code{a} is not an open input annotator.}

\apientry{getiaorigfreq}{Time resolution recorded in an annotation file}
\begin{lstlisting}[language=C,style=signature]
WFDB_Frequency getiaorigfreq(WFDB_Annotator a);
\end{lstlisting}
Returns the time resolution recorded in the file of input annotator \code{a}
by the program that wrote it (see \code{setafreq}).
\returns{The time resolution in Hz; 0 if the file records none (its times are
then in sample intervals of the record's frames); -2 if \code{a} is not an
open input annotator.}

\apientry{annstr}{Mnemonic of an annotation code}
\begin{lstlisting}[language=C,style=signature]
char *annstr(int annotation_code);
\end{lstlisting}
Returns the mnemonic of an annotation code (\code{"N"} for \code{NORMAL}),
from the table that \code{setannstr} and the modification labels of the input
annotation files change (\cref{chap:formats} lists the standard mnemonics).
The codes of the range 0 to \code{ACMAX} that are not defined have mnemonics
such as \code{"[15]"}.
\begin{params}
\item[annotation\_code] The annotation code.
\end{params}
\returns{The mnemonic; for a code outside 0 to \code{ACMAX}, the code in
brackets (\code{"[55]"}).}

\apientry{strann}{Annotation code of a mnemonic}
\begin{lstlisting}[language=C,style=signature]
int strann(const char *annotation_mnemonic_string);
\end{lstlisting}
Returns the code (1 to \code{ACMAX}) whose mnemonic, in the table of
\code{annstr}, is \code{annotation\_mnemonic\_string}; if several codes have
it, the lowest.
\begin{params}
\item[annotation\_mnemonic\_string] The mnemonic.
\end{params}
\returns{The annotation code; \code{NOTQRS} (0) if no code has that mnemonic.}

\apientry{setannstr}{Change the mnemonic of an annotation code}
\begin{lstlisting}[language=C,style=signature]
int setannstr(int annotation_code, const char *annotation_mnemonic_string);
\end{lstlisting}
Sets the mnemonic of a code in the table of \code{annstr} and \code{strann}
(\code{NULL} is taken as \code{""}). With a positive code, the change is also
marked to be recorded, as a modification label, in the output annotation
files opened afterwards, so that other programs read the files with these
mnemonics. With a negative code, the mnemonic of \code{-code} is changed for
this program only. \code{wfdbquit} does not reset the table.
\begin{params}
\item[annotation\_code] The code (or its negation).
\item[annotation\_mnemonic\_string] The new mnemonic.
\end{params}
\returns{0 on success; -1 if \code{|code|} is greater than \code{ACMAX}.}
\raises{\code{setannstr: illegal annotation code N}.}

\apientry{anndesc}{Description of an annotation code}
\begin{lstlisting}[language=C,style=signature]
char *anndesc(int annotation_code);
\end{lstlisting}
Returns the description of an annotation code (\code{"Normal beat"} for
\code{NORMAL}), from the table that \code{setanndesc} and the modification
labels of the input annotation files change.
\begin{params}
\item[annotation\_code] The annotation code.
\end{params}
\returns{The description; \code{NULL} for a code of the range 0 to
\code{ACMAX} that has none; \code{"illegal annotation code"} for a code outside
that range (the Programmer's Guide says \code{NULL}; the C code returns this
string).}

\apientry{setanndesc}{Change the description of an annotation code}
\begin{lstlisting}[language=C,style=signature]
int setanndesc(int annotation_code, const char *annotation_description);
\end{lstlisting}
As \code{setannstr}, for the table of \code{anndesc}: a positive code also
marks the description to be recorded in the output annotation files opened
afterwards, a negative code changes the description of \code{-code} for this
program only.
\begin{params}
\item[annotation\_code] The code (or its negation).
\item[annotation\_description] The new description.
\end{params}
\returns{0 on success; -1 if \code{|code|} is greater than \code{ACMAX}.}
\raises{\code{setanndesc: illegal annotation code N}.}

\apientry{ecgstr}{Standard mnemonic of an annotation code}
\begin{lstlisting}[language=C,style=signature]
char *ecgstr(int annotation_code);
\end{lstlisting}
As \code{annstr}, from a separate table that only \code{setecgstr} changes
(not the modification labels of annotation files), for programs that need a
fixed set of mnemonics.
\begin{params}
\item[annotation\_code] The annotation code.
\end{params}
\returns{The mnemonic; the code in brackets for a code outside 0 to
\code{ACMAX}.}

\apientry{strecg}{Annotation code of a standard mnemonic}
\begin{lstlisting}[language=C,style=signature]
int strecg(const char *annotation_mnemonic_string);
\end{lstlisting}
As \code{strann}, with the table of \code{ecgstr}.
\begin{params}
\item[annotation\_mnemonic\_string] The mnemonic.
\end{params}
\returns{The annotation code; \code{NOTQRS} (0) if no code has that mnemonic.}

\apientry{setecgstr}{Change a standard mnemonic}
\begin{lstlisting}[language=C,style=signature]
int setecgstr(int annotation_code, const char *annotation_mnemonic_string);
\end{lstlisting}
Sets the mnemonic of a code in the table of \code{ecgstr} and \code{strecg}
(\code{NULL} is taken as \code{""}). Nothing is recorded in annotation files,
and negative codes are illegal.
\begin{params}
\item[annotation\_code] The code, 0 to \code{ACMAX}.
\item[annotation\_mnemonic\_string] The new mnemonic.
\end{params}
\returns{0 on success; -1 if the code is outside 0 to \code{ACMAX}.}
\raises{\code{setecgstr: illegal annotation code N}.}

\subsection{Classification and mapping of annotation codes}
\label{sec:library-ecgmap}

\file{<wfdb/ecgmap.h>} defines macros that classify annotation codes and map
them to smaller sets, with tables that a program can change; the C library
provides each of them also as a function, whose name begins with
\code{wfdb\_} (\code{wfdb\_isqrs} for \code{isqrs}). WFDB-x provides both
names in every language, with the same results; the macro names are
functions too (in C++, inline functions of \file{<wfdbx/ecgmap.hpp>}). A
setter stores the new value as a C \code{char}, and returns it as stored:
\code{setisqrs(2, 300)} returns 44. An argument that is not a legal code
(1 to \code{ACMAX}) changes nothing. In R, these functions accept vectors of
codes. \code{isann} (and \code{wfdb\_isann}) returns a logical value in
Python, R and MATLAB/Octave (\code{True}, \code{TRUE}, \code{true}), and 1 or
0 elsewhere.

\apientry{isann}{Whether a code is a legal annotation code}
\begin{lstlisting}[language=C,style=signature]
int isann(int code);              /* a macro of <wfdb/ecgmap.h> */
\end{lstlisting}
Tells whether \code{code} is a legal annotation code, 1 to \code{ACMAX} (49).
\begin{params}
\item[code] The annotation code.
\end{params}
\returns{1 if it is legal, 0 otherwise.}

\apientry{wfdb\_isann}{Function form of isann}
\begin{lstlisting}[language=C,style=signature]
int wfdb_isann(int code);
\end{lstlisting}
The function form of \code{isann}, with the same result.
\begin{params}
\item[code] The annotation code.
\end{params}
\returns{1 if it is legal, 0 otherwise.}

\apientry{isqrs}{Whether a code is a beat}
\begin{lstlisting}[language=C,style=signature]
int isqrs(int code);              /* a macro of <wfdb/ecgmap.h> */
\end{lstlisting}
Tells whether annotation code \code{code} denotes a QRS complex (a beat), from
a table that \code{setisqrs} changes. Initially, the beat codes are 1 to 13,
25, 30, 31, 34, 35, 38 and 41.
\begin{params}
\item[code] The annotation code.
\end{params}
\returns{The value of the table (initially 1 for a beat, 0 otherwise); 0 for
an illegal code.}

\apientry{wfdb\_isqrs}{Function form of isqrs}
\begin{lstlisting}[language=C,style=signature]
int wfdb_isqrs(int code);
\end{lstlisting}
The function form of \code{isqrs}, with the same result.
\begin{params}
\item[code] The annotation code.
\end{params}
\returns{As \code{isqrs}.}

\apientry{setisqrs}{Change the result of isqrs}
\begin{lstlisting}[language=C,style=signature]
int setisqrs(int code, int newval);   /* a macro of <wfdb/ecgmap.h> */
\end{lstlisting}
Sets the value that \code{isqrs(code)} returns.
\begin{params}
\item[code] The annotation code.
\item[newval] The new value.
\end{params}
\returns{The value stored (\code{newval} as a C \code{char}); 0 for an
illegal code.}

\apientry{wfdb\_setisqrs}{Function form of setisqrs}
\begin{lstlisting}[language=C,style=signature]
int wfdb_setisqrs(int code, int newval);
\end{lstlisting}
The function form of \code{setisqrs}, with the same result.
\begin{params}
\item[code] The annotation code.
\item[newval] The new value.
\end{params}
\returns{As \code{setisqrs}.}

\apientry{map1}{Map a code to the five basic classes}
\begin{lstlisting}[language=C,style=signature]
int map1(int code);               /* a macro of <wfdb/ecgmap.h> */
\end{lstlisting}
Maps an annotation code to one of \code{NOTQRS}, \code{NORMAL}, \code{PVC},
\code{FUSION} and \code{LEARN}, from a table that \code{setmap1} changes (for
example, the bundle branch block beats map to \code{NORMAL}, ventricular
escape beats to \code{PVC}).
\begin{params}
\item[code] The annotation code.
\end{params}
\returns{The mapped code; \code{NOTQRS} for an illegal code.}

\apientry{wfdb\_map1}{Function form of map1}
\begin{lstlisting}[language=C,style=signature]
int wfdb_map1(int code);
\end{lstlisting}
The function form of \code{map1}, with the same result.
\begin{params}
\item[code] The annotation code.
\end{params}
\returns{As \code{map1}.}

\apientry{setmap1}{Change the result of map1}
\begin{lstlisting}[language=C,style=signature]
int setmap1(int code, int newval);    /* a macro of <wfdb/ecgmap.h> */
\end{lstlisting}
Sets the value that \code{map1(code)} returns.
\begin{params}
\item[code] The annotation code.
\item[newval] The new value.
\end{params}
\returns{The value stored (as a C \code{char}); \code{NOTQRS} for an illegal
code.}

\apientry{wfdb\_setmap1}{Function form of setmap1}
\begin{lstlisting}[language=C,style=signature]
int wfdb_setmap1(int code, int newval);
\end{lstlisting}
The function form of \code{setmap1}, with the same result.
\begin{params}
\item[code] The annotation code.
\item[newval] The new value.
\end{params}
\returns{As \code{setmap1}.}

\apientry{map2}{Map a code to the six basic classes}
\begin{lstlisting}[language=C,style=signature]
int map2(int code);               /* a macro of <wfdb/ecgmap.h> */
\end{lstlisting}
As \code{map1}, with \code{SVPB} as a sixth class: the supraventricular
premature beats (\code{ABERR}, \code{NPC}, \code{APC}, \code{SVPB}) map to
\code{SVPB} instead of \code{NORMAL}. \code{setmap2} changes its table.
\begin{params}
\item[code] The annotation code.
\end{params}
\returns{The mapped code; \code{NOTQRS} for an illegal code.}

\apientry{wfdb\_map2}{Function form of map2}
\begin{lstlisting}[language=C,style=signature]
int wfdb_map2(int code);
\end{lstlisting}
The function form of \code{map2}, with the same result.
\begin{params}
\item[code] The annotation code.
\end{params}
\returns{As \code{map2}.}

\apientry{setmap2}{Change the result of map2}
\begin{lstlisting}[language=C,style=signature]
int setmap2(int code, int newval);    /* a macro of <wfdb/ecgmap.h> */
\end{lstlisting}
Sets the value that \code{map2(code)} returns.
\begin{params}
\item[code] The annotation code.
\item[newval] The new value.
\end{params}
\returns{The value stored (as a C \code{char}); \code{NOTQRS} for an illegal
code.}

\apientry{wfdb\_setmap2}{Function form of setmap2}
\begin{lstlisting}[language=C,style=signature]
int wfdb_setmap2(int code, int newval);
\end{lstlisting}
The function form of \code{setmap2}, with the same result.
\begin{params}
\item[code] The annotation code.
\item[newval] The new value.
\end{params}
\returns{As \code{setmap2}.}

\apientry{ammap}{Map an AHA annotation code to an MIT code}
\begin{lstlisting}[language=C,style=signature]
int ammap(int code);              /* a macro of <wfdb/ecgmap.h> */
\end{lstlisting}
Maps an AHA annotation code, a character (its character code, such as that of
\code{'V'}), to an MIT annotation code: \code{'N'} to \code{NORMAL},
\code{'V'} to \code{PVC}, \code{'F'} to \code{FUSION}, \code{'E'} to
\code{VESC}, \code{'P'} to \code{PACE}, \code{'Q'} to \code{UNKNOWN},
\code{'R'} to \code{RONT}, \code{'O'} to \code{NOTE}, \code{'U'} to
\code{NOISE}, \code{'['} to \code{VFON}, \code{']'} to \code{VFOFF}. The table
is fixed.
\begin{params}
\item[code] The AHA code.
\end{params}
\returns{The MIT code; \code{NOTQRS} for any other character.}

\apientry{wfdb\_ammap}{Function form of ammap}
\begin{lstlisting}[language=C,style=signature]
int wfdb_ammap(int code);
\end{lstlisting}
The function form of \code{ammap}, with the same result.
\begin{params}
\item[code] The AHA code.
\end{params}
\returns{As \code{ammap}.}

\apientry{mamap}{Map an MIT annotation code to an AHA code}
\begin{lstlisting}[language=C,style=signature]
int mamap(int code, int subtype);     /* a macro of <wfdb/ecgmap.h> */
\end{lstlisting}
Maps an MIT annotation code, with its subtype, to an AHA code (a character
code), from a fixed table: \code{NOISE} maps to \code{'U'} only with subtype
-1 (unreadable signals), and otherwise to \code{'O'}, as do the codes without
an AHA equivalent.
\begin{params}
\item[code] The MIT annotation code.
\item[subtype] Its subtype.
\end{params}
\returns{The AHA code; \code{'O'} for an illegal code.}

\apientry{wfdb\_mamap}{Function form of mamap}
\begin{lstlisting}[language=C,style=signature]
int wfdb_mamap(int code, int subtype);
\end{lstlisting}
The function form of \code{mamap}, with the same result.
\begin{params}
\item[code] The MIT annotation code.
\item[subtype] Its subtype.
\end{params}
\returns{As \code{mamap}.}

\apientry{annpos}{Display position of an annotation code}
\begin{lstlisting}[language=C,style=signature]
int annpos(int code);             /* a macro of <wfdb/ecgmap.h> */
\end{lstlisting}
Returns where programs that display annotations place those of code
\code{code}, from a table that \code{setannpos} changes: \code{APSTD},
\code{APHIGH}, \code{APLOW}, \code{APATT}, \code{APAHIGH}, \code{APALOW} or
\code{APUNDEF} (\cref{tab:library-constants}).
\begin{params}
\item[code] The annotation code.
\end{params}
\returns{The position code; \code{APUNDEF} for an illegal code.}

\apientry{wfdb\_annpos}{Function form of annpos}
\begin{lstlisting}[language=C,style=signature]
int wfdb_annpos(int code);
\end{lstlisting}
The function form of \code{annpos}, with the same result.
\begin{params}
\item[code] The annotation code.
\end{params}
\returns{As \code{annpos}.}

\apientry{setannpos}{Change the result of annpos}
\begin{lstlisting}[language=C,style=signature]
int setannpos(int code, int newval);  /* a macro of <wfdb/ecgmap.h> */
\end{lstlisting}
Sets the value that \code{annpos(code)} returns.
\begin{params}
\item[code] The annotation code.
\item[newval] The new position code.
\end{params}
\returns{The value stored (as a C \code{char}); \code{APUNDEF} for an
illegal code.}

\apientry{wfdb\_setannpos}{Function form of setannpos}
\begin{lstlisting}[language=C,style=signature]
int wfdb_setannpos(int code, int newval);
\end{lstlisting}
The function form of \code{setannpos}, with the same result.
\begin{params}
\item[code] The annotation code.
\item[newval] The new position code.
\end{params}
\returns{As \code{setannpos}.}

% ==========================================================================
\section{Calibration}
\label{sec:library-calibration}

The calibration functions manage the \emph{calibration list}, kept in
memory: for types of signals, the levels and shape of their calibration
pulses and their customary plotting scale (\code{WFDB\_Calinfo},
\cref{sec:library-structures}). A calibration file has one entry per line:
the signal type (\code{*} for any), a tab, then the low and high levels
(\code{-} for an AC-coupled signal's low level, or for a high level equal to
the low level), the pulse shape (\code{square}, \code{sine},
\code{sawtooth}, or anything else for an undefined shape), the scale and the
units, separated by spaces or tabs; lines that begin with \code{\#} and
incomplete lines are ignored.

\apientry{calopen}{Read a calibration file}
\begin{lstlisting}[language=C,style=signature]
int calopen(const char *calibration_filename);
\end{lstlisting}
Reads a calibration file into the calibration list, which is emptied first,
unless the name begins with \code{+}, which is then removed (the entries are
appended). \code{NULL} reads the file named by the environment variable
\code{WFDBCAL}, or else \file{wfdbcal}; \code{-} reads the standard input.
\begin{params}
\item[calibration\_filename] The name of the calibration file, or
  \code{NULL}.
\end{params}
\returns{0 on success; -2 if the file cannot be opened. (The Programmer's
Guide also lists -1, for a lack of memory.)}
\raises{\code{calopen: can't read calibration file FILE}.}

WFDB-x: the file is read only where its name says, under the input directory
for a relative name, not searched for in the WFDB path (\cref{chap:behavior}).

\apientry{getcal}{Find the calibration of a signal}
\begin{lstlisting}[language=C,style=signature]
int getcal(const char *description, const char *units, WFDB_Calinfo *cal);
\end{lstlisting}
Copies into \code{*cal} the first entry of the calibration list whose signal
type is a prefix of \code{description} (or equal to it), and whose units are
\code{units}. A \code{NULL} argument matches every entry.
\begin{params}
\item[description] The signal description (\code{desc} of a
  \code{WFDB\_Siginfo}), or \code{NULL}.
\item[units] The physical units, or \code{NULL}.
\item[cal] The structure filled by the function.
\end{params}
\returns{0 if an entry matches; -1 otherwise (also if the list is empty).}

In R and MATLAB/Octave, \code{getcal} takes \code{description} and
\code{units} only, and returns the entry with the status: \code{list(stat,
cal)} in R, \code{[stat, cal]} in MATLAB/Octave.

\apientry{putcal}{Add an entry to the calibration list}
\begin{lstlisting}[language=C,style=signature]
int putcal(const WFDB_Calinfo *cal);
\end{lstlisting}
Appends a copy of \code{*cal} to the end of the calibration list.
\begin{params}
\item[cal] The entry.
\end{params}
\returns{0. (The Programmer's Guide also lists -1, for a lack of memory.)}

\apientry{newcal}{Write the calibration list to a file}
\begin{lstlisting}[language=C,style=signature]
int newcal(const char *calibration_filename);
\end{lstlisting}
Writes the calibration list, which is not changed, to a new file in the
format read by \code{calopen}, with \code{undefined} for an undefined pulse
shape. The name may contain only letters, digits, \code{\_},
\code{\textasciitilde}, \code{-} and \code{/}.
\begin{params}
\item[calibration\_filename] The name of the file.
\end{params}
\returns{0 on success; -1 if the name is illegal, the file cannot be created,
or it cannot be written completely.}
\raises{\code{newcal: can't create calibration file FILE}; \code{newcal: write
error in calibration file}; \code{init: illegal character N in calibration
file name}.}

WFDB-x: the file is created under the output directory; a write error is
counted (\code{wfdbx\_write\_errors}).

\apientry{flushcal}{Empty the calibration list}
\begin{lstlisting}[language=C,style=signature]
void flushcal(void);
\end{lstlisting}
Empties the calibration list. \code{wfdbquit} does not.

% ==========================================================================
\section{Header and info strings}
\label{sec:library-header}

A header (\file{.hea} file) describes a record: \code{newheader} writes it
for the signals written with \code{putvec}, \code{setheader} from signal
specifications, \code{setmsheader} for a multi-segment record. The header
being written stays open, so that \code{putinfo} can add info strings to it,
until another header is written or \code{osigfopen(NULL, 0)} (or
\code{wfdbquit}) closes it. Info strings are the comment lines of a header
that follow the record line, and the lines of an optional \file{.info} file;
\cref{chap:metadata} describes the \code{Key: value} entries that WFDB-x
reads and writes in them.

WFDB-x (\cref{chap:behavior}): headers are created under the output
directory; a header written under a path (\code{out/rec}) names the record,
its signal files and its segments relative to its own directory
(\code{rec}, \code{rec.dat}), and only the record's own name (the last
component) must be made of letters, digits, \code{\_}, \code{-} and
\code{\textasciitilde}. Legacy mode checks the whole name and writes the
names as given.

\apientry{newheader}{Write the header of the output signals}
\begin{lstlisting}[language=C,style=signature]
int newheader(char *record);
\end{lstlisting}
Writes the header of \code{record} for the open output signals, as
\code{setheader} does, with the numbers of samples and the checksums that
\code{putvec} computed: it is called after the samples are written, before the
output signals are closed. The frequencies and the base time are those set
by \code{setsampfreq}, \code{setcfreq}, \code{setbasecount} and
\code{setbasetime}, and the skews and prolog lengths those set by
\code{wfdbsetskew} and \code{wfdbsetstart}. The record name \code{-} writes
the header to the standard output. A record name that begins with
\code{+} is not accepted (the Programmer's Guide says that the \code{+} is
removed; the code rejects the character).
\begin{params}
\item[record] The record name.
\end{params}
\returns{0 on success; -1 if the record name is illegal or the header cannot
be created.}
\raises{\code{init: illegal character N in record name}; \code{newheader:
can't create header for record REC}.}

\apientry{setheader}{Write a header from signal specifications}
\begin{lstlisting}[language=C,style=signature]
int setheader(char *record, const WFDB_Siginfo *siarray, unsigned int nsig);
\end{lstlisting}
Writes (or rewrites) the header of \code{record} for the \code{nsig} signals
described by \code{siarray}, with the frequencies, base time, skews and prolog
lengths as for \code{newheader}; the length is the \code{nsamp} of the first
signal. It serves to change a header, for example the gain of a signal: the
specifications are read with \code{isigopen(record, siarray, -nsig)}, changed
and written back; the info strings must be copied explicitly
(\code{getinfo}, \code{putinfo}). As in C, and although the argument is
declared \code{const}, the \code{units} of each structure are cut at their
first blank and the \code{desc} at their first line break, in the caller's
structures (in R and MATLAB/Octave, the changed structures are returned:
\code{list(stat, si)}, \code{[stat, si]}).
\begin{params}
\item[record] The record name.
\item[siarray] The signal specifications.
\item[nsig] The number of signals.
\end{params}
\returns{0 on success; -1 if the record name is illegal or the header cannot
be created.}
\raises{\code{newheader: can't create header for record REC}.}

\apientry{setmsheader}{Write the header of a multi-segment record}
\begin{lstlisting}[language=C,style=signature]
int setmsheader(char *record, char **segnames, unsigned int nsegments);
\end{lstlisting}
Closes the input signals, and writes the header of the multi-segment record
\code{record}, made of the \code{nsegments} existing records named in
\code{segnames}, in that order. The segments' signal files are neither
copied nor changed; the readers of the record (\code{getvec},
\code{isigsettime}, ...) go from one segment to the next. The segments must
have the same number of signals and the same sampling frequency, and headers
that give their length; the record takes the counter frequency, base counter
value and base time and date of the first segment.
\begin{params}
\item[record] The name of the multi-segment record.
\item[segnames] The names of the segments.
\item[nsegments] The number of segments.
\end{params}
\returns{0 on success; -1 if the record name is illegal, there is no segment,
or the header cannot be created; -2 if a segment name is longer than
\code{WFDB\_MAXRNL} (50) characters; -3 if the header of a segment cannot be
read (or is itself that of a multi-segment record); -4 if a segment's length
is not given, or its number of signals or sampling frequency differs from
those of the first segment.}
\raises{\code{setmsheader: record must contain at least one segment};
\code{setmsheader: can't read segment SEG header}; \code{setmsheader: length
of segment SEG must be specified}; \code{setmsheader: incorrect number of
signals in segment SEG}; \code{setmsheader: incorrect sampling frequency in
segment SEG}; \code{setmsheader: can't create header file for record REC}.}

\apientry{getseginfo}{The segments of the open record}
\begin{lstlisting}[language=C,style=signature]
int getseginfo(WFDB_Seginfo **segments);
\end{lstlisting}
Makes \code{*segments} point to the array of \code{WFDB\_Seginfo} structures
(\cref{sec:library-structures}) that describe the segments of the open
multi-segment record: the name, the length and the first sample number of
each. A segment named \code{\textasciitilde} is a gap, in which no signal is
available. If no multi-segment record is open, \code{*segments} is not
changed.
\begin{params}
\item[segments] The address of a pointer, set by the function.
\end{params}
\returns{The number of segments; 0 if the open record is not a
multi-segment record.}

The form differs between languages: in C++, \code{getseginfo(const
WFDB\_Seginfo **)} gives a pointer to the library's table; in Rust, Julia and
Kotlin, the argument (\code{\&mut Option<Vec<Seginfo>>}, a \code{Ref}, a
\code{MutableList<Seginfo>?}) receives a copy, and the number of segments is
returned; in Python, R and MATLAB/Octave, \code{getseginfo()} takes no
argument and returns the list (a struct array in MATLAB/Octave) of segments,
empty if the record is not a multi-segment record.

\apientry{getinfo}{Read the info strings of a record}
\begin{lstlisting}[language=C,style=signature]
char *getinfo(char *record);
\end{lstlisting}
With a record name, reads the info strings of \code{record} (the comment
lines of its header after the record line, then the lines of its
\file{.info} file that begin with \code{\#}, if it has one) and returns the
first; with \code{NULL}, returns the next one, reading those of the current
record (the last one opened or named) on the first call. An info string is
the line without its \code{\#} and its line ending.
\begin{params}
\item[record] The record name, or \code{NULL}.
\end{params}
\returns{The info string; \code{NULL} if there are no more.}
\raises{\code{getinfo: caller did not specify record name}.}

WFDB-x: the header and the \file{.info} file are read only where the record
name says (\cref{chap:behavior}).

\apientry{putinfo}{Write an info string}
\begin{lstlisting}[language=C,style=signature]
int putinfo(const char *info);
\end{lstlisting}
Writes \code{info} as an info string (a line \code{\#}\emph{info}, with a
CR-LF line ending): to the \file{.info} file opened by \code{setinfo}, or
else to the header being written (after \code{newheader}, \code{setheader}
or \code{setmsheader}). It should not contain a line break; WFDB-x does not
limit its length (the Programmer's Guide gives 254 characters).
\begin{params}
\item[info] The info string.
\end{params}
\returns{0 on success; -1 if no info file or header is open.}
\raises{\code{putinfo: caller has not specified a record name}.}

\apientry{setinfo}{Open the info file of a record}
\begin{lstlisting}[language=C,style=signature]
int setinfo(char *record);
\end{lstlisting}
Closes the info file opened by an earlier \code{setinfo}, if any, and opens
the \file{.info} file of \code{record} for writing with \code{putinfo},
appending to it if it exists. Such a file holds information about the record
without changing its header. \code{setinfo(NULL)} only closes the info file.
\begin{params}
\item[record] The record name, or \code{NULL}.
\end{params}
\returns{0 on success; -1 if the record name is illegal or the file cannot be
opened; with \code{NULL}, -2 if the info file could not be written
completely.}
\raises{\code{setinfo: can't create info file for record REC};
\code{setinfo: write error in info file}.}

\apientry{wfdb\_freeinfo}{Release the info strings}
\begin{lstlisting}[language=C,style=signature]
void wfdb_freeinfo(void);
\end{lstlisting}
Releases the info strings read by \code{getinfo}: the next \code{getinfo}
reads them again, from the first. \code{wfdbquit} calls it. (The Programmer's
Guide says that it also closes the file opened by \code{setinfo}; the code
does not, and neither does WFDB-x: \code{setinfo(NULL)} and \code{wfdbquit}
close it.)

\apientry{wfdbgetskew}{Skew of an input signal}
\begin{lstlisting}[language=C,style=signature]
int wfdbgetskew(WFDB_Signal s);
\end{lstlisting}
Returns the skew of input signal \code{s} recorded in the header, in frames
(the header gives it in samples). \code{getvec} and \code{getframe} already
correct the skew; the function serves programs that rewrite a header and
must keep the skews (with \code{wfdbsetskew}).
\begin{params}
\item[s] The input signal number.
\end{params}
\returns{The skew in frames; 0 if \code{s} is not an input signal.}

\apientry{wfdbsetiskew}{Change the skew of an input signal}
\begin{lstlisting}[language=C,style=signature]
void wfdbsetiskew(WFDB_Signal s, int skew);
\end{lstlisting}
Changes the skew, in frames, with which input signal \code{s} is read (a
function of the C library that the Programmer's Guide does not document). The
new skew is applied only if it is between 0 and the largest skew of the
record's header; otherwise nothing changes.
\begin{params}
\item[s] The input signal number.
\item[skew] The skew, in frames.
\end{params}

\apientry{wfdbsetskew}{Set the skew of an output signal}
\begin{lstlisting}[language=C,style=signature]
void wfdbsetskew(WFDB_Signal s, int skew);
\end{lstlisting}
Sets the skew, in frames, that \code{newheader} and \code{setheader} write
for output signal \code{s} (in samples). It does not affect the reading of
input signals. Nothing is done if \code{s} is not an output signal.
\begin{params}
\item[s] The output signal number.
\item[skew] The skew, in frames.
\end{params}

\apientry{wfdbgetstart}{Prolog length of an input signal file}
\begin{lstlisting}[language=C,style=signature]
long wfdbgetstart(WFDB_Signal s);
\end{lstlisting}
Returns the length of the prolog (the bytes before sample 0) of the signal
file of input signal \code{s}, as the header records it; the readers already
skip it. With no input signals open, signal 0 gives the prolog of the first
signal of the last header read.
\begin{params}
\item[s] The input signal number.
\end{params}
\returns{The prolog length in bytes; 0 if \code{s} is not an input signal.}

\apientry{wfdbsetstart}{Set the prolog length of an output signal}
\begin{lstlisting}[language=C,style=signature]
void wfdbsetstart(WFDB_Signal s, long bytes);
\end{lstlisting}
Sets the prolog length that \code{newheader} and \code{setheader} write for
the signal file of output signal \code{s}; it does not affect the reading of
input signals. The value also applies to every signal of the next header
written by \code{setheader} whose file has no prolog length of its own (as
when no output signals are open).
\begin{params}
\item[s] The output signal number.
\item[bytes] The prolog length, in bytes.
\end{params}

\apientry{wfdbputprolog}{Write a prolog to an output signal file}
\begin{lstlisting}[language=C,style=signature]
int wfdbputprolog(const char *prolog, long bytes, WFDB_Signal s);
\end{lstlisting}
Writes the first \code{bytes} bytes of \code{prolog} to the signal file of
output signal \code{s} (normally before any sample), and records the number
of bytes written with \code{wfdbsetstart}.
\begin{params}
\item[prolog] The bytes of the prolog.
\item[bytes] Their number.
\item[s] The output signal number.
\end{params}
\returns{0 on success; -1 if fewer bytes were written.}
\raises{\code{wfdbputprolog: only N of M bytes written}.}

% ==========================================================================
\section{Error handling and settings}
\label{sec:library-errors}

The library reports an error by its return value and by a message, printed on
the standard error output unless \code{wfdbquiet} was called; the last
message is kept for \code{wfdberror}.

\apientry{wfdbquiet}{Stop printing error messages}
\begin{lstlisting}[language=C,style=signature]
void wfdbquiet(void);
\end{lstlisting}
Stops the printing of the library's messages on the standard error output;
they are still kept for \code{wfdberror}.

\apientry{wfdbverbose}{Print error messages again}
\begin{lstlisting}[language=C,style=signature]
void wfdbverbose(void);
\end{lstlisting}
Restores the printing of the library's messages (the default).

\apientry{wfdberror}{The last error message}
\begin{lstlisting}[language=C,style=signature]
char *wfdberror(void);
\end{lstlisting}
Returns the last message of the library, printed or not.
\returns{The message, with its final newline; if there has been none, the
library version, \code{"WFDB library version 10.7.0 (Jun 10 2022).\textbackslash
n"} (the build date of the C release reproduced).}

\apientry{wfdb\_error}{Report an error as the library does}
\begin{lstlisting}[language=C,style=signature]
void wfdb_error(const char *format_string, ...);
\end{lstlisting}
Formats a message as \code{printf} does, keeps it for \code{wfdberror} and
prints it on the standard error output unless \code{wfdbquiet} was called.
The library reports its own errors with it; a program may use it for its
own messages.
\begin{params}
\item[format\_string] A \code{printf} format.
\item[...] The values to format.
\end{params}

In Python, \code{wfdb\_error(fmt, *args)}; in Rust and Kotlin,
\code{wfdb\_error(msg)} takes a message already formatted (Rust also has
\code{wfdb\_error\_fmt(fmt, args)}); in Julia, \code{WFDBX.wfdb\_error(fmt,
args...)} (not exported). The formatting is that of the GNU C library in
every language.

\apientry{wfdbmemerr}{Behavior on memory errors}
\begin{lstlisting}[language=C,style=signature]
void wfdbmemerr(int exit_on_error);
\end{lstlisting}
In C, chooses whether a failure to allocate memory ends the program (nonzero,
the default) or lets the function continue (0). WFDB-x keeps the value, for
\code{wfdb\_me\_fatal}, but has no other use for it: the languages report
memory exhaustion themselves.
\begin{params}
\item[exit\_on\_error] Nonzero to end the program on a memory error.
\end{params}

\apientry{wfdb\_me\_fatal}{Whether memory errors are fatal}
\begin{lstlisting}[language=C,style=signature]
int wfdb_me_fatal(void);
\end{lstlisting}
Returns the value set by \code{wfdbmemerr}.
\returns{The value; 1 by default.}

\apientry{wfdbversion}{Version of the library}
\begin{lstlisting}[language=C,style=signature]
const char *wfdbversion(void);
\end{lstlisting}
Returns the version of the WFDB library that WFDB-x reproduces (the version
of WFDB-x itself is that of its packages, \wfdbxversion).
\returns{\code{"10.7.0"}.}

\apientry{wfdbldflags}{Linker flags of the C library}
\begin{lstlisting}[language=C,style=signature]
const char *wfdbldflags(void);
\end{lstlisting}
In C, returns the linker options with which programs are linked with the
library (fixed when it is built). WFDB-x has none.
\returns{\code{""}.}

\apientry{wfdbcflags}{Compiler flags of the C library}
\begin{lstlisting}[language=C,style=signature]
const char *wfdbcflags(void);
\end{lstlisting}
In C, returns the compiler options with which programs that use the library
are compiled. WFDB-x has none.
\returns{\code{""}.}

\apientry{wfdbdefwfdb}{Default WFDB path}
\begin{lstlisting}[language=C,style=signature]
const char *wfdbdefwfdb(void);
\end{lstlisting}
Returns the default WFDB path of the C library, used in legacy mode when the
environment variable \code{WFDB} is not set (that of a C release installed in
\file{/usr/local}).
\returns{\code{". /usr/local/database http://physionet.org/physiobank/database"}.}

\apientry{wfdbdefwfdbcal}{Default calibration file}
\begin{lstlisting}[language=C,style=signature]
const char *wfdbdefwfdbcal(void);
\end{lstlisting}
Returns the name of the calibration file that \code{calopen(NULL)} reads when
the environment variable \code{WFDBCAL} is not set.
\returns{\code{"wfdbcal"}.}

% ==========================================================================
\section{File and path functions}
\label{sec:library-files}

The C library finds its input files by searching the WFDB path, a list of
directories and URLs. WFDB-x finds a file only where its name says, under the
input directory for a relative name (\cref{chap:behavior,chap:files}): the
functions of this section keep their C names and arguments, and set and
return the input directory; in legacy mode, they work on the WFDB path as in
C, without network access.

\apientry{setwfdb}{Set the input directory (the WFDB path)}
\begin{lstlisting}[language=C,style=signature]
void setwfdb(const char *database_path_string);
\end{lstlisting}
WFDB-x: sets the input directory. The string is split as a WFDB path is (at
spaces, tabs, newlines, \code{;}, and \code{:} outside a URL), and must have at
most one nonempty component, which becomes the input directory (a directory,
a \code{\textasciitilde} path or a URL); with none (\code{""}, \code{";"},
\code{NULL}), the default is restored (\code{WFDBX\_INPUT\_DIR}, or the
current directory). With more than one, nothing changes and an error is
reported. A directory whose name contains a space must be set with
\code{wfdbx\_set\_input\_dir}.

In legacy mode, as in C: sets the WFDB path, the list of locations searched
for input files (\code{@file} reads it from a file, up to ten levels deep);
\code{NULL} restores the value of the environment variable \code{WFDB}, or
the default (\code{wfdbdefwfdb}); the value is also exported as the
environment variable \code{WFDB}, for the programs started afterwards; URL
components are not used.
\begin{params}
\item[database\_path\_string] The input directory (legacy mode: the WFDB
  path), or \code{NULL}.
\end{params}
\raises{\code{setwfdb: WFDB-x reads from a single input directory; use
wfdbx\_set\_input\_dir (or legacy mode for the WFDB search path)}.}

\apientry{getwfdb}{The input directory (the WFDB path)}
\begin{lstlisting}[language=C,style=signature]
char *getwfdb(void);
\end{lstlisting}
WFDB-x: returns the input directory. In legacy mode, as in C: returns the
WFDB path, read on the first call from the environment variable \code{WFDB}
(or the default), \code{@file} indirections evaluated. The environment
variable \code{WFDB} is ignored outside legacy mode, and the library says so
once (\code{WFDB-x: the WFDB environment variable is ignored; set
WFDBX\_INPUT\_DIR instead (or use legacy mode)}).
\returns{The input directory, or \code{"."} if none is set (legacy mode: the
WFDB path).}

\apientry{resetwfdb}{Restore the input directory (the WFDB path)}
\begin{lstlisting}[language=C,style=signature]
void resetwfdb(void);
\end{lstlisting}
WFDB-x: restores the input directory set by the last \code{setwfdb}, if any
(\code{wfdbx\_set\_input\_dir} may have changed it since). In legacy mode, as
in C: restores the WFDB path to the value that the first \code{getwfdb} of the
process returned, if \code{getwfdb} was called before \code{setwfdb};
\code{wfdbquit} calls it (in legacy mode only, in WFDB-x).

\apientry{wfdbfile}{The path of a WFDB file}
\begin{lstlisting}[language=C,style=signature]
char *wfdbfile(const char *file_type, char *record);
\end{lstlisting}
Returns the name under which the library opens the file of \code{record}
with suffix \code{file\_type} (\code{"hea"}, \code{"atr"}, ...), if it can be
opened: \code{wfdbfile("hea", "100s")} gives \file{100s.hea}, or the path
under the input directory. A file whose name does not contain a record name
is given by \code{file\_type} with \code{record} \code{NULL}. As in C, a
second name, with the suffix cut to 3 characters, is tried. With both
arguments \code{NULL}, it returns the name of the file most recently opened
by the library.
\begin{params}
\item[file\_type] The file type (suffix), or a file name.
\item[record] The record name, or \code{NULL}.
\end{params}
\returns{The path name or URL; \code{NULL} if the file cannot be opened.}

WFDB-x: the name is the one location that the name designates (an absolute
name or URL as given, a relative name under the input directory), not the
first found in the WFDB path; a URL is opened only with network access
(\cref{chap:behavior}).

% ==========================================================================
\section{WFDB-x additions}
\label{sec:library-additions}

These functions are not in the C library. Their names begin with
\code{wfdbx\_}, except the reader of many sample vectors; they are the same
in the seven languages, and the signatures below are those of the C++
implementation (\file{cpp/include/wfdbx/wfdb.hpp}). \Cref{chap:behavior}
explains the settings they control; the environment variables
\code{WFDBX\_INPUT\_DIR}, \code{WFDBX\_OUTPUT\_DIR} and \code{WFDBX\_NETWORK}
give their initial values, which a call overrides.

\apientry{wfdbx\_set\_input\_dir}{Set the input directory}
\begin{lstlisting}[language=C++,style=signature]
std::optional<std::string> wfdbx_set_input_dir(const char *dir);
\end{lstlisting}
Sets the input directory: every file named by a relative name (\code{100s},
\code{mitdb/100}, \code{../other/100}) is read from it; an absolute name, a
\code{\textasciitilde} name or a URL is used as given. It may be a local
directory, a \code{\textasciitilde} path, or the URL of a directory, such as
\code{https://physionet.org/files/mitdb/1.0.0/} (which also needs network
access, \code{wfdbx\_set\_network}). The name is used as it is (it may contain
spaces). The setting applies to the library functions, the high-level
functions (whose \code{input\_dir} option sets it for the duration of a call)
and the tools, also for the files that they read themselves
(\code{wfdbx\_input\_path}); \code{wfdbquit} does not change it; legacy mode
does not use it (it searches the WFDB path).
\begin{params}
\item[dir] The input directory; \code{NULL} or \code{""} restores the
  default: the value of \code{WFDBX\_INPUT\_DIR} if it is set, otherwise the
  current directory.
\end{params}
\returns{The previous input directory (that of \code{WFDBX\_INPUT\_DIR} if
it was in effect); \code{NULL} (\code{std::nullopt}, \code{None},
\code{nothing}, ...) if none was set.}

\apientry{wfdbx\_get\_input\_dir}{The input directory}
\begin{lstlisting}[language=C++,style=signature]
std::optional<std::string> wfdbx_get_input_dir();
\end{lstlisting}
Returns the input directory set by \code{wfdbx\_set\_input\_dir},
\code{setwfdb} or \code{WFDBX\_INPUT\_DIR}.
\returns{The input directory; \code{NULL} if none is set (files are then read
from the current directory; \code{getwfdb} returns \code{"."}).}

\apientry{wfdbx\_set\_output\_dir}{Set the output directory}
\begin{lstlisting}[language=C++,style=signature]
std::optional<std::string> wfdbx_set_output_dir(const char *dir);
\end{lstlisting}
Sets the output directory: every file written under a relative name (a
signal file, a header, an annotation, info or calibration file) is created
in it; missing directories are created. An absolute name or a
\code{\textasciitilde} name is used as given. The names recorded in headers
do not change. The setting applies to the library functions, the high-level
writers (whose \code{dir} option sets it for the duration of a call) and the
tools, also for the files that they create themselves
(\code{wfdbx\_output\_path}); \code{wfdbquit} does not change it. In legacy
mode it is a plain prefix of every output name, absolute names included.
\begin{params}
\item[dir] A local directory or a \code{\textasciitilde} path; \code{NULL} or
  \code{""} restores the default: the value of \code{WFDBX\_OUTPUT\_DIR} if it
  is set, otherwise the current directory. A URL is refused.
\end{params}
\returns{The previous output directory; \code{NULL} if none was set. For a
URL, the current output directory, unchanged.}
\raises{\code{wfdbx\_set\_output\_dir: the output directory must be a local
directory}.}

\apientry{wfdbx\_get\_output\_dir}{The output directory}
\begin{lstlisting}[language=C++,style=signature]
std::optional<std::string> wfdbx_get_output_dir();
\end{lstlisting}
Returns the output directory set by \code{wfdbx\_set\_output\_dir} or
\code{WFDBX\_OUTPUT\_DIR}.
\returns{The output directory; \code{NULL} if none is set (files are then
written in the current directory).}

\apientry{wfdbx\_set\_network}{Allow or forbid network access}
\begin{lstlisting}[language=C++,style=signature]
int wfdbx_set_network(int on);
\end{lstlisting}
Allows or forbids reading files from \code{http://} and \code{https://}
URLs; it is forbidden by default. A remote file is read in full the first
time it is opened and kept for the rest of the process; a download gives up
if it cannot connect, or receives nothing, for 30 seconds. With network
access forbidden, the first attempt to open a URL prints \code{WFDB-x:
network access is disabled; set WFDBX\_NETWORK=1 or call wfdbx\_set\_network(1)
to open URL}, and the file cannot be opened. Legacy mode reads no remote
files.
\begin{params}
\item[on] 1 (or any positive value; \code{True} in Python) to allow, 0 to
  forbid; a negative value (or \code{None} in Python) restores the default:
  allowed if \code{WFDBX\_NETWORK} is a nonzero number.
\end{params}
\returns{The previous setting: 1 (allowed) or 0.}

\apientry{wfdbx\_get\_network}{Whether network access is allowed}
\begin{lstlisting}[language=C++,style=signature]
int wfdbx_get_network();
\end{lstlisting}
\returns{1 if network access is allowed, 0 otherwise.}

\apientry{wfdbx\_set\_legacy}{Turn legacy mode on or off}
\begin{lstlisting}[language=C++,style=signature]
int wfdbx_set_legacy(int on);
\end{lstlisting}
Turns legacy mode on or off for the whole library. In legacy mode, every
deliberate change of WFDB-x is undone (\cref{chap:behavior}): the built-in
sampling frequency, gain and ADC resolution, the WFDB search path (without
network access), sorting by the \code{sortann} program, output files that may
be input files, and C's checks and names in headers. It is off by default,
and is meant only to reproduce results of the original package.
\begin{params}
\item[on] Nonzero to turn legacy mode on, 0 to turn it off.
\end{params}
\returns{The previous setting: 1 (on) or 0.}

\apientry{wfdbx\_get\_legacy}{Whether legacy mode is on}
\begin{lstlisting}[language=C++,style=signature]
int wfdbx_get_legacy();
\end{lstlisting}
\returns{1 if legacy mode is on, 0 otherwise.}

\apientry{wfdbx\_legacy\_args}{Recognize the --legacy option of a tool}
\begin{lstlisting}[language=C++,style=signature]
int wfdbx_legacy_args(int &argc, char **&argv);
\end{lstlisting}
The function that every tool calls first: if the first argument after the
program name (\code{argv[1]}) is exactly \code{--legacy}, it is removed (the
program name stays first) and legacy mode is turned on. \code{--legacy}
elsewhere is not recognized. In C++, Rust and Kotlin, the call also arranges
the exit status of the program as for the tools: at exit, outside legacy
mode, a failure to write the standard output is reported (\code{PROGRAM:
write error on standard output}), and an exit status of 0 becomes 1 after a
write error (see \code{wfdbx\_write\_errors}).
\begin{params}
\item[argc] The number of arguments, decremented if \code{--legacy} is
  removed.
\item[argv] The arguments, program name first.
\end{params}
\returns{1 if \code{--legacy} was found (and removed), 0 otherwise.}

The argument list is that of each language: in Python, a list
(\code{sys.argv}), changed in place; in Rust and Kotlin, a mutable list of
byte strings; in R, a character vector, and the result is \code{list(argv,
legacy)}; in MATLAB/Octave, a cell array, with \code{[r, argv] =
wfdbx.wfdbx\_legacy\_args(argv)}. In Julia, \code{wfdbx\_legacy\_args(args)}
takes the arguments \emph{without} the program name (Julia's \code{ARGS}), and
looks at \code{args[1]}.

\apientry{wfdbx\_time\_needs\_freq}{Whether strtim needs the sampling frequency}
\begin{lstlisting}[language=C++,style=signature]
int wfdbx_time_needs_freq(const char *time_string);
\end{lstlisting}
Tells whether \code{strtim} needs a sampling frequency to convert
\code{time\_string}, so that a program can report a missing frequency in its
own terms.
\begin{params}
\item[time\_string] The time string.
\end{params}
\returns{0 if the string (after leading spaces, tabs, carriage returns and
newlines) begins with \code{e}, \code{i}, \code{o} or \code{s}, or with
\code{c} while no counter frequency is defined; 1 otherwise.}

\apientry{wfdbx\_write\_errors}{Number of write errors}
\begin{lstlisting}[language=C++,style=signature]
int wfdbx_write_errors();
\end{lstlisting}
Returns the number of write errors since the program started: the output
files (signal files, headers, annotation, info and calibration files, and
the standard output opened as \code{-}) that the library could not write,
flush or close, found when they are closed (by \code{wfdbquit},
\code{osigfopen(NULL, 0)}, \code{oannclose}, ...), and those counted by
\code{wfdbx\_note\_write\_error}. The count belongs to the process, not to a
library state object (Rust, Julia, Kotlin). The tools exit with status 1
when it is not 0, and the high-level writers raise an error when they meet a
write error (\cref{chap:behavior}).
\returns{The number of write errors.}

\apientry{wfdbx\_note\_write\_error}{Count a write error}
\begin{lstlisting}[language=C++,style=signature]
void wfdbx_note_write_error();
\end{lstlisting}
Adds one to the count of \code{wfdbx\_write\_errors}, for a program that
writes a file itself and could not write it. (Rust and Kotlin also have
\code{wfdbx\_fclose\_output(file, pname, name)}, which closes such a file,
and reports \code{PNAME: write error on NAME} and counts it on failure.)

\apientry{wfdbx\_input\_path}{Where a program reads a file itself}
\begin{lstlisting}[language=C++,style=signature]
std::optional<std::string> wfdbx_input_path(const char *name);
\end{lstlisting}
Returns the path under which a program opens a file that it reads itself,
rather than through the library, so that the file is found where
\code{wfdb\_open} finds a file of that name (with an empty file type). In
normal mode, a relative name is found under the input directory, if one is
set (\code{wfdbx\_set\_input\_dir}), a \code{\textasciitilde} name is
expanded, and an absolute name is used as given; in legacy mode, the name is
used as given, as the original programs did. The tools use it for the files
named on their command line that they read themselves, such as the \opt{-i}
file of \cmd{wrsamp} (\cref{sec:tool-files}). For example, with the input
directory \file{data}, \code{wfdbx\_input\_path("samples.txt")} returns
\code{data/samples.txt}.
\begin{params}
\item[name] A file name. \code{""} and \code{"-"} (the standard input, if the
  program treats it so) are returned as given.
\end{params}
\returns{The path to open; \code{NULL} (\code{std::nullopt}, \code{None},
\code{nothing}, ...) if \code{name} is \code{NULL} or a URL, or becomes a URL
when the input directory is prefixed to it: such a file cannot be opened in
this way.}

\apientry{wfdbx\_output\_path}{Where a program creates a file itself}
\begin{lstlisting}[language=C++,style=signature]
std::optional<std::string> wfdbx_output_path(const char *name);
\end{lstlisting}
Returns the path under which a program creates a file that it writes
itself, rather than through the library, so that the file is created where
\code{wfdb\_open} creates a file of that name (with an empty file type). In
normal mode, a relative name is created under the output directory, if one
is set (\code{wfdbx\_set\_output\_dir}), a \code{\textasciitilde} name is
expanded, and an absolute name is used as given; in legacy mode, the output
directory, if one is set, is a plain prefix of any name, absolute names
included. When the output directory has been prefixed to the name, the
directories missing in the resulting path are created, as for the output
files of the library. The tools use it for the files that they create
themselves, such as the WAV file of \cmd{mit2wav} (\cref{sec:tool-files}).
For example, with the output directory \file{out},
\code{wfdbx\_output\_path("notes.txt")} returns \code{out/notes.txt}.
\begin{params}
\item[name] A file name. \code{""} and \code{"-"} (the standard output, if the
  program treats it so) are returned as given.
\end{params}
\returns{The path to create; \code{NULL} (\code{std::nullopt}, \code{None},
\code{nothing}, ...) if \code{name} is \code{NULL} or a URL (which cannot be
written).}

\apientry{getframes}{Read many sample vectors at once}
\begin{lstlisting}[language=C++,style=signature]
long long getvecs(WFDB_Sample *out, long long n, int *status = nullptr);
// The same function in each language:
//   C++, Rust, Julia, Kotlin:  getvecs(out, n), and in Rust, Julia and Kotlin
//                              getvecs_status(out, n), which also returns status
//   Python:                    frames, status = getframes(n)
//   R:                         r <- getframes(n, minstat = 0)   # r$v, r$stat
//   MATLAB/Octave:             [v, status, last] = wfdbx.getframes(n, minstat)
\end{lstlisting}
Reads up to \code{n} sample vectors, with exactly the result of \code{n}
calls of \code{getvec} that stop at the first negative status, but much
faster: whole blocks of the signal files are decoded at once, for every
format, for multi-frequency records in both modes of \code{getvec}, and for
multi-segment records; the frames that need special handling (the one at
which a checksum is verified, a change of segment, skewed or resampled
signals) go through \code{getvec}. The vectors follow one another in
\code{out}, one value per signal in each (in Python, R and MATLAB/Octave, a
matrix with one row per vector and one column per signal).
\begin{params}
\item[out] (C++, Rust, Julia, Kotlin) An array with room for \code{n} times
  the number of values of a vector, filled by the function.
\item[n] The maximum number of vectors to read (\code{Inf} in R and
  MATLAB/Octave reads to the end).
\item[status] (C++, optional) Receives the value returned by the last
  \code{getvec}: the negative status of the first vector that could not be
  read (-1 at the end of the data, -3 for an unexpected end of file, ...), or
  the number of signals if all \code{n} were read.
\item[minstat] (R, MATLAB/Octave, optional) Stop at the first vector for
  which \code{getvec} would return less than \code{minstat} (as a C loop
  \code{while (getvec(v) >= minstat)}); 0 by default.
\end{params}
\returns{C++, Rust, Julia, Kotlin: the number of vectors read (\code{>= 0}),
or the negative status of \code{getvec} if none could be read;
\code{getvecs\_status} returns the number and the status. Python: a NumPy
array (\code{int64}, one row per vector) and the status. R:
\code{list(v, stat)}, a matrix and the status. MATLAB/Octave: the matrix, the
status, and the vector filled by the \code{getvec} call that stopped the
reading (the last frame of the record for a checksum error, -4).}

% ==========================================================================
\section{Constants and structures}
\label{sec:library-structures}

The constants of \file{<wfdb/wfdb.h>} have their C names and values in every
language (\cref{tab:library-constants}); in MATLAB/Octave they are functions
without arguments (\code{wfdbx.WFDB\_READ()}), and the calibration constants
and \code{WFDB\_FMT\_LIST} are not defined there. The annotation codes
(\code{NORMAL}, \code{PVC}, ..., \code{ACMAX}) of \file{<wfdb/ecgcodes.h>}
are listed in \cref{chap:formats}.

\begin{xltabular}{\linewidth}{@{}P{4.4cm}P{2.4cm}L@{}}
\caption{The constants of the library.}
\label{tab:library-constants}\\
\toprule
{\tabhead Constant} & {\tabhead Value} & {\tabhead Meaning}\\
\midrule
\endfirsthead
\toprule
{\tabhead Constant} & {\tabhead Value} & {\tabhead Meaning}\\
\midrule
\endhead
\bottomrule
\endfoot
\code{WFDB\_READ} & 0 & \code{stat} of an MIT-format input annotation file\\
\code{WFDB\_WRITE} & 1 & \code{stat} of an MIT-format output annotation file\\
\code{WFDB\_AHA\_READ} & 2 & \code{stat} of an AHA-format input annotation
file\\
\code{WFDB\_AHA\_WRITE} & 3 & \code{stat} of an AHA-format output annotation
file\\
\code{WFDB\_APPEND} & 4 & mode of an output info file (opened for appending)\\
\code{WFDB\_INVALID\_SAMPLE} & -32768 & value of an invalid or missing
sample, in \code{getvec}, \code{getframe} and \code{sample}, unless
\code{WFDB\_GVPAD} is set; written by \code{putvec} as the invalid value of
the format\\
\code{WFDB\_LOWRES} & 0 & \code{getvec} mode: one sample per signal per
frame\\
\code{WFDB\_HIGHRES} & 1 & \code{getvec} mode: every sample of the
oversampled signals\\
\code{WFDB\_GVPAD} & 2 & \code{getvec} flag: invalid samples replaced by the
previous valid sample\\
\code{WFDB\_DEFFREQ} & 250.0 & C's default sampling frequency (Hz); not used
by WFDB-x outside legacy mode\\
\code{WFDB\_DEFGAIN} & 200.0 & C's default gain (adu per physical unit); not
used outside legacy mode\\
\code{WFDB\_DEFRES} & 12 & C's default ADC resolution (bits); not used
outside legacy mode\\
\code{WFDB\_MAXRNL} & 50 & maximum length of a record name (of a segment
name); not exported in R\\
\code{WFDB\_MAXUSL} & 50 & maximum length of a \code{units} string\\
\code{WFDB\_MAXDSL} & 100 & maximum length of a \code{desc} string\\
\code{WFDB\_FMT\_LIST} & 0, 8, 16, 61, 80, 160, 212, 310, 311, 24, 32, 508,
516, 524 & the signal storage formats (\cref{chap:formats})\\
\code{WFDB\_AC\_COUPLED} & 0 & \code{caltype}: AC-coupled signal\\
\code{WFDB\_DC\_COUPLED} & 1 & \code{caltype}: DC-coupled signal\\
\code{WFDB\_CAL\_SQUARE} & 2 & \code{caltype}: square calibration pulse\\
\code{WFDB\_CAL\_SINE} & 4 & \code{caltype}: sine calibration pulse\\
\code{WFDB\_CAL\_SAWTOOTH} & 6 & \code{caltype}: sawtooth calibration pulse\\
\code{WFDB\_CAL\_UNDEF} & 8 & \code{caltype}: undefined pulse shape\\
\code{APUNDEF}, \code{APSTD}, \code{APHIGH}, \code{APLOW}, \code{APATT},
\code{APAHIGH}, \code{APALOW} & 0 to 6 & display positions (\code{annpos}):
undefined, standard, above, below, attached to the signal \code{chan}, above
and below that\\
\end{xltabular}

The simple types of the C library are integers and floating-point numbers:
\code{WFDB\_Sample} (\code{int}, 32 bits) for samples in ADC units,
\code{WFDB\_Time} (64 bits on the reference platform) for times in sample
intervals, \code{WFDB\_Date} (\code{long}) for Julian dates,
\code{WFDB\_Frequency} and \code{WFDB\_Gain} (\code{double}), and
\code{WFDB\_Group}, \code{WFDB\_Signal} and \code{WFDB\_Annotator}
(\code{unsigned int}). \Cref{sec:library-calling} gives the corresponding
types of each language (doubles in R and MATLAB/Octave).

The structures have the C field names in every language
(\cref{tab:library-siginfo,tab:library-anninfo,tab:library-annotation,tab:library-calinfo,tab:library-seginfo});
their names are \code{WFDB\_Siginfo}, ..., in C++ (also \code{Siginfo}, ...),
\code{Siginfo}, \code{Anninfo}, \code{Annotation}, \code{Calinfo} and
\code{Seginfo} in Python, Rust, Julia and Kotlin, and lists or structures made
by \code{siginfo()}, \code{anninfo()}, \code{annotation()}, \code{calinfo()}
and \code{seginfo()} in R and MATLAB/Octave (\code{wfdbx.siginfo(n)}, ...). A
string field that C may leave \code{NULL} is absent in the language's terms
(\code{None}, \code{std::nullopt}, \code{nothing}, \code{NULL}, \code{[]},
\code{null}).

\begin{xltabular}{\linewidth}{@{}P{3.4cm}L@{}}
\caption{\code{WFDB\_Siginfo}: the specification of a signal.}
\label{tab:library-siginfo}\\
\toprule
{\tabhead Field} & {\tabhead Meaning}\\
\midrule
\endfirsthead
\toprule
{\tabhead Field} & {\tabhead Meaning}\\
\midrule
\endhead
\bottomrule
\endfoot
\code{char *fname} & the name of the signal file, relative to the header's
directory (WFDB-x); \code{-} is the standard input or output\\
\code{char *desc} & the description of the signal (\code{MLII}), at most
\code{WFDB\_MAXDSL} characters, without line breaks; a header that gives none
gives \code{record REC, signal N}\\
\code{char *units} & the physical units, at most \code{WFDB\_MAXUSL}
characters, without blanks; \code{NULL} if not specified (WFDB-x: never
taken to be millivolts)\\
\code{WFDB\_Gain gain} & ADC units per physical unit; 0 if uncalibrated\\
\code{WFDB\_Sample initval} & the value of sample 0 (used by format 8)\\
\code{WFDB\_Group group} & the signal group: the signals of a group are stored
in one file, and have consecutive numbers\\
\code{int fmt} & the storage format (\code{WFDB\_FMT\_LIST},
\cref{chap:formats})\\
\code{int spf} & samples per frame (more than 1 for an oversampled signal)\\
\code{int bsize} & the block size in bytes, for character special files (0
for ordinary files)\\
\code{int adcres} & the ADC resolution in bits; 0 if not specified (WFDB-x:
not inferred from the format)\\
\code{int adczero} & the ADC output for an input at the center of its range\\
\code{int baseline} & the ADC value of 0 physical units\\
\code{long nsamp} & the number of samples per signal (of frames, for an
oversampled signal); 0 if not specified\\
\code{int cksum} & the 16-bit checksum of the samples, verified by
\code{getvec} at the end of the record\\
\end{xltabular}

The constructors give C's zero values, except that \code{spf} is 1 in Python,
R and MATLAB/Octave (0 in C++, Rust, Julia and Kotlin; \code{osigfopen} takes
an \code{spf} below 1 as 1).

\begin{xltabular}{\linewidth}{@{}P{3.4cm}L@{}}
\caption{\code{WFDB\_Anninfo}: an annotator to open.}
\label{tab:library-anninfo}\\
\toprule
{\tabhead Field} & {\tabhead Meaning}\\
\midrule
\endfirsthead
\toprule
{\tabhead Field} & {\tabhead Meaning}\\
\midrule
\endhead
\bottomrule
\endfoot
\code{char *name} & the annotator name, the suffix of the annotation file
(\code{atr} is the reference annotation file of a database record);
\code{-} is the standard input or output\\
\code{int stat} & \code{WFDB\_READ}, \code{WFDB\_WRITE},
\code{WFDB\_AHA\_READ} or \code{WFDB\_AHA\_WRITE}\\
\end{xltabular}

\begin{xltabular}{\linewidth}{@{}P{3.4cm}L@{}}
\caption{\code{WFDB\_Annotation}: an annotation.}
\label{tab:library-annotation}\\
\toprule
{\tabhead Field} & {\tabhead Meaning}\\
\midrule
\endfirsthead
\toprule
{\tabhead Field} & {\tabhead Meaning}\\
\midrule
\endhead
\bottomrule
\endfoot
\code{WFDB\_Time time} & the time, in sample intervals from the beginning of
the record (or in the units of \code{setiafreq} and \code{setafreq})\\
\code{char anntyp} & the annotation code, 1 to \code{ACMAX}
(\cref{chap:formats})\\
\code{signed char subtyp} & the subtype (-128 to 127); for noise
annotations, the signals affected\\
\code{unsigned char chan} & the channel (0 to 255), usually the signal to
which the annotation refers\\
\code{signed char num} & the annotator number (-128 to 127), free for the
program's use\\
\code{unsigned char *aux} & \code{NULL}, or the auxiliary data: its first
byte is the number of bytes that follow (at most 255), such as the rhythm of
a rhythm change \code{(N}, or the text of a comment; it may contain any
bytes\\
\end{xltabular}

\begin{xltabular}{\linewidth}{@{}P{3.4cm}L@{}}
\caption{\code{WFDB\_Calinfo}: an entry of the calibration list.}
\label{tab:library-calinfo}\\
\toprule
{\tabhead Field} & {\tabhead Meaning}\\
\midrule
\endfirsthead
\toprule
{\tabhead Field} & {\tabhead Meaning}\\
\midrule
\endhead
\bottomrule
\endfoot
\code{double low} & the low level of the calibration pulse, in physical
units (0 for an AC-coupled signal)\\
\code{double high} & the high level of the pulse (the pulse amplitude for an
AC-coupled signal)\\
\code{double scale} & the customary plotting scale, in physical units per
centimeter\\
\code{char *sigtype} & the signal type, matched against the beginning of a
signal description (empty: any signal)\\
\code{char *units} & the physical units\\
\code{int caltype} & \code{WFDB\_AC\_COUPLED} or \code{WFDB\_DC\_COUPLED},
plus \code{WFDB\_CAL\_SQUARE}, \code{WFDB\_CAL\_SINE} or
\code{WFDB\_CAL\_SAWTOOTH}\\
\end{xltabular}

\begin{xltabular}{\linewidth}{@{}P{3.4cm}L@{}}
\caption{\code{WFDB\_Seginfo}: a segment of a multi-segment record.}
\label{tab:library-seginfo}\\
\toprule
{\tabhead Field} & {\tabhead Meaning}\\
\midrule
\endfirsthead
\toprule
{\tabhead Field} & {\tabhead Meaning}\\
\midrule
\endhead
\bottomrule
\endfoot
\code{char recname[WFDB\_MAXRNL+1]} & the name of the segment (a record), or
\code{\textasciitilde} for a gap without signals\\
\code{WFDB\_Time nsamp} & the length of the segment, in samples\\
\code{WFDB\_Time samp0} & the number of samples before the segment in the
multi-segment record (the sample number of its first sample)\\
\end{xltabular}

\setlength{\emergencystretch}{0pt}

% The command-line tools: how to run them, their conventions, and one reference entry per
% tool. The option lists are those of the Python implementation (python/src/wfdbx/apps),
% which check_manual.py compares with the usage text of each tool.
\chapter{The command-line tools}
\label{chap:tools}

\wfdbx{} includes the 33 command-line tools of the WFDB Software Package \wfdbversion{}
that read, write and convert records and annotation files. Each of the seven
implementations has all of them, with the options, output, messages and exit status of the
C tools, so that a script written for the WFDB tools runs unchanged with \wfdbx{}. The
exceptions are the deliberate divergences of \cref{chap:behavior}: no built-in signal
parameters (hence the options \opt{-F}, \opt{-G}, \opt{-R} and \opt{-U} added to some
tools), files found only where they are named, annotation files sorted without loss, an
output signal file never written over an input file, write errors reported in the exit
status, and the record line of a header found by \code{getinfo} as the header parser finds
it (the \code{Notes} of \cmd{wfdbdesc}). The original behavior remains available with \opt{-{}-legacy}
(\cref{sec:tools-legacy}).

Every implementation is checked against the C tools by the same conformance suite (1155
cases): the standard output, the standard error, the exit status and the files written are
compared byte for byte. The examples of this chapter were produced by the Python
implementation, run in a copy of the directory \file{tests/data} of the repository, which
holds the sample record \code{100s} (the first minute of record 100 of the MIT-BIH
Arrhythmia Database) and the other test inputs; every other implementation gives the same
output.

\section{Running the tools}
\label{sec:tools-running}

The tools are separate programs in C++ and Rust, and subcommands of a launcher in the other
languages, as in \cref{tab:tools-launchers}. The arguments that follow the tool's name are
those of the C tool: \cmd{wfdbx rdsamp -r 100s -t 10} (Python, Julia, R, MATLAB/Octave and
Kotlin, with the alias \cmd{wfdbx}) and \cmd{rdsamp -r 100s -t 10} (C++ and Rust) do the
same thing. The tool sees its own name as its program name, so its messages begin with
\code{rdsamp:}, as those of the C tool do.

\begin{xltabular}{\linewidth}{@{}P{2.5cm}L@{}}
\caption{How each implementation runs a tool (here \code{rdsamp}).}
\label{tab:tools-launchers}\\
\toprule
{\leavevmode\tabhead Language} & {\leavevmode\tabhead Command} \\
\midrule
\endfirsthead
\toprule
{\leavevmode\tabhead Language} & {\leavevmode\tabhead Command} \\
\midrule
\endhead
\bottomrule
\endfoot
Python & \cmd{wfdbx rdsamp ...}: the command \cmd{wfdbx} is installed with the package
  (\cmd{pip install}); without it, \cmd{python -m wfdbx.cli rdsamp ...} or
  \cmd{python -m wfdbx.apps.rdsamp ...}. In a program, \code{wfdbx.apps.rdsamp.main(argv)}
  runs the tool (\code{argv} is the whole argument list, the program name first) and
  returns its exit status. \\
C++ & \cmd{rdsamp ...}: one program per tool, built in \file{cpp/build/bin/} (or the
  \file{bin} directory of the installation). \\
Rust & \cmd{rdsamp ...}: one program per tool, installed by
  \cmd{cargo install -{}-path rust}. \\
Julia & \cmd{julia -{}-startup-file=no -{}-project=JULIADIR JULIADIR/bin/wfdbx.jl rdsamp ...},
  where \code{JULIADIR} is the directory \file{julia} of the repository; a shell script
  named \cmd{wfdbx} that runs this command makes it \cmd{wfdbx rdsamp ...}. In a session,
  \code{WFDBX.run\_tool("rdsamp", ["-r", "100s"]; dir=...)} returns the exit status and the
  bytes written to the standard output and the standard error, and \code{WFDBX.tools()}
  the names of the tools. \\
R & \cmd{Rscript r/inst/bin/wfdbx.R rdsamp ...} (in an installed package, the script is
  \code{system.file("bin", "wfdbx.R", package = "wfdbx")}); with the alias
  \cmd{alias wfdbx='Rscript .../wfdbx.R'}, \cmd{wfdbx rdsamp ...}. In a session,
  \code{wfdbx\_tool("rdsamp", c("-r", "100s"))} runs the tool and returns its exit status,
  and \code{wfdbx\_tools()} lists the tools. \\
MATLAB/Octave & Octave: \cmd{octave-cli -{}-no-gui -{}-norc -q matlab/bin/wfdbx\_cli.m rdsamp ...}
  (with an alias \cmd{wfdbx}, \cmd{wfdbx rdsamp ...}); MATLAB:
  \cmd{matlab/bin/wfdbx\_matlab.sh rdsamp ...}, which runs \cmd{matlab -batch} (or the
  program named by the environment variable \code{WFDBX\_MATLAB}). In a session,
  \code{status = wfdbx.apps.rdsamp(\{'-r', '100s'\})}; \cmd{help wfdbx.apps.rdsamp} shows
  the usage. \\
Kotlin & \cmd{kotlin/bin/wfdbx rdsamp ...}, a POSIX shell script that runs
  \cmd{java -cp build/libs/wfdbx-all.jar wfdbx.cli.MainKt rdsamp ...} (the jar is built by
  \cmd{./gradlew fatJar}); a link to \file{bin/wfdbx} named after a tool runs that tool. \\
\end{xltabular}

The launchers also take a few arguments of their own, in place of the tool's name:
\opt{-{}-list} prints the names of the tools, one per line (Python, R, MATLAB/Octave and
Kotlin); \opt{-{}-version} prints the version of WFDB-x, as \code{WFDB-x} followed by the
version number (Python, Julia, R, MATLAB/Octave and Kotlin); \opt{-h} or \opt{-{}-help}
prints the launcher's usage (Python, R, MATLAB/Octave and Kotlin), and the Python launcher
also has \opt{-{}-examples} \textit{DIR}, which copies the example notebooks to \textit{DIR}
(by default \file{wfdbx-examples}). Without arguments, or with an unknown tool name, a
launcher prints a message on the standard error and exits with status 1. The programs of
C++ and Rust, like the C tools, have no \opt{-{}-version} option.

\section{Conventions}
\label{sec:tools-conventions}

\subsection{Records and files}
\label{sec:tools-files}

A record is named as in \cref{chap:records}, without the suffix of its header
(\code{100s}, not \code{100s.hea}), and an annotation file by its record and its
annotator (\opt{-r 100s -a atr} is the file \file{100s.atr}). The tools find every file
that they read through the library, as \cref{chap:files} describes: a relative name is
taken relative to the \emph{input directory}, set by the environment variable
\code{WFDBX\_INPUT\_DIR} (by default the current directory), an absolute name, a \code{\textasciitilde}
name or a URL is used as given, and nothing is searched for. A URL (only \code{http://} and
\code{https://} URLs are read) is opened only when network access is turned on, by
\code{WFDBX\_NETWORK=1}. Files written through the library
(headers, signal files, annotation files) are created in the \emph{output directory},
\code{WFDBX\_OUTPUT\_DIR} (by default the current directory). For example, to read record
100 of the MIT-BIH Arrhythmia Database from PhysioNet:
\begin{lstlisting}[style=shell]
export WFDBX_INPUT_DIR=https://physionet.org/files/mitdb/1.0.0/
export WFDBX_NETWORK=1
wfdbx rdsamp -r 100 -t 0:05
\end{lstlisting}

A few files are opened by the tools themselves, not through the library: the input files
named by \opt{-i} in \cmd{wrsamp}, \cmd{edf2mit}, \cmd{wav2mit}, \cmd{a2m} and \cmd{ad2m},
the input files of \cmd{ahaecg2mit}, the script of \cmd{xform -S}, and the signal file
whose size \cmd{wfdbdesc} measures; and the output files that a tool writes itself: the WAV
file of \cmd{mit2wav}, the EDF file of \cmd{mit2edf}, the master header of \cmd{snip -m},
and the \file{.des}, \file{.ecg}, \file{.key} and \file{.txt} files of \cmd{parsescp}.
They follow the same rules as the files of the library (\cref{sec:tool-files}): a relative
name is read from the input directory, or created in the output directory (with its
missing directories), a \code{\textasciitilde} name is expanded, an absolute name is used
as given, and \code{-} keeps its meaning. A URL is not accepted for these files, also when
the input directory is a URL: the tool reports that it cannot read, or create, the file. In legacy mode, the input files are
opened as named, relative to the current directory, and the output directory, if one is
set, is a plain prefix of the names of the output files.

\subsection{Time arguments}
\label{sec:tools-times}

Options that take a time (\opt{-f}, \opt{-t}, \opt{-l} and others) accept the time strings
of the library function \code{strtim}:
\begin{itemize}
\item an elapsed time from the beginning of the record, \code{hh:mm:ss}, \code{mm:ss} or
  \code{ss}, with an optional fraction of a second (\code{1:30.5}, \code{10});
\item a sample number, \code{s} followed by the number (\code{s21600});
\item a time of day, in brackets: \code{[hh:mm:ss]}, \code{[hh:mm:ss dd/mm/yyyy]}, or
  \code{[hh:mm:ss d]} with a number of days after the start (it needs the base time of the
  record);
\item \code{e}, the end of the record; \code{c} followed by a counter value; \code{f}
  followed by a frame number.
\end{itemize}
Elapsed times and times of day need the sampling frequency of the record. When a record
does not give one, \wfdbx{} reports an error instead of assuming 250\,Hz
(\cref{chap:behavior}); sample numbers and \code{e} need no frequency. The annotation tools
\cmd{rdann}, \cmd{ann2rr}, \cmd{sortann}, \cmd{a2m} and \cmd{m2a} take \opt{-F}
\textit{FREQ} for that case, and \cmd{rr2ann} already had it.

\subsection{Help, options and arguments}
\label{sec:tools-help}

The options are those of the C tools: one letter after a hyphen, given separately
(\opt{-v -p}, not \opt{-vp}), their arguments as the next command-line arguments. As in
C, only the first letter after the hyphen is read in most tools, so that \opt{-pd} is
\opt{-p} with the time format \code{d}, and an argument that is not an option where an
option is expected is an error (\code{unrecognized argument}). Options that take a list
(\opt{-s} \textit{SIGNAL} \code{[}\textit{SIGNAL} \code{...]}, \opt{-p} \textit{TYPE} \code{...}) take
every following argument that does not begin with a hyphen.

\opt{-h} prints the tool's usage summary on the standard error. A tool that needs an
argument that was not given (often the record name) prints the same summary. The exit
status after \opt{-h} is that of the C tool, 0 or 1 depending on the tool (each entry says
which); as in C, \cmd{wfdbdesc}, \cmd{sampfreq} and \cmd{revise} take \opt{-h} alone as a
record name. The help texts of \cmd{ann2rr} and \cmd{wrsamp} pause after 23
lines, print \code{-{}-More-{}-} and read a line from the standard input, as the C tools do. The usage
line of each tool is given at the start of its entry below.

\subsection{Legacy mode}
\label{sec:tools-legacy}

\opt{-{}-legacy}, as the \emph{first} argument of a tool (\cmd{rdsamp -{}-legacy -r 100s},
\cmd{wfdbx rdann -{}-legacy -r 100s -a atr}), is removed before the options are read, and the
tool then behaves exactly as the original one: the options added by \wfdbx{} (such as
\opt{-F} of \cmd{rdann} or \opt{-R} of \cmd{wrsamp}) are unrecognized, the help texts and
messages are those of the original, the built-in values of the original (250\,Hz,
200\,adu/mV, ...) are used, files are found through the WFDB path (the environment variable
\code{WFDB}; the files that a tool reads itself are opened as named), and the exit status
is the original one. \opt{-{}-legacy} anywhere else is not
recognized: it is treated as any other argument. Legacy mode is meant for reproducing results of the original package;
\cref{chap:behavior} explains why it is not the default.

\subsection{Standard input and output}
\label{sec:tools-stdio}

Text output goes to the standard output, and messages, prompts and help texts to the
standard error. Some tools read the standard input (\cmd{wrsamp} without \opt{-i},
\cmd{edf2mit -i -}, \cmd{wrann}, \cmd{rr2ann}, \cmd{parsescp}, \cmd{readid}) and some write
binary data to the standard output (\cmd{wfdbcat}, \cmd{makeid}, \cmd{parsescp -a}, and
\cmd{wrsamp} without \opt{-o}).

\subsection{Exit status}
\label{sec:tools-status}

The exit status is that of the C tool: 0 on success, 1 for an error in the arguments
(including a required option that is missing), 2 when the record or a file cannot be read,
and, in some tools, other values that their entries give. In addition, \wfdbx{} follows the
convention of the standard Unix tools: when an output file or the standard output cannot be
written (a full disk, for example), a tool that would have exited with status 0 exits with
status 1, and reports a standard output that could not be written as
\code{PROGRAM: write error on standard output}, and a file that \cmd{snip}, \cmd{mit2edf}
or \cmd{parsescp} writes itself as \code{PROGRAM: write error on FILE}
(\cref{chap:behavior}). The original tools exited with status 0 in both cases; legacy
mode does too.

\section{Summary of the tools}
\label{sec:tools-summary}

\begin{xltabular}{\linewidth}{@{}P{2.3cm}L@{}}
\caption{The 33 tools of \wfdbx{}.}
\label{tab:tools}\\
\toprule
{\leavevmode\tabhead Tool} & {\leavevmode\tabhead What it does} \\
\midrule
\endfirsthead
\toprule
{\leavevmode\tabhead Tool} & {\leavevmode\tabhead What it does} \\
\midrule
\endhead
\bottomrule
\endfoot
\multicolumn{2}{@{}l}{\textit{Signal tools} (\cref{sec:tools-signal})} \\
\code{rdsamp} & print the samples of a record (text, CSV or XML; ADC or physical units) \\
\code{wrsamp} & write a record from a text file of samples \\
\code{wfdbdesc} & describe the signals of a record (the contents of its header) \\
\code{wfdbtime} & convert times to sample numbers, elapsed times and times of day \\
\code{time2sec} & convert a time to seconds \\
\code{sampfreq} & print the sampling frequency of a record \\
\code{signame} & print the names of the signals of a record \\
\code{signum} & print the numbers of the signals with given names \\
\code{sigamp} & measure signal amplitudes (peak-to-peak near beats, or RMS) \\
\code{snip} & copy a segment of a record to a new record \\
\code{xform} & convert a record: storage format, sampling frequency, gain, signals \\
\code{wfdbcat} & copy WFDB files to the standard output \\
\code{wfdbwhich} & print the path of a WFDB file \\
\midrule
\multicolumn{2}{@{}l}{\textit{Annotation tools} (\cref{sec:tools-annotation})} \\
\code{rdann} & print the annotations of an annotation file \\
\code{wrann} & write an annotation file from text in the format of \code{rdann} \\
\code{sortann} & sort an annotation file into canonical order \\
\code{ann2rr} & print the intervals between annotations (RR intervals) \\
\code{rr2ann} & write an annotation file from a list of RR intervals \\
\midrule
\multicolumn{2}{@{}l}{\textit{Format converters} (\cref{sec:tools-converters})} \\
\code{edf2mit} & convert an EDF or EDF+ file to a WFDB record \\
\code{mit2edf} & convert a WFDB record to an EDF or EDF+ file \\
\code{rdedfann} & print the annotations of an EDF+ file as \code{rdann} text \\
\code{wav2mit} & make a WAV file readable as a WFDB record \\
\code{mit2wav} & convert a WFDB record to a WAV file \\
\code{wfdb2mat} & convert a WFDB record to a MATLAB/Octave \file{.mat} file \\
\code{a2m} & convert an AHA-format annotation file to WFDB (MIT) format \\
\code{m2a} & convert a WFDB annotation file to AHA format \\
\code{ad2m} & convert an AHA-format signal file (tape or CD-ROM) to a WFDB record \\
\code{md2a} & convert a WFDB record to an AHA-format signal file \\
\code{ahaecg2mit} & convert AHA DVD files (binary or text) to WFDB records \\
\code{makeid} & write an AHA-format record ID block \\
\code{readid} & read an AHA-format record ID block \\
\code{revise} & convert an obsolete-format header to the current format \\
\code{parsescp} & read an SCP-ECG file, describe it, and convert it to a record \\
\end{xltabular}

% ==========================================================================================
\section{Signal tools}
\label{sec:tools-signal}

% ------------------------------------------------------------------------------------------
\apientry{rdsamp}{Print the samples of a record}
\begin{lstlisting}[style=signature,language={}]
usage: rdsamp -r RECORD [OPTIONS ...]
\end{lstlisting}
Reads the signals of a record and prints one line per sample interval: the sample number
(or, with \opt{-p}, the time) and the sample of each selected signal, in ADC units, or in
physical units with \opt{-p}. The output is plain text in columns (tab-separated), CSV
(\opt{-c}) or WFDB-XML (\opt{-X}), which \cmd{wrsamp} reads back. An invalid sample is
printed as \code{-} in physical units, and as the value $-32768$ in ADC units.
\begin{params}
\item[-r] The input record, \textit{RECORD} (required).
\item[-c] Use the CSV (comma-separated value) output format: the column headings of
  \opt{-v} are quoted (\code{'sample \#','MLII','V5'}), and the columns are separated by
  commas.
\item[-f] Begin at the time that follows (\cref{sec:tools-times}); by default, the beginning
  of the record.
\item[-h] Print the usage summary, and exit with status 0.
\item[-H] Read multifrequency records in high-resolution mode: one line per sample of the
  signals with the most samples per frame, the samples of the other signals repeated (by
  default, one line per frame, the samples of a frame averaged).
\item[-l] Stop after the time interval that follows (\code{hh:mm:ss}, or \code{s} followed
  by a number of samples), counted from the beginning time; a \opt{-t} time before the end
  of that interval still applies.
\item[-p] Print times and samples in physical units, with 3 decimal places (times in
  seconds by default). A letter appended to the option chooses the time column:
  \opt{-pd} the time of day and the date, when the record has a base time (elapsed time
  otherwise), \opt{-pe} the elapsed time as \code{hh:mm:ss.mmm}, \opt{-ph} hours,
  \opt{-pm} minutes, \opt{-ps} seconds (the default), \opt{-pS} the sample number.
\item[-P] As \opt{-p}, with 8 decimal places (columns 15 characters wide); \opt{-Pd},
  \opt{-Pe}, \opt{-Ph}, \opt{-Pm}, \opt{-Ps} and \opt{-PS} choose the time column.
\item[-s] Print only the signals that follow, given by number (0, 1, ...) or by name, in
  the order given: \opt{-s} \textit{SIGNAL} \code{[}\textit{SIGNAL} \code{...]}. A signal that
  the record does not have is an error.
\item[-S] Begin at the first valid sample of the signal that follows (a number or a name),
  at or after the \opt{-f} time.
\item[-t] Stop at the time that follows (by default, the end of the record).
\item[-v] Print column headings (the signal names, and with \opt{-p} the units).
\item[-X] Print WFDB-XML (an XML prolog, the sampling frequency, the number of signals,
  the headings and the samples); implies \opt{-c} and \opt{-v}.
\end{params}
The first four samples of \code{100s}, in ADC units, then four samples from 1 second in
physical units (\code{100s} gives no units, so they are shown as \code{nd}), and the CSV
format:
\begin{lstlisting}[style=text]
$ rdsamp -r 100s -t s4
              0     995    1011
              1     995    1011
              2     995    1011
              3     995    1011
$ rdsamp -r 100s -f 1 -t s364 -p -v
   Elapsed time    MLII      V5
      (seconds)    (nd)    (nd)
          1.000  -0.535  -0.205
          1.003  -0.505  -0.080
          1.006  -0.415   0.015
          1.008  -0.300   0.065
$ rdsamp -r 100s -t s3 -c -v
'sample #','MLII','V5'
0,995,1011
1,995,1011
2,995,1011
\end{lstlisting}
\returns{Exit status 0 when the samples have been printed (also when a read error or a
checksum error ends the output early, as in C).}
\raises{Exit status 1 for an error in the arguments (or a missing \opt{-r}); 2 if the record
cannot be opened, a signal of \opt{-s} or \opt{-S} cannot be found, or the sampling
frequency or a gain that the output needs is not specified.}
\textit{Differences from the C tool} (\cref{chap:behavior}): the sampling frequency is
required for time arguments other than sample numbers, for time columns (\opt{-p} without
\opt{-pS}, or with \opt{-v}) and for \opt{-X}; the gain of each selected signal is required
for physical units (\code{rdsamp: signal N of record R is uncalibrated}, status 2);
unspecified units are shown as \code{nd} (C: \code{mV}). The C tool assumed 250\,Hz and
200\,adu/mV. The quotes and backslashes of signal names are escaped in CSV headings, where
the C tool did not terminate.

% ------------------------------------------------------------------------------------------
\apientry{wrsamp}{Write a record from a text file of samples}
\begin{lstlisting}[style=signature,language={}]
usage: wrsamp [OPTIONS ...] [COLUMN ...]
\end{lstlisting}
Reads a text file with a sample of each signal on each line, and writes the samples to a
signal file and, with \opt{-o}, a header. Fields are separated by tabs, spaces, commas or
any combination of these (consecutive separators count as one). If the first line contains
letters, it gives the signal names; if the second line also contains letters, it gives the
units (with or without parentheses). WFDB-XML, as written by \cmd{rdsamp -X}, is read as
well. Non-numeric fields are written as invalid samples. Samples are rounded to the nearest
integer after scaling (\opt{-x}).
\begin{params}
\item[COLUMN] The fields to copy, by number (the leftmost field is column 0), in the order
  of the output signals; by default, all fields in order (see \opt{-z}).
\item[-c] Accepted for compatibility; it has no effect (as in the C tool).
\item[-d] Add dither (random noise of up to one unit) to the samples before rounding.
\item[-f] Start copying at the input line that follows (a number, $\ge 1$; the first line
  of samples is line 0).
\item[-F] The sampling frequency to write in the header (Hz). Required with \opt{-o}.
\item[-G] The gain (adu per physical unit) to write in the header; a quoted list of values
  (\opt{-G "100 50"}) gives one per signal, the last applying to the remaining signals. A
  gain of 0 declares the signals uncalibrated. Required with \opt{-o}.
\item[-h] Print the usage summary, and exit with status 1.
\item[-i] Read the input from the file that follows (by default, or with \code{-}, the
  standard input). The file is read from the input directory (\cref{sec:tools-files}).
\item[-o] Write the record that follows: the signal file \textit{RECORD}\file{.dat} and the
  header \textit{RECORD}\file{.hea}. Without \opt{-o}, the samples are written to the
  standard output (in format 16 unless \opt{-O} is given) and no header is written.
\item[-O] The storage format of the output (default 16); any format of
  \cref{tab:formats}.
\item[-r] The character that follows separates input lines (records), instead of the
  line feed.
\item[-R] The ADC resolution in bits; a quoted list gives one per signal, as for \opt{-G}.
  Required with \opt{-o}.
\item[-s] Any character of the string that follows separates input fields (by default:
  space, tab, carriage return or comma; only tab if the first line contains a tab).
\item[-t] Stop copying at the input line that follows (a number, $\ge 1$).
\item[-U] The units of the signals; a quoted list gives one per signal, as for \opt{-G}.
  Required with \opt{-o} for calibrated signals, unless the input gives the units.
\item[-x] Multiply the input values by the scale factor that follows; a quoted list gives
  one per signal.
\item[-z] Do not copy column 0 (usually a sample number), unless it is selected by a
  \textit{COLUMN} argument.
\end{params}
The obsolete option \opt{-l} (and its argument) is ignored, with a message, as in C. For
example, a record written from the output of \cmd{rdsamp} (column 0, the sample number, is
skipped by \opt{-z}):
\begin{lstlisting}[style=text]
$ rdsamp -r 100s -t s5 > s.txt
$ wrsamp -i s.txt -o rec -F 360 -G 200 -R 11 -U mV -z
$ cat rec.hea
rec 2 360 5
rec.dat 16 200/mV 11 0 995 4975 0 col 1
rec.dat 16 200/mV 11 0 1011 5055 0 col 2
$ wrsamp -i s.txt -o rec2
wrsamp: the sampling frequency must be specified (-F)
\end{lstlisting}
\returns{Exit status 0 when the record has been written.}
\raises{Exit status 1 for an error in the arguments, a missing \opt{-F}, \opt{-G}, \opt{-R}
or \opt{-U} (\code{wrsamp: the ... must be specified (-X)}), or an XML input without
samples; 2 if the input file cannot be read or the output cannot be opened; 3 if the input
contains no line separator.}
\textit{Differences from the C tool} (\cref{chap:behavior}): when a record is written
(\opt{-o}), the sampling frequency (\opt{-F}), the gains (\opt{-G}), the ADC resolutions
(the new option \opt{-R}) and, for calibrated signals, the units (the new option \opt{-U})
must be given; the C tool assumed 250\,Hz, 200\,adu/mV, a resolution based on the format,
and millivolts. A header written with \opt{-U mV} says \code{200/mV} where the C tool
wrote \code{200}. The \opt{-i} file is read from the input directory (C: from the current
directory).

% ------------------------------------------------------------------------------------------
\apientry{wfdbdesc}{Describe the signals of a record}
\begin{lstlisting}[style=signature,language={}]
usage: wfdbdesc RECORD [-readable]
\end{lstlisting}
Prints the contents of a record's header in readable form: the info strings (\code{Notes}),
the starting time, the length, the sampling frequency, and for each signal its file,
description, gain and units, initial value, storage format, ADC resolution and zero,
baseline and checksum. For a multi-segment record, it also gives the layout (fixed or
variable) and, for a fixed layout, the length of the first segment, whose signals it
describes.
\begin{params}
\item[RECORD] The record to describe.
\item[-readable] Describe only the signals whose files can be opened (by default, every
  signal named in the header, without opening the files). Any abbreviation of
  \opt{-readable} is accepted.
\end{params}
\begin{lstlisting}[style=text]
$ wfdbdesc 100s
Record 100s
Notes
=====
 69 M 1085 1629 x1
 Aldomet, Inderal
Produced by xform from record 100, beginning at 0:0
=====

Starting time: not specified
Length:     1:00.000 (21600 sample intervals)
Sampling frequency: 360 Hz
2 signals
Group 0, Signal 0:
 File: 100s.dat
 Description: MLII
 Gain: 200 adu/nd
 Initial value: 995
 Storage format: 212
 I/O: can be unbuffered
 ADC resolution: 11 bits
 ADC zero: 1024
 Baseline: 1024
 Checksum: 21537
Group 0, Signal 1:
 ...
\end{lstlisting}
\returns{Exit status 0.}
\raises{Exit status 1 without a record name; 2 if the header cannot be read, or the record
has no signals.}
\textit{Differences from the C tool} (\cref{chap:behavior}): an unspecified sampling
frequency or ADC resolution is shown as \code{not specified} (lengths are then given in
sample intervals only), a gain of 0 as \code{uncalibrated} (C:
\code{uncalibrated; assume 200}), and unspecified units as \code{adu/nd}; the C tool assumed
250\,Hz, 12 bits and millivolts. For a record whose header gives no length, the size of the
first signal file, from which the length is computed, is measured in the header's
directory, where the library reads the file (C: in the current directory). The notes are
the info strings that \code{getinfo} returns, after the record line found as the header
parser finds it (\cref{sec:record-line-divergence}): for a header whose first line is
empty, the comments before the record line are not printed (the C tool, and
\cmd{wfdbdesc -{}-legacy}, print them as notes).

% ------------------------------------------------------------------------------------------
\apientry{wfdbtime}{Convert times and sample numbers}
\begin{lstlisting}[style=signature,language={}]
usage: wfdbtime [ -r RECORD [ TIME ] ]
\end{lstlisting}
Converts each time argument, at the sampling frequency of the record, and prints three
columns: the sample number (\code{s} followed by the number), the elapsed time from the
beginning of the record, and the time of day (the elapsed time again if the record has no
base time).
\begin{params}
\item[-r] The record, \textit{RECORD}, whose sampling frequency and base time are used. A
  later \opt{-r} switches to another record for the times that follow it.
\item[TIME] A time to convert, in any form of \cref{sec:tools-times}; any number of times
  may follow the record.
\end{params}
\begin{lstlisting}[style=text]
$ wfdbtime -r 100s 1:00 s100
         s21600        1:00.000        1:00.000
           s100        0:00.278        0:00.278
\end{lstlisting}
\returns{Exit status 0.}
\raises{Exit status 1 without \opt{-r} (also for \opt{-h}, which prints the usage); 2 if
the record cannot be opened or its header gives no sampling frequency.}
\textit{Differences from the C tool} (\cref{chap:behavior}): the record must give its
sampling frequency (C: 250\,Hz assumed).

% ------------------------------------------------------------------------------------------
\apientry{time2sec}{Convert a time to seconds}
\begin{lstlisting}[style=signature,language={}]
usage: time2sec [-r RECORD] TIME
\end{lstlisting}
Prints a time interval in seconds. Without \opt{-r}, \textit{TIME} is converted at
1000 samples per second and printed with up to 12 significant digits; with \opt{-r}, it is
converted at the sampling frequency of the record and printed with at most three decimals
(milliseconds). A time of day in brackets (\code{[9:0:0]}) gives the elapsed time from the
beginning of the record.
\begin{params}
\item[-r] The record, \textit{RECORD}, whose sampling frequency (and base time) are used.
\item[TIME] The time to convert (\cref{sec:tools-times}).
\end{params}
\begin{lstlisting}[style=text]
$ time2sec 1:30.5
90.5
$ time2sec -r 100s s1000
2.778
\end{lstlisting}
\returns{Exit status 0.}
\raises{Exit status 1 for other arguments (also for \opt{-h}), after the usage summary.
Exit status 2 with \opt{-r} if the header of the record cannot be read, or gives no
sampling frequency (\code{time2sec: sampling frequency of record RECORD is not
specified}).}
\textit{Differences from the C tool}: with \opt{-r}, a record whose header gives no
sampling frequency (or cannot be read) is an error. The C tool converted the time at its
default frequency of 250\,Hz, a result unrelated to the record, because \wfdbx{} assumes no
frequency that a record does not give (\cref{sec:no-defaults}); until version 1.2.1,
\wfdbx{} divided by the unspecified frequency (0) and printed \code{-9223372036854775808}
with the exit status 0. With \opt{-{}-legacy}, the C behavior (250\,Hz).

% ------------------------------------------------------------------------------------------
\apientry{sampfreq}{Print the sampling frequency of a record}
\begin{lstlisting}[style=signature,language={}]
usage: sampfreq [OPTION] RECORD
\end{lstlisting}
Prints the sampling frequency (frame rate) of a record, in Hz. With \opt{-a}, prints the
name and sampling frequency of each signal, tab-separated; with \opt{-H}, the highest
sampling frequency of a multifrequency record.
\begin{params}
\item[RECORD] The record.
\item[-a] List all signals and their sampling frequencies.
\item[-h] Listed in the usage summary, which is printed (status 1) when no argument is
  given; as in C, \cmd{sampfreq -h} alone takes \code{-h} as a record name.
\item[-H] Print the maximum sampling frequency of a multifrequency record.
\end{params}
\begin{lstlisting}[style=text]
$ sampfreq 100s
360
$ sampfreq -a mfreq
slow    100
medium  200
fast    400
$ sampfreq -H mfreq
400
\end{lstlisting}
\returns{Exit status 0; a header that cannot be read prints the library's message and
\code{-1}.}
\raises{Exit status 1 without arguments.}
\textit{Differences from the C tool} (\cref{chap:behavior}): for a header that does not
give a sampling frequency, the library reports 0, which is printed (C: 250).

% ------------------------------------------------------------------------------------------
\apientry{signame}{Print the names of the signals of a record}
\begin{lstlisting}[style=signature,language={}]
usage: signame -r RECORD [OPTIONS ...]
\end{lstlisting}
Prints the name (description) of each signal of the record, one per line, or of the signals
whose numbers follow \opt{-s}. A number that is not that of a signal prints
\code{[INVALID]}.
\begin{params}
\item[-r] The record, \textit{RECORD} (required).
\item[-h] Print the usage summary, and exit with status 0.
\item[-s] Print the names of the signals whose numbers follow:
  \opt{-s} \textit{SIGNAL} \code{[}\textit{SIGNAL} \code{...]}.
\end{params}
\begin{lstlisting}[style=text]
$ signame -r 100s -s 1 0 2
V5
MLII
[INVALID]
\end{lstlisting}
\returns{Exit status 0.}
\raises{Exit status 1 for an error in the arguments; 2 if the record cannot be opened or
has no signals.}
\textit{Differences from the C tool}: none, apart from \opt{-{}-legacy} (the C tool also
read past the end of its array for the number equal to the number of signals; \wfdbx{}
prints \code{[INVALID]}).

% ------------------------------------------------------------------------------------------
\apientry{signum}{Print the numbers of the signals with given names}
\begin{lstlisting}[style=signature,language={}]
usage: signum -r RECORD [OPTIONS ...]
\end{lstlisting}
For each name that follows \opt{-s}, prints the numbers of the signals of the record with
that exact name (tab-separated, when several have it), or \code{X} if none has it, one line
per name.
\begin{params}
\item[-r] The record, \textit{RECORD} (required).
\item[-h] Print the usage summary, and exit with status 0.
\item[-s] The signal names to look for (required):
  \opt{-s} \textit{SIGNAL} \code{[}\textit{SIGNAL} \code{...]}.
\end{params}
\begin{lstlisting}[style=text]
$ signum -r 100s -s V5 MLII ABP
1
0
X
\end{lstlisting}
\returns{Exit status 0.}
\raises{Exit status 1 for an error in the arguments, or without \opt{-r} or \opt{-s};
2 if the record cannot be opened.}
\textit{Differences from the C tool}: none, apart from \opt{-{}-legacy}.

% ------------------------------------------------------------------------------------------
\apientry{sigamp}{Measure signal amplitudes}
\begin{lstlisting}[style=signature,language={}]
usage: sigamp -r RECORD [OPTIONS ...]
\end{lstlisting}
Measures the amplitude of each signal, up to \textit{NMAX} times, and prints the trimmed
mean of the measurements (the largest and smallest 5\% excluded), one tab-separated column
per signal. With an annotator (\opt{-a}), each measurement is the peak-to-peak amplitude in
a window around a beat annotation (an annotation that \code{map1} maps to \code{NORMAL});
without one, it is the RMS amplitude over consecutive windows of 1 second (\opt{-w}). The
amplitudes are in ADC units, or in physical units with \opt{-p}.
\begin{params}
\item[-r] The record, \textit{RECORD} (required).
\item[-a] Measure peak-to-peak amplitudes near the beat annotations of the annotator that
  follows. \opt{-a} and \opt{-w} cannot be used together (\opt{-w} is then ignored, with a
  message).
\item[-d] The window of the peak-to-peak measurements, \textit{DT1} \textit{DT2}: the time
  before and the time after each annotation (default 0.05 seconds each); used only with
  \opt{-a}.
\item[-f] Begin at the time that follows.
\item[-h] Print the usage summary, and exit with status 0.
\item[-H] Read multifrequency records in high-resolution mode.
\item[-n] Make up to the number of measurements that follows per signal (default 300).
\item[-p] Print the amplitudes in physical units. A letter appended chooses the time
  format of the individual measurements (\opt{-q}, \opt{-v}): \opt{-pd} the time of day,
  \opt{-pe} \code{hh:mm:ss}, \opt{-ph} hours, \opt{-pm} minutes, \opt{-ps} seconds (the
  default), \opt{-pS} sample numbers.
\item[-q] Quick mode: print the individual measurements only, each with its time.
\item[-t] Stop at the time that follows.
\item[-v] Verbose mode: print the individual measurements, then the trimmed means
  (\code{Trimmed mean}).
\item[-w] The window of the RMS measurements, \textit{DTW} (default 1 second).
\end{params}
The environment variable \code{WFDBGVMODE} sets the resolution mode when \opt{-H} is not
given, as in C.
\begin{lstlisting}[style=text]
$ sigamp -r 100s
        32.97   23.7155
$ sigamp -r 100s -a atr -n 4 -v
0.213889        265     176
1.027778        295     193
1.838889        306     216
2.627778        301     216
Trimmed mean    291.75  200.25
\end{lstlisting}
\returns{Exit status 0.}
\raises{Exit status 1 for an error in the arguments; 2 if the record or the annotator cannot
be opened, the record gives no sampling frequency, or a signal is uncalibrated with
\opt{-p}.}
\textit{Differences from the C tool} (\cref{chap:behavior}): the record must give its
sampling frequency (the windows are in seconds), and, with \opt{-p}, the gain of each signal
(C: 250\,Hz and 200\,adu/mV assumed).

% ------------------------------------------------------------------------------------------
\apientry{snip}{Copy a segment of a record to a new record}
\begin{lstlisting}[style=signature,language={}]
usage: snip -i IREC -n NREC [OPTIONS ...]
\end{lstlisting}
Copies the samples of a record between two times to a new record, with the same signals,
storage format and signal specifications, and optionally its annotations, with times
relative to the beginning of the excerpt. The header of the new record keeps the info
strings of the input and adds \code{Produced by snip from record IREC, beginning at TIME}.
To change the sampling frequency, gains or signals, \cmd{xform} is the tool to use.
\begin{params}
\item[-i] The input record, \textit{IREC} (required; \opt{-r} is accepted as a synonym).
\item[-n] The new record, \textit{NREC} (required). It must differ from \textit{IREC}: the
  header of the new record is written over that of a record of the same name.
\item[-a] Copy the annotations of the annotators that follow:
  \opt{-a} \textit{ANNOTATOR} \code{[}\textit{ANNOTATOR} \code{...]}.
\item[-f] Begin at the time that follows (default 0).
\item[-h] Print the usage summary, and exit with status 0.
\item[-l] Copy only the duration that follows (\code{hh:mm:ss} or \code{s} followed by a
  number of samples); overrides \opt{-t}.
\item[-m] Keep the segments of a multi-segment input record, if possible. The master
  header of the new record is created in the output directory, as the other files are.
\item[-O] Write the samples in the storage format that follows (by default, the format of
  the input signals, or, if they differ, the best fit to the largest ADC resolution).
\item[-s] Do not write info strings in the new header (neither those of the input nor the
  \code{Produced by snip} line).
\item[-t] Stop at the time that follows (default: the end of the record).
\end{params}
\begin{lstlisting}[style=text]
$ snip -i 100s -n ex -f 10 -t 20 -a atr
$ cat ex.hea
ex 2 360 3600
ex.dat 212 200 11 1024 946 -16262 0 MLII
ex.dat 212 200 11 1024 969 18142 0 V5
# 69 M 1085 1629 x1
# Aldomet, Inderal
#Produced by xform from record 100, beginning at 0:0
#Produced by snip from record 100s, beginning at 10
\end{lstlisting}
\returns{Exit status 0.}
\raises{Exit status 1 for an error in the arguments or an improper interval; 2 if a header
cannot be read or written, or the input record gives no sampling frequency.}
\textit{Differences from the C tool} (\cref{chap:behavior}): the input record must give its
sampling frequency, for the new header (C: 250\,Hz assumed); an output signal file is never
one of the input files; when the last frame to copy cannot be read and no frame was read
before it, 0 is written (C: uninitialized values); the master header of \opt{-m} is created
in the output directory (C: in the current directory).

% ------------------------------------------------------------------------------------------
\apientry{xform}{Convert a record (format, frequency, gain, signals)}
\begin{lstlisting}[style=signature,language={}]
usage: xform -i IREC [OPTIONS ...]
\end{lstlisting}
Copies the signals of a record to a new one, converting them: to another storage format,
sampling frequency (by interpolation), gain, ADC resolution or zero level, and to a subset
or another order of the signals; annotations can be copied too, with their times converted.
The output is described either by an existing header (\opt{-o}) or by answers to prompts
(printed on the standard error and read from the terminal, the standard input, or a script
given by \opt{-S}): the output record name (unless \opt{-n} gives it), the number of
signals, the sampling frequency, the directory of the signal files, the storage format,
whether all signals go in one file, and for each signal its gain, ADC resolution and ADC
zero (and with \opt{-N} its description and units). Samples out of the range of the output
format wrap around unless \opt{-c} is given. While it converts, it prints
\code{Checking input signals ...} and progress dots.
\begin{params}
\item[-i] The input record, \textit{IREC} (required).
\item[-a] Copy the annotations of the annotators that follow:
  \opt{-a} \textit{ANNOTATOR} \code{[}\textit{ANNOTATOR} \code{...]}.
\item[-c] Clip the output samples to the range of the format (by default, they wrap around).
\item[-d] Add dither to the input if the sampling frequency or the gain changes.
\item[-f] Begin at the time that follows.
\item[-h] Print the usage summary, and exit with status 0.
\item[-H] Open the input record in high-resolution mode.
\item[-M] Process a multifrequency record without changing the frequencies of its signals.
\item[-n] Create a header for the record \textit{NREC} that follows, with the signal
  descriptions of \textit{IREC}.
\item[-N] As \opt{-n}, but with the signal descriptions of \textit{OREC} (or asked for).
\item[-o] Write the output signal files as the existing header of the record \textit{OREC}
  that follows describes them (no prompts).
\item[-s] Write only the input signals that follow, by number or name, in that order:
  \opt{-s} \textit{SIGNAL} \code{[}\textit{SIGNAL} \code{...]}. A signal that the record does not
  have is skipped, with a warning.
\item[-S] Take the answers to the prompts from the text file \textit{SCRIPT} that follows,
  one answer per line (an empty line takes the default answer). The file is read from the
  input directory (\cref{sec:tools-files}).
\item[-t] Stop at the time that follows.
\item[-u] Adjust annotation times, if needed, to make them unique.
\end{params}
For example, a copy of the first 2 seconds of \code{100s} in format 16, with the answers in
the file \file{ans.txt} (empty lines, except \code{16} on the fourth: the default number
of signals, frequency and directory, then format 16, then the default answers):
\begin{lstlisting}[style=text]
$ xform -i 100s -n x16 -S ans.txt -t 2 2>/dev/null
$ cat x16.hea
x16 2 360 720
x16.dat 16 200 11 1024 995 -26770 0 MLII
x16.dat 16 200 11 1024 1011 -12260 0 V5
# 69 M 1085 1629 x1
# Aldomet, Inderal
#Produced by xform from record 100, beginning at 0:0
#Produced by xform from record 100s, beginning at 0:0
\end{lstlisting}
\returns{Exit status 0.}
\raises{Exit status 1 for an error in the arguments, identical \opt{-n} and \opt{-o} names,
\opt{-M} with a change of frequency, an output frequency of 0, a missing ADC resolution, a
\opt{-t} time before the \opt{-f} time, or the end of a script before the last answer; 2 if
a record cannot be opened or written, a sampling frequency is not specified, or a requested
rescaling is impossible.}
\textit{Differences from the C tool} (\cref{chap:behavior}): the input record (and the
record of \opt{-o}) must give a sampling frequency, and an output frequency of 0 is refused
(C: 250\,Hz); rescaling an uncalibrated signal to a given gain, and rescaling by ADC
resolution when a resolution is unspecified, are refused; an ADC resolution is offered as
the default only if it comes from the input (or the previous output signal) and suits the
output format, and must be given otherwise; unspecified units are shown as \code{nd} in the
prompts; an output signal file that would be one of the input files is refused
(\code{osigfopen: output file NAME is also an input file}). With \opt{-M} and an output
frame larger than the input frame, the samples beyond the input frame are 0 (C: a crash).
The script of \opt{-S} is read from the input directory (C: from the current directory).

% ------------------------------------------------------------------------------------------
\apientry{wfdbcat}{Copy WFDB files to the standard output}
\begin{lstlisting}[style=signature,language={}]
usage: wfdbcat FILE [FILE ...]
\end{lstlisting}
Copies each file, found as the library finds files (\cref{sec:tools-files}: relative to the
input directory, or as a full path or a URL), to the standard output, byte for byte. A file
that cannot be found is reported (\code{wfdbcat: `FILE' not found}) and skipped.
\begin{params}
\item[FILE] A file to copy, named with its suffix (\code{100s.hea}).
\end{params}
\begin{lstlisting}[style=text]
$ wfdbcat 100s.hea
100s 2 360 21600
100s.dat 212 200 11 1024 995 21537 0 MLII
100s.dat 212 200 11 1024 1011 -3962 0 V5
# 69 M 1085 1629 x1
# Aldomet, Inderal
#Produced by xform from record 100, beginning at 0:0
\end{lstlisting}
\returns{Exit status 0, also when a file was not found.}
\raises{Exit status 1 without arguments, or with \opt{-h} alone.}
\textit{Differences from the C tool} (\cref{chap:behavior}): each file is read from where it
is named, not found by a search of the WFDB path.

% ------------------------------------------------------------------------------------------
\apientry{wfdbwhich}{Print the path of a WFDB file}
\begin{lstlisting}[style=signature,language={}]
usage: wfdbwhich [-r RECORD ] FILENAME
 -or-  wfdbwhich FILE-TYPE RECORD
\end{lstlisting}
Prints the path by which the library opens a file: a file named by \textit{FILENAME}, or
the file of type \textit{FILE-TYPE} (a suffix such as \code{hea}, \code{dat} or \code{atr})
of a record. A file that cannot be found is reported with the input directory
(\code{(WFDB path = ...)}).
\begin{params}
\item[-r] Open the record \textit{RECORD} first.
\item[FILENAME] The name of the file to find.
\item[FILE-TYPE] The type (suffix) of the file of the record to find.
\item[RECORD] The record.
\end{params}
\begin{lstlisting}[style=text]
$ wfdbwhich hea 100s
100s.hea
$ wfdbwhich 100x.dat
`100x.dat' not found
(WFDB path = .)
\end{lstlisting}
\returns{Exit status 0 when the file is found.}
\raises{Exit status 1 for other arguments (also \opt{-h}); 2 if the file is not found.}
\textit{Differences from the C tool} (\cref{chap:behavior}): the file is looked for only
where it is named (the input directory, a full path or a URL), not by a search of the WFDB
path, and the path printed in the message is the input directory.

% ==========================================================================================
\section{Annotation tools}
\label{sec:tools-annotation}

% ------------------------------------------------------------------------------------------
\apientry{rdann}{Print the annotations of an annotation file}
\begin{lstlisting}[style=signature,language={}]
usage: rdann -r RECORD -a ANNOTATOR [OPTIONS ...]
\end{lstlisting}
Prints one line per annotation: the time (elapsed time, or time of day in brackets when the
record has a base time), the sample number, the mnemonic of the annotation type, the
subtype, channel and number fields, and the auxiliary information (after a tab), if any.
Options select the annotations to print and the time format. This text is the input format
of \cmd{wrann}. The old syntax of the C tool, \cmd{rdann} \textit{ANNOTATOR RECORD}
\code{[}\textit{FROM} \code{[}\textit{TO} \code{[}\textit{TYPE} \code{...]]]}, is accepted too.
\begin{params}
\item[-r] The record, \textit{RECORD} (required).
\item[-a] The annotator, \textit{ANNOTATOR} (required).
\item[-c] Print only the annotations with the channel number that follows (from $-128$
  to 127).
\item[-e] Print the times as elapsed times, even for a record with a base time. Only one of
  \opt{-e} and \opt{-x} can be used (the last one given applies).
\item[-f] Start at the time that follows; also \code{a} followed by a number (from that
  annotation number), or \code{b} or \code{\#} followed by a number (from that beat number).
\item[-F] The sampling frequency, used if the record has no header or its header gives
  none.
\item[-h] Print the usage summary, and exit with status 0.
\item[-n] Print only the annotations with the \code{num} value that follows.
\item[-p] Print only the annotations of the types that follow, given by their mnemonics:
  \opt{-p} \textit{TYPE} \code{[}\textit{TYPE} \code{...]} (\opt{-p N V}).
\item[-s] Print only the annotations with the subtype that follows.
\item[-t] Stop at the time that follows; also \code{a}, \code{b} or \code{\#} followed by a
  number, as for \opt{-f} (\code{\#}: that number of beats).
\item[-v] Print column headings.
\item[-x] Print the times in seconds, minutes and hours (with enough decimal places to
  resolve the time resolution of the file).
\end{params}
\begin{lstlisting}[style=text]
$ rdann -r 100s -a atr -t 3 -v
      Time   Sample #  Type  Sub Chan  Num      Aux
    0:00.050       18     +    0    0    0      (N
    0:00.214       77     N    0    0    0
    0:01.028      370     N    0    0    0
    0:01.839      662     N    0    0    0
    0:02.628      946     N    0    0    0
$ rdann -r 100s -a atr -t 2 -x
    0.050   0.00083 0.0000139     +    0    0    0      (N
    0.214   0.00356 0.0000594     N    0    0    0
    1.028   0.01713 0.0002855     N    0    0    0
    1.839   0.03065 0.0005108     N    0    0    0
\end{lstlisting}
\returns{Exit status 0.}
\raises{Exit status 1 for an error in the arguments, or if the sampling frequency is not
specified (\code{rdann: sampling frequency of record R is not specified (use -F)}); 2 if the
annotation file cannot be opened or a starting annotation or beat number is beyond its end.}
\textit{Differences from the C tool} (\cref{chap:behavior}): the new option \opt{-F} gives
the sampling frequency of a record without a header (or whose header gives none), and is
then required; the C tool assumed 250\,Hz.

% ------------------------------------------------------------------------------------------
\apientry{wrann}{Write an annotation file from text}
\begin{lstlisting}[style=signature,language={}]
usage: wrann -r RECORD -a ANNOTATOR <TEXT-FILE
\end{lstlisting}
Reads annotations from the standard input, in the format printed by \cmd{rdann}, and writes
the annotation file \textit{RECORD}\file{.}\textit{ANNOTATOR}. The time column is skipped and
the sample number is used, except on lines that begin with a time of day in brackets, whose
time is converted. A mnemonic that is not an annotation code is written as a \code{NOTE}
(comment) annotation whose auxiliary information begins with that mnemonic. Text after a
tab is the auxiliary information. The time resolution recorded in the file is the sampling
frequency of the record, if its header gives one. Annotations that are not in canonical
order are sorted when the file is closed (\cref{chap:behavior}).
\begin{params}
\item[-r] The record, \textit{RECORD} (required).
\item[-a] The output annotator, \textit{ANNOTATOR} (required).
\end{params}
\opt{-h}, not listed in the usage, prints it and exits with status 0.
\begin{lstlisting}[style=text]
$ rdann -r 100s -a atr | wrann -r 100s -a copy
\end{lstlisting}
\returns{Exit status 0.}
\raises{Exit status 1 for an error in the arguments; 2 if the output file cannot be opened.}
\textit{Differences from the C tool} (\cref{chap:behavior}): an out-of-order file is sorted
by the library itself, without losing annotations, instead of by running \cmd{sortann}.

% ------------------------------------------------------------------------------------------
\apientry{sortann}{Sort an annotation file}
\begin{lstlisting}[style=signature,language={}]
usage: sortann -r RECORD -a ANNOTATOR [OPTIONS ...]
\end{lstlisting}
Rearranges the annotations of a file in canonical order: by time, then \code{num}, then
\code{chan}. A later annotation with the same time, \code{num} and \code{chan} replaces an
earlier one, and \code{NOTQRS} annotations are dropped. Without \opt{-o}, the sorted file is
written under the name of the input (in the output directory, so that it replaces the
input when the input and output directories are the same), unless the input is already
sorted: that is reported (\code{sortann: input is already ordered -{}- no output written})
and nothing is written.
\begin{params}
\item[-r] The record, \textit{RECORD} (required). As in C, \opt{-r}\textit{DIR} (the
  directory attached to the option) sets the input directory before the record name.
\item[-a] The input annotator, \textit{ANNOTATOR} (required).
\item[-f] Copy only the annotations from the time that follows (requires \opt{-o}).
\item[-F] The sampling frequency, used if the record has no header or its header gives
  none; then required for \opt{-f} and \opt{-t}.
\item[-h] Print the usage summary, and exit with status 0.
\item[-o] Write the output as the annotator \textit{OUTANN} that follows (it must differ
  from the input annotator); required with \opt{-f} or \opt{-t}.
\item[-t] Copy only the annotations before the time that follows (requires \opt{-o}).
\end{params}
\begin{lstlisting}[style=text]
$ sortann -r 100s -a atr
sortann: input is already ordered -- no output written
\end{lstlisting}
\returns{Exit status 0, also for an input that is already sorted.}
\raises{Exit status 1 for an error in the arguments, a missing \opt{-o}, identical input
and output annotators, or a sampling frequency that is needed and not specified; 2 if an
annotation file cannot be opened.}
\textit{Differences from the C tool} (\cref{chap:behavior}): the new option \opt{-F}; an
out-of-order annotation at time 0 (with number and channel 0) is no longer lost, nor are
those inserted after it (the C tool dropped them; legacy mode does too).

% ------------------------------------------------------------------------------------------
\apientry{ann2rr}{Print the intervals between annotations}
\begin{lstlisting}[style=signature,language={}]
usage: ann2rr -r RECORD -a ANNOTATOR [OPTIONS ...]
\end{lstlisting}
Prints the intervals between consecutive beat (QRS) annotations, one per line, in sample
intervals by default, optionally with the times and types of the annotations that begin and
end each interval. The first interval begins at the starting time.
\begin{params}
\item[-r] The record, \textit{RECORD} (required).
\item[-a] The annotator, \textit{ANNOTATOR} (required).
\item[-A] Print all intervals between annotations, not only RR intervals (overrides \opt{-c}
  and \opt{-p}).
\item[-c] Print only intervals between consecutive valid annotations.
\item[-f] Start at the time that follows.
\item[-F] The sampling frequency, used if the record has no header or its header gives
  none; then required for times and for intervals in units of time.
\item[-h] Print the usage summary, and exit with status 0.
\item[-i] Print the intervals in the format that follows: \code{h} (hours), \code{m}
  (minutes), \code{s} (seconds), each optionally followed by a number of decimal places from
  0 to 15 (default 3; \opt{-i s8}), or \code{t} (\code{hh:mm:ss}).
\item[-p] Print only intervals that end with annotations of the types that follow:
  \opt{-p} \textit{TYPE} \code{[}\textit{TYPE} \code{...]}, given by mnemonics such as
  \code{N} or \code{V}.
\item[-P] Print only intervals that begin with annotations of the types that follow.
\item[-t] Stop at the time that follows.
\item[-v] Print the time of the end of each interval, in the format that follows, if any:
  \code{h}, \code{m} or \code{s} (with decimal places, as for \opt{-i}), \code{t}, or
  \code{T} (date and time, if available); without a format, in sample intervals.
\item[-V] Print the time of the beginning of each interval, as for \opt{-v}.
\item[-w] Print the annotations that end the intervals.
\item[-W] Print the annotations that begin the intervals.
\end{params}
\begin{lstlisting}[style=text]
$ ann2rr -r 100s -a atr -t 4 -i s -V s -w
0.000   0.214   N
0.214   0.814   N
1.028   0.811   N
1.839   0.789   N
2.628   0.792   N
\end{lstlisting}
\returns{Exit status 0.}
\raises{Exit status 1 for an error in the arguments, or a sampling frequency that is needed
and not specified; 2 if the annotation file cannot be opened.}
\textit{Differences from the C tool} (\cref{chap:behavior}): the new option \opt{-F}; the C
tool assumed 250\,Hz.

% ------------------------------------------------------------------------------------------
\apientry{rr2ann}{Write an annotation file from RR intervals}
\begin{lstlisting}[style=signature,language={}]
usage: rr2ann -r RECORD -a ANNOTATOR [OPTIONS ...] <TEXT-FILE
\end{lstlisting}
Reads a column of RR intervals (or of times of occurrence, with \opt{-T}) in sample
intervals from the standard input, and writes an annotation file of \code{N} (normal beat)
annotations at the corresponding times, or of the types of a second column (\opt{-w}). The
first token of a line is the interval; empty lines, non-numeric tokens, and negative or zero
intervals are skipped. If the record has no header, one is created.
\begin{params}
\item[-r] The record, \textit{RECORD} (required).
\item[-a] The output annotator, \textit{ANNOTATOR} (required).
\item[-F] The sampling frequency, used (and required) if the header does not exist or does
  not give one; otherwise ignored, with a warning.
\item[-h] Print the usage summary, and exit with status 0.
\item[-T] The input contains times of occurrence, not RR intervals.
\item[-x] Multiply the input values by the factor that follows (default 1) to obtain sample
  intervals.
\item[-w] Take the annotation type of each line from its second column (a mnemonic).
\end{params}
\begin{lstlisting}[style=text]
$ ann2rr -r 100s -a atr | rr2ann -r 100s -a rr
$ rdann -r 100s -a rr -t 2
    0:00.214       77     N    0    0    0
    0:01.028      370     N    0    0    0
    0:01.839      662     N    0    0    0
\end{lstlisting}
\returns{Exit status 0.}
\raises{Exit status 1 for an error in the arguments, or a sampling frequency that is needed
and not given; 2 if the header cannot be created; 3 if the annotation file cannot be
opened.}
\textit{Differences from the C tool} (\cref{chap:behavior}): \opt{-F} is required when the
record has no header or its header gives no frequency (C: 250\,Hz assumed); a header is
created only when there is none.

% ==========================================================================================
\section{Format converters}
\label{sec:tools-converters}

% ------------------------------------------------------------------------------------------
\apientry{edf2mit}{Convert an EDF or EDF+ file to a WFDB record}
\begin{lstlisting}[style=signature,language={}]
usage: edf2mit [OPTIONS ...]
\end{lstlisting}
Reads a European Data Format file (EDF or EDF+) and writes a WFDB record with the same
data: a header and a signal file in format 16, with the sampling frequencies, gains,
baselines, units and start time of the EDF header. The signals \code{EDF Annotations} of an
EDF+ file are copied as signals (\cmd{rdedfann} reads them). Unknown options are ignored.
\begin{params}
\item[-i] The EDF file to read, \textit{EDFILE} (required); \code{-} reads the standard
  input. It is read from the input directory (\cref{sec:tools-files}).
\item[-b] The input is in big-endian byte order (by default, little-endian).
\item[-h] Print the usage summary, and exit with status 0.
\item[-r] The record to create, \textit{RECORD} (by default, the input file name up to its
  first dot, with every character other than a letter or digit replaced by \code{\_}).
\item[-s] Copy only the signals that follow, by number (from 0):
  \opt{-s} \textit{SIGNAL} \code{[}\textit{SIGNAL} \code{...]}.
\item[-v] Print the fields of the EDF header (debugging output) on the standard output.
\end{params}
\begin{lstlisting}[style=text]
$ edf2mit -i edfp.edf -r edfp2
$ cat edfp2.hea
edfp2 2 10 20 13:24:35 02/03/2002
edfp2.dat 16x10 4.095/uV 11 0 0 26823 0 EEG Fpz-Cz
edfp2.dat 16x3 32767.5 15 0 12331 14629 0 EDF Annotations
\end{lstlisting}
\begin{phibox}
The patient and recording identification fields of an EDF header can identify a person;
\opt{-v} prints them.
\end{phibox}
\returns{Exit status 0.}
\raises{Exit status 1 for a missing or unreadable input file or an invalid record name;
2 if the input is not EDF, or is shorter than the 256-byte fixed part of an EDF header
(\code{edf2mit: incomplete EDF header in FILE}; no file is created).}
\textit{Differences from the C tool} (\cref{chap:behavior}): the \opt{-i} file is read from
the input directory (C: from the current directory); otherwise none, apart from
\opt{-{}-legacy}, except for input on which the C tool fails: a file shorter than the fixed
part of the header (C: fields from uninitialized memory), and signals with 0 samples per
data record (C: a crash).

% ------------------------------------------------------------------------------------------
\apientry{mit2edf}{Convert a WFDB record to an EDF file}
\begin{lstlisting}[style=signature,language={}]
usage: mit2edf -r RECORD [OPTIONS ...]
\end{lstlisting}
Reads a record and writes an EDF file with the same data: 16-bit samples in data records
(blocks) of 10 seconds, or shorter if a block would exceed 61440 bytes. The record name is
written as the patient identification, and the base date and time, if any, as the start
date and time (\code{Startdate not recorded} otherwise). An empty line is printed on the
standard output when the file is complete. Unknown options are ignored.
\begin{params}
\item[-r] The input record, \textit{RECORD} (required).
\item[-h] Print the usage summary, and exit with status 0.
\item[-o] The EDF file to write, \textit{EDFILE} (default: \textit{RECORD}\file{.edf},
  without the directory part of the record name), created in the output directory
  (\cref{sec:tools-files}).
\item[-p] Write an EDF+ file (its name must end in \file{.edf} or \file{.EDF}); warnings say
  what EDF+ requires that the record does not give.
\item[-v] Verbose mode: print the sizes of the header, the blocks and the file, and a dot
  per minute of data.
\end{params}
\begin{lstlisting}[style=text]
$ mit2edf -r 100s -o 100s.edf -v
Converting record 100s to 100s.edf (EDF mode)
 Header block size: 768 bytes
 Data block size: 10 seconds (3600 frames or 14400 bytes)
 Recording length: 1:00 (7 data blocks, 25200 frames, 100800 bytes)
 Total length of file to be written: 101568 bytes
.
\end{lstlisting}
\returns{Exit status 0.}
\raises{Exit status 1 for an error in the arguments or an invalid EDF+ file name; 2 if the
record cannot be opened; 3 if its signals cannot be read, or the sampling frequency, a gain
or an ADC resolution is not specified; 4 if the output file cannot be created; 5 if a frame
does not fit in an EDF block.}
\textit{Differences from the C tool} (\cref{chap:behavior}): the record must give the
sampling frequency and the gain and ADC resolution of every signal (an EDF file stores
physical and digital ranges); the C tool assumed 250\,Hz and 200\,adu/mV and guessed the
resolution from the storage format. Unspecified units are written as \code{nd} (C:
\code{mV}). The EDF file is created in the output directory (C: in the current directory).

% ------------------------------------------------------------------------------------------
\apientry{rdedfann}{Print the annotations of an EDF+ file}
\begin{lstlisting}[style=signature,language={}]
usage: rdedfann -r RECORD [OPTIONS ...]
\end{lstlisting}
Reads the \code{EDF Annotations} signal of an EDF+ file, read as a record
(\cref{chap:records}), and prints its annotations as text in the format of \cmd{rdann},
which \cmd{wrann} can convert to an annotation file. The annotation text is printed in place
of the mnemonic (by default) or as the auxiliary information (\opt{-x}), followed by the
duration, if any.
\begin{params}
\item[-r] The EDF+ file, read as a record, \textit{RECORD} (required): \opt{-r edfp.edf}.
\item[-F] Compute the sample numbers at the frequency \textit{FREQ} that follows, in Hz
  (a value that is not positive gives 1).
\item[-h] Print the usage summary, and exit with status 0.
\item[-v] Print column headings.
\item[-x] Print the annotation text in the auxiliary information, with the type
  \code{\textquotedbl} (\code{NOTE}).
\end{params}
\begin{lstlisting}[style=text]
$ rdedfann -r edfp.edf -v
      Time   Sample #  Type  Sub Chan  Num      Aux
    0:00.500       50 Sleep_stage_W    0    0    0              duration: 1.5
    0:01.250      125 Arousal    0    0    0            duration: 1.5
    0:01.250      125 Spindle_onset    0    0    0              duration: 1.5
\end{lstlisting}
\returns{Exit status 0.}
\raises{Exit status 1 for an error in the arguments; 2 if the record cannot be opened; 3 if
it has no EDF annotations.}
\textit{Differences from the C tool}: none, apart from \opt{-{}-legacy}.

% ------------------------------------------------------------------------------------------
\apientry{wav2mit}{Make a WAV file readable as a record}
\begin{lstlisting}[style=signature,language={}]
usage: wav2mit -i FILE.wav -r RECORD [ OPTIONS ... ]
\end{lstlisting}
Writes a header for a WAV file (PCM, 8 or 16 bits per sample), so that the WAV file itself
is read as the signal file of a record: format 16 (or 80 for 8-bit files) after a 44-byte
prolog, with the sampling frequency, number of signals and resolution of the WAV file.
\begin{params}
\item[-i] The WAV file, \textit{FILE.wav} (required; the name must end in \file{.wav}),
  read from the input directory (\cref{sec:tools-files}).
\item[-G] The gain in adu per physical unit (required; 0 declares the signals
  uncalibrated).
\item[-h] Print the usage summary, and exit with status 1.
\item[-r] The record to create, \textit{RECORD}: its header is \textit{RECORD}\file{.hea}
  (default: the WAV file name without \file{.wav}).
\item[-U] The physical units (required unless the gain is 0).
\end{params}
For example, for the WAV file written by the example of \cmd{mit2wav} below:
\begin{lstlisting}[style=text]
$ wav2mit -i 100s.wav -r w2 -G 200 -U mV
$ cat w2.hea
w2 2 360 21600
100s.wav 16+44 200/mV 16 0 0 0 0
100s.wav 16+44 200/mV 16 0 0 0 0
\end{lstlisting}
\returns{Exit status 0.}
\raises{Exit status 1 for an error in the arguments or a missing \opt{-G} or \opt{-U};
2 if the WAV file cannot be opened; 3 if it is not a supported WAV file (not PCM, more
than 16 bits per sample, inconsistent chunk lengths); 4 if the header cannot be written.}
\textit{Differences from the C tool} (\cref{chap:behavior}): the new options \opt{-G}
(required) and \opt{-U} (required unless the gain is 0); the C tool wrote 12.5\,adu/mV for
8-bit and 6400\,adu/mV for 16-bit files, in millivolts. The \opt{-i} file is read from the
input directory (C: from the current directory).

% ------------------------------------------------------------------------------------------
\apientry{mit2wav}{Convert a record to a WAV file}
\begin{lstlisting}[style=signature,language={}]
usage: mit2wav -o FILE.wav -r RECORD [ OPTIONS ... ]
\end{lstlisting}
Writes the signals of a record to a WAV file (16-bit PCM, one channel per signal, at the
sampling frequency of the record), the samples shifted left to 16 bits (the WAV header gives
the largest ADC resolution as the number of bits per sample). The samples are written
through the library, and the WAV header is then put in front of them, in the same file.
With \opt{-n}, it also writes a header that describes the WAV file as a record, with the
gains and baselines adjusted to the shift.
\begin{params}
\item[-o] The WAV file to write, \textit{FILE.wav} (required; the name must end in
  \file{.wav}), created in the output directory (\cref{sec:tools-files}).
\item[-r] The input record, \textit{RECORD} (required).
\item[-h] Print the usage summary, and exit with status 1.
\item[-n] Write a header for the WAV file, as the record \textit{NEWREC} that follows (it
  must differ from \textit{RECORD}).
\end{params}
\begin{lstlisting}[style=text]
$ mit2wav -r 100s -o 100s.wav -n wavrec
$ cat wavrec.hea
wavrec 2 360 21600
100s.wav 16+44 6400 16 0 -928 -31712 0 MLII
100s.wav 16+44 6400 16 0 -416 4288 0 V5
\end{lstlisting}
\returns{Exit status 0.}
\raises{Exit status 1 for an error in the arguments; 2 if the record cannot be opened; 3 if
its signals cannot be read, it gives no sampling frequency, or the header of \opt{-n}
cannot be written; 4 if the output file cannot be opened.}
\textit{Differences from the C tool} (\cref{chap:behavior}): the record must give its
sampling frequency (C: 250\,Hz assumed); a signal of more than 16 bits is shifted right to
16 bits and gets a gain multiplier of 0 (C: undefined shift, samples written as 0); the WAV
header is put in front of the samples in the file that the library created, in the output
directory (C: in the file of that name in the current directory).

% ------------------------------------------------------------------------------------------
\apientry{wfdb2mat}{Convert a record to a MATLAB/Octave file}
\begin{lstlisting}[style=signature,language={}]
usage: wfdb2mat -r RECORD [OPTIONS ...]
\end{lstlisting}
Writes the samples of a record (in ADC units, less the ADC zero) as a MATLAB/Octave
\file{.mat} file (MATLAB 5 format) holding a matrix \code{val} of one row per signal, and a
header that describes the \file{.mat} file as a record: format 80, 16 or 32 (the smallest
that holds the ADC resolutions of the signals, or their storage formats where the resolution
is not specified) after a 192-byte prolog. Both are named after the record
with an \code{m} appended: \file{100sm.mat} and \file{100sm.hea} for \code{100s} (the help
text, that of the C tool, says \textit{RECORD}\file{.mat}). A summary of the conversion is
printed on the standard output.
\begin{params}
\item[-r] The input record, \textit{RECORD} (required).
\item[-f] Begin at the time that follows.
\item[-h] Print the usage summary, and exit with status 0.
\item[-H] Read multifrequency records in high-resolution mode.
\item[-l] Stop after the time interval that follows (as in \cmd{rdsamp}).
\item[-s] Convert only the signals that follow, by number or name:
  \opt{-s} \textit{SIGNAL} \code{[}\textit{SIGNAL} \code{...]}.
\item[-S] Begin at the first valid sample of the signal that follows, at or after the
  \opt{-f} time.
\item[-t] Stop at the time that follows.
\end{params}
\begin{lstlisting}[style=text]
$ wfdb2mat -r 100s -t 2
Source: record 100s
val has 2 rows (signals) and 720 columns (samples/signal)
Duration:     0:02
Sampling frequency: 360 Hz  Sampling interval: 0.002777777778 sec
Row     Signal  Gain    Base    Units
1       MLII    200     0       nd
2       V5      200     0       nd

To convert from raw units to the physical units shown
above, call the 'rdmat.m' function from the wfdb-matlab
toolbox: https://physionet.org/physiotools/matlab/wfdb-app-matlab/
\end{lstlisting}
\returns{Exit status 0.}
\raises{Exit status 1 for an error in the arguments, an empty interval
(\code{wfdb2mat: no output}), an output larger than 2\,GB, or an output file that cannot be
opened; 2 if the record or a signal cannot be opened, or the
record gives no sampling frequency.}
\textit{Differences from the C tool} (\cref{chap:behavior}): the record must give its
sampling frequency (C: 250\,Hz); uncalibrated signals stay uncalibrated in the output
header (C: 200\,adu/mV); unspecified units are written as \code{nd} (C: \code{mV}); a
one-character record name \code{m} gives the output \code{mm} (C wrote over the input
header).

% ------------------------------------------------------------------------------------------
\apientry{a2m}{Convert AHA-format annotations to WFDB format}
\begin{lstlisting}[style=signature,language={}]
usage: a2m -i AHAFILE [ OPTIONS ... ]
\end{lstlisting}
Converts an AHA-format annotation file (\cref{sec:formats-aha-annotations}) to an MIT-format
annotation file. Four kinds of input are read: files in AHA format written by WFDB tools
(type 0, times in sample intervals), short-format and long-format AHA tape files (types 1
and 2, times in milliseconds), and the compressed files of the AHA Database CD-ROMs and
diskettes (type 3, 6 bytes per annotation, times in milliseconds). The annotation times are
shifted by an offset (\opt{-s}).
\begin{params}
\item[-i] The AHA-format input file, \textit{AHAFILE} (required). It is read from the
  input directory (\cref{sec:tools-files}).
\item[-a] The output annotator (default \code{atr}).
\item[-F] The sampling frequency, used (and then required) if the record header does not
  give one (AHA Database: 250).
\item[-h] Print the usage summary, and exit with status 1.
\item[-r] The output record, \textit{RECORD} (default: the input file name without its
  directory and suffix).
\item[-s] Shift the annotations by the time that follows (default: 0 for type 0, 5:0 for
  type 1, 2:30:0 for type 2, and for type 3 5:0 or 2:30:0 for the records \code{x2xx} and
  \code{x3xx}, or \code{x0xx} and \code{x1xx}).
\item[-t] The type of the input, from 0 to 3, as above (default: 3 if the file name ends
  in \file{.ANO} or \file{.ano}, 0 otherwise).
\end{params}
\begin{lstlisting}[style=text]
$ a2m -i 1201.ANO -F 250
init: can't open header for record 1201
$ rdann -r 1201 -a atr -F 250 -t 5:3
init: can't open header for record 1201
    5:01.000    75250     N    0    0    0
    5:01.812    75453     N    0    0    0
    5:02.628    75657     V    0    0    0
\end{lstlisting}
\returns{Exit status 0.}
\raises{Exit status 1 for an error in the arguments or a sampling frequency that is not
specified; 2 if the input or the annotation files cannot be opened.}
\textit{Differences from the C tool} (\cref{chap:behavior}): the new option \opt{-F}
(the C tool assumed 250\,Hz with a header, or 1 sample per second without one), and
annotation times in milliseconds (types 1 to 3) are converted at the record's sampling
frequency (C: 4\,ms per sample, whatever the frequency); the \opt{-i} file is read from the
input directory (C: from the current directory).

% ------------------------------------------------------------------------------------------
\apientry{m2a}{Convert WFDB annotations to AHA format}
\begin{lstlisting}[style=signature,language={}]
usage: m2a -r RECORD -a MIT-ANNOTATOR AHA-ANNOTATOR [ OPTIONS ... ]
\end{lstlisting}
Copies an MIT-format annotation file to an AHA-format annotation file of the same record
(\cref{sec:formats-aha-annotations}), the annotation types mapped by \code{mamap}. With
\opt{-s}, the annotations before the shift are dropped, and the times of the others are
shifted back and converted to milliseconds, as in the AHA tape files.
\begin{params}
\item[-r] The record, \textit{RECORD} (required).
\item[-a] The input (MIT) and output (AHA) annotators, \textit{MIT-ANNOTATOR}
  \textit{AHA-ANNOTATOR} (required).
\item[-F] The sampling frequency, used if the record header gives none; then required
  for \opt{-s}.
\item[-h] Print the usage summary, and exit with status 1.
\item[-s] Shift the annotation times back by the time that follows, and convert them to
  milliseconds.
\end{params}
\begin{lstlisting}[style=text]
$ m2a -r 100s -a atr aha
\end{lstlisting}
\returns{Exit status 0 (the file \file{100s.aha} is written).}
\raises{Exit status 1 for an error in the arguments or a sampling frequency that is needed
and not specified; 2 if an annotation file cannot be opened.}
\textit{Differences from the C tool} (\cref{chap:behavior}): the new option \opt{-F}; the C
tool assumed 250\,Hz.

% ------------------------------------------------------------------------------------------
\apientry{ad2m}{Convert an AHA-format signal file to a WFDB record}
\begin{lstlisting}[style=signature,language={}]
usage: ad2m -i AHAFILE [ OPTIONS ... ]
\end{lstlisting}
Converts the two signals of an AHA Database signal file, in the old tape format (format 16
with an end-of-data marker) or in the compressed format of the CD-ROMs (\opt{-c}), to a
WFDB record in format 212: \textit{RECORD}\file{.dat} and \textit{RECORD}\file{.hea}.
\begin{params}
\item[-i] The AHA-format input file, \textit{AHAFILE} (required). It is read from the
  input directory (\cref{sec:tools-files}).
\item[-c] The input is compressed (the default if the name ends in \file{.CMP} or
  \file{.cmp}; otherwise the old tape format is assumed).
\item[-f] Begin at the time that follows (default 0, the beginning of the input file).
\item[-F] The sampling frequency in Hz (required; AHA Database: 250).
\item[-G] The gain in adu per physical unit (required; AHA Database: 400; 0 declares the
  signals uncalibrated).
\item[-h] Print the usage summary, and exit with status 1.
\item[-r] The output record, \textit{RECORD} (default: the input file name without its
  directory and suffix).
\item[-R] The ADC resolution in bits, from 1 to 32 (required; AHA Database: 12).
\item[-t] Stop at the time that follows (default: 35 minutes after the beginning for the
  records \code{x2xx} and \code{x3xx}, 3 hours for \code{x0xx} and \code{x1xx}, otherwise
  the end of the input, with a warning).
\item[-U] The physical units (required unless the gain is 0; AHA: \code{mV}).
\end{params}
\begin{lstlisting}[style=text]
$ ad2m -i 1201.CMP -F 250 -G 400 -R 12 -U mV
$ cat 1201.hea
1201 2 250 800
1201.dat 212 400/mV 12 0 0 32104 0
1201.dat 212 400/mV 12 0 90 -31586 0
\end{lstlisting}
\returns{Exit status 0.}
\raises{Exit status 1 for an error in the arguments, a missing \opt{-F}, \opt{-G},
\opt{-R} or \opt{-U}, or an end time not after the start time; 2 if the input cannot be
read; 3 if the output signal file cannot be opened.}
\textit{Differences from the C tool} (\cref{chap:behavior}): the new options \opt{-F},
\opt{-G}, \opt{-R} and \opt{-U}, all required (\opt{-U} unless the gain is 0); the C tool
assumed the parameters of the AHA Database (250\,Hz, 400\,adu/mV, 12 bits, millivolts).
The \opt{-f} and \opt{-t} times are converted at the \opt{-F} frequency. The \opt{-i} file
is read from the input directory (C: from the current directory).

% ------------------------------------------------------------------------------------------
\apientry{md2a}{Convert a record to an AHA-format signal file}
\begin{lstlisting}[style=signature,language={}]
usage: md2a -o AHAFILE -r RECORD [ OPTIONS ... ]
\end{lstlisting}
Writes the signals of a record to an AHA-format signal file: format 16, the signals
interleaved, followed by an end-of-data marker. With \opt{-n}, it also writes a header that
describes the output file as a record.
\begin{params}
\item[-o] The output file, \textit{AHAFILE} (required).
\item[-r] The input record, \textit{RECORD} (required).
\item[-h] Print the usage summary, and exit with status 1.
\item[-n] Write a header for the output file, as the record \textit{NEWREC} that follows.
\end{params}
\begin{lstlisting}[style=text]
$ md2a -r 100s -o 100s.ah -n ahrec
$ cat ahrec.hea
ahrec 2 360 21600
100s.ah 16 200 11 1024 995 21537 0
100s.ah 16 200 11 1024 1011 -3962 0
\end{lstlisting}
\returns{Exit status 0.}
\raises{Exit status 1 for an error in the arguments; 2 if the record cannot be opened; 3 if
the header of \opt{-n} cannot be created; 4 if the output file cannot be opened.}
\textit{Differences from the C tool}: none, apart from \opt{-{}-legacy} (the samples of a
frame beyond those of the record are 0, where the C tool used uninitialized values).

% ------------------------------------------------------------------------------------------
\apientry{ahaecg2mit}{Convert AHA DVD files to WFDB records}
\begin{lstlisting}[style=signature,language={}]
usage: ahaecg2mit [-s] -F FREQ -G GAIN -R BITS -U UNITS RECORD.XXX [RECORD.XXX]...
\end{lstlisting}
Converts the files of the AHA Database DVD, in binary (\file{.ecg}) or text (\file{.txt})
form, to WFDB records: for each input file \textit{RECORD}\file{.XXX}, it writes
\textit{RECORD}\file{.hea}, \textit{RECORD}\file{.dat} (format 212) and
\textit{RECORD}\file{.atr}, and reports \code{wrote ...} on the standard error. Files with
another suffix are ignored, with a message. The options come first, in any order.
\begin{params}
\item[-s] Make short-format (35-minute) records from long-format (2.5-hour) files: the
  first 145 minutes are skipped (located from the \opt{-F} frequency), and the second digit
  of the output record name is incremented by 2.
\item[-F] The sampling frequency in Hz (required; AHA DVD: 250).
\item[-G] The gain in adu per physical unit (required; AHA DVD: 400).
\item[-R] The ADC resolution in bits, from 1 to 32 (required; AHA DVD: 12).
\item[-U] The physical units (required unless the gain is 0; AHA DVD: \code{mV}).
\item[RECORD.XXX] An input file; \textit{RECORD} is a 4-digit AHA Database record name, and
  \textit{XXX} is \code{ecg} or \code{txt}. Several files may be converted in one run. The
  files are read from the input directory (\cref{sec:tools-files}).
\end{params}
\begin{lstlisting}[style=text]
$ ahaecg2mit -F 250 -G 400 -R 12 -U mV 1001.ecg
Rearranging annotations for output annotator atr ...done!
wrote 1001.atr, 1001.dat, and 1001.hea
$ cat 1001.hea
1001 2 250 700
1001.dat 212 400/mV 12 0 0 1213 0 ECG0
1001.dat 212 400/mV 12 0 300 -1026 0 ECG1
\end{lstlisting}
\returns{Exit status 0 (also when an input file is ignored or cannot be opened).}
\raises{Exit status 1 without arguments, with \opt{-h} as the first argument (the usage), or
for a missing or invalid \opt{-F}, \opt{-G}, \opt{-R} or \opt{-U}.}
\textit{Differences from the C tool} (\cref{chap:behavior}): the new options \opt{-F},
\opt{-G}, \opt{-R} and \opt{-U}, required (\opt{-U} unless the gain is 0), which precede the
input files in any order (the C tool had only \opt{-s}, as the first argument, and assumed
250\,Hz, 400\,adu/mV, 12 bits and millivolts); the input files are read from the input
directory (C: from the current directory).

% ------------------------------------------------------------------------------------------
\apientry{makeid}{Write an AHA-format record ID block}
\begin{lstlisting}[style=signature,language={}]
usage: makeid annotator record
\end{lstlisting}
Writes the 512-byte ID block of the AHA tape format for a record (with two signals) and an
annotator to the standard output: the record ID (8 bytes), the number of annotations, the
first and last samples of the annotated segment, the number of 512-byte blocks of samples,
the times of the first and last annotations, and the number of blocks of annotations, as
16- and 32-bit integers. It describes these values on the standard error. The record name
must have at most 8 characters.
\begin{params}
\item[annotator] The annotator.
\item[record] The record.
\end{params}
\begin{lstlisting}[style=text]
$ makeid atr 100s > 100s.id
100s
75 annotations
Annotated segment begins at sample 0
Annotated segment ends at sample 21600
169 blocks of sample data
First annotation is at sample 18
Last annotation is at sample 21423
3 blocks of annotations
\end{lstlisting}
\returns{Exit status 0.}
\raises{Exit status 1 with fewer than two arguments or a record name longer than 8
characters; 2 if the record or the annotator cannot be opened.}
\textit{Differences from the C tool}: none, apart from \opt{-{}-legacy}.

% ------------------------------------------------------------------------------------------
\apientry{readid}{Read an AHA-format record ID block}
\begin{lstlisting}[style=signature,language={}]
usage: readid data-block-size annot-block-size
\end{lstlisting}
Reads an ID block (as written by \cmd{makeid}) from the standard input, and prints the
record ID, the number of data blocks and the number of annotation blocks, recounted in
blocks of the given sizes (in bytes).
\begin{params}
\item[data-block-size] The size of a block of samples, in bytes (at least 1).
\item[annot-block-size] The size of a block of annotations, in bytes (at least 1).
\end{params}
\begin{lstlisting}[style=text]
$ readid 512 16 < 100s.id
100s 169 96
\end{lstlisting}
\returns{Exit status 0.}
\raises{Exit status 1 with fewer than two arguments.}
\textit{Differences from the C tool}: none, apart from \opt{-{}-legacy} (a record ID of
more than 8 characters, which overflows a buffer in the C tool, is printed whole).

% ------------------------------------------------------------------------------------------
\apientry{revise}{Convert an obsolete-format header}
\begin{lstlisting}[style=signature,language={}]
usage: revise record
\end{lstlisting}
Reads the obsolete-format header of a record, \textit{record}\file{.header} (one value per
line, as written by early versions of the WFDB library), and writes a header in the current
format, \textit{record}\file{.hea}.
\begin{params}
\item[record] The record (a \code{header.} prefix is removed).
\end{params}
\returns{Exit status 0.}
\raises{Exit status 1 without an argument; 2 if the old header cannot be read.}
\textit{Differences from the C tool} (\cref{chap:behavior}): a missing ADC resolution stays
unspecified (C: 12 bits).

% ------------------------------------------------------------------------------------------
\apientry{parsescp}{Read and convert an SCP-ECG file}
\begin{lstlisting}[style=signature,language={}]
usage: parsescp -o RECORD [OPTIONS ...] <SCPFILE
\end{lstlisting}
Reads an SCP-ECG file from the standard input, checks its CRCs, decodes its sections
(patient and acquisition data, Huffman tables, leads, QRS locations, reference beats,
rhythm data, measurements, diagnoses) and the ECG, derives the missing limb leads, and
writes: \textit{RECORD}\file{.des}, a description (the patient data, acquisition data and
measurements), \textit{RECORD}\file{.key}, the patient's name and medical record number,
and \textit{RECORD}\file{.ecg}, the 12-lead ECG in a binary layout of its own; with
\opt{-w}, a WFDB record of 12 signals, \textit{RECORD}\file{.hea} and
\textit{RECORD}\file{.dat} (format 16, 200\,adu/mV); with \opt{-v}, the samples as text,
\textit{RECORD}\file{.txt}, and an analysis of the file on the standard output. All these
files are created in the output directory (\cref{sec:tools-files}).
\begin{params}
\item[-o] The record name, \textit{RECORD}, of at most 11 characters (default \code{ecg}).
\item[-a] Anonymize the input and copy it to the standard output (the CRC recomputed);
  no other output.
\item[-f] Parse a file with CRC errors.
\item[-h] Print the usage summary, and exit with status 1.
\item[-l] Low-pass filter (smooth) the output waveforms.
\item[-v] Print a (very) detailed analysis of the input on the standard output, and write
  \textit{RECORD}\file{.txt}; given twice, the text file has column headings.
\item[-w] Write a WFDB-format record.
\item[-z] Make the signals zero-mean.
\item[-b] (Debugging) Show the baseline (residuals), without the reference beats.
\item[-s] (Debugging) Shift the reference beats by the number of samples that follows
  before adding them.
\item[-S] (Debugging) Skip the parsing of the section whose number follows (\code{-} for
  all sections).
\item[-t] (Debugging) Show the reference beats only, without the baseline.
\item[-x] (Debugging) Show hex dumps (implies \opt{-v}).
\end{params}
\begin{lstlisting}[style=text]
$ parsescp -o scp -w < example.scp
$ head -3 scp.hea
scp 12 500 5000     0:22
scp.dat 16 200/mV 16 0 -1 -2500 0 I
scp.dat 16 200/mV 16 0 -3 -2034 0 II
\end{lstlisting}
\begin{phibox}
The \file{.key} file holds the patient's name and medical record number, the \file{.des}
file and the header the age and sex, and \opt{-v} prints all the patient data of the file;
\opt{-a} writes a copy of the input with these fields removed.
\end{phibox}
\returns{Exit status 0.}
\raises{Exit status 1 for an error in the arguments, an input shorter than 6 bytes, a CRC
error (without \opt{-f}), or a reference-beat subtraction zone larger than the template; 2 if the declared length of the
input cannot be allocated; 3 if the input is shorter than its declared length.}
\textit{Differences from the C tool} (\cref{chap:behavior}): the record is written with the
sampling frequency given in the SCP file, and is not written if the file gives none
(\code{parsescp: sampling frequency not specified; record R not written}); the C tool
always wrote 500\,Hz. The \file{.des}, \file{.ecg}, \file{.key} and \file{.txt} files are
created in the output directory, as the record is (C: in the current directory). Signals and
reference beats that cannot be allocated, and samples
beyond the end of a signal shorter than lead 0, are read as absent or as zeros (undefined
behavior in the C tool).

% Storage formats and annotation codes: the byte layouts of the signal formats, the FLAC
% formats, the MIT and AHA annotation file formats, and the tables of annotation codes, as
% implemented by python/src/wfdbx/_signal.py, _flac.py, _annot.py and _ecgcodes.py (and
% identically by the other implementations).
\chapter{Storage formats and annotation codes}
\label{chap:formats}

This chapter is the reference of the binary files of a record: the storage formats of the
signal files, named by the format numbers of the header (\cref{chap:records}), the layout
of the annotation files, and the annotation codes with their mnemonics, descriptions and
mapping tables. These are the formats of the WFDB library \wfdbversion{}; the files that
\wfdbx{} writes are read by the C library, and those of the C library by \wfdbx{}, in all
seven implementations.

\section{Signal storage formats}
\label{sec:formats-signal}

A signal file holds the samples of one \emph{signal group}: the signals of the header that
name the same file. The samples are stored frame by frame: for each frame, the samples of
the first signal of the group (one, or \textit{n} for a signal with \textit{n} samples per
frame, written \code{16x}\textit{n} in the header), then those of the second signal, and so
on. The formats that pack several samples into a group of bytes (212, 310 and 311) pack
consecutive samples of this sequence, whatever signals they belong to. The header can also
give a byte offset (\code{16+44}: the samples begin after 44 bytes of prolog, which the
library skips) and a skew for each signal (\cref{chap:records}).

Each format represents a range of integer values (ADC units). The most negative value of a
format is reserved: it marks an \emph{invalid} sample (a sample that is missing or was not
recorded). \code{putvec} writes the invalid sample \code{WFDB\_INVALID\_SAMPLE} ($-32768$) as
that value, and \code{getvec} reads that value back as \code{WFDB\_INVALID\_SAMPLE} (or, in
the \code{WFDB\_GVPAD} mode of \code{setgvmode}, as the previous valid sample of the signal).
A sample outside the range of the format is written modulo $2^{\mathit{bits}}$ (its lowest
bits), except in format 8, which limits the change from one sample to the next. The formats
are listed by the library constant \code{WFDB\_FMT\_LIST}, in the order of
\cref{tab:formats}; every implementation reads and writes all of them.

\begin{xltabular}{\linewidth}{@{}rrP{5.4cm}Lc@{}}
\caption{The signal storage formats (bits per sample, the bytes of a group of samples, the
range of the values, and whether \wfdbx{} reads (R) and writes (W) them).}
\label{tab:formats}\\
\toprule
{\leavevmode\tabhead Format} & {\leavevmode\tabhead Bits} & {\leavevmode\tabhead Layout} & {\leavevmode\tabhead Range (invalid value)} &
  {\leavevmode\tabhead R/W} \\
\midrule
\endfirsthead
\toprule
{\leavevmode\tabhead Format} & {\leavevmode\tabhead Bits} & {\leavevmode\tabhead Layout} & {\leavevmode\tabhead Range (invalid value)} &
  {\leavevmode\tabhead R/W} \\
\midrule
\endhead
\bottomrule
\endfoot
0 & -- & no signal file (a null signal) & every sample invalid & R/W \\
8 & 8 & 1 byte: the difference from the previous sample, two's complement & differences
  $-128$ to 127 (no invalid value) & R/W \\
16 & 16 & 2 bytes, two's complement, little-endian & $-32768$ to 32767 ($-32768$) & R/W \\
61 & 16 & 2 bytes, two's complement, big-endian & $-32768$ to 32767 ($-32768$) & R/W \\
80 & 8 & 1 byte, offset binary (value + 128) & $-128$ to 127 ($-128$) & R/W \\
160 & 16 & 2 bytes, offset binary (value + 32768), little-endian & $-32768$ to 32767
  ($-32768$) & R/W \\
212 & 12 & 2 samples in 3 bytes & $-2048$ to 2047 ($-2048$) & R/W \\
310 & 10 & 3 samples in two 16-bit little-endian words (15 bits of each) & $-512$ to 511
  ($-512$) & R/W \\
311 & 10 & 3 samples in one 32-bit little-endian word (30 bits) & $-512$ to 511 ($-512$) &
  R/W \\
24 & 24 & 3 bytes, two's complement, little-endian & $-2^{23}$ to $2^{23}-1$ ($-2^{23}$) &
  R/W \\
32 & 32 & 4 bytes, two's complement, little-endian & $-2^{31}$ to $2^{31}-1$ ($-2^{31}$) &
  R/W \\
508 & 8 & FLAC stream, 8 bits per sample & $-128$ to 127 ($-128$) & R/W \\
516 & 16 & FLAC stream, 16 bits per sample & $-32768$ to 32767 ($-32768$) & R/W \\
524 & 24 & FLAC stream, 24 bits per sample & $-2^{23}$ to $2^{23}-1$ ($-2^{23}$) & R/W \\
\end{xltabular}

The ADC resolution given in the header (\cref{chap:records}) may be smaller than the
storage width of the format (11-bit samples in format 212, as in \code{100s}); the format
only has to hold the range of the values. The precise layout of each format follows. In
the bit diagrams, bit 0 is the least significant bit, and $s_0$, $s_1$, $s_2$ are
consecutive samples of the group.

\subsection{Format 0: null signals}
\label{sec:formats-0}

A signal of format 0 has no samples on disk (its file name in the header is usually
\code{\textasciitilde}). Every sample read is invalid (\code{WFDB\_INVALID\_SAMPLE}, or, in
\code{WFDB\_GVPAD} mode, the signal's initial value from the header); samples written to it
are discarded. Format 0 is used by the layout segment and the null segments of
multi-segment records.

\subsection{Format 8: first differences}
\label{sec:formats-8}

Each byte is the difference between a sample and the previous sample of the same signal,
as a signed 8-bit integer; the first difference is taken from the initial value given in
the header. The samples themselves can take any value, but they change by $-128$ to 127 from one
sample to the next: \code{putvec} limits a larger difference to $-128$ or 127 and returns 0 (the slew
rate is too high), and the following samples are then relative to the value
actually written. The format has no invalid value, and an error in the file affects every
following sample.

\subsection{Formats 16 and 61: 16-bit amplitudes}
\label{sec:formats-16}

Each sample is a 16-bit two's complement integer, least significant byte first (format 16)
or most significant byte first (format 61).

\subsection{Formats 80 and 160: offset binary}
\label{sec:formats-80}

Each sample is stored as an unsigned integer, the value plus $2^{\mathit{bits}-1}$: one byte,
the value plus 128 (format 80), or two bytes, the value plus 32768, least significant byte
first (format 160). Format 80 is that of 8-bit WAV files (\cmd{wav2mit}).

\subsection{Format 212: 12-bit amplitudes}
\label{sec:formats-212}

Two samples are packed into three bytes $b_0 b_1 b_2$:
\begin{itemize}
\item $s_0$: the low 8 bits in $b_0$, the high 4 bits in bits 0--3 of $b_1$;
\item $s_1$: the high 4 bits in bits 4--7 of $b_1$, the low 8 bits in $b_2$.
\end{itemize}
Each sample is a 12-bit two's complement integer. When the number of samples written is
odd, the last sample is written as the first two bytes of a group ($b_0 b_1$).

\subsection{Format 310: 10-bit amplitudes in 16-bit words}
\label{sec:formats-310}

Three samples are packed into two 16-bit little-endian words $w_0$ and $w_1$, bit 0 of each
word being unused (0):
\begin{itemize}
\item $s_0$: bits 1--10 of $w_0$;
\item $s_1$: bits 1--10 of $w_1$;
\item $s_2$: its low 5 bits in bits 11--15 of $w_0$, its high 5 bits in bits 11--15 of
  $w_1$.
\end{itemize}
Each sample is a 10-bit two's complement integer. An incomplete group at the end of the
file is written as the words that hold its samples ($w_0$ for one sample, $w_0 w_1$ for
two).

\subsection{Format 311: 10-bit amplitudes in 32-bit words}
\label{sec:formats-311}

Three samples are packed into one 32-bit little-endian word $w$ (four bytes), bits 30 and
31 being unused (0):
\begin{itemize}
\item $s_0$: bits 0--9 of $w$;
\item $s_1$: bits 10--19;
\item $s_2$: bits 20--29.
\end{itemize}
Each sample is a 10-bit two's complement integer. An incomplete group at the end of the
file is written as its first two bytes (one sample) or its four bytes (two samples).

\subsection{Formats 24 and 32: 24- and 32-bit amplitudes}
\label{sec:formats-24}

Each sample is a two's complement integer of 3 bytes (format 24) or 4 bytes (format 32),
least significant byte first. In format 32, the invalid value $-2^{31}$ differs from
\code{WFDB\_INVALID\_SAMPLE}: an invalid sample is written as $-2^{31}$, which is read back
as \code{WFDB\_INVALID\_SAMPLE}.

\section{FLAC formats}
\label{sec:formats-flac}

Formats 508, 516 and 524 store the samples of a signal group losslessly compressed, as a
FLAC stream (Free Lossless Audio Codec) of $\mathit{format} - 500$ bits per sample (8, 16
or 24). Each signal of the group is a channel of the stream, so a FLAC file holds at most 8
signals, and all of them must have the same number of samples per frame (the library
reports \code{cannot store N signals in a single FLAC file} or
\code{every signal in a FLAC file must have the same sampling frequency}). The samples of a
frame of a signal with several samples per frame are consecutive samples of its channel.
The FLAC sample rate is nominal (96\,kHz); the sampling frequency is the one given in the
header.

The C library reads and writes these formats with libFLAC. \wfdbx{} needs no external
library:
\begin{itemize}
\item \textbf{Reading} uses a translation of the stream decoder of libFLAC 1.4.3, in each
  implementation: every subframe type (constant, verbatim, fixed and linear prediction),
  stereo decorrelation, wasted bits, variable block sizes and CRC checks. Damaged files are
  decoded, resynchronized and reported as by the C library, and seeking
  (\code{isigsettime}) works as in C. The MD5 signature of the stream is checked when the
  input is closed if the environment variable \code{WFDB\_FLAC\_CHECK\_MD5} is set
  (\code{isigclose: warning: incorrect MD5 hash in FLAC input}).
\item \textbf{Writing} uses an encoder of \wfdbx{}'s own: the stream begins with
  \code{fLaC}, a \code{STREAMINFO} block (block size 4096, the frame sizes, the number of
  channels and bits per sample, the total number of samples and the MD5 signature, all
  completed when the file is closed) and an empty \code{VORBIS\_COMMENT} block (vendor
  \code{wfdb-x}). The samples are coded in frames of 4096 samples per channel (the last
  frame shorter), the channels independently, each as a constant, verbatim or
  fixed-predictor (order 0 to 4) subframe, whichever is smallest, with Rice-coded residuals
  (partition orders 0 to 6). A sample is first reduced to the width of the format.
\end{itemize}
The files written are standard FLAC streams: libFLAC, and therefore the C library, reads
them back with identical sample values. Their bytes differ from those that libFLAC writes
(which depend on its version and on the processor), so a record written in a FLAC format
by \wfdbx{} has the same samples and checksums as one written by the C library, but not the
same signal file.

\section{Annotation files}
\label{sec:formats-annotations}

An annotation file belongs to a record and an annotator: \code{100s.atr} holds the
annotations of the annotator \code{atr} for the record \code{100s}. Each annotation has a
time (a sample number, or a number of ticks of the file's time resolution), an annotation
code (\code{anntyp}, \cref{sec:formats-codes}), three small integer fields (\code{subtyp},
\code{chan} and \code{num}) and optional auxiliary information (\code{aux}, a string of up
to 255 bytes). The library reads two formats, MIT (the WFDB format, the default for
writing) and AHA; it recognizes the format of a file when it opens it for reading.

\subsection{MIT format}
\label{sec:formats-mit-annotations}

An MIT-format annotation file is a sequence of 16-bit words, each stored least significant
byte first. The 6 high bits of a word are a code $A$ (\code{(w >> 10) \& 63}), and the 10
low bits a value $I$ (\code{w \& 1023}). An annotation is an \emph{annotation word}, whose
code $A$ (1 to 49) is the annotation code and whose value $I$ is the interval, in ticks,
from the time of the previous annotation (from 0 for the first), followed by any number of
\emph{pseudo-annotation words} that modify it, with the codes 59 to 63 of
\cref{tab:mit-pseudo}. The 32-bit integers of the file are stored as two 16-bit words, the
most significant word first, each least significant byte first (the byte order of the
PDP-11: the bytes 2, 3, 0, 1 of a little-endian integer). The word 0 ($A = 0$, $I = 0$)
ends the file.

\begin{xltabular}{\linewidth}{@{}rlL@{}}
\caption{The pseudo-annotation codes of MIT-format annotation files.}
\label{tab:mit-pseudo}\\
\toprule
{\leavevmode\tabhead $A$} & {\leavevmode\tabhead Name} & {\leavevmode\tabhead Meaning} \\
\midrule
\endfirsthead
\toprule
{\leavevmode\tabhead $A$} & {\leavevmode\tabhead Name} & {\leavevmode\tabhead Meaning} \\
\midrule
\endhead
\bottomrule
\endfoot
59 & \code{SKIP} & $I$ is 0; the next two words are a signed 32-bit interval, added to
  the time. It precedes an annotation word whose interval does not fit in 10 bits (more
  than 1023 ticks, or a negative interval, in a file that is out of order), which then has
  the interval 0. Intervals beyond $\pm(2^{31}-1)$ take several \code{SKIP}s. \\
60 & \code{NUM} & $I$ is the \code{num} field of this annotation and of the following
  ones (written only when it changes; initially 0). \\
61 & \code{SUB} & $I$ is the \code{subtyp} field of this annotation only (written only when
  it is not 0). \\
62 & \code{CHN} & $I$ is the \code{chan} field of this annotation and of the following
  ones (written only when it changes; initially 0). \\
63 & \code{AUX} & The low 8 bits of $I$ are a byte count $n$; the next $n$ bytes are the
  \code{aux} string of this annotation, followed by a null byte if $n$ is odd, so that the
  file stays word-aligned. \\
\end{xltabular}

The \code{subtyp} and \code{num} fields are signed 8-bit integers and \code{chan} an
unsigned 8-bit integer, stored in the 10 bits of $I$. Codes 50 to 58 are not used. The
codes 59 to 63 are reserved for the pseudo-annotation words, and cannot be annotation
codes: an annotation word with one of them would be read as a control word, and would
corrupt the rest of the file. The times \code{WFDB\_TIME\_MIN} and \code{WFDB\_TIME\_MAX}
($-2^{63}$ and $2^{63}-1$) cannot be written either. The library function \code{putann}
refuses these times (\code{time overflow}) but does not check the code, as in C (a code
above 63 also loses its high bits); the high-level \code{write\_annotations} accepts only
the codes 0 to 58 and refuses the others, and these times, before it writes anything
(\cref{sec:additions}). An annotation of code 0 (\code{NOTQRS}), which a word of value 0 could not represent, is
written as a \code{SKIP} of its interval minus 1 followed by a word of code 0 and interval
1.

At the beginning of a file, \code{NOTE} annotations at time 0 (with \code{subtyp},
\code{chan} and \code{num} 0) whose \code{aux} begins with \code{\#\#} carry information
for the library, written by \code{putann} before the first annotation and read by
\code{annopen}:
\begin{itemize}
\item \code{\#\# time resolution: }\textit{F}: the times of the file are in ticks of
  \textit{F} per second, not in sample intervals of the record (written when
  \code{setafreq} has set a resolution, as \cmd{wrann} does with the sampling frequency of
  the record; \code{getiaorigfreq} returns it);
\item \code{\#\# annotation type definitions}, then one annotation per code whose mnemonic
  or description has been redefined (by \code{setannstr} or \code{setanndesc} with a
  positive code), with the \code{aux} \textit{code mnemonic description}, then
  \code{\#\# end of definitions}. When it reads such a file, the library takes these
  definitions as the mnemonics and descriptions of \code{annstr}, \code{strann} and
  \code{anndesc};
\item after these, an annotation of code 0 at time 0 ends the table.
\end{itemize}

An MIT-format file written by \code{putann} need not be in canonical order (by time, then
\code{num}, then \code{chan}); \wfdbx{} sorts an out-of-order file itself when it is closed,
without losing annotations (\cref{chap:behavior}), unless the environment variable
\code{WFDBANNSORT} is 0.

\subsection{AHA format}
\label{sec:formats-aha-annotations}

An AHA-format annotation file, as used by the AHA Database and written for it by
\cmd{m2a}, has 16 bytes per annotation:

\begin{xltabular}{\linewidth}{@{}lL@{}}
\caption{An annotation of an AHA-format annotation file.}
\label{tab:aha-annotation}\\
\toprule
{\leavevmode\tabhead Bytes} & {\leavevmode\tabhead Contents} \\
\midrule
\endfirsthead
\toprule
{\leavevmode\tabhead Bytes} & {\leavevmode\tabhead Contents} \\
\midrule
\endhead
\bottomrule
\endfoot
0 & 0 \\
1 & the AHA annotation code, a character (\cref{tab:aha-codes}); mapped to an annotation
  code by \code{ammap} when read, and from one by \code{mamap} when written \\
2--5 & the time, a 32-bit integer in PDP-11 order (most significant word first) \\
6--7 & the serial number of the annotation, a 16-bit little-endian integer, from 1 \\
8 & the \code{subtyp} field \\
9 & the annotation code (\code{anntyp}) when written; read as the \code{chan} field \\
10--15 & the \code{aux} string, up to 6 bytes, padded with null bytes \\
\end{xltabular}

The file is padded with bytes 255 (\code{0377}) to a multiple of 1024 bytes; a byte 255 in
place of the first byte of an annotation ends it. When it reads an AHA file, the library
gives an AHA code \code{U} with a \code{subtyp} of 0 the \code{subtyp} $-1$, and
\code{mamap} maps \code{NOISE} to \code{U} only with that subtype. A file opened for
reading is taken to be in AHA format when its first byte is 0 and its second byte is an AHA
code other than \code{[} and \code{]}; the library warns (\code{file appears to be in AHA
format}) when the format it finds is not the one requested. The times are in the units
written: sample intervals for files written by \cmd{m2a} without \opt{-s}, and milliseconds
in the AHA Database tape files (\cmd{a2m} \opt{-t 1} and \opt{-t 2}). The compressed
annotation files of the AHA Database CD-ROMs (\file{.ANO}), read only by \cmd{a2m}
(\opt{-t 3}), have 6 bytes per annotation: the AHA code, a 24-bit time in milliseconds (most
significant byte first; 255 in its high byte marks the end), and a 16-bit serial number.

\section{Annotation codes}
\label{sec:formats-codes}

Annotation codes are integers from 1 to \code{ACMAX} (49); 0, \code{NOTQRS}, is not an
annotation code (\code{isann} is false for it), and codes 42 to 49 have no standard meaning
(they can be defined by \code{setannstr} and \code{setanndesc}, and by the definitions of
an annotation file). Each code has a symbolic name (a constant of the library: \code{NORMAL},
\code{PVC}, ...), a mnemonic (the text of \cmd{rdann}, returned by \code{annstr} and
\code{ecgstr}) and a description (\code{anndesc}). \Cref{tab:anncodes} gives the standard
values, those of the library when no annotation file has redefined them.

\begin{xltabular}{\linewidth}{@{}rP{3.3cm}lL@{}}
\caption{The annotation codes: number, symbolic name, mnemonic and description.}
\label{tab:anncodes}\\
\toprule
{\leavevmode\tabhead Code} & {\leavevmode\tabhead Name} & {\leavevmode\tabhead Mnemonic} & {\leavevmode\tabhead Description} \\
\midrule
\endfirsthead
\toprule
{\leavevmode\tabhead Code} & {\leavevmode\tabhead Name} & {\leavevmode\tabhead Mnemonic} & {\leavevmode\tabhead Description} \\
\midrule
\endhead
\bottomrule
\endfoot
0 & \code{NOTQRS} & (space) & (not an annotation code) \\
1 & \code{NORMAL} & \code{N} & Normal beat \\
2 & \code{LBBB} & \code{L} & Left bundle branch block beat \\
3 & \code{RBBB} & \code{R} & Right bundle branch block beat \\
4 & \code{ABERR} & \code{a} & Aberrated atrial premature beat \\
5 & \code{PVC} & \code{V} & Premature ventricular contraction \\
6 & \code{FUSION} & \code{F} & Fusion of ventricular and normal beat \\
7 & \code{NPC} & \code{J} & Nodal (junctional) premature beat \\
8 & \code{APC} & \code{A} & Atrial premature beat \\
9 & \code{SVPB} & \code{S} & Supraventricular premature or ectopic beat \\
10 & \code{VESC} & \code{E} & Ventricular escape beat \\
11 & \code{NESC} & \code{j} & Nodal (junctional) escape beat \\
12 & \code{PACE} & \code{/} & Paced beat \\
13 & \code{UNKNOWN} & \code{Q} & Unclassifiable beat \\
14 & \code{NOISE} & \code{\textasciitilde{}} & Change in signal quality \\
15 & -- & \code{[15]} & (none) \\
16 & \code{ARFCT} & \code{\textbar{}} & Isolated QRS-like artifact \\
17 & -- & \code{[17]} & (none) \\
18 & \code{STCH} & \code{s} & ST segment change \\
19 & \code{TCH} & \code{T} & T-wave change \\
20 & \code{SYSTOLE} & \code{*} & Systole \\
21 & \code{DIASTOLE} & \code{D} & Diastole \\
22 & \code{NOTE} & \code{\textquotedbl{}} & Comment annotation \\
23 & \code{MEASURE} & \code{=} & Measurement annotation \\
24 & \code{PWAVE} & \code{p} & P-wave peak \\
25 & \code{BBB} & \code{B} & Bundle branch block beat (unspecified) \\
26 & \code{PACESP} & \code{\textasciicircum{}} & (Non-captured) pacemaker artifact \\
27 & \code{TWAVE} & \code{t} & T-wave peak \\
28 & \code{RHYTHM} & \code{+} & Rhythm change \\
29 & \code{UWAVE} & \code{u} & U-wave peak \\
30 & \code{LEARN} & \code{?} & Beat not classified during learning \\
31 & \code{FLWAV} & \code{!} & Ventricular flutter wave \\
32 & \code{VFON} & \code{[} & Start of ventricular flutter/fibrillation \\
33 & \code{VFOFF} & \code{]} & End of ventricular flutter/fibrillation \\
34 & \code{AESC} & \code{e} & Atrial escape beat \\
35 & \code{SVESC} & \code{n} & Supraventricular escape beat \\
36 & \code{LINK} & \code{@} & Link to external data (aux contains URL) \\
37 & \code{NAPC} & \code{x} & Non-conducted P-wave (blocked APC) \\
38 & \code{PFUS} & \code{f} & Fusion of paced and normal beat \\
39 & \code{WFON}, \code{PQ} & \code{(} & Waveform onset \\
40 & \code{WFOFF}, \code{JPT} & \code{)} & Waveform end \\
41 & \code{RONT} & \code{r} & R-on-T premature ventricular contraction \\
42 & -- & \code{[42]} & (none) \\
43 & -- & \code{[43]} & (none) \\
44 & -- & \code{[44]} & (none) \\
45 & -- & \code{[45]} & (none) \\
46 & -- & \code{[46]} & (none) \\
47 & -- & \code{[47]} & (none) \\
48 & -- & \code{[48]} & (none) \\
49 & -- & \code{[49]} & (none) \\
\end{xltabular}

\subsection{Mapping tables}
\label{sec:formats-maps}

The library classifies the annotation codes with tables that programs can query and change
(\cref{chap:library}): \code{isqrs(a)} is 1 for the codes of QRS complexes (beats);
\code{map1(a)} maps a code to one of \code{NOTQRS}, \code{NORMAL}, \code{PVC},
\code{FUSION} and \code{LEARN}; \code{map2(a)} to one of those or \code{SVPB}; \code{mamap(a, s)} gives the AHA
code (a character) of an annotation code with the subtype \code{s}; and \code{annpos(a)}
the position at which an annotation is drawn (\code{APSTD} at the signal, \code{APHIGH}
above it, \code{APLOW} below, \code{APUNDEF} undefined; the position codes \code{APATT} (4,
attached to the signal given by \code{chan}), \code{APAHIGH} (5) and \code{APALOW} (6, above
and below it) are defined but used by no standard code).
\code{setisqrs}, \code{setmap1}, \code{setmap2} and \code{setannpos} change the tables. A
code out of range gives 0, \code{NOTQRS}, \code{NOTQRS}, \code{O} and \code{APUNDEF}.
\Cref{tab:annmaps} gives the tables for the codes 0 to 41 (codes 42 to 49 give 0,
\code{NOTQRS}, \code{NOTQRS}, \code{O} and \code{APUNDEF}).

\begin{xltabular}{\linewidth}{@{}rlrlllL@{}}
\caption{The mapping tables of the annotation codes 0 to 41 (\code{mamap} of code 14 is
\code{U} with the subtype $-1$, \code{O} otherwise).}
\label{tab:annmaps}\\
\toprule
{\leavevmode\tabhead Code} & {\leavevmode\tabhead Mnemonic} & {\leavevmode\tabhead isqrs} & {\leavevmode\tabhead map1} & {\leavevmode\tabhead map2} & {\leavevmode\tabhead mamap} & {\leavevmode\tabhead annpos} \\
\midrule
\endfirsthead
\toprule
{\leavevmode\tabhead Code} & {\leavevmode\tabhead Mnemonic} & {\leavevmode\tabhead isqrs} & {\leavevmode\tabhead map1} & {\leavevmode\tabhead map2} & {\leavevmode\tabhead mamap} & {\leavevmode\tabhead annpos} \\
\midrule
\endhead
\bottomrule
\endfoot
0 & \code{\ } & 0 & \code{NOTQRS} & \code{NOTQRS} & \code{O} & \code{APUNDEF} \\
1 & \code{N} & 1 & \code{NORMAL} & \code{NORMAL} & \code{N} & \code{APSTD} \\
2 & \code{L} & 1 & \code{NORMAL} & \code{NORMAL} & \code{N} & \code{APSTD} \\
3 & \code{R} & 1 & \code{NORMAL} & \code{NORMAL} & \code{N} & \code{APSTD} \\
4 & \code{a} & 1 & \code{NORMAL} & \code{SVPB} & \code{N} & \code{APSTD} \\
5 & \code{V} & 1 & \code{PVC} & \code{PVC} & \code{V} & \code{APSTD} \\
6 & \code{F} & 1 & \code{FUSION} & \code{FUSION} & \code{F} & \code{APSTD} \\
7 & \code{J} & 1 & \code{NORMAL} & \code{SVPB} & \code{N} & \code{APSTD} \\
8 & \code{A} & 1 & \code{NORMAL} & \code{SVPB} & \code{N} & \code{APSTD} \\
9 & \code{S} & 1 & \code{NORMAL} & \code{SVPB} & \code{N} & \code{APSTD} \\
10 & \code{E} & 1 & \code{PVC} & \code{PVC} & \code{E} & \code{APSTD} \\
11 & \code{j} & 1 & \code{NORMAL} & \code{NORMAL} & \code{N} & \code{APSTD} \\
12 & \code{/} & 1 & \code{NORMAL} & \code{NORMAL} & \code{P} & \code{APSTD} \\
13 & \code{Q} & 1 & \code{NORMAL} & \code{NORMAL} & \code{Q} & \code{APSTD} \\
14 & \code{\textasciitilde{}} & 0 & \code{NOTQRS} & \code{NOTQRS} & \code{U} & \code{APHIGH} \\
15 & \code{[15]} & 0 & \code{NOTQRS} & \code{NOTQRS} & \code{O} & \code{APUNDEF} \\
16 & \code{\textbar{}} & 0 & \code{NOTQRS} & \code{NOTQRS} & \code{O} & \code{APHIGH} \\
17 & \code{[17]} & 0 & \code{NOTQRS} & \code{NOTQRS} & \code{O} & \code{APUNDEF} \\
18 & \code{s} & 0 & \code{NOTQRS} & \code{NOTQRS} & \code{O} & \code{APHIGH} \\
19 & \code{T} & 0 & \code{NOTQRS} & \code{NOTQRS} & \code{O} & \code{APHIGH} \\
20 & \code{*} & 0 & \code{NOTQRS} & \code{NOTQRS} & \code{O} & \code{APHIGH} \\
21 & \code{D} & 0 & \code{NOTQRS} & \code{NOTQRS} & \code{O} & \code{APHIGH} \\
22 & \code{\textquotedbl{}} & 0 & \code{NOTQRS} & \code{NOTQRS} & \code{O} & \code{APHIGH} \\
23 & \code{=} & 0 & \code{NOTQRS} & \code{NOTQRS} & \code{O} & \code{APHIGH} \\
24 & \code{p} & 0 & \code{NOTQRS} & \code{NOTQRS} & \code{O} & \code{APHIGH} \\
25 & \code{B} & 1 & \code{NORMAL} & \code{NORMAL} & \code{N} & \code{APSTD} \\
26 & \code{\textasciicircum{}} & 0 & \code{NOTQRS} & \code{NOTQRS} & \code{O} & \code{APHIGH} \\
27 & \code{t} & 0 & \code{NOTQRS} & \code{NOTQRS} & \code{O} & \code{APHIGH} \\
28 & \code{+} & 0 & \code{NOTQRS} & \code{NOTQRS} & \code{O} & \code{APLOW} \\
29 & \code{u} & 0 & \code{NOTQRS} & \code{NOTQRS} & \code{O} & \code{APHIGH} \\
30 & \code{?} & 1 & \code{LEARN} & \code{LEARN} & \code{Q} & \code{APSTD} \\
31 & \code{!} & 1 & \code{PVC} & \code{PVC} & \code{O} & \code{APSTD} \\
32 & \code{[} & 0 & \code{NOTQRS} & \code{NOTQRS} & \code{[} & \code{APSTD} \\
33 & \code{]} & 0 & \code{NOTQRS} & \code{NOTQRS} & \code{]} & \code{APSTD} \\
34 & \code{e} & 1 & \code{NORMAL} & \code{NORMAL} & \code{N} & \code{APSTD} \\
35 & \code{n} & 1 & \code{NORMAL} & \code{NORMAL} & \code{N} & \code{APSTD} \\
36 & \code{@} & 0 & \code{NOTQRS} & \code{NOTQRS} & \code{O} & \code{APHIGH} \\
37 & \code{x} & 0 & \code{NOTQRS} & \code{NOTQRS} & \code{O} & \code{APHIGH} \\
38 & \code{f} & 1 & \code{NORMAL} & \code{NORMAL} & \code{N} & \code{APSTD} \\
39 & \code{(} & 0 & \code{NOTQRS} & \code{NOTQRS} & \code{O} & \code{APHIGH} \\
40 & \code{)} & 0 & \code{NOTQRS} & \code{NOTQRS} & \code{O} & \code{APHIGH} \\
41 & \code{r} & 1 & \code{PVC} & \code{PVC} & \code{R} & \code{APSTD} \\
\end{xltabular}

\code{ammap(c)}, the reverse of \code{mamap}, maps an AHA code (the character \code{c},
from \code{E} to \code{]}) to an annotation code, and gives \code{NOTQRS} for any other
character:

\begin{xltabular}{\linewidth}{@{}lL@{}}
\caption{The AHA annotation codes and the annotation codes that \code{ammap} gives them.}
\label{tab:aha-codes}\\
\toprule
{\leavevmode\tabhead AHA code} & {\leavevmode\tabhead ammap} \\
\midrule
\endfirsthead
\toprule
{\leavevmode\tabhead AHA code} & {\leavevmode\tabhead ammap} \\
\midrule
\endhead
\bottomrule
\endfoot
\code{E} & \code{VESC} \\
\code{F} & \code{FUSION} \\
\code{G} & \code{NOTQRS} \\
\code{H} & \code{NOTQRS} \\
\code{I} & \code{NOTQRS} \\
\code{J} & \code{NOTQRS} \\
\code{K} & \code{NOTQRS} \\
\code{L} & \code{NOTQRS} \\
\code{M} & \code{NOTQRS} \\
\code{N} & \code{NORMAL} \\
\code{O} & \code{NOTE} \\
\code{P} & \code{PACE} \\
\code{Q} & \code{UNKNOWN} \\
\code{R} & \code{RONT} \\
\code{S} & \code{NOTQRS} \\
\code{T} & \code{NOTQRS} \\
\code{U} & \code{NOISE} \\
\code{V} & \code{PVC} \\
\code{W} & \code{NOTQRS} \\
\code{X} & \code{NOTQRS} \\
\code{Y} & \code{NOTQRS} \\
\code{Z} & \code{NOTQRS} \\
\code{[} & \code{VFON} \\
\code{\textbackslash{}} & \code{NOTQRS} \\
\code{]} & \code{VFOFF} \\
\end{xltabular}


\end{document}
